\documentclass[UTF8,12pt,a4paper,fontset=fandol]{ctexbook}

% 支持从项目根目录或当前笔记目录加载共享样式。
\makeatletter
\def\input@path{{../}{../../}{../../../}{../../../../}}
\makeatother
\usepackage{researchnotes}

% 编译章节列表：{0} 编译全文；{3} 只编译第三章；{2,3} 或 {2,5} 编译多章。
% 正文章号范围为 1--22，多个章号用半角逗号分隔，仍按原书顺序输出。
% 列表中含 0 时编译全文。日常只需修改下面这一行。
\newcommand{\CompileChapter}{0}

% 以下条件由章节列表自动确定，无须手动切换，也无须拆分章节文件。
\makeatletter
\@for\notes@chapter:=\CompileChapter\do{%
    \expandafter\let\csname notes@selected@\number\notes@chapter\endcsname\@empty
}
\newif\ifcompileallchapters
\@ifundefined{notes@selected@0}
    {\compileallchaptersfalse}
    {\compileallchapterstrue}
\newif\ifcompilechapter
% 可选参数是本章在正文或附录内的序号，必选参数是章节列表中的序号。
\newcommand{\selectchapter}[2][0]{%
    \compilechapterfalse
    \ifcompileallchapters
        \compilechaptertrue
    \else
        \@ifundefined{notes@selected@#2}{}{%
            \compilechaptertrue
            \ifnum#1=0\relax
                \setcounter{chapter}{\numexpr#2-1\relax}%
            \else
                \setcounter{chapter}{\numexpr#1-1\relax}%
            \fi
        }%
    \fi
}
% 参数依次为目标章号、编号引用、文字指引；目标章入选时保留编号引用。
\newcommand{\crosschapterref}[3]{%
    \ifcompileallchapters
        #2%
    \else
        \@ifundefined{notes@selected@#1}{#3}{#2}%
    \fi
}
\makeatother



\title{机器学习教材（李航）笔记}
\author{}
\date{}

\begin{document}
\ifcompileallchapters


\maketitle
\tableofcontents

% 本章依据：机器学习/textbooks/机器学习（分章原文）/机器学习-1.pdf，
% PDF 第 1--32 页（原书第 3--34 页）。节与小节沿用教材顺序；
% 推导、条件说明与另拟例题为学习整理，不将教材示意图视为数值实验。


\fi

\selectchapter{1}
\ifcompilechapter
\chapter{统计学习及监督学习概论}
\label{chap:lh-learning-overview}

给定一批邮件及其类别，怎样判断新邮件是否为垃圾邮件；给定若干输入与测量结果，怎样预测新的结果；没有标记时，又能从数据中发现什么结构？这些问题都要求从有限经验中提取可以用于新情形的规律。统计学习关注的核心问题是：用什么模型表示规律，按照什么准则从数据中确定模型，以及怎样检验所得规律能否推广到未见数据。

本章先建立数据、学习任务与模型的共同语言，再以\emph{模型、策略和算法}组织监督学习方法，随后讨论评估、正则化、交叉验证与泛化，最后比较生成方法和判别方法，说明分类、标注和回归三类应用。阅读时应始终区分数据生成的未知规律、学习者选择的模型族、样本上的优化目标与新数据上的预测表现；这四者相互联系，却不是同一个对象。

\section{统计学习}
\label{sec:lh-statistical-learning}

\emph{统计学习}（statistical learning），也称统计机器学习（statistical machine learning），研究计算机如何依据数据构建统计模型，并利用模型进行预测和分析。这里的“学习”体现为系统利用经验改进在给定任务上的表现，而不是单纯保存已经出现的样本。预测侧重对未知结果作出判断，分析侧重揭示数据中的关联或潜在结构；二者都需要明确研究对象、数据表示和评价标准。

数据可以是数值、文本、图像、声音或它们的组合，进入模型前通常需要表示为变量或特征。用随机变量描述观测的不确定性，用概率分布描述其统计规律，是统计学习的基本建模方式。\emph[caution]{数据具有可利用的统计规律是建模前提，不是有限样本自动保证的结论。}样本是否具有代表性、训练与应用环境是否一致、特征是否包含必要信息，都会影响学习结果；本章主要讨论模型学习，不展开数据采集问题。

一个典型学习过程是：取得有限训练数据，规定候选模型的范围，确定选择模型的准则，通过计算求得模型，再将其用于新数据。研究方法时要回答模型怎样构造，研究理论时要回答方法何时有效、需要多少数据和计算，研究应用时则要将任务转化为可检验的学习问题。这一过程同时涉及概率统计、信息论、计算与优化；本章使用的独立同分布假设适用于基础监督学习分析，并不覆盖所有在线、序列或交互数据。

\section{统计学习的分类}
\label{sec:lh-learning-categories}

学习方法可以按反馈信息、模型形式、数据到达方式和所用技术分类。这些划分考察的是不同方面，彼此可以交叉：例如，一个方法可以同时是监督的、概率的、参数化的，并采用批量训练。

\subsection{基本分类}
\label{subsec:lh-basic-categories}

\emph{监督学习}（supervised learning）从带有输入与输出对应关系的\emph{标注数据}（labeled data）中学习预测规则。为明确对象，记原始输入空间为 $\mathcal U$，特征映射为 $\phi:\mathcal U\to\mathcal X$，其中 $\mathcal X\subseteq\mathbb R^d$ 是特征空间（feature space），$\mathcal Y$ 是输出空间（output space）。若不区分原始输入和特征，就直接将 $\mathcal X$ 称为输入空间（input space）。一个实例（instance）的特征向量写为
\[
    x=(x^{(1)},\ldots,x^{(d)})^\top,
\]
其中 $x^{(j)}$ 是第 $j$ 个特征；第 $i$ 个样本的第 $j$ 个特征则记为 $x_i^{(j)}$，不要把样本下标与特征上标混淆。

随机输入、输出分别记为 $X$、$Y$，其观测值为 $x$、$y$。训练集（training set）记为
\begin{equation}
    T=\{(x_i,y_i)\}_{i=1}^N,
    \qquad x_i\in\mathcal X,\quad y_i\in\mathcal Y.
    \label{eq:lh-training-set}
\end{equation}
基本分析假设随机样本对 $(X_i,Y_i)$ 独立同分布（independent and identically distributed，i.i.d.），共同分布记为 $P_{XY}$，且独立的新样本也来自该分布。\emph[caution]{独立同分布指样本对之间独立且分布相同，不要求一个样本内部的 $X$ 与 $Y$ 独立。}恰恰是二者的统计联系，使利用输入预测输出成为可能。$P_{XY}$ 对学习者未知，训练数据只提供关于它的有限信息。

\begin{definition}
\label{def:lh-hypothesis-space}
\emph{假设空间}（hypothesis space）是学习者预先规定的候选模型集合，记为 $\mathcal F$。决策模型通常写成 $f:\mathcal X\to\mathcal Y$；概率预测模型通常写成条件分布 $p_\theta(\cdot\mid x)$。训练算法将训练集 $T$ 映射为选定的模型 $\widehat f$ 或 $p_{\widehat\theta}$；预测时，使用该模型处理新的输入。
\end{definition}

决策函数产生预测值 $\widehat y=\widehat f(x)$；概率分类模型产生各类别的条件概率，在 $0$--$1$ 损失下可按
\[
    \widehat y\in\arg\max_{y\in\mathcal Y}p_{\widehat\theta}(y\mid x)
\]
作出决策，最大值并列时预先规定选择规则。这里的 $\widehat f(X)$ 是对真实输出 $Y$ 的预测，不能把 $Y=\widehat f(X)$ 当成一般数据必然满足的无噪声关系。训练的目标是在有限数据中学习可推广的关系，而不只是复现训练标签。

\emph{无监督学习}（unsupervised learning）使用无标注数据 $U=\{x_i\}_{i=1}^N$，研究数据中的潜在结构或概率规律。聚类为相似实例分组，降维将实例表示为较低维坐标，密度估计近似数据分布。若隐结构记为 $Z$，相应模型可以是映射 $z=g(x)$、条件分布 $p_\theta(z\mid x)$，或联合模型 $p_\theta(x,z)$；在联合模型中，对 $Z$ 求和或积分得到 $p_\theta(x)$。不同目标对应不同评价准则，“没有人工标签”并不意味着没有优化目标，也不意味着所得簇必然对应预先设想的语义类别。

\emph{强化学习}（reinforcement learning）研究智能体在与环境交互时如何选择行动，以改进长期累积奖励。采用有限马尔可夫决策过程（Markov decision process，MDP）建模时，记其为 $(\mathcal S,\mathcal A,P,r,\gamma)$：$\mathcal S$、$\mathcal A$ 分别为有限状态集、动作集，$P(s'\mid s,a)$ 为转移概率，$r(s,a)$ 为期望即时奖励，$\gamma$ 为折扣因子。令 $S_t,A_t,R_{t+1}$ 分别表示时刻 $t$ 的状态、动作和随后收到的奖励，则
\[
    P(s'\mid s,a)=\Pr(S_{t+1}=s'\mid S_t=s,A_t=a),
    \qquad
    r(s,a)=\mathbb E[R_{t+1}\mid S_t=s,A_t=a].
\]
马尔可夫假设要求给定当前状态与动作后，下一状态和奖励的条件规律不再依赖更早历史。$r(s,a)$ 只描述奖励的条件均值，并不表示奖励必然是确定值。

\emph{策略}（policy）$\pi(a\mid s)$ 规定给定状态时选择动作的概率，确定性策略是其特殊情形。为使无限时域回报有定义，本章采用 $0\leq\gamma<1$，并假设存在有限常数 $C$ 使 $|R_t|\leq C$ 几乎处处成立；此时折扣回报绝对有界，状态价值与动作价值分别定义为
\begin{align}
    V^\pi(s)
    &=\mathbb E_\pi\!\left[\sum_{k=0}^\infty\gamma^kR_{t+k+1}
      \,\middle|\,S_t=s\right],
      \label{eq:lh-state-value}\\
    Q^\pi(s,a)
    &=\mathbb E_\pi\!\left[\sum_{k=0}^\infty\gamma^kR_{t+k+1}
      \,\middle|\,S_t=s,A_t=a\right].
      \label{eq:lh-action-value}
\end{align}
第二式先执行给定动作 $a$，以后按 $\pi$ 行动；期望来自环境转移、奖励及随机策略。若用 $\gamma=1$，则还需有限时域或保证回报可积的终止等条件。学习最优策略关注长期价值，因此即时奖励最大的动作未必最好；行动还会改变随后获得的数据，这与静态监督学习有明显区别。

强化学习可按是否显式利用环境模型分为有模型（model-based）和无模型（model-free）方法；前者利用已知或学习到的转移、奖励模型规划。另一条划分是直接优化策略的基于策略（policy-based）方法与估计价值后导出决策的基于价值（value-based）方法。两条划分考察不同性质，不能据此断言所有基于策略或价值的方法都只能是无模型方法。

\emph{半监督学习}（semi-supervised learning）联合利用少量标注数据和较多无标注数据；无标注数据能否改善预测，取决于其结构与目标标签之间是否有可利用的联系。\emph{主动学习}（active learning）则由学习者选择希望获得标签的实例，向标注者查询后更新模型，关注有限标注预算下的效果。前者改变可利用的数据组成，后者改变获得标注的方式，二者也可以结合。

\subsection{按模型分类}
\label{subsec:lh-model-categories}

\emph{概率模型}（probabilistic model）显式描述概率分布，例如类别条件概率 $p_\theta(y\mid x)$；\emph{非概率模型}（non-probabilistic model）或确定性预测模型（deterministic prediction model）直接给出函数值、类别或评分。对分类概率取最大值可以得到决策函数，但决策函数通常丢失概率信息，不能唯一还原原来的分布。若给出类别评分 $s_y(x)$，可另外定义
\[
    q(y\mid x)=\frac{\exp(s_y(x))}{\sum_{c\in\mathcal Y}\exp(s_c(x))},
\]
但这是对有限类别评分所作的新概率建模，不能自动保证所得概率与真实频率吻合。条件随机场和逻辑斯谛回归都可直接建模条件概率，因而\emph[caution]{“概率模型”与“生成模型”并不是同义概念}；后者的区分见第~\ref{sec:lh-generative-discriminative} 节。

概率图模型（probabilistic graphical model）利用图描述分布的分解及条件独立关系。最基本的计算仍是边缘化与乘法规则：对离散变量，$p(x)=\sum_y p(x,y)$、$p(x,y)=p(x)p(y\mid x)$，后式的条件概率在 $p(x)>0$ 处使用；连续情形相应用积分与密度表示。有向图和无向图采用不同的因子分解形式，图本身并不能免除归一化与概率一致性的要求。

\emph{线性模型}（linear model）和\emph{非线性模型}（nonlinear model）的区分需要说明相对于什么变量。预测函数 $f_w(x)=w^\top x+b$ 关于输入是仿射的，机器学习中通常也称线性模型；$f_w(x)=w^\top\phi(x)$ 对固定特征 $\phi(x)$ 是线性的，而对原始输入可以非线性。多项式回归关于输入非线性，却关于系数线性；近邻方法的预测函数一般也不是输入的线性函数。因此，不宜仅凭算法名称判定线性与非线性。

\emph{参数化模型}（parametric model）由维数固定、不随样本量增长的参数向量刻画，例如固定特征数的线性回归或固定结构的神经网络。\emph{非参数化模型}（nonparametric model）不以固定有限维参数族作为整体限制，其表示规模可随数据量或所选复杂度增长，例如通常保存训练实例的近邻方法。非参数化不等于没有参数；实际计算仍会使用有限的数据结构。模型的参数化性质取决于具体设定，例如固定维数的线性支持向量机不能仅因名为支持向量机就一概归为非参数模型。

\subsection{按算法分类}
\label{subsec:lh-algorithm-categories}

\emph{批量学习}（batch learning）先取得一批训练数据，再训练模型并用于预测；\emph{在线学习}（online learning）面向依次到达的数据，在反复预测、获得反馈与更新之间推进。在线监督学习在时刻 $t$ 先用当前模型 $f_t$ 对 $x_t$ 预测，取得 $y_t$ 后计算损失，再更新为 $f_{t+1}$。若模型为 $f_\theta$、损失对参数可微，一个可能的更新是
\[
    \theta_{t+1}=\theta_t-\eta_t
    \left.\nabla_\theta L\bigl(y_t,f_\theta(x_t)\bigr)
      \right|_{\theta=\theta_t},
    \qquad \eta_t>0,
\]
其中 $L$ 度量真实输出与预测的差异，$\eta_t$ 是步长；非光滑目标需要相应的次梯度等方法。此式只是更新示例，其收敛性还需要另外规定目标、数据与步长条件。

在线方式适合需要及时反馈、难以保存全部数据或分布随时间变化的情形，但“每次用一个样本计算梯度”与“在线学习任务”不能完全等同：在固定数据上多次随机抽样训练也可使用随机梯度。批量或在线描述的是数据和更新的组织方式，不能单凭这一点断言哪一种预测必然更准。

\subsection{按技巧分类}
\label{subsec:lh-technique-categories}

\emph{贝叶斯学习}（Bayesian learning）将未知参数表示为随机变量 $\Theta$，以先验分布（prior distribution）表达观测前的信息，以后验分布（posterior distribution）表达观测后的信息。设观测数据为 $D$、参数取值为 $\theta$，用 $\pi(\theta)$ 表示先验密度，$p(D\mid\theta)$ 表示给定参数时的数据密度。若边缘似然 $p(D)$ 有限且为正，则贝叶斯公式为
\begin{equation}
    \pi(\theta\mid D)
    =\frac{p(D\mid\theta)\pi(\theta)}{p(D)},
    \qquad
    p(D)=\int p(D\mid\theta)\pi(\theta)\,\mathrm d\theta.
    \label{eq:lh-bayes-posterior}
\end{equation}
离散参数使用求和。$p(D\mid\theta)$ 对固定 $\theta$ 是关于数据的概率模型，对固定 $D$ 则是关于参数的\emph{似然函数}（likelihood function），不因此成为参数的概率分布。

给出单一参数时，可采用\emph{最大后验估计}（maximum a posteriori estimation，MAP）；\emph{极大似然估计}（maximum likelihood estimation，MLE）则只优化似然：
\[
    \widehat\theta_{\mathrm{MAP}}\in
      \arg\max_{\theta}\pi(\theta\mid D),
    \qquad
    \widehat\theta_{\mathrm{ML}}\in
      \arg\max_{\theta}p(D\mid\theta).
\]
两式均以相应最优解存在为前提。在同一参数化及同一约束区域上，若先验密度为正常数，则 MAP 与 MLE 的解集相同；无界参数空间上未必存在可归一化的均匀先验。完整的贝叶斯预测还要平均参数不确定性：若新观测 $Z$ 在给定 $\Theta$ 后与 $D$ 条件独立，则
\begin{equation}
    p(z\mid D)=\int p(z\mid\theta)\pi(\theta\mid D)\,\mathrm d\theta.
    \label{eq:lh-posterior-predictive}
\end{equation}
这与直接把 $\widehat\theta_{\mathrm{MAP}}$ 代入模型是不同的预测方式。

\emph{核方法}（kernel method）通过内积表示在特征空间中的计算。若映射 $\phi:\mathcal X\to\mathcal H$ 的值属于实内积空间，则
\begin{equation}
    K(x,x')=\langle\phi(x),\phi(x')\rangle_{\mathcal H}
    \label{eq:lh-kernel-inner-product}
\end{equation}
称为相应的核函数（kernel function）。当算法只通过内积使用特征时，可用 $K(x,x')$ 代替显式构造 $\phi(x)$ 后的内积，这就是核技巧（kernel trick）。由内积性质，对任意有限样本和实数 $c_1,\ldots,c_m$，必有
\[
    \sum_{i,j=1}^m c_ic_jK(x_i,x_j)
    =\left\|\sum_{i=1}^m c_i\phi(x_i)\right\|_{\mathcal H}^2\geq0.
\]
因此任意定义的“相似度”不能未经核查就当成内积核使用；对称性与上述正半定条件是必要的，反向构造属于后续核方法的理论内容。特征变换可以增加表达能力，但不自动保证给定数据一定可分。

\begin{example}
\label{ex:lh-quadratic-kernel}
设 $x=(x^{(1)},x^{(2)})^\top\in\mathbb R^2$。构造特征映射，使 $K(x,x')=(x^\top x')^2$ 表示映射后的内积。
\end{example}
\begin{solve}
取 $\phi(x)=((x^{(1)})^2,\sqrt2x^{(1)}x^{(2)},(x^{(2)})^2)^\top$，则
\[
    \phi(x)^\top\phi(x')
    =(x^{(1)}x'^{(1)})^2
      +2x^{(1)}x^{(2)}x'^{(1)}x'^{(2)}
      +(x^{(2)}x'^{(2)})^2
    =(x^\top x')^2.
\]
在该特征上取线性组合就得到关于原始输入的二次函数，而计算核值时不必逐项生成这些二次特征。
\end{solve}

\section{统计学习方法三要素}
\label{sec:lh-three-elements}

“模型、策略和算法”分别回答\emph{学什么、按什么标准选择、怎样计算}。模型族相同而损失或正则项不同，会形成不同的学习问题；优化目标相同而求解算法不同，则可能得到不同精度和计算代价的数值结果。

\subsection{模型}
\label{subsec:lh-model}

模型是对输入输出关系的具体假设，假设空间则规定所有允许的候选。参数化决策模型可写为
\[
    \mathcal F=\{f_\theta:\mathcal X\to\mathcal Y\mid\theta\in\Omega\},
    \qquad \Omega\subseteq\mathbb R^p;
\]
概率模型族相应写为 $\mathcal P=\{p_\theta(\cdot\mid x):\theta\in\Omega\}$。$\Omega$ 称为参数空间（parameter space），$p$ 是参数维数，未必等于输入特征数 $d$；不同参数也可能表示同一个模型。真实数据规律不一定包含在选定模型族中，因此学习可以是在受限范围内寻找较好近似，而不一定是恢复一个确实存在的“真实参数”。

\subsection{策略}
\label{subsec:lh-strategy}

\begin{definition}
\label{def:lh-loss-risk}
\emph{损失函数}（loss function）$L(y,a)\geq0$ 度量真实输出 $y$ 与预测 $a$ 的差异。本笔记沿用教材次序，将真实输出放在第一项。以下假设模型与损失的复合函数可测。对决策函数 $f$，其\emph{期望风险}（expected risk）定义为
\begin{equation}
    R(f)=\mathbb E_{(X,Y)\sim P_{XY}}[L(Y,f(X))]
        =\int L(y,f(x))\,\mathrm dP_{XY}(x,y).
    \label{eq:lh-expected-risk}
\end{equation}
给定训练集 $T$，其\emph{经验风险}（empirical risk）定义为
\begin{equation}
    \widehat R_T(f)=\frac1N\sum_{i=1}^N L(y_i,f(x_i)).
    \label{eq:lh-empirical-risk}
\end{equation}
比较有限风险时假设损失可积。概率预测则将 $f(x)$ 换为所预测的分布，损失也相应定义在概率分布上。
\end{definition}

分类常用 $0$--$1$ 损失 $L(y,a)=\mathbf1_{\{y\ne a\}}$，其中 $\mathbf1_A$ 是事件 $A$ 的指示函数；实值回归常用平方损失 $L(y,a)=(y-a)^2$ 和绝对损失 $L(y,a)=|y-a|$。对离散输出的概率预测 $q(\cdot\mid x)$，\emph{对数损失}（logarithmic loss）为
\begin{equation}
    L(y,q(\cdot\mid x))=-\log q(y\mid x),
    \label{eq:lh-log-loss}
\end{equation}
约定对真实发生的类别预测概率为零时损失为 $+\infty$。本章 $\log$ 表示自然对数。连续输出的负对数密度也可作为优化目标，但密度可以大于 $1$，其负对数不一定非负，不能直接套用离散概率损失的所有性质。

期望风险衡量总体表现，经验风险只衡量样本表现。由于 $P_{XY}$ 未知，通常不能直接计算式~\eqref{eq:lh-expected-risk}。对一个\emph{不依赖训练样本而固定的} $f$，若样本独立同分布且损失可积，大数定律给出 $\widehat R_T(f)\to R(f)$ 几乎必然成立。训练所得 $\widehat f$ 却依赖同一训练集，不能把固定函数的大数定律直接作为其泛化保证；还需控制候选模型集合或学习过程的复杂度。

\begin{definition}
\label{def:lh-erm-srm}
在经验最小值可以达到时，\emph{经验风险最小化}（empirical risk minimization，ERM）选择
\begin{equation}
    \widehat f\in\arg\min_{f\in\mathcal F}\widehat R_T(f).
    \label{eq:lh-erm}
\end{equation}
加入正则项 $J(f)\geq0$ 后，可按惩罚形式选择
\begin{equation}
    \widehat f_\lambda\in\arg\min_{f\in\mathcal F}
      \{\widehat R_T(f)+\lambda J(f)\},\qquad\lambda\geq0.
    \label{eq:lh-penalized-risk}
\end{equation}
其中 $J$ 表达选定的复杂度惩罚或偏好，$\lambda$ 控制惩罚强度。若最小值不能达到，或只使用数值算法，应另外说明近似求解的含义。
\end{definition}

教材用式~\eqref{eq:lh-penalized-risk} 说明\emph{结构风险最小化}（structural risk minimization，SRM）的思想：同时考虑拟合误差和复杂度。更一般的结构风险最小化可在一系列复杂度递增的假设空间中，依据风险上界选择模型；惩罚目标是常用实现形式，不能不加条件地将任意罚项都视为严格的泛化误差上界。正则项也不是总体风险本身，“目标值更小”与“实际泛化更好”仍需区别。

对条件概率模型，使用式~\eqref{eq:lh-log-loss} 时，经验风险最小化等价于极大化条件似然，推导见例~\ref{ex:lh-erm-mle}。若再加入先验密度 $\pi(\theta)$，最大后验估计等价于最小化
\begin{equation}
    -\frac1N\sum_{i=1}^N\log p_\theta(y_i\mid x_i)
    -\frac1N\log\pi(\theta).
    \label{eq:lh-map-penalty}
\end{equation}
这给出正则化的一种概率解释；加上与参数无关的常数不改变最优解。\emph[caution]{平均损失与总损失的尺度不同，先验项或正则项的系数必须同步换算。}

\subsection{算法}
\label{subsec:lh-algorithm}

算法（algorithm）是求解模型选择目标的具体计算过程。若存在适用的解析解，可以直接计算；更常见的是用迭代算法求近似解。描述算法时需要交代输入数据和超参数、初始化、更新规则、终止条件以及输出对象。停止可依据目标变化、梯度或次梯度条件、误差容限及计算预算，但达到停止条件并不自动意味着到达全局最优。

目标的凸性、可微性、约束和数值条件会影响可用算法及其保证。即使每步目标值不增，也不必然推出参数收敛；参数收敛也不必然意味着得到全局最优；优化得很精确仍不能代替泛化分析。模型参数由训练估计，而正则化强度、模型阶数等通常作为超参数，在训练过程之外利用验证信息选择。

\section{模型评估与模型选择}
\label{sec:lh-model-evaluation}

\subsection{训练误差与测试误差}
\label{subsec:lh-training-test-error}

设 $\widehat f$ 已由训练过程确定，另有未参与训练和模型选择的测试集 $T_{\mathrm{test}}=\{(x'_j,y'_j)\}_{j=1}^{N_{\mathrm{test}}}$。在给定评价损失下，训练误差是 $\widehat R_T(\widehat f)$，测试误差为
\begin{equation}
    \widehat R_{\mathrm{test}}(\widehat f)
      =\frac1{N_{\mathrm{test}}}\sum_{j=1}^{N_{\mathrm{test}}}
        L(y'_j,\widehat f(x'_j)).
    \label{eq:lh-test-risk}
\end{equation}
当测试样本独立于整个训练与选择过程、来自同一总体且损失可积时，给定所得模型的测试误差期望等于其总体风险。训练误差一般没有这样的无偏性，因为模型就是根据训练数据选出的。一个有限测试集仍有抽样误差，不能把一次比较当作对方法优劣的绝对证明。

若评价采用 $0$--$1$ 损失，测试误差就是错误率（error rate），准确率（accuracy）为
\[
    \mathrm{Acc}=\frac1{N_{\mathrm{test}}}
      \sum_{j=1}^{N_{\mathrm{test}}}\mathbf1_{\{y'_j=\widehat f(x'_j)\}}
      =1-\widehat R_{\mathrm{test}}(\widehat f).
\]
训练损失与评价指标可以不同，例如分类模型按对数损失训练、按准确率评价；这种选择常出于可优化性与任务需求，但必须明确两者各自测量什么。

\subsection{过拟合与模型选择}
\label{subsec:lh-overfitting}

\emph{模型选择}（model selection）是在不同模型族、结构或超参数之间作出选择。\emph{过拟合}（overfitting）指学习过度适应有限训练数据中的偶然性或噪声，表现为训练拟合很好，而对新数据的预测较差；\emph{欠拟合}（underfitting）则指模型或学习过程不足以刻画主要规律。过拟合并不要求模型参数数目超过某个已知“真实参数数目”，其判断依赖目标分布下的预测表现。

对嵌套空间 $\mathcal F_0\subseteq\mathcal F_1\subseteq\cdots$，在同一数据和损失下精确最小化经验风险，有
\[
    \inf_{f\in\mathcal F_{m+1}}\widehat R_T(f)
       \leq\inf_{f\in\mathcal F_m}\widehat R_T(f),
\]
因为扩大集合不会提高下确界。这只说明最优训练误差不增，不说明测试误差也不增。测试误差随复杂度先降后升是常见示意，不能作为所有数据、模型与训练方式必须满足的定律。

\begin{example}
\label{ex:lh-polynomial-fitting}
教材例 1.1 用多项式拟合说明模型选择。给定实数样本 $(x_i,y_i)$，在次数不超过 $M$ 的多项式
\[
    f_M(x;w)=\sum_{j=0}^M w_jx^j,
    \qquad w=(w_0,\ldots,w_M)^\top\in\mathbb R^{M+1}
\]
中最小化平方误差。说明如何求系数、解何时唯一，以及增加次数为何可能导致过拟合。
\end{example}
\begin{solve}
记 $y=(y_1,\ldots,y_N)^\top$，设计矩阵 $V\in\mathbb R^{N\times(M+1)}$ 的元素为 $V_{i,j+1}=x_i^j$，其中 $j=0,\ldots,M$。目标可写成
\begin{equation}
    E_M(w)=\frac1{2N}\sum_{i=1}^N(f_M(x_i;w)-y_i)^2
      =\frac1{2N}\|Vw-y\|_2^2.
    \label{eq:lh-polynomial-least-squares}
\end{equation}
令梯度为零得到正规方程（normal equations）
\begin{equation}
    V^\top V\widehat w=V^\top y.
    \label{eq:lh-normal-equations}
\end{equation}
该目标是凸二次函数，所以正规方程的解都是全局最小点。最小点存在：在有限维空间中，将 $y$ 正交投影到 $V$ 的列空间，即得到某个 $V\widehat w$；残差与列空间正交，正好满足正规方程。若 $V$ 满列秩，则 $V^\top V$ 正定，唯一解为
\[
    \widehat w=(V^\top V)^{-1}V^\top y.
\]
若 $V$ 不满列秩，则存在非零 $h$ 使 $Vh=0$，$\widehat w+h$ 与 $\widehat w$ 有相同拟合值，因而系数不唯一；训练点上的最优拟合向量仍由正交投影唯一确定。

当 $x_i$ 两两不同时，若 $N\geq M+1$，$Vh=0$ 意味着一个次数不超过 $M$ 的多项式有至少 $M+1$ 个不同零点，因此该多项式恒为零，$V$ 满列秩。特别地，$M=N-1$ 时可唯一插值任意给定的 $y_i$，训练误差为零；$M>N-1$ 时则可插值但系数不唯一。若观测中含噪声，这样的插值也会拟合噪声，不能仅由零训练误差推断新点预测准确。教材比较 $10$ 个样本下 $M=0,1,3,9$ 的曲线，其作用是展示欠拟合与过拟合的可能性，不是在任何数据上都推荐三次多项式。实际选择 $M$ 应使用验证数据或其他有依据的复杂度控制。
\end{solve}

\section{正则化与交叉验证}
\label{sec:lh-regularization-cv}

\subsection{正则化}
\label{subsec:lh-regularization}

\emph{正则化}（regularization）在拟合目标之外加入限制或偏好，以控制模型对有限数据的适应程度。沿用式~\eqref{eq:lh-polynomial-least-squares} 的模型，对全部系数施加 $L_2$ 或 $L_1$ 惩罚可分别写为
\begin{align}
    E_{2,\lambda}(w)
      &=\frac1{2N}\|Vw-y\|_2^2+\frac\lambda2\|w\|_2^2,
      \label{eq:lh-ridge-objective}\\
    E_{1,\lambda}(w)
      &=\frac1{2N}\|Vw-y\|_2^2+\lambda\|w\|_1,
      \label{eq:lh-lasso-objective}
\end{align}
其中 $\|w\|_2^2=\sum_{j=0}^M w_j^2$、$\|w\|_1=\sum_{j=0}^M|w_j|$，$\lambda\geq0$。$L_2$ 惩罚连续收缩系数，$L_1$ 惩罚在零点具有非光滑性，常产生部分系数恰为零的稀疏解，但不保证每个问题都出现零系数。这里惩罚了截距 $w_0$；实际建模常将截距排除，届时应相应改变惩罚矩阵和解的条件。

对式~\eqref{eq:lh-ridge-objective}，当 $\lambda>0$ 时有
\[
    \nabla E_{2,\lambda}(w)=\frac1N V^\top(Vw-y)+\lambda w,
    \qquad
    \widehat w_\lambda=(V^\top V+N\lambda I_{M+1})^{-1}V^\top y.
\]
其中 $I_{M+1}$ 是单位矩阵；任意非零 $u$ 都满足 $u^\top(V^\top V+N\lambda I_{M+1})u=\|Vu\|_2^2+N\lambda\|u\|_2^2>0$，所以逆存在且解唯一。这说明正则化既改变统计偏好，也可能改善求解的确定性；它不是在原最小二乘解之后任意缩小几个系数。

从式~\eqref{eq:lh-map-penalty} 看，负对数先验提供惩罚：独立高斯先验对应平方 $L_2$ 项，独立拉普拉斯先验对应 $L_1$ 项，比例系数须结合似然尺度计算。奥卡姆剃刀（Occam's razor）在此体现为，在解释数据的能力相近时偏向较简单的模型；这是一种选择原则，而非保证简单模型总是正确的定理。特征单位会影响系数惩罚，因此缩放与 $\lambda$ 的选择需要一并考虑；具体参数选择仍依赖验证或额外理论。

\subsection{交叉验证}
\label{subsec:lh-cross-validation}

\emph{交叉验证}（cross-validation）通过反复划分数据，比较候选训练方案在未参与该次拟合的数据上的表现。首先区分三种用途：训练集拟合参数，验证集（validation set）选择模型或超参数，测试集（test set）用于最终评价。教材“简单交叉验证”中用于选择模型的留出部分，在本笔记中称为验证集，以免与最终测试集混淆；反复依据某组数据修改模型后，这组数据就已经参与了选择过程。

当数据足够时，可先划分训练、验证、测试三部分：在训练集上拟合各候选，比较验证误差，确定方案后可在训练与验证数据的合并集上重新训练，最后在保留的测试集上评估。若数据较少，可先保留测试集，再在开发数据上做 $S$ 折交叉验证（$S$-fold cross-validation），其中 $2\leq S\leq N_{\mathrm{dev}}$，$N_{\mathrm{dev}}$ 是开发样本数。将开发样本索引划分为非空、互不相交且大小尽量接近的 $I_1,\ldots,I_S$；对候选配置 $h$，第 $s$ 次仅在 $I_s$ 之外训练，所得模型记为 $\widehat f_{h,-s}$。按样本加权的验证误差为
\begin{equation}
    \widehat R_{\mathrm{CV}}(h)
       =\frac1{N_{\mathrm{dev}}}\sum_{s=1}^S\sum_{i\in I_s}
          L(y_i,\widehat f_{h,-s}(x_i)),
    \qquad
    \widehat h\in\arg\min_{h\in\mathcal H}\widehat R_{\mathrm{CV}}(h),
    \label{eq:lh-cross-validation}
\end{equation}
其中 $\mathcal H$ 是预先规定的有限候选集合。等大折时，这也等于各折平均损失的算术平均；折大小不等时应说明采用的加权方式。选出的是\emph{训练配置}，随后应按 $\widehat h$ 在全部开发数据上重新拟合，而不是直接把某一折的模型当作最终模型。

\emph{留一交叉验证}（leave-one-out cross-validation）取 $S=N_{\mathrm{dev}}$，每次仅留出一个样本，其余样本用于训练。它能让每次训练使用较多数据，但通常需要更多次拟合，而且不同折训练集高度重叠，不能把折结果当成相互独立的观测。留出法、$S$ 折法与留一法的选择取决于数据规模、拟合代价和评价目的，不存在单凭折数就保证更准确的规则。

\begin{notealert}[type=caution,title={验证对象应包含整个训练流程}]
标准化、特征筛选等凡需从数据估计的步骤，都应只在每折训练部分拟合，再用于该折验证部分。先用全部数据筛选特征再交叉验证，会将验证信息泄露到训练过程。随机划分适用于样本近似独立且使用目标相符的情形；同一主体的重复记录应考虑按主体分组，时间预测应保持先训练、后验证的时间关系。交叉验证分数参与了模型选择，不能代替独立的最终测试；需要评估完整调参流程而又没有独立测试集时，可采用外层评估、内层调参的嵌套划分。
\end{notealert}

\section{泛化能力}
\label{sec:lh-generalization}

\subsection{泛化误差}
\label{subsec:lh-generalization-error}

\emph{泛化能力}（generalization ability）是模型对未见数据进行预测的能力。给定由训练集 $T$ 学到的模型 $\widehat f_T$，其\emph{泛化误差}（generalization error）就是独立新样本上的期望风险
\begin{equation}
    R(\widehat f_T)
      =\mathbb E_{(X,Y)\sim P_{XY}}
          [L(Y,\widehat f_T(X))\mid T],
    \label{eq:lh-generalization-error}
\end{equation}
这里先将训练结果固定，期望只针对新样本；若算法自身随机，还应固定其随机结果或另行对其平均。由于训练集随机，$R(\widehat f_T)$ 也是随机量；再对训练样本取期望，才得到训练过程的平均泛化误差。相应的泛化间隙为 $R(\widehat f_T)-\widehat R_T(\widehat f_T)$，它与泛化误差本身不同。

测试误差是总体风险的样本估计，泛化界则给出某个抽样概率下的理论保证。一个上界较松，不表示实际误差一定大；两个方法的上界大小，也未必能决定其真实风险大小。若训练与应用分布改变，本章同分布条件下的结论需要重新核查。

\subsection{泛化误差上界}
\label{subsec:lh-generalization-bound}

先说明本节使用的概率工具。\emph{Hoeffding 不等式}指出：若独立随机变量 $Z_i$ 满足 $a_i\leq Z_i\leq b_i$ 几乎处处成立，且 $\sum_i(b_i-a_i)^2>0$，记 $\overline Z=N^{-1}\sum_{i=1}^N Z_i$，则对任意 $t>0$，
\begin{align}
    \Pr(\overline Z-\mathbb E\overline Z\geq t)
       &\leq\exp\!\left(-\frac{2N^2t^2}{\sum_{i=1}^N(b_i-a_i)^2}\right),
       \label{eq:lh-hoeffding-upper}\\
    \Pr(\mathbb E\overline Z-\overline Z\geq t)
       &\leq\exp\!\left(-\frac{2N^2t^2}{\sum_{i=1}^N(b_i-a_i)^2}\right).
       \label{eq:lh-hoeffding-lower}
\end{align}
这里按教材引用该不等式，不证明其本身；以下完整给出由它推出的有限假设空间泛化界。注意变量只需独立而不必同分布；在学习问题中使用独立同分布，另为保证各项损失具有相同的总体均值。

\begin{theorem}[有限假设空间的泛化界]
\label{thm:lh-finite-generalization}
设二分类输出空间为 $\{-1,+1\}$，训练样本对 $(X_i,Y_i)$ 独立同分布于 $P_{XY}$。设 $\mathcal F=\{f_1,\ldots,f_m\}$ 是\emph{不依赖本次训练样本而预先固定的} 非空有限分类函数集合，各 $f_j:\mathcal X\to\{-1,+1\}$ 可测，采用 $0$--$1$ 损失。对任意 $0<\delta<1$，以至少 $1-\delta$ 的概率，对所有 $f\in\mathcal F$ 同时成立
\begin{equation}
    R(f)\leq\widehat R_T(f)+\varepsilon_{N,m,\delta},
    \qquad
    \varepsilon_{N,m,\delta}
       =\sqrt{\frac{\log m+\log(1/\delta)}{2N}}.
    \label{eq:lh-finite-generalization-bound}
\end{equation}
概率针对训练样本的联合抽样而言。
\end{theorem}

\begin{proof}
对一个固定的 $f$，令 $Z_i=L(Y_i,f(X_i))$。这些变量独立、取值于 $[0,1]$，且 $\mathbb E Z_i=R(f)$、$N^{-1}\sum_iZ_i=\widehat R_T(f)$。由式~\eqref{eq:lh-hoeffding-lower}，对任意 $\varepsilon>0$，
\[
    \Pr\bigl(R(f)-\widehat R_T(f)\geq\varepsilon\bigr)
       \leq e^{-2N\varepsilon^2}.
\]
对所有候选使用并集上界，不要求不同候选的损失相互独立，得到
\begin{align*}
    \Pr\!\left(\exists f\in\mathcal F:
          R(f)-\widehat R_T(f)\geq\varepsilon\right)
    &\leq\sum_{f\in\mathcal F}
       \Pr\bigl(R(f)-\widehat R_T(f)\geq\varepsilon\bigr)\\
    &\leq m e^{-2N\varepsilon^2}.
\end{align*}
令右端等于 $\delta$，解得式~\eqref{eq:lh-finite-generalization-bound} 中的 $\varepsilon_{N,m,\delta}$。在补事件上，全部候选同时满足严格小于该偏差，因而也满足所述非严格不等式。
\end{proof}

该结论的关键是“对所有候选同时成立”，因此同一事件上也可将依赖训练数据选出的经验最小化模型 $\widehat f_T$ 代入。若只对一个事先固定的 $f$ 建立结论，就不能不加论证地替换成训练后挑出的模型。有限集合上的 $\log m$ 项反映了在多个候选中进行选择的代价；无限集合不能把 $m$ 直接写成无穷后仍得到有用结论，需要其他容量工具。

\begin{corollary}[经验风险最小化的风险比较]
\label{cor:lh-erm-risk-comparison}
在定理~\ref{thm:lh-finite-generalization} 的条件下，记
\[
    \widetilde\varepsilon
       =\sqrt{\frac{\log(2m/\delta)}{2N}},
    \quad
    \widehat f_T\in\arg\min_{f\in\mathcal F}\widehat R_T(f),
    \quad
    f_{\mathcal F}^*\in\arg\min_{f\in\mathcal F}R(f).
\]
则以至少 $1-\delta$ 的概率，
\begin{equation}
    R(\widehat f_T)\leq R(f_{\mathcal F}^*)+2\widetilde\varepsilon.
    \label{eq:lh-erm-oracle-bound}
\end{equation}
这是对教材有限空间界的补充推论，比较的是学习结果与同一假设空间内的最优风险。
\end{corollary}

\begin{proof}
对每个 $f$ 的两个偏差方向分别使用 Hoeffding 不等式，再对 $m$ 个函数取并集上界，得
\[
    \Pr\!\left(\exists f\in\mathcal F:
       |R(f)-\widehat R_T(f)|\geq\widetilde\varepsilon\right)
       \leq 2m e^{-2N\widetilde\varepsilon^2}=\delta.
\]
在其补事件上，由经验最小化性质，
\[
    R(\widehat f_T)
      \leq\widehat R_T(\widehat f_T)+\widetilde\varepsilon
      \leq\widehat R_T(f_{\mathcal F}^*)+\widetilde\varepsilon
      \leq R(f_{\mathcal F}^*)+2\widetilde\varepsilon.
      \qedhere
\]
\end{proof}

固定 $m$ 和 $\delta$ 时，随样本量增加而趋于零的是上述偏差项，而不是必然整个泛化误差；噪声和模型限制仍可能使 $R(f_{\mathcal F}^*)>0$。若希望定理中的单侧偏差项不超过 $\eta>0$，充分条件是 $N\geq\log(m/\delta)/(2\eta^2)$。模型规模也随 $N$ 改变时，应同时检查 $\log m$ 的增长，不能只考察分母的样本量。

\section{生成模型与判别模型}
\label{sec:lh-generative-discriminative}

\begin{definition}
\label{def:lh-generative-discriminative}
\emph{生成方法}（generative approach）学习输入与输出的联合分布 $p_\theta(x,y)$，再据此推导预测所需的条件分布或决策，其模型称为\emph{生成模型}（generative model）。\emph{判别方法}（discriminative approach）直接学习条件分布 $p_\theta(y\mid x)$、决策函数或用于决策的评分，其模型称为\emph{判别模型}（discriminative model）。
\end{definition}

对有限类别，在边缘概率或密度为正的输入处，生成模型给出
\begin{equation}
    p_\theta(y\mid x)
       =\frac{p_\theta(x,y)}{\sum_{c\in\mathcal Y}p_\theta(x,c)}
       =\frac{p_\theta(x\mid y)p_\theta(y)}
          {\sum_{c\in\mathcal Y}p_\theta(x\mid c)p_\theta(c)}.
    \label{eq:lh-generative-posterior}
\end{equation}
这相当于同时描述类别比例与各类别中的输入分布。朴素贝叶斯和隐马尔可夫模型是典型例子；逻辑斯谛回归、条件随机场直接建模条件分布，感知机、支持向量机等直接学习决策或评分，属于判别方法。是否输出概率与是否建模联合分布是两个问题，概率输出本身不足以判断生成或判别。

从联合分布可以求出条件分布，但单独给定 $p(y\mid x)$ 一般不能唯一恢复 $p(x,y)$，因为还缺少 $p(x)$。生成建模适合需要描述数据分布、采样或对隐变量求和的任务；判别建模将重点放在给定输入的预测上，允许直接采用针对任务的特征。生成假设正确与否、模型复杂度及样本量都会影响两类方法的表现，不能普遍断言生成方法总收敛更快或判别方法总更准确。含隐变量也不是生成方法独有的性质，判别模型同样可以引入隐变量。

\begin{example}
\label{ex:lh-generative-binary}
设类别 $Y\in\{0,1\}$，先验概率为 $\Pr(Y=1)=0.4$；二元特征 $X\in\{0,1\}$ 满足 $\Pr(X=1\mid Y=1)=0.8$、$\Pr(X=1\mid Y=0)=0.2$。由该生成模型判断观察到 $X=1$ 时的类别。
\end{example}
\begin{solve}
联合概率分别为 $\Pr(X=1,Y=1)=0.8\times0.4=0.32$ 与 $\Pr(X=1,Y=0)=0.2\times0.6=0.12$，所以
\[
    \Pr(Y=1\mid X=1)=\frac{0.32}{0.32+0.12}=\frac8{11}.
\]
在两类误判代价相同的 $0$--$1$ 损失下，应预测为类别 $1$；若误判代价改变，最优决策阈值也会改变，不能仅凭概率较大就断言在任意损失下都应作相同决策。
\end{solve}

\section{监督学习应用}
\label{sec:lh-supervised-applications}

\subsection{分类问题}
\label{subsec:lh-classification}

\emph{分类}（classification）的输出空间是有限类别集合 $\mathcal Y=\{1,\ldots,C\}$，其中 $C\geq2$；输入既可以是连续变量，也可以是离散变量或其组合。训练所得预测规则称为\emph{分类器}（classifier）。例如文本分类可用词是否出现、词频等表示一篇文本，将其分入预先给定的主题或邮件类别；类别标签来自训练数据，区别于无监督聚类自行形成分组。

二分类评价需先指定关注的正类（positive class），其余为负类（negative class）。记真正例（true positive）数为 $\mathrm{TP}$、假负例（false negative）数为 $\mathrm{FN}$、假正例（false positive）数为 $\mathrm{FP}$、真负例（true negative）数为 $\mathrm{TN}$。它们分别对应真实为正且预测为正、真实为正却预测为负、真实为负却预测为正、真实为负且预测为负。样本总数为四者之和，常用指标为
\begin{align}
    \mathrm{Acc}
       &=\frac{\mathrm{TP}+\mathrm{TN}}
          {\mathrm{TP}+\mathrm{FN}+\mathrm{FP}+\mathrm{TN}},
          \label{eq:lh-accuracy}\\
    \mathrm{Prec}
       &=\frac{\mathrm{TP}}{\mathrm{TP}+\mathrm{FP}},
    \qquad
    \mathrm{Rec}=\frac{\mathrm{TP}}{\mathrm{TP}+\mathrm{FN}},
          \label{eq:lh-precision-recall}\\
    F_1
       &=\frac{2\,\mathrm{Prec}\,\mathrm{Rec}}{\mathrm{Prec}+\mathrm{Rec}}
        =\frac{2\mathrm{TP}}{2\mathrm{TP}+\mathrm{FP}+\mathrm{FN}}.
          \label{eq:lh-f1}
\end{align}
\emph{精确率}（precision，也称查准率）考察预测正例中有多少是真的，\emph{召回率}（recall，也称查全率）考察真实正例中有多少被找到，$F_1$ 是二者的调和平均。各分式只在分母非零时定义；特别地，两率都为零时首个 $F_1$ 分式需作约定，若后一个计数分式分母为正，可据此赋值为零。评价前应说明如何处理不存在预测正例或真实正例的情形。

\begin{example}
\label{ex:lh-imbalanced-classification}
某测试集有 $100$ 个正例、$900$ 个负例。分类器给出 $\mathrm{TP}=60$、$\mathrm{FN}=40$、$\mathrm{FP}=20$、$\mathrm{TN}=880$。计算指标，并与全部预测为负类的规则比较。
\end{example}
\begin{solve}
该分类器的准确率为 $940/1000=0.94$，精确率为 $60/80=0.75$，召回率为 $60/100=0.60$，$F_1=120/(120+20+40)=2/3$。全部预测为负的规则也有 $0.90$ 的准确率，却没有找到任何正例，召回率为零，精确率因没有预测正例而未定义。这说明类别不平衡时只看准确率可能掩盖任务关心的失败；仍应结合正类含义和误判代价评价。
\end{solve}

\subsection{标注问题}
\label{subsec:lh-tagging}

\emph{标注}（tagging）是为输入观测序列预测相应标记序列的监督学习任务，也是\emph{结构化预测}（structured prediction）的一种形式。为区别样本索引与序列位置，第 $i$ 个样本写成
\[
    x_i=(x_{i,1},\ldots,x_{i,n_i}),
    \qquad
    y_i=(y_{i,1},\ldots,y_{i,n_i}),\qquad y_{i,t}\in\mathcal C,
\]
其中 $n_i$ 是该序列长度，$\mathcal C$ 是有限标记集；不同样本的长度可以不同，与训练序列总数 $N$ 之间不存在固定大小关系。词性标注和命名实体识别都是典型任务，每个位置的预测往往依赖周围观测与相邻标记，不能一概当作相互独立的分类。

若对长度为 $n$ 的序列建模为 $p_\theta(y_{1:n}\mid x_{1:n})$，在整条序列的 $0$--$1$ 损失下，预测可写为
\begin{equation}
    \widehat y_{1:n}\in\arg\max_{y_{1:n}\in\mathcal C^n}
        p_{\widehat\theta}(y_{1:n}\mid x_{1:n}).
    \label{eq:lh-sequence-decoding}
\end{equation}
若每个位置有 $K=|\mathcal C|$ 个候选标记，无约束序列共有 $K^n$ 种，因此常需利用隐马尔可夫模型或条件随机场等模型的分解结构进行有效推断。逐位置取边缘概率最大的标签与取整条序列联合概率最大的标签不一定相同，它们对应不同决策目标。

例如用 BIO 规则表示实体片段，$B$ 表示一个片段的开始，$I$ 表示其后续位置，$O$ 表示片段之外；“北京大学位于北京”若按“北京大学、位于、北京”三个词输入，可标为 $(B,O,B)$，表示两个单词级实体片段。标注评价可以计算词元准确率，也可以按完整实体片段计算精确率、召回率与 $F_1$；必须说明评价单位、边界和类型是否都要匹配。词元标对较多并不保证实体片段识别同样准确。

\subsection{回归问题}
\label{subsec:lh-regression}

\emph{回归}（regression）预测实值或实向量输出，本章主要取 $Y\in\mathbb R$。特征既可以包含连续量，也可以包含编码后的离散量，不要求输入输出都必须是连续随机变量。输入特征数不同可区分一元与多元回归；预测函数的形式不同可区分线性与非线性回归。拟合的是统计关系，单凭回归关系不能推断改变某个输入就必然导致输出发生相应变化。

在平方损失下，预测目标与条件均值有直接联系。假设 $\mathbb E[Y^2]<\infty$，记 $\mu(X)=\mathbb E[Y\mid X]$；对二阶可积的预测 $f(X)$，将 $Y-f(X)=(Y-\mu(X))+(\mu(X)-f(X))$ 平方并取条件期望，由 $\mathbb E[Y-\mu(X)\mid X]=0$ 得
\begin{equation}
    \mathbb E[(Y-f(X))^2\mid X]
       =\operatorname{Var}(Y\mid X)+(\mu(X)-f(X))^2.
    \label{eq:lh-regression-conditional-risk}
\end{equation}
再取期望可知，在所有二阶可积预测函数中，条件均值最小化平方风险；若只能在受限假设空间中选择，则是在其中逼近条件均值，不保证能精确表示它。方差项也说明即使预测了条件均值，风险仍可能因噪声而为正。绝对损失则以条件中位数为最优预测之一，条件中位数不唯一时最优解也可能不唯一；因此同一数据上的“最佳回归函数”需要连同损失说明。

实际回归可用均方误差或平均绝对误差评价。若平方损失与关于系数线性的模型结合，就得到最小二乘问题，例如例~\ref{ex:lh-polynomial-fitting}；一般非线性模型的平方损失优化则未必有解析解或全局唯一解。时间预测还应使用预测时确实可得的特征，并采用符合时间顺序的评价方式。

\section*{本章概要}
\addcontentsline{toc}{section}{本章概要}

统计学习以有限数据为依据，在假设空间内按给定准则选择模型，再用算法实现求解。监督、无监督和强化学习的主要区别在于可利用的反馈与任务目标；概率性、线性、参数化和在线性则考察方法的不同侧面。理解一个具体方法，应分别指出数据形式、模型族、损失与正则项、求解方式和评价目标。

训练误差是拟合依据，泛化误差才刻画给定分布下的新样本表现；验证用于选择方案，独立测试用于评价最终流程。正则化与交叉验证帮助控制选择，但效果仍取决于数据和模型假设。有限假设空间的 Hoeffding 界说明如何同时控制全部候选的偏差，以及为什么训练后选择的模型需要统一保证。生成与判别的区分在于是否建模联合分布，分类、标注与回归的区分则主要取决于输出结构和损失。

\section*{习题与解答}
\addcontentsline{toc}{section}{习题与解答}

以下对应教材本章习题 1.1、1.2，重点将估计公式放回模型、策略与算法的框架。

\begin{example}
\label{ex:lh-bernoulli-estimation}
教材习题 1.1。设 $X_1,\ldots,X_n$ 独立同分布于伯努利分布（Bernoulli distribution），$n\geq1$，$\Pr(X_i=1)=\theta$，$\theta\in[0,1]$。观测值为 $x_i\in\{0,1\}$，其中 $k=\sum_{i=1}^n x_i$。求极大似然估计与一种贝叶斯估计，并分别说明统计学习方法三要素。
\end{example}
\begin{solve}
模型是伯努利分布族 $\{p_\theta(x)=\theta^x(1-\theta)^{1-x}:\theta\in[0,1]\}$，各端点概率按伯努利分布直接定义。对已观测的有序样本，似然为
\[
    \mathcal L(\theta)=\prod_{i=1}^n p_\theta(x_i)
       =\theta^k(1-\theta)^{n-k}.
\]
若只观测到成功次数 $k$，二项似然多出与 $\theta$ 无关的组合系数，最优解相同。对 $0<k<n$，对数似然在 $(0,1)$ 上满足
\[
    \ell'(\theta)=\frac{k}{\theta}-\frac{n-k}{1-\theta},
    \qquad
    \ell''(\theta)=-\frac{k}{\theta^2}
                     -\frac{n-k}{(1-\theta)^2}<0,
\]
故其唯一最大点为 $k/n$。若 $k=0$，似然 $(1-\theta)^n$ 在 $\theta=0$ 达到最大；若 $k=n$，似然 $\theta^n$ 在 $\theta=1$ 达到最大。因此统一得到
\begin{equation}
    \widehat\theta_{\mathrm{ML}}=\frac{k}{n}.
    \label{eq:lh-bernoulli-mle}
\end{equation}
这里模型是上述分布族，策略是极大化似然或最小化平均负对数似然，算法是统计 $k$ 后使用解析式；读取全部样本需 $O(n)$ 时间、$O(1)$ 额外空间。

贝叶斯方法需另外指定先验。取 $\Theta\sim\operatorname{Beta}(\alpha,\beta)$，其中 $\alpha,\beta>0$，密度为
\[
    \pi(\theta)=\frac{\theta^{\alpha-1}(1-\theta)^{\beta-1}}
                         {B(\alpha,\beta)},\quad 0<\theta<1,
    \qquad
    B(a,b)=\int_0^1 t^{a-1}(1-t)^{b-1}\,\mathrm dt.
\]
似然乘先验后，各幂次相加，因此
\begin{equation}
    \Theta\mid x_{1:n}\sim\operatorname{Beta}(k+\alpha,n-k+\beta).
    \label{eq:lh-bernoulli-posterior}
\end{equation}
当 $k+\alpha>1$ 且 $n-k+\beta>1$ 时，对数后验严格凹且内部模态存在，其最大后验估计为
\begin{equation}
    \widehat\theta_{\mathrm{MAP}}
       =\frac{k+\alpha-1}{n+\alpha+\beta-2}.
    \label{eq:lh-bernoulli-map}
\end{equation}
若条件不满足，需另行讨论端点或无界密度，不能直接套用内部模态公式。例如均匀先验 $\alpha=\beta=1$ 下，按闭区间上的连续密度版本取模态，MAP 与 MLE 相同，包括 $k=0,n$ 的端点情形。

“贝叶斯估计”不唯一指 MAP。若参数估计采用平方损失，则后验风险满足
\[
    \mathbb E[(\Theta-a)^2\mid x_{1:n}]
       =\operatorname{Var}(\Theta\mid x_{1:n})
          +(\mathbb E[\Theta\mid x_{1:n}]-a)^2,
\]
故最优点估计是后验均值。由 $\operatorname{Beta}(A,B)$ 的均值为 $A/(A+B)$，得到
\begin{equation}
    \widehat\theta_{\mathrm{Bayes}}
       =\mathbb E[\Theta\mid x_{1:n}]
       =\frac{k+\alpha}{n+\alpha+\beta}.
    \label{eq:lh-bernoulli-posterior-mean}
\end{equation}
该均值公式也可直接由 $B(A+1,B)/B(A,B)=A/(A+B)$ 算出。新一次试验在给定 $\Theta$ 后与已有试验独立时，后验预测成功概率同样等于这个后验均值。均匀先验下为 $(k+1)/(n+2)$，与 MAP 的 $k/n$ 不同。

因此贝叶斯建模由伯努利似然与 Beta 先验共同规定；策略可以是完整后验推断后进行预测，也可以是在明确损失下选择点估计，或取 MAP。算法利用共轭形式更新两个形状参数，再计算所需的均值、模态或预测概率。只有明确所求对象后，才能完整说明其“策略”，不能把后验分布、MAP 和后验均值视为同一结果。
\end{solve}

\begin{example}
\label{ex:lh-erm-mle}
教材习题 1.2。设监督学习模型为条件概率分布族 $\{p_\theta(y\mid x):\theta\in\Omega\}$。证明在对数损失下，经验风险最小化与条件极大似然估计等价，并说明因子分解的条件。
\end{example}
\begin{solve}
假设给定参数时样本对独立同分布，且在给定全部输入后，各输出的条件分布按 $p_\theta(y_i\mid x_i)$ 分解，则条件似然为
\[
    \mathcal L_c(\theta;T)
       =\prod_{i=1}^N p_\theta(y_i\mid x_i).
\]
对数损失下的经验风险为
\[
    \widehat R_T(\theta)
       =-\frac1N\sum_{i=1}^N\log p_\theta(y_i\mid x_i)
       =-\frac1N\log\mathcal L_c(\theta;T).
\]
由于对数严格递增、乘以负常数反转大小关系，在相同参数域上且最优解存在时，
\[
    \arg\min_{\theta\in\Omega}\widehat R_T(\theta)
       =\arg\max_{\theta\in\Omega}\mathcal L_c(\theta;T).
\]
对零条件概率按扩展值 $-\log0=+\infty$ 处理，此关系仍保持；若最优值只能逼近，则应比较相应上确界、下确界及近似解，而不能声称存在一个已达到的估计。若建模的是完整联合似然 $\prod_i p_\theta(x_i,y_i)$，还会包含输入分布；只有输入分布不含待估参数，或明确只采用条件似然时，才能直接使用上述条件形式。序列内部存在依赖时，也应对整个序列的模型似然使用相应对数损失，不能任意分解为位置间独立的乘积。
\end{solve}

% 教材来源：机器学习/textbooks/机器学习（分章原文）/机器学习-2.pdf，正文第 35--47 页。


\fi

\selectchapter{2}
\ifcompilechapter
\chapter{感知机}
\label{chap:lh-perceptron}

给定带有两类标签的训练样本，最直接的分类思路是寻找一个超平面，使两类样本分别位于它的两侧。\emph{感知机（perceptron）}采用这一线性判别规则，并通过逐次纠正误分样本来确定超平面。它提供了理解监督学习的一条完整路线：先规定预测模型，再设计可计算的损失，最后研究更新算法在什么条件下能够终止。这里的收敛性是对有限训练集的分离保证，不直接给出未知数据上的分类误差保证。

全章使用训练集 $T=\{(x_i,y_i)\}_{i=1}^{N}$，其中 $N\geq 1$，$x_i\in\mathbb R^d$ 为第 $i$ 个样本的特征向量，$y_i\in\{-1,+1\}$ 为其类别标签。向量均视为列向量，$u^\top v$ 表示欧氏内积，$\|u\|$ 表示欧氏范数。样本下标写在右下角，向量坐标写为 $x^{(j)}$，以免混淆样本编号与特征编号。

\section{感知机模型}
\label{sec:lh-perceptron-model}

\begin{definition}[感知机模型]
\label{def:lh-perceptron}
设输入空间 $\mathcal X\subseteq\mathbb R^d$，输出空间 $\mathcal Y=\{-1,+1\}$。给定权重向量（weight vector）$w\in\mathbb R^d$ 和偏置（bias）$b\in\mathbb R$，感知机的得分函数和分类函数分别为
\begin{equation}
\label{eq:lh-perceptron-model}
s_{w,b}(x)=w^\top x+b,
\qquad
f_{w,b}(x)=\operatorname{sign}\bigl(s_{w,b}(x)\bigr).
\end{equation}
本笔记将符号函数在零点的取值约定为
\[
\operatorname{sign}(z)=
\begin{cases}
+1,&z\geq 0,\\
-1,&z<0.
\end{cases}
\]
这一约定只决定恰落在决策边界上的输入如何预测；训练时另行规定边界样本仍须触发更新。
\end{definition}

感知机直接由特征给出类别，属于\emph{判别模型（discriminative model）}。其假设空间是
\[
\mathcal H=
\left\{x\longmapsto\operatorname{sign}(w^\top x+b):
w\in\mathbb R^d,\ b\in\mathbb R\right\}.
\]
其中 $s_{w,b}$ 是实值得分，$f_{w,b}$ 才是取两个类别值的分类器；二者不能混作同一个函数。得分本身没有被建模为概率，也不能直接解释为分类置信概率。

当 $w\neq 0$ 时，方程
\begin{equation}
\label{eq:lh-perceptron-hyperplane}
S=\{x\in\mathbb R^d:w^\top x+b=0\}
\end{equation}
定义一个超平面（hyperplane），其法向量为 $w$，两个开半空间分别满足 $s_{w,b}(x)>0$ 和 $s_{w,b}(x)<0$。正方向由 $w$ 指向得分增加的一侧。对任意 $c>0$，$(w,b)$ 与 $(cw,cb)$ 给出同一分类器，因此参数通常不唯一。$w=0$ 时得分为常数，不能将式~\eqref{eq:lh-perceptron-hyperplane} 解释为通常的超平面，但这样的参数可作为算法的初始状态。

训练阶段利用标签确定 $w,b$；预测阶段则固定所得参数，对新输入 $x$ 计算式~\eqref{eq:lh-perceptron-model}。训练集能否被分开，首先取决于它相对于当前特征表示是否线性可分，而不是取决于训练算法运行了多少次。

\section{感知机学习策略}
\label{sec:lh-perceptron-strategy}

\subsection{数据集的线性可分性}
\label{subsec:lh-perceptron-separability}

\begin{definition}[线性可分性]
\label{def:lh-linear-separability}
若存在 $w\in\mathbb R^d$、$b\in\mathbb R$，使训练集中每个样本都满足
\begin{equation}
\label{eq:lh-perceptron-separable}
y_i(w^\top x_i+b)>0,
\qquad i=1,\ldots,N,
\end{equation}
则称 $T$ 为\emph{线性可分数据集（linearly separable data set）}；否则称为线性不可分。满足上述不等式的超平面称为分离超平面（separating hyperplane）。
\end{definition}

式~\eqref{eq:lh-perceptron-separable} 将两类样本的条件合为一个式子：$y_i=+1$ 时要求得分严格为正，$y_i=-1$ 时要求得分严格为负。这里强调严格不等式，因而没有训练样本落在分离超平面上。当两类样本均存在时，分离向量必有 $w\neq0$；若训练集只有一类，则可以用常数偏置完成分类。

\begin{definition}[函数间隔与几何间隔]
\label{def:lh-perceptron-margin}
对参数 $(w,b)$，第 $i$ 个样本的\emph{函数间隔（functional margin）}为
\begin{equation}
\label{eq:lh-perceptron-functional-margin}
m_i(w,b)=y_i(w^\top x_i+b).
\end{equation}
当 $w\neq0$ 时，其有符号几何间隔（geometric margin）为 $m_i(w,b)/\|w\|$，点 $x_i$ 到超平面的无符号距离为 $|w^\top x_i+b|/\|w\|$。
\end{definition}

上述距离公式可由正交分解得到：从 $x_i$ 沿法向量方向移动至超平面，垂足为
\[
x_i-\frac{w^\top x_i+b}{\|w\|^2}w,
\]
移动长度就是 $|w^\top x_i+b|/\|w\|$。函数间隔为正表示样本处于正确一侧，为负表示误分，为零表示落在边界上。正比例缩放参数会改变函数间隔，但不改变几何间隔。

\subsection{从误分类到感知机损失}
\label{subsec:lh-perceptron-loss}

直接最小化误分类数需要优化不连续的类别判断，不利于利用梯度信息。教材从误分点到分离超平面的距离引出替代目标：对于 $m_i(w,b)<0$ 的点，其到超平面的距离为 $-m_i(w,b)/\|w\|$。去掉公共归一化因子 $1/\|w\|$ 后，可得到只惩罚负函数间隔的感知机损失。\emph[caution]{由于 $\|w\|$ 也随参数改变，去掉该因子是选择新的损失函数，并非保持原距离最小化问题不变的代数化简。}

\begin{definition}[感知机损失]
\label{def:lh-perceptron-loss}
单个样本的感知机损失和训练集的总损失分别定义为
\begin{equation}
\label{eq:lh-perceptron-loss}
\ell_i(w,b)=\max\{0,-y_i(w^\top x_i+b)\},
\qquad
L(w,b)=\sum_{i=1}^{N}\ell_i(w,b).
\end{equation}
若将误分或位于边界的样本指标集记为
\[
M(w,b)=\{i:y_i(w^\top x_i+b)\leq0\},
\]
则等价地有
\begin{equation}
\label{eq:lh-perceptron-loss-active}
L(w,b)=-\sum_{i\in M(w,b)}y_i(w^\top x_i+b).
\end{equation}
边界样本在和式中的贡献为零。采用平均损失 $L/N$ 只改变整体尺度，不改变其最小值点集合。
\end{definition}

若 $w\neq0$，全部误分点到超平面的距离之和为 $L(w,b)/\|w\|$，而 $L$ 本身是负函数间隔之和。不同参数之间不能只凭误分点个数判断 $L$ 的大小：误分点较少时，某个点的负间隔仍可能很大。反过来，参数整体缩小时，$L$ 可以减小而分类结果完全不变。

\begin{proposition}[损失的基本性质]
\label{prop:lh-perceptron-loss-properties}
函数 $L$ 非负、连续、凸且分段线性，但一般不处处可导。它满足
\begin{equation}
\label{eq:lh-perceptron-zero-loss}
L(w,b)=0
\quad\Longleftrightarrow\quad
y_i(w^\top x_i+b)\geq0\quad\text{对所有 }i.
\end{equation}
严格分离蕴含零损失，零损失一般不蕴含严格分离。
\end{proposition}

\begin{proof}
每个 $\ell_i$ 是常数函数 $0$ 与仿射函数 $-y_i(w^\top x_i+b)$ 的最大值，因此非负、连续、凸且分段线性，有限求和保留这些性质。两个分支交界处的斜率一般不同，所以不能据此声称函数处处可导。有限个非负项之和为零当且仅当每项均为零，这就得到式~\eqref{eq:lh-perceptron-zero-loss}。严格正间隔当然满足非负间隔；反例是 $(w,b)=(0,0)$，此时任何训练集都满足 $L(0,0)=0$，但所有间隔均为零。
\end{proof}

教材将学习策略写为
\begin{equation}
\label{eq:lh-perceptron-objective}
\min_{w\in\mathbb R^d,\ b\in\mathbb R} L(w,b).
\end{equation}
理解这一表达时必须同时保留算法的更新和停止规则：由于零参数总是式~\eqref{eq:lh-perceptron-objective} 的最小值点，单独求解这个无约束最小化问题并不能保证找到严格分离器。感知机在间隔为零时仍更新，直至所有训练间隔严格为正。因此，\emph{零损失条件与算法终止条件并不相同}，线性可分条件也不会因为损失函数是凸函数而自动成立。

\section{感知机学习算法}
\label{sec:lh-perceptron-algorithm}

\subsection{原始形式与逐次更新}
\label{subsec:lh-perceptron-primal}

在没有样本跨越边界的局部区域内，负间隔样本集保持不变，此时损失的梯度为
\[
\nabla_w L=-\sum_{i:m_i<0}y_i x_i,
\qquad
\frac{\partial L}{\partial b}=-\sum_{i:m_i<0}y_i.
\]
若一次只选择一个间隔非正的样本 $(x_i,y_i)$，就得到原始形式（primal form）的更新
\begin{equation}
\label{eq:lh-perceptron-primal-update}
w\leftarrow w+\eta y_i x_i,
\qquad
b\leftarrow b+\eta y_i,
\qquad \eta>0.
\end{equation}
其中 $\eta$ 为固定学习率（learning rate）。教材采用不大于 $1$ 的步长，并在例题中取 $1$；本章的有限更新证明只要求它是固定正常数。

从优化角度看，这是逐样本的\emph{次梯度（subgradient）}更新。对凸函数 $h$，若向量 $g$ 满足对所有 $v$ 都有 $h(v)\geq h(u)+g^\top(v-u)$，则称 $g$ 为 $h$ 在 $u$ 处的次梯度，所有这样的向量组成次微分 $\partial h(u)$。令扩充参数 $\widetilde w=(w^\top,b)^\top$，扩充输入 $\widetilde x_i=(x_i^\top,1)^\top$，则
\[
\partial\ell_i(\widetilde w)=
\begin{cases}
\{-y_i\widetilde x_i\},&m_i<0,\\
\{-\lambda y_i\widetilde x_i:0\leq\lambda\leq1\},&m_i=0,\\
\{0\},&m_i>0.
\end{cases}
\]
边界处选择 $\lambda=1$ 就得到式~\eqref{eq:lh-perceptron-primal-update}，而选择 $\lambda=0$ 会允许算法停在零参数。教材用随机梯度下降（stochastic gradient descent）描述逐点修正过程；严格地说，在不可导位置应使用次梯度，且固定顺序遍历样本也能得到后面的有限更新保证，无须依赖随机抽样。

原始形式的输入是训练集 $T$、学习率 $\eta>0$ 和初值 $(w_0,b_0)$，输出是分类参数与式~\eqref{eq:lh-perceptron-model} 中的分类器。为使选点与停止条件明确，可采用下面的循环遍历实现。
\begin{enumerate}
  \item 初始化 $w=w_0$、$b=b_0$，常用 $w_0=0$、$b_0=0$。
  \item 开始一轮训练，将本轮更新次数置为零，依次访问 $i=1,\ldots,N$。
  \item 每访问一个样本，计算当前间隔 $m_i=y_i(w^\top x_i+b)$；若 $m_i\leq0$，执行式~\eqref{eq:lh-perceptron-primal-update} 并累计一次更新；否则保持参数不变。
  \item 完整一轮没有发生更新时停止；否则开始下一轮。在线性可分性未知时另设最大轮数，达到上限只能报告尚未找到严格分离器，不能判定数据必定不可分。
\end{enumerate}

一轮无更新意味着参数在整轮中固定，故所有样本在最终参数下的间隔均为正；仅仅在一轮中让每个点都曾经被正确分类，并不足以停止，因为后续更新可能破坏先前点的分类。教材允许直接选择任意间隔非正的点；循环遍历则保证只要还有这样的点，就会继续发现并处理它。

更新对当前点的作用可以精确写成
\begin{equation}
\label{eq:lh-perceptron-margin-increase}
y_i\bigl((w+\eta y_i x_i)^\top x_i+b+\eta y_i\bigr)
=m_i+\eta(\|x_i\|^2+1)>m_i.
\end{equation}
但一次更新未必把该点的间隔直接变为正。对其他样本 $(x_j,y_j)$，间隔改变量为 $\eta y_i y_j(x_i^\top x_j+1)$，可正可负，所以不能由式~\eqref{eq:lh-perceptron-margin-increase} 推出总损失每步下降，或所有样本同时得到改善。

\begin{example}
\label{ex:lh-perceptron-primal}
教材例 2.1：正例为 $x_1=(3,3)^\top$、$x_2=(4,3)^\top$，负例为 $x_3=(1,1)^\top$。从 $w_0=0$、$b_0=0$ 出发，取 $\eta=1$，用原始形式求一个感知机分类器。
\end{example}

\begin{solve}
每次从样本 $1$ 开始检查，选择第一个间隔非正的点更新，再重新检查。这样得到的更新点顺序为 $x_1,x_3,x_3,x_3,x_1,x_3,x_3$。下表中 $k$ 只计发生参数更新的次数。
\begin{center}
\begin{tabular}{cccc}
\toprule
$k$ & 更新所用样本 & $w_k^\top$ & $b_k$ \\
\midrule
0 & 初值 & $(0,0)$ & $0$ \\
1 & $x_1$ & $(3,3)$ & $1$ \\
2 & $x_3$ & $(2,2)$ & $0$ \\
3 & $x_3$ & $(1,1)$ & $-1$ \\
4 & $x_3$ & $(0,0)$ & $-2$ \\
5 & $x_1$ & $(3,3)$ & $-1$ \\
6 & $x_3$ & $(2,2)$ & $-2$ \\
7 & $x_3$ & $(1,1)$ & $-3$ \\
\bottomrule
\end{tabular}
\end{center}
例如第 $4$ 次更新前，$x_3$ 的间隔为 $-((1,1)(1,1)^\top-1)=-1$，故仍需修正；更新后 $w=0$、$b=-2$，两个正例又需要重新检查。这说明中间状态不必始终对应非退化的超平面，也不必保持先前的分类结果。

最终三个样本的间隔依次为 $3,4,1$，均严格为正，因此得到分离直线 $x^{(1)}+x^{(2)}-3=0$ 和分类器
\[
f(x)=\operatorname{sign}\bigl(x^{(1)}+x^{(2)}-3\bigr).
\]
这一解不是唯一的。例如 $2x^{(1)}+x^{(2)}-5=0$ 也能严格分离相同数据，其三个间隔分别为 $4,6,2$。感知机没有规定从多个分离器中选择最大间隔者，不同选点顺序或初值可能给出不同结果。
\end{solve}

\begin{notealert}[type=caution,title={损失值不能替代训练停止检查}]
例~\ref{ex:lh-perceptron-primal} 中初始总损失为 $0$，第一次更新后负例的得分为 $7$，总损失反而变为 $7$。这并不违反后面的收敛定理：该定理控制的是间隔非正所触发的更新次数，并不要求 $L$ 沿迭代单调递减。实现时应检查所有训练间隔是否严格为正，不能以“总损失已为零”替代算法的停止规则。
\end{notealert}

\subsection{算法的收敛性}
\label{subsec:lh-perceptron-convergence}

为统一处理权重与偏置，沿用扩充表示
\begin{equation}
\label{eq:lh-perceptron-augmentation}
\widetilde x_i=(x_i^\top,1)^\top\in\mathbb R^{d+1},
\qquad
\widetilde w=(w^\top,b)^\top\in\mathbb R^{d+1}.
\end{equation}
于是 $\widetilde w^\top\widetilde x_i=w^\top x_i+b$，更新成为 $\widetilde w\leftarrow\widetilde w+\eta y_i\widetilde x_i$。\emph[caution]{扩充参数的范数为 $\sqrt{\|w\|^2+b^2}$，与输入空间超平面距离公式中使用的 $\|w\|$ 不同。}

\begin{theorem}[感知机有限更新定理，Novikoff]
\label{thm:lh-perceptron-convergence}
设有限训练集 $T$ 线性可分，使用固定学习率 $\eta>0$，从 $\widetilde w_0=0$ 开始，只在 $y_i\widetilde w^\top\widetilde x_i\leq0$ 时更新。则存在单位向量 $\widetilde u\in\mathbb R^{d+1}$ 及常数 $\gamma>0$，使
\begin{equation}
\label{eq:lh-perceptron-separator-margin}
\|\widetilde u\|=1,
\qquad
y_i\widetilde u^\top\widetilde x_i\geq\gamma,
\quad i=1,\ldots,N.
\end{equation}
令 $R=\max_{1\leq i\leq N}\|\widetilde x_i\|$，则算法发生的更新总次数 $K$ 满足
\begin{equation}
\label{eq:lh-perceptron-mistake-bound}
K\leq\left(\frac{R}{\gamma}\right)^2.
\end{equation}
只要选点规则保证持续检查整个训练集，算法就会在有限次检查后得到所有训练间隔严格为正的参数。
\end{theorem}

\begin{proof}
先说明单位分离向量及正间隔的存在性。由线性可分性，存在非零扩充参数 $\widetilde v$，使 $y_i\widetilde v^\top\widetilde x_i>0$ 对所有 $i$ 成立。令 $\widetilde u=\widetilde v/\|\widetilde v\|$，并取
\[
\gamma=\min_{1\leq i\leq N}y_i\widetilde u^\top\widetilde x_i.
\]
因为样本数有限，有限个严格正数的最小值仍严格为正，故式~\eqref{eq:lh-perceptron-separator-margin} 成立。$\widetilde u$ 在这里仅是用于分析的一个单位分离向量，无须预先求得，更没有假设它使间隔达到最大。

用 $i_k$ 表示第 $k$ 次更新选中的样本编号。更新前有
\[
y_{i_k}\widetilde w_{k-1}^\top\widetilde x_{i_k}\leq0,
\qquad
\widetilde w_k=\widetilde w_{k-1}+\eta y_{i_k}\widetilde x_{i_k}.
\]
证明的关键是比较两个增长速度：参数在正确分离方向上的投影至少按 $k$ 线性增长，而参数范数至多按 $\sqrt{k}$ 增长。

对投影，由式~\eqref{eq:lh-perceptron-separator-margin} 得
\[
\widetilde w_k^\top\widetilde u
=\widetilde w_{k-1}^\top\widetilde u
+\eta y_{i_k}\widetilde u^\top\widetilde x_{i_k}
\geq\widetilde w_{k-1}^\top\widetilde u+\eta\gamma.
\]
从零初值递推可得
\begin{equation}
\label{eq:lh-perceptron-projection-bound}
\widetilde w_k^\top\widetilde u\geq k\eta\gamma.
\end{equation}
另一方面，由更新条件及 $\|\widetilde x_{i_k}\|\leq R$，
\[
\begin{aligned}
\|\widetilde w_k\|^2
&=\|\widetilde w_{k-1}\|^2
+2\eta y_{i_k}\widetilde w_{k-1}^\top\widetilde x_{i_k}
+\eta^2\|\widetilde x_{i_k}\|^2\\
&\leq\|\widetilde w_{k-1}\|^2+\eta^2R^2.
\end{aligned}
\]
继续递推得到
\begin{equation}
\label{eq:lh-perceptron-norm-bound}
\|\widetilde w_k\|^2\leq k\eta^2R^2.
\end{equation}
若 $k\geq1$，结合式~\eqref{eq:lh-perceptron-projection-bound}、式~\eqref{eq:lh-perceptron-norm-bound} 和 Cauchy--Schwarz 不等式，有
\[
k\eta\gamma
\leq\widetilde w_k^\top\widetilde u
\leq\|\widetilde w_k\|\|\widetilde u\|
\leq\eta R\sqrt{k}.
\]
除以正数 $\eta\sqrt{k}\gamma$ 并平方，得到 $k\leq(R/\gamma)^2$；$k=0$ 时同样满足该界。因此不可能出现无限次更新。若循环遍历的某一阶段仍有非正间隔样本，则至多再检查一整轮就会发现可更新样本；故更新次数有限必然导致某一整轮无更新，并在严格分离状态停止。
\end{proof}

式~\eqref{eq:lh-perceptron-mistake-bound} 中的 $K$ 是更新次数，包含落在边界时的更新，不是遍历轮数，也不是程序运行时间。证明不依赖误分点的具体选择顺序；对于给定的单位分离向量，$R/\gamma$ 越大，所得最坏情况上界越宽松。这里的 $\gamma$ 是扩充空间中相对于单位参数的间隔，不能直接与输入空间中除以 $\|w\|$ 的几何间隔混用。

上界不含固定学习率 $\eta$，这也可以从零初值的缩放性质理解：在相同选点规则下，将学习率由 $1$ 改为 $\eta>0$，每一步参数都恰为原参数的 $\eta$ 倍，间隔符号与更新点顺序不变。该结论针对本章没有额外阈值或正则项的更新规则，并不表示一般优化算法的行为都与学习率无关。

对任意固定非零初值，有限更新性也可由同一思路得到，但不能原样套用式~\eqref{eq:lh-perceptron-mistake-bound}。此时两条递推分别成为
\[
\widetilde w_k^\top\widetilde u
\geq\widetilde w_0^\top\widetilde u+k\eta\gamma,
\qquad
\|\widetilde w_k\|^2
\leq\|\widetilde w_0\|^2+k\eta^2R^2.
\]
若更新次数能够无限增长，左侧投影的线性下界最终会超过由范数给出的平方根上界，产生矛盾。因此仍会有限终止，只是次数界还依赖初值。

若训练集线性不可分，就不存在满足式~\eqref{eq:lh-perceptron-separator-margin} 的正数 $\gamma$，上述证明失去前提。持续遍历并坚持严格正间隔的停止条件时，算法无法终止；参数可能反复变化，不能把有限可分情形的结论推广过去。最大轮数只是计算预算，若要处理不可分数据，需要另行选择允许误差的目标或算法。即使训练集可分，所得分类器对未见数据的表现仍需独立评价。

\subsection{对偶形式与 Gram 矩阵}
\label{subsec:lh-perceptron-dual}

从 $w_0=0$、$b_0=0$ 出发，每次只增加 $\eta y_i x_i$ 和 $\eta y_i$。若第 $i$ 个样本共触发 $n_i$ 次更新，令 $\alpha_i=\eta n_i\geq0$，则任一迭代时刻都有
\begin{equation}
\label{eq:lh-perceptron-dual-representation}
w=\sum_{j=1}^{N}\alpha_jy_jx_j,
\qquad
b=\sum_{j=1}^{N}\alpha_jy_j.
\end{equation}
将其代入得分可得
\begin{equation}
\label{eq:lh-perceptron-dual-score}
s(x)=\sum_{j=1}^{N}\alpha_jy_jx_j^\top x+b.
\end{equation}
这就是\emph{对偶形式（dual form）}的基本想法：以样本线性组合的系数表示权重，用内积完成训练和预测。这里的“对偶”首先是一种等价参数表示，不需要另行声称已经推导出某个拉格朗日对偶优化问题。

\begin{definition}[Gram 矩阵]
\label{def:lh-perceptron-gram}
训练样本的\emph{Gram 矩阵（Gram matrix）}定义为
\begin{equation}
\label{eq:lh-perceptron-gram}
G=(G_{ij})_{N\times N},
\qquad G_{ij}=x_i^\top x_j.
\end{equation}
它是对称半正定矩阵。事实上，对任意 $c\in\mathbb R^N$，
\[
c^\top Gc
=\left\|\sum_{i=1}^{N}c_i x_i\right\|^2\geq0.
\]
若采用扩充输入，则相应矩阵为 $\widetilde G_{ij}=G_{ij}+1$。
\end{definition}

对偶算法输入训练集 $T$ 和固定学习率 $\eta>0$，输出系数 $\alpha$、偏置 $b$ 及式~\eqref{eq:lh-perceptron-dual-score} 对应的分类器。先计算 $G$，并初始化 $\alpha_1=\cdots=\alpha_N=0$、$b=0$；按原始形式相同的规则访问样本。当
\begin{equation}
\label{eq:lh-perceptron-dual-condition}
y_i\left(\sum_{j=1}^{N}\alpha_jy_jG_{ji}+b\right)\leq0
\end{equation}
时，执行
\begin{equation}
\label{eq:lh-perceptron-dual-update}
\alpha_i\leftarrow\alpha_i+\eta,
\qquad
b\leftarrow b+\eta y_i,
\end{equation}
其余系数不变。完整一轮没有更新时停止；可分性未知时使用与原始形式相同的轮数上限及未完成说明。由于 $G$ 对称，式~\eqref{eq:lh-perceptron-dual-condition} 中也可写 $G_{ij}$，但必须保留求和系数中的标签 $y_j$。

\begin{proposition}[两种形式的逐步等价性]
\label{prop:lh-perceptron-primal-dual}
若原始形式与对偶形式都从零初值开始，使用相同学习率，并按相同规则访问和选择样本，则每次更新后两者的 $w,b$ 都满足式~\eqref{eq:lh-perceptron-dual-representation}，得分及更新条件完全一致。
\end{proposition}

\begin{proof}
初始时所有参数为零，等式成立。假设某一步成立，则原始形式在样本 $x_i$ 上的得分为
\[
w^\top x_i+b
=\sum_{j=1}^{N}\alpha_jy_jG_{ji}+b,
\]
因而两种形式决定是否更新的条件相同。若更新，对偶形式只令 $\alpha_i$ 增加 $\eta$，故所表示的权重增加 $\eta y_i x_i$，偏置增加 $\eta y_i$，恰好与原始形式一致；不更新时等式保持。由归纳法，结论对所有迭代成立。
\end{proof}

偏置既可以作为独立变量递推，也可以每次由 $b=\sum_j\alpha_jy_j$ 恢复，两者不能重复累加。若直接使用扩充 Gram 矩阵，则得分应写成 $\sum_j\alpha_jy_j\widetilde G_{ji}$，此时已经包含偏置贡献，不能再加一个 $b$。零初始化是式~\eqref{eq:lh-perceptron-dual-representation} 的前提；一般初值时还需保留 $w_0,b_0$ 项。

\begin{example}
\label{ex:lh-perceptron-dual}
教材例 2.2：对例~\ref{ex:lh-perceptron-primal} 的三个样本，使用同一更新点顺序，以对偶形式求感知机模型。
\end{example}

\begin{solve}
直接计算内积得到
\[
G=
\begin{pmatrix}
18&21&6\\
21&25&7\\
6&7&2
\end{pmatrix}.
\]
取 $\eta=1$，按更新点顺序 $x_1,x_3,x_3,x_3,x_1,x_3,x_3$ 计算，有
\begin{center}
\begin{tabular}{ccccc}
\toprule
$k$ & $\alpha_1$ & $\alpha_2$ & $\alpha_3$ & $b$ \\
\midrule
0 & 0 & 0 & 0 & 0 \\
1 & 1 & 0 & 0 & 1 \\
2 & 1 & 0 & 1 & 0 \\
3 & 1 & 0 & 2 & $-1$ \\
4 & 1 & 0 & 3 & $-2$ \\
5 & 2 & 0 & 3 & $-1$ \\
6 & 2 & 0 & 4 & $-2$ \\
7 & 2 & 0 & 5 & $-3$ \\
\bottomrule
\end{tabular}
\end{center}
因此
\[
w=2x_1-5x_3=(1,1)^\top,
\qquad
b=2-5=-3,
\]
所得模型与例~\ref{ex:lh-perceptron-primal} 相同。$\alpha_2=0$ 表示该点在这次训练轨迹中没有触发更新，并不表示它的标签没有约束训练结果；算法停止时仍须核查它的间隔为正。
\end{solve}

对稠密 $d$ 维特征，原始形式一次得分计算和一次权重更新各需 $O(d)$ 运算。预计算 Gram 矩阵通常需 $O(N^2d)$ 运算、$O(N^2)$ 存储，之后按式~\eqref{eq:lh-perceptron-dual-condition} 计算一个训练点的得分需 $O(N)$ 运算，而系数更新只涉及一个分量。因此对偶形式并非无条件更快，是否合算取决于样本数、维数及存储条件。训练完成后，也可以先恢复 $w$ 再用原始形式预测；若保留系数表示，则预测只需对非零 $\alpha_j$ 对应的样本求和。

系数 $\alpha_i/\eta$ 记录样本在特定训练轨迹中的更新次数，反映其对这条轨迹的累积贡献，但不能单凭更新次数断定该点到最终超平面的距离更小。它还依赖样本访问顺序及中间参数。对偶形式揭示了内积的作用，为以后引入核函数提供入口；本章的等价性与收敛性论证仍针对这里明确给出的线性内积。

\section{章末习题与进一步理解}
\label{sec:lh-perceptron-exercises}

\subsection{异或问题的不可分性}
\label{subsec:lh-perceptron-xor}

\begin{example}
\label{ex:lh-perceptron-xor}
教材习题 2.1：将二元异或（exclusive OR，XOR）的输出 $1$ 编为标签 $+1$，输出 $0$ 编为标签 $-1$，得到
\[
(0,0)\mapsto-1,\qquad
(1,0)\mapsto+1,\qquad
(0,1)\mapsto+1,\qquad
(1,1)\mapsto-1.
\]
证明单个感知机不能表示这一分类规则。
\end{example}

\begin{solve}
若存在严格分离参数 $w=(w_1,w_2)^\top,b$，四个样本依次要求
\[
b<0,\qquad w_1+b>0,\qquad w_2+b>0,\qquad w_1+w_2+b<0.
\]
两个正例的约束相加得 $w_1+w_2+2b>0$，两个负例的约束相加则得 $w_1+w_2+2b<0$，矛盾。因此数据不满足定义~\ref{def:lh-linear-separability}。

即使允许正例落在边界，并按定义~\ref{def:lh-perceptron} 将零得分判为 $+1$，也不能表示 XOR：此时两个正例仍要求 $w_1+w_2+2b\geq0$，两个负例仍要求该量严格小于零，矛盾不变。这个限制来自原始特征空间中单个线性边界的表达能力，增加训练轮数无法消除它。
\end{solve}

\subsection{自行构造训练例子}
\label{subsec:lh-perceptron-constructed-example}

\begin{example}
\label{ex:lh-perceptron-constructed}
教材习题 2.2 要求构造从训练集求感知机模型的例子。取一维数据 $x_1=2,y_1=+1$，$x_2=1,y_2=-1$，从 $w_0=b_0=0$ 开始，令 $\eta=1$，每次重新按 $x_1,x_2$ 的顺序寻找间隔非正的样本。求迭代结果，并核查是否严格分离。
\end{example}

\begin{solve}
一维时参数 $w$ 为实数。逐步计算得到
\begin{center}
\begin{tabular}{cccc}
\toprule
$k$ & 更新所用样本 & $w_k$ & $b_k$ \\
\midrule
0 & 初值 & 0 & 0 \\
1 & $x_1$ & 2 & 1 \\
2 & $x_2$ & 1 & 0 \\
3 & $x_2$ & 0 & $-1$ \\
4 & $x_1$ & 2 & 0 \\
5 & $x_2$ & 1 & $-1$ \\
6 & $x_2$ & 0 & $-2$ \\
7 & $x_1$ & 2 & $-1$ \\
8 & $x_2$ & 1 & $-2$ \\
9 & $x_1$ & 3 & $-1$ \\
10 & $x_2$ & 2 & $-2$ \\
11 & $x_2$ & 1 & $-3$ \\
12 & $x_1$ & 3 & $-2$ \\
13 & $x_2$ & 2 & $-3$ \\
\bottomrule
\end{tabular}
\end{center}
例如 $(w,b)=(1,-2)$ 时，正例的得分为零，必须继续更新；$(w,b)=(2,-2)$ 时，负例的得分为零，也必须更新。最终得到 $f(x)=\operatorname{sign}(2x-3)$，两个训练间隔均为 $1$，严格分离点为 $x=3/2$。这一例子突出边界处理：若将判据误写为 $m_i<0$，从零初值出发就不会发生任何更新。
\end{solve}

\subsection{线性可分与两类凸包不相交}
\label{subsec:lh-perceptron-convex-hulls}

\begin{definition}[有限点集的凸包]
\label{def:lh-perceptron-convex-hull}
对非空有限点集 $A=\{a_1,\ldots,a_r\}\subset\mathbb R^d$，其\emph{凸包（convex hull）}为所有凸组合构成的集合：
\[
\operatorname{conv}(A)=
\left\{\sum_{j=1}^{r}\lambda_j a_j:
\lambda_j\geq0,\ \sum_{j=1}^{r}\lambda_j=1\right\}.
\]
它是非空紧凸集：非负且和为 $1$ 的系数集合是有限维闭有界单纯形，凸包是该紧集在线性映射下的像；两组凸组合再作凸组合仍为凸组合。
\end{definition}

\begin{theorem}[两类凸包判据]
\label{thm:lh-perceptron-convex-hulls}
教材习题 2.3：设有限训练集的正例点集 $A=\{x_i:y_i=+1\}$ 与负例点集 $B=\{x_i:y_i=-1\}$ 均非空。则训练集线性可分，当且仅当
\begin{equation}
\label{eq:lh-perceptron-disjoint-hulls}
\operatorname{conv}(A)\cap\operatorname{conv}(B)=\varnothing.
\end{equation}
\end{theorem}

\begin{proof}
记 $C_+=\operatorname{conv}(A)$、$C_-=\operatorname{conv}(B)$。先证必要性。若 $s(x)=w^\top x+b$ 严格分离两个点集，则任意正例凸组合 $a=\sum_j\lambda_j a_j$ 满足
\[
s(a)=\sum_j\lambda_j s(a_j)>0,
\]
这里用了系数非负、和为 $1$，且至少有一个系数为正。同理，每个 $c\in C_-$ 都满足 $s(c)<0$。一个点不可能同时具有正得分和负得分，故两个凸包不相交。

再证充分性。假设两个凸包不相交。由于它们是非空紧集，连续函数 $(a,c)\mapsto\|a-c\|$ 在 $C_+\times C_-$ 上取得最小值。取最近点对 $p\in C_+$、$q\in C_-$，则 $p\neq q$，令
\[
w=p-q\neq0,
\qquad
b=-\frac12 w^\top(p+q).
\]
下面证明垂直于 $p-q$ 且经过两点中点的超平面严格分离两个凸包。

任取 $a\in C_+$，凸性保证 $p+t(a-p)\in C_+$ 对 $0\leq t\leq1$ 成立。由最近点性质，
\[
\|p+t(a-p)-q\|^2\geq\|p-q\|^2.
\]
展开后对 $t>0$ 除以 $t$，再令 $t\downarrow0$，得到 $w^\top(a-p)\geq0$。同样，对任意 $c\in C_-$，沿 $q+t(c-q)$ 变化有
\[
\|p-q-t(c-q)\|^2\geq\|p-q\|^2,
\]
从而 $w^\top(c-q)\leq0$。因此
\[
\begin{aligned}
w^\top a+b
&\geq w^\top p-\tfrac12w^\top(p+q)
=\tfrac12\|p-q\|^2>0,\\
w^\top c+b
&\leq w^\top q-\tfrac12w^\top(p+q)
=-\tfrac12\|p-q\|^2<0.
\end{aligned}
\]
两个凸包中的全部点均被严格分开，原训练样本作为凸包中的点也被严格分开，故训练集线性可分。
\end{proof}

有限性在充分性证明中保证了凸包紧性及最近点对的存在；不能不加条件地把这一论证照搬到任意不相交集合。对 XOR，两个正例的凸包与两个负例的凸包都包含 $(1/2,1/2)^\top$，由定理~\ref{thm:lh-perceptron-convex-hulls} 可再次看出不可分。若同一个特征向量在训练集中具有相反标签，它本身就同时属于两个凸包，也立即排除了严格线性分离。

本章的核心联系是：模型用仿射得分的符号分类，学习规则以非正函数间隔触发更新，有限可分性通过单位分离方向上的正间隔给出更新次数界，对偶形式则将相同轨迹表示为样本内积及非负更新系数。损失的连续凸性、算法有限终止和未知样本上的泛化是三个不同层面的性质，使用时应分别核对其定义与条件。

% 教材来源：机器学习/textbooks/机器学习（分章原文）/机器学习-3.pdf，第 49--58 页。


\fi

\selectchapter{3}
\ifcompilechapter
\chapter{$k$ 近邻法}
\label{chap:lihang-knn}

给定已经标注类别的样本，一个直接的分类想法是：相近的输入往往具有相近的输出，因此可用新输入附近的已知样本推断它的类别。\emph{$k$ 近邻法}（$k$-nearest neighbor，$k$-NN）把这一想法落实为先找邻居、再作决策的过程。本章依据教材第 3 章，依次讨论近邻分类算法、决定模型的三个基本要素，以及用 kd 树加速近邻搜索的方法；重点区分分类规则本身与实现该规则的搜索算法。近邻法也可用于回归，本章与教材一致，只讨论分类。

\section{$k$ 近邻算法}
\label{sec:lihang-knn-algorithm}

设训练集为 $T=\{(x_i,y_i)\}_{i=1}^{N}$，其中 $N\geq1$，$x_i\in\mathcal X\subseteq\mathbb R^d$，$y_i\in\mathcal Y=\{c_1,\ldots,c_C\}$，$d$ 为特征维数，$C$ 为类别数。给定查询点 $x\in\mathcal X$、距离度量 $\rho$ 和整数 $1\leq k\leq N$，先按 $\rho(x,x_i)$ 从小到大选取 $k$ 个训练样本，再根据它们的标签确定预测类别。为使算法在距离相等时仍有唯一结果，约定距离并列时优先选择索引较小的样本，并把选出的样本索引集合记为 $N_k(x)$；使用索引可以保留坐标相同的不同观测，而不把它们合并成一个点。

\begin{definition}
\label{def:lihang-knn-rule}
在上述约定下，采用\emph{多数表决规则}（majority voting rule）的 $k$ 近邻分类器定义为
\begin{equation}
  \widehat f_T(x)
  =\underset{c_j\in\mathcal Y}{\operatorname{arg\,max}}
    \sum_{i\in N_k(x)}\mathbf 1\{y_i=c_j\},
  \label{eq:lihang-knn-vote}
\end{equation}
其中 $\mathbf 1\{A\}$ 是事件或条件 $A$ 的指示函数，条件成立时取 $1$，否则取 $0$。若多个类别获得相同的最高票数，固定选择类别索引 $j$ 最小者。这里的“多数”指票数最多，不要求某一类别的票数严格超过 $k/2$。
\end{definition}

算法的输入是训练集 $T$、查询点 $x$、近邻数 $k$、距离 $\rho$ 及并列处理规则，输出是 $\widehat f_T(x)$。具体执行时，计算或检索距离最近的 $k$ 个样本，统计这 $k$ 个样本在各类上的票数，再按式~\eqref{eq:lihang-knn-vote} 输出类别。$k=1$ 时得到\emph{最近邻法}（nearest neighbor）：把最近训练样本的类别直接赋给查询点。奇数 $k$ 可以排除二分类中的票数平局，但多分类中即使 $k$ 为奇数也仍可能并列，因此并列规则不能省略。

教材所说的“没有显式的学习过程”，是指基本近邻法不必先拟合一个有限维参数化判别函数；训练样本本身保留了决定预测所需的信息。这并不表示训练阶段完全没有工作：特征变换、$k$ 的选择和搜索索引的建立仍需要数据与计算。分类器随训练集而改变，其预测能力也不能由“记住了训练样本”直接推出。

\section{$k$ 近邻模型的三个基本要素}
\label{sec:lihang-knn-model}

\subsection{由训练样本确定的空间划分}
\label{subsec:lihang-knn-partition}

当训练集、距离、$k$、表决及并列规则都确定后，每个输入都有唯一的预测类别，因此分类器对应于特征空间的一个划分：预测为同一类的点组成一个决策区域。这里的区域一般不是线性子空间；同一类别的区域也可能不连通。近邻法的模型隐含在样本与规则共同形成的划分中，而不必写成一个简单的解析方程。

最近邻法的几何结构最直接。对每个样本点 $x_i$，先定义闭区域
\begin{equation}
  \widetilde V_i
  =\{x\in\mathcal X:\rho(x,x_i)\leq\rho(x,x_j),\ j=1,\ldots,N\}.
  \label{eq:lihang-knn-cell}
\end{equation}
对于互异的样本点和欧氏距离，这些区域是通常的\emph{Voronoi 单元}（Voronoi cell）。它们覆盖特征空间，但可能在等距边界上重叠；把边界按固定的距离并列规则归属后，才得到互不重叠的单元。最近邻法把第 $i$ 个单元中的输入标记为 $y_i$，同类样本的单元可合并为该类别的决策区域。若不同样本具有相同坐标，仍需按索引或另行约定的规则处理，不能仅凭几何距离断言最近邻唯一。

当 $k>1$ 时，先由最近的 $k$ 个样本组合确定局部邻域，再由标签表决决定类别；近邻集合发生变化的边界与最终类别发生变化的边界不一定相同。由此可见，\emph{距离度量、$k$ 值和分类决策规则}共同决定模型，改变其中任何一项都可能改变预测。

\subsection{距离度量}
\label{subsec:lihang-knn-distance}

距离把“相似”转化为可计算的比较，因此也是建模假设的一部分。设 $u=(u^{(1)},\ldots,u^{(d)})^\mathsf T$、$v=(v^{(1)},\ldots,v^{(d)})^\mathsf T\in\mathbb R^d$。对 $1\leq p<\infty$，\emph{$L_p$ 距离}（$L_p$ distance），又称闵可夫斯基距离（Minkowski distance），定义为
\begin{equation}
  L_p(u,v)=\left(\sum_{\ell=1}^{d}|u^{(\ell)}-v^{(\ell)}|^p\right)^{1/p}.
  \label{eq:lihang-knn-lp}
\end{equation}
常用的三个特例是\emph{曼哈顿距离}（Manhattan distance）、\emph{欧氏距离}（Euclidean distance）和 $L_\infty$ 距离：
\begin{align}
  L_1(u,v)&=\sum_{\ell=1}^{d}|u^{(\ell)}-v^{(\ell)}|,
  \label{eq:lihang-knn-l1}\\
  L_2(u,v)&=\left(\sum_{\ell=1}^{d}|u^{(\ell)}-v^{(\ell)}|^2\right)^{1/2},
  \label{eq:lihang-knn-l2}\\
  L_\infty(u,v)&=\max_{1\leq\ell\leq d}|u^{(\ell)}-v^{(\ell)}|.
  \label{eq:lihang-knn-linf}
\end{align}
若 $a=\max_\ell|u^{(\ell)}-v^{(\ell)}|$，则 $a\leq L_p(u,v)\leq d^{1/p}a$，令 $p\to\infty$ 即得到式~\eqref{eq:lihang-knn-linf}。二维空间中，以原点为中心的 $L_1$、$L_2$、$L_\infty$ 单位球分别为菱形、圆形和边平行于坐标轴的正方形，说明不同距离对各方向上的偏离赋予了不同含义。

条件 $p\geq1$ 不宜省略。例如 $0<p<1$ 时，令 $u=(1,0)^\mathsf T$、$v=(0,1)^\mathsf T$、$o=(0,0)^\mathsf T$，有 $L_p(u,v)=2^{1/p}>2=L_p(u,o)+L_p(o,v)$，因而式~\eqref{eq:lihang-knn-lp} 在此范围内不满足距离的三角不等式。另一方面，即使选定欧氏距离，特征的量纲和尺度仍会影响近邻次序；某一坐标的数值尺度很大时，可能主导距离。需要标准化时，应只用训练部分估计变换，再把同一变换用于验证样本与查询点。

\begin{example}
\label{ex:lihang-knn-distance-change}
教材例 3.1 给出 $x_1=(1,1)^\mathsf T$、$x_2=(5,1)^\mathsf T$、$x_3=(4,4)^\mathsf T$。以 $x_1$ 为查询点，在另外两个点中寻找 $L_p$ 距离下的最近邻，并说明结果如何随 $p$ 改变。
\end{example}
\begin{solve}
对于 $1\leq p<\infty$，
\[
  L_p(x_1,x_2)=4,\qquad L_p(x_1,x_3)=3\cdot2^{1/p}.
\]
设 $p_0=\log 2/\log(4/3)\approx2.4094$。比较两式可知：$1\leq p<p_0$ 时最近邻为 $x_2$；$p=p_0$ 时两点等距，按本章约定选择 $x_2$；$p>p_0$ 时最近邻为 $x_3$。特别地，$p=1,2$ 时选 $x_2$，整数 $p\geq3$ 时选 $x_3$；$p=\infty$ 时两距离分别为 $4$ 和 $3$，也选 $x_3$。这说明更换距离会改变局部邻域，因而可能改变同一查询点的预测类别。
\end{solve}

\subsection{$k$ 值的选择}
\label{subsec:lihang-knn-k}

$k$ 控制预测所使用的局部信息范围。较小的 $k$ 使决策更依赖查询点附近的少数样本，能够保留较细的局部结构，却也更容易受噪声或偶然采样波动影响；较大的 $k$ 把更多样本纳入表决，通常减弱局部波动，却可能把距离较远、类别分布不同的样本混在一起，抹平真实的分类边界。教材将其解释为\emph{近似误差}（approximation error）与\emph{估计误差}（estimation error）之间的权衡：局部规则过于粗糙可能造成系统性偏差，局部样本过少则使从有限数据得到的判断不稳定。这里描述的是常见机制，不能理解为任意固定数据集上的两类误差或测试误差都随 $k$ 严格单调变化。

两个极端有助于理解这种权衡。若训练样本坐标互异，$k=1$ 且预测训练样本时允许它自身参与搜索，则每个训练样本都以自身为最近邻，训练误分类率为零；这并不保证对新样本预测准确。若存在坐标相同而标签冲突的样本，上述零训练误差结论也不再自动成立。$k=N$ 时，所有查询点使用同一个训练集表决，预测退化为训练集中的最多数类；即使类别票数并列，固定的表决并列规则也使预测与 $x$ 无关。

实际选择 $k$ 时，可先给定候选集合，再用验证集或交叉验证比较各候选值的误分类率。交叉验证的每一折都只能从该折训练部分中寻找邻居，$k$ 不能超过该部分的样本数；标准化等预处理也应在该训练部分内重新估计。留一验证时应排除被验证样本自身。选定 $k$ 及其他超参数后，再用全部训练数据建立最终分类器，并用未参与选择的测试集评价。所谓“最优 $k$”在这里是给定候选集合与验证方案下的经验选择，不等同于已知总体分布下的理论最优值。

\subsection{多数表决与局部经验风险}
\label{subsec:lihang-knn-decision}

设 $(X,Y)$ 是新样本的随机输入与标签，取分类损失为 $0$--$1$ 损失 $\ell(c,y)=\mathbf 1\{y\ne c\}$。给定训练集 $T$ 后，分类器的总体风险为
\begin{equation}
  R(\widehat f_T)
  =\mathbb E\bigl[\ell(\widehat f_T(X),Y)\bigr]
  =\mathbb P\{Y\ne\widehat f_T(X)\},
  \label{eq:lihang-knn-population-risk}
\end{equation}
期望与概率针对新样本的分布；训练集在此视为固定。这个总体量通常未知。近邻表决直接使用的是查询点附近的样本频数，而不是式~\eqref{eq:lihang-knn-population-risk} 的精确值。

\begin{proposition}
\label{prop:lihang-knn-local-erm}
固定查询点 $x$ 及其近邻索引集 $N_k(x)$。若对这 $k$ 个样本统一预测为某一类别 $c\in\mathcal Y$，则多数表决最小化局部经验风险
\begin{equation}
  \widehat R_x(c)
  =\frac1k\sum_{i\in N_k(x)}\mathbf 1\{y_i\ne c\}.
  \label{eq:lihang-knn-local-risk}
\end{equation}
\end{proposition}
\begin{proof}
对任意 $c\in\mathcal Y$，指示函数恒等式给出
\[
  \widehat R_x(c)
  =1-\frac1k\sum_{i\in N_k(x)}\mathbf 1\{y_i=c\}.
\]
因此最小化 $\widehat R_x(c)$ 等价于最大化类别 $c$ 的票数；票数并列时，各并列类别都达到相同的最小经验风险，再按既定规则选择其一。
\end{proof}

也可以把 $\widehat\eta_j(x)=k^{-1}\sum_{i\in N_k(x)}\mathbf 1\{y_i=c_j\}$ 看作对局部类别比例的估计，表决即选择 $\widehat\eta_j(x)$ 最大的类别。\emph[caution]{局部经验风险最小不等于总体风险最小，也不等于在所有分类函数中最小化整个训练集的经验风险。}要把局部频数解释为对条件类别概率的良好估计，还需要样本与预测任务相匹配，以及所选距离能够反映有意义的局部性；上述代数推导本身没有证明这些统计性质。

\section{$k$ 近邻法的实现：kd 树}
\label{sec:lihang-knn-kdtree}

分类规则确定后，计算上的主要问题是如何找到最近的 $k$ 个样本。\emph{线性扫描}（linear scan）逐一计算查询点到全部样本的距离；对 $L_p$ 距离，计算这些距离需 $O(Nd)$ 次基本运算，再选出最小的 $k$ 个距离。直接全排序需 $O(N\log N)$ 的比较，也可以用大小为 $k$ 的候选堆或选择算法减少筛选开销。线性扫描是明确、可靠的实现基线，其是否足够快取决于数据规模与查询需求，不能笼统断言它不可用。

\emph{kd 树}（kd tree）是用轴对齐切分组织多维样本的二叉树。每个结点对应一个区域，并记录切分轴、切分位置及相关样本；子结点进一步划分该区域。它利用空间距离下界排除不可能包含更优近邻的子树，从而可能减少实际距离计算。名称中的 kd 沿用其多维空间树结构的习惯称呼；本章始终用 $d$ 表示维数、$k$ 表示邻居数，二者不可混淆。\emph{精确搜索时，更换搜索结构不改变分类器，只改变取得同一组近邻的计算过程。}

\subsection{构造平衡 kd 树}
\label{subsec:lihang-knn-kdtree-build}

从包含全部训练样本的根区域开始，每次选择一个坐标轴，用垂直于该轴的超平面切分当前区域。若根结点深度为 $h=0$，循环选择切分轴
\begin{equation}
  \ell(h)=1+(h\bmod d).
  \label{eq:lihang-knn-split-axis}
\end{equation}
于是二维情形依次使用第一、第二、第一坐标轴，三维情形则循环使用三条坐标轴。不同切分轴与切分位置会得到不同的树，但只要搜索与剪枝正确，精确近邻结果应保持一致。

为明确偶数个样本及重复坐标的处理，下面采用与教材例子相容的约定。当前非空索引集为 $M$，$m=|M|$；按第 $\ell$ 个坐标排序，取第 $\lfloor m/2\rfloor+1$ 个样本的该坐标作为切分值 $s$，即偶数个样本时取靠右的中间样本坐标，而不取两个坐标的平均数。把所有满足 $x_i^{(\ell)}=s$ 的样本存放在当前结点，其余样本分别进入左右子树：
\begin{equation}
  B=\{i\in M:x_i^{(\ell)}=s\},\quad
  M_L=\{i\in M:x_i^{(\ell)}<s\},\quad
  M_R=\{i\in M:x_i^{(\ell)}>s\}.
  \label{eq:lihang-knn-kdtree-split}
\end{equation}
$B$ 是当前结点的样本桶；坐标相等的样本都保留各自的索引和标签。

\textbf{构造算法。}输入为带索引的样本坐标、当前索引集 $M$ 和深度 $h$，初始取 $M=\{1,\ldots,N\}$、$h=0$；输出为根结点及递归生成的子树。
\begin{enumerate}
  \item 若 $M$ 为空，返回空树；否则按式~\eqref{eq:lihang-knn-split-axis} 选择切分轴，按上述中间次序统计量确定 $s$。
  \item 建立当前结点，记录切分轴 $\ell$、切分值 $s$ 及样本桶 $B$，并按式~\eqref{eq:lihang-knn-kdtree-split} 分出 $M_L$、$M_R$。
  \item 对 $M_L$、$M_R$ 分别以深度 $h+1$ 递归构造左右子树。若两者均为空，当前结点就是叶结点；完成后返回当前结点。
\end{enumerate}
每个非空结点至少保存一个样本，因此递归必定终止；同一样本恰好保存在一个结点。每个子集的大小不超过 $\lfloor m/2\rfloor$，故树高为 $O(\log N)$。若许多样本在所选坐标上相同，它们可能集中于一个较大的样本桶；树高较小并不意味着每个结点的处理工作都很少。

\begin{example}
\label{ex:lihang-knn-kdtree-build}
教材例 3.2 的二维样本坐标为
\[
  x_1=(2,3)^\mathsf T,\quad x_2=(5,4)^\mathsf T,\quad
  x_3=(9,6)^\mathsf T,\quad x_4=(4,7)^\mathsf T,\quad
  x_5=(8,1)^\mathsf T,\quad x_6=(7,2)^\mathsf T.
\]
按循环坐标轴及上述中间样本坐标约定构造 kd 树。
\end{example}
\begin{solve}
根结点按第一坐标排序，次序为 $2,4,5,7,8,9$，取第 $4$ 个坐标 $7$，故根保存 $x_6=(7,2)^\mathsf T$，切分直线为 $u^{(1)}=7$。左子集为 $\{x_1,x_2,x_4\}$，右子集为 $\{x_5,x_3\}$。左子集按第二坐标排序为 $3,4,7$，保存中间样本 $x_2=(5,4)^\mathsf T$，其左右子结点分别保存 $x_1$、$x_4$。右子集按第二坐标排序为 $1,6$，保存靠右的中间样本 $x_3=(9,6)^\mathsf T$，其左子结点保存 $x_5$，右子树为空。树的父子关系可写为
\[
  \begin{array}{c|c|c}
    \text{当前结点}&\text{左子结点}&\text{右子结点}\\ \hline
    (7,2)^\mathsf T&(5,4)^\mathsf T&(9,6)^\mathsf T\\
    (5,4)^\mathsf T&(2,3)^\mathsf T&(4,7)^\mathsf T\\
    (9,6)^\mathsf T&(8,1)^\mathsf T&\text{空}
  \end{array}
\]
通常统计学定义下，这六个第一坐标的中位数是 $(5+7)/2=6$；教材例子采用的切分值 $7$ 是某个实际样本的坐标。明确使用靠右中间样本，才能同时解释根结点的选择和每个结点保存样本的实现方式。两种中位数约定均可用于设计空间划分，但不能在同一算法描述中混用。
\end{solve}

若每个结点都重新排序当前子集，平衡递归给出的朴素构造时间为 $O(N\log^2 N)$；若用线性时间的次序统计量选择并线性划分，则固定维数下可做到 $O(N\log N)$。这些是具体实现对应的界，而不是仅凭“取中位数”就自动得到的同一个复杂度。保存全部样本及结点记录需要 $O(N)$ 存储；若另存各结点的 $d$ 维包围盒，还需计入相应的维度开销。

\subsection{最近邻搜索、回溯与剪枝}
\label{subsec:lihang-knn-kdtree-search}

沿切分方向下降，只能找到一个局部候选，不能保证真正最近的样本就在第一次进入的叶结点中。正确的搜索必须保留当前已访问样本中的最佳候选，并在回溯时判断其他子树是否仍可能包含更近的点。以下先用欧氏距离说明：若当前候选索引为 $i_\ast$，半径 $r=L_2(x,x_{i_\ast})$，则所有可能改进距离的样本都在以 $x$ 为中心、半径为 $r$ 的开球内；若还要按索引处理等距情形，也应检查球面上的候选。

设某结点用超平面 $u^{(\ell)}=s$ 切分。查询点到该平面的距离为
\begin{equation}
  \Delta=|x^{(\ell)}-s|.
  \label{eq:lihang-knn-plane-bound}
\end{equation}
对于位于查询点另一侧子区域中的任意样本 $u$，都有 $L_2(x,u)\geq|x^{(\ell)}-u^{(\ell)}|\geq\Delta$。因此 $\Delta>r$ 时可跳过另一侧子树；$\Delta\leq r$ 时仍可能存在更优候选，需要继续搜索，或再使用更紧的区域下界判断。条件 $\Delta\leq r$ 只是不能据此剪枝，并不保证其中一定存在更近样本。同一坐标下界对本章的其他 $L_p$ 距离也成立。

\textbf{精确最近邻搜索算法。}输入为非空训练集构造的 kd 树和查询点 $x$，输出为使 $(L_2(x,x_i),i)$ 按字典序最小的索引，即距离最小、等距时索引最小的样本。初始化时无候选，令 $r=+\infty$；从根结点调用下述递归过程。
\begin{enumerate}
  \item 若当前子树为空，立即返回。否则计算当前结点样本桶中各样本到 $x$ 的距离；若距离更小，或距离相同而索引更小，就更新 $i_\ast$ 和 $r$。
  \item 若 $x^{(\ell)}<s$，先递归搜索左子树，否则先递归搜索右子树。把先搜索的一侧称为近侧，另一侧称为远侧；近侧为空时跳过该递归调用，继续后续判断。
  \item 近侧搜索结束后，使用更新后的 $r$ 与式~\eqref{eq:lihang-knn-plane-bound} 比较。若 $\Delta\leq r$，再递归搜索远侧；若 $\Delta>r$，剪去远侧。完成后返回父调用；根调用全部结束时，输出 $i_\ast$ 及其坐标。
\end{enumerate}
这种写法在下降时检查结点样本，在递归返回时完成回溯，与先到达叶结点再向上检查的思路一致，并能直接处理空子树和一个结点保存多个样本的情形。\emph[caution]{某一子树被剪去只表示不再访问这棵子树，仍须返回上层检查其他尚未处理的分支；不能在第一次遇到不相交区域时结束整个搜索。}

\begin{proposition}
\label{prop:lihang-knn-kdtree-correct}
对有限非空训练集，按上述构造与搜索算法执行，最终返回欧氏距离最近、等距时索引最小的训练样本。
\end{proposition}
\begin{proof}
样本桶及左右子树使训练样本被不重不漏地保存。搜索始终把已访问样本中按距离与索引排序的最优者作为当前候选，故 $r$ 在获得第一个候选后只会减小或保持不变。若某子树被剪去，则其中每个样本的距离都大于剪枝当时的 $r$，由于半径不增，也严格大于最终半径，因而不能成为更优者。距离下界等于 $r$ 时没有剪枝，等距且索引更小的候选不会因此遗漏。有限树的递归结束后，每个样本或者已检查，或者处于已证明不能改进结果的子树中，因此当前候选就是全体样本中的最优者。
\end{proof}

若已知某子树样本均包含于轴对齐闭包围盒 $R=\prod_{j=1}^{d}[a_j,b_j]$，还可使用更紧的下界。令
\[
  \delta_j=\max\{a_j-x^{(j)},\ 0,\ x^{(j)}-b_j\},
\]
其中 $\delta_j$ 是查询点第 $j$ 个坐标到区间 $[a_j,b_j]$ 的距离，则
\begin{equation}
  \operatorname{dist}_2(x,R)
  =\inf_{u\in R}L_2(x,u)
  =\left(\sum_{j=1}^{d}\delta_j^2\right)^{1/2}.
  \label{eq:lihang-knn-box-bound}
\end{equation}
等式成立是因为平方距离按坐标相加，每个坐标都可独立投影到对应区间。若这个下界大于 $r$，整棵子树即可剪去；否则继续检查。对 $L_p$ 距离相应使用 $(\sum_j\delta_j^p)^{1/p}$，对 $L_\infty$ 距离使用 $\max_j\delta_j$。剪枝下界必须与实际选用的距离一致。

\begin{example}
\label{ex:lihang-knn-kdtree-query}
在例~\ref{ex:lihang-knn-kdtree-build} 的 kd 树中，按欧氏距离寻找 $x=(3,4.5)^\mathsf T$ 的最近邻。这也是教材习题 3.2 的查询点。
\end{example}
\begin{solve}
先检查根样本 $(7,2)^\mathsf T$，平方距离为 $89/4$。查询点第一坐标小于 $7$，进入左子树；检查 $(5,4)^\mathsf T$ 后，当前最小平方距离变为 $17/4$。查询点第二坐标大于 $4$，先到 $(4,7)^\mathsf T$，其平方距离为 $29/4$，不能改进候选。

回到以第二坐标 $4$ 切分的结点，查询点到切分线的距离为 $|4.5-4|=1/2$，小于当前半径 $\sqrt{17}/2$，故还需检查另一侧的 $(2,3)^\mathsf T$。该点的平方距离为
\[
  (3-2)^2+(4.5-3)^2=\frac{13}{4},
\]
于是更新当前候选为 $(2,3)^\mathsf T$，半径为 $\sqrt{13}/2$。再回到根结点，查询点到切分线 $u^{(1)}=7$ 的距离为 $4>\sqrt{13}/2$，因此右子树可整体剪去。最终最近邻为 $(2,3)^\mathsf T$，距离为 $\sqrt{13}/2$。第一次进入的叶结点保存 $(4,7)^\mathsf T$，而最终结果来自另一分支，说明回溯不可省略。
\end{solve}

从最近邻推广到 $k$ 近邻时，用容量为 $k$ 的\emph{最大堆}（max-heap）保存目前最好的候选，按 $(L_2(x,x_i),i)$ 排序，堆顶是当前候选中最差者。每检查一个样本，若堆未满就加入；若已满且新样本优于堆顶，就替换堆顶。堆未满时令剪枝半径 $r=+\infty$，堆满后令 $r$ 为堆顶距离。近侧优先访问、远侧下界比较和递归返回的过程均沿用上述算法；远侧下界大于 $r$ 时才能剪枝，等于 $r$ 时仍需保留索引并列比较。根调用结束后，堆内即为所需的 $k$ 个近邻，按距离及索引排序即可输出 $N_k(x)$，再交给式~\eqref{eq:lihang-knn-vote} 分类。这给出了教材习题 3.3 所要求的完整推广，$k=1$ 时退化为前述搜索。

最大堆搜索的正确性来自同一不变式：堆中始终保留已访问样本中最优的至多 $k$ 个；堆满后的半径不会增大，剪去的样本也不能进入最终前 $k$ 名。平衡 kd 树的树高虽为 $O(\log N)$，精确查询却可能回溯访问许多分支；在固定低维和有利数据分布下通常能够明显减少访问，但不能仅由“平衡”或“随机分布”断言所有查询都有 $O(\log N)$ 的复杂度。最近邻搜索最坏仍需检查全部样本，距离计算为 $O(Nd)$；维护 $k\geq2$ 个候选时，最坏另有 $O(N\log k)$ 的堆操作。维数较高或空间划分不利时，剪枝可能失效，性能接近线性扫描，因此应结合维数、分布和查询需求选择实现方式。

\section{学习要点与习题研讨}
\label{sec:lihang-knn-review}

本章的逻辑链是：距离确定邻域，$k$ 确定邻域规模，决策规则把局部标签转化为预测；三者与训练集共同形成分类模型。多数表决的局部经验风险解释见命题~\ref{prop:lihang-knn-local-erm}，它不替代对泛化误差的评价。kd 树的空间切分则服务于检索，与最近邻模型的类别区域是两种不同的划分：前者按坐标组织计算，后者由距离和标签决定预测。核查一次 kd 树搜索时，应逐一说明当前候选、半径、待查子树及剪枝下界，而不能只报告下降路径。

\begin{example}
\label{ex:lihang-knn-partition-k}
按教材习题 3.1 的要求自选二维样本，比较 $k=1$ 与 $k=2$ 的决策区域。取 $x_1=(0,0)^\mathsf T$、$y_1=c_1$，$x_2=(2,0)^\mathsf T$、$y_2=c_2$，$x_3=(0,2)^\mathsf T$、$y_3=c_2$，使用欧氏距离，并采用定义~\ref{def:lihang-knn-rule} 的距离和类别并列规则。
\end{example}
\begin{solve}
写查询点为 $x=(u,v)^\mathsf T$，三条平方距离为
\[
  D_1=u^2+v^2,\qquad
  D_2=(u-2)^2+v^2,\qquad
  D_3=u^2+(v-2)^2.
\]
$k=1$ 时，仅 $x_1$ 的类别是 $c_1$，所以预测为 $c_1$ 当且仅当 $D_1\leq D_2$ 且 $D_1\leq D_3$，即 $u\leq1$ 且 $v\leq1$；其余点预测为 $c_2$。等距时索引 $1$ 优先，因此这两条边界在满足另一条件时也归入 $c_1$ 区域。

$k=2$ 时，只要 $x_1$ 位于两个最近邻中，两类便各得一票，按类别索引优先预测为 $c_1$。只有两个近邻恰为 $x_2,x_3$ 时才预测为 $c_2$，这要求 $D_1>D_2$ 且 $D_1>D_3$，即 $u>1$ 且 $v>1$。因此 $k=2$ 时 $c_2$ 区域是这个开象限，其余区域均预测为 $c_1$。若改成票数并列时优先 $c_2$，则 $k=2$ 的预测将处处为 $c_2$，可见并列规则也是模型定义的一部分。

两种 $k$ 的边界不同，但仅凭这三个训练样本及几何划分，无法断定哪种规则对未知分布的新样本更准确；还需独立验证或额外的分布假设。教材习题 3.2 的数值搜索已在例~\ref{ex:lihang-knn-kdtree-query} 完成，习题 3.3 的多近邻搜索则由第~\ref{subsec:lihang-knn-kdtree-search} 小节的候选最大堆给出。
\end{solve}


\fi

\selectchapter{4}
\ifcompilechapter
\chapter{朴素贝叶斯法}
\label{chap:lihang-naive-bayes}

% 研读来源：机器学习/textbooks/机器学习（分章原文）/机器学习-4.pdf，教材第 59--66 页。
% 下文沿用教材的两节、五个小节；推导、适用条件与例题后验概率为学习补充。

分类可以直接学习输入到类别的决策规则，也可以先描述不同类别如何产生特征，再由观测到的特征反推类别。\emph{朴素贝叶斯法}（naive Bayes method）采用后一条路线：学习类别的先验概率和各类中的特征分布，利用贝叶斯定理计算后验概率，最后作出分类。困难在于，高维特征的联合分布需要大量参数；本方法通过\emph{特征条件独立假设}降低估计难度。本章先说明这一模型及其决策依据，再讨论离散特征的极大似然估计、学习与分类算法，以及用于处理零频数的贝叶斯估计。

\section{朴素贝叶斯法的学习与分类}
\label{sec:lihang-nb-learning-classification}

\subsection{基本方法}
\label{subsec:lihang-nb-basic-method}

设类别随机变量 $Y$ 取值于 $\mathcal{Y}=\{c_1,\ldots,c_K\}$，其中 $K\geq2$；特征随机向量为 $X=(X^{(1)},\ldots,X^{(n)})^{\mathsf T}$，其中 $n\geq1$。本章具体估计公式针对有限离散特征：第 $j$ 个特征的取值集合为 $\mathcal{A}_j=\{a_{j1},\ldots,a_{jS_j}\}$，输入空间为 $\mathcal{X}=\mathcal{A}_1\times\cdots\times\mathcal{A}_n$。类别和特征值的集合预先给定；数值形式的类别名称与特征编码不表示大小或距离。训练集
\[
    T=\{(x_i,y_i)\}_{i=1}^{N},\qquad
    x_i=(x_i^{(1)},\ldots,x_i^{(n)})^{\mathsf T},\quad N\geq1,
\]
由未知联合分布 $P$ 独立同分布抽样得到。本章的 $\log$ 表示自然对数。$X,Y$ 表示随机对象，$x_i,y_i$ 表示已经观测到的取值，$i$ 是样本索引，$j$ 是特征索引，$k$ 是类别索引，$l$ 是某一特征的取值索引。

记先验概率（prior probability）为 $\pi_k=P(Y=c_k)$，类条件概率（class-conditional probability）为 $\theta_{jlk}=P(X^{(j)}=a_{jl}\mid Y=c_k)$。这些条件概率在 $\pi_k>0$ 时由联合分布确定；对先验概率为零的类别，可任意选取规范化的条件分布而不改变联合分布。在不加结构限制时，每个类别中的完整特征向量有 $\prod_{j=1}^{n}S_j$ 种可能取值，需要同样数量的概率表项。扣除概率和为 $1$ 的约束，所有类条件分布共有 $K(\prod_{j=1}^{n}S_j-1)$ 个自由参数，另有 $K-1$ 个先验自由参数。当各 $S_j=S\geq2$ 时，联合表的规模随 $n$ 按 $S^n$ 增长，有限样本很难为每种组合提供足够观测。

\begin{definition}
\label{def:lihang-nb-conditional-independence}
若对每个 $P(Y=c_k)>0$ 的类别及所有 $x\in\mathcal{X}$，都有
\begin{equation}
\label{eq:lihang-nb-conditional-independence}
    P(X=x\mid Y=c_k)
    =\prod_{j=1}^{n}P(X^{(j)}=x^{(j)}\mid Y=c_k),
\end{equation}
则称各特征在给定类别后\emph{相互条件独立}（mutually conditionally independent）。以此分解形式作为模型假设的生成分类方法称为朴素贝叶斯法。
\end{definition}

这里要求给定类别后所有特征的联合分布可以分解，特征数大于 $2$ 时仅有两两条件独立一般不够；也不要求特征在未给定类别时相互独立。例如，设 $P(Y=c_1)=P(Y=c_2)=1/2$，两个二元特征在每个类别内独立，且
\[
    P(X^{(j)}=1\mid Y=c_1)=0.9,\qquad
    P(X^{(j)}=1\mid Y=c_2)=0.1,\qquad j=1,2.
\]
此时 $P(X^{(1)}=1)=P(X^{(2)}=1)=0.5$，但
\[
    P(X^{(1)}=1,X^{(2)}=1)
    =\tfrac12(0.9)^2+\tfrac12(0.1)^2=0.41\neq0.25.
\]
类别混合可以造成边缘相关，因此不能用原始样本中两个特征相关这一事实直接否定条件独立。另一方面，实际数据未必满足式~\eqref{eq:lihang-nb-conditional-independence}，使用朴素贝叶斯模型时，这一分解是一项建模假设，不是由贝叶斯定理推出的结论。

令 $\Theta=(\pi,\theta)$ 表示模型全部参数。对 $x^{(j)}=a_{jl}$，约定 $\theta_{jk}(x^{(j)})=\theta_{jlk}$，于是模型联合概率为
\begin{equation}
\label{eq:lihang-nb-joint-model}
    P_\Theta(X=x,Y=c_k)=\pi_k\prod_{j=1}^{n}\theta_{jk}(x^{(j)}),\qquad
    \pi_k\geq0,\quad\sum_{k=1}^{K}\pi_k=1,\quad
    \theta_{jlk}\geq0,\quad\sum_{l=1}^{S_j}\theta_{jlk}=1.
\end{equation}
它确实定义了概率分布，因为对 $x$ 求和时，各特征的有限和可以分开，得到 $\sum_x P_\Theta(X=x,Y=c_k)=\pi_k$，再对 $k$ 求和即为 $1$。模型只需 $K\sum_j S_j$ 个条件概率表项，连同先验共有 $(K-1)+K\sum_j(S_j-1)$ 个自由参数；特征取值数有统一上界时，规模随特征数线性增长。它先建模联合分布再导出分类规则，因此属于\emph{生成模型}（generative model）。条件独立带来估计与计算的简化，也限制了模型能够表达的类内依赖关系。

对满足 $P_\Theta(X=x)>0$ 的输入，贝叶斯定理给出后验概率（posterior probability）
\begin{equation}
\label{eq:lihang-nb-posterior}
    P_\Theta(Y=c_k\mid X=x)
    =\frac{\pi_k\prod_{j=1}^{n}\theta_{jk}(x^{(j)})}
    {\displaystyle\sum_{r=1}^{K}\pi_r\prod_{j=1}^{n}\theta_{jr}(x^{(j)})}.
\end{equation}
分母使用独立的求和索引 $r$，且对待比较的类别 $k$ 相同。记各类的未归一化得分为
\begin{equation}
\label{eq:lihang-nb-score-rule}
    s_k(x)=\pi_k\prod_{j=1}^{n}\theta_{jk}(x^{(j)}),\qquad
    k^*(x)\in\arg\max_{1\leq k\leq K}s_k(x),\qquad f_\Theta(x)=c_{k^*(x)}.
\end{equation}
若有并列最大值，固定选取索引最小的类别，使分类器成为单值函数。\emph[caution]{得分是联合概率 $P_\Theta(X=x,Y=c_k)$，不是已经归一化的后验概率}；分类只需比较得分，报告后验概率则必须除以 $\sum_r s_r(x)$。若所有得分都为零，模型认为该输入不可能出现，式~\eqref{eq:lihang-nb-posterior} 的分母为零，不能把任意的并列选择解释为有后验依据的分类。

\subsection{后验概率最大化的含义}
\label{subsec:lihang-nb-posterior-risk}

为何选择后验概率最大的类别，需要由损失函数说明。设 $f:\mathcal{X}\to\mathcal{Y}$ 为分类决策函数，采用\emph{零一损失}（zero-one loss）
\[
    L(y,a)=\mathbf{1}\{y\neq a\}
    =\begin{cases}0,&y=a,\\1,&y\neq a,\end{cases}
\]
其中 $a$ 是预测类别，$\mathbf{1}\{A\}$ 是事件或命题 $A$ 的指示函数。相应的\emph{期望风险}（expected risk）为 $R_P(f)=\mathbb{E}_{(X,Y)\sim P}[L(Y,f(X))]=P(Y\neq f(X))$，即总体误分类概率。

\begin{theorem}
\label{thm:lihang-nb-bayes-decision}
对有限输入空间和类别空间上的任意联合分布 $P$，在所有 $P(X=x)>0$ 的输入处选取后验概率最大的类别，可以使零一损失下的期望风险在所有分类函数中达到最小。该结论不需要特征条件独立假设。
\end{theorem}

\begin{proof}
对固定输入 $x$ 和候选预测类别 $a$，条件风险为
\begin{align*}
    r_P(a\mid x)
    &=\mathbb{E}_P[L(Y,a)\mid X=x]\\
    &=\sum_{k=1}^{K}\mathbf{1}\{c_k\neq a\}P(Y=c_k\mid X=x)\\
    &=1-P(Y=a\mid X=x).
\end{align*}
因此，最小化条件风险等价于最大化 $P(Y=a\mid X=x)$。由全期望公式，
\[
    R_P(f)=\mathbb{E}_X[r_P(f(X)\mid X)].
\]
令 $f^*(x)$ 在每个正概率输入处选取条件风险的最小值，则对任意 $f$，都有 $r_P(f^*(x)\mid x)\leq r_P(f(x)\mid x)$，乘以 $P(X=x)$ 并求和得到 $R_P(f^*)\leq R_P(f)$。零概率输入处的取值不影响风险。最小风险为
\[
    R_P^*=1-\mathbb{E}_X\!\left[\max_{1\leq k\leq K}P(Y=c_k\mid X)\right].\qedhere
\]
\end{proof}

按真实后验概率分类所得的规则称为\emph{贝叶斯分类器}（Bayes classifier）。实际训练只能得到参数估计 $\widehat\Theta$，进而得到估计分布 $\widehat P=P_{\widehat\Theta}$。将定理~\ref{thm:lihang-nb-bayes-decision} 应用于 $\widehat P$，说明朴素贝叶斯决策在\emph{估计模型下}最小化零一风险；若真实分布不满足条件独立假设，或参数存在估计误差，便不能据此断言它已达到真实分布 $P$ 下的最小风险。对错误类别赋予不同代价时，应改为最小化 $\sum_k L(c_k,a)\widehat P(Y=c_k\mid X=x)$，一般不能继续仅比较后验概率大小。

\section{朴素贝叶斯法的参数估计}
\label{sec:lihang-nb-parameter-estimation}

\subsection{极大似然估计}
\label{subsec:lihang-nb-mle}

训练样本已包含类别标签，因而可以直接按类别及特征值计数。定义
\begin{equation}
\label{eq:lihang-nb-counts}
    N_k=\sum_{i=1}^{N}\mathbf{1}\{y_i=c_k\},\qquad
    N_{jlk}=\sum_{i=1}^{N}\mathbf{1}\{x_i^{(j)}=a_{jl},\ y_i=c_k\}.
\end{equation}
其中 $N_k$ 是第 $k$ 类的样本数，$N_{jlk}$ 是该类中第 $j$ 个特征取值为 $a_{jl}$ 的样本数。由于每个样本属于且仅属于一个类别、每个特征具有且仅具有一个取值，有 $\sum_k N_k=N$ 和 $\sum_l N_{jlk}=N_k$。训练数据完整观测是这些计数恒等式的前提。

\begin{proposition}
\label{prop:lihang-nb-mle}
在式~\eqref{eq:lihang-nb-joint-model} 的有限离散模型中，$N\geq1$ 时先验概率的\emph{极大似然估计}（maximum likelihood estimation，MLE）为
\begin{equation}
\label{eq:lihang-nb-mle-prior}
    \widehat\pi_k=\frac{N_k}{N},\qquad k=1,\ldots,K.
\end{equation}
对 $N_k>0$ 的类别，条件概率的极大似然估计为
\begin{equation}
\label{eq:lihang-nb-mle-conditional}
    \widehat\theta_{jlk}=\frac{N_{jlk}}{N_k},\qquad
    j=1,\ldots,n,\quad l=1,\ldots,S_j.
\end{equation}
若 $N_k=0$，该类的条件分布不能由似然唯一确定，式~\eqref{eq:lihang-nb-mle-conditional} 不适用；任取满足概率约束的条件分布，配合 $\widehat\pi_k=0$，均给出同样的最大似然值。
\end{proposition}

\begin{proof}
给定独立同分布的已标注样本，参数似然可以按计数整理为
\begin{equation}
\label{eq:lihang-nb-likelihood}
    \mathcal{L}(\Theta;T)
    =\prod_{i=1}^{N}\left[\pi_{k_i}\prod_{j=1}^{n}\theta_{jk_i}(x_i^{(j)})\right]
    =\prod_{k=1}^{K}\pi_k^{N_k}
      \prod_{j=1}^{n}\prod_{k=1}^{K}\prod_{l=1}^{S_j}\theta_{jlk}^{N_{jlk}},
\end{equation}
其中 $k_i$ 是满足 $y_i=c_{k_i}$ 的索引，约定零次幂为 $1$。取对数得到
\[
    \ell(\Theta;T)=\sum_{k=1}^{K}N_k\log\pi_k
    +\sum_{j=1}^{n}\sum_{k=1}^{K}\sum_{l=1}^{S_j}N_{jlk}\log\theta_{jlk},
\]
其中 $0\log0$ 按连续延拓记为 $0$，正计数对应的参数若为零则对数似然为 $-\infty$。先验向量和每个条件概率向量分别受到和为 $1$ 的约束，因此可以分别求解。

一般地，设计数 $m_1,\ldots,m_q\geq0$，$m=\sum_r m_r>0$，在概率单纯形 $p_r\geq0,\sum_r p_r=1$ 上最大化 $\sum_r m_r\log p_r$。令 $\bar p_r=m_r/m$，并记 $A=\{r:m_r>0\}$。若某个 $r\in A$ 对应 $p_r=0$，该目标不可能达到最大；否则由 $\log t\leq t-1$，
\[
    \sum_{r\in A}m_r\log\frac{p_r}{\bar p_r}
    \leq m\sum_{r\in A}\bar p_r\!\left(\frac{p_r}{\bar p_r}-1\right)
    =m\left(\sum_{r\in A}p_r-1\right)\leq0.
\]
取 $p_r=\bar p_r$ 达到等号；等号要求所有正计数项均取此值，零计数项均取零。因此归一化计数给出唯一最大值点，也覆盖最优解在边界上的情形。将该结论分别用于 $m_r=N_k$ 和 $m_r=N_{jlk}$，即得两式。若 $N_k=0$，则相应的所有 $N_{jlk}=0$，似然中与该类条件参数有关的因子恒为 $1$，故这些参数没有唯一估计。
\end{proof}

极大似然训练的目标是联合样本似然最大化，并非直接最小化训练分类错误数。对 $N_k>0$ 的类别，未观察到的特征值给出 $\widehat\theta_{jlk}=0$；这符合极大似然解的边界性质，却可能把一个本来可能出现的特征值视为不可能。若新输入在某一类下有任意一个零条件概率，该类的整个乘积得分便为零。这是有限样本计数与乘积结构共同造成的问题，不能从零频数直接推断真实概率为零。

\subsection{学习与分类算法}
\label{subsec:lihang-nb-algorithm}

训练与预测可以分为两个阶段。输入为训练集 $T$、已知的类别及特征取值集合，以及待分类的 $x\in\mathcal{X}$；输出为预测类别，必要时同时输出估计后验概率。教材的基本算法采用极大似然估计，要求对需要估计条件概率的类别有 $N_k>0$。完整步骤如下。
\begin{enumerate}
    \item 将所有 $N_k,N_{jlk}$ 初始化为 $0$，逐个读取样本；对其所属类别累加 $N_k$，并对每个特征的对应取值累加 $N_{jlk}$。按式~\eqref{eq:lihang-nb-mle-prior}、\eqref{eq:lihang-nb-mle-conditional} 归一化计数，得到 $\widehat\pi_k,\widehat\theta_{jlk}$；零样本类别的条件分布作任意规范化约定，其分类得分为零。
    \item 对给定 $x$，以估计参数代入式~\eqref{eq:lihang-nb-score-rule}，计算 $\widehat s_k(x)=\widehat\pi_k\prod_j\widehat\theta_{jk}(x^{(j)})$。若 $\sum_k\widehat s_k(x)>0$，选择得分最大的类别，并按固定规则处理并列；后验概率为 $\widehat s_k(x)/\sum_r\widehat s_r(x)$。若所有得分为零，则基本模型无法提供该输入的后验分类，需采用下节的正平滑等处理。
\end{enumerate}
这是一遍计数加有限次比较的闭式算法，没有迭代收敛或学习率的要求。在特征值可用常数时间索引的前提下，计数耗时 $O(Nn)$，完整概率表初始化与归一化耗时 $O(K\sum_j S_j)$，模型存储量为 $O(K+K\sum_j S_j)$，单个输入的预测耗时 $O(Kn)$。

实现时常用对数得分
\begin{equation}
\label{eq:lihang-nb-log-score}
    \widehat g_k(x)=\log\widehat\pi_k+
    \sum_{j=1}^{n}\log\widehat\theta_{jk}(x^{(j)}).
\end{equation}
约定 $\log0=-\infty$；至少一个得分为正时，对数的单调性保证最大值类别不变，可避免多个小概率相乘造成浮点下溢。若还需要后验概率，令 $m=\max_k\widehat g_k(x)$，则
\[
    \widehat P(Y=c_k\mid X=x)
    =\frac{\exp(\widehat g_k(x)-m)}{\sum_{r=1}^{K}\exp(\widehat g_r(x)-m)}.
\]
这一写法要求 $m$ 有限；所有类均为 $-\infty$ 时，数值变换不能修复模型的零分母问题。

\begin{example}
\label{ex:lihang-nb-mle-classification}
教材例题的训练数据见表~\ref{tab:lihang-nb-training}。设 $\mathcal{A}_1=\{1,2,3\}$、$\mathcal{A}_2=\{S,M,L\}$，类别集合为 $\{1,-1\}$。用极大似然估计训练朴素贝叶斯分类器，并预测 $x=(2,S)^{\mathsf T}$ 的类别。
\end{example}

\begin{table}[htbp]
\centering
\caption{朴素贝叶斯分类的训练数据（据教材本章训练数据表整理）}
\label{tab:lihang-nb-training}
\begin{tabular}{c|ccccccccccccccc}
\toprule
样本 $i$ & 1 & 2 & 3 & 4 & 5 & 6 & 7 & 8 & 9 & 10 & 11 & 12 & 13 & 14 & 15 \\
\midrule
$x_i^{(1)}$ & 1 & 1 & 1 & 1 & 1 & 2 & 2 & 2 & 2 & 2 & 3 & 3 & 3 & 3 & 3 \\
$x_i^{(2)}$ & $S$ & $M$ & $M$ & $S$ & $S$ & $S$ & $M$ & $M$ & $L$ & $L$ & $L$ & $M$ & $M$ & $L$ & $L$ \\
$y_i$ & $-1$ & $-1$ & 1 & 1 & $-1$ & $-1$ & $-1$ & 1 & 1 & 1 & 1 & 1 & 1 & 1 & $-1$ \\
\bottomrule
\end{tabular}
\end{table}

\begin{solve}
为避免反复展开同类计算，以类别值 $+$、$-$ 分别作下标，表示 $Y=1$、$Y=-1$。计数为 $N_+=9,N_-=6$；按取值顺序 $(1,2,3)$ 与 $(S,M,L)$，两个特征在各类中的计数向量分别为
\[
    (N_{1l,+})_{l=1}^{3}=(2,3,4),\qquad
    (N_{1l,-})_{l=1}^{3}=(3,2,1),\qquad
    (N_{2l,+})_{l=1}^{3}=(1,4,4),\qquad
    (N_{2l,-})_{l=1}^{3}=(3,2,1).
\]
因此 $\widehat\pi_+=9/15,\widehat\pi_-=6/15$，四个条件概率向量依次为 $(2,3,4)/9$、$(3,2,1)/6$、$(1,4,4)/9$、$(3,2,1)/6$。对于输入 $(2,S)^{\mathsf T}$，只需取相应的两个条件概率：
\begin{align*}
    \widehat s_+(x)&=\frac{9}{15}\cdot\frac{3}{9}\cdot\frac{1}{9}=\frac{1}{45},\\
    \widehat s_-(x)&=\frac{6}{15}\cdot\frac{2}{6}\cdot\frac{3}{6}=\frac{1}{15}.
\end{align*}
负类得分较大，故预测为 $-1$。将得分归一化，模型的后验概率为
\[
    \widehat P(Y=1\mid X=x)=\frac14,\qquad
    \widehat P(Y=-1\mid X=x)=\frac34.
\]
这些概率依赖训练计数及条件独立模型，不能解释为已经核实的真实条件概率。
\end{solve}

\subsection{贝叶斯估计}
\label{subsec:lihang-nb-bayesian-estimation}

为缓解零频数问题，可以对每一个预先规定的可能取值加入同样的正\emph{伪计数}（pseudocount）$\lambda>0$，再进行归一化。教材给出的贝叶斯估计为
\begin{equation}
\label{eq:lihang-nb-smoothed-probabilities}
    \widehat\pi_k^{(\lambda)}=\frac{N_k+\lambda}{N+K\lambda},\qquad
    \widehat\theta_{jlk}^{(\lambda)}=\frac{N_{jlk}+\lambda}{N_k+S_j\lambda}.
\end{equation}
每个分母增加的量等于该概率向量所有伪计数的总和，因此
\[
    \sum_k\widehat\pi_k^{(\lambda)}=1,\qquad
    \sum_l\widehat\theta_{jlk}^{(\lambda)}=1,
\]
且各分量严格为正。即使 $N_k=0$，条件估计也有定义，并等于 $1/S_j$。当 $\lambda=1$ 时称为\emph{拉普拉斯平滑}（Laplace smoothing）；当 $\lambda\downarrow0$ 时，先验估计趋于极大似然估计，$N_k>0$ 类别的条件估计也趋于其极大似然估计。对 $N_k=0$ 的类别，不能在条件公式中直接令 $\lambda=0$。

平滑的作用也可由下式看出：对 $N_k>0$，
\begin{equation}
\label{eq:lihang-nb-smoothing-mixture}
    \widehat\theta_{jlk}^{(\lambda)}
    =\frac{N_k}{N_k+S_j\lambda}\frac{N_{jlk}}{N_k}
     +\frac{S_j\lambda}{N_k+S_j\lambda}\frac{1}{S_j}.
\end{equation}
估计是经验频率与均匀分布的加权平均；$\lambda$ 越大，或该类样本越少，均匀分布所占权重越大。固定 $\lambda>0$ 且某类样本数增长时，这一权重趋于零；这说明伪计数的相对作用减小，并不单独构成对真实分类风险的保证。\emph[caution]{平滑保证预定有限取值集合中的概率为正，不会自动为集合外的新取值建立概率}；应用中若要使用“未知值”类别，应在确定取值集合时就将它纳入，并在训练及预测时使用相同编码。

\begin{proposition}
\label{prop:lihang-nb-dirichlet-estimation}
将先验概率向量 $\pi$ 及每个条件概率向量 $\theta_{jk}=(\theta_{j1k},\ldots,\theta_{jS_jk})$ 视为未知随机参数，分别赋予相互独立的对称 Dirichlet（狄利克雷）先验
\[
    \pi\sim\operatorname{Dir}_K(\lambda,\ldots,\lambda),\qquad
    \theta_{jk}\sim\operatorname{Dir}_{S_j}(\lambda,\ldots,\lambda),\qquad\lambda>0.
\]
则在完整观测的训练集 $T$ 给定后，各参数向量的后验均值等于式~\eqref{eq:lihang-nb-smoothed-probabilities}，它们也是逐个参数分量在平方损失下的贝叶斯估计。
\end{proposition}

\begin{proof}
先考虑 $q\geq2$ 类的概率向量 $p=(p_1,\ldots,p_q)$。对参数 $\alpha_r>0$，Dirichlet 分布在概率单纯形上的密度为
\[
    h(p)=\frac{1}{B(\alpha)}\prod_{r=1}^{q}p_r^{\alpha_r-1},\qquad
    B(\alpha)=\frac{\prod_{r=1}^{q}\Gamma(\alpha_r)}{\Gamma(\sum_{r=1}^{q}\alpha_r)}.
\]
这里 $\Gamma$ 为伽马函数。若观测的分类计数为 $m_1,\ldots,m_q$，似然关于 $p$ 正比于 $\prod_r p_r^{m_r}$，乘上先验密度可得后验分布 $\operatorname{Dir}_q(\alpha_1+m_1,\ldots,\alpha_q+m_q)$，这就是\emph{共轭先验}（conjugate prior）的性质。

令 $\beta_r=\alpha_r+m_r$，$\beta_0=\sum_r\beta_r$，$e_r$ 为第 $r$ 个坐标为 $1$、其余为 $0$ 的向量。由归一化积分及 $\Gamma(t+1)=t\Gamma(t)$，
\[
    \mathbb{E}[p_r\mid T]
    =\frac{B(\beta+e_r)}{B(\beta)}=\frac{\beta_r}{\beta_0}
    =\frac{m_r+\alpha_r}{\sum_s m_s+\sum_s\alpha_s}.
\]
由于式~\eqref{eq:lihang-nb-likelihood} 按参数向量分解，且这些向量的先验独立，它们的后验也分别分解为
\[
    \pi\mid T\sim\operatorname{Dir}_K(N_1+\lambda,\ldots,N_K+\lambda),\qquad
    \theta_{jk}\mid T\sim\operatorname{Dir}_{S_j}(N_{j1k}+\lambda,\ldots,N_{jS_jk}+\lambda).
\]
代入后验均值公式即得所需估计。若某个 $S_j=1$，该条件概率恒为 $1$，不需要非退化的密度描述，式~\eqref{eq:lihang-nb-smoothed-probabilities} 仍成立。对任意参数分量 $p_r$ 及点估计 $d_r$，后验平方损失满足
\[
    \mathbb{E}[(p_r-d_r)^2\mid T]
    =\operatorname{Var}(p_r\mid T)+(\mathbb{E}[p_r\mid T]-d_r)^2,
\]
故在 $d_r=\mathbb{E}[p_r\mid T]$ 处最小，且这些均值合起来仍是规范化的概率向量。
\end{proof}

这一推导需要明确参数先验、似然及点估计的损失，不能仅将“分子分母加常数”当作贝叶斯定理本身。由于已标注完整数据下参数向量的后验独立，对一个新样本的后验预测联合概率还满足
\[
    P(X_{\mathrm{new}}=x,Y_{\mathrm{new}}=c_k\mid T)
    =\mathbb{E}[\pi_k\prod_j\theta_{jk}(x^{(j)})\mid T]
    =\widehat\pi_k^{(\lambda)}\prod_j\widehat\theta_{jk}^{(\lambda)}(x^{(j)}),
\]
其中新样本在参数给定后服从同一朴素贝叶斯模型，并与训练样本条件独立。对这一联合概率按类别归一化，就得到平滑后的分类后验。因此，在上述先验设定下，用后验均值代入分类公式也等价于对一个新样本先积分参数再分类。

\begin{notealert}[type=caution,title={两种“后验最大化”的对象}]
分类时最大化的是给定输入后的\emph{类别后验概率}；参数的\emph{最大后验估计}（maximum a posteriori estimation，MAP）则最大化参数的后验密度。式~\eqref{eq:lihang-nb-smoothed-probabilities} 是后验均值。对 $q$ 类概率向量，若后验参数满足 $m_r+\alpha_r>1$，其内部众数为 $(m_r+\alpha_r-1)/(m+\sum_s\alpha_s-q)$，通常不同于后验均值。$\lambda=1$ 的对称均匀先验给出的后验均值是拉普拉斯平滑，参数 MAP 则退回极大似然解；$\lambda=0$ 不对应合法的对称 Dirichlet 概率先验。
\end{notealert}

\begin{example}
\label{ex:lihang-nb-laplace-classification}
沿用例~\ref{ex:lihang-nb-mle-classification} 的数据，对先验和条件概率均取 $\lambda=1$，重新预测 $x=(2,S)^{\mathsf T}$ 的类别。
\end{example}

\begin{solve}
类别数为 $K=2$，两个特征的取值数均为 $3$，故先验的分母为 $15+2=17$，两个类别的条件概率分母分别为 $9+3=12$、$6+3=9$。平滑后的先验为 $\widehat\pi_+^{(1)}=10/17,\widehat\pi_-^{(1)}=7/17$；按前例的取值顺序，四个条件概率向量为
\[
    \widehat\theta_{1,+}^{(1)}=\frac{(3,4,5)}{12},\qquad
    \widehat\theta_{1,-}^{(1)}=\frac{(4,3,2)}{9},\qquad
    \widehat\theta_{2,+}^{(1)}=\frac{(2,5,5)}{12},\qquad
    \widehat\theta_{2,-}^{(1)}=\frac{(4,3,2)}{9}.
\]
因此
\begin{align*}
    \widehat s_+^{(1)}(x)&=\frac{10}{17}\cdot\frac{4}{12}\cdot\frac{2}{12}
    =\frac{5}{153}\approx0.03268,\\
    \widehat s_-^{(1)}(x)&=\frac{7}{17}\cdot\frac{3}{9}\cdot\frac{4}{9}
    =\frac{28}{459}\approx0.06100.
\end{align*}
负类得分仍较大，故预测仍为 $-1$；归一化后有
\[
    \widehat P_1(Y=1\mid X=x)=\frac{15}{43},\qquad
    \widehat P_1(Y=-1\mid X=x)=\frac{28}{43}.
\]
这里 $\widehat P_1$ 表示 $\lambda=1$ 的估计分布。平滑改变了估计的后验概率，但本例的类别排序没有改变；这并不意味着平滑在其他样本上也保持类别不变。
\end{solve}

\section{本章概要与习题}
\label{sec:lihang-nb-summary-exercises}

朴素贝叶斯法把高维类条件联合分布化为各个特征条件概率的乘积，通过较少的参数描述生成过程。训练阶段按类别与特征取值计数，预测阶段比较先验与类条件概率的乘积；需要报告后验概率时再归一化。零一损失解释了为何选择后验概率最大的类，但最优性始终相对于所用分布而言。极大似然估计用经验频率拟合概率，贝叶斯估计通过参数先验引入伪计数；拉普拉斯平滑能够处理预定取值集合中的零频数，却不能消除条件独立假设造成的模型偏差。

教材章末两道习题分别要求推导极大似然公式与贝叶斯估计公式。下面将它们与正文推导对应，便于复习时独立完成。

\begin{example}
\label{ex:lihang-nb-exercise-mle}
教材习题一：从极大似然原理推出先验概率及条件概率的估计公式，并说明零频数和零样本类别如何处理。
\end{example}

\begin{solve}
先按式~\eqref{eq:lihang-nb-likelihood} 写出有监督样本的联合似然，再按概率向量分组。在所有计数严格为正的概率向量内部，用拉格朗日乘子法最大化 $\sum_r m_r\log p_r$，约束为 $\sum_r p_r=1$，驻点条件为 $m_r/p_r=\mu$；两边乘以 $p_r$ 后求和，得到 $\mu=\sum_r m_r=m$，故 $p_r=m_r/m$。分别代入类别计数和类内特征计数，便得到式~\eqref{eq:lihang-nb-mle-prior} 与式~\eqref{eq:lihang-nb-mle-conditional}。存在零计数时，内部驻点推导不充分，应使用命题~\ref{prop:lihang-nb-mle} 证明中的对数不等式处理边界；若 $N_k=0$，条件概率不受该组数据的似然约束，不能写成 $0/0$。
\end{solve}

\begin{example}
\label{ex:lihang-nb-exercise-bayes}
教材习题二：在适当的先验与估计准则下推出加法平滑公式，解释拉普拉斯平滑为何在每个可能取值上加 $1$。
\end{example}

\begin{solve}
对每个概率向量独立采用各分量参数均为 $\lambda>0$ 的 Dirichlet 先验，乘上分类计数似然后，后验参数由 $\lambda$ 更新为“频数 $+\lambda$”。由命题~\ref{prop:lihang-nb-dirichlet-estimation}，平方损失下的贝叶斯估计等于后验均值，即每个更新后的参数除以该向量的参数总和；对先验向量与类条件向量分别应用，便得式~\eqref{eq:lihang-nb-smoothed-probabilities}。取 $\lambda=1$ 即为拉普拉斯平滑；它来自均匀 Dirichlet 先验的后验均值，而不是将后验密度最大化。
\end{solve}

教材的继续阅读将朴素贝叶斯与更一般的贝叶斯网络联系起来：本章模型对应于类别指向各特征的有向图，给定类别后没有额外的特征依赖。若希望表达类内特征间的依赖，需要进一步规定有向无环图结构及相应的条件概率分解，而不能只取消乘积假设就认为已经得到完整的新模型。这里仅指出后续方向，不展开图模型的学习与推断。

% 教材来源：机器学习/textbooks/机器学习（分章原文）/机器学习-5.pdf。
% 覆盖教材第 5 章正文及习题；章、节编号随全文自动生成。

\fi

\selectchapter{5}
\ifcompilechapter
\chapter{决策树}
\label{ch:lihang-decision-tree}

本章对应教材第 5 章。决策树（decision tree）通过一系列特征测试，把输入空间逐步划分为若干区域，并在每个区域给出类别或数值预测。它把复杂的判断拆成可解释的条件分支，预测时只需沿一条路径进行测试。学习的困难在于：怎样选择划分、何时停止细分，以及怎样在拟合训练数据与控制模型复杂度之间取得平衡。以下先讨论分类树的三种理解方式，再建立信息增益、信息增益比和基尼指数等划分准则，最后说明树的生成、剪枝及分类与回归树算法。

\section{决策树模型与学习}
\label{sec:lihang-tree-model}

\subsection{决策树模型}

设输入空间为 $\mathcal X\subseteq\mathbb R^d$，类别集合为 $\mathcal Y=\{C_1,\ldots,C_K\}$，训练数据为
\[
  D=\{(x_i,y_i)\}_{i=1}^N,\qquad x_i\in\mathcal X,\quad y_i\in\mathcal Y.
\]
离散属性也可用有限集合中的取值表示；这里的向量记号并不要求所有属性都有数值大小的意义。$X$ 和 $Y$ 分别表示随机输入与随机类别，$x_i$ 和 $y_i$ 表示已经观测到的样本取值。

\begin{definition}[分类决策树]
分类决策树是一棵有限的有根有向树。内部结点规定对输入特征的测试，各条出边规定互斥的测试结果，叶结点规定一个类别。对输入 $x$，从根结点开始按测试结果逐层选择分支，到达叶结点后输出该叶的类别。每个内部结点的分支应覆盖所有允许到达该结点的输入。
\end{definition}

离散特征 $A$ 的测试可以按 $A(x)=a_1,\ldots,a_r$ 产生多个分支；数值特征的测试可以按 $x^{(j)}\leq s$ 与 $x^{(j)}>s$ 产生两个分支。测试的是输入的属性，输出的是类别，两者不能混淆。若一条预测路径包含 $h$ 个内部结点且每次测试为常数时间，则该输入的预测时间为 $O(h)$；树越深，预测所需的测试通常越多。

\subsection{决策树与条件规则}

从根到某个叶的路径可写为一条条件规则：若路径上的所有测试条件同时成立，则预测该叶的类别。例如，“有自己的房子，或者没有自己的房子但有工作”可以分别由两条路径表达。不同路径可以输出相同类别，但它们对应的输入区域仍然互斥。

\begin{proposition}
\label{prop:lihang-tree-partition}
若每个内部结点的分支互斥且完备，则所有根到叶路径对应的非空区域构成 $\mathcal X$ 的一个划分；相应规则对每个输入恰好覆盖一次。
\end{proposition}
\begin{proof}
任取 $x\in\mathcal X$。在根结点，完备性保证至少有一个分支可走，互斥性保证至多有一个分支可走。逐层应用同样的论证，由树的有限性可知 $x$ 最终到达唯一的叶。因此所有叶区域的并是 $\mathcal X$，任意两个不同叶区域的交为空。\qedhere
\end{proof}

\subsection{决策树与条件概率分布}

设树的叶区域为 $R_1,\ldots,R_M$，其中第 $m$ 个区域包含 $N_m>0$ 个训练样本，第 $k$ 类的样本数为 $N_{mk}$。可以在该区域上定义分段常数的条件概率估计
\begin{equation}
  \widehat P_D(Y=C_k\mid X=x)=\frac{N_{mk}}{N_m},
  \qquad x\in R_m,\quad k=1,\ldots,K.
  \label{eq:lihang-tree-probability}
\end{equation}
这表示模型对同一区域中的输入使用同一个类别分布，并不声称真实分布在该区域内一定恒定。对没有训练样本的分支，经验比例没有定义，需要另行指定规则，例如沿用父结点的分布；后文的生成算法只对非空训练子集继续递归。

在 $0$--$1$ 损失下，若预测类别为 $C_k$，模型给出的条件期望损失是 $1-\widehat P_D(Y=C_k\mid X=x)$，所以最优决策是选择概率最大的类别，即叶结点中的多数类。若多个类别并列最大，应事先固定一种打破平局的规则。概率估计、由概率导出的决策，以及真实的条件概率，是三个不同的对象。

\subsection{决策树学习}

一棵树同时规定划分结构和各叶的预测。若两个样本的全部可用特征相同而标签不同，任何确定性的特征测试树都不能把它们同时分对；即使数据可以完全分开，也可能有许多训练误差为零而泛化表现不同的树。因此，学习目标不能只描述为“把训练集完全分对”，还应约束模型复杂度。

决策树学习通常包括\emph{特征选择、树的生成和剪枝}。在当前结点，特征选择比较不同测试的划分效果；树的生成从根向下递归使用局部最优测试；剪枝则在已经生成的树中回缩子树，以整体代价比较不同复杂度的模型。后文的典型目标为
\begin{equation}
  C_\alpha(T)=C(T)+\alpha|T|,\qquad \alpha\geq0,
  \label{eq:lihang-tree-regularization}
\end{equation}
其中 $C(T)$ 是训练拟合代价，$|T|$ 始终表示叶结点数，而非全部结点数。$\alpha$ 越大，叶结点带来的复杂度代价越高。生成时的贪心选择一般不保证全局最优；剪枝的最优性也只针对固定初始树的可剪子树，不包括重新设计测试所能得到的全部决策树。

\section{特征选择}
\label{sec:lihang-tree-feature-selection}

\subsection{特征选择问题}

一个合适的测试应使子结点中的类别比父结点更集中。表~\ref{tab:lihang-loan-data} 给出教材使用的贷款申请数据，其中年龄、有工作、有自己的房子和信贷情况依次记为 $A_1,A_2,A_3,A_4$，类别表示是否同意贷款。这里的“是”“否”是数据中的取值；这个小样本只用于演示算法的计算。

\begin{table}
  \centering
  \caption{贷款申请训练数据}
  \label{tab:lihang-loan-data}
  \begin{tabular}{cccccc}
    \toprule
    编号 & 年龄 & 有工作 & 有自己的房子 & 信贷情况 & 同意贷款 \\
    \midrule
    1 & 青年 & 否 & 否 & 一般 & 否 \\
    2 & 青年 & 否 & 否 & 好 & 否 \\
    3 & 青年 & 是 & 否 & 好 & 是 \\
    4 & 青年 & 是 & 是 & 一般 & 是 \\
    5 & 青年 & 否 & 否 & 一般 & 否 \\
    6 & 中年 & 否 & 否 & 一般 & 否 \\
    7 & 中年 & 否 & 否 & 好 & 否 \\
    8 & 中年 & 是 & 是 & 好 & 是 \\
    9 & 中年 & 否 & 是 & 非常好 & 是 \\
    10 & 中年 & 否 & 是 & 非常好 & 是 \\
    11 & 老年 & 否 & 是 & 非常好 & 是 \\
    12 & 老年 & 否 & 是 & 好 & 是 \\
    13 & 老年 & 是 & 否 & 好 & 是 \\
    14 & 老年 & 是 & 否 & 非常好 & 是 \\
    15 & 老年 & 否 & 否 & 一般 & 否 \\
    \bottomrule
  \end{tabular}
\end{table}

若按年龄划分，会得到三个子集；若按是否有房划分，会得到两个子集。分支数多并不自动意味着划分好，需要用同一准则度量划分前后的类别不确定性。

\subsection{信息增益}

\begin{definition}[熵与条件熵]
设有限离散随机变量 $Y$ 的分布为 $p_k=P(Y=C_k)$。其\emph{熵}（entropy）定义为
\begin{equation}
  H(Y)=-\sum_{k=1}^K p_k\log_2p_k,\qquad 0\log_20:=0.
  \label{eq:lihang-tree-entropy}
\end{equation}
设另一个有限离散随机变量 $A$ 的取值为 $a_1,\ldots,a_r$，则\emph{条件熵}（conditional entropy）为
\[
  H(Y\mid A)=\sum_{v:P(A=a_v)>0}P(A=a_v)H(Y\mid A=a_v).
\]
本章统一使用以 $2$ 为底的对数，熵的单位为比特（bit）。
\end{definition}

熵只依赖概率分布，不依赖类别的名称。它满足 $0\leq H(Y)\leq\log_2K$：类别确定时取下界，$K$ 类等概率时取上界。非负性来自 $-p_k\log_2p_k\geq0$；上界由 $\log_2K-H(Y)=\sum_kp_k\log_2(p_k/(1/K))\geq0$ 得到，后一不等式可由 $\ln u\leq u-1$ 推出。二分类时记
\[
  h(p)=-p\log_2p-(1-p)\log_2(1-p),\qquad 0\leq p\leq1.
\]
于是 $h(0)=h(1)=0$，$h(1/2)=1$；两类各占一半时最难仅凭类别比例作出判断。

对当前非空训练子集 $D$，记 $N=|D|$，第 $k$ 类样本数为 $N_k$。特征 $A$ 将 $D$ 划分为非空子集 $D_1,\ldots,D_r$，其中 $N_v=|D_v|$，$N_{vk}$ 为 $D_v$ 中第 $k$ 类样本数。以经验频率代替真实概率，得到\emph{经验熵}与\emph{经验条件熵}
\begin{align}
  H(D)&=-\sum_{k=1}^K\frac{N_k}{N}\log_2\frac{N_k}{N},
  \label{eq:lihang-tree-empirical-entropy}\\
  H(D\mid A)&=\sum_{v=1}^r\frac{N_v}{N}H(D_v)
  =-\sum_{v=1}^r\sum_{k=1}^K\frac{N_{vk}}{N}\log_2\frac{N_{vk}}{N_v}.
  \label{eq:lihang-tree-conditional-entropy}
\end{align}
其中零计数项按 $0\log_20=0$ 处理。条件熵是各分支熵的样本比例加权平均，不能直接把分支熵相加，也不能一律取等权平均。

\begin{definition}[信息增益]
特征 $A$ 对当前数据集 $D$ 的\emph{信息增益}（information gain）定义为
\begin{equation}
  g(D,A)=H(D)-H(D\mid A).
  \label{eq:lihang-tree-information-gain}
\end{equation}
信息增益准则在所有候选特征中选择 $g(D,A)$ 最大者。
\end{definition}

信息增益是经验分布下类别 $Y$ 与特征 $A$ 的\emph{互信息}（mutual information）。记 $\widehat p_{vk}=N_{vk}/N$、$\widehat p_v=N_v/N$、$\widehat p_k=N_k/N$，展开定义可得
\begin{equation}
  g(D,A)=\sum_{v,k:\widehat p_{vk}>0}\widehat p_{vk}
  \log_2\frac{\widehat p_{vk}}{\widehat p_v\widehat p_k}\geq0.
  \label{eq:lihang-tree-mutual-information}
\end{equation}
为核查非负性，令 $p$ 为联合经验分布、$q$ 为边缘分布的乘积；在 $p>0$ 的位置有 $q>0$，并且
\[
  -\sum_{p>0}p\ln(q/p)\geq\sum_{p>0}(p-q)\geq0.
\]
等号成立当且仅当两个经验随机变量独立。因此 $0\leq g(D,A)\leq H(D)$；若每个分支都是纯类，条件熵为零，增益达到 $H(D)$。这些结论针对当前样本的经验分布，不能据此断言新样本上的分类误差一定下降。

\begin{example}
\label{ex:lihang-loan-gain}
对表~\ref{tab:lihang-loan-data} 的数据计算四个特征的信息增益，并选择根结点。
\end{example}
\begin{solve}
训练集中有 $9$ 个同意贷款样本、$6$ 个不同意贷款样本，故 $H(D)=h(9/15)\approx0.970951$。按年龄划分后，青年、中年、老年组的同意贷款比例分别为 $2/5,3/5,4/5$，所以
\[
  H(D\mid A_1)=\frac13\bigl[h(2/5)+h(3/5)+h(4/5)\bigr]\approx0.887943.
\]
有工作的 $5$ 个样本全为同意，无工作的 $10$ 个样本中有 $4$ 个同意；有房的 $6$ 个样本全为同意，无房的 $9$ 个样本中有 $3$ 个同意。因此
\[
  H(D\mid A_2)=\frac{10}{15}h(4/10),\qquad
  H(D\mid A_3)=\frac9{15}h(3/9).
\]
信贷一般、好、非常好的组大小分别为 $5,6,4$，同意贷款数分别为 $1,4,4$，故
\[
  H(D\mid A_4)=\frac5{15}h(1/5)+\frac6{15}h(4/6)+\frac4{15}h(1).
\]
相应信息增益为
\[
  g(D,A_1)\approx0.083008,\quad g(D,A_2)\approx0.323650,\quad
  g(D,A_3)\approx0.419973,\quad g(D,A_4)\approx0.362990.
\]
最大者是 $A_3$，因此选择“有自己的房子”作为根结点的测试。
\end{solve}

\subsection{信息增益比}

信息增益偏向取值较多的特征。例如，在同一训练集中，若把唯一的样本编号作为特征，每个编号对应一个样本，所有分支熵均为零，增益达到 $H(D)$，但这种划分并没有提供对新编号的可靠预测规则。

\begin{definition}[信息增益比]
用特征 $A$ 的取值划分 $D$，定义划分本身的经验熵
\[
  H_A(D)=-\sum_{v=1}^r\frac{N_v}{N}\log_2\frac{N_v}{N}.
\]
当 $H_A(D)>0$ 时，\emph{信息增益比}（information gain ratio）为
\begin{equation}
  g_R(D,A)=\frac{g(D,A)}{H_A(D)}.
  \label{eq:lihang-tree-gain-ratio}
\end{equation}
\end{definition}

分母度量特征自身的取值分散程度，不是 $H(D)$。若 $H_A(D)=0$，该特征在当前样本中为常量，不能产生非空的不同分支，应排除，而不能计算 $0/0$。由互信息的对称性及条件熵非负性，$0\leq g(D,A)\leq H_A(D)$，所以合法增益比在 $[0,1]$ 内。增益比对多取值偏好作了归一化修正，但比值较大仍不等于绝对增益较大；分母很小时应同时查看实际分支样本数和增益。

\section{决策树的生成}
\label{sec:lihang-tree-generation}

\subsection{ID3 算法}

ID3 在每个结点选择信息增益最大的离散特征，并沿其取值递归生成子树。设当前非空数据为 $D$，剩余特征集合为 $\mathcal A$，停止阈值为 $\varepsilon>0$。输入为 $(D,\mathcal A,\varepsilon)$，输出为带类别标记的决策树；类别和特征并列时均按固定规则选择。
\begin{enumerate}
  \item 若 $D$ 中标签全相同，输出以该类别为标记的叶；若 $\mathcal A=\varnothing$，输出以 $D$ 的多数类为标记的叶。
  \item 否则计算所有能产生至少两个非空分支的候选特征的信息增益。若不存在这样的特征，或最大增益小于 $\varepsilon$，输出多数类叶。
  \item 选择最大增益特征 $A^*$，建立测试结点。按 $A^*$ 的取值划分 $D$，对每个非空子集 $D_v$，以 $(D_v,\mathcal A\setminus\{A^*\},\varepsilon)$ 为输入递归调用上述步骤，并连接相应子树。
\end{enumerate}
训练中未出现的合法取值可预设为返回当前结点的多数类，以保证预测规则完整。在这里的离散多叉形式中，每沿一层路径至少删除一个特征，故路径深度至多为 $d$，算法有限终止。若用散列表计数且将基本计数与对数计算视为常数时间，一个有 $m$ 个样本、$q$ 个剩余特征的结点可用 $O(mq)$ 时间计算候选增益；深度为 $h$ 时，逐层累计有 $O(Ndh)$ 的上界，因此 $h\leq d$ 时可给出 $O(Nd^2)$ 的粗略上界。

\begin{example}
\label{ex:lihang-loan-id3}
在例~\ref{ex:lihang-loan-gain} 的基础上继续生成 ID3 决策树，取足够小的正阈值。
\end{example}
\begin{solve}
根结点按有房、无房分成两个子集。有房组的 $6$ 个样本全部同意贷款，直接成为“是”叶。无房组 $D_-$ 含 $9$ 个样本，其中 $3$ 个同意、$6$ 个不同意，故 $H(D_-)=h(1/3)\approx0.918296$。在该组重新计算，得到
\[
  g(D_-,A_1)\approx0.251629,\qquad
  g(D_-,A_2)\approx0.918296,\qquad
  g(D_-,A_4)\approx0.473851.
\]
因此第二层选择“有工作”：无房但有工作的 $3$ 个样本全部同意，无房且无工作的 $6$ 个样本全部不同意。最终预测规则为
\begin{equation}
  \widehat y(x)=
  \begin{cases}
    \text{是},& A_3(x)=\text{是},\\
    \text{是},& A_3(x)=\text{否},\ A_2(x)=\text{是},\\
    \text{否},& A_3(x)=\text{否},\ A_2(x)=\text{否}.
  \end{cases}
  \label{eq:lihang-loan-tree}
\end{equation}
该树有 $2$ 个内部结点、$3$ 个叶结点，对这 $15$ 个训练样本全部分对。
\end{solve}

最大信息增益只表示当前一步的熵下降最大，不能保证最终树最小或测试误差最小。一个直接的反例是四个等频二进制输入 $(0,0),(0,1),(1,0),(1,1)$，标签为两个特征的异或：两个单独特征的信息增益均为零，采用正阈值会在根停止，但深度为 $2$ 的树可以完全分类。这也说明仅按单个特征的当前增益筛除特征，可能丢失特征交互信息。

\subsection{C4.5 的生成算法}

教材给出的 C4.5 离散特征生成过程保留上述递归结构，将特征评分换成式~\eqref{eq:lihang-tree-gain-ratio} 的信息增益比：只在 $H_A(D)>0$ 的特征中选最大者，若最大比值小于预设阈值则输出多数类叶，否则按该特征分支并在各子集重新计算。纯类、无剩余特征以及无有效划分时同样停止。信息增益与增益比量纲和数值范围不同，停止阈值应针对所用评分设置。

\begin{example}
\label{ex:lihang-loan-c45}
用信息增益比为表~\ref{tab:lihang-loan-data} 的数据生成决策树，对应教材习题 5.1。
\end{example}
\begin{solve}
根结点的四个分母和比值如下；表中数值保留六位小数。
\begin{center}
  \begin{tabular}{cccc}
    \toprule
    特征 & $g(D,A_j)$ & $H_{A_j}(D)$ & $g_R(D,A_j)$ \\
    \midrule
    年龄 & 0.083008 & 1.584963 & 0.052372 \\
    有工作 & 0.323650 & 0.918296 & 0.352447 \\
    有自己的房子 & 0.419973 & 0.970951 & 0.432538 \\
    信贷情况 & 0.362990 & 1.565596 & 0.231854 \\
    \bottomrule
  \end{tabular}
\end{center}
最大者仍是“有自己的房子”。有房分支为纯类；在无房子集 $D_-$ 中，年龄、有工作、信贷情况的增益比分别约为 $0.164411,1,0.340374$，因此第二层选择“有工作”。当阈值足够小时，得到的树仍为式~\eqref{eq:lihang-loan-tree}。本例中两种准则选出同一棵树，只是数据计算的结果，并不说明它们通常等价。
\end{solve}

\section{决策树的剪枝}
\label{sec:lihang-tree-pruning}

生成过程不断降低训练不纯度，过深的树可能把有限样本中的偶然差异也当作规则。\emph{剪枝}（pruning）把一个内部结点及其全部后代替换为一个叶结点，并在该结点重新确定类别或预测值。以下只考虑这种回缩操作，不允许改变保留结点原有的测试。

设 $\mathcal L(T)$ 为树 $T$ 的叶集合。对叶 $t$，记其样本数为 $N_t$，第 $k$ 类样本数为 $N_{tk}$，经验熵为 $H_t=-\sum_k(N_{tk}/N_t)\log_2(N_{tk}/N_t)$。教材使用的熵代价为
\begin{equation}
  C_H(T)=\sum_{t\in\mathcal L(T)}N_tH_t,
  \qquad C_{H,\alpha}(T)=C_H(T)+\alpha|T|.
  \label{eq:lihang-tree-entropy-cost}
\end{equation}
这里用的是样本数加权的熵之和，不是各叶熵的简单平均，也不是误分类样本数。若改用 $C_H(T)/N$，为得到相同的最优树，应同时把复杂度系数改为 $\alpha/N$。

\begin{proposition}[熵代价的似然解释]
固定树的划分，对各叶分别用类别频率估计类别概率，则 $C_H(T)$ 恰好等于拟合后条件似然的以 $2$ 为底的负对数。
\end{proposition}
\begin{proof}
叶 $t$ 的类别概率记为 $q_{t1},\ldots,q_{tK}$。在非负且和为 $1$ 的概率向量中，经验频率 $\widehat q_{tk}=N_{tk}/N_t$ 最大化该叶的多项式似然，因为
\[
  -\sum_k\widehat q_{tk}\log_2q_{tk}
  =H_t+\sum_{k:\widehat q_{tk}>0}\widehat q_{tk}
  \log_2\frac{\widehat q_{tk}}{q_{tk}}\geq H_t.
\]
若在正计数类别上令 $q_{tk}=0$，左端为正无穷，也不能改善目标。将最优概率代入全部样本的条件似然，得到
\[
  -\log_2\prod_{i=1}^N\widehat P_D(Y=y_i\mid X=x_i)
  =-\sum_{t,k:N_{tk}>0}N_{tk}\log_2\frac{N_{tk}}{N_t}
  =C_H(T).\qedhere
\]
\end{proof}

因此，式~\eqref{eq:lihang-tree-entropy-cost} 的最小化是一种正则化的极大似然模型选择。若父结点 $t$ 的一次划分按特征 $A$ 产生各子结点，划分前后的熵代价之差为 $N_tg(D_t,A)$，这里 $D_t$ 为到达 $t$ 的训练样本集。这解释了信息增益与局部似然改善的联系；逐步贪心最大化该改善，仍不同于在全部树结构中求整体最优。

设当前以 $t$ 为根的子树为 $T_t$，有 $L_t\geq2$ 个叶，回缩后单叶的未惩罚代价为 $C_H(t)=N_tH_t$。树外的代价在比较时抵消，剪去 $T_t$ 的条件为
\begin{equation}
  C_H(t)-C_H(T_t)\leq\alpha(L_t-1).
  \label{eq:lihang-tree-local-pruning}
\end{equation}
左端是合并叶结点造成的拟合代价增加，右端是节省的复杂度代价。两边相等时取较小的树。这种局部比较要自下而上进行，并在比较父结点前先处理其子树。

\begin{proposition}[固定初始树的最优剪枝]
\label{prop:lihang-tree-pruning-dp}
给定有限初始树 $T_0$、$\alpha\geq0$ 和对叶可加的代价 $C(T)=\sum_{t\in\mathcal L(T)}C(t)$。其中 $C(t)$ 是将结点 $t$ 作为叶时用到达该结点的固定训练样本重新拟合得到的代价。令 $F_\alpha(t)$ 为初始子树 $T_{0,t}$ 的所有可剪子树中的最小惩罚代价，则
\begin{equation}
  F_\alpha(t)=
  \begin{cases}
    C(t)+\alpha,&t\text{ 是初始叶},\\
    \displaystyle\min\left\{C(t)+\alpha,
      \sum_{u\in\operatorname{ch}(t)}F_\alpha(u)\right\},&t\text{ 是内部结点},
  \end{cases}
  \label{eq:lihang-tree-pruning-dp}
\end{equation}
其中 $\operatorname{ch}(t)$ 表示 $t$ 的直接子结点集合。最小值存在；若代价相等时选择回缩为叶，则得到唯一的最小最优子树，它是每棵最优子树的可剪子树。
\end{proposition}
\begin{proof}
初始叶没有其他选择。对内部结点，任何可剪子树或者仅保留 $t$ 作为叶，或者保留 $t$ 的测试并在每个子结点下独立选择一棵可剪子树。代价可加，故保留测试时的最小代价为各 $F_\alpha(u)$ 之和，得到递推式。初始树有限，因此候选数有限，最小值存在。

再从叶向上归纳最小最优子树的唯一性。若回缩的代价严格更小，所有最优树都必须回缩；若保留测试的代价严格更小，每个子结点都必须选择局部最优子树，而归纳假设给出唯一的最小选择；若两种代价相等，单叶本身最优且可由每棵保留测试的候选树剪得。三种情形都给出唯一且包含于所有最优树的最小选择。\qedhere
\end{proof}

该命题给出教材剪枝过程的动态规划解释：先计算结点的单叶代价，再按后序遍历计算递推式，并记录回缩或保留的选择。已有各结点代价时，一次遍历所需时间为 $O(|V(T_0)|)$，其中 $V(T_0)$ 表示初始树的全部结点集合。\emph[caution]{唯一的是最小最优子树；若不规定相等时选较小树，达到最小代价的树可能不唯一。}这也回答了教材习题 5.3 中“唯一的最小子树”的含义。

\section{CART 算法}
\label{sec:lihang-cart}

\emph{分类与回归树}（classification and regression tree，CART）统一采用二叉划分。分类树通常用基尼指数选择测试，回归树通常用平方误差选择测试；树生成后再以代价复杂度剪枝产生候选序列，并用验证数据选择复杂度。内部结点回答一个“是或否”的问题，并不要求原特征只有两个取值。

\subsection{CART 生成}

\subsubsection{最小二乘回归树}

回归时 $y_i\in\mathbb R$。设叶区域 $R_1,\ldots,R_M$ 构成输入空间的划分，各区域输出常数 $c_m$，则回归函数为
\begin{equation}
  f(x)=\sum_{m=1}^Mc_m\mathbf1\{x\in R_m\},
  \label{eq:lihang-regression-tree}
\end{equation}
其中 $\mathbf1\{\cdot\}$ 为示性函数。固定划分后的平方误差目标为
\[
  \sum_{i=1}^N(y_i-f(x_i))^2
  =\sum_{m=1}^M\sum_{i:x_i\in R_m}(y_i-c_m)^2.
\]
不同区域的参数可分别求解。

\begin{proposition}[回归叶的最优常数]
\label{prop:lihang-regression-leaf-mean}
若区域 $R$ 中有 $n_R>0$ 个训练样本，则平方误差的唯一最优常数是其响应均值 $\overline y_R=n_R^{-1}\sum_{i:x_i\in R}y_i$。
\end{proposition}
\begin{proof}
利用 $\sum_{i:x_i\in R}(y_i-\overline y_R)=0$，有
\[
  \sum_{i:x_i\in R}(y_i-c)^2
  =\sum_{i:x_i\in R}(y_i-\overline y_R)^2+n_R(c-\overline y_R)^2.
\]
后一项非负，且仅在 $c=\overline y_R$ 时为零。\qedhere
\end{proof}

在当前区域内，数值特征 $j$ 和切分点 $s$ 产生两个子区域 $R_L=\{x:x^{(j)}\leq s\}$、$R_R=\{x:x^{(j)}>s\}$，这里两者都与当前区域取交。只比较两侧训练样本数均非零的划分，求解
\begin{equation}
  (j^*,s^*)\in\operatorname*{arg\,min}_{j,s}
  \left\{\sum_{i:x_i\in R_L}(y_i-\overline y_L)^2
  +\sum_{i:x_i\in R_R}(y_i-\overline y_R)^2\right\}.
  \label{eq:lihang-cart-regression-split}
\end{equation}
选定切分后在两个子区域重复这一过程，直到达到预设最大深度、最小结点样本数、最小叶样本数，或没有满足条件的有效切分；若另设最小误差下降量，也应在生成前确定。输出为各叶区域及其响应均值，而非某个训练样本的标签。

\begin{proposition}[平方误差的划分收益]
设当前结点分为两个非空子结点，样本数分别为 $n_L,n_R$，均值分别为 $\overline y_L,\overline y_R$，则划分前后平方误差之差为
\begin{equation}
  \Delta_{\mathrm{SSE}}=
  \frac{n_Ln_R}{n_L+n_R}(\overline y_L-\overline y_R)^2\geq0.
  \label{eq:lihang-cart-sse-gain}
\end{equation}
\end{proposition}
\begin{proof}
父结点均值为 $\overline y=(n_L\overline y_L+n_R\overline y_R)/(n_L+n_R)$。分别对两组应用命题~\ref{prop:lihang-regression-leaf-mean} 证明中的平方和分解，误差下降等于 $n_L(\overline y_L-\overline y)^2+n_R(\overline y_R-\overline y)^2$，代入父均值并化简即得结论。\qedhere
\end{proof}

对一个有 $m$ 个样本的结点，将第 $j$ 个数值特征的不同值排序为 $v_1<\cdots<v_q$，只需比较相邻值之间的阈值，例如 $s=(v_\ell+v_{\ell+1})/2$，因为同一间隔内的阈值在训练集上产生相同分组。用某一侧的样本数 $n$、响应和 $S$、响应平方和 $Q$ 可计算其最小平方误差 $Q-S^2/n$，因此排序后可用前缀和线性扫描所有阈值。对 $d$ 个数值特征逐一排序并扫描，单结点的直接实现用时 $O(dm\log m)$；这是该实现的结点代价，不能当作整棵树的总训练代价。

\begin{example}
\label{ex:lihang-cart-regression-exercise}
教材习题 5.2 给出 $x_i=i$，$i=1,\ldots,10$，响应依次为
\[
  4.50,\ 4.75,\ 4.91,\ 5.34,\ 5.80,\ 7.05,\ 7.90,\ 8.23,\ 8.70,\ 9.00.
\]
用平方误差准则生成二叉回归树。
\end{example}
\begin{solve}
根结点的响应均值为 $6.618$，不切分时平方误差为 $27.63236$。枚举九个中点，各切分的两侧平方误差之和如下，表中误差保留四位小数。
\begin{center}
  \begin{tabular}{cccccc}
    \toprule
    $s$ & 误差 & $s$ & 误差 & $s$ & 误差 \\
    \midrule
    1.5 & 22.6480 & 4.5 & 7.3787 & 7.5 & 10.0525 \\
    2.5 & 17.7022 & 5.5 & 3.3587 & 8.5 & 15.1778 \\
    3.5 & 12.1935 & 6.5 & 5.0740 & 9.5 & 21.3280 \\
    \bottomrule
  \end{tabular}
\end{center}
因此第一刀取 $s=5.5$，两侧均值为 $5.060$ 和 $8.176$，只保留这一刀时
\[
  f_2(x)=
  \begin{cases}
    5.060,&x\leq5.5,\\
    8.176,&x>5.5.
  \end{cases}
\]
继续递归，左、右两组的最优切分分别为 $3.5$ 和 $7.5$。若在深度为 $2$ 时停止，得到四个区域，输出依次为 $4.720,5.570,7.475,8.643333\ldots$。若约定一直生长到每叶仅含一个样本，则所有内部结点可用下表描述，其中索引范围指到达该结点的训练样本，而非预先限定的输入定义域。
\begin{center}
  \begin{tabular}{cc}
    \toprule
    样本索引 & 切分点 \\
    \midrule
    $1$--$10$ & 5.5 \\
    $1$--$5$ & 3.5 \\
    $6$--$10$ & 7.5 \\
    $1$--$3$ & 1.5 \\
    $4$--$5$ & 4.5 \\
    $6$--$7$ & 6.5 \\
    $8$--$10$ & 8.5 \\
    $2$--$3$ & 2.5 \\
    $9$--$10$ & 9.5 \\
    \bottomrule
  \end{tabular}
\end{center}
最终十个叶的输出分别为原来的十个 $y_i$，训练平方误差为零。题目未指定停止条件时，树的最终大小并不唯一；上述约定给出了可复现的完整生成过程，实际预测仍应通过验证选择停止条件或剪枝强度。
\end{solve}

\subsubsection{分类树与基尼指数}

\begin{definition}[基尼指数]
设 $K$ 类的概率向量为 $p=(p_1,\ldots,p_K)$，定义\emph{基尼指数}（Gini index）为
\begin{equation}
  \operatorname{Gini}(p)=\sum_{k=1}^Kp_k(1-p_k)
  =1-\sum_{k=1}^Kp_k^2.
  \label{eq:lihang-gini}
\end{equation}
对非空样本集 $D$，用类别频率 $N_k/N$ 代入得到 $\operatorname{Gini}(D)$。若一个二叉测试将 $D$ 分成非空的 $D_L,D_R$，其划分后不纯度为
\begin{equation}
  G(D;D_L,D_R)=\frac{|D_L|}{|D|}\operatorname{Gini}(D_L)
  +\frac{|D_R|}{|D|}\operatorname{Gini}(D_R).
  \label{eq:lihang-cart-gini-split}
\end{equation}
\end{definition}

基尼指数等于从分布 $p$ 中独立抽取两个类别时二者不同的概率，取值范围为 $[0,1-1/K]$；纯类时为零，均匀分布时最大。二分类时 $\operatorname{Gini}(p)=2p(1-p)$，多数类分类器的错误率则为 $\min\{p,1-p\}$。两者都反映混杂程度，但基尼指数并非多数类的实际误分类率，熵也不是该误分类率。

记两子结点的类别比例向量为 $p_L,p_R$，$w=|D_L|/|D|$，则父结点的比例为 $p=wp_L+(1-w)p_R$。展开平方和得到
\begin{equation}
  \operatorname{Gini}(D)-G(D;D_L,D_R)
  =w(1-w)\sum_{k=1}^K(p_{Lk}-p_{Rk})^2\geq0.
  \label{eq:lihang-cart-gini-gain}
\end{equation}
因此，在同一父结点上最小化划分后的加权基尼指数，等价于最大化基尼下降。与回归划分收益一样，收益同时取决于两侧比例和两侧的差异。

CART 分类树的输入为当前训练集与预先规定的停止条件。在每个结点枚举候选二叉测试，选择式~\eqref{eq:lihang-cart-gini-split} 最小者，建立两个子结点并递归；纯类、达到深度或样本数限制、无有效切分等情形输出多数类叶。对数值特征可比较 $x^{(j)}\leq s$；对离散特征，教材采用 $A=a$ 与 $A\ne a$ 的测试。更一般的离散二叉测试也可按取值子集与其补集划分，所允许的候选集合应事先确定。在二叉树中，一个多值或数值特征可能沿同一路径再次提供有效切分，因此不能像离散多叉 ID3 那样在第一次使用后统一删除该特征。

\begin{example}
\label{ex:lihang-loan-cart}
对表~\ref{tab:lihang-loan-data} 的数据，按教材的离散二叉测试生成 CART 分类树。
\end{example}
\begin{solve}
父结点的基尼指数为 $1-(9/15)^2-(6/15)^2=0.48$。四个特征的候选切分结果为
\begin{center}
  \begin{tabular}{ccc}
    \toprule
    特征 & 测试取值（另一分支为其补集） & 划分后的加权基尼指数 \\
    \midrule
    年龄 & 青年；中年；老年 & $0.44;\ 0.48;\ 0.44$ \\
    有工作 & 是 & $0.32$ \\
    有自己的房子 & 是 & $0.266667$ \\
    信贷情况 & 一般；好；非常好 & $0.32;\ 0.474074;\ 0.363636$ \\
    \bottomrule
  \end{tabular}
\end{center}
例如，有房组为纯类，无房组的同意比例为 $1/3$，故该切分的加权指数为
\[
  \frac6{15}\times0+\frac9{15}\times2\times\frac13\times\frac23
  =\frac4{15}\approx0.266667.
\]
它是所有候选中的最小者。在无房子集中再按“有工作”切分，两侧都变为纯类，加权指数降到零。因此本例的 CART 树也与式~\eqref{eq:lihang-loan-tree} 一致。
\end{solve}

\subsection{CART 剪枝}
\label{subsec:lihang-cart-pruning}

\subsubsection{代价复杂度与有效剪枝阈值}

给定已生成的二叉树 $T_0$，记 $T\preceq T_0$ 表示 $T$ 可由 $T_0$ 回缩子树得到。\emph{代价复杂度剪枝}（cost-complexity pruning）求解
\begin{equation}
  \min_{T\preceq T_0}\bigl\{C(T)+\alpha|T|\bigr\},\qquad\alpha\geq0.
  \label{eq:lihang-cart-cost-complexity}
\end{equation}
这里 $C$ 必须是统一约定的可加拟合代价。例如，回归树可取各叶平方误差之和，分类树可取 $\sum_tN_t\operatorname{Gini}(D_t)$ 或误分类数 $\sum_t(N_t-\max_kN_{tk})$。生成时选择测试所用的准则与剪枝代价是可以区分的，但一次剪枝序列中不能混用不同代价或归一化方式。上述代价都满足继续细分不会增加最优训练拟合代价。

在当前树 $T$ 中，对任意内部结点 $t$，其当前子树为 $T_t$，若把它回缩为单叶，则两方案的代价差为
\begin{equation}
  C_\alpha(t)-C_\alpha(T_t)
  =C(t)-C(T_t)-\alpha(|T_t|-1).
  \label{eq:lihang-cart-pruning-difference}
\end{equation}
因此定义该结点的有效剪枝阈值
\begin{equation}
  \eta_T(t)=\frac{C(t)-C(T_t)}{|T_t|-1}.
  \label{eq:lihang-cart-effective-alpha}
\end{equation}
分母是减少的叶数，分子是增加的训练代价，所以 $\eta_T(t)$ 表示每减少一个叶所付出的平均拟合代价。在 $\alpha<\eta_T(t)$ 时保留当前子树更好，在 $\alpha>\eta_T(t)$ 时回缩更好，等号处两者并列。\emph[caution]{该阈值不是整体损失下降量；它是使两种方案代价相等的复杂度系数。}

\subsubsection{最弱环节剪枝与最优子树序列}

\emph{最弱环节剪枝}（weakest-link pruning）每轮计算当前树所有内部结点的 $\eta_T(t)$，取最小值作为下一临界参数，剪去达到该值的子树。若并列结点存在祖先与后代关系，剪去祖先即已包含对后代的处理；等价地，在临界值处按命题~\ref{prop:lihang-tree-pruning-dp} 选择最小最优子树。每次剪枝后须重新计算剩余结点的子树代价与叶数，不能沿用初始树的全部阈值。

\begin{theorem}[最小最优子树的嵌套性与剪枝路径]
\label{thm:lihang-cart-pruning-path}
设 $T_0$ 为有限二叉树，$C(T)$ 对叶可加，且细分不增加未惩罚代价。对每个 $\alpha\geq0$，在最小化式~\eqref{eq:lihang-cart-cost-complexity} 的树中取最小者，记为 $T_\alpha$。则 $T_\alpha$ 唯一，且 $\beta>\alpha$ 时有 $T_\beta\preceq T_\alpha$。先将 $T_0$ 在 $\alpha=0$ 时作零代价回缩，所得最小最优树记为 $S_0$；从 $S_0$ 开始的最弱环节剪枝产生有限嵌套序列
\[
  S_0\succ S_1\succ\cdots\succ S_J
\]
及严格递增的临界值 $0=a_0<a_1<\cdots<a_J$，其中 $S_J$ 仅含根结点。令 $a_{J+1}=+\infty$，则对 $\alpha\in[a_j,a_{j+1})$，最小最优树为 $S_j$。若 $S_0$ 已为单叶，结论退化为只有一个区间。
\end{theorem}
\begin{proof}
存在性和最小最优树的唯一性由命题~\ref{prop:lihang-tree-pruning-dp} 给出。为证明嵌套性，设某个结点 $t$ 在参数 $\alpha$ 下被最小最优子树回缩为叶。对初始树中以 $t$ 为根的任意真正展开的候选树 $U$，局部最优性给出
\[
  C(U)+\alpha|U|\geq C(t)+\alpha.
\]
当 $\beta>\alpha$ 时，
\[
  C(U)+\beta|U|-C(t)-\beta
  \geq(\beta-\alpha)(|U|-1)>0.
\]
因此参数增大后，若仍到达 $t$，该处依然必须取叶，不会重新展开。对所有已回缩结点应用这一论证，得到 $T_\beta\preceq T_\alpha$。

现设当前最小最优树为 $S$，将其剪成任意候选 $U\preceq S$。存在一组互不为祖先后代的内部结点 $B$，使得 $U$ 恰好通过回缩各 $S_t$、$t\in B$ 得到。因此
\begin{equation}
  C_\alpha(U)-C_\alpha(S)
  =\sum_{t\in B}(|S_t|-1)\bigl(\eta_S(t)-\alpha\bigr).
  \label{eq:lihang-cart-subtree-comparison}
\end{equation}
记 $a=\min_t\eta_S(t)$。若 $\alpha<a$，任何真回缩都会严格增加代价；在 $\alpha=a$ 时，所有阈值等于 $a$ 的回缩均不增加代价，取它们共同给出的最小树即得下一棵最小最优树。由已证嵌套性，后续全局最优树都在 $S$ 的可剪子树内，因此这种比较没有遗漏需要重新展开的更优候选。

在每个临界值处必须处理完全部并列最优的回缩；这样得到的树若仍有内部结点，其所有下一阈值都严格大于当前值，否则还可继续作同代价回缩，与当前树最小矛盾。每轮至少减少一个叶，故过程有限终止。最终只能留下根结点；在相邻临界值之间，式~\eqref{eq:lihang-cart-subtree-comparison} 表明最小最优树不变，得到所述区间。\qedhere
\end{proof}

该定理同时给出教材习题 5.4 的证明。若初始树不存在任何零代价回缩，就有 $S_0=T_0$；若存在零阈值，先消除它们才能把后续临界值写成严格递增序列。在临界值上，上一棵较大的树和剪后较小的树可以同时达到最小代价，因此以左闭右开区间记录较小的树，与“唯一的最小最优树”的约定一致。

\begin{example}
\label{ex:lihang-loan-pruning}
对例~\ref{ex:lihang-loan-id3} 的三叶树使用熵代价剪枝，计算子树的临界参数。
\end{example}
\begin{solve}
三叶都为纯类，故整树 $T$ 的 $C_H(T)=0$、$|T|=3$。无房子树有两个纯叶，回缩成单叶的代价为 $9h(1/3)\approx8.264663$，其有效阈值约为 $8.264663$。根结点回缩代价为 $15h(3/5)\approx14.564259$，其有效阈值则为
\[
  \frac{15h(3/5)-0}{3-1}\approx7.282129.
\]
因此先剪去根结点以下的整棵子树：$0\leq\alpha<7.282129\ldots$ 时最小最优树是三叶树，达到该临界值后最小最优树是仅含根的多数类叶。只把无房子树回缩所得的两叶树从未成为最优，这说明剪枝路径可以跳过某些叶数。
\end{solve}

\subsubsection{用验证选择最终模型}

剪枝路径由训练集上的 $C(T)$ 计算得到；确定最终复杂度还需要验证。若采用独立验证集 $D_{\mathrm{val}}$，回归可比较
\[
  \widehat R_{\mathrm{val}}(S_j)
  =\frac1{|D_{\mathrm{val}}|}\sum_{(x,y)\in D_{\mathrm{val}}}(y-f_{S_j}(x))^2,
\]
分类可比较平均 $0$--$1$ 损失或预先选定的概率损失。验证时保持训练得到的叶预测；若使用基尼等不纯度评价，也应明确其计算规则及与最终评价目标的关系。验证误差最小与训练惩罚代价最小是两次不同的选择，前者用于在候选复杂度之间估计泛化表现。

若采用 $V$ 折交叉验证（$V$ 为折数），每一折都应仅用该折的训练部分重新生成树和剪枝路径，用留出部分比较共同候选复杂度参数的预测损失，最后在全体训练数据上重新拟合选定复杂度的模型。不能先在全数据上确定树结构，再把其中一部分称作独立验证数据。验证选出的树是给定候选集合和评价规则下的选择，并不保证在未知总体上达到真正的最小风险。

\section{本章联系与学习检查}
\label{sec:lihang-tree-review}

决策树的三种表示互相补充：树结构给出测试顺序，路径规则给出局部判断条件，叶区域上的概率分布说明预测依据。ID3、教材所述 C4.5 生成过程与 CART 的主要区别如下。
\begin{center}
  \begin{tabular}{cccc}
    \toprule
    方法 & 主要任务 & 当前测试的选择准则 & 本章讨论的分支形式 \\
    \midrule
    ID3 & 分类 & 最大信息增益 & 离散特征多叉划分 \\
    C4.5 生成 & 分类 & 最大信息增益比 & 离散特征多叉划分 \\
    CART 分类 & 分类 & 最小加权基尼指数 & 二叉测试 \\
    CART 回归 & 回归 & 最小两侧平方误差和 & 二叉测试 \\
    \bottomrule
  \end{tabular}
\end{center}

信息增益、基尼下降与平方误差下降都评价当前划分的训练效果；它们不能替代复杂度控制或独立验证。剪枝则在固定初始树上比较完整子树的拟合代价与叶数，动态规划说明固定 $\alpha$ 时怎样取最小最优子树，最弱环节算法说明怎样随 $\alpha$ 增大生成嵌套候选序列。

复习时应能从样本计数写出经验熵、条件熵与各候选测试的加权不纯度，解释常量特征为何不能计算增益比，区分回归叶的均值与分类叶的多数类，并由式~\eqref{eq:lihang-cart-pruning-difference} 推出有效阈值。教材习题 5.1、5.2 的计算分别见例~\ref{ex:lihang-loan-c45} 与例~\ref{ex:lihang-cart-regression-exercise}；习题 5.3、5.4 的论证分别见命题~\ref{prop:lihang-tree-pruning-dp} 与定理~\ref{thm:lihang-cart-pruning-path}。这些练习共同检验从模型、局部划分到整体模型选择的逻辑链条。

% 教材来源：机器学习/textbooks/机器学习（分章原文）/机器学习-6.pdf。
% 对应教材第 6 章，原书第 91--110 页；本文件采用自动章节编号。


\fi

\selectchapter{6}
\ifcompilechapter
\chapter{逻辑斯谛回归与最大熵模型}
\label{ch:lh-logistic-maxent}

分类不仅要给出类别，还常常需要表达预测的不确定性。本章研究条件概率模型 $P(Y\mid X)$：先把输入的线性得分转换为二分类或多分类概率，得到逻辑斯谛回归；再从满足经验特征约束的条件分布中选择熵最大的一个，推导最大熵模型。两条路线最终得到相同类型的对数线性模型，并把学习归结为光滑凸目标的最小化。

本章对应教材第6章，保留其“逻辑斯谛回归模型、最大熵模型、模型学习的最优化算法”三部分及主要小节。阅读需要条件概率、微分、矩阵运算和拉格朗日乘子的基础；涉及熵、凸性与对偶性的关键步骤在使用时说明。本章采用自然对数，$\mathbb R$ 表示实数域，$\mathbf 1\{A\}$ 表示事件或条件 $A$ 成立时取 $1$、否则取 $0$ 的示性函数，$\|\cdot\|_2$ 表示欧氏范数。

\section{逻辑斯谛回归模型}
\label{sec:lh-logistic}

\subsection{逻辑斯谛分布}
\label{subsec:lh-logistic-distribution}

\begin{definition}[逻辑斯谛分布]
\label{def:lh-logistic-distribution}
若连续随机变量 $Z$ 的分布函数为
\begin{equation}
F(z)=P(Z\leq z)=\frac{1}{1+\exp(-(z-\mu)/\gamma)},
\qquad z\in\mathbb R,
\label{eq:lh-logistic-cdf}
\end{equation}
则称 $Z$ 服从\emph{逻辑斯谛分布}（logistic distribution），其中 $\mu\in\mathbb R$ 为位置参数，$\gamma>0$ 为尺度参数。其概率密度为
\begin{equation}
f_Z(z)=F'(z)=\frac{\exp(-(z-\mu)/\gamma)}
{\gamma\bigl(1+\exp(-(z-\mu)/\gamma)\bigr)^2}.
\label{eq:lh-logistic-density}
\end{equation}
\end{definition}

令 $\sigma(t)=1/(1+e^{-t})$，称其为逻辑斯谛函数；它是一种 S 形函数（sigmoid function），而“S 形函数”并不专指这一函数。由定义可得
\begin{equation}
\sigma(-t)=1-\sigma(t),\qquad
\sigma'(t)=\sigma(t)\bigl(1-\sigma(t)\bigr)>0.
\label{eq:lh-sigmoid-identities}
\end{equation}
因此 $F(\mu+t)+F(\mu-t)=1$，分布函数的图像关于 $(\mu,1/2)$ 中心对称，密度关于 $z=\mu$ 对称。分布函数在中心的斜率为 $F'(\mu)=1/(4\gamma)$，且两端分别趋于 $0$ 和 $1$；$\gamma$ 越小，中心附近的变化越陡。教材将 $\gamma$ 称为形状参数，从位置尺度族的角度看，它直接控制横向尺度。

\subsection{二项逻辑斯谛回归模型}
\label{subsec:lh-binary-logistic}

设输入随机向量 $X$ 的取值 $x\in\mathbb R^d$，输出 $Y\in\{0,1\}$。线性得分 $\beta^{\mathsf T}x+b$ 可以取任意实数，不能直接作为概率；用 $\sigma$ 将其映射到 $(0,1)$，便得到二项逻辑斯谛回归。

\begin{definition}[二项逻辑斯谛回归]
\label{def:lh-binary-logistic}
\emph{二项逻辑斯谛回归模型}（binomial logistic regression model）是条件分布
\begin{equation}
P_{\beta,b}(Y=1\mid X=x)=\sigma(\beta^{\mathsf T}x+b),\qquad
P_{\beta,b}(Y=0\mid X=x)=1-\sigma(\beta^{\mathsf T}x+b),
\label{eq:lh-binary-logistic}
\end{equation}
其中 $\beta\in\mathbb R^d$ 为权值向量，$b\in\mathbb R$ 为偏置。以下令 $\widetilde x=(x^{\mathsf T},1)^{\mathsf T}\in\mathbb R^{d+1}$、$w=(\beta^{\mathsf T},b)^{\mathsf T}\in\mathbb R^{d+1}$，统一写成 $p_w(x)=\sigma(w^{\mathsf T}\widetilde x)$。
\end{definition}

这里回归的是概率的变换，最终任务仍为分类。对于 $0<p<1$，\emph{几率}（odds）为 $p/(1-p)$，\emph{对数几率}（log odds）又称 logit，为 $\operatorname{logit}(p)=\log(p/(1-p))$。将式~\eqref{eq:lh-binary-logistic} 代入可得
\begin{equation}
\frac{p_w(x)}{1-p_w(x)}=\exp(w^{\mathsf T}\widetilde x),\qquad
\operatorname{logit}(p_w(x))=w^{\mathsf T}\widetilde x.
\label{eq:lh-log-odds}
\end{equation}
所以模型的核心假设是\emph{对数几率为输入的仿射函数}。若其他坐标固定，$x_j$ 增加一个单位使几率乘以 $e^{\beta_j}$，而不是使概率增加固定数值。概率变化还取决于原有得分，因为 $\partial p_w(x)/\partial x_j=\beta_jp_w(x)(1-p_w(x))$。

在两类误判代价相同、以误分类概率为损失时，选择条件概率较大的类别；约定相等时选择 $1$，便有
\begin{equation}
\widehat y(x)=\mathbf 1\{p_w(x)\geq 1/2\}
=\mathbf 1\{\beta^{\mathsf T}x+b\geq 0\}.
\label{eq:lh-logistic-decision}
\end{equation}
当 $\beta\ne0$ 时，决策边界为超平面。若另取概率阈值 $\tau\in(0,1)$，对应的得分阈值为 $\log(\tau/(1-\tau))$；因此预测概率与选定分类阈值是两个不同环节。

\begin{example}
\label{ex:lh-logistic-odds}
设一维模型为 $p(x)=\sigma((\log 2)x)$。求 $x=0,1,2$ 时的概率与几率，并说明两者的变化。
\end{example}
\begin{solve}
几率为 $e^{(\log2)x}=2^x$，故三处几率分别为 $1,2,4$，概率分别为 $1/2,2/3,4/5$。输入每增加一个单位，几率都翻倍，而概率增量依次为 $1/6$ 与 $2/15$。这说明参数具有对数几率上的线性解释，概率本身随输入非线性变化。
\end{solve}

\begin{notealert}[type=caution,title={分布与条件模型的区别}]
式~\eqref{eq:lh-logistic-cdf} 描述连续随机变量的分布，式~\eqref{eq:lh-binary-logistic} 描述给定输入后离散标签的条件分布。后者只使用前者的函数形状，并不要求输入 $X$ 服从逻辑斯谛分布，也没有把二值标签当作连续随机变量。
\end{notealert}

\subsection{模型参数估计}
\label{subsec:lh-logistic-estimation}

给定观测数据 $T=\{(x_i,y_i)\}_{i=1}^N$，其中 $x_i\in\mathbb R^d$、$y_i\in\{0,1\}$。若样本独立同分布，则给定输入后标签的条件似然可以分解；也可以直接以给定输入后的标签相互独立为条件建模假设。记 $p_i=p_w(x_i)$，条件似然与对数似然为
\begin{align}
\mathcal L(w)&=\prod_{i=1}^N p_i^{y_i}(1-p_i)^{1-y_i},\notag\\
\ell_{\mathrm{bin}}(w)&=\sum_{i=1}^N\bigl[y_i\log p_i+(1-y_i)\log(1-p_i)\bigr]\notag\\
&=\sum_{i=1}^N\bigl[y_iw^{\mathsf T}\widetilde x_i
-\log(1+\exp(w^{\mathsf T}\widetilde x_i))\bigr].
\label{eq:lh-logistic-loglik}
\end{align}
极大化 $\ell_{\mathrm{bin}}$ 等价于最小化平均负对数似然，即二元交叉熵（binary cross-entropy）
\begin{equation}
J_{\mathrm{bin}}(w)=-\frac1N\ell_{\mathrm{bin}}(w)
=\frac1N\sum_{i=1}^N\bigl[\log(1+e^{w^{\mathsf T}\widetilde x_i})
-y_iw^{\mathsf T}\widetilde x_i\bigr].
\label{eq:lh-logistic-objective}
\end{equation}
这一步把概率估计转化为参数优化；求和与平均只相差正常数，最优解相同，但梯度尺度和学习率相应不同。

\begin{proposition}[梯度与凸性]
\label{prop:lh-logistic-convexity}
对有限训练集和任意有限 $w$，$J_{\mathrm{bin}}$ 光滑，且
\begin{align}
\nabla J_{\mathrm{bin}}(w)&=\frac1N\sum_{i=1}^N(p_i-y_i)\widetilde x_i,
\label{eq:lh-logistic-gradient}\\
\nabla^2J_{\mathrm{bin}}(w)&=\frac1N\sum_{i=1}^Np_i(1-p_i)
\widetilde x_i\widetilde x_i^{\mathsf T}\succeq0.
\label{eq:lh-logistic-hessian}
\end{align}
因此目标为凸函数，每个驻点都是全局最小点。若增广设计矩阵 $A=(\widetilde x_1,\ldots,\widetilde x_N)^{\mathsf T}$ 列满秩，则目标严格凸，有限最小点若存在便唯一。
\end{proposition}
\begin{proof}
对单个得分 $z_i=w^{\mathsf T}\widetilde x_i$，有 $\frac{\mathrm d}{\mathrm dz_i}\log(1+e^{z_i})=p_i$，再用式~\eqref{eq:lh-sigmoid-identities} 和链式法则即得梯度与海森矩阵。对任意 $v\in\mathbb R^{d+1}$，
\[
v^{\mathsf T}\nabla^2J_{\mathrm{bin}}(w)v
=\frac1N\sum_{i=1}^Np_i(1-p_i)(v^{\mathsf T}\widetilde x_i)^2\geq0.
\]
有限参数下各 $p_i(1-p_i)>0$。若 $A$ 列满秩，非零 $v$ 不可能同时满足所有 $v^{\mathsf T}\widetilde x_i=0$，故上式严格为正，得到严格凸性。
\end{proof}

\emph[caution]{凸性不保证有限最优参数一定存在。}若存在 $v$ 使所有 $(2y_i-1)v^{\mathsf T}\widetilde x_i>0$，称训练集严格线性可分。取 $w=tv$、$t\to+\infty$，每项损失均趋于 $0$，但有限 $w$ 给出的概率严格介于 $0$ 与 $1$，损失仍为正，因而只有下确界而没有有限最小点。例如 $(x_1,y_1)=(-1,0)$、$(x_2,y_2)=(1,1)$，取偏置为 $0$、斜率为 $t$，有 $J_{\mathrm{bin}}=\log(1+e^{-t})\downarrow0$。这时设计矩阵仍列满秩，足见严格凸性与最小值的存在性必须分别讨论。

一种补充做法是加入二次正则化，求解
\begin{equation}
\min_{w\in\mathbb R^{d+1}}
J_{\mathrm{bin},\lambda}(w)=J_{\mathrm{bin}}(w)+\frac\lambda2\|w\|_2^2,
\qquad\lambda>0.
\label{eq:lh-logistic-regularized}
\end{equation}
此处对包括偏置在内的全部参数惩罚，目标连续、强凸，且当 $\|w\|_2\to\infty$ 时趋于无穷，因此存在唯一有限最小点。若实现中不惩罚偏置，则不能直接照用这个强凸性与存在性的结论，需要另查数据和未惩罚方向的条件。

由式~\eqref{eq:lh-logistic-gradient} 立即得到教材习题所需的批量梯度下降算法：输入 $T$、初值 $w^{(0)}$、容差 $\varepsilon>0$ 与迭代上限，每轮计算全部 $p_i$ 和梯度；梯度范数不超过 $\varepsilon$ 时返回当前参数，否则取 $\eta_k>0$，更新
\begin{equation}
w^{(k+1)}=w^{(k)}-\frac{\eta_k}{N}
\sum_{i=1}^N\bigl(p_{w^{(k)}}(x_i)-y_i\bigr)\widetilde x_i.
\label{eq:lh-logistic-gd}
\end{equation}
达到迭代上限则返回近似解并注明未满足停止条件。因为 $p_i(1-p_i)\leq1/4$，梯度的一个利普希茨常数为 $L_g=\lambda_{\max}(A^{\mathsf T}A)/(4N)$；$\lambda_{\max}$ 表示最大特征值。当有限最小点存在时，固定步长 $0<\eta\leq1/L_g$ 可用于标准凸梯度下降。对式~\eqref{eq:lh-logistic-regularized}，梯度加 $\lambda w$，常数可取 $L_g+\lambda$。单轮稠密计算为 $O(N(d+1))$；这些结论针对批量精确梯度，不直接给出随机梯度方法的保证。

实际计算不宜先得到接近 $0$ 或 $1$ 的概率再取对数，而应利用
\[
\log(1+e^z)=\max\{0,z\}+\log(1+e^{-|z|})
\]
求损失，并按 $z$ 的正负选用等价的 sigmoid 表达式，避免指数溢出。数值稳定的计算与优化目标的数学定义相互对应，但有限迭代输出仍是参数的数值近似。

\subsection{多项逻辑斯谛回归}
\label{subsec:lh-multinomial-logistic}

当 $Y\in\{1,\ldots,K\}$、$K\geq2$ 时，可为每类设置得分 $a_k(x)=w_k^{\mathsf T}\widetilde x$，并定义\emph{多项逻辑斯谛回归}（multinomial logistic regression）
\begin{equation}
P_W(Y=k\mid X=x)=\frac{e^{a_k(x)}}{\sum_{r=1}^Ke^{a_r(x)}},
\qquad k=1,\ldots,K.
\label{eq:lh-softmax}
\end{equation}
这里 $W=(w_1,\ldots,w_K)$ 为参数组，右端称为 softmax。所有得分同时加上同一个数不改变概率，特别是所有 $w_k$ 同时加上相同向量不会改变模型。因此可选第 $K$ 类为基准，固定 $w_K=0$，得到教材形式
\begin{equation}
\begin{aligned}
P_W(Y=k\mid X=x)&=\frac{e^{w_k^{\mathsf T}\widetilde x}}
{1+\sum_{r=1}^{K-1}e^{w_r^{\mathsf T}\widetilde x}},\quad k<K,\\
P_W(Y=K\mid X=x)&=\frac1{1+\sum_{r=1}^{K-1}e^{w_r^{\mathsf T}\widetilde x}}.
\end{aligned}
\label{eq:lh-baseline-softmax}
\end{equation}
相对基准类的对数几率满足 $\log(P_W(Y=k\mid x)/P_W(Y=K\mid x))=w_k^{\mathsf T}\widetilde x$；当只有两个类别时，它退化为二项模型。

极大似然估计等价于最小化多类交叉熵
\begin{equation}
J_{\mathrm{multi}}(W)=\frac1N\sum_{i=1}^N
\left[\log\left(\sum_{r=1}^Ke^{w_r^{\mathsf T}\widetilde x_i}\right)
-w_{y_i}^{\mathsf T}\widetilde x_i\right],\qquad w_K=0.
\label{eq:lh-multiclass-objective}
\end{equation}
对 $k<K$，梯度为 $N^{-1}\sum_i(P_W(Y=k\mid x_i)-\mathbf1\{y_i=k\})\widetilde x_i$。其凸性将在最大熵模型的统一推导中得到解释；选定基准去除了公共平移造成的冗余，但其他特征冗余或可分性问题仍可能存在。计算归一化因子时可令 $a_{\max}=\max_r a_r$，使用 $\log\sum_re^{a_r}=a_{\max}+\log\sum_re^{a_r-a_{\max}}$。

\section{最大熵模型}
\label{sec:lh-maxent}

\subsection{最大熵原理}
\label{subsec:lh-maxent-principle}

仅有部分统计信息时，通常有许多分布与这些信息相容。\emph{最大熵原理}（maximum entropy principle）是在预先指定的样本空间和约束下，选择熵最大的分布；它把已知信息写成约束，同时避免额外加入没有依据的偏好。这是一项建模准则，并不宣称最大熵分布必然等于未知的真实分布，也不自动保证分类准确率最高。

若离散随机变量 $Y$ 的取值集为有限集合 $\mathcal Y$，$|\mathcal Y|=K$，其分布为 $p$，则\emph{熵}（entropy）为
\begin{equation}
H(p)=-\sum_{y\in\mathcal Y}p(y)\log p(y),\qquad 0\log0:=0.
\label{eq:lh-entropy}
\end{equation}
它满足 $0\leq H(p)\leq\log K$，右侧等号当且仅当 $p(y)=1/K$ 对所有 $y$ 成立。为核查上界，令均匀分布为 $u(y)=1/K$，则相对熵
\[
D(p\|u)=\sum_{y:p(y)>0}p(y)\log\frac{p(y)}{u(y)}
=\log K-H(p)\geq0.
\]
这里相对熵的一般非负性可由 $-\log t\geq1-t$ 得到：若 $r(y)>0$，则 $D(p\|r)\geq\sum_{y:p(y)>0}(p(y)-r(y))\geq0$，等号仅当 $p=r$ 成立。熵的下界则来自各项 $-p(y)\log p(y)\geq0$。

概率向量 $(p_1,\ldots,p_K)$ 位于概率单纯形 $\{p_k\geq0,\ \sum_kp_k=1\}$ 中；线性统计约束把它截成凸可行集，熵在这个集合上选出一个分布。无约束时均匀分布具有最大熵，有约束时一般不再全局均匀；只有在约束无法区分的状态之间，才有对称性带来的等概率关系。

\begin{example}
\label{ex:lh-five-outcomes}
设 $Y$ 有五个可能取值 $\{A,B,C,D,E\}$。先在只有归一化约束时求最大熵分布，再增加约束 $p(A)+p(B)=3/10$，求新的最大熵分布。这对应教材的五状态例子。
\end{example}
\begin{solve}
只有归一化约束时，上述熵上界给出五个概率均为 $1/5$。增加分组质量约束后，记 $G=\{A,B\}$、$G^{\mathrm c}=\{C,D,E\}$；两组质量分别为 $3/10$ 和 $7/10$。用条件概率展开熵可得
\[
H(p)=h(3/10)+\frac3{10}H(p(\cdot\mid G))
+\frac7{10}H(p(\cdot\mid G^{\mathrm c})),
\]
其中 $h(t)=-t\log t-(1-t)\log(1-t)$。第一项固定，后两项分别在组内均匀时最大，故唯一解为
\begin{equation}
p(A)=p(B)=\frac3{20},\qquad
p(C)=p(D)=p(E)=\frac7{30}.
\label{eq:lh-five-outcome-solution}
\end{equation}
组内等概率来自约束的对称性，两组之间并不等概率。
\end{solve}

若进一步加入教材提出的 $p(A)+p(C)=1/2$，原来的分组对称性就被打破，不能继续沿用式~\eqref{eq:lh-five-outcome-solution}。作为补充推导，令 $a=p(A)$，则 $p(B)=3/10-a$、$p(C)=1/2-a$、$p(D)+p(E)=1/5+a$；给定 $a$ 时，熵在 $p(D)=p(E)=1/10+a/2$ 处最大。把它代回熵，在 $0<a<3/10$ 上求导得
\[
\frac{\mathrm dH}{\mathrm da}
=\log\frac{(3/10-a)(1/2-a)}{a(1/10+a/2)}.
\]
导数为零等价于 $a^2-(9/5)a+3/10=0$，区间内唯一根为 $a=(9-\sqrt{51})/10$。熵沿这条非平凡仿射路径严格凹，因此该根给出新增约束下的唯一最大熵解。这也说明重叠约束必须联合处理。

\subsection{最大熵模型的定义}
\label{subsec:lh-maxent-definition}

回到分类问题，设 $\mathcal Y=\{1,\ldots,K\}$，训练集为 $T=\{(x_i,y_i)\}_{i=1}^N$。令 $\mathcal S=\{x_i:1\leq i\leq N\}$ 表示训练集中不同输入组成的有限集合。经验联合分布与经验输入边缘分布分别为
\begin{equation}
\widehat p(x,y)=\frac1N\sum_{i=1}^N\mathbf1\{x_i=x,\ y_i=y\},\qquad
\widehat p(x)=\sum_{y\in\mathcal Y}\widehat p(x,y),\quad x\in\mathcal S.
\label{eq:lh-empirical-distributions}
\end{equation}
即使原输入空间连续，有限样本的经验分布仍只在 $\mathcal S$ 上赋予质量。以下未标明范围的 $\sum_{x,y}$ 均指 $x\in\mathcal S$、$y\in\mathcal Y$ 上的有限求和，且 $\widehat p(x)>0$。

\emph{特征函数}（feature function）$f_j(x,y)$ 用来描述输入与候选输出之间的统计关系。例如 $f_j(x,y)=\mathbf1\{x\in A_j,\ y=k_j\}$ 表示输入落在某个预先指定集合 $A_j$ 且输出为 $k_j$。教材首先采用二值特征，模型推导可推广到任意有限实值特征；但后文的 IIS 推导另需非负性。给定 $m$ 个特征，记 $f(x,y)=(f_1(x,y),\ldots,f_m(x,y))^{\mathsf T}$，经验特征期望为
\begin{equation}
c_j=\widehat{\mathbb E}[f_j]
=\sum_{x,y}\widehat p(x,y)f_j(x,y)
=\frac1N\sum_{i=1}^Nf_j(x_i,y_i).
\label{eq:lh-empirical-moments}
\end{equation}
若候选条件分布为 $q(y\mid x)$，则模型特征期望是
\begin{equation}
\mathbb E_q[f_j]=\sum_{x,y}\widehat p(x)q(y\mid x)f_j(x,y).
\label{eq:lh-model-moments}
\end{equation}
两个期望分别使用经验联合分布与联合分布 $\widehat p(x)q(y\mid x)$。条件模型不另外学习输入的边缘分布，而是把经验边缘固定下来。约束 $\mathbb E_q[f_j]=c_j$ 要求模型重现选定特征的样本均值；这是训练时采用的匹配条件，不是关于总体真实期望的无误差断言。

\begin{definition}[条件最大熵模型]
\label{def:lh-conditional-maxent}
令 $\mathcal P$ 为 $\mathcal S$ 上全部条件分布组成的集合，即每个 $x$ 均满足 $q(y\mid x)\geq0$ 和 $\sum_yq(y\mid x)=1$。定义
\begin{equation}
\mathcal C=\{q\in\mathcal P:\mathbb E_q[f_j]=c_j,\ j=1,\ldots,m\},\qquad
H(q)=-\sum_{x,y}\widehat p(x)q(y\mid x)\log q(y\mid x).
\label{eq:lh-maxent-feasible-set}
\end{equation}
在 $\mathcal C$ 中使条件熵 $H(q)$ 最大的分布 $q^*$ 称为\emph{条件最大熵模型}（conditional maximum entropy model）。
\end{definition}

这里的条件熵始终相对于输入分布 $\widehat p(x)$ 计算。经验条件分布 $\widehat p(y\mid x)=\widehat p(x,y)/\widehat p(x)$ 本身满足全部经验特征约束，所以 $\mathcal C$ 非空。它是有限个概率单纯形的乘积与仿射约束的交，因而紧且凸；熵连续并在这个条件分布空间上严格凹，因此最大熵分布在 $\mathcal S$ 上存在且唯一。这个结论不要求最优概率都为正，也不等于参数表示唯一。对训练中未出现的输入，原始经验熵目标没有施加约束，预测依靠随后选定的参数化特征模型作延拓。

\subsection{最大熵模型的学习}
\label{subsec:lh-maxent-learning}

为推导模型形式，先将最大熵问题改写成凸最小化问题：
\begin{equation}
\begin{aligned}
\min_{q}\quad&-H(q)=\sum_{x,y}\widehat p(x)q(y\mid x)\log q(y\mid x),\\
\text{满足}\quad&\mathbb E_q[f_j]=c_j,\quad j=1,\ldots,m,\\
&\sum_yq(y\mid x)=1,\quad x\in\mathcal S,\qquad q(y\mid x)\geq0.
\end{aligned}
\label{eq:lh-maxent-primal}
\end{equation}
以下先讨论存在严格正可行分布的情形。在有限空间上，这给出使用有限拉格朗日乘子和内部驻点推导的充分条件。引入特征约束乘子 $w\in\mathbb R^m$，并为\emph{每个输入的归一化约束}分别引入 $\alpha_x\in\mathbb R$，定义
\begin{equation}
\begin{aligned}
\mathscr L(q,w,\alpha)={}&\sum_{x,y}\widehat p(x)q(y\mid x)\log q(y\mid x)\\
&+\sum_{j=1}^mw_j\bigl(c_j-\mathbb E_q[f_j]\bigr)
+\sum_x\widehat p(x)\alpha_x\left(1-\sum_yq(y\mid x)\right).
\end{aligned}
\label{eq:lh-maxent-lagrangian}
\end{equation}
因为各 $\widehat p(x)>0$，在归一化乘子前加入这个权重不会改变约束。对单个坐标 $q(y\mid x)>0$ 求偏导，得到
\[
\frac{\partial\mathscr L}{\partial q(y\mid x)}
=\widehat p(x)\bigl[\log q(y\mid x)+1-w^{\mathsf T}f(x,y)-\alpha_x\bigr].
\]
令其为零，再逐个 $x$ 归一化，得到 $\alpha_x=1-\log Z_w(x)$ 和
\begin{equation}
p_w(y\mid x)=\frac{\exp(w^{\mathsf T}f(x,y))}{Z_w(x)},\qquad
Z_w(x)=\sum_{y\in\mathcal Y}\exp(w^{\mathsf T}f(x,y)).
\label{eq:lh-maxent-model}
\end{equation}
$Z_w(x)$ 称为\emph{规范化因子}或配分函数（partition function），它一般随输入变化。模型称为\emph{对数线性模型}（log-linear model），因为对数概率由特征的线性组合减去对数规范化因子构成。任意参数都给出合法分布，但只有满足所需最优性条件的参数才对应当前约束下的最大熵解。

为明确对偶问题的定义，先消去归一化约束，令 $q\in\mathcal P$，只对特征等式作拉格朗日松弛。其对偶函数为
\begin{equation}
\begin{aligned}
\Psi(w)&=\inf_{q\in\mathcal P}
\left[-H(q)+w^{\mathsf T}\bigl(c-\mathbb E_q[f]\bigr)\right]\\
&=w^{\mathsf T}c-\sum_x\widehat p(x)\log Z_w(x),
\qquad c=(c_1,\ldots,c_m)^{\mathsf T}.
\end{aligned}
\label{eq:lh-maxent-dual}
\end{equation}
第二步可直接核查：对任意 $q\in\mathcal P$，方括号内的量等于
\[
w^{\mathsf T}c-\sum_x\widehat p(x)\log Z_w(x)
+\sum_x\widehat p(x)D\bigl(q(\cdot\mid x)\|p_w(\cdot\mid x)\bigr),
\]
最后一项非负且在 $q=p_w$ 时为零。\emph[caution]{求对偶函数的内层极小值时，只保留概率归一化与非负约束，不能同时把已松弛的特征匹配约束留在可行域中。}

因此对偶问题是 $\max_{w\in\mathbb R^m}\Psi(w)$。在上述严格可行条件下，强对偶性成立且存在有限最优乘子；也可以用下面的相对熵恒等式直接验证候选解的最优性，而不把“原始目标凸”单独当作强对偶的充分说明。

\begin{proposition}[最大熵解的相对熵验证]
\label{prop:lh-maxent-certificate}
若某个有限参数 $w^*$ 使 $p_{w^*}$ 满足全部特征约束，则它是式~\eqref{eq:lh-maxent-primal} 的唯一最优条件分布。对任意 $q\in\mathcal C$，有
\begin{equation}
H(q)=H(p_{w^*})-\sum_x\widehat p(x)
D\bigl(q(\cdot\mid x)\|p_{w^*}(\cdot\mid x)\bigr).
\label{eq:lh-maxent-kl-certificate}
\end{equation}
\end{proposition}
\begin{proof}
由 $\log p_w(y\mid x)=w^{\mathsf T}f(x,y)-\log Z_w(x)$ 和可行性，
\[
\sum_x\widehat p(x)D(q\|p_w)
=-H(q)-w^{\mathsf T}c+\sum_x\widehat p(x)\log Z_w(x)
=-H(q)-\Psi(w).
\]
这里 $D(q\|p_w)$ 的两个参数均指对应输入处的条件分布。再取 $w=w^*$、$q=p_{w^*}$，得到 $\Psi(w^*)=-H(p_{w^*})$；代回即为所需恒等式。相对熵的非负性给出最大性，等号条件和各 $\widehat p(x)>0$ 给出唯一性。
\end{proof}

\begin{example}
\label{ex:lh-five-outcomes-dual}
用式~\eqref{eq:lh-maxent-model} 重新求解例~\ref{ex:lh-five-outcomes} 中带分组质量约束的情形。
\end{example}
\begin{solve}
把输入视为只有一个固定值，取唯一特征 $f(y)=\mathbf1\{y\in\{A,B\}\}$，经验期望为 $c=3/10$。模型为
\[
p_w(A)=p_w(B)=\frac{e^w}{2e^w+3},\qquad
p_w(C)=p_w(D)=p_w(E)=\frac1{2e^w+3}.
\]
对偶函数为 $\Psi(w)=(3/10)w-\log(2e^w+3)$。令 $\Psi'(w)=3/10-2e^w/(2e^w+3)=0$，解得 $e^{w^*}=9/14$。代入即得到式~\eqref{eq:lh-five-outcome-solution}，且模型特征期望确为 $3/10$，所以命题~\ref{prop:lh-maxent-certificate} 保证它为唯一最大熵分布。
\end{solve}

若约束迫使某些最优概率为零，式~\eqref{eq:lh-maxent-model} 的有限参数表示可能不存在，因为指数形式在有限参数下总为正。此时原始问题的最大熵分布仍存在，但相应参数可能只能沿发散序列逼近；“分布存在”与“有限最优参数存在”应当区分。另一方面，若不同权值产生同一条件分布，唯一的最大熵分布也不会带来唯一权值。

\subsection{极大似然估计}
\label{subsec:lh-maxent-mle}

对式~\eqref{eq:lh-maxent-model} 作条件极大似然估计，平均对数似然为
\begin{equation}
\begin{aligned}
\ell(w)&=\frac1N\sum_{i=1}^N\log p_w(y_i\mid x_i)
=\sum_{x,y}\widehat p(x,y)\log p_w(y\mid x)\\
&=w^{\mathsf T}c-\sum_x\widehat p(x)\log Z_w(x)=\Psi(w).
\end{aligned}
\label{eq:lh-maxent-loglik}
\end{equation}
因此最大熵问题的对偶优化与这一对数线性条件模型的极大似然优化具有完全相同的目标。这里取的是平均对数似然，通常的样本对数似然为 $N\ell(w)$。以下统一最小化 $J(w)=-\ell(w)$，避免在算法中混淆极大化与极小化的梯度方向。

对有限求和逐项求导，先得
\[
\frac{\partial\log Z_w(x)}{\partial w_j}
=\sum_y p_w(y\mid x)f_j(x,y),
\]
从而
\begin{equation}
\frac{\partial J(w)}{\partial w_j}=\mathbb E_{p_w}[f_j]-c_j,
\qquad
\nabla^2J(w)=\sum_x\widehat p(x)
\operatorname{Cov}_{p_w(\cdot\mid x)}\bigl(f(x,Y)\bigr)\succeq0.
\label{eq:lh-maxent-gradient-hessian}
\end{equation}
这里 $\operatorname{Cov}$ 是协方差矩阵；给定 $x$ 时，随机性仅来自 $Y\sim p_w(\cdot\mid x)$。对任意 $v\in\mathbb R^m$，海森矩阵的二次型为 $\sum_x\widehat p(x)\operatorname{Var}_{p_w(\cdot\mid x)}(v^{\mathsf T}f(x,Y))\geq0$，因而 $J$ 凸、$\ell$ 凹。有限无约束最优参数满足 $\mathbb E_{p_{w^*}}[f_j]=c_j$，即梯度为零正好等价于特征期望匹配。若存在非零 $v$，使 $v^{\mathsf T}f(x,y)$ 对每个固定 $x$ 都不随 $y$ 变化，沿该方向改变权值不改变模型，目标就不会严格凸。

逻辑斯谛回归是这个框架的特例：二分类时取 $f_j(x,y)=y\widetilde x_j$，$j=1,\ldots,d+1$，便有 $Z_w(x)=1+e^{w^{\mathsf T}\widetilde x}$；多分类时取 $f_{k,j}(x,y)=\mathbf1\{y=k\}\widetilde x_j$，$k<K$，便得到基准类 softmax。这里逻辑斯谛回归的特征可以为负，仍适用式~\eqref{eq:lh-maxent-gradient-hessian} 的凸性结论，但不能不加检查地代入后文依赖非负特征的 IIS 公式。

若改为最小化 $J_\lambda(w)=J(w)+(\lambda/2)\|w\|_2^2$、$\lambda>0$，则目标存在唯一有限最小点，梯度为 $\mathbb E_{p_w}[f]-c+\lambda w$。此时最优性条件为 $\mathbb E_{p_w}[f]-c=-\lambda w$，一般不再精确匹配经验矩；正则化估计因此不能直接称为原来硬约束问题的同一个解。

\begin{remark}
\label{rem:lh-exponential-family}
对固定输入 $x$，式~\eqref{eq:lh-maxent-model} 的形式为 $\exp(w^{\mathsf T}f(x,y)-\log Z_w(x))$，属于以 $f(x,y)$ 为充分统计量的条件指数族。特别地，二项逻辑斯谛回归的标签分布为
\[
p(y\mid x)=\exp\bigl(y\theta-\log(1+e^\theta)\bigr),\qquad
y\in\{0,1\},\quad\theta=w^{\mathsf T}\widetilde x,
\]
这验证的是伯努利条件分布的指数族表示。教材习题中关于“逻辑斯谛分布”的表述需要先区分对象；上述计算并未证明定义~\ref{def:lh-logistic-distribution} 中连续位置尺度分布族的有限维指数族性质，不能把两种分布混为一谈。
\end{remark}

\section{模型学习的最优化算法}
\label{sec:lh-maxent-optimization}

本节继续使用式~\eqref{eq:lh-maxent-loglik} 的平均对数似然 $\ell(w)$ 及负对数似然 $J(w)=-\ell(w)$。IIS 直接构造使 $\ell$ 不下降的更新，拟牛顿法则最小化 $J$。两者优化的是同一模型的同一估计准则，但使用的特征条件和每轮计算方式不同；不能把算法名称当作另一种概率模型。

\subsection{改进的迭代尺度法}
\label{subsec:lh-iis}

\emph{改进的迭代尺度法}（improved iterative scaling，IIS）的思路是构造对数似然改变量的可分离下界，从而把一个向量更新拆成若干一元方程。以下假设 $f_j(x,y)\geq0$，教材的二值特征是其特例。令
\[
\mu_w(x,y)=\widehat p(x)p_w(y\mid x),\qquad
f^\#(x,y)=\sum_{j=1}^mf_j(x,y).
\]
给定当前参数 $w$ 和增量 $\delta\in\mathbb R^m$，由规范化因子的比值可得
\begin{equation}
\ell(w+\delta)-\ell(w)
=\delta^{\mathsf T}c-\sum_x\widehat p(x)
\log\left(\sum_yp_w(y\mid x)e^{\delta^{\mathsf T}f(x,y)}\right).
\label{eq:lh-iis-difference}
\end{equation}
对每个输入使用 $-\log a\geq1-a$，并利用 $\sum_x\widehat p(x)=1$，得到第一层下界
\begin{equation}
\ell(w+\delta)-\ell(w)\geq A(\delta\mid w)
:=\delta^{\mathsf T}c+1-\sum_{x,y}\mu_w(x,y)e^{\delta^{\mathsf T}f(x,y)}.
\label{eq:lh-iis-first-bound}
\end{equation}
它仍然耦合各个 $\delta_j$。若 $f^\#(x,y)>0$，则 $f_j/f^\#$ 是非负且和为 $1$ 的权值，指数函数的凸性给出 Jensen 不等式
\begin{equation}
e^{\delta^{\mathsf T}f(x,y)}
=\exp\left(\sum_j\frac{f_j(x,y)}{f^\#(x,y)}\delta_jf^\#(x,y)\right)
\leq\sum_j\frac{f_j(x,y)}{f^\#(x,y)}e^{\delta_jf^\#(x,y)}.
\label{eq:lh-iis-jensen}
\end{equation}
若 $f^\#(x,y)=0$，非负性保证所有 $f_j(x,y)=0$，相应指数项恒为 $1$，应单独保留，不能计算 $0/0$。把这些常数项抵消后，得到对零特征总和也有定义的第二层下界
\begin{equation}
\begin{aligned}
B(\delta\mid w):={}&\sum_{j=1}^m\delta_jc_j\\
&+\sum_{\substack{x,y\\f^\#(x,y)>0}}\mu_w(x,y)
\sum_{j=1}^m\frac{f_j(x,y)}{f^\#(x,y)}
\left(1-e^{\delta_jf^\#(x,y)}\right),
\end{aligned}
\label{eq:lh-iis-separable-bound}
\end{equation}
满足 $\ell(w+\delta)-\ell(w)\geq A(\delta\mid w)\geq B(\delta\mid w)$，且 $B(0\mid w)=0$。这个下界可以按 $\delta_j$ 分开最大化，因为每一项只含一个增量坐标。

求偏导并将零特征项重新纳入求和，得到
\begin{align}
\frac{\partial B}{\partial\delta_j}
&=c_j-\sum_{x,y}\mu_w(x,y)f_j(x,y)e^{\delta_jf^\#(x,y)},\notag\\
\frac{\partial^2 B}{\partial\delta_j^2}
&=-\sum_{x,y}\mu_w(x,y)f_j(x,y)f^\#(x,y)
e^{\delta_jf^\#(x,y)}\leq0.
\label{eq:lh-iis-bound-derivatives}
\end{align}
因此每个非平凡坐标的最大化归结为一元方程
\begin{equation}
\sum_{x,y}\widehat p(x)p_w(y\mid x)f_j(x,y)
e^{\delta_jf^\#(x,y)}=c_j.
\label{eq:lh-iis-update-equation}
\end{equation}
当 $c_j>0$ 时，特征在训练支持上必非恒零；对有限 $w$，左端随 $\delta_j$ 严格递增，从 $0$ 趋于无穷，故存在唯一有限根。若所有坐标均满足这一条件，各自取最大化点便有 $B(\delta\mid w)\geq B(0\mid w)=0$，从而保证一次完整更新不降低对数似然。

\textbf{IIS 的迭代步骤。}输入经验分布、有限非负特征、容差 $\varepsilon>0$ 与最大迭代次数；先删除在训练支持与全部候选标签上恒为零的特征，并检查剩余特征是否存在零经验期望。对于每个剩余 $c_j>0$ 的通常情形，算法如下。
\begin{enumerate}
\item 初始化 $w^{(0)}=0$，计算经验期望 $c$。此时每个输入处的初始条件分布均为 $1/K$。
\item 在第 $k$ 轮固定 $p_{w^{(k)}}$，对每个 $j$ 求解式~\eqref{eq:lh-iis-update-equation}，得到 $\delta_j$。本轮所有一元方程均使用同一轮的旧分布。
\item 同时更新 $w^{(k+1)}=w^{(k)}+\delta$，重新计算模型概率和特征期望。若 $\|\mathbb E_{p_{w^{(k+1)}}}[f]-c\|_2\leq\varepsilon$，返回当前参数和模型；否则进入下一轮，直到迭代上限。
\end{enumerate}
这里把一轮中各增量求完后再统一更新，与可分离下界的推导直接对应。数值近似求根时应控制方程残差，并检查目标是否确有改善；不能把精确求根下的单调性无条件转用于任意近似更新。若 $c_j=0$ 但 $f_j$ 在某些候选输入输出对上为正，则式~\eqref{eq:lh-iis-update-equation} 没有有限根；相应约束迫使这些位置概率为零，涉及边界解，不能直接计算 $\log0$ 或把该特征当作恒零特征删除。

若特征总和在所有 $(x,y)$ 上恒为常数 $M>0$，则更新方程化为 $e^{M\delta_j}\mathbb E_{p_w}[f_j]=c_j$，故
\begin{equation}
\delta_j=\frac1M\log\frac{c_j}{\mathbb E_{p_w}[f_j]}.
\label{eq:lh-iis-constant-sum}
\end{equation}
当模型低估某特征的期望时，更新增大该特征的权值；高估时则减小。这个方向解释依赖当前模型期望与经验期望均为正。

\begin{example}
\label{ex:lh-iis-five-outcomes}
对例~\ref{ex:lh-five-outcomes} 的分组约束，取 $f_1(y)=\mathbf1\{y\in\{A,B\}\}$、$f_2(y)=1-f_1(y)$，从零权值开始进行一轮 IIS。
\end{example}
\begin{solve}
两个特征之和恒为 $M=1$，经验期望为 $(3/10,7/10)$，初始均匀分布下的模型期望为 $(2/5,3/5)$，所以
\[
\delta_1=\log(3/4),\qquad\delta_2=\log(7/6).
\]
更新后前两个状态的未归一化权为 $3/4$，后三个为 $7/6$，总和为 $2(3/4)+3(7/6)=5$。归一化后恰为 $3/20,3/20,7/30,7/30,7/30$，已经满足约束。这一例子具有分组对称性，不能据此认为一般 IIS 都在一轮内收敛。两个互补特征还带来公共平移冗余，但不妨碍得到唯一概率分布。
\end{solve}

一般情况下，$f^\#$ 随输入输出对变化。令
\begin{equation}
\begin{aligned}
g_j(a)&=\sum_{x,y}\mu_w(x,y)f_j(x,y)e^{af^\#(x,y)}-c_j,\\
g_j'(a)&=\sum_{x,y}\mu_w(x,y)f_j(x,y)f^\#(x,y)e^{af^\#(x,y)},
\end{aligned}
\label{eq:lh-iis-root-function}
\end{equation}
可用牛顿迭代 $a^{(t+1)}=a^{(t)}-g_j(a^{(t)})/g_j'(a^{(t)})$ 求根。在上述正经验期望条件下 $g_j'>0$，可先找到函数异号区间，再采用保留区间的牛顿步或二分步。唯一简单根能支持根附近的牛顿局部二次收敛，并不保证任意初值的无保护牛顿迭代全局收敛。

精确 IIS 使 $\ell(w^{(k)})$ 单调不减，而每个观测对数概率不超过 $0$，所以目标值序列有上界并收敛。这还不是参数收敛定理；若参数序列确实收敛到有限 $w^*$，则增量趋于零，在式~\eqref{eq:lh-iis-update-equation} 中取极限便得到特征期望匹配，由命题~\ref{prop:lh-maxent-certificate} 可知极限模型最优。一次稠密模型期望计算约需 $O(|\mathcal S|Km)$ 运算，额外的一元求根还会增加开销；稀疏特征可按非零项计算。

\subsection{拟牛顿法}
\label{subsec:lh-maxent-quasi-newton}

\emph{拟牛顿法}（quasi-Newton method）利用相邻迭代点的梯度差近似曲率，避免每次显式构造精确海森矩阵。这里最小化 $J(w)=-\ell(w)$，其梯度由式~\eqref{eq:lh-maxent-gradient-hessian} 给出；若采用正则化，则把目标及梯度一并替换为 $J_\lambda$ 与 $\nabla J_\lambda$。这类方法允许一般实值特征，不依赖 IIS 的非负性条件。

先以牛顿法作为参照：在 $\nabla^2J(w)$ 可逆时，用 $\nabla^2J(w)p=-\nabla J(w)$ 求搜索方向，再选步长。对数线性模型的海森矩阵可能因特征冗余而奇异，所以不能仅由半正定性就直接取逆。BFGS 法以对称正定矩阵 $B_k\in\mathbb R^{m\times m}$ 近似海森矩阵，记 $g_k=\nabla J(w^{(k)})$，解
\begin{equation}
B_kp_k=-g_k,\qquad
w^{(k+1)}=w^{(k)}+\alpha_kp_k,\quad\alpha_k>0.
\label{eq:lh-bfgs-direction}
\end{equation}
当 $g_k\ne0$ 时，$g_k^{\mathsf T}p_k=-g_k^{\mathsf T}B_k^{-1}g_k<0$，因此 $p_k$ 为下降方向。定义参数步长向量 $s_k=w^{(k+1)}-w^{(k)}$ 和梯度差向量 $r_k=g_{k+1}-g_k$，BFGS 更新为
\begin{equation}
B_{k+1}=B_k-\frac{B_ks_ks_k^{\mathsf T}B_k}{s_k^{\mathsf T}B_ks_k}
+\frac{r_kr_k^{\mathsf T}}{r_k^{\mathsf T}s_k}.
\label{eq:lh-bfgs-update}
\end{equation}
它满足割线条件 $B_{k+1}s_k=r_k$，即用本轮观测到的梯度变化拟合沿实际步长的曲率。当 $r_k^{\mathsf T}s_k>0$ 且 $B_k\succ0$ 时，新矩阵仍对称正定，两个分母也均为正。

教材以沿搜索方向精确最小化描述步长选择；实际可使用满足 Wolfe 条件的线搜索。取 $0<c_1<c_2<1$，要求
\begin{align}
J(w^{(k)}+\alpha_kp_k)&\leq J(w^{(k)})
+c_1\alpha_kg_k^{\mathsf T}p_k,\notag\\
\nabla J(w^{(k)}+\alpha_kp_k)^{\mathsf T}p_k
&\geq c_2g_k^{\mathsf T}p_k.
\label{eq:lh-wolfe-conditions}
\end{align}
第一式要求充分下降，第二式控制更新后的方向导数。由第二式和下降方向可得
\[
r_k^{\mathsf T}s_k
=\alpha_k(g_{k+1}-g_k)^{\mathsf T}p_k
\geq\alpha_k(c_2-1)g_k^{\mathsf T}p_k>0,
\]
这解释了线搜索与保持近似海森矩阵正定之间的联系。

\textbf{最大熵模型学习的 BFGS 算法。}输入训练集或经验分布、特征、目标 $J$ 及其梯度、容差 $\varepsilon>0$ 与迭代上限；选定 $w^{(0)}$，取 $B_0=I_m$ 或其他对称正定矩阵，$I_m$ 为 $m$ 阶单位矩阵。
\begin{enumerate}
\item 计算 $g_k=\nabla J(w^{(k)})$。若 $\|g_k\|_2\leq\varepsilon$，返回当前参数与 $p_{w^{(k)}}$；否则解式~\eqref{eq:lh-bfgs-direction} 得到 $p_k$。
\item 沿 $p_k$ 搜索满足式~\eqref{eq:lh-wolfe-conditions} 的正步长，更新参数并计算 $g_{k+1}$。若已满足梯度容差则返回。
\item 计算 $s_k,r_k$，检查曲率量 $r_k^{\mathsf T}s_k$；在其可靠为正时使用式~\eqref{eq:lh-bfgs-update}，否则重新检查线搜索，必要时跳过该次矩阵更新或重置正定矩阵。继续迭代，达到上限或线搜索失败时报告当前近似解及停止原因。
\end{enumerate}
精确驻点由凸性保证全局最优，但数值容差下的点只是近似驻点；没有额外条件时，小梯度也不等于给出了明确的参数误差界。若对全部参数作 $\lambda>0$ 的二次正则化，则强凸性给出 $\|w-w^*\|_2\leq\|\nabla J_\lambda(w)\|_2/\lambda$。标准 BFGS 的局部超线性收敛需要有限最优点、其附近海森矩阵正定及适当光滑性和线搜索等条件，不能只凭目标凸便宣称任意初始化都具有这一速度。

实现时也可维护逆海森矩阵近似 $H_k=B_k^{-1}$，直接令 $p_k=-H_kg_k$。逆 BFGS 更新为
\begin{equation}
H_{k+1}=(I_m-\rho_ks_kr_k^{\mathsf T})H_k
(I_m-\rho_kr_ks_k^{\mathsf T})+\rho_ks_ks_k^{\mathsf T},
\qquad\rho_k=\frac1{r_k^{\mathsf T}s_k}.
\label{eq:lh-inverse-bfgs}
\end{equation}
海森近似 $B_k$ 与逆海森近似 $H_k$ 的更新公式不同，不能交换使用。维护稠密 $H_k$ 时，方向与矩阵更新需 $O(m^2)$ 时间、$O(m^2)$ 存储，另加线搜索中的目标和梯度计算；若直接维护 $B_k$，每轮求解线性方程的成本还取决于所用矩阵分解方式。

教材习题还要求写出最大熵模型学习的 DFP（Davidon--Fletcher--Powell）算法。它同样维护逆海森近似，目标、梯度、初始化、下降方向、线搜索及停止准则均可采用上述流程，只需将逆矩阵更新替换为
\begin{equation}
H_{k+1}=H_k+\frac{s_ks_k^{\mathsf T}}{s_k^{\mathsf T}r_k}
-\frac{H_kr_kr_k^{\mathsf T}H_k}{r_k^{\mathsf T}H_kr_k}.
\label{eq:lh-dfp-update}
\end{equation}
在 $H_k\succ0$、$s_k^{\mathsf T}r_k>0$ 时，DFP 更新保持对称正定并满足逆割线条件 $H_{k+1}r_k=s_k$。因此，教材的梯度下降与 DFP 算法练习分别由式~\eqref{eq:lh-logistic-gd} 及式~\eqref{eq:lh-dfp-update} 连同各自流程给出了可执行的学习步骤；关于指数族的对象区别见备注~\ref{rem:lh-exponential-family}。

回看本章，逻辑斯谛回归从线性对数几率定义条件模型，最大熵方法从经验特征约束推导指数形式，两者在对数线性模型上汇合。最大熵与极大似然的等价发生在明确指定的特征、经验边缘和未正则化目标下；算法则负责数值求解这些目标。学习时应依次检查概率是否归一化、经验量与模型量是否区分、负对数似然的梯度方向是否正确，以及有限最优参数和相应收敛条件是否成立。



% 教材来源：机器学习/textbooks/机器学习（分章原文）/机器学习-7.pdf。
% 对应教材第 7 章，PDF 第 1--44 页，书页 111--154；本笔记独立编排章节编号。


\fi

\selectchapter{7}
\ifcompilechapter
\chapter{支持向量机}
\label{chap:lihang-svm}

分类正确并不能唯一确定一个分类器：即使训练数据线性可分，通常也有许多超平面能够把两类样本分开。\emph{支持向量机}（support vector machine，SVM）进一步要求分离边界与最靠近它的训练样本之间保留较大的距离，由此把选择分类器的问题转化为有明确约束的优化问题。本章依次讨论硬间隔、软间隔、核函数和求解对偶问题的序列最小最优化算法；所需基础是向量内积与范数、凸函数、拉格朗日对偶及最优性条件，相关条件将在使用时说明。

记训练集为 $T=\{(x_i,y_i)\}_{i=1}^{N}$，其中 $x_i\in\mathbb R^d$，$y_i\in\{-1,+1\}$，并假设两类样本均非空。随机输入与标签记为 $X,Y$，$x_i,y_i$ 是其观测值。在线性模型中，分数函数与分类决策分别为
\begin{equation}
  g_{w,b}(x)=w^\top x+b,
  \qquad f_{w,b}(x)=\operatorname{sign}(g_{w,b}(x)),
  \label{eq:lihang-svm-score}
\end{equation}
其中 $w\in\mathbb R^d$，$b\in\mathbb R$。为使决策有定义，约定零分数判为 $+1$；涉及严格分离时仍要求 $y_i g_{w,b}(x_i)>0$，不会依靠零分数的并列规则。本章的优化结论针对给定训练集，不能仅由训练间隔大就推出任意测试分布上的错误率小。

% 教材 7.1。
\section{线性可分支持向量机与硬间隔最大化}
\label{sec:lihang-svm-hard}

\subsection{函数间隔与几何间隔}

\begin{definition}[线性可分与两种间隔]
若存在 $w\ne0,b$，使所有 $i$ 都满足 $y_i(w^\top x_i+b)>0$，则称训练集线性可分（linearly separable）。对任意 $w\ne0$，定义超平面 $(w,b)$ 关于样本 $(x_i,y_i)$ 的\emph{函数间隔}（functional margin）和\emph{几何间隔}（geometric margin）为
\begin{equation}
  \widehat\gamma_i=y_i(w^\top x_i+b),
  \qquad
  \gamma_i=\frac{y_i(w^\top x_i+b)}{\lVert w\rVert_2}.
  \label{eq:lihang-svm-margins}
\end{equation}
相对于整个训练集的间隔分别为 $\widehat\gamma=\min_i\widehat\gamma_i$ 和 $\gamma=\min_i\gamma_i$。
\end{definition}

点 $x_i$ 到超平面 $w^\top x+b=0$ 的距离是 $|w^\top x_i+b|/\lVert w\rVert_2$。因此 $\gamma_i$ 是按真实标签定向的带符号距离：严格分类正确时为正，误分类时为负，落在超平面上时为零。若把 $(w,b)$ 同时乘以正数 $a$，超平面和分类方向不变，函数间隔变为原来的 $a$ 倍，几何间隔保持不变；只最大化未经规范化的函数间隔会得到无意义的尺度放大。

线性可分时可以选择使最小几何间隔最大的超平面，写成
\[
  \max_{w\ne0,b,\gamma}\ \gamma,
  \qquad
  y_i(w^\top x_i+b)\ge\gamma\lVert w\rVert_2,
  \quad i=1,\ldots,N.
\]
有限训练集上的严格分离保证 $\widehat\gamma>0$，故可将参数除以 $\widehat\gamma$，使最小函数间隔为 $1$。在这个规范下，最大化 $1/\lVert w\rVert_2$ 等价于最小化 $\lVert w\rVert_2^2/2$，得到\emph{硬间隔支持向量机}（hard-margin SVM）的原始问题（primal problem）：
\begin{equation}
  \begin{aligned}
    \min_{w,b}\quad &\frac12\lVert w\rVert_2^2,\\
    \text{s.t.}\quad &y_i(w^\top x_i+b)\ge1,
      \quad i=1,\ldots,N.
  \end{aligned}
  \label{eq:lihang-svm-hard-primal}
\end{equation}
目标为凸二次函数，约束为仿射不等式，所以这是凸二次规划（convex quadratic programming）。其最优解必有至少一个约束取等号；否则把 $(w,b)$ 同时适当缩小仍然可行，却能减小目标值。

\subsection{最大间隔解的存在性、唯一性与支持向量}

\begin{theorem}[硬间隔解的存在唯一性]
\label{thm:lihang-svm-hard-unique}
设有限训练集线性可分，且正负两类均非空，则问题~\eqref{eq:lihang-svm-hard-primal} 存在唯一最优解 $(w^*,b^*)$，其中 $w^*\ne0$。其超平面严格分开两类训练样本，并使最小几何间隔最大。
\end{theorem}
\begin{proof}
严格分离参数经正比例放大后满足全部约束，故可行集非空。取目标值趋于下确界的可行序列 $(w_k,b_k)$，并令其目标值有共同上界，于是 $w_k$ 有界。固定一个正例 $x_p$ 和一个负例 $x_n$，可行性给出
\[
  1-w_k^\top x_p\le b_k\le-1-w_k^\top x_n.
\]
故 $b_k$ 也有界。这一序列存在收敛子列，闭约束和目标的连续性保证极限达到最小值。由于两类均非空，$w=0$ 不可能同时满足正负类约束，所以 $w^*\ne0$。这里使用了有界子列与闭性，不能只从目标有下界推出最小值一定能取到。

若有两个最优解的法向量不同，则它们的中点仍可行，而平方范数关于 $w$ 严格凸，导致中点的目标严格更小，矛盾。因此最优法向量唯一。对这一法向量记
\[
  a=\min_{i:y_i=+1}(w^*)^\top x_i,
  \qquad c=\max_{i:y_i=-1}(w^*)^\top x_i.
\]
可行截距满足 $1-a\le b\le-1-c$，所以 $a-c\ge2$。若 $a-c>2$，取 $t=2/(a-c)<1$，则 $w'=tw^*$、$b'=-t(a+c)/2$ 仍满足所有间隔约束，却有更小范数，矛盾。故 $a-c=2$，截距区间缩为一点，$b^*=-(a+c)/2$ 也唯一。约束保证严格分离，与几何间隔最大化的等价性已由规范化说明。
\end{proof}

在最优规范下，两条\emph{间隔边界}（margin boundaries）为
\[
  H_+:(w^*)^\top x+b^*=1,
  \qquad H_-:(w^*)^\top x+b^*=-1.
\]
分离超平面位于两者中央，每侧到分离面的距离为 $1/\lVert w^*\rVert_2$，两边界之间的总宽度为 $2/\lVert w^*\rVert_2$。本章用 $\gamma$ 表示单侧最小几何间隔，避免把这两个相差一倍的量混用。训练集中满足 $y_i g_{w^*,b^*}(x_i)=1$ 的点是几何意义上的边界支持向量；它们约束了间隔继续扩大。

\subsection{拉格朗日对偶与最优性条件}
\label{subsec:lihang-svm-hard-dual}

给每个约束 $1-y_i(w^\top x_i+b)\le0$ 配置乘子 $\alpha_i\ge0$，拉格朗日函数为
\begin{equation}
  L(w,b,\alpha)=\frac12\lVert w\rVert_2^2
    +\sum_{i=1}^{N}\alpha_i\{1-y_i(w^\top x_i+b)\}.
  \label{eq:lihang-svm-hard-lagrangian}
\end{equation}
先对 $w,b$ 求下确界。关于 $w$ 的驻点条件给出 $w=\sum_i\alpha_i y_i x_i$；关于不受约束的 $b$，只有 $\sum_i\alpha_i y_i=0$ 时下确界才有限，否则沿适当方向令 $b$ 发散即可使 $L$ 趋于负无穷。代回得到对偶问题（dual problem）
\begin{equation}
  \begin{aligned}
    \max_{\alpha\in\mathbb R^N}\quad
      &D(\alpha)=\sum_{i=1}^{N}\alpha_i
       -\frac12\sum_{i,j=1}^{N}\alpha_i\alpha_j y_i y_j x_i^\top x_j,\\
    \text{s.t.}\quad
      &\alpha_i\ge0\ (i=1,\ldots,N),\qquad
       \sum_{i=1}^{N}\alpha_i y_i=0.
  \end{aligned}
  \label{eq:lihang-svm-hard-dual}
\end{equation}
若写成最小化形式，只需对 $D$ 取负号。矩阵 $Q_{ij}=y_i y_j x_i^\top x_j$ 半正定，因为对任意 $v\in\mathbb R^N$，有 $v^\top Qv=\lVert\sum_i v_i y_i x_i\rVert_2^2\ge0$，故 $D$ 为凹函数。

严格分离参数可继续放大，使每个约束都严格满足，即满足 Slater 严格可行条件。结合凸性与原始最优解的存在性，原始与对偶最优值相等，并存在最优乘子；\emph{Karush--Kuhn--Tucker 条件}（KKT conditions）在这里是最优性的充要条件：
\begin{equation}
  \begin{gathered}
    w=\sum_i\alpha_i y_i x_i,\qquad \sum_i\alpha_i y_i=0,\\
    \alpha_i\ge0,\qquad y_i(w^\top x_i+b)\ge1,\\
    \alpha_i\{y_i(w^\top x_i+b)-1\}=0,
       \qquad i=1,\ldots,N.
  \end{gathered}
  \label{eq:lihang-svm-hard-kkt}
\end{equation}
最后一行是互补松弛（complementary slackness）：若 $\alpha_i>0$，对应样本必在间隔边界上；若 $y_i g(x_i)>1$，则必有 $\alpha_i=0$。

给定对偶最优解 $\alpha^*$，可以恢复
\begin{equation}
  w^*=\sum_i\alpha_i^*y_i x_i,
  \qquad
  b^*=y_j-\sum_i\alpha_i^*y_i x_i^\top x_j
       \quad(\alpha_j^*>0).
  \label{eq:lihang-svm-hard-recovery}
\end{equation}
这样的 $j$ 必存在，否则 $\alpha^*=0$ 将给出不可行的 $w^*=0$。因此学习时可以先求对偶最优乘子，再恢复截距，最后用
\begin{equation}
  g^*(x)=\sum_i\alpha_i^*y_i x_i^\top x+b^*,
  \qquad f^*(x)=\operatorname{sign}(g^*(x))
  \label{eq:lihang-svm-hard-prediction}
\end{equation}
进行预测。所有训练输入只通过内积出现，这是后续引入核函数的关键。

\begin{notealert}[type=caution,title={边界点与正乘子不能无条件等同}]
对一个选定的对偶最优解，通常把 $\alpha_i^*>0$ 的样本称为对偶表示中的支持向量。KKT 条件只保证“正乘子推出位于边界”，不保证反向蕴含；有重复样本或其他退化情况时，某个边界点的乘子可以为零。原始参数唯一也不意味着乘子唯一。删除一个严格位于间隔外、乘子为零的点会保留原有最优证书；但不能据此断言每一个边界点都不可删除，或每次移动支持向量都必然改变分类器。
\end{notealert}

\begin{example}
\label{ex:lihang-svm-three-points}
教材例 7.1 与例 7.2 的训练数据为正例 $x_1=(3,3)^\top$、$x_2=(4,3)^\top$，负例 $x_3=(1,1)^\top$。求最大间隔分类器，并核对原始解与对偶解。
\end{example}
\begin{solve}
原始约束为 $3w_1+3w_2+b\ge1$、$4w_1+3w_2+b\ge1$、$-w_1-w_2-b\ge1$。对偶等式给出 $\alpha_3=\alpha_1+\alpha_2$，代入负对偶目标后得到
\[
  F(\alpha_1,\alpha_2)
    =4\alpha_1^2+\frac{13}{2}\alpha_2^2
      +10\alpha_1\alpha_2-2\alpha_1-2\alpha_2,
  \qquad \alpha_1,\alpha_2\ge0.
\]
在 $(\alpha_1,\alpha_2)=(1/4,0)$ 处，两个偏导数分别为 $0$ 和 $1/2$，满足非负象限上凸函数最小化的最优性条件。因此 $\alpha^*=(1/4,0,1/4)^\top$，并由式~\eqref{eq:lihang-svm-hard-recovery} 得
\[
  w^*=\frac14(3,3)^\top-\frac14(1,1)^\top
      =\left(\frac12,\frac12\right)^\top,
  \qquad b^*=-2.
\]
三个样本的函数间隔分别为 $1,3/2,1$，因此 $x_1,x_3$ 是支持向量。分离面与决策函数为
\[
  \frac12x^{(1)}+\frac12x^{(2)}-2=0,
  \qquad
  f^*(x)=\operatorname{sign}\left(\frac12x^{(1)}+\frac12x^{(2)}-2\right),
\]
其中 $x^{(k)}$ 表示输入的第 $k$ 个坐标。原始最优值与对偶最优值均为 $1/4$，单侧间隔为 $\sqrt2$，两边界间距为 $2\sqrt2$。
\end{solve}

% 教材 7.2。
\section{线性支持向量机与软间隔最大化}
\label{sec:lihang-svm-soft}

\subsection{松弛变量与惩罚参数}

训练样本可能因噪声、类别重叠或特征表达不足而不可分，硬间隔约束此时无可行解。\emph{软间隔支持向量机}（soft-margin SVM）引入\emph{松弛变量}（slack variable）$\xi_i\ge0$，允许一部分样本未达到单位函数间隔，并通过目标函数支付代价：
\begin{equation}
  \begin{aligned}
    \min_{w,b,\xi}\quad
      &\frac12\lVert w\rVert_2^2+C\sum_{i=1}^{N}\xi_i,\\
    \text{s.t.}\quad
      &y_i(w^\top x_i+b)\ge1-\xi_i,\qquad
       \xi_i\ge0,\quad i=1,\ldots,N,
  \end{aligned}
  \qquad C>0.
  \label{eq:lihang-svm-soft-primal}
\end{equation}
对固定的 $w,b$，最优松弛量为 $\xi_i=[1-y_i g_{w,b}(x_i)]_+$，其中 $[u]_+=\max(0,u)$。因此 $\xi_i$ 衡量单位函数间隔的缺失程度：$0<\xi_i<1$ 时仍严格分类正确，$\xi_i=1$ 时分数为零，$\xi_i>1$ 时严格误分类。\emph[caution]{$\sum_i\xi_i$ 惩罚的是间隔违反的总量，不是误分类样本数。}

在式~\eqref{eq:lihang-svm-soft-primal} 的求和约定下，$C$ 越大，目标越重视压低间隔损失；$C$ 越小，平方范数的相对作用越强。这是目标权重的比较，不能推出每个样本的误分类状态都随 $C$ 单调变化。软间隔对可分和不可分数据都适用，但可分时也不必与硬间隔给出相同解，尤其在 $C$ 较小时。

该问题总有可行点，例如 $w=b=0,\xi_i=1$。最优解也存在：有界目标值同时控制 $\lVert w\rVert_2$ 与全部 $\xi_i$，正负类约束进而给出 $b$ 的上下界，再用闭性取极限即可。严格凸性保证最优 $w^*$ 唯一，但目标不直接惩罚 $b$，所以 $b^*$ 可能不唯一。极端情况下 $w^*=0$，模型成为常数分数函数，此时不再有法向量非零的分离超平面，也不能使用 $1/\lVert w^*\rVert_2$ 解释间隔。

\subsection{软间隔的对偶与支持向量的分类}
\label{subsec:lihang-svm-soft-kkt}

分别给间隔约束与非负约束配置乘子 $\alpha_i\ge0$、$\mu_i\ge0$，拉格朗日函数为
\[
  L=\frac12\lVert w\rVert_2^2+C\sum_i\xi_i
    +\sum_i\alpha_i\{1-\xi_i-y_i(w^\top x_i+b)\}
    -\sum_i\mu_i\xi_i.
\]
对 $w,b,\xi$ 的驻点条件分别给出
\[
  w=\sum_i\alpha_i y_i x_i,
  \qquad\sum_i\alpha_i y_i=0,
  \qquad C-\alpha_i-\mu_i=0.
\]
由 $\mu_i\ge0$ 得到盒约束 $0\le\alpha_i\le C$。消去原始变量后，对偶为
\begin{equation}
  \begin{aligned}
    \max_\alpha\quad
      &D(\alpha)=\sum_i\alpha_i
        -\frac12\sum_{i,j}\alpha_i\alpha_j y_i y_j x_i^\top x_j,\\
    \text{s.t.}\quad
      &\sum_i\alpha_i y_i=0,\qquad
       0\le\alpha_i\le C\quad(i=1,\ldots,N).
  \end{aligned}
  \label{eq:lihang-svm-soft-dual}
\end{equation}
相对于硬间隔，对偶目标不变，只给每个乘子加了上界。原始问题可取 $w=b=0,\xi_i=2$ 得到严格可行点，故强对偶成立。除原始可行性、盒约束与上述驻点条件外，KKT 条件还包含
\begin{equation}
  \alpha_i\{y_i g(x_i)-1+\xi_i\}=0,
  \qquad (C-\alpha_i)\xi_i=0.
  \label{eq:lihang-svm-soft-complementarity}
\end{equation}
令 $m_i=y_i g(x_i)$，可将这些条件等价地整理为
\begin{equation}
  \begin{array}{c|c|c}
    \text{乘子状态}&\text{函数间隔}&\text{最优松弛量}\\ \hline
    \alpha_i=0&m_i\ge1&\xi_i=0\\
    0<\alpha_i<C&m_i=1&\xi_i=0\\
    \alpha_i=C&m_i\le1&\xi_i=1-m_i
  \end{array}
  \label{eq:lihang-svm-soft-status}
\end{equation}
因此 $0<\alpha_i<C$ 的\emph{自由支持向量}（free support vector）位于间隔边界上；$\alpha_i=C$ 的支持向量可能位于边界、间隔内部或误分类一侧。特别地，$\alpha_i=C$ 不必意味着误分类，甚至不必有 $\xi_i>0$。

\subsection{截距恢复与没有自由支持向量的情形}
\label{subsec:lihang-svm-intercept}

对偶最优解给出唯一的 $w^*=\sum_i\alpha_i^*y_i x_i$。若存在 $0<\alpha_j^*<C$，由 $y_jg(x_j)=1$ 得
\begin{equation}
  b^*=y_j-\sum_i\alpha_i^*y_i x_i^\top x_j.
  \label{eq:lihang-svm-soft-b-free}
\end{equation}
精确最优解下，所有自由支持向量给出的截距相同；数值近似时可对这些候选值取平均，再检查全部 KKT 条件。

不能预先假设自由支持向量必存在。令 $h_i=\sum_j\alpha_j^*y_j x_j^\top x_i$，若每个乘子都在 $0$ 或 $C$ 上，则式~\eqref{eq:lihang-svm-soft-status} 要求
\begin{align}
  b_{\mathrm{lo}}
    &=\max\bigl(\{1-h_i:y_i=+1,\alpha_i^*=0\}
           \cup\{-1-h_i:y_i=-1,\alpha_i^*=C\}\bigr),\\
  b_{\mathrm{hi}}
    &=\min\bigl(\{-1-h_i:y_i=-1,\alpha_i^*=0\}
           \cup\{1-h_i:y_i=+1,\alpha_i^*=C\}\bigr),\\
  b_{\mathrm{lo}}&\le b^*\le b_{\mathrm{hi}}.
  \label{eq:lihang-svm-b-interval}
\end{align}
在本章两类非空、$C>0$ 的约定下，无自由分量时上下界集合均非空：否则盒约束的端点状态将与 $\sum_i\alpha_i^*y_i=0$ 冲突。精确对偶最优解保证区间非空，任选其中一点即可；它可能是一点，也可能有正长度。这一求法同样适用于后文的核模型，只需把内积换成核函数。

\begin{example}
\label{ex:lihang-svm-nonunique-b}
训练集只有两个观测 $x_1=x_2=0$，标签分别为 $+1,-1$。说明软间隔截距可能不唯一。
\end{example}
\begin{solve}
最优 $w^*=0$，其余目标是 $C([1-b]_++[1+b]_+)$，在任意 $b\in[-1,1]$ 上均取最小值 $2C$。对偶约束要求 $\alpha_1=\alpha_2$，对偶目标为 $2\alpha_1$，故 $\alpha_1^*=\alpha_2^*=C$，没有自由支持向量。式~\eqref{eq:lihang-svm-b-interval} 恰给出 $[-1,1]$，展示了唯一法向量与非唯一截距可以同时出现。
\end{solve}

\subsection{合页损失与正则化经验风险}

\begin{definition}[合页损失]
\emph{合页损失}（hinge loss）定义为 $\ell_{\mathrm{hinge}}(g(x),y)=[1-yg(x)]_+$。它不仅要求分类方向正确，还要求函数间隔达到 $1$ 才给出零损失。
\end{definition}

\begin{proposition}[消去松弛变量]
\label{prop:lihang-svm-hinge}
软间隔原始问题~\eqref{eq:lihang-svm-soft-primal} 等价于无约束凸优化问题
\begin{equation}
  \min_{w,b}\ \frac12\lVert w\rVert_2^2
      +C\sum_{i=1}^{N}[1-y_i(w^\top x_i+b)]_+.
  \label{eq:lihang-svm-hinge-objective}
\end{equation}
这里的等价指最优 $(w,b)$ 相同，最优松弛量由 $\xi_i^*=[1-y_i g_{w^*,b^*}(x_i)]_+$ 恢复。
\end{proposition}
\begin{proof}
固定 $w,b$ 后，约束对第 $i$ 个松弛量的要求为 $\xi_i\ge\max(0,1-y_i g(x_i))$。由于 $C>0$，目标随每个 $\xi_i$ 严格递增，故其最优值必取这个下界。逐个消去 $\xi_i$ 即得式~\eqref{eq:lihang-svm-hinge-objective}，反向代回也满足全部约束。
\end{proof}

对任意实数间隔 $m$，有 $\mathbf1\{m\le0\}\le[1-m]_+$，故合页损失是把零间隔也计入的分类错误指标的凸上界，是一种\emph{代理损失}（surrogate loss）。相比之下，感知机损失为 $[-m]_+$，在 $m=0$ 时已经为零，不会像合页损失那样要求正的安全间隔。合页损失连续且凸，但在 $m=1$ 处不可微；使用凸优化并不要求目标处处可微。

不同的正则化记法需要先统一尺度。将式~\eqref{eq:lihang-svm-hinge-objective} 除以 $C$，可写成 $\sum_i[1-m_i]_++\lambda\lVert w\rVert_2^2$，此时 $\lambda=1/(2C)$；若采用平均损失
\begin{equation}
  \frac1N\sum_i[1-m_i]_+
       +\frac{\lambda}{2}\lVert w\rVert_2^2,
  \qquad \lambda=\frac1{NC}.
  \label{eq:lihang-svm-average-risk}
\end{equation}
两种 $\lambda$ 的定义不同，不能混用。在上述模型中截距 $b$ 不受平方范数惩罚；若把常数特征并入 $w$ 且一并惩罚，便改变了优化问题及其对偶约束。

% 教材 7.3。
\section{非线性支持向量机与核函数}
\label{sec:lihang-svm-kernel}

\subsection{从非线性变换到核技巧}

输入空间中的非线性边界可以在适当的特征空间中变成超平面。例如对 $x=(x^{(1)},x^{(2)})^\top$，取映射 $\phi(x)=((x^{(1)})^2,(x^{(2)})^2)^\top$，则边界 $a(x^{(1)})^2+c(x^{(2)})^2+b=0$ 在新坐标中成为直线。这个例子说明的是边界表示方式的变化，并不意味着任何映射都能把任意带噪数据变为可分数据。

\begin{definition}[核函数]
设 $\mathcal X$ 为非空输入集合。若存在实希尔伯特空间 $\mathcal H$（Hilbert space，即关于内积范数完备的内积空间）及映射 $\phi:\mathcal X\to\mathcal H$，使
\begin{equation}
  K(x,z)=\langle\phi(x),\phi(z)\rangle_{\mathcal H},
  \qquad x,z\in\mathcal X,
  \label{eq:lihang-svm-kernel-definition}
\end{equation}
则称 $K$ 为\emph{核函数}（kernel function），$\mathcal H$ 为相应的特征空间，$\phi$ 为特征映射。这里的核是实对称半正定核，教材也称其为正定核（positive definite kernel）。
\end{definition}

把软间隔模型中的输入改为 $\phi(x_i)$，其对偶与预测只需计算 $\langle\phi(x_i),\phi(x_j)\rangle_{\mathcal H}$。因此只要直接计算 $K(x_i,x_j)$，就可在不显式构造特征坐标的情况下完成学习与预测，这称为\emph{核技巧}（kernel trick）。给定同一个核，特征映射不唯一，例如在特征空间中作正交变换不会改变内积；核确定的是相关内积结构。

\begin{example}
\label{ex:lihang-svm-quadratic-kernel}
教材例 7.3 给定 $x,z\in\mathbb R^2$ 及 $K(x,z)=(x^\top z)^2$，求一个特征映射。
\end{example}
\begin{solve}
展开可得
\[
  (x^\top z)^2=(x^{(1)})^2(z^{(1)})^2
       +2x^{(1)}x^{(2)}z^{(1)}z^{(2)}
       +(x^{(2)})^2(z^{(2)})^2.
\]
故可取 $\mathcal H=\mathbb R^3$，以及
\[
  \phi(x)=\bigl((x^{(1)})^2,\sqrt2x^{(1)}x^{(2)},(x^{(2)})^2\bigr)^\top.
\]
中间的 $\sqrt2$ 保证内积展开时交叉项的系数为 $2$。也可在 $\mathbb R^4$ 中取 $((x^{(1)})^2,x^{(1)}x^{(2)},x^{(1)}x^{(2)},(x^{(2)})^2)^\top$，说明映射与特征空间的表示均可不同。
\end{solve}

\subsection{正定核的判别与再生核希尔伯特空间}
\label{subsec:lihang-svm-psd-kernel}

\begin{theorem}[核函数的 Gram 矩阵判别]
\label{thm:lihang-svm-gram}
实对称函数 $K:\mathcal X\times\mathcal X\to\mathbb R$ 是式~\eqref{eq:lihang-svm-kernel-definition} 意义下的核，当且仅当对任意正整数 $m$、任意 $z_1,\ldots,z_m\in\mathcal X$，\emph{Gram 矩阵}（Gram matrix）$G=[K(z_i,z_j)]_{i,j=1}^{m}$ 均为半正定，即对所有 $c\in\mathbb R^m$ 有 $c^\top Gc\ge0$。
\end{theorem}
\begin{proof}
必要性来自内积：若有特征映射，则
\[
  \sum_{i,j}c_i c_j K(z_i,z_j)
    =\left\lVert\sum_i c_i\phi(z_i)\right\rVert_{\mathcal H}^2\ge0.
\]
为证充分性，考虑函数空间
\[
  S=\left\{f(\cdot)=\sum_{i=1}^{m}a_iK(\cdot,z_i):
      m<\infty,\ a_i\in\mathbb R,\ z_i\in\mathcal X\right\}.
\]
对 $f=\sum_i a_iK(\cdot,z_i)$、$h=\sum_j b_jK(\cdot,u_j)$，定义
\begin{equation}
  B(f,h)=\sum_{i,j}a_i b_jK(z_i,u_j).
  \label{eq:lihang-svm-rkhs-inner}
\end{equation}
先核对这一运算不依赖表示：右端等于 $\sum_j b_j f(u_j)$，故表示零函数的任何 $f$ 都与所有 $h$ 配对为零；再利用 $K$ 的对称性可得关于 $h$ 的同样结论。因此改变任一函数的有限和表示不会改变 $B$。

由 Gram 矩阵半正定性，$B(f,f)\ge0$。对任意 $t\in\mathbb R$，$B(f+th,f+th)\ge0$，由此得到半内积的 Cauchy--Schwarz 不等式 $|B(f,h)|^2\le B(f,f)B(h,h)$；当 $B(h,h)=0$ 时，直接由该一次式对全部 $t$ 非负可得 $B(f,h)=0$。又有
\[
  B(f,K(\cdot,x))=f(x),
  \qquad |f(x)|^2\le B(f,f)K(x,x).
\]
若 $B(f,f)=0$，则对所有 $x$ 都有 $f(x)=0$，故 $f$ 是零函数。于是 $B$ 实际上是 $S$ 上的内积。

将此内积空间完备化得到 $\mathcal H$。上述点值不等式保证每个求值映射 $f\mapsto f(x)$ 有界，因而能唯一连续延拓到 $\mathcal H$，并保留\emph{再生性质}（reproducing property）
\begin{equation}
  f(x)=\langle f,K(\cdot,x)\rangle_{\mathcal H}.
  \label{eq:lihang-svm-reproducing}
\end{equation}
这些延拓仍把元素识别为函数：若某个元素在所有 $x$ 的取值都为零，它与稠密子空间 $S$ 正交，故只能为零。取 $\phi(x)=K(\cdot,x)$，便有 $\langle\phi(x),\phi(z)\rangle_{\mathcal H}=K(x,z)$，完成构造。
\end{proof}

以上构造的空间称为\emph{再生核希尔伯特空间}（reproducing kernel Hilbert space，RKHS）。教材使用的“正定核”允许 Gram 矩阵仅半正定，不要求对每个非零 $c$ 都严格为正；例如 $K(x,z)=1$ 是合法核，却在多个样本上产生秩为一的矩阵。判别定理要求对\emph[caution]{任意有限输入集合}成立，仅检查某一个训练集的核矩阵半正定，不能证明候选函数在整个输入域上都是核。

\subsection{多项式核、高斯核与字符串核}

\emph{多项式核}（polynomial kernel）的常用形式为
\begin{equation}
  K(x,z)=(x^\top z+1)^p,
  \qquad p\in\{1,2,\ldots\}.
  \label{eq:lihang-svm-polynomial-kernel}
\end{equation}
令 $\widetilde x=(1,x^\top)^\top$，则该式等于 $(\widetilde x^\top\widetilde z)^p$。取所有 $p$ 重坐标乘积作为特征，就得到一个显式内积表示；合并重复单项式后，对应于不超过 $p$ 次的多项式特征。齐次形式 $(x^\top z)^p$ 的详细展开见例~\ref{ex:lihang-svm-power-kernel}。这里限定 $p$ 为正整数；不能随意改成实数幂而仍假定它在整个 $\mathbb R^d$ 上是核。

\emph{高斯核}（Gaussian kernel）为
\begin{equation}
  K_\sigma(x,z)=\exp\left(-\frac{\lVert x-z\rVert_2^2}{2\sigma^2}\right),
  \qquad \sigma>0,
  \label{eq:lihang-svm-gaussian-kernel}
\end{equation}
也称高斯径向基核（Gaussian radial basis function kernel，Gaussian RBF kernel）。它的半正定性可以从级数展开核对：
\[
  K_\sigma(x,z)
   =e^{-\lVert x\rVert_2^2/(2\sigma^2)}
    e^{-\lVert z\rVert_2^2/(2\sigma^2)}
      \sum_{r=0}^{\infty}\frac{(x^\top z)^r}{r!\sigma^{2r}}.
\]
每个整数次幂都是核，非负线性组合仍是核；对核乘以 $a(x)a(z)$，相应的 Gram 二次型只需把系数 $c_i$ 换为 $c_i a(x_i)$，所以仍半正定。上式级数在每对固定输入上收敛，在任意有限 Gram 二次型中可逐项取极限，因此高斯核也半正定。它一般对应无限维特征空间，但计算一个核值只需计算原空间距离。

当固定的两个输入不同时，$\sigma\downarrow0$ 使核值趋于零；$\sigma\to\infty$ 则使有限样本间的核值趋于一。前者使相似性更加局部，后者使样本在核中趋于相似。核宽度与 $C$ 共同影响拟合行为，训练时应根据验证数据选择，不能仅比较训练分类率。输入的尺度会改变距离，因而特征缩放也是模型选择的一部分，其统计量应只由训练数据估计。

核的输入不必是实向量。设 $\Sigma$ 是有限字母表，$\Sigma^*$ 为全部有限字符串的集合；给定正整数 $q$ 和衰减参数 $0<\rho<1$。对字符串 $s=s_1\cdots s_{|s|}$，用严格递增的指标组 $I=(i_1,\ldots,i_q)$ 表示子序列 $s[I]=s_{i_1}\cdots s_{i_q}$，跨度定义为 $\ell(I)=i_q-i_1+1$。子序列允许字符之间有间隔；它不同于必须连续的子串。

定义特征映射 $\phi_q(s)\in\mathbb R^{|\Sigma|^q}$，其坐标由 $u\in\Sigma^q$ 索引，且
\begin{equation}
  [\phi_q(s)]_u=\sum_{I:s[I]=u}\rho^{\ell(I)}.
  \label{eq:lihang-svm-string-feature}
\end{equation}
若没有这样的指标组，和为零。于是\emph{字符串子序列核}（string subsequence kernel）为
\begin{equation}
  K_q(s,t)=\sum_{u\in\Sigma^q}[\phi_q(s)]_u[\phi_q(t)]_u
    =\sum_{u\in\Sigma^q}\sum_{I,J:s[I]=t[J]=u}
           \rho^{\ell(I)+\ell(J)}.
  \label{eq:lihang-svm-string-kernel}
\end{equation}
这是显式特征向量的内积，故必为核。相同的长度为 $q$ 的子序列越多、出现跨度越短，对核值的贡献通常越大；实际计算可利用动态规划组织重叠子问题，无须先列出全部特征坐标。

例如取 $q=3,u=\texttt{asd}$，$\texttt{Nasdaq}$ 中相应位置为 $(2,3,4)$，故该坐标为 $\rho^3$；$\texttt{lass das}$ 中的空格也占一个字符位置，相应位置为 $(2,3,6)$、$(2,4,6)$，故该坐标为 $2\rho^5$。它们对核值的这一坐标贡献是 $2\rho^8$，核值本身还须对其他共同子序列求和。式~\eqref{eq:lihang-svm-string-kernel} 是内积；当两边特征均非零时，只有进一步归一化为
\[
  \widetilde K_q(s,t)
    =\frac{K_q(s,t)}{\sqrt{K_q(s,s)K_q(t,t)}}
\]
才是特征向量的余弦相似度（cosine similarity）。

\subsection{非线性支持向量机的目标与预测}

给定核 $K$，特征空间中的原始模型是 $g(x)=\langle w,\phi(x)\rangle_{\mathcal H}+b$，并最小化
\[
  \frac12\lVert w\rVert_{\mathcal H}^2+C\sum_i\xi_i,
  \qquad y_i(\langle w,\phi(x_i)\rangle_{\mathcal H}+b)\ge1-\xi_i,
  \quad\xi_i\ge0.
\]
其最优 $w$ 可以取在训练特征张成的有限维子空间中：将任意 $w$ 分解为该子空间上的投影 $w_{\parallel}$ 与正交分量 $w_{\perp}$，训练分数只依赖前者，而 $\lVert w\rVert^2=\lVert w_{\parallel}\rVert^2+\lVert w_{\perp}\rVert^2$，故去掉正交分量不会增加目标。即使 $\mathcal H$ 无限维，求解也能归结为有限个系数。

以核替换内积，得到统一的对偶最小化形式
\begin{equation}
  \begin{aligned}
    \min_\alpha\quad
      &F(\alpha)=\frac12\sum_{i,j=1}^{N}
          \alpha_i\alpha_j y_i y_j K(x_i,x_j)-\sum_i\alpha_i,\\
    \text{s.t.}\quad
      &\sum_i\alpha_i y_i=0,\qquad
       0\le\alpha_i\le C\quad(i=1,\ldots,N).
  \end{aligned}
  \label{eq:lihang-svm-kernel-dual}
\end{equation}
由定理~\ref{thm:lihang-svm-gram}，$Q_{ij}=y_i y_jK(x_i,x_j)$ 半正定，故这是凸二次规划。可行集非空且紧，连续目标必能取到最小值；乘子解可能不唯一。求得 $\alpha^*$ 后，用自由支持向量恢复
\begin{equation}
  b^*=y_j-\sum_i\alpha_i^*y_i K(x_i,x_j)
       \quad(0<\alpha_j^*<C),
  \qquad
  f^*(x)=\operatorname{sign}\left(\sum_i\alpha_i^*y_iK(x_i,x)+b^*\right).
  \label{eq:lihang-svm-kernel-prediction}
\end{equation}
若没有自由支持向量，则用第~\ref{subsec:lihang-svm-intercept} 小节的截距区间，并令其中 $h_i=\sum_j\alpha_j^*y_jK(x_j,x_i)$。从输入空间看决策边界可以是非线性的；从特征空间看仍是带正则化的线性分数模型。

训练时存储完整核矩阵需要 $O(N^2)$ 个数，计算它需要 $O(N^2)$ 次核求值；也可按需计算并缓存核行。预测若有 $M$ 个正乘子，只需 $M$ 次核求值及相应求和。支持向量带来的是这种表示上的稀疏性，并没有保证 $M$ 总比 $N$ 小很多；在线性核情形，还可先显式合并 $w^*$，把每次预测化为一次 $d$ 维内积。

% 教材 7.4。
\section{序列最小最优化算法}
\label{sec:lihang-svm-smo}

\emph{序列最小最优化}（sequential minimal optimization，SMO）把式~\eqref{eq:lihang-svm-kernel-dual} 分解为一系列只更新两个乘子的子问题，并用解析公式求解每个子问题。之所以不能只改变一个乘子，是因为当其他分量固定时，等式 $\sum_i\alpha_i y_i=0$ 已经唯一确定了剩下的分量；选择两个分量，才留有一个自由度。以下取一对不同下标，并为简便将它们记为 $1,2$。

\subsection{两个乘子的可行区间}

记更新前乘子为 $\alpha_1,\alpha_2$，更新后为 $\alpha_1^+,\alpha_2^+$，其他乘子保持不变。等式约束要求
\begin{equation}
  \alpha_1^+=\alpha_1+y_1y_2(\alpha_2-\alpha_2^+).
  \label{eq:lihang-svm-smo-equality}
\end{equation}
因此子问题是在盒子 $[0,C]^2$ 与一条直线的交线上优化。将 $\alpha_2^+$ 作为自由变量，两者的盒约束等价于 $L\le\alpha_2^+\le H$，其中
\begin{equation}
  \begin{array}{c|c|c}
    &L&H\\ \hline
    y_1\ne y_2&\max(0,\alpha_2-\alpha_1)
              &\min(C,C+\alpha_2-\alpha_1)\\
    y_1=y_2&\max(0,\alpha_1+\alpha_2-C)
            &\min(C,\alpha_1+\alpha_2)
  \end{array}
  \label{eq:lihang-svm-smo-bounds}
\end{equation}
这是分别消去 $\alpha_1^+$ 后，由 $0\le\alpha_1^+,\alpha_2^+\le C$ 得到的交集，不能只把第二个乘子截断到 $[0,C]$。若 $L=H$，这一对没有可行移动方向，应另选工作变量。

\subsection{解析更新、剪辑与零曲率}

记 $K_{ij}=K(x_i,x_j)$，使用当前参数定义
\begin{equation}
  g(x)=\sum_j\alpha_j y_jK(x_j,x)+b,
  \qquad E_i=g(x_i)-y_i.
  \label{eq:lihang-svm-smo-error}
\end{equation}
$E_i$ 是实值分数与数值标签之差，不是分类错误指示，也不是概率误差。令
\begin{equation}
  \eta=K_{11}+K_{22}-2K_{12}
       =\lVert\phi(x_1)-\phi(x_2)\rVert_{\mathcal H}^2\ge0.
  \label{eq:lihang-svm-smo-curvature}
\end{equation}
以 $t=\alpha_2^+$、$\delta=t-\alpha_2$ 表示这一可行方向上的变化，把式~\eqref{eq:lihang-svm-smo-equality} 代入二次目标可得
\begin{equation}
  F(\alpha^+)-F(\alpha)
      =\frac12\eta\delta^2+y_2(E_2-E_1)\delta.
  \label{eq:lihang-svm-smo-reduced}
\end{equation}
其中 $E_2-E_1$ 中的 $b$ 抵消，符合对偶目标本身不含截距的事实。式~\eqref{eq:lihang-svm-smo-reduced} 也给出了更新公式的直接推导：若 $\eta>0$，令关于 $t$ 的导数为零，得到未经约束的极小点
\[
  \widetilde\alpha_2=\alpha_2+\frac{y_2(E_1-E_2)}{\eta}.
\]
由于单变量目标是凸二次函数，在闭区间上的极小点就是该点到区间的投影：
\begin{equation}
  \alpha_2^+=\min\{H,\max\{L,\widetilde\alpha_2\}\},
  \qquad
  \alpha_1^+=\alpha_1+y_1y_2(\alpha_2-\alpha_2^+).
  \label{eq:lihang-svm-smo-update}
\end{equation}
第二个等式恢复等式可行性，式~\eqref{eq:lihang-svm-smo-bounds} 同时保证两个新乘子都满足盒约束。

若 $\eta=0$，不能使用含 $1/\eta$ 的公式。这时式~\eqref{eq:lihang-svm-smo-reduced} 是线性函数，比较 $t=L,H$ 的目标值即可选出最优端点；若线性系数也为零，则整条线段目标相同，可保留原点并另选一对。重复输入或不同输入具有相同特征映射时都可能出现这一情况。合法核在精确运算中保证 $\eta\ge0$，实现时遇到显著负曲率需要检查核矩阵或数值计算。

\subsection{截距与误差缓存的更新}

令 $\Delta\alpha_i=\alpha_i^+-\alpha_i$，根据更新后的第一个或第二个样本达到单位间隔的要求，构造
\begin{align}
  b_1&=b-E_1-y_1\Delta\alpha_1K_{11}
                   -y_2\Delta\alpha_2K_{12},\\
  b_2&=b-E_2-y_1\Delta\alpha_1K_{12}
                   -y_2\Delta\alpha_2K_{22}.
  \label{eq:lihang-svm-smo-b-candidates}
\end{align}
若 $0<\alpha_1^+<C$，取 $b^+=b_1$；否则若 $0<\alpha_2^+<C$，取 $b^+=b_2$；两个乘子都在内部时，精确算术下两个候选值相同。若两者均在盒边界，可取 $b^+=(b_1+b_2)/2$；对子问题而言，该取法满足这一对样本对截距的 KKT 要求，但尚不能保证其他全部样本的 KKT 条件，仍须继续迭代。

对每一个样本 $k$，误差可增量更新为
\begin{equation}
  E_k^+=E_k+y_1\Delta\alpha_1K(x_1,x_k)
              +y_2\Delta\alpha_2K(x_2,x_k)+(b^+-b).
  \label{eq:lihang-svm-smo-error-update}
\end{equation}
与重新对所有正乘子求和相比，缓存 $E_k$ 并使用两条核行可减少重复计算；但缓存必须随乘子与截距一起更新，否则后续变量选择与停止判断会使用过时分数。

\subsection{变量选择、停止条件与完整流程}

由式~\eqref{eq:lihang-svm-soft-status}，令 $r_i=y_i g(x_i)-1$，精确 KKT 条件为 $\alpha_i=0$ 时 $r_i\ge0$、$0<\alpha_i<C$ 时 $r_i=0$、$\alpha_i=C$ 时 $r_i\le0$。可用违反量
\begin{equation}
  v_i=
  \begin{cases}
    \max(0,-r_i),&\alpha_i=0,\\
    |r_i|,&0<\alpha_i<C,\\
    \max(0,r_i),&\alpha_i=C
  \end{cases}
  \label{eq:lihang-svm-kkt-residual}
\end{equation}
衡量当前误差，选取 $v_i$ 较大的分量作为第一个工作变量。数值实现中应另设判断乘子是否贴近 $0,C$ 的边界容差，并同时监控盒约束和 $\sum_i\alpha_i y_i=0$ 的可行性误差。

固定第一个变量后，常从自由乘子中选择使 $|E_1-E_2|$ 较大的第二个变量。其动机来自未经剪辑的位移公式，但实际下降量还取决于 $\eta$ 和可行区间，所以这只是启发式选择，不保证每次选到全体变量对中的最优者。若所选变量对的 $L=H$、位移过小或目标无足够下降，应依次尝试其他自由乘子和其他训练样本，必要时更换第一个变量；扫描自由乘子后还必须回到全训练集检查，避免遗漏边界变量上的违反。

对给定训练集、核函数、$C>0$、KKT 容差 $\varepsilon>0$ 和迭代预算，一次完整求解可以按下述顺序组织。
\begin{enumerate}
  \item 初始化 $\alpha=0,b=0$，于是等式与盒约束均满足，并令 $E_i=-y_i$；准备核值与误差缓存。
  \item 检查当前 KKT 违反，选取两个不同工作变量，计算 $L,H$；若该对无法产生有效移动，则按上述规则继续选取。
  \item 按正曲率解析公式或零曲率端点比较精确求解子问题，更新两个乘子，再更新 $b$ 与全部误差缓存。
  \item 交替扫描自由乘子与全部样本，只有全体样本均满足 $v_i\le\varepsilon$，且可行性误差在容差内时，才按该精度结束；否则继续迭代。
  \item 输出当前 $\alpha,b$ 及核决策函数；若预算先耗尽，报告尚未满足的 KKT 违反量，所得结果只能称为当前近似解。
\end{enumerate}

每次精确求解子问题都不会增大 $F$，而可行域紧且目标有下界，所以目标值序列收敛。\emph[caution]{这并不单独证明乘子序列收敛，更不保证任意变量选择规则都达到全局最优。}达到全局最优的充分证书是全部 KKT 条件成立；算法的收敛分析还需要变量选择与有效下降的相应条件，不能仅从目标值单调推出。教材的启发式过程也不提供与数据无关的固定迭代次数界。

若需要直接评估最优性差距，可以利用当前可行 $\alpha$ 令 $w=\sum_i\alpha_i y_i\phi(x_i)$、$\xi_i=[1-y_i g(x_i)]_+$，并记 $Q_{ij}=y_i y_jK_{ij}$，则
\begin{equation}
  \begin{aligned}
    P&=\frac12\alpha^\top Q\alpha+C\sum_i\xi_i,\\
    D&=\sum_i\alpha_i-\frac12\alpha^\top Q\alpha,\\
    P-D&=\alpha^\top Q\alpha+C\sum_i\xi_i-\sum_i\alpha_i\ge0.
  \end{aligned}
  \label{eq:lihang-svm-duality-gap}
\end{equation}
由弱对偶，原始最优值位于 $[D,P]$，因此对偶间隙给出目标值精度的可核查界。这一数值证书仍只针对训练优化目标，不是测试风险的保证。若核行已缓存，一步更新全部误差需要 $O(N)$ 次标量运算；未缓存时还需相应核求值，不能把两个乘子的解析更新误认为整轮训练只有常数代价。

\begin{example}
\label{ex:lihang-svm-smo-step}
对例~\ref{ex:lihang-svm-three-points} 的三个样本取 $C=1$，从 $\alpha=0,b=0$ 出发，选择样本 $1,3$ 进行一次 SMO 更新。
\end{example}
\begin{solve}
两样本异类，$L=0,H=1$。线性核下 $\eta=18+2-2\times6=8$，初始误差为 $E_1=-1,E_3=1$，故
\[
  \alpha_3^+=0+\frac{(-1)(-1-1)}8=\frac14,
  \qquad \alpha_1^+=\frac14.
\]
两个新乘子都严格位于 $(0,1)$，两个截距候选值均为 $-2$。更新后 $w=(1/2,1/2)^\top$，三个函数间隔为 $1,3/2,1$，其余乘子 $\alpha_2=0$，所以全部 KKT 条件均成立，对偶间隙为零。这组特殊数据一次更新即达到最优，不能据此推断一般数据的迭代次数。
\end{solve}

% 教材章末习题 7.1--7.4；不复现其中要求的作图，仅保留可核查的解析结果。
\section{教材习题与进一步理解}
\label{sec:lihang-svm-exercises}

\begin{example}
\label{ex:lihang-svm-perceptron-dual}
教材习题 7.1：比较感知机的对偶表示与线性可分支持向量机的对偶表示。
\end{example}
\begin{solve}
从 $w=b=0$ 出发，感知机对满足 $y_i(w^\top x_i+b)\le0$ 的样本执行 $w\leftarrow w+\tau y_i x_i$、$b\leftarrow b+\tau y_i$，其中 $\tau>0$ 是固定步长。若第 $i$ 个样本被更新 $n_i$ 次，令 $\beta_i=\tau n_i$，则
\[
  w=\sum_i\beta_i y_i x_i,
  \qquad b=\sum_i\beta_i y_i,
  \qquad g(x)=\sum_i\beta_i y_i x_i^\top x+b.
\]
因此两者都可把分数表示为训练输入的线性组合，并通过 Gram 内积计算训练分数。但感知机的 $\beta_i$ 记录误分更新次数，不是式~\eqref{eq:lihang-svm-hard-dual} 的最优拉格朗日乘子，也不要求 $\sum_i\beta_i y_i=0$；其截距由更新规则累积。支持向量机的乘子则由最大间隔凸优化确定，满足对偶等式和互补条件，截距由支持向量约束恢复。在线性可分条件下，感知机寻找某个分离面，所得结果可依赖扫描顺序；硬间隔支持向量机的规范化原始解由定理~\ref{thm:lihang-svm-hard-unique} 保证唯一。二者共享内积表示，不共享同一学习目标。
\end{solve}

\begin{example}
\label{ex:lihang-svm-five-points}
教材习题 7.2：正例为 $x_1=(1,2)^\top$、$x_2=(2,3)^\top$、$x_3=(3,3)^\top$，负例为 $x_4=(2,1)^\top$、$x_5=(3,2)^\top$。求最大间隔分离面、间隔边界和支持向量。
\end{example}
\begin{solve}
考虑由正类边界点 $x_1,x_3$ 确定法向方向，并令负类点 $x_5$ 位于另一条间隔边界。解
\[
  w^\top x_1+b=w^\top x_3+b=1,
  \qquad w^\top x_5+b=-1
\]
得到 $w=(-1,2)^\top,b=-2$。不能仅凭这三个等式断定最优，还须检查全部约束及对偶证书：五个样本的函数间隔依次为 $1,2,1,2,1$，均满足硬间隔约束；取
\[
  \alpha=\left(\frac12,0,2,0,\frac52\right)^\top,
  \qquad \sum_i\alpha_i y_i=0,
  \qquad\sum_i\alpha_i y_i x_i=(-1,2)^\top,
\]
并且正乘子只对应间隔为 $1$ 的样本，故全部 KKT 条件成立。由此分离面为 $-x^{(1)}+2x^{(2)}-2=0$，决策为其左端的符号；两条边界分别为 $-x^{(1)}+2x^{(2)}=3$ 和 $-x^{(1)}+2x^{(2)}=1$，支持向量为 $x_1,x_3,x_5$。单侧间隔为 $1/\sqrt5$，总宽度为 $2/\sqrt5$；原始与对偶最优值均为 $5/2$。这些方程也完整确定了教材作图题中需要标出的几何对象。
\end{solve}

\begin{example}
\label{ex:lihang-svm-squared-slack}
教材习题 7.3：把线性松弛惩罚改成平方惩罚，求下列模型的对偶：
\begin{equation}
  \min_{w,b,\xi}\ \frac12\lVert w\rVert_2^2+C\sum_i\xi_i^2,
  \qquad y_i(w^\top x_i+b)\ge1-\xi_i,
  \quad\xi_i\ge0,\quad C>0.
  \label{eq:lihang-svm-squared-primal}
\end{equation}
\end{example}
\begin{solve}
对固定 $w,b$，最优松弛量仍为 $[1-y_i g(x_i)]_+$；即使不另列 $\xi_i\ge0$，在半直线 $\xi_i\ge1-y_i g(x_i)$ 上最小化平方也给出同一结果。因此可以先用与之等价、允许 $\xi_i\in\mathbb R$ 的问题推导对偶。给间隔约束配置 $\alpha_i\ge0$，拉格朗日函数为
\[
  L=\frac12\lVert w\rVert_2^2+C\sum_i\xi_i^2
     +\sum_i\alpha_i\{1-\xi_i-y_i(w^\top x_i+b)\}.
\]
分别对 $w,b,\xi$ 求极小，得到
\[
  w=\sum_i\alpha_i y_i x_i,
  \qquad\sum_i\alpha_i y_i=0,
  \qquad\xi_i=\frac{\alpha_i}{2C}.
\]
代回即得对偶最大化形式
\begin{equation}
  \begin{aligned}
    \max_\alpha\quad
      &\sum_i\alpha_i
        -\frac12\sum_{i,j}\alpha_i\alpha_j y_i y_j x_i^\top x_j
        -\frac1{4C}\sum_i\alpha_i^2,\\
    \text{s.t.}\quad
      &\alpha_i\ge0\ (i=1,\ldots,N),\qquad
       \sum_i\alpha_i y_i=0.
  \end{aligned}
  \label{eq:lihang-svm-squared-dual}
\end{equation}
平方惩罚把原来的乘子上界替换成了对偶目标中的二次项；此时没有 $\alpha_i\le C$ 的盒上界。相应最小化目标的 Hessian 为 $Q+I/(2C)$，严格正定，所以对偶最优乘子唯一。消去原始松弛变量则得到平方合页损失 $[1-y_i g(x_i)]_+^2$，不能沿用线性松弛模型的全部 SMO 边界公式。
\end{solve}

\begin{example}
\label{ex:lihang-svm-power-kernel}
教材习题 7.4：证明 $K(x,z)=(x^\top z)^p$ 是核，其中 $x,z\in\mathbb R^d$，$p$ 为正整数。
\end{example}
\begin{solve}
取多重指标 $\beta=(\beta_1,\ldots,\beta_d)$，各分量为非负整数，并记 $|\beta|=\sum_k\beta_k$、$\beta!=\prod_k\beta_k!$、$x^\beta=\prod_k(x^{(k)})^{\beta_k}$。多项式定理给出
\[
  (x^\top z)^p
    =\sum_{|\beta|=p}\frac{p!}{\beta!}x^\beta z^\beta.
\]
令特征向量由所有 $|\beta|=p$ 的坐标 $[\phi_p(x)]_\beta=\sqrt{p!/\beta!}\,x^\beta$ 组成，便有 $K(x,z)=\phi_p(x)^\top\phi_p(z)$，特征维数为 $\binom{d+p-1}{p}$。于是对任意有限输入和系数 $c_i$，
\[
  \sum_{i,j}c_i c_jK(x_i,x_j)
    =\left\lVert\sum_i c_i\phi_p(x_i)\right\rVert_2^2\ge0,
\]
由定义或定理~\ref{thm:lihang-svm-gram} 即知它是正定核。$p=2,d=2$ 时正好恢复例~\ref{ex:lihang-svm-quadratic-kernel} 的三维映射；$p=0$ 的常数函数 $K=1$ 也可单独作为核处理。
\end{solve}

硬间隔通过严格分离约束确定最大几何间隔，软间隔以合页损失允许违反，核函数改变特征空间的内积结构，SMO 则负责数值求解同一个有限维对偶问题。学习时应分别检查数据是否满足建模条件、核是否半正定、参数与截距是否满足 KKT 条件，以及最终模型在独立数据上的表现；其中任何一项都不能由另一项自动代替。

% ===== 李航教材第 8 章：提升方法（对应 机器学习-8.pdf） =====

\fi

\selectchapter{8}
\ifcompilechapter
\chapter{提升方法}
\label{chap:lh-boosting}

提升（boosting）研究如何把一系列较简单的预测规则逐步组合成较强的模型。单个规则可能只能识别数据的一部分结构；后续规则根据当前模型的不足继续学习，最终由这些规则共同作出预测。本章沿着两条相互联系的线索展开：AdaBoost 通过调整样本权重学习二分类器，提升树则通过拟合残差或损失的负梯度逐步改进预测。二者都可以放在\emph{加法模型与前向分步优化}的框架中理解。本章对应教材第 8 章，先说明算法如何运行，再证明训练误差界，最后解释损失函数如何决定每轮学习的对象。

本章记训练集为 $\mathcal D=\{(x_i,y_i)\}_{i=1}^{N}$，其中 $x_i\in\mathcal X\subseteq\mathbb R^d$，$N$ 为样本数，$d$ 为特征数；$M$ 为预先给定的最大迭代轮数。分类部分取 $y_i\in\{-1,+1\}$，回归部分取 $y_i\in\mathbb R$。损失写为 $L(y,t)$，第二个变量 $t$ 是模型在输入处的实数输出；$I(A)$ 表示事件或条件 $A$ 的示性函数，成立时为 $1$，否则为 $0$。下文关于训练误差的代数结论在固定训练集上成立，本身不需要随机抽样假设，也不直接给出新样本上的误差保证。

\section{提升方法与 AdaBoost 算法}
\label{sec:lh-boost-adaboost}

\subsection{从弱学习到逐轮加权}
\label{subsec:lh-boost-weak-learning}

在二分类问题中，随机猜测的错误率为 $1/2$，因而错误率略低于 $1/2$ 的规则也含有可以利用的信息。教材以概率近似正确（probably approximately correct，PAC）学习中的\emph{弱可学习性}（weak learnability）与\emph{强可学习性}（strong learnability）说明提升的动机：前者要求学习结果以高概率比随机猜测好一个受控的幅度，后者要求对任意给定精度和置信要求达到相应保证，并满足相应的多项式资源要求。教材介绍的弱学习提升为强学习，是在该学习框架及其条件下的结论；某个规则只在一份训练集上碰巧优于随机猜测，并不足以得到这一结论。本章不证明 PAC 提升定理，而把每轮加权训练误差小于 $1/2$ 作为具体算法的工作条件。

提升需要解决两个相连的问题：下一轮关注哪些样本，以及各轮规则怎样组合。AdaBoost 的做法是提高当前基本分类器误分样本的相对权重，使下一轮学习更多地关注这些样本；组合时采用加权表决，让当前加权误差较小的规则获得较大的系数。这里改变的是训练样本在优化目标中的权重，观测值与标签保持原样。权重集中反映当前加法模型拟合的困难程度，不表示数据生成分布本身发生了变化，也不意味着被赋予大权重的样本一定没有标注错误。

\subsection{算法、权重更新与边界情形}
\label{subsec:lh-boost-algorithm}

设基本分类器来自给定集合 $\mathcal H$，每个 $G\in\mathcal H$ 都将 $\mathcal X$ 映射到 $\{-1,+1\}$。第 $m$ 轮的权重分布为 $D_m=(w_{m1},\ldots,w_{mN})$，满足 $w_{mi}>0$ 和 $\sum_iw_{mi}=1$。这里的概率分布定义在训练样本的\emph{索引集合}上：若随机索引 $K$ 满足 $\Pr_{D_m}(K=i)=w_{mi}$，则基本分类器的加权误差是
\begin{equation}
 e_m=\sum_{i=1}^{N}w_{mi}I\bigl(G_m(x_i)\ne y_i\bigr)
     =\Pr_{K\sim D_m}\bigl(G_m(x_K)\ne y_K\bigr).
 \label{eq:lh-boost-weighted-error}
\end{equation}
这一写法明确区分了对样本索引取概率与对给定样本计算示性函数，不能把每个固定样本的误分条件本身误写成未经定义的概率再求和。

\textbf{AdaBoost 算法。}输入为训练集 $\mathcal D$、能够接受样本权重的基本学习算法及最大轮数 $M$，输出为实值打分函数 $f$ 与二分类器 $G$。采用以下步骤：
\begin{enumerate}
 \item 初始化 $f_0(x)=0$，$w_{1i}=1/N$，$i=1,\ldots,N$。
 \item 对 $m=1,\ldots,M$，在权重 $D_m$ 下训练基本分类器 $G_m$，计算式~\eqref{eq:lh-boost-weighted-error} 中的 $e_m$。在 $0<e_m<1/2$ 时，取
 \begin{equation}
  \alpha_m=\frac12\log\frac{1-e_m}{e_m},\qquad
  f_m(x)=f_{m-1}(x)+\alpha_mG_m(x),
  \label{eq:lh-boost-alpha}
 \end{equation}
 并按
 \begin{equation}
  Z_m=\sum_{i=1}^{N}w_{mi}e^{-\alpha_my_iG_m(x_i)},\qquad
  w_{m+1,i}=\frac{w_{mi}e^{-\alpha_my_iG_m(x_i)}}{Z_m}
  \label{eq:lh-boost-weight-update}
 \end{equation}
 更新下一轮权重；对边界情形按下文说明处理。
 \item 达到最大轮数或满足另行规定的停止条件后，设实际完成的轮数为 $\widehat M$，输出
 \begin{equation}
  f(x)=f_{\widehat M}(x)=\sum_{m=1}^{\widehat M}\alpha_mG_m(x),\qquad
  G(x)=\operatorname{sign}_{+}\bigl(f(x)\bigr),
  \label{eq:lh-boost-final-classifier}
 \end{equation}
 其中 $\operatorname{sign}_{+}(t)=+1$ 当 $t\geq0$，否则为 $-1$。这一约定规定了表决平局时的输出。
\end{enumerate}

式~\eqref{eq:lh-boost-alpha} 中的对数是自然对数。系数 $\alpha_m$ 在 $0<e_m<1/2$ 时为正，并随 $e_m$ 减小而增大；这些系数不要求和为 $1$。比较不同轮的系数时，各轮的误差是在各自的分布 $D_m$ 下计算的，不能把它们当成同一个固定分布下的精度排名。$f(x)$ 是表决得分，$|f(x)|$ 可以作为该模型的表决强度，但未经额外论证或校准，不能直接解释为预测正确的概率。

将式~\eqref{eq:lh-boost-weight-update} 按误分与正确分类分组，得到
\begin{equation}
 Z_m=(1-e_m)e^{-\alpha_m}+e_me^{\alpha_m}
     =2\sqrt{e_m(1-e_m)}.
 \label{eq:lh-boost-normalizer}
\end{equation}
代入系数的表达式，权重更新可进一步写为
\begin{equation}
 w_{m+1,i}=
 \begin{cases}
  \dfrac{w_{mi}}{2(1-e_m)},&G_m(x_i)=y_i,\\
  \dfrac{w_{mi}}{2e_m},&G_m(x_i)\ne y_i.
 \end{cases}
 \label{eq:lh-boost-weight-factors}
\end{equation}
因此正确分类样本的权重确实下降，误分样本的权重确实上升；误分样本相对于正确分类样本的放大倍数为 $(1-e_m)/e_m=e^{2\alpha_m}$，并不是每个误分样本自身的绝对放大倍数。更新后，$G_m$ 原来误分的样本总权重与原来正确分类的样本总权重各为 $1/2$。这说明下一轮分布重新平衡了这两类样本，而不是简单地重复学习同一问题。

\begin{notealert}[type=caution,title={有限系数公式的适用范围}]
$e_m=0$ 时，式~\eqref{eq:lh-boost-alpha} 给出无穷大的系数，不能作为普通有限数值更新；可停止并输出这个在训练集上完全正确的基本分类器，或另行规定有限系数使它在训练样本上支配此前的表决。$e_m=1/2$ 时 $\alpha_m=0$，模型与权重均不变，应停止或重新寻找规则。$e_m>1/2$ 时，若允许使用 $-G_m$，可先反转预测，其误差变为 $1-e_m$；否则不能套用正系数弱学习步骤。若反转后误差为零，仍按零误差情形处理。有限系数的训练误差分析均假设各轮满足 $0<e_m<1/2$。
\end{notealert}

算法不要求每次都求出 $\mathcal H$ 内加权误差的全局最小值；找到满足弱学习条件的规则即可应用后面的误差界。但将它严格等同于逐步\emph{精确极小化}指数损失时，需要额外假设每轮确实选择最优基本分类器。若第 $m$ 轮基本学习的计算成本为 $C_m$，且每个基本分类器在单个样本上的求值成本有常数上界，则加权误差、得分和权重更新的额外总成本为 $O(MN)$，总训练成本为 $O(\sum_mC_m+MN)$；基本分类器本身较复杂时，应另计其求值成本。

\subsection{用决策树桩完成三轮提升}
\label{subsec:lh-boost-classification-example}

\begin{example}
\label{ex:lh-boost-stumps}
取教材中的一维二分类数据，按 $x=0,1,\ldots,9$ 排列，其标签为
\[
 (y_1,\ldots,y_{10})=(1,1,1,-1,-1,-1,1,1,1,-1).
\]
基本分类器为决策树桩（decision stump），即只作一次阈值判断的二叶决策树。允许阈值两侧的预测方向互换；在每轮选择加权误差最小的树桩，误差相同时允许任选一个最优者。求三轮 AdaBoost 的分类结果。
\end{example}
\begin{solve}
初始权重均为 $1/10$。第一轮可以选择
\[
 G_1(x)=\begin{cases}+1,&x<2.5,\\-1,&x\geq2.5.\end{cases}
\]
它误分 $x=6,7,8$，故 $e_1=3/10$，$\alpha_1=\tfrac12\log(7/3)\approx0.423649$。由式~\eqref{eq:lh-boost-weight-factors} 得
\[
 D_2=\left(\frac1{14},\frac1{14},\frac1{14},\frac1{14},\frac1{14},\frac1{14},
            \frac16,\frac16,\frac16,\frac1{14}\right).
\]
第二轮在 $D_2$ 下选择
\[
 G_2(x)=\begin{cases}+1,&x<8.5,\\-1,&x\geq8.5.\end{cases}
\]
这次误分 $x=3,4,5$，其总权重为 $e_2=3/14$，因此 $\alpha_2=\tfrac12\log(11/3)\approx0.649641$，并有
\[
 D_3=\left(\frac1{22},\frac1{22},\frac1{22},\frac16,\frac16,\frac16,
            \frac7{66},\frac7{66},\frac7{66},\frac1{22}\right).
\]
第三轮选择相反预测方向的树桩
\[
 G_3(x)=\begin{cases}-1,&x<5.5,\\+1,&x\geq5.5.\end{cases}
\]
其误分点为 $x=0,1,2,9$，故
\[
 e_3=\frac4{22}=\frac2{11},\qquad
 \alpha_3=\frac12\log\frac92\approx0.752039.
\]
于是 $f_3=\alpha_1G_1+\alpha_2G_2+\alpha_3G_3$，在四个区间上的取值分别为
\[
 f_3(x)\approx
 \begin{cases}
  0.321252,&x<2.5,\\
  -0.526046,&2.5\leq x<5.5,\\
  0.978031,&5.5\leq x<8.5,\\
  -0.321252,&x\geq8.5.
 \end{cases}
\]
它在全部十个训练样本上的符号均与标签一致，训练错误率为零。第一轮与第二轮组合后的训练错误率都为 $3/10$，但误分对象已经改变；第三轮的组合利用前两轮积累的信息使全部训练样本分类正确。以上计算在轮间保留精确权重，仅在展示系数和得分时取近似值；由精确的 $e_3=2/11$ 计算得到 $\alpha_3\approx0.752039$，采用这一复算值而不沿用教材中的 $0.7514$。
\end{solve}

\section{AdaBoost 的训练误差分析}
\label{sec:lh-boost-training-error}

指数损失（exponential loss）$L(y,t)=e^{-yt}$ 把分类方向与分类强度联系起来。对样本 $(x_i,y_i)$，量 $y_if(x_i)$ 称为该样本的\emph{分类间隔}（classification margin，这里指未归一化的表决间隔）：它为正表示分类方向正确，为负表示方向错误，而其增大使指数损失减小。这一用法不同于支持向量机中经过参数范数归一化的几何间隔；此处的数值会随所有表决系数的共同缩放而改变。

\begin{theorem}[训练误差的指数损失上界]
\label{thm:lh-boost-training-bound}
在均匀初始权重 $w_{1i}=1/N$ 下，设算法完成 $M$ 轮有限系数更新，$f_M$、$G$ 与 $Z_m$ 分别按式~\eqref{eq:lh-boost-final-classifier} 和式~\eqref{eq:lh-boost-weight-update} 定义。则经验分类错误率满足
\begin{equation}
 \widehat R_{01}(G):=\frac1N\sum_{i=1}^{N}I\bigl(G(x_i)\ne y_i\bigr)
 \leq\frac1N\sum_{i=1}^{N}e^{-y_if_M(x_i)}
 =\prod_{m=1}^{M}Z_m.
 \label{eq:lh-boost-training-bound}
\end{equation}
\end{theorem}
\begin{proof}
若 $G(x_i)\ne y_i$，按所规定的平局规则仍有 $y_if_M(x_i)\leq0$，故 $e^{-y_if_M(x_i)}\geq1$；若分类正确，示性函数为零而指数函数为正。因此逐点成立
\[
 I\bigl(G(x_i)\ne y_i\bigr)\leq e^{-y_if_M(x_i)}.
\]
对 $i$ 求平均即得不等式。另一方面，逐次展开权重递推式，利用 $f_0=0$，得到
\begin{equation}
 w_{M+1,i}=\frac{(1/N)\exp\bigl(-y_i\sum_{m=1}^{M}\alpha_mG_m(x_i)\bigr)}
                    {\prod_{m=1}^{M}Z_m}
           =\frac{e^{-y_if_M(x_i)}}{N\prod_{m=1}^{M}Z_m}.
 \label{eq:lh-boost-expanded-weight}
\end{equation}
每个 $D_{m+1}$ 均被规范化，故对式~\eqref{eq:lh-boost-expanded-weight} 求和并使用 $\sum_iw_{M+1,i}=1$，即得平均指数损失等于 $\prod_mZ_m$。
\end{proof}

式~\eqref{eq:lh-boost-expanded-weight} 也解释了权重的另一层含义：第 $m+1$ 轮权重正比于当前加法模型在样本上的指数损失。因此，大权重不仅可以由上一轮误分造成，也反映了此前各轮表决累积形成的较小间隔。该恒等式对任意有限实系数和取值为 $\{-1,+1\}$ 的基本分类器成立；只有代入 AdaBoost 的特定系数后，才能进一步得到下一结论。

\begin{theorem}[弱学习优势与指数型误差界]
\label{thm:lh-boost-edge-bound}
若每轮均有 $0<e_m<1/2$ 并采用式~\eqref{eq:lh-boost-alpha} 中的系数，记弱分类器相对随机猜测的优势为 $\gamma_m=1/2-e_m$，则
\begin{equation}
 \prod_{m=1}^{M}Z_m
 =\prod_{m=1}^{M}2\sqrt{e_m(1-e_m)}
 =\prod_{m=1}^{M}\sqrt{1-4\gamma_m^2}
 \leq\exp\left(-2\sum_{m=1}^{M}\gamma_m^2\right).
 \label{eq:lh-boost-edge-bound}
\end{equation}
\end{theorem}
\begin{proof}
前两个等式分别来自式~\eqref{eq:lh-boost-normalizer} 与 $e_m=1/2-\gamma_m$。对 $0\leq u<1$，令 $h(u)=\log(1-u)+u$，则 $h(0)=0$ 且 $h'(u)=-u/(1-u)\leq0$，所以 $\log(1-u)\leq-u$。取 $u=4\gamma_m^2$，有
\[
 \log Z_m=\frac12\log(1-4\gamma_m^2)\leq-2\gamma_m^2.
\]
对各轮求和再取指数，得到式~\eqref{eq:lh-boost-edge-bound}。
\end{proof}

\begin{corollary}
\label{cor:lh-boost-uniform-edge}
在定理~\ref{thm:lh-boost-edge-bound} 的条件下，若还存在与轮次无关的 $\gamma>0$，使每轮 $\gamma_m\geq\gamma$，则
\begin{equation}
 \widehat R_{01}(G)\leq e^{-2M\gamma^2}.
 \label{eq:lh-boost-uniform-edge}
\end{equation}
特别地，对 $0<\varepsilon<1$，取 $M\geq\log(1/\varepsilon)/(2\gamma^2)$ 足以使此上界不超过 $\varepsilon$。
\end{corollary}
\begin{proof}
由 $\sum_{m=1}^{M}\gamma_m^2\geq M\gamma^2$，结合定理~\ref{thm:lh-boost-training-bound} 与定理~\ref{thm:lh-boost-edge-bound} 即得。最后解不等式 $e^{-2M\gamma^2}\leq\varepsilon$，得到轮数条件。
\end{proof}

算法本身使用各轮实际误差 $e_m$ 计算系数，不需要事先知道统一下界 $\gamma$。若无限迭代且始终满足有限系数条件，更一般地，只要 $\sum_m\gamma_m^2=\infty$，式~\eqref{eq:lh-boost-edge-bound} 中的上界就趋于零；仅有每轮 $e_m<1/2$ 而优势衰减过快，则不能从该界推出趋零。

\begin{notealert}[type=caution,title={下降的量与结论的范围}]
在 $0<e_m<1/2$ 时 $Z_m<1$，所以平均指数损失及其给出的训练错误率上界随轮数严格下降。上界下降并不推出实际的零一训练错误率每轮下降，后者取决于各个间隔是否穿过零。上述分析也没有给出测试错误率的单调性；迭代轮数与基本模型复杂度仍需结合独立验证数据等方式选择。对于难以拟合的异常点或错误标签，指数损失可能持续赋予较大权重，这需要与训练误差界分开考虑。
\end{notealert}

\section{加法模型与 AdaBoost 的前向分步解释}
\label{sec:lh-boost-stagewise}

\subsection{前向分步优化的目标}
\label{subsec:lh-boost-additive}

\begin{definition}
\label{def:lh-boost-additive}
设 $b(x;\theta)$ 是参数为 $\theta\in\Theta$ 的实值基函数。形如
\begin{equation}
 f_M(x)=f_0(x)+\sum_{m=1}^{M}\beta_mb(x;\theta_m)
 \label{eq:lh-boost-additive-model}
\end{equation}
的模型称为\emph{加法模型}（additive model），其中 $\beta_m\in\mathbb R$ 为组合系数，$f_0$ 为初始函数。这里的“加法”指对基函数求和，基函数本身可以是输入 $x$ 的非线性函数。
\end{definition}

给定损失 $L$，联合学习所有参数的经验风险目标为
\begin{equation}
 \min_{\{\beta_m,\theta_m\}_{m=1}^{M}}
 \sum_{i=1}^{N}L\left(y_i,f_0(x_i)+\sum_{m=1}^{M}\beta_mb(x_i;\theta_m)\right).
 \label{eq:lh-boost-joint-objective}
\end{equation}
\emph{前向分步算法}（forward stagewise algorithm）以逐轮添加基函数代替同时优化：先选择 $f_0$，然后对 $m=1,\ldots,M$，固定已有模型，求
\begin{equation}
 (\beta_m,\theta_m)\in\operatorname*{arg\,min}_{\beta\in\mathbb R,\,\theta\in\Theta}
 \sum_{i=1}^{N}L\bigl(y_i,f_{m-1}(x_i)+\beta b(x_i;\theta)\bigr),\qquad
 f_m=f_{m-1}+\beta_mb(\cdot;\theta_m).
 \label{eq:lh-boost-stagewise-update}
\end{equation}
输入是训练集、损失、基函数集合、初始模型和轮数或停止条件，输出是最后一个加法模型。式~\eqref{eq:lh-boost-stagewise-update} 假设每步最小值可达到；实际采用近似搜索时，应说明所得只是近似解。该方法不重新调整此前各轮参数，所以是贪心构造，并不自动求出式~\eqref{eq:lh-boost-joint-objective} 的全局最优解。若允许 $\beta=0$ 且精确求解每步子问题，经验损失至少不会增加。

\subsection{从指数损失推导 AdaBoost}
\label{subsec:lh-boost-exponential}

\begin{proposition}
\label{prop:lh-boost-stagewise-adaboost}
设 $y_i\in\{-1,+1\}$，基本分类器集合 $\mathcal H$ 关于取负封闭，$f_0=0$。以 $L(y,t)=e^{-yt}$ 为损失，在每轮精确最小化加权误差，且最优误差满足 $0<e_m<1/2$，则式~\eqref{eq:lh-boost-stagewise-update} 的前向分步更新等同于 AdaBoost 更新。
\end{proposition}
\begin{proof}
给定 $f_{m-1}$，定义未规范化权重及其总和
\[
 \widetilde w_{mi}=e^{-y_if_{m-1}(x_i)},\qquad
 S_m=\sum_{i=1}^{N}\widetilde w_{mi},\qquad
 w_{mi}=\widetilde w_{mi}/S_m.
\]
它们相对于本轮待优化的 $(\alpha,G)$ 都是已知量，但在不同样本之间不相同，不能从加权求和中删去。本轮目标写为
\[
 \sum_{i=1}^{N}e^{-y_i(f_{m-1}(x_i)+\alpha G(x_i))}
 =S_m\sum_{i=1}^{N}w_{mi}e^{-\alpha y_iG(x_i)}.
\]
只有共同的正常数 $S_m$ 可以在确定极小点时省略。记 $e(G)=\sum_iw_{mi}I(G(x_i)\ne y_i)$，则规范化后的目标为
\begin{equation}
 Q_m(\alpha,G)=(1-e(G))e^{-\alpha}+e(G)e^{\alpha}
              =e^{-\alpha}+(e^{\alpha}-e^{-\alpha})e(G).
 \label{eq:lh-boost-stagewise-exponential}
\end{equation}
由于 $\mathcal H$ 对取负封闭，任意 $\alpha<0$ 的增量都可以写为 $(-\alpha)(-G)$，故只需考虑 $\alpha\geq0$。对每个固定的 $\alpha>0$，式~\eqref{eq:lh-boost-stagewise-exponential} 随 $e(G)$ 严格增加，最优基本分类器因此恰为加权错误率最小者 $G_m$。

固定 $G_m$ 后，目标关于 $\alpha$ 的一、二阶导数分别为
\[
 \frac{\partial Q_m}{\partial\alpha}=-(1-e_m)e^{-\alpha}+e_me^{\alpha},\qquad
 \frac{\partial^2Q_m}{\partial\alpha^2}=(1-e_m)e^{-\alpha}+e_me^{\alpha}>0.
\]
令一阶导数为零，唯一极小点就是 $\alpha_m=\tfrac12\log((1-e_m)/e_m)>0$，其目标值 $2\sqrt{e_m(1-e_m)}<1=Q_m(0,G)$。最后，由模型相加可得
\[
 \widetilde w_{m+1,i}=\widetilde w_{mi}e^{-\alpha_my_iG_m(x_i)}.
\]
两侧分别规范化后得到式~\eqref{eq:lh-boost-weight-update}。初始 $f_0=0$ 又给出 $w_{1i}=1/N$，故全部更新一致。
\end{proof}

这一推导把模型、损失与算法分开说明：模型是二值基函数的加权和，损失是指数损失，算法是在已有模型固定时选择下一项。若基本学习器只找到一个满足 $e_m<1/2$ 的规则，则系数仍精确最小化沿该规则方向的损失，但未必精确解出对所有 $G\in\mathcal H$ 的联合子问题；训练误差界仍然适用。

\begin{remark}
\label{rem:lh-boost-conditional-risk}
可以从总体风险进一步理解指数损失的分类方向。令 $p(x)=\Pr(Y=+1\mid X=x)$，对满足 $0<p(x)<1$ 的输入，若暂时允许在该点自由选择实值输出 $t$，则条件期望损失为
\[
 r_x(t)=p(x)e^{-t}+(1-p(x))e^t.
\]
由 $r_x'(t)=-p(x)e^{-t}+(1-p(x))e^t$ 及 $r_x''(t)>0$，唯一最优值为
\[
 t^*(x)=\frac12\log\frac{p(x)}{1-p(x)}.
\]
因此 $t^*(x)$ 的符号与 $p(x)-1/2$ 一致，在 $p(x)=1/2$ 时任一类别都是最优决策；当 $p(x)=0$ 或 $1$ 时，下确界分别在 $t\to-\infty$ 或 $t\to+\infty$ 时逼近。这里讨论的是不受模型约束的总体条件风险，有限样本上有限轮数的 AdaBoost 得分一般不等于 $t^*(x)$，不能由此宣称其概率预测已经校准。
\end{remark}

\section{提升树与梯度提升}
\label{sec:lh-boost-trees}

\subsection{树的加法模型}
\label{subsec:lh-boost-tree-model}

以分类树或回归树为基函数的提升方法称为\emph{提升树}（boosting tree）。分类树输出离散类别；回归树在每个叶区域上输出一个常数。当用取值为 $\{-1,+1\}$ 的分类树作为 AdaBoost 的基本分类器时，得到的是分类提升树；对一般损失进行梯度提升时，即使最终任务是分类，也可以用输出实数的回归树逐步修正分类得分。这两种构造的输出形式与学习目标需要分别说明。

设第 $m$ 棵回归树把输入空间划分为 $J_m$ 个互不相交且并集为 $\mathcal X$ 的叶区域 $R_{m1},\ldots,R_{mJ_m}$，每个叶区域均包含训练样本，其函数形式为
\begin{equation}
 T(x;\Theta_m)=\sum_{j=1}^{J_m}c_{mj}I(x\in R_{mj}),\qquad
 \Theta_m=\{(R_{mj},c_{mj})\}_{j=1}^{J_m}.
 \label{eq:lh-boost-regression-tree}
\end{equation}
区域由树中的递归划分产生，不能当成任意不受约束的集合来优化。树的叶值可以吸收加法系数，故模型写成
\begin{equation}
 f_M(x)=f_0(x)+\sum_{m=1}^{M}T(x;\Theta_m).
 \label{eq:lh-boost-tree-ensemble}
\end{equation}
每轮固定 $f_{m-1}$，希望选择
\begin{equation}
 \Theta_m\in\operatorname*{arg\,min}_{\Theta}
 \sum_{i=1}^{N}L\bigl(y_i,f_{m-1}(x_i)+T(x_i;\Theta)\bigr).
 \label{eq:lh-boost-tree-stage}
\end{equation}
树结构的全局搜索一般很困难，通常以递归贪心划分近似求解，不能把这种构造直接视为对所有允许树结构的全局极小化。叶数、深度与树的棵数共同影响表达能力；每棵树简单并不意味着整个加法模型也简单。

\subsection{平方损失下的残差拟合}
\label{subsec:lh-boost-residual}

对回归任务，取平方损失 $L(y,t)=(y-t)^2$。定义第 $m$ 轮拟合前的\emph{残差}（residual）为
\begin{equation}
 r_{mi}=y_i-f_{m-1}(x_i),\qquad i=1,\ldots,N.
 \label{eq:lh-boost-residual}
\end{equation}
本轮优化目标随即成为
\begin{equation}
 \sum_{i=1}^{N}\bigl[y_i-f_{m-1}(x_i)-T(x_i;\Theta)\bigr]^2
 =\sum_{i=1}^{N}\bigl[r_{mi}-T(x_i;\Theta)\bigr]^2.
 \label{eq:lh-boost-residual-objective}
\end{equation}
所以\emph[tip]{拟合当前残差，就是求解平方损失下下一棵树的子问题}，不需要重新用原始标签拟合一棵独立的树。若区域划分已经确定，记 $n_{mj}=\#\{i:x_i\in R_{mj}\}>0$，则最优叶值为
\begin{equation}
 c_{mj}=\frac1{n_{mj}}\sum_{x_i\in R_{mj}}r_{mi}.
 \label{eq:lh-boost-residual-leaf}
\end{equation}
这是因为对该叶内的损失求导得到 $-2\sum_{x_i\in R_{mj}}(r_{mi}-c)$，令其为零即得平均值，而二阶导数 $2n_{mj}>0$ 保证它是唯一极小点。

\textbf{平方损失提升树算法。}输入为实值标签训练集、允许的回归树复杂度以及轮数 $M$。按教材取 $f_0=0$；对每一轮，先按式~\eqref{eq:lh-boost-residual} 计算残差，再以 $(x_i,r_{mi})$ 为新训练数据构造平方误差回归树 $T_m$，并用式~\eqref{eq:lh-boost-residual-leaf} 确定其叶值，最后更新 $f_m=f_{m-1}+T_m$。达到给定轮数或预设停止条件后输出 $f_M$。也可以从最优常数 $f_0=\bar y=N^{-1}\sum_i y_i$ 出发，但它会改变每轮树的数值；以下教材例题统一采用 $f_0=0$。

\begin{proposition}
\label{prop:lh-boost-squared-decrease}
在第 $m$ 轮，若各叶非空且叶值取该叶残差平均值，则无论划分是否全局最优，都有
\begin{equation}
 \sum_{i=1}^{N}\bigl(y_i-f_{m-1}(x_i)\bigr)^2
 -\sum_{i=1}^{N}\bigl(y_i-f_m(x_i)\bigr)^2
 =\sum_{j=1}^{J_m}n_{mj}c_{mj}^2\geq0.
 \label{eq:lh-boost-squared-decrease}
\end{equation}
\end{proposition}
\begin{proof}
在叶区域 $R_{mj}$ 上，展开两平方之差并使用 $\sum_{x_i\in R_{mj}}r_{mi}=n_{mj}c_{mj}$，得到
\[
 \sum_{x_i\in R_{mj}}\bigl[r_{mi}^2-(r_{mi}-c_{mj})^2\bigr]
 =2c_{mj}\sum_{x_i\in R_{mj}}r_{mi}-n_{mj}c_{mj}^2
 =n_{mj}c_{mj}^2.
\]
对所有叶求和即得。
\end{proof}

\begin{example}
\label{ex:lh-boost-regression}
取教材的十个回归样本 $x_i=i$，$i=1,\ldots,10$，其观测值为
\[
 (y_1,\ldots,y_{10})=(5.56,5.70,5.91,6.40,6.80,7.05,8.90,8.70,9.00,9.05).
\]
在输入范围 $[0.5,10.5]$ 上，以回归树桩为基函数，从 $f_0=0$ 开始进行提升。
\end{example}
\begin{solve}
对候选阈值 $s\in\{1.5,2.5,\ldots,9.5\}$，令 $R_1(s)=\{x:x<s\}$、$R_2(s)=\{x:x\geq s\}$，在两个区域内分别用响应值平均数 $c_1(s),c_2(s)$ 拟合。第一轮比较
\[
 E(s)=\sum_{x_i<s}\bigl(y_i-c_1(s)\bigr)^2
      +\sum_{x_i\geq s}\bigl(y_i-c_2(s)\bigr)^2.
\]
例如 $E(6.5)\approx1.930008$，而 $E(5.5)\approx3.911320$、$E(7.5)\approx8.009810$；枚举全部九个候选点，最小值在 $s=6.5$ 取得。两叶平均值为
\[
 c_{11}=\frac{37.42}{6}=6.236666\ldots,\qquad
 c_{12}=\frac{35.65}{4}=8.9125,
\]
故 $T_1(x)=c_{11}$ 当 $x<6.5$，否则为 $c_{12}$，且 $f_1=T_1$。下一轮拟合前的残差约为
\[
 (r_{21},\ldots,r_{2,10})\approx
 (-0.676667,-0.536667,-0.326667,0.163333,0.563333,0.813333,
 -0.012500,-0.212500,0.087500,0.137500).
\]
对这些残差重复阈值枚举，第二轮最优阈值为 $3.5$，得到
\[
 T_2(x)=\begin{cases}-0.513333\ldots,&x<3.5,\\0.22,&x\geq3.5.\end{cases}
\]
于是 $f_2=f_1+T_2$ 在 $x<3.5$、$3.5\leq x<6.5$、$x\geq6.5$ 三段上的值分别为 $5.723333\ldots$、$6.456666\ldots$、$9.1325$，平方误差总和降为 $0.800675$。继续按同一规则迭代，结果见表~\ref{tab:lh-boost-regression-rounds}；表中 $c_{mL}$、$c_{mR}$ 分别表示阈值左侧与右侧的叶值。
\begin{table}[htbp]
 \centering
 \caption{以精确残差逐轮拟合回归树桩的结果，显示值取六位小数}
 \label{tab:lh-boost-regression-rounds}
 \begin{tabular}{ccrrr}
  \toprule
  轮次 $m$ & 阈值 $s_m$ & $c_{mL}$ & $c_{mR}$ & $\sum_i(y_i-f_m(x_i))^2$\\
  \midrule
  1 & 6.5 & 6.236667 & 8.912500 & 1.930008\\
  2 & 3.5 & -0.513333 & 0.220000 & 0.800675\\
  3 & 6.5 & 0.146667 & -0.220000 & 0.478008\\
  4 & 4.5 & -0.160833 & 0.107222 & 0.305559\\
  5 & 6.5 & 0.071481 & -0.107222 & 0.228915\\
  6 & 2.5 & -0.150648 & 0.037662 & 0.172178\\
  \bottomrule
 \end{tabular}
\end{table}

把六棵树相加，得到
\[
 f_6(x)\approx
 \begin{cases}
  5.630000,&0.5\leq x<2.5,\\
  5.818310,&2.5\leq x<3.5,\\
  6.551644,&3.5\leq x<4.5,\\
  6.819699,&4.5\leq x<6.5,\\
  8.950162,&6.5\leq x\leq10.5.
 \end{cases}
\]
它是比任何单棵树桩更细的分段常数函数。若事先约定平方误差总和不超过 $0.18$ 即停止，则第六轮首次达到这一标准。所有内部运算均保留原始数据精度，表中的显示值不作为下一轮输入；教材给出的部分近似数值与这一复算结果不同，例如第二轮误差总和应约为 $0.800675$。这里的停止标准只针对本例的训练拟合，并不是测试性能保证。
\end{solve}

\subsection{一般可微损失下的梯度提升}
\label{subsec:lh-boost-gradient}

对于一般损失，直接求解式~\eqref{eq:lh-boost-tree-stage} 可能不方便。\emph{梯度提升}（gradient boosting）先在训练样本的预测值空间中寻找损失下降方向，再用回归树近似该方向。记当前预测向量为 $\boldsymbol t=(t_1,\ldots,t_N)^{\mathsf T}$，经验损失总和为 $\mathcal L(\boldsymbol t)=\sum_iL(y_i,t_i)$。假设每个 $L(y_i,\cdot)$ 在当前预测值处可微，则第 $m$ 轮的\emph{伪残差}（pseudo-residual）定义为
\begin{equation}
 r_{mi}=-\left.\frac{\partial L(y_i,t)}{\partial t}\right|_{t=f_{m-1}(x_i)}.
 \label{eq:lh-boost-pseudo-residual}
\end{equation}
这里的 $r_{mi}$ 是损失对预测值的负导数，不是对输入 $x_i$ 求导；只有特定损失下它才等于通常的残差 $y_i-f_{m-1}(x_i)$。

设 $\boldsymbol r_m=(r_{m1},\ldots,r_{mN})^{\mathsf T}$。可微性给出对任意固定方向 $\boldsymbol h$ 的一阶展开
\[
 \mathcal L(\boldsymbol t+\eta\boldsymbol h)
 =\mathcal L(\boldsymbol t)-\eta\boldsymbol r_m^{\mathsf T}\boldsymbol h+o(\eta),
 \qquad \eta\to0.
\]
在 $\boldsymbol r_m\ne0$ 时，单位范数方向中负梯度方向使一阶下降最大，但任意预测向量未必能表示为一棵受限回归树。对伪残差作最小二乘拟合，就是在可表示的树函数中近似这个方向。若按固定非空叶划分取伪残差平均值，拟合值向量 $\boldsymbol h$ 满足
\[
 \boldsymbol r_m^{\mathsf T}\boldsymbol h
 =\sum_{j=1}^{J_m}n_{mj}\bar r_{mj}^2\geq0,
 \qquad
 \bar r_{mj}=\frac1{n_{mj}}\sum_{x_i\in R_{mj}}r_{mi}.
\]
因此只有在右侧严格为正时，这个方向才给出严格负的方向导数，并保证足够小的正步长下降；若所有叶平均值都为零，该树划分没有捕获非零的一阶下降分量。

\textbf{梯度提升树算法。}输入为训练集、对第二变量可微的损失 $L$、允许的树复杂度及最大轮数 $M$；以下假设所写的一维最小化都有有限极小点，选择其中任一个即可。算法输出一个树的加法模型，步骤如下：
\begin{enumerate}
 \item 以最优常数初始化，取
 \begin{equation}
  c_0\in\operatorname*{arg\,min}_{c\in\mathbb R}\sum_{i=1}^{N}L(y_i,c),\qquad
  f_0(x)=c_0.
  \label{eq:lh-boost-gradient-initial}
 \end{equation}
 \item 对 $m=1,\ldots,M$，依次进行本轮更新：按式~\eqref{eq:lh-boost-pseudo-residual} 计算伪残差，以 $(x_i,r_{mi})$ 为数据拟合平方误差回归树，得到其叶区域 $R_{mj}$，$j=1,\ldots,J_m$；再在每个叶区域内回到\emph{原损失}，分别作一维搜索：
 \begin{equation}
  c_{mj}\in\operatorname*{arg\,min}_{c\in\mathbb R}
  \sum_{x_i\in R_{mj}}L\bigl(y_i,f_{m-1}(x_i)+c\bigr).
  \label{eq:lh-boost-gradient-leaf-search}
 \end{equation}
 随后更新
 \begin{equation}
  f_m(x)=f_{m-1}(x)+\sum_{j=1}^{J_m}c_{mj}I(x\in R_{mj}).
  \label{eq:lh-boost-gradient-update}
 \end{equation}
 \item 达到最大轮数或预定停止条件时，输出
 \begin{equation}
  f_M(x)=c_0+\sum_{m=1}^{M}\sum_{j=1}^{J_m}c_{mj}I(x\in R_{mj}),
  \label{eq:lh-boost-gradient-output}
 \end{equation}
 提前停止时把 $M$ 替换为实际完成轮数。
\end{enumerate}

每轮先通过拟合伪残差选择区域划分，再用原损失确定这些区域上的实际增量；最小二乘拟合伪残差得到的叶均值一般不等于原损失下的最优叶增量。因为各叶互不相交，式~\eqref{eq:lh-boost-gradient-leaf-search} 中的优化可以分别进行；精确求解时，每个叶都可以选择 $c=0$，所以式~\eqref{eq:lh-boost-gradient-update} 不会使训练损失增加。这一结论依赖叶内优化确实求到不劣于零增量的值，并不保证得到全局最优树结构或整个函数空间中的最优模型。

对于 $L(y,t)=\tfrac12(y-t)^2$，负梯度恰为 $y-t$，叶内搜索也给出平均残差，因而恢复平方损失的残差拟合。若按教材使用 $L(y,t)=(y-t)^2$，负梯度实际为 $2(y-t)$；它与残差仅相差公共因子 $2$，不是逐字相等。在精确平方误差拟合、叶数等约束不变且约束不依赖响应尺度时，这个因子不改变最优划分，随后的原损失叶内搜索仍给出残差平均值。实践中若划分停止阈值等依赖响应尺度，则也需要相应调整，不能忽略该因子对具体实现的影响。

以二分类指数损失为例，$-\partial L(y,t)/\partial t=ye^{-yt}$，所以伪残差是带符号的实数；回归树拟合的是这些实数，而 AdaBoost 直接最小化二值规则的加权误差。因此，二者共有逐步降低损失的思路，但使用一般回归树的梯度提升步骤不必逐轮等于 AdaBoost。对于绝对损失 $|y-t|$，在 $t=y$ 处不可微，不能直接无条件套用上述可微公式；若另行规定次梯度与相应的优化步骤，则叶内最优增量可取当前叶残差的中位数。

\begin{notealert}[type=caution,title={初始化、有限极小点与输出形式}]
最优常数 $c_0$ 一般不为零，最终模型必须保留式~\eqref{eq:lh-boost-gradient-output} 中的初始常数项。某些损失与数据组合的一维目标只有下确界而没有有限极小点，例如单一类别样本上的指数损失；这时需另行限制搜索范围或规定数值处理，不能把不存在的有限 $\operatorname*{arg\,min}$ 当作已求得。算法输出的是多棵树的加法模型；若用于分类，尚需明确从得分到类别的决策规则。
\end{notealert}

以上为教材所述的基本梯度提升步骤。常见的扩展是在更新中乘收缩系数（shrinkage factor）$0<\nu\leq1$，即把增量改为 $\nu\sum_jc_{mj}I(x\in R_{mj})$，再通过验证损失选择树数。$\nu<1$ 时，不再是沿各叶增量一次走到该轮的极小点；若叶内损失关于 $c$ 凸，且 $c_{mj}$ 是有限极小点，则由凸性可知这一步仍不增加训练损失，对一般非凸损失则需另行检查步长。这一扩展不能由前面的 AdaBoost 误差界直接获得泛化保证。

\section{方法联系与习题讨论}
\label{sec:lh-boost-review}

教材的招聘分类习题要求把一维例题推广到多个属性。使用树桩时，先明确允许的划分规则：二值属性只有一种非平凡二分，有序属性可以在相邻等级间设阈值，并枚举两侧预测类别；若把属性当作无序类别，则应另外规定允许哪些类别子集划分。每轮用当前 $D_m$ 计算所有候选树桩的加权错误率，选择最优者，再按式~\eqref{eq:lh-boost-alpha} 与式~\eqref{eq:lh-boost-weight-update} 更新。手算时应逐轮记录误分样本、加权错误率、系数及组合模型的错误率；基本分类器的当前误差与组合模型的当前误差是两个不同的量，不能相互代用。

比较支持向量机、AdaBoost 与逻辑斯谛回归时，可先统一二分类记号 $y\in\{-1,+1\}$，设实值分类得分为 $t=f(x)$。三者的典型分类损失分别是合页损失（hinge loss）$\max\{0,1-yt\}$、指数损失 $e^{-yt}$ 与逻辑斯谛损失（logistic loss）$\log(1+e^{-yt})$。软间隔支持向量机通常把合页损失与控制参数范数的正则项结合，通过凸二次规划、对偶或相应优化方法学习；逻辑斯谛回归把概率模型的负对数似然最小化，在线性得分下形成光滑凸问题，可用梯度法、牛顿法或拟牛顿法；本章 AdaBoost 则在基本分类器集合上逐轮构造加法模型，以重加权和一维系数优化实现下降。比较时还要固定或明确各自的函数空间与正则化方式，不能把模型表达能力的差异全部归因于损失函数。这里的损失对照采用同一得分尺度；备注~\ref{rem:lh-boost-conditional-risk} 中指数损失的最优得分是对数优势的一半，逻辑斯谛模型通常采用的得分尺度则是完整的对数优势。

本章的统一关系是：选择基本函数集合确定模型可表达什么，选择损失确定“尚未拟合好”的含义，前向分步策略确定每轮如何添加修正。AdaBoost 用指数损失产生样本权重；平方损失提升树把下一轮目标化为残差；梯度提升进一步以一般损失的负梯度选择树划分，并用原损失确定叶增量。训练损失下降、训练分类正确与新样本预测可靠是三个需要分别判断的问题，阅读证明和实现算法时都应保持这一区分。
% ===== 李航教材第 8 章笔记结束 =====

% 研读来源：机器学习/textbooks/机器学习（分章原文）/机器学习-9.pdf，全部 18 页。
% 本章对应教材第 9 章；笔记中的章节、公式和定理均按本文结构自动编号。

\fi

\selectchapter{9}
\ifcompilechapter
\chapter{EM 算法及其推广}
\label{chap:lihang-em}

在混合分布、缺失数据和潜在类别等问题中，观测数据的分布往往需要对未观测变量求和或积分。直接极大化其似然可能很困难，而假如这些变量也已知，参数估计反而容易。\emph{期望极大算法}（expectation maximization algorithm，EM 算法）利用这一差别：先根据当前模型计算隐变量的条件分布，再极大化完全数据对数似然的条件期望。本章对应教材第 9 章，依次讨论算法的构造、单调性与收敛条件、高斯混合模型中的具体更新，以及基于辅助函数的推广。

本章以 $Y$ 表示观测随机变量或观测随机向量，$y$ 表示已经取得并固定的数据；以 $Z$ 表示隐变量，$z$ 表示其可能取值。参数为 $\theta\in\Theta\subseteq\mathbb R^d$，$p_\theta(y,z)$ 统一表示联合概率质量函数或概率密度函数。除特别说明外，求和形式的理论推导先在有限隐状态、关于参数不变的正支撑上进行，并要求观测似然为正、涉及的期望有限；连续隐变量的情形需将求和换成积分，并另行保证可积性及交换微分与积分的条件。

\section{EM 算法的引入}
\label{sec:lihang-em-introduction}

\subsection{EM 算法}
\label{subsec:lihang-em-algorithm}

\emph{观测变量}（observable variable）可由数据取得，\emph{隐变量或潜在变量}（hidden variable or latent variable）则在所用模型中存在而没有被观测。二者共同组成的 $(y,z)$ 称为\emph{完全数据}（complete data），只有 $y$ 时称为\emph{不完全数据}（incomplete data）。这些名称相对于所选模型而言，并不表示观测数据本身存在记录错误。参数估计的目标仍然是观测数据的极大似然估计（maximum likelihood estimation，MLE）：
\begin{equation}
  p_\theta(y)=\sum_z p_\theta(y,z),\qquad
  \ell(\theta)=\log p_\theta(y),\qquad
  \widehat\theta\in\operatorname*{arg\,max}_{\theta\in\Theta}\ell(\theta).
  \label{eq:lihang-em-observed-objective}
\end{equation}
主要计算困难常来自对数内部的求和；这并不意味着一切含隐变量的模型都没有解析解，也不意味着把隐变量直接补成一个确定值就能得到正确估计。

\begin{example}
\label{ex:lihang-em-three-coins}
教材例 9.1 的三硬币模型。硬币 A、B、C 的正面概率分别为 $\pi,p,q$。每次先掷 A，正面选 B，反面选 C，再掷所选硬币，只记录第二次是否为正面。独立重复 $n=10$ 次，观测序列为
\[
  y=(1,1,0,1,0,0,1,0,1,1).
\]
建立概率模型，推导 EM 更新，并说明这些数据能够确定哪些参数信息。
\end{example}

\begin{solve}
令 $Z_i=1$ 表示第 $i$ 次选 B，$Z_i=0$ 表示选 C，$Y_i\in\{0,1\}$ 表示记录结果。模型假设为
\[
  P_\theta(Z_i=1)=\pi,\qquad
  P_\theta(Y_i=1\mid Z_i=1)=p,\qquad
  P_\theta(Y_i=1\mid Z_i=0)=q,
\]
其中 $\theta=(\pi,p,q)\in[0,1]^3$。边缘化 $Z_i$，得到
\begin{equation}
  p_\theta(y_i)=\pi p^{y_i}(1-p)^{1-y_i}
  +(1-\pi)q^{y_i}(1-q)^{1-y_i},\qquad y_i\in\{0,1\}.
  \label{eq:lihang-em-coin-marginal}
\end{equation}
若每次选了哪枚硬币也已知，完全数据的对数似然就是
\[
  \ell_{\mathrm c}(\theta;y,z)
  =\sum_{i=1}^n\bigl\{
  z_i[\log\pi+y_i\log p+(1-y_i)\log(1-p)]
  +(1-z_i)[\log(1-\pi)+y_i\log q+(1-y_i)\log(1-q)]\bigr\}.
\]
它将三个参数分别放入伯努利分布的对数似然中，因而容易极大化。EM 用条件期望替代这里未知的 $z_i$。在第 $t$ 轮参数 $\theta^{(t)}$ 下，根据贝叶斯公式计算
\begin{equation}
  r_i^{(t)}=P_{\theta^{(t)}}(Z_i=1\mid Y_i=y_i)
  =\frac{\pi^{(t)}(p^{(t)})^{y_i}(1-p^{(t)})^{1-y_i}}
  {\pi^{(t)}(p^{(t)})^{y_i}(1-p^{(t)})^{1-y_i}
   +(1-\pi^{(t)})(q^{(t)})^{y_i}(1-q^{(t)})^{1-y_i}}.
  \label{eq:lihang-em-coin-responsibility}
\end{equation}
由于样本对 $(Y_i,Z_i)$ 相互独立，给定全部观测时该条件概率也只依赖 $y_i$。固定这些概率，将 $\ell_{\mathrm c}$ 中的 $z_i$ 换成 $r_i^{(t)}$，便得到待极大化的条件期望。对 $\pi$ 求导的方程为
\[
  \frac{\sum_i r_i^{(t)}}{\pi}
  -\frac{\sum_i(1-r_i^{(t)})}{1-\pi}=0;
\]
对 $p,q$ 分别求导同样得到加权频率。因此新参数为
\begin{equation}
  \begin{aligned}
    \pi^{(t+1)}&=\frac1n\sum_{i=1}^n r_i^{(t)},\\
    p^{(t+1)}&=\frac{\sum_{i=1}^n r_i^{(t)}y_i}{\sum_{i=1}^n r_i^{(t)}},\qquad
    q^{(t+1)}=\frac{\sum_{i=1}^n(1-r_i^{(t)})y_i}{\sum_{i=1}^n(1-r_i^{(t)})}.
  \end{aligned}
  \label{eq:lihang-em-coin-update}
\end{equation}
三个加权伯努利对数似然均为凹函数；相应的成功与失败权重都为正时，上述驻点就是各自唯一的最大点。若某一成分的总权重为零，该成分的正面概率无法由本轮目标确定；若只有成功或只有失败权重，则最大值在参数边界上。取三个参数的初值都严格位于 $(0,1)$ 内，可使本例两种观测下的响应概率均有定义。

从 $(0.5,0.5,0.5)$ 出发，所有 $r_i^{(0)}=0.5$，第一轮得到 $(0.5,0.6,0.6)$，此后保持不变。从 $(0.4,0.6,0.7)$ 出发，第一轮得到约 $(0.406417,0.536842,0.643243)$，也已达到本例的一个固定点。理解这些不同答案，需要注意观测分布实际上只是正面概率为
\[
  \rho=\pi p+(1-\pi)q
\]
的伯努利分布。这十次观测中有六次正面，故
\[
  \ell(\theta)=6\log\rho+4\log(1-\rho),\qquad
  \frac{\mathrm d\ell}{\mathrm d\rho}=\frac6\rho-\frac4{1-\rho}.
\]
该函数在 $0<\rho<1$ 上严格凹，最大点为 $\widehat\rho=0.6$。所以满足 $\pi p+(1-\pi)q=0.6$ 的全部可行三元组都是极大似然解；观测分布不能唯一确定 $\pi,p,q$，称为\emph{参数不可辨识}（non-identifiability）。这里可直接写出极大似然解集，采用 EM 的意义在于展示隐变量迭代机制。特别地，$\widehat\pi=0.5$ 是某次初始化所得解的一部分，不能据此断言数据证明了硬币 A 均匀。
\end{solve}

\begin{definition}[$Q$ 函数]
\label{def:lihang-em-q-function}
固定观测 $y$ 和当前参数 $\theta^{(t)}$，定义
\begin{equation}
  Q(\theta\mid\theta^{(t)})
  =\mathbb E_{\theta^{(t)}}[\log p_\theta(y,Z)\mid Y=y]
  =\sum_z p_{\theta^{(t)}}(z\mid y)\log p_\theta(y,z).
  \label{eq:lihang-em-q-function}
\end{equation}
它是完全数据对数似然关于当前隐变量后验分布的期望。第一个参数是待优化变量，竖线后的参数只确定本轮使用的条件分布。
\end{definition}

一般 EM 算法的输入是观测 $y$、联合模型 $p_\theta(y,z)$、可行参数域 $\Theta$，以及计算隐变量条件分布的方法；隐变量的实际取值不是已知输入。选择使 $p_{\theta^{(0)}}(y)>0$ 且期望存在的初值，规定正容差 $\varepsilon$ 与最大迭代次数 $T$，然后进行如下迭代：
\begin{enumerate}
  \item \emph{E 步}（expectation step）：利用当前后验 $p_{\theta^{(t)}}(z\mid y)$，构造式~\eqref{eq:lihang-em-q-function} 中的 $Q(\theta\mid\theta^{(t)})$。
  \item \emph{M 步}（maximization step）：在可行域内求
  \begin{equation}
    \theta^{(t+1)}\in\operatorname*{arg\,max}_{\theta\in\Theta}
    Q(\theta\mid\theta^{(t)}).
    \label{eq:lihang-em-m-step}
  \end{equation}
  \item 计算观测对数似然。若
  \[
    |\ell(\theta^{(t+1)})-\ell(\theta^{(t)})|
    \leq\varepsilon\bigl(1+|\ell(\theta^{(t)})|\bigr),
  \]
  则按所定精度停止；达到 $T$ 时也终止，并记录是否满足容差。输出最后的参数及其目标值。
\end{enumerate}
式~\eqref{eq:lihang-em-m-step} 的极大值需要存在且能够求得；不能达到精确极大值时，可使用第~\ref{subsec:lihang-em-gem} 小节的广义更新。E 步求的是关于隐变量的条件期望，通常仍是 $\theta$ 的函数，并非只计算一个数；它也不是先把隐变量固定成概率最大的类别。停止条件用于控制数值计算，似然变化很小并不构成全局最优的证明。若采用参数变化准则，还需考虑各参数的尺度。

\begin{notealert}[type=tip,title={从完全数据构造算法}]
推导具体 EM 算法时，先选定隐变量并写出联合模型，再计算完全数据对数似然，最后求其中隐变量相关项的后验期望。若目标是极大后验估计（maximum a posteriori estimation，MAP），对给定先验密度 $p(\theta)$，M 步改为极大化 $Q(\theta\mid\theta^{(t)})+\log p(\theta)$；E 步仍使用给定当前参数的隐变量后验，单调性所针对的目标则是 $\ell(\theta)+\log p(\theta)$。
\end{notealert}

\subsection{EM 算法的导出}
\label{subsec:lihang-em-derivation}

EM 的关键不是让完全数据和观测数据的似然成为同一个函数，而是构造一个在当前参数处与目标相等的下界。记本轮隐变量后验为 $w_t(z)=p_{\theta^{(t)}}(z\mid y)$。在共同正支撑的假设下，用对数函数的凹性，即 Jensen 不等式，得到
\begin{equation}
  \begin{aligned}
    \ell(\theta)
    &=\log\sum_z w_t(z)\frac{p_\theta(y,z)}{w_t(z)}\\
    &\geq\sum_z w_t(z)\log\frac{p_\theta(y,z)}{w_t(z)}
    =:B(\theta\mid\theta^{(t)}).
  \end{aligned}
  \label{eq:lihang-em-jensen-bound}
\end{equation}
这里 $w_t(z)\geq0$ 且 $\sum_z w_t(z)=1$。令
\[
  H(w_t)=-\sum_z w_t(z)\log w_t(z),
\]
则下界可以写成
\begin{equation}
  B(\theta\mid\theta^{(t)})=Q(\theta\mid\theta^{(t)})+H(w_t),\qquad
  B(\theta^{(t)}\mid\theta^{(t)})=\ell(\theta^{(t)}).
  \label{eq:lihang-em-touching-bound}
\end{equation}
第二个等式来自 $p_{\theta^{(t)}}(y,z)/w_t(z)=p_{\theta^{(t)}}(y)$：在当前参数处，Jensen 不等式中的被取平均量不随 $z$ 变化，因而取等号。固定 $w_t$ 后，$H(w_t)$ 与待优化参数 $\theta$ 无关，所以极大化 $B$ 正好等价于 M 步极大化 $Q$。

每一轮都先在当前点建立一个接触目标的下界，再提高这个下界，随后用新参数重新建立下界。E 步的期望若可精确求出，就不涉及随机近似；算法对于原目标的局限在于这种逐轮提高下界的过程通常只找到某个候选解，未必完成原似然的全局极大化。用 $B$ 的几何解释时，必须同时保留下界关系与当前点取等号这两个条件，单独提高一个任意下界不足以推出目标值提高。

\subsection{EM 算法在无监督学习中的应用}
\label{subsec:lihang-em-unsupervised}

无监督学习没有给出每个观测对应的类别或其他目标输出。设输入为 $X$，潜在类别为 $C$，生成模型为 $p_\theta(x,c)$。假设各样本独立同分布，则由观测样本 $x_1,\ldots,x_n$ 估计参数的目标是
\[
  \sum_{i=1}^n\log p_\theta(x_i)
  =\sum_{i=1}^n\log\sum_c p_\theta(x_i,c).
\]
将未给出的类别看作隐变量，就得到 EM 可以处理的形式。E 步给每个可能类别分配后验权重，M 步用这些权重重新估计生成模型；高斯混合模型和朴素贝叶斯模型的无监督学习都是具体例子。EM 是参数估计方法，其适用范围也包括缺失数据等问题；无监督任务能否采用 EM，取决于所设联合模型以及 E、M 两步是否可计算。

\section{EM 算法的收敛性}
\label{sec:lihang-em-convergence}

\subsection{观测似然的单调性}

\begin{theorem}[提高 $Q$ 函数足以保持似然单调]
\label{thm:lihang-em-monotonicity}
在本章下界构造的条件下，若一次参数更新满足
\[
  Q(\theta^{(t+1)}\mid\theta^{(t)})\geq Q(\theta^{(t)}\mid\theta^{(t)}),
\]
则 $\ell(\theta^{(t+1)})\geq\ell(\theta^{(t)})$，从而观测似然也不减。精确的 EM 更新满足上述条件。
\end{theorem}

\begin{proof}
由式~\eqref{eq:lihang-em-jensen-bound} 和式~\eqref{eq:lihang-em-touching-bound}，
\begin{equation}
  \begin{aligned}
    \ell(\theta^{(t+1)})
    &\geq B(\theta^{(t+1)}\mid\theta^{(t)})\\
    &=Q(\theta^{(t+1)}\mid\theta^{(t)})+H(w_t)\\
    &\geq Q(\theta^{(t)}\mid\theta^{(t)})+H(w_t)
    =\ell(\theta^{(t)}).
  \end{aligned}
  \label{eq:lihang-em-monotonicity-chain}
\end{equation}
精确 M 步的可行候选中包含 $\theta^{(t)}$，所以其最大值不会小于当前值。指数函数严格递增，再将对数似然的不等式取指数即可。
\end{proof}

还可以把每次增量分解，以看清何时取等号。对两个有限状态分布 $a,b$，定义\emph{相对熵或 KL 散度}（Kullback--Leibler divergence）
\[
  D_{\mathrm{KL}}(a\Vert b)=\sum_z a(z)\log\frac{a(z)}{b(z)}.
\]
约定 $0\log(0/b)=0$，若 $a(z)>0$ 而 $b(z)=0$，则散度为 $+\infty$。由 $\log u\leq u-1$，在散度有限时有
\[
  -D_{\mathrm{KL}}(a\Vert b)
  =\sum_{z:a(z)>0}a(z)\log\frac{b(z)}{a(z)}
  \leq\sum_{z:a(z)>0}b(z)-1\leq0.
\]
因此散度非负，且当且仅当两个分布相同才为零。定义教材证明所用的条件对数概率期望
\[
  H(\theta\mid\theta^{(t)})
  =\sum_z w_t(z)\log p_\theta(z\mid y).
\]
它不是前面的分布熵 $H(w_t)$；两者在当前点满足 $H(\theta^{(t)}\mid\theta^{(t)})=-H(w_t)$。由联合概率与条件概率的分解，
\[
  \ell(\theta)=Q(\theta\mid\theta^{(t)})-H(\theta\mid\theta^{(t)}),
\]
所以在各项有限时得到增量恒等式
\begin{equation}
  \begin{aligned}
    \ell(\theta^{(t+1)})-\ell(\theta^{(t)})
    &=Q(\theta^{(t+1)}\mid\theta^{(t)})-Q(\theta^{(t)}\mid\theta^{(t)})\\
    &\quad+D_{\mathrm{KL}}\!\left(w_t\Vert p_{\theta^{(t+1)}}(\,\cdot\mid y)\right).
  \end{aligned}
  \label{eq:lihang-em-increment-identity}
\end{equation}
在提高 $Q$ 的更新中，右边两项均非负；似然保持不变当且仅当 $Q$ 没有提高且新旧隐变量后验相同。因此 EM 的普遍保证是\emph{单调不减}，并不是每一步都严格增大，也不是一旦不变就已经得到唯一的极大值。

\subsection{目标值收敛、参数收敛与驻点}
\label{subsec:lihang-em-stationarity}

\begin{proposition}[目标值的收敛]
\label{prop:lihang-em-value-convergence}
设迭代从有限的 $\ell(\theta^{(0)})$ 出发，且存在有限常数 $M$ 使所有迭代点都满足 $\ell(\theta^{(t)})\leq M$。则 EM 产生的对数似然序列收敛到某个有限值 $\ell^*$。
\end{proposition}

\begin{proof}
由定理~\ref{thm:lihang-em-monotonicity}，该序列单调不减；再由给定上界，实数单调有界定理保证其收敛。
\end{proof}

这个命题只涉及函数值。相同似然可能对应不同参数，参数序列也可能没有唯一极限；例~\ref{ex:lihang-em-three-coins} 中的不可辨识性已经说明不能从最优似然唯一反推参数。对于连续分布，密度值没有概率质量不超过 $1$ 的限制，观测似然也未必有上界；高斯混合模型中就会出现这一问题。

教材关于参数极限为稳定点的结论需要额外的正则条件。下面给出一组便于核查的充分条件，只证明\emph{已经收敛的参数序列}具有什么性质。

\begin{proposition}[参数极限为驻点的一组充分条件]
\label{prop:lihang-em-stationary-limit}
设 $\Theta\subseteq\mathbb R^d$ 为开集，隐状态有限，$p_\theta(y,z)>0$ 且关于 $\theta$ 连续可微，$\nabla_\theta Q(\theta\mid\eta)$ 关于 $(\theta,\eta)$ 连续。若精确 EM 序列满足 $\theta^{(t)}\to\theta^*\in\Theta$，且各 M 步在域内取得最大值，则 $\nabla\ell(\theta^*)=0$。
\end{proposition}

\begin{proof}
域内的 M 步最大点满足 $\nabla_\theta Q(\theta^{(t+1)}\mid\theta^{(t)})=0$。令 $t\to\infty$，由梯度连续性得到 $\nabla_\theta Q(\theta^*\mid\theta^*)=0$。另一方面，有限求和允许逐项微分，故
\begin{equation}
  \begin{aligned}
    \nabla\ell(\theta)
    &=\frac{\sum_z\nabla p_\theta(y,z)}{p_\theta(y)}\\
    &=\sum_z p_\theta(z\mid y)\nabla\log p_\theta(y,z)
    =\left.\nabla_\vartheta Q(\vartheta\mid\theta)\right|_{\vartheta=\theta}.
  \end{aligned}
  \label{eq:lihang-em-score-identity}
\end{equation}
取 $\theta=\theta^*$，即得结论。
\end{proof}

这里的驻点（stationary point）仅表示梯度为零，可能是局部极大点、鞍点或退化的平坦点，不能直接理解为全局最优点。若参数受约束，应在可行方向上陈述相应的一阶必要条件；若隐变量连续，还必须增加使式~\eqref{eq:lihang-em-score-identity} 中微分与积分交换成立的条件。实际计算可采用多个可行初值，比较同一目标下得到的有限解，并检查约束、退化情形与停止状态；这些操作提高找到较好解的机会，但不构成一般的全局最优保证。

\section{EM 算法在高斯混合模型学习中的应用}
\label{sec:lihang-em-gmm}

\subsection{高斯混合模型}
\label{subsec:lihang-em-gmm-model}

\begin{definition}[一维高斯混合模型]
\label{def:lihang-em-gmm}
给定正整数 $K$，\emph{高斯混合模型}（Gaussian mixture model，GMM）的密度为
\begin{equation}
  p_\theta(y)=\sum_{k=1}^K\alpha_k\phi(y;\mu_k,\sigma_k^2),\qquad
  \phi(y;\mu_k,\sigma_k^2)
  =\frac1{\sqrt{2\pi\sigma_k^2}}
   \exp\left(-\frac{(y-\mu_k)^2}{2\sigma_k^2}\right),
  \label{eq:lihang-em-gmm-density}
\end{equation}
其中 $\alpha_k\geq0$、$\sum_{k=1}^K\alpha_k=1$、$\mu_k\in\mathbb R$、$\sigma_k^2>0$。参数 $\theta$ 包含各混合系数、均值与方差；成分数 $K$ 在本章的参数估计中预先给定。
\end{definition}

模型可用两步生成过程解释：先以概率 $\alpha_k$ 选择成分 $Z_i=k$，再按该成分的高斯分布生成 $Y_i$。假设 $n$ 个样本独立同分布，则
\[
  p_\theta(y_i,Z_i=k)=\alpha_k\phi(y_i;\mu_k,\sigma_k^2),\qquad
  \ell(\theta)=\sum_{i=1}^n\log\sum_{k=1}^K
  \alpha_k\phi(y_i;\mu_k,\sigma_k^2).
\]
若直接对观测对数似然求导，各参数会被成分求和耦合。引入成分指示随机变量 $W_{ik}=\mathbf1_{\{Z_i=k\}}$，其取值记为 $w_{ik}$；每个样本满足 $\sum_{k=1}^K w_{ik}=1$。完全数据的对数似然因此变为
\begin{equation}
  \ell_{\mathrm c}(\theta;y,w)
  =\sum_{i=1}^n\sum_{k=1}^K w_{ik}
  \left[\log\alpha_k-\frac12\log(2\pi\sigma_k^2)
  -\frac{(y_i-\mu_k)^2}{2\sigma_k^2}\right].
  \label{eq:lihang-em-gmm-complete-loglik}
\end{equation}
以下推导在混合系数为正、各方差为正的迭代点进行，因而观测密度及后验分母均为正。

\subsection{高斯混合模型参数估计的 EM 算法}
\label{subsec:lihang-em-gmm-updates}

E 步计算每个成分对每个观测的\emph{响应度}（responsibility）：
\begin{equation}
  \gamma_{ik}^{(t)}
  =\mathbb E_{\theta^{(t)}}[W_{ik}\mid Y=y]
  =\frac{\alpha_k^{(t)}\phi(y_i;\mu_k^{(t)},(\sigma_k^2)^{(t)})}
  {\sum_{h=1}^K\alpha_h^{(t)}\phi(y_i;\mu_h^{(t)},(\sigma_h^2)^{(t)})},\qquad
  \sum_{k=1}^K\gamma_{ik}^{(t)}=1.
  \label{eq:lihang-em-gmm-responsibility}
\end{equation}
这里 $\gamma_{ik}^{(t)}$ 是确定的条件概率，与尚未观测到的随机指示量 $W_{ik}$ 不同。记本轮的成分有效样本数为 $N_k^{(t)}=\sum_{i=1}^n\gamma_{ik}^{(t)}$，有 $\sum_kN_k^{(t)}=n$；它通常不是整数。将式~\eqref{eq:lihang-em-gmm-complete-loglik} 中的指示量换成其条件期望，得到
\begin{equation}
  Q(\theta\mid\theta^{(t)})
  =\sum_{k=1}^K\left[
  N_k^{(t)}\log\alpha_k-\frac{N_k^{(t)}}2\log(2\pi\sigma_k^2)
  -\frac1{2\sigma_k^2}\sum_{i=1}^n\gamma_{ik}^{(t)}(y_i-\mu_k)^2\right].
  \label{eq:lihang-em-gmm-q}
\end{equation}

M 步首先对混合系数进行约束极大化。对 $\sum_kN_k^{(t)}\log\alpha_k$ 引入约束 $\sum_k\alpha_k=1$ 的拉格朗日乘子，使用
\[
  \mathcal L(\alpha,\lambda)
  =\sum_kN_k^{(t)}\log\alpha_k+\lambda\left(1-\sum_k\alpha_k\right),
\]
一阶条件给出 $N_k^{(t)}/\alpha_k=\lambda$。求和得 $\lambda=n$，从而 $\alpha_k^{(t+1)}=N_k^{(t)}/n$。当所有 $N_k^{(t)}>0$ 时，目标关于混合系数严格凹，所以这是单纯形内的唯一最大点。

其次，对第 $k$ 个成分记 $S_k(\mu)=\sum_i\gamma_{ik}^{(t)}(y_i-\mu)^2$。固定方差时，极大化 $Q$ 等价于极小化 $S_k$，求导得到新均值
\[
  \overline y_k^{(t)}=\frac{\sum_i\gamma_{ik}^{(t)}y_i}{N_k^{(t)}}.
\]
恒等式 $S_k(\mu)=S_k(\overline y_k^{(t)})+N_k^{(t)}(\mu-\overline y_k^{(t)})^2$ 说明，对任意正方差，这个均值都使对应成分的 $Q$ 最大。再令 $v=\sigma_k^2$，在这个均值处有
\[
  \frac{\partial Q}{\partial v}
  =-\frac{N_k^{(t)}}{2v}+\frac{S_k(\overline y_k^{(t)})}{2v^2}.
\]
若 $S_k(\overline y_k^{(t)})>0$，导数在 $v=S_k/N_k^{(t)}$ 两侧由正变负，该点就是关于 $v>0$ 的唯一最大点。合并以上结果，得到
\begin{equation}
  \begin{aligned}
    \alpha_k^{(t+1)}&=\frac{N_k^{(t)}}n,\\
    \mu_k^{(t+1)}&=\frac{\sum_{i=1}^n\gamma_{ik}^{(t)}y_i}{N_k^{(t)}},\\
    (\sigma_k^2)^{(t+1)}
    &=\frac{\sum_{i=1}^n\gamma_{ik}^{(t)}(y_i-\mu_k^{(t+1)})^2}{N_k^{(t)}},
    \qquad k=1,\ldots,K.
  \end{aligned}
  \label{eq:lihang-em-gmm-m-step}
\end{equation}
\emph[caution]{新方差必须围绕本轮更新后的均值计算}。所有响应度均来自旧参数，但在同一次 M 步中，方差公式使用 $\mu_k^{(t+1)}$。分母是有效样本数 $N_k^{(t)}$，因为这里求极大似然估计，不能换成无偏样本方差中的自由度修正。若 $N_k^{(t)}=0$，该成分的均值和方差在本轮不可由 $Q$ 确定；若 $S_k=0$，正方差域内的最大值不存在，不能把方差零当作正常高斯分布。

具体计算时，输入样本、给定的 $K$、正混合系数与正方差的初值、容差和最大轮数；交替使用式~\eqref{eq:lihang-em-gmm-responsibility} 与式~\eqref{eq:lihang-em-gmm-m-step}，每轮计算观测对数似然并检查约束及停止条件，最后输出参数、似然与停止状态。一维情形下，一轮计算全部响应度和加权统计量需要 $O(nK)$ 次标量运算；若存储整个响应度矩阵，需要 $O(nK)$ 存储空间。密度很小时，可在对数域计算 $b_{ik}=\log\alpha_k^{(t)}+\log\phi(y_i;\mu_k^{(t)},(\sigma_k^2)^{(t)})$，令 $m_i=\max_k b_{ik}$，使用
\[
  \log\sum_k e^{b_{ik}}=m_i+\log\sum_k e^{b_{ik}-m_i}
\]
归一化响应度并求似然，以减轻直接计算很小密度造成的数值下溢。

\begin{example}
\label{ex:lihang-em-gmm-one-step}
对样本 $y=(-1,0,1)$ 拟合两个高斯成分。取初值 $\alpha_1^{(0)}=\alpha_2^{(0)}=1/2$、$\mu_1^{(0)}=-1$、$\mu_2^{(0)}=1$、两个方差都为 $1$，计算一轮 EM 更新。
\end{example}

\begin{solve}
两个初始密度的比值为 $\phi(y;-1,1)/\phi(y;1,1)=e^{-2y}$，所以第一成分的响应度依次为
\[
  \gamma_{11}^{(0)}=\frac{e^2}{1+e^2},\qquad
  \gamma_{21}^{(0)}=\frac12,\qquad
  \gamma_{31}^{(0)}=\frac1{1+e^2}.
\]
两个有效样本数均为 $3/2$，故混合系数不变。令 $a=(e^2-1)/(e^2+1)$，更新后的均值和方差为
\[
  \mu_1^{(1)}=-\frac{2a}{3}\approx-0.507729,\qquad
  \mu_2^{(1)}=\frac{2a}{3}\approx0.507729,
\]
\[
  (\sigma_1^2)^{(1)}=(\sigma_2^2)^{(1)}
  =\frac23-\frac{4a^2}{9}\approx0.408877.
\]
方差计算使用加权二阶矩减去新均值平方。该轮中间样本同时参与两个成分的估计，具体体现了响应度作为软分配权重的含义。
\end{solve}

\begin{notealert}[type=caution,title={高斯混合模型的似然可能无上界}]
当 $K\geq2$ 且各成分方差可任意趋近于零时，可以固定一个正的混合系数，将对应均值放在某个观测 $y_i$ 上，并令其标准差 $\sigma\to0$。该点的高斯密度为 $1/(\sqrt{2\pi}\sigma)\to\infty$，而另一个保持固定的正权重高斯成分可使其他观测的混合密度保持正的下界。因此总对数似然可趋于 $+\infty$，命题~\ref{prop:lihang-em-value-convergence} 的上界假设在无约束 GMM 中并不自动成立。

若模型明确约束 $\sigma_k^2\geq v_{\min}>0$，其中 $v_{\min}$ 是事先给定的方差下界，则对应的约束 M 步取 $\max\{v_{\min},S_k(\mu_k^{(t+1)})/N_k^{(t)}\}$，仍可按同一约束目标证明单调性；采用先验则对应另一种 MAP 目标。随意重置退化成分的参数未必提高原来的 $Q$，需要重新核查。即使排除了方差退化，交换成分编号仍不改变混合密度，因此类别编号本身也不具有唯一的统计含义。
\end{notealert}

\section{EM 算法的推广}
\label{sec:lihang-em-generalizations}

\subsection{$F$ 函数的极大--极大算法}
\label{subsec:lihang-em-f-function}

Jensen 下界中的后验权重也可以视为待优化的分布变量，从而将 EM 写成两个变量组之间的交替极大化。以下性质仍先在有限隐状态的概率单纯形上陈述；涉及局部极值的连续性结论不直接从有限情形推广到任意分布空间。

\begin{definition}[$F$ 函数]
\label{def:lihang-em-f-function}
对隐状态上的任意概率分布 $r=(r(z))$，定义其熵 $H(r)=-\sum_zr(z)\log r(z)$，并令
\begin{equation}
  F(r,\theta)=\mathbb E_r[\log p_\theta(y,Z)]+H(r)
  =\sum_zr(z)\log\frac{p_\theta(y,z)}{r(z)}.
  \label{eq:lihang-em-f-definition}
\end{equation}
在联合概率为正的有限状态情形，约定 $r(z)=0$ 的项为零。$r$ 是辅助分布，$\theta$ 是生成模型参数，二者承担不同的角色。
\end{definition}

\begin{lemma}[$F$ 的最大值与取等条件]
\label{lem:lihang-em-f-gap}
设 $0<p_\theta(y)<\infty$ 且联合概率在有限隐状态上为正，则
\begin{equation}
  \ell(\theta)-F(r,\theta)
  =D_{\mathrm{KL}}\!\left(r\Vert p_\theta(\,\cdot\mid y)\right)\geq0.
  \label{eq:lihang-em-f-gap}
\end{equation}
因此固定 $\theta$ 时，唯一使 $F$ 达到最大值的分布为 $r_\theta(z)=p_\theta(z\mid y)$，且
\begin{equation}
  F(r_\theta,\theta)=\ell(\theta)=\max_rF(r,\theta).
  \label{eq:lihang-em-f-profile}
\end{equation}
若 $p_\theta(y,z)$ 关于 $\theta$ 连续且分母保持为正，则 $r_\theta$ 关于 $\theta$ 连续。
\end{lemma}

\begin{proof}
代入 $p_\theta(y,z)=p_\theta(z\mid y)p_\theta(y)$，并利用 $\sum_zr(z)=1$，得到
\[
  F(r,\theta)=\ell(\theta)+\sum_zr(z)\log\frac{p_\theta(z\mid y)}{r(z)}.
\]
由 KL 散度的非负性及取等条件即得两式。后验分布每个坐标都是连续联合概率除以有限个连续联合概率之和，分母为正，因此连续。
\end{proof}

引理~\ref{lem:lihang-em-f-gap} 同时给出了教材两个引理的结论。也可用拉格朗日乘子核查最大分布：对 $F(r,\theta)+\lambda(1-\sum_zr(z))$ 关于正的 $r(z)$ 求导，得
\[
  \log p_\theta(y,z)-\log r(z)-1-\lambda=0.
\]
因而 $r(z)$ 与 $p_\theta(y,z)$ 成比例，归一化即得后验。熵项使 $F$ 关于 $r$ 严格凹，保证该解是唯一的最大分布，而不仅是一个满足一阶条件的候选点。

\begin{proposition}[联合极大点与似然极大点]
\label{prop:lihang-em-f-extrema}
在引理~\ref{lem:lihang-em-f-gap} 的有限状态、正联合概率及后验连续条件下，若 $(r^*,\theta^*)$ 是 $F$ 的联合局部极大点，则 $r^*=r_{\theta^*}$，且 $\theta^*$ 是 $\ell$ 的局部极大点。反过来，若 $\theta^*$ 是 $\ell$ 的局部极大点，则 $(r_{\theta^*},\theta^*)$ 是 $F$ 的联合局部极大点。把局部极大换成全局最大，两方向仍成立。
\end{proposition}

\begin{proof}
固定 $\theta^*$，$F$ 关于 $r$ 为严格凹函数。若 $r^*$ 不是其唯一最大点，沿着连接 $r^*$ 与 $r_{\theta^*}$ 的线段就会有任意近的点提高 $F$，与联合局部极大矛盾，所以 $r^*=r_{\theta^*}$。若存在趋向 $\theta^*$ 的参数序列使 $\ell$ 更大，由后验连续性，相应的 $(r_\theta,\theta)$ 趋向 $(r^*,\theta^*)$，但 $F(r_\theta,\theta)=\ell(\theta)>\ell(\theta^*)$，仍矛盾。反方向由 $F(r,\theta)\leq\ell(\theta)\leq\ell(\theta^*)=F(r_{\theta^*},\theta^*)$ 在相应邻域内成立得到。全局情形直接将邻域换成整个可行域。
\end{proof}

EM 的一次迭代恰好可以写为
\begin{equation}
  \begin{aligned}
    r^{(t+1)}&=p_{\theta^{(t)}}(\,\cdot\mid y)
      \in\operatorname*{arg\,max}_rF(r,\theta^{(t)}),\\
    \theta^{(t+1)}&\in\operatorname*{arg\,max}_{\theta\in\Theta}F(r^{(t+1)},\theta)
      =\operatorname*{arg\,max}_{\theta\in\Theta}Q(\theta\mid\theta^{(t)}).
  \end{aligned}
  \label{eq:lihang-em-max-max}
\end{equation}
第二行成立是因为固定 $r^{(t+1)}$ 后，$F=Q+H(r^{(t+1)})$ 中的熵不随 $\theta$ 变化。这就是教材的\emph{极大--极大算法}（maximization--maximization algorithm），亦即其 GEM 算法 1 所写的精确两步。M 步完成后，$r^{(t+1)}$ 仍是旧参数的后验，要到下一次 E 步才与新参数匹配。命题~\ref{prop:lihang-em-f-extrema} 谈的是两个变量同时变化下的联合极大点；仅仅分别对每组变量无法继续提高，并不自动满足该命题的联合局部极大假设。

\subsection{GEM 算法}
\label{subsec:lihang-em-gem}

\emph{广义期望极大算法}（generalized expectation maximization algorithm，GEM）保留 E 步的当前后验，但放宽精确 M 步：只要求新参数在可行域内满足
\begin{equation}
  Q(\theta^{(t+1)}\mid\theta^{(t)})
  \geq Q(\theta^{(t)}\mid\theta^{(t)}).
  \label{eq:lihang-em-gem-condition}
\end{equation}
未停止时通常寻求严格提高。由定理~\ref{thm:lihang-em-monotonicity}，这种更新仍使观测似然不减；但放宽条件也可能容许没有实际进展的更新，所以仅凭该不等式不能证明参数收敛或其极限为驻点。教材的 GEM 算法 2 就是每轮计算 $Q$ 后，寻找满足上述提高条件的参数，继而检查停止准则。

例如在无约束开域中，若本轮 $Q$ 在当前点可微且 $g_t=\nabla_\theta Q(\theta^{(t)}\mid\theta^{(t)})\ne0$，则沿该方向的一阶展开为
\[
  Q(\theta^{(t)}+\eta g_t\mid\theta^{(t)})
  =Q(\theta^{(t)}\mid\theta^{(t)})+\eta\|g_t\|^2+o(\eta).
\]
因而充分小的正步长会提高 $Q$；可通过逐步缩短步长并实际检查目标增量来选取更新。存在参数约束时，还需使整个更新可行，不能直接沿任意梯度方向越出概率单纯形或正方差域。

教材的 GEM 算法 3 将 M 步分成若干次\emph{条件极大化}（conditional maximization）。设参数分为 $d$ 个坐标或参数块，令 $\theta^{(t,0)}=\theta^{(t)}$。在第 $j$ 次内层更新中，只改变第 $j$ 块，求
\begin{equation}
  \theta_j^{(t,j)}\in\operatorname*{arg\,max}_{u:
    (\theta_1^{(t,j-1)},\ldots,u,\ldots,\theta_d^{(t,j-1)})\in\Theta}
  Q(\theta_1^{(t,j-1)},\ldots,u,\ldots,\theta_d^{(t,j-1)}\mid\theta^{(t)}),
  \label{eq:lihang-em-conditional-maximization}
\end{equation}
其余块保持 $\theta_h^{(t,j)}=\theta_h^{(t,j-1)}$，$h\ne j$。依次进行 $j=1,\ldots,d$，再置 $\theta^{(t+1)}=\theta^{(t,d)}$。每个条件可行域都包含更新前的该块值，因此每次内层更新不降低 $Q$，串联即得式~\eqref{eq:lihang-em-gem-condition}。\emph[caution]{整轮条件极大化固定同一个旧后验}；已经更新的参数块用于后续条件优化，但不能把后验权重也无说明地逐块换掉。各条件最大值必须存在，带有耦合约束的坐标若无法单独改变，需选择合适的参数块。

上述放宽直接作用于 M 步。若同时限制辅助分布的形式、令 E 步不能取得真实后验，则通常只能从交替提高 $F$ 得到辅助下界不减；原似然的单调性还需要额外论证。原因是此时 $F$ 与 $\ell$ 之间的 KL 差距不一定在每轮起点为零。这也再次说明，EM 的似然保证依赖当前后验给出的接触条件。

\section*{本章回顾与教材习题}
\addcontentsline{toc}{section}{本章回顾与教材习题}

EM 将观测似然的优化分解为后验计算和完全数据条件期望的优化；Jensen 下界、增量恒等式与 $F$ 函数是同一机制的三种表达。具体模型中的推导应首先说明联合分布与隐变量，其次计算后验权重，最后验证带约束的 M 步是否真正存在并提高目标。单调性、目标值收敛、参数收敛及全局最优是不同层次的结论。教材习题 9.2 要求的 $F(r_\theta,\theta)=\ell(\theta)$ 已在引理~\ref{lem:lihang-em-f-gap} 中证明，以下补充其余三题的计算与构造。

\begin{example}
\label{ex:lihang-em-coin-exercise}
教材习题 9.1。沿用例~\ref{ex:lihang-em-three-coins} 的观测数据，以 $(\pi^{(0)},p^{(0)},q^{(0)})=(0.46,0.55,0.67)$ 为初值求参数估计，并解释其含义。
\end{example}

\begin{solve}
对于正面和反面观测，分别使用
\[
  r_1=\frac{0.46\times0.55}{0.46\times0.55+0.54\times0.67},\qquad
  r_0=\frac{0.46\times0.45}{0.46\times0.45+0.54\times0.33}.
\]
代入式~\eqref{eq:lihang-em-coin-update}，第一轮得到
\[
  (\pi^{(1)},p^{(1)},q^{(1)})
  \approx(0.461863,0.534595,0.656135).
\]
更一般地，只要更新分母为正，令 $\overline y=n^{-1}\sum_i y_i$，同一公式直接给出
\[
  \pi^{(t+1)}p^{(t+1)}
  +(1-\pi^{(t+1)})q^{(t+1)}
  =\frac{\sum_i r_i^{(t)}y_i+\sum_i(1-r_i^{(t)})y_i}{n}
  =\overline y.
\]
本题第一轮已使混合正面概率等于 $0.6$，故到达极大似然解集。由于这类观测的响应度只取 $r_1,r_0$ 两个值，下一轮正面响应度为 $\pi^{(1)}p^{(1)}/\overline y=r_1$，反面响应度同样为 $r_0$，所以参数随后保持不变。不同初值所得参数可以不同，却有相同的最优观测似然。
\end{solve}

\begin{example}
\label{ex:lihang-em-gmm-exercise}
教材习题 9.3。对数据
\[
  -67,-48,6,8,14,16,23,24,28,29,41,49,56,60,75
\]
拟合两个成分的一维高斯混合模型，给出一次具有明确初值和停止条件的数值估计。
\end{example}

\begin{solve}
取初始混合系数 $(0.5,0.5)$、均值 $(-50,30)$、方差 $(400,400)$，按式~\eqref{eq:lihang-em-gmm-responsibility} 和式~\eqref{eq:lihang-em-gmm-m-step} 进行精确形式的 EM 更新。在对数域求响应度及观测似然，设最大迭代次数为 $1000$，连续两轮观测对数似然绝对差小于 $10^{-10}$ 时停止。本次计算在第 $5$ 轮满足该条件，结果约为
\[
  \begin{aligned}
    (\widehat\alpha_1,\widehat\alpha_2)&=(0.133172,0.866828),\\
    (\widehat\mu_1,\widehat\mu_2)&=(-57.511077,32.984887),\\
    (\widehat\sigma_1^2,\widehat\sigma_2^2)&=(90.249878,429.458341),
  \end{aligned}
\]
相应观测对数似然约为 $-71.063362$，计算中各轮目标值不减。这里报告的是该初值和停止规则下的有限数值迭代结果，其全局最优性未经证明；无约束模型的似然退化问题仍然存在。交换两个成分的编号会得到同一个混合密度，比较不同运行时应先对齐编号。
\end{solve}

\begin{example}
\label{ex:lihang-em-naive-bayes}
教材习题 9.4。设样本 $x_1,\ldots,x_n$ 没有类别标记，各 $x_i=(x_{i1},\ldots,x_{im})$ 由离散特征组成，第 $j$ 个特征的有限取值集合为 $\mathcal A_j$。写出给定类别数 $K$ 时，朴素贝叶斯模型无监督参数估计的 EM 算法。
\end{example}

\begin{solve}
令潜在类别为 $C_i\in\{1,\ldots,K\}$，类别先验为 $\pi_k$，特征条件概率为 $\beta_{jka}=P(X_{ij}=a\mid C_i=k)$。模型假设各样本独立同分布，且给定类别后各特征条件独立，于是
\[
  p_\theta(x_i,C_i=k)=\pi_k\prod_{j=1}^m\beta_{jkx_{ij}},\qquad
  \sum_k\pi_k=1,\qquad \sum_{a\in\mathcal A_j}\beta_{jka}=1.
\]
输入观测样本、类别数、各特征取值集合及正的初始概率，E 步计算
\[
  \tau_{ik}^{(t)}
  =\frac{\pi_k^{(t)}\prod_{j=1}^m\beta_{jkx_{ij}}^{(t)}}
  {\sum_{h=1}^K\pi_h^{(t)}\prod_{j=1}^m\beta_{jhx_{ij}}^{(t)}}.
\]
完全数据对数似然中的类别指示量取条件期望后，得到
\[
  Q(\theta\mid\theta^{(t)})
  =\sum_{i=1}^n\sum_{k=1}^K\tau_{ik}^{(t)}
  \left[\log\pi_k+\sum_{j=1}^m\log\beta_{jkx_{ij}}\right].
\]
对每组概率施加归一化约束，M 步即为后验加权频率：
\[
  N_k^{(t)}=\sum_i\tau_{ik}^{(t)},\qquad
  \pi_k^{(t+1)}=\frac{N_k^{(t)}}n,\qquad
  \beta_{jka}^{(t+1)}
  =\frac{\sum_i\tau_{ik}^{(t)}\mathbf1_{\{x_{ij}=a\}}}{N_k^{(t)}}.
\]
在 $N_k^{(t)}>0$ 时，这些式子分别是相应多项分布似然的最大点；若有效样本数为零，该类别的条件参数不由本轮目标确定。交替 E、M 两步，按观测对数似然容差或最大轮数停止，输出类别先验与各特征条件概率。应避免所有类别使用完全相同的初始参数，否则对称性可能一直保持；没有真实类别标记时，潜在类别的编号不能直接解释为预先指定的语义类别。
\end{solve}

% 教材来源：机器学习/textbooks/机器学习（分章原文）/机器学习-10.pdf。
% 对应教材第 10 章，原书第 193--214 页；本文件采用自动章节编号。

\fi

\selectchapter{10}
\ifcompilechapter
\chapter{隐马尔可夫模型}
\label{chap:lh-hmm}

序列标注需要把观测序列与不可直接观测的状态序列联系起来。例如，语音信号可以被观测，而相应的音素状态需要推断。\emph{隐马尔可夫模型}（hidden Markov model，HMM）用一条隐藏的马尔可夫链描述状态的演化，再由当前状态生成观测，因此直接建模的是状态与观测的联合分布。本章依据教材第 10 章展开：先说明模型的生成假设，再用前向与后向递推计算概率和状态后验，继而推导参数学习方法，最后用维特比算法求解整条最可能的状态路径。所需基础是条件概率、全概率公式、极大似然估计和动态规划；Baum--Welch 算法所用的期望最大化思想在推导中一并说明。

\section{隐马尔可夫模型的基本概念}
\label{sec:lh-hmm-model}

\subsection{隐马尔可夫模型的定义}

设状态集合为 $\mathcal S=\{1,\ldots,N\}$，观测符号集合为 $\mathcal V=\{v_1,\ldots,v_M\}$，其中 $N,M$ 为正整数。用大写 $Z_t$ 表示时刻 $t$ 的隐藏状态随机变量，$X_t$ 表示观测随机变量；小写 $z_t,x_t$ 表示相应取值。长度为 $T\geq1$ 的随机序列记为 $Z_{1:T}$ 和 $X_{1:T}$，给定的观测序列记为 $x_{1:T}=(x_1,\ldots,x_T)$。教材使用 $I,O$ 表示状态序列和观测序列，这里改用上述记号以区分随机变量与取值。

\begin{definition}
\label{def:lh-hmm}
有限状态、离散观测的齐次隐马尔可夫模型由参数 $\lambda=(A,B,\pi)$ 确定，其中
\begin{equation}
\label{eq:lh-hmm-parameters}
 a_{ij}=P_\lambda(Z_{t+1}=j\mid Z_t=i),\qquad
 b_i(k)=P_\lambda(X_t=v_k\mid Z_t=i),\qquad
 \pi_i=P_\lambda(Z_1=i).
\end{equation}
这里 $A=(a_{ij})\in\mathbb R^{N\times N}$ 为状态转移概率矩阵（transition probability matrix），$B=(b_i(k))\in\mathbb R^{N\times M}$ 为观测概率矩阵，亦称发射概率矩阵（emission probability matrix），$\pi=(\pi_1,\ldots,\pi_N)^{\mathsf T}$ 为初始状态概率向量。各元素非负，并满足
\[
 \sum_{j=1}^N a_{ij}=1,\qquad
 \sum_{k=1}^M b_i(k)=1,\qquad
 \sum_{i=1}^N\pi_i=1.
\]
模型假定 $Z_t$ 构成以 $A$ 为转移矩阵的齐次一阶马尔可夫链；给定整个状态序列后，各时刻观测条件独立，且 $X_t$ 的分布仅由 $Z_t$ 决定。
\end{definition}

上述假设的可操作表述是：对具有正概率的条件事件，
\begin{align}
\label{eq:lh-hmm-markov}
 P_\lambda(Z_{t+1}=j\mid Z_{1:t}=z_{1:t},X_{1:t}=x_{1:t})&=a_{z_tj},
 &&1\leq t<T,\\
\label{eq:lh-hmm-emission}
 P_\lambda(X_{1:T}=x_{1:T}\mid Z_{1:T}=z_{1:T})&=\prod_{t=1}^T b_{z_t}(x_t).
\end{align}
当 $x=v_k$ 时，约定 $b_i(x)=b_i(k)$，避免把观测符号本身误当作矩阵列索引。齐次性要求 $A,B$ 不随 $t$ 改变；观测独立性是\emph[caution]{给定隐藏状态后的条件独立}，并不意味着 $X_1,\ldots,X_T$ 在边缘分布下相互独立，也不意味着观测序列自身必是一阶马尔可夫链。

由概率链式法则和模型假设，完整路径与观测的联合概率为
\begin{equation}
\label{eq:lh-hmm-joint}
 p_\lambda(z_{1:T},x_{1:T})
 =\pi_{z_1}b_{z_1}(x_1)\prod_{t=1}^{T-1}a_{z_tz_{t+1}}b_{z_{t+1}}(x_{t+1}).
\end{equation}
其中 $p_\lambda$ 表示相应的概率质量函数，$T=1$ 时空乘积取 $1$。这个分解既给出生成过程，也为后续所有递推和学习公式提供依据。状态数、观测符号集合和序列长度在本章中视为已知；$T$ 固定，不额外设置终止状态或终止概率。

\begin{example}
\label{ex:lh-hmm-four-boxes}
教材的四盒抽球模型中，四个盒子内红、白球数分别为 $(5,5)$、$(3,7)$、$(6,4)$、$(8,2)$。初始等概率选择盒子；此后盒子 $1$ 必转到盒子 $2$，盒子 $2,3$ 分别以 $0.4,0.6$ 的概率转到左、右相邻盒子，盒子 $4$ 以各 $0.5$ 的概率转到盒子 $3$ 或留在原处。每次从当前盒子均匀抽一球、记录颜色并放回，只能观察球色。写出该模型。
\end{example}
\begin{solve}
令 $Z_t$ 为第 $t$ 次所选盒子的编号，$X_t$ 为球色，按红、白排列观测符号，则
\[
 A=\begin{pmatrix}
 0&1&0&0\\
 0.4&0&0.6&0\\
 0&0.4&0&0.6\\
 0&0&0.5&0.5
 \end{pmatrix},\qquad
 B=\begin{pmatrix}
 0.5&0.5\\0.3&0.7\\0.6&0.4\\0.8&0.2
 \end{pmatrix},\qquad
 \pi=\begin{pmatrix}0.25\\0.25\\0.25\\0.25\end{pmatrix}.
\]
每行概率之和均为 $1$。放回抽样保证每个盒子的球色分布不因之前的抽球结果改变，在每次抽球使用独立随机性的条件下符合观测假设。长度为 $5$ 的一次观测可以是 $(\text{红},\text{红},\text{白},\text{白},\text{红})$，但这组颜色通常不能唯一确定盒子的路径。
\end{solve}

\subsection{观测序列的生成过程}

输入合法参数 $\lambda$ 与长度 $T$，先按 $\pi$ 抽取 $Z_1$。对 $t=1,\ldots,T$，在已得到 $Z_t=i$ 时按 $B$ 的第 $i$ 行生成 $X_t$；若 $t<T$，再按 $A$ 的第 $i$ 行抽取 $Z_{t+1}$，否则终止。输出观测序列 $X_{1:T}$，需要模拟完整数据时同时保留 $Z_{1:T}$。这个顺序先生成当前观测再决定是否转移，因而恰好生成 $T$ 个观测和 $T-1$ 次转移，$T=1$ 时也成立。

\subsection{隐马尔可夫模型的三个基本问题}

\emph{概率计算}给定 $\lambda,x_{1:T}$，求观测似然 $p_\lambda(x_{1:T})$；\emph{学习}由训练序列估计 $\lambda$，通常最大化观测数据的对数似然；\emph{解码}（decoding）给定 $\lambda,x_{1:T}$，求最可能的状态路径
\begin{equation}
\label{eq:lh-hmm-decoding-objective}
 z_{1:T}^*\in\operatorname*{arg\,max}_{z_{1:T}\in\mathcal S^T}
 P_\lambda(Z_{1:T}=z_{1:T}\mid X_{1:T}=x_{1:T}).
\end{equation}
后验概率要求 $p_\lambda(x_{1:T})>0$。此时其分母与候选路径无关，解码等价于最大化式~\eqref{eq:lh-hmm-joint} 的联合概率。教材称解码为“预测”，这里特指对\emph{已观测序列的隐藏状态}作推断；预测未来观测或未来状态是另一个条件概率问题。

\section{概率计算算法}
\label{sec:lh-hmm-inference}

\subsection{直接计算法}

对所有隐藏路径边缘化可得
\begin{equation}
\label{eq:lh-hmm-likelihood}
 p_\lambda(x_{1:T})
 =\sum_{z_1=1}^N\cdots\sum_{z_T=1}^N
 \pi_{z_1}b_{z_1}(x_1)\prod_{t=1}^{T-1}a_{z_tz_{t+1}}b_{z_{t+1}}(x_{t+1}).
\end{equation}
逐条计算 $N^T$ 条路径、每条使用 $O(T)$ 次运算，需要 $O(TN^T)$ 时间。困难在于反复计算共享的路径前缀。\emph[tip]{把具有同一末状态的全部前缀合并}后，未来只需知道这个末状态及前缀的总概率，而无须保存每条前缀，这就得到前向算法。

\subsection{前向算法}

\begin{definition}
\label{def:lh-hmm-forward}
给定 $\lambda,x_{1:T}$，前向概率（forward probability）定义为
\begin{equation}
\label{eq:lh-hmm-alpha-def}
 \alpha_t(i)=P_\lambda(X_{1:t}=x_{1:t},Z_t=i),\qquad 1\leq t\leq T.
\end{equation}
它是已观测前缀与当前状态的联合概率，通常不满足 $\sum_i\alpha_t(i)=1$。
\end{definition}

\begin{proposition}
\label{prop:lh-hmm-forward}
前向概率满足
\begin{align}
\label{eq:lh-hmm-alpha-init}
 \alpha_1(i)&=\pi_i b_i(x_1),\\
\label{eq:lh-hmm-alpha-rec}
 \alpha_{t+1}(j)&=b_j(x_{t+1})\sum_{i=1}^N\alpha_t(i)a_{ij},
 &&1\leq t<T,\\
\label{eq:lh-hmm-alpha-end}
 p_\lambda(x_{1:T})&=\sum_{i=1}^N\alpha_T(i).
\end{align}
\end{proposition}
\begin{proof}
初值直接来自联合概率分解。对递推式，按 $Z_t$ 的可能取值划分事件，得到
\[
 \begin{aligned}
 \alpha_{t+1}(j)
 &=\sum_{i=1}^N P_\lambda(X_{1:t}=x_{1:t},Z_t=i,Z_{t+1}=j,X_{t+1}=x_{t+1})\\
 &=\sum_{i=1}^N\alpha_t(i)a_{ij}b_j(x_{t+1}).
 \end{aligned}
\]
第二步使用马尔可夫性和观测条件独立性。最后对终点状态求和，即把所有路径分成 $N$ 类，得到观测似然。
\end{proof}

前向算法以 $\lambda,x_{1:T}$ 为输入，按式~\eqref{eq:lh-hmm-alpha-init} 初始化，依次完成 $T-1$ 次递推后按式~\eqref{eq:lh-hmm-alpha-end} 输出似然。稠密 $A$ 下每步计算 $N$ 个各含 $N$ 项的和，总时间为 $O(TN^2)$；只计算似然时滚动保存相邻时刻的向量即可，额外空间为 $O(N)$，若保留全部前向概率则为 $O(TN)$。$\alpha_t(i)$ 的下标 $t$ 同时标明观测前缀的终点和当前状态所在时刻。

\begin{example}
\label{ex:lh-hmm-three-boxes}
教材的三盒模型采用
\begin{equation}
\label{eq:lh-hmm-three-boxes}
 A=\begin{pmatrix}0.5&0.2&0.3\\0.3&0.5&0.2\\0.2&0.3&0.5\end{pmatrix},\qquad
 B=\begin{pmatrix}0.5&0.5\\0.4&0.6\\0.7&0.3\end{pmatrix},\qquad
 \pi=\begin{pmatrix}0.2\\0.4\\0.4\end{pmatrix},
\end{equation}
其中 $B$ 的两列依次对应红、白。求 $x_{1:3}=(\text{红},\text{白},\text{红})$ 的概率。
\end{example}
\begin{solve}
把前向量按行写出，依次得到
\[
 \begin{aligned}
 \alpha_1&=(0.1,\,0.16,\,0.28),\\
 \alpha_2&=(0.077,\,0.1104,\,0.0606),\\
 \alpha_3&=(0.04187,\,0.035512,\,0.052836).
 \end{aligned}
\]
例如 $\alpha_2(1)=(0.1\times0.5+0.16\times0.3+0.28\times0.2)\times0.5=0.077$。因此
\[
 p_\lambda(x_{1:3})=0.04187+0.035512+0.052836=0.130218\approx0.13022.
\]
这里取教材参数的小数为精确值计算，最后再舍入；所得似然汇总了全部 $3^3$ 条路径，不能解释为某一条路径的概率。
\end{solve}

\subsection{后向算法}

\begin{definition}
\label{def:lh-hmm-backward}
后向概率（backward probability）为
\begin{equation}
\label{eq:lh-hmm-beta-def}
 \beta_t(i)=P_\lambda(X_{t+1:T}=x_{t+1:T}\mid Z_t=i),\qquad1\leq t\leq T.
\end{equation}
对不可达状态也可按模型指定的转移与发射机制定义右侧的后续生成分布。$\beta_t(i)$ 不含 $x_t$；$t=T$ 时后缀为空，故 $\beta_T(i)=1$。
\end{definition}

按下一个状态分解后续事件，得到后向递推与终止公式
\begin{align}
\label{eq:lh-hmm-beta-rec}
 \beta_t(i)&=\sum_{j=1}^N a_{ij}b_j(x_{t+1})\beta_{t+1}(j),
 &&t=T-1,\ldots,1,\\
\label{eq:lh-hmm-beta-end}
 p_\lambda(x_{1:T})&=\sum_{i=1}^N\pi_i b_i(x_1)\beta_1(i).
\end{align}
每项依次对应“转到 $j$、生成下一观测、生成剩余后缀”。因此，算法输入 $\lambda,x_{1:T}$，初始化 $\beta_T$ 为全 $1$ 向量，逆序递推到 $1$ 后输出式~\eqref{eq:lh-hmm-beta-end}。稠密情形时间复杂度同为 $O(TN^2)$；只求似然时额外空间为 $O(N)$。给定 $Z_t$ 后过去与未来分离，故
\begin{equation}
\label{eq:lh-hmm-cut}
 P_\lambda(X_{1:T}=x_{1:T},Z_t=i)=\alpha_t(i)\beta_t(i),\qquad
 p_\lambda(x_{1:T})=\sum_{i=1}^N\alpha_t(i)\beta_t(i).
\end{equation}
该恒等式在任意 $t=1,\ldots,T$ 都成立，可用于核对两种递推的一致性。在例~\ref{ex:lh-hmm-three-boxes} 中，$\beta_3=(1,1,1)$、$\beta_2=(0.54,0.49,0.57)$、$\beta_1=(0.2451,0.2622,0.2277)$，因而 $0.1\times0.2451+0.16\times0.2622+0.28\times0.2277=0.130218$。

\subsection{状态后验概率与期望计数}

以下假设观测似然 $L=p_\lambda(x_{1:T})>0$，所有后验都针对同一整段观测。定义单状态后验与相邻状态后验为
\begin{align}
\label{eq:lh-hmm-gamma}
 \gamma_t(i)&=P_\lambda(Z_t=i\mid X_{1:T}=x_{1:T})
 =\frac{\alpha_t(i)\beta_t(i)}{L},\\
\label{eq:lh-hmm-xi}
 \xi_t(i,j)&=P_\lambda(Z_t=i,Z_{t+1}=j\mid X_{1:T}=x_{1:T})
 =\frac{\alpha_t(i)a_{ij}b_j(x_{t+1})\beta_{t+1}(j)}{L},
 \quad t<T.
\end{align}
前一式由式~\eqref{eq:lh-hmm-cut} 和贝叶斯公式得到；后一式的分子把路径在相邻两个状态之间切开，依次相乘前缀概率、转移概率、下一观测概率和剩余后缀概率。它们满足
\begin{equation}
\label{eq:lh-hmm-posterior-consistency}
 \sum_i\gamma_t(i)=1,\qquad
 \sum_{i,j}\xi_t(i,j)=1,\qquad
 \sum_j\xi_t(i,j)=\gamma_t(i),\qquad
 \sum_i\xi_t(i,j)=\gamma_{t+1}(j).
\end{equation}
这些等式来自同一后验联合分布的边缘化，也说明 $\gamma_t$ 与 $\xi_t$ 不能任意独立指定。

为把后验推断与参数学习连接起来，记 $\mathbf1\{E\}$ 为条件 $E$ 成立时取 $1$、否则取 $0$ 的示性函数。由条件期望的线性性，
\begin{align}
\label{eq:lh-hmm-counts}
 \mathbb E_\lambda\!\left[\sum_{t=1}^T\mathbf1\{Z_t=i\}\,\middle|\,X_{1:T}=x_{1:T}\right]
 &=\sum_{t=1}^T\gamma_t(i),\\
\label{eq:lh-hmm-transition-counts}
 \mathbb E_\lambda\!\left[\sum_{t=1}^{T-1}\mathbf1\{Z_t=i,Z_{t+1}=j\}\,\middle|\,X_{1:T}=x_{1:T}\right]
 &=\sum_{t=1}^{T-1}\xi_t(i,j).
\end{align}
状态 $i$ 作为转移出发点的期望次数为 $\sum_{t=1}^{T-1}\gamma_t(i)$，因为最后一个时刻没有后续转移；状态 $i$ 发射 $v_k$ 的期望次数为 $\sum_{t=1}^T\gamma_t(i)\mathbf1\{x_t=v_k\}$。这些量是\emph{给定观测后的计数条件期望}，不是实际已经观测到的状态频数。另需区分滤波（filtering）概率 $P_\lambda(Z_t=i\mid X_{1:t}=x_{1:t})=\alpha_t(i)/\sum_j\alpha_t(j)$ 与平滑（smoothing）概率 $\gamma_t(i)$：后者还利用 $t$ 之后的观测。

\begin{notealert}[type=tip,title={长序列的缩放计算}]
长序列的概率连乘容易在浮点运算中下溢。令 $\widetilde\alpha_1(i)=\pi_i b_i(x_1)$，$c_1=\sum_i\widetilde\alpha_1(i)$，$\widehat\alpha_1=\widetilde\alpha_1/c_1$；此后递推
\[
 \widetilde\alpha_{t+1}(j)=b_j(x_{t+1})\sum_i\widehat\alpha_t(i)a_{ij},\qquad
 c_{t+1}=\sum_j\widetilde\alpha_{t+1}(j),\qquad
 \widehat\alpha_{t+1}(j)=\frac{\widetilde\alpha_{t+1}(j)}{c_{t+1}}.
\]
若各 $c_t>0$，则 $c_t=P_\lambda(X_t=x_t\mid X_{1:t-1}=x_{1:t-1})$（$t=1$ 时无条件前缀），从而 $\log L=\sum_t\log c_t$。配套的后向量须采用相同缩放：$\widehat\beta_T(i)=1$，
\[
 \widehat\beta_t(i)=\frac{\sum_j a_{ij}b_j(x_{t+1})\widehat\beta_{t+1}(j)}{c_{t+1}},\qquad
 \gamma_t(i)=\widehat\alpha_t(i)\widehat\beta_t(i).
\]
这可由 $\widehat\alpha_t=\alpha_t/\prod_{s\leq t}c_s$ 和 $\widehat\beta_t=\beta_t/\prod_{s>t}c_s$ 验证，不能把已缩放的前向量与未缩放的后向量直接代入原公式。精确运算中若某个 $c_t=0$，则 $L=0$，该参数下观测不可能发生，后验无定义。
\end{notealert}

\section{学习算法}
\label{sec:lh-hmm-learning}

\subsection{监督学习方法}

若训练数据包括相互独立的完整序列 $\{(x^{(s)}_{1:T_s},z^{(s)}_{1:T_s})\}_{s=1}^S$，且各序列共享同一参数，就可以直接统计初始状态、转移和发射次数。记
\[
 \begin{aligned}
 C_i^{\mathrm{init}}&=\sum_{s=1}^S\mathbf1\{z_1^{(s)}=i\},\\
 C_{ij}^{\mathrm{tr}}&=\sum_{s=1}^S\sum_{t=1}^{T_s-1}
 \mathbf1\{z_t^{(s)}=i,z_{t+1}^{(s)}=j\},\\
 C_{ik}^{\mathrm{em}}&=\sum_{s=1}^S\sum_{t=1}^{T_s}
 \mathbf1\{z_t^{(s)}=i,x_t^{(s)}=v_k\}.
 \end{aligned}
\]
由式~\eqref{eq:lh-hmm-joint}，忽略与参数无关的项，完整数据对数似然为
\begin{equation}
\label{eq:lh-hmm-supervised-loglik}
 \ell_{\mathrm{comp}}(\lambda)
 =\sum_i C_i^{\mathrm{init}}\log\pi_i
 +\sum_{i,j}C_{ij}^{\mathrm{tr}}\log a_{ij}
 +\sum_{i,k}C_{ik}^{\mathrm{em}}\log b_i(k).
\end{equation}
每行参数分别受概率归一化约束，因而极大似然估计为
\begin{equation}
\label{eq:lh-hmm-supervised-mle}
 \widehat\pi_i=\frac{C_i^{\mathrm{init}}}{S},\qquad
 \widehat a_{ij}=\frac{C_{ij}^{\mathrm{tr}}}{\sum_{r=1}^N C_{ir}^{\mathrm{tr}}},\qquad
 \widehat b_i(k)=\frac{C_{ik}^{\mathrm{em}}}{\sum_{r=1}^M C_{ir}^{\mathrm{em}}}.
\end{equation}
后两式要求相应分母为正。这个结论可归结为：给定非负计数 $c_r$、$C=\sum_r c_r>0$，在 $p_r\geq0,\sum_rp_r=1$ 上最大化 $\sum_rc_r\log p_r$，解为 $p_r=c_r/C$。正计数部分的拉格朗日条件为 $c_r/p_r=\eta$，求和得 $\eta=C$；零计数部分在最优解中分配零概率即可。约定 $0\log0=0$，正计数乘以 $\log0$ 则表示目标为 $-\infty$。若整行计数均为零，似然不约束该行，任何合法概率行都达到相同值，不能用 $0/0$ 定义估计。

\subsection{Baum--Welch 算法的 EM 推导}

只有观测序列而没有状态标注时，极大化目标为
\begin{equation}
\label{eq:lh-hmm-observed-loglik}
 \ell(\lambda)=\sum_{s=1}^S\log p_\lambda(x^{(s)}_{1:T_s}).
\end{equation}
每个对数内都包含隐藏路径求和，无法直接用已知频数作归一化。\emph{Baum--Welch 算法}（Baum--Welch algorithm）是期望最大化（expectation-maximization，EM）算法在 HMM 中的实现：先按当前模型计算隐藏路径的后验期望计数，再用这些计数重新估计参数。以下先推导单条序列情形；本小节用 $Q$ 表示 EM 的辅助函数，与状态集合 $\mathcal S$ 区别。

固定当前参数 $\bar\lambda$，要求 $p_{\bar\lambda}(x_{1:T})>0$，令 $q(z_{1:T})=P_{\bar\lambda}(Z_{1:T}=z_{1:T}\mid X_{1:T}=x_{1:T})$。\emph{E 步}定义
\begin{equation}
\label{eq:lh-hmm-Q-def}
 Q(\lambda\mid\bar\lambda)
 =\sum_{z_{1:T}}q(z_{1:T})\log p_\lambda(z_{1:T},x_{1:T}).
\end{equation}
这里使用归一化的后验概率。教材推导中去掉的因子 $1/p_{\bar\lambda}(x_{1:T})$ 对单条固定观测、固定 $\bar\lambda$ 而言是正常数，故不改变 M 步的极大值点；保留它更便于将 $Q$ 解释为条件期望，并推广到多条独立序列。

把式~\eqref{eq:lh-hmm-joint} 取对数后代入式~\eqref{eq:lh-hmm-Q-def}，各求和均有限，可以交换顺序。用当前参数计算的后验记为 $\bar\gamma_t(i),\bar\xi_t(i,j)$，得到
\begin{equation}
\label{eq:lh-hmm-Q-expanded}
 \begin{aligned}
 Q(\lambda\mid\bar\lambda)
 &=\sum_i\bar\gamma_1(i)\log\pi_i
 +\sum_{i,j}\left(\sum_{t=1}^{T-1}\bar\xi_t(i,j)\right)\log a_{ij}\\
 &\quad+\sum_{i,k}\left(\sum_{t=1}^T\bar\gamma_t(i)\mathbf1\{x_t=v_k\}\right)\log b_i(k).
 \end{aligned}
\end{equation}
例如转移项的系数来自 $\mathbb E_{\bar\lambda}[\mathbf1\{Z_t=i,Z_{t+1}=j\}\mid x_{1:T}]=\bar\xi_t(i,j)$。关键是后验全部由 $\bar\lambda$ 计算，在对候选参数 $\lambda$ 极大化时保持不变；不能把候选参数代回后验并在 M 步中一并求导。

\emph{M 步}分别在概率单纯形上极大化三组对数项。以 $A$ 的第 $i$ 行为例，令 $d_{ij}=\sum_{t<T}\bar\xi_t(i,j)$、$d_i=\sum_jd_{ij}=\sum_{t<T}\bar\gamma_t(i)>0$，拉格朗日函数为
\[
 \mathcal L_i=\sum_jd_{ij}\log a_{ij}+\eta_i\left(1-\sum_j a_{ij}\right).
\]
正计数坐标上 $d_{ij}/a_{ij}=\eta_i$，归一化得 $\eta_i=d_i$，因而 $a_{ij}^{\mathrm{new}}=d_{ij}/d_i$。初始分布和发射分布使用相同论证，所以
\begin{align}
\label{eq:lh-hmm-bw-a}
 a_{ij}^{\mathrm{new}}&=\frac{\sum_{t=1}^{T-1}\bar\xi_t(i,j)}{\sum_{t=1}^{T-1}\bar\gamma_t(i)},\\
\label{eq:lh-hmm-bw-b}
 b_i^{\mathrm{new}}(k)&=\frac{\sum_{t=1}^T\bar\gamma_t(i)\mathbf1\{x_t=v_k\}}{\sum_{t=1}^T\bar\gamma_t(i)},\\
\label{eq:lh-hmm-bw-pi}
 \pi_i^{\mathrm{new}}&=\bar\gamma_1(i).
\end{align}
零计数和整行零计数的处理与监督学习相同；若分母为零，可保留该行的旧值作为一个合法的 M 步选择。特别地，$T=1$ 时没有任何转移观测，$A$ 无法从该序列估计。式~\eqref{eq:lh-hmm-bw-a} 的分母只到 $T-1$，而式~\eqref{eq:lh-hmm-bw-b} 的分母到 $T$，因为每条长度为 $T$ 的路径有 $T-1$ 次转移、$T$ 次发射。

\begin{proposition}
\label{prop:lh-hmm-em-monotone}
固定状态数和模型结构，若当前参数下观测似然为正，E 步精确计算后验，M 步选择使 $Q$ 不减的合法参数，则更新后的观测对数似然不减。该结论不保证得到全局极大值，也不单独保证参数序列收敛。
\end{proposition}
\begin{proof}
只需证明单序列情形，多序列结果可逐项相加。对候选参数令
\[
 F(\lambda,q)=\sum_z q(z)\log p_\lambda(z,x)-\sum_zq(z)\log q(z),
\]
其中 $z,x$ 简写完整序列。若候选参数对 $q(z)>0$ 的路径均给出正联合概率，记其后验为 $q_\lambda(z)=p_\lambda(z,x)/p_\lambda(x)$，则
\[
 \log p_\lambda(x)-F(\lambda,q)
 =\sum_zq(z)\log\frac{q(z)}{q_\lambda(z)}\geq0.
\]
右侧为非负的 Kullback--Leibler 散度；其非负性由 $\log u\leq u-1$ 得到。若候选参数在某个 $q(z)>0$ 的路径上给零概率，则 $F=-\infty$，不可能使有限的当前 $Q$ 不减。当前参数 $\bar\lambda$ 下 $q=q_{\bar\lambda}$，上式取等号；而固定 $q$ 时 $F$ 与 $Q$ 仅差不依赖 $\lambda$ 的后验熵。因此
\[
 \log p_{\lambda^{\mathrm{new}}}(x)
 \geq F(\lambda^{\mathrm{new}},q)
 \geq F(\bar\lambda,q)
 =\log p_{\bar\lambda}(x).\qedhere
\]
\end{proof}

\subsection{参数重估、训练流程与多序列边界}

实际训练输入观测序列、状态数 $N$、观测字典、合法初值 $\lambda^{(0)}$、容差 $\varepsilon>0$ 和最大迭代次数 $K$，输出最后一次合法参数与相应对数似然。对 $n=0,1,\ldots,K-1$，在 $\lambda^{(n)}$ 下用前向与后向算法计算后验，再按式~\eqref{eq:lh-hmm-bw-a}、式~\eqref{eq:lh-hmm-bw-b} 和式~\eqref{eq:lh-hmm-bw-pi} 同时更新全部参数。计算新的观测对数似然后，若增量满足 $0\leq\ell(\lambda^{(n+1)})-\ell(\lambda^{(n)})\leq\varepsilon(1+|\ell(\lambda^{(n)})|)$，可停止；达到 $K$ 也停止。数值实现中若出现超出舍入误差的似然下降，应先检查缩放、后验归一化和更新公式；若似然为零，则不能继续按这些后验公式作 E 步。稠密单序列的每轮主要计算量为 $O(TN^2)$，在同时遍历 $M$ 个观测类别初始化或输出参数时另计 $O(NM)$。

对 $S$ 条独立序列，分别计算各自的 $\bar\gamma_t^{(s)},\bar\xi_t^{(s)}$ 后汇总期望计数，更新公式为
\begin{equation}
\label{eq:lh-hmm-bw-multiple}
 \begin{aligned}
 a_{ij}^{\mathrm{new}}
 &=\frac{\sum_{s=1}^S\sum_{t=1}^{T_s-1}\bar\xi_t^{(s)}(i,j)}
 {\sum_{s=1}^S\sum_{t=1}^{T_s-1}\bar\gamma_t^{(s)}(i)},\\
 b_i^{\mathrm{new}}(k)
 &=\frac{\sum_{s=1}^S\sum_{t=1}^{T_s}\bar\gamma_t^{(s)}(i)\mathbf1\{x_t^{(s)}=v_k\}}
 {\sum_{s=1}^S\sum_{t=1}^{T_s}\bar\gamma_t^{(s)}(i)},\\
 \pi_i^{\mathrm{new}}&=\frac1S\sum_{s=1}^S\bar\gamma_1^{(s)}(i).
 \end{aligned}
\end{equation}
每条序列分别初始化和终止。\emph[caution]{不能直接拼接独立序列}后套用单序列更新，否则会在前一条末尾与后一条开头之间增加不存在的转移，并漏掉各序列的初始状态计数。上述公式是汇总期望计数后归一化，通常也不同于先对每条序列估计参数再作简单平均。

观测对数似然一般不是凹函数，初始化会影响结果，可用多个合法初值比较最终似然。若允许的转移或发射一开始设成零，则相应路径的后验计数也可能始终为零；结构性禁止项应保持零，而没有结构依据的零初值可能限制后续搜索。添加伪计数可用于平滑，但会改变纯极大似然的更新问题，应明确其先验或正则化含义。无状态标注时，对全部状态同时重命名不改变观测分布，因此即便取得相同最优似然，也不能仅据此断言参数具有唯一的语义解释。

\section{预测与解码算法}
\label{sec:lh-hmm-decoding}

\subsection{逐时刻后验最大化的近似算法}

当 $L>0$ 时，最直接的方法是在各时刻分别选择
\begin{equation}
\label{eq:lh-hmm-marginal-decode}
 \widetilde z_t\in\operatorname*{arg\,max}_{1\leq i\leq N}\gamma_t(i),
 \qquad t=1,\ldots,T.
\end{equation}
这称为逐位置的后验解码（posterior decoding）。它确实最小化不限制输出路径时的条件期望逐位置错误数：对任意候选 $\widehat z_{1:T}$，
\[
 \mathbb E_\lambda\!\left[\sum_t\mathbf1\{\widehat z_t\ne Z_t\}\,\middle|\,x_{1:T}\right]
 =\sum_t[1-\gamma_t(\widehat z_t)],
\]
故各项可独立最小化。但式~\eqref{eq:lh-hmm-decoding-objective} 最大化的是\emph{整条路径同时正确的概率}，两个优化目标不同；逐位置解甚至可能包含零概率转移。

例如只有一个观测符号、所有状态均以概率 $1$ 发射该符号，取
\[
 \pi=(0.5,0.26,0.24)^{\mathsf T},\qquad
 A=\begin{pmatrix}0.6&0&0.4\\0&1&0\\0&1&0\end{pmatrix},\qquad T=2.
\]
观测不提供额外信息，故两时刻后验分别为 $(0.5,0.26,0.24)$ 和 $(0.3,0.5,0.2)$；逐位置最大化给出 $(1,2)$，但 $a_{12}=0$。这表明各时刻的边缘概率不能代替路径的联合概率。即使所有转移均为正，两种目标的最优解仍可能不同。

\subsection{维特比算法}

\emph{维特比算法}（Viterbi algorithm）用动态规划计算概率最大的整条路径。定义
\begin{equation}
\label{eq:lh-hmm-delta-def}
 \delta_t(j)=\max_{z_1,\ldots,z_{t-1}}
 p_\lambda(z_1,\ldots,z_{t-1},j,x_{1:t}).
\end{equation}
它保存以状态 $j$ 结尾的最佳前缀的联合概率；与 $\alpha_t(j)$ 对所有前缀求和不同，$\delta_t(j)$ 只取最大值。对两个均以 $j$ 结尾的前缀，任何固定未来路径带来的后续乘积相同且非负，故保留较大概率的前缀不会丢失全局最优解。这就是本模型的最优子结构，依赖于未来只通过当前状态与过去相联系。

算法输入 $\lambda,x_{1:T}$，输出最优状态序列和其联合概率，采用以下递推与回溯：
\begin{align}
\label{eq:lh-hmm-viterbi-init}
 \delta_1(j)&=\pi_jb_j(x_1),\qquad \psi_1(j)=0,\\
\label{eq:lh-hmm-viterbi-rec}
 \delta_t(j)&=b_j(x_t)\max_{1\leq i\leq N}\{\delta_{t-1}(i)a_{ij}\},
 &&2\leq t\leq T,\\
\label{eq:lh-hmm-viterbi-pointer}
 \psi_t(j)&\in\operatorname*{arg\,max}_{1\leq i\leq N}\{\delta_{t-1}(i)a_{ij}\},
 &&2\leq t\leq T,\\
\label{eq:lh-hmm-viterbi-end}
 p^*&=\max_j\delta_T(j),\qquad
 z_T^*\in\operatorname*{arg\,max}_{j}\delta_T(j),\\
\label{eq:lh-hmm-viterbi-backtrack}
 z_t^*&=\psi_{t+1}(z_{t+1}^*),
 &&t=T-1,\ldots,1.
\end{align}
$\psi_t(j)$ 记录最佳前缀在前一时刻的状态；出现并列时按固定规则选取一个即可，所得路径仍为最优路径。$T=1$ 时无需回溯。若 $p^*=0$，则所有路径的联合概率均为零，从而 $L=0$，不能解释为有效的后验解码。若 $L>0$，应注意
\[
 p^*=p_\lambda(z_{1:T}^*,x_{1:T}),\qquad
 P_\lambda(Z_{1:T}=z_{1:T}^*\mid x_{1:T})=\frac{p^*}{L};
\]
$\max_j\delta_T(j)$ 本身不是路径的后验概率。

稠密模型下时间复杂度为 $O(TN^2)$，滚动得分需要 $O(N)$ 空间，保存回溯指针需要 $O(TN)$ 空间。长序列可在对数域计算，约定 $\log0=-\infty$，则
\[
 d_1(j)=\log\pi_j+\log b_j(x_1),\qquad
 d_t(j)=\log b_j(x_t)+\max_i\{d_{t-1}(i)+\log a_{ij}\}.
\]
这里 $d_t(j)=\log\delta_t(j)$；指针保存最大值的来源，对数域运算保持最优路径不变，并避免直接连乘导致的下溢。前向算法在对数域对应的是对指数项求和再取对数，不能也把求和改成最大值，否则所解问题已经改变。

\begin{example}
\label{ex:lh-hmm-viterbi}
对例~\ref{ex:lh-hmm-three-boxes} 的三盒模型及观测 $(\text{红},\text{白},\text{红})$，用维特比算法求最优路径。
\end{example}
\begin{solve}
初始化为 $\delta_1=(0.1,0.16,0.28)$，递推得到
\[
 \begin{aligned}
 \delta_2&=(0.028,\,0.0504,\,0.042),&\qquad\psi_2&=(3,3,3),\\
 \delta_3&=(0.00756,\,0.01008,\,0.0147),&\qquad\psi_3&=(2,2,3).
 \end{aligned}
\]
例如 $\delta_2(1)=\max\{0.1\times0.5,0.16\times0.3,0.28\times0.2\}\times0.5=0.028$，取最大值的前驱为 $3$。终点 $z_3^*=3$，回溯得 $z_2^*=\psi_3(3)=3$、$z_1^*=\psi_2(3)=3$，故最优路径为 $(3,3,3)$，联合概率为
\[
 p^*=0.4\times0.7\times0.5\times0.3\times0.5\times0.7=0.0147.
\]
其后验概率为 $0.0147/0.130218\approx0.11289$。由前向与后向概率还可算出逐位置后验最大化路径为 $(3,2,3)$，其联合概率只有 $0.007056$。两个结果均为可行路径，但维特比解最大化整条路径的概率，体现了两种解码目标的差别。
\end{solve}

\section{习题与方法对照}
\label{sec:lh-hmm-exercises}

以下保留教材五道习题的核心条件，并给出可核查的解答。前三题分别练习后向似然、状态平滑和最优路径，后两题把前向、后向与维特比递推统一到路径分解中。

\begin{example}
\label{ex:lh-hmm-ex-backward}
教材习题 10.1：沿用式~\eqref{eq:lh-hmm-three-boxes} 的模型，设 $T=4$，$x_{1:4}=(\text{红},\text{白},\text{红},\text{白})$，用后向算法求观测似然。
\end{example}
\begin{solve}
从 $\beta_4=(1,1,1)$ 逆序递推，得到
\[
 \begin{aligned}
 \beta_3&=(0.46,\,0.51,\,0.43),\\
 \beta_2&=(0.2461,\,0.2312,\,0.2577),\\
 \beta_1&=(0.112462,\,0.121737,\,0.104881).
 \end{aligned}
\]
例如 $\beta_3(1)=0.5\times0.5+0.2\times0.6+0.3\times0.3=0.46$，因为下一观测是白色。代入初始状态分布和第一次红色观测的概率，得
\[
 p_\lambda(x_{1:4})
 =0.1\times0.112462+0.16\times0.121737+0.28\times0.104881
 =0.0600908.\qedhere
\]
\end{solve}

\begin{example}
\label{ex:lh-hmm-ex-smoothing}
教材习题 10.2：设观测矩阵仍为式~\eqref{eq:lh-hmm-three-boxes} 中的 $B$，但转移矩阵及初始分布改为
\[
 A=\begin{pmatrix}0.5&0.1&0.4\\0.3&0.5&0.2\\0.2&0.2&0.6\end{pmatrix},\qquad
 \pi=(0.2,0.3,0.5)^{\mathsf T}.
\]
对观测 $x_{1:8}=(\text{红},\text{白},\text{红},\text{红},\text{白},\text{红},\text{白},\text{白})$，求 $P_\lambda(Z_4=3\mid X_{1:8}=x_{1:8})$。
\end{example}
\begin{solve}
先向前递推到 $4$，再从 $8$ 向后递推到 $4$。为便于核算，列出两侧中间量：
\[
 \begin{aligned}
 \alpha_1&=(0.1,\,0.12,\,0.35),\\
 \alpha_2&=(0.078,\,0.084,\,0.0822),\\
 \alpha_3&=(0.04032,\,0.026496,\,0.068124),\\
 \alpha_4&=(0.0208668,\,0.01236192,\,0.04361112),\\
 \beta_8&=(1,\,1,\,1),\\
 \beta_7&=(0.43,\,0.51,\,0.4),\\
 \beta_6&=(0.1861,\,0.2415,\,0.1762),\\
 \beta_5&=(0.105521,\,0.100883,\,0.111934),\\
 \beta_4&=(0.04586531,\,0.05280909,\,0.04280618).
 \end{aligned}
\]
分子 $\alpha_4(3)\beta_4(3)=0.0018668254527216$，分母必须汇总全部状态，而不能只使用第 $3$ 个状态的概率。于是
\[
 \begin{aligned}
 L&=\sum_{i=1}^3\alpha_4(i)\beta_4(i)=0.0034767094492824,\\
 P_\lambda(Z_4=3\mid x_{1:8})
 &=\frac{0.0018668254527216}{0.0034767094492824}\approx0.536952.
 \end{aligned}
\]
这利用了整条八时刻序列的信息；只把 $\alpha_4$ 归一化会得到滤波概率，不能回答本题的平滑问题。
\end{solve}

\begin{example}
\label{ex:lh-hmm-ex-path}
教材习题 10.3：对例~\ref{ex:lh-hmm-ex-backward} 的模型及四时刻观测，用维特比算法求最优路径。
\end{example}
\begin{solve}
前三时刻的得分与例~\ref{ex:lh-hmm-viterbi} 相同，因为维特比前向递推在时刻 $t$ 只使用前 $t$ 个观测。对第四次白色观测继续计算，得到
\[
 \delta_4=(0.00189,\,0.003024,\,0.002205),\qquad \psi_4=(1,2,3).
\]
最大终点为 $z_4^*=2$，再由 $\psi_4(2)=2$、$\psi_3(2)=2$、$\psi_2(2)=3$ 回溯，得到
\[
 z_{1:4}^*=(3,2,2,2),\qquad
 p_\lambda(z_{1:4}^*,x_{1:4})=0.003024.
\]
其后验概率为 $0.003024/0.0600908\approx0.050324$。追加一次观测后，最优路径的早期状态也可能改变，因此不能只在三时刻最优路径 $(3,3,3)$ 后添加一个状态就声称得到四时刻最优路径。
\end{solve}

\begin{example}
\label{ex:lh-hmm-ex-two-cut}
教材习题 10.4：对 $1\leq t<T$，用前向与后向概率证明
\[
 p_\lambda(x_{1:T})=\sum_{i=1}^N\sum_{j=1}^N
 \alpha_t(i)a_{ij}b_j(x_{t+1})\beta_{t+1}(j).
\]
\end{example}
\begin{solve}
事件 $\{Z_t=i,Z_{t+1}=j\}$ 随 $i,j$ 变化组成样本空间的一个有限划分。给定相邻状态后，按式~\eqref{eq:lh-hmm-joint} 将其余路径分成左侧前缀与右侧后缀并分别求和，得到
\[
 \begin{aligned}
 p_\lambda(x_{1:T})
 &=\sum_{i,j}P_\lambda(X_{1:T}=x_{1:T},Z_t=i,Z_{t+1}=j)\\
 &=\sum_{i,j}\alpha_t(i)a_{ij}b_j(x_{t+1})\beta_{t+1}(j).
 \end{aligned}
\]
也可把式~\eqref{eq:lh-hmm-beta-rec} 代入式~\eqref{eq:lh-hmm-cut} 立即得到同一结果。此式是对联合概率求和，不需要 $L>0$；只有进一步除以 $L$ 求后验时才需要正似然条件。
\end{solve}

\begin{example}
\label{ex:lh-hmm-ex-sum-max}
教材习题 10.5：比较前向算法的 $\alpha$ 与维特比算法的 $\delta$ 的计算，说明两种输出的关系。
\end{example}
\begin{solve}
两者初值相同，都是 $\pi_jb_j(x_1)$，递推都先结合转移概率再乘当前发射概率。前向递推对前驱状态求和，保留全部前缀的总概率；维特比递推对前驱状态取最大值，并记录其来源以便回溯。因此前者终止时求和得到 $L$，后者终止时取最大值得到 $p^*$。在路径层面，这个差别正是
\[
 L=\sum_{z_{1:T}}p_\lambda(z_{1:T},x_{1:T}),\qquad
 p^*=\max_{z_{1:T}}p_\lambda(z_{1:T},x_{1:T}),\qquad 0\leq p^*\leq L.
\]
若 $L>0$，等号 $p^*=L$ 当且仅当只有一条路径具有正联合概率。Baum--Welch 的 E 步需要对全部路径取后验期望，故使用前向与后向算法；若改成只用维特比最优路径计数，就改变了训练算法，不再是这里推导的 EM 更新。
\end{solve}

本章的主要依赖关系是：联合分布分解决定前向与后向递推，两者给出状态与转移的后验，从而把有标注数据的频数估计推广为无标注数据的期望计数估计；将前缀聚合中的求和换为最大化，则得到求整条最佳路径的维特比算法。进一步阅读时，可把 HMM 看作动态贝叶斯网络（dynamic Bayesian network）的一种特例；教材还以概率上下文无关文法（probabilistic context-free grammar）提示从隐藏状态序列到隐藏树结构的推广。这些拓展不改变本章始终需要区分的三类对象：观测似然、单状态后验与整条路径的联合或后验概率。




% 教材来源：机器学习/textbooks/机器学习（分章原文）/机器学习-12.pdf。
% 覆盖教材第 12 章正文（原书第 237--241 页）；随后空白页与无监督学习篇扉页不另立正文。

\fi

\selectchapter{11}
\ifcompilechapter
\chapter{条件随机场}
\label{chap:lihang-crf}

% 研读来源：机器学习/textbooks/机器学习（分章原文）/机器学习-11.pdf，教材第 215--236 页。
% 覆盖教材第 11 章全部小节、三个例题及四道习题；补充正性条件、边界约定与关键证明。

序列标注要求根据一整段观测为各个位置选择标记。例如，词性标注既需要利用词语本身及其上下文，也需要考虑相邻词性之间的搭配。若逐位置独立分类，就不能直接约束输出标记的组合；若对所有标记序列逐一计算，长度为 $n$、每个位置有 $m$ 个候选标记时又要比较 $m^n$ 种可能。\emph{条件随机场}（conditional random field，CRF）直接建模给定输入后的输出联合分布，并利用输出图结构表达条件独立性。\emph{线性链条件随机场}（linear-chain CRF）把标记之间的作用限制在相邻位置，从而兼顾序列级建模与动态规划计算。本章先建立无向图模型与因子分解，再依次讨论 CRF 的表示、概率计算、参数学习和最优序列预测。

\section{概率无向图模型}
\label{sec:lihang-crf-undirected}

\subsection{模型定义与三种马尔可夫性质}
\label{subsec:lihang-crf-markov}

设 $G=(V,E)$ 是有限无向图，$V$ 为结点集，$E$ 为无向边集。每个结点 $v$ 对应一个取值于有限集合 $\mathcal Y_v$ 的随机变量 $Y_v$；记 $Y=(Y_v)_{v\in V}$、$Y_A=(Y_v)_{v\in A}$，并以小写 $y,y_A$ 表示取值。以下条件独立关系只在条件事件具有正概率时解释。记 $N(v)=\{u\in V:\{u,v\}\in E\}$ 为 $v$ 的邻居集合，$A\perp B\mid C$ 表示随机对象 $A,B$ 在给定 $C$ 后条件独立。

\begin{definition}
\label{def:lihang-crf-markov-properties}
联合分布 $p(y)$ 相对于 $G$ 的三种\emph{马尔可夫性质}（Markov properties）分别如下。
\begin{enumerate}
  \item \emph{成对马尔可夫性质}（pairwise Markov property）：对任意没有边相连的不同结点 $u,v$，有
  \[
    Y_u\perp Y_v\mid Y_{V\setminus\{u,v\}}.
  \]
  \item \emph{局部马尔可夫性质}（local Markov property）：对任意结点 $v$，有
  \[
    Y_v\perp Y_{V\setminus(\{v\}\cup N(v))}\mid Y_{N(v)}.
  \]
  \item \emph{全局马尔可夫性质}（global Markov property）：对任意两两不交的结点子集 $A,B,C$，若从 $A$ 到 $B$ 的每条路径都经过 $C$，即 $C$ 分离 $A,B$，则
  \[
    Y_A\perp Y_B\mid Y_C.
  \]
\end{enumerate}
满足全局马尔可夫性质的分布称为相对于 $G$ 的\emph{概率无向图模型}（probabilistic undirected graphical model），也称\emph{马尔可夫随机场}（Markov random field，MRF）。在严格正的有限离散情形，三种性质等价，可以用其中任一种定义。
\end{definition}

这里无边表示给定其余变量后的条件独立，不表示两个变量在未给定其他变量时独立；有边则允许二者存在直接作用，但不保证特定参数下确实相关。对链 $1-2-3$，全局性质给出 $Y_1\perp Y_3\mid Y_2$，而 $Y_1,Y_3$ 的边缘分布仍可相关。全局性质蕴含局部性质，是因为邻居集合隔开了一个结点与其余结点；局部性质又蕴含成对性质，可由条件独立的分解与弱联合性质得到。反方向需要附加条件，下节给出的严格正性保证了这一点，不能对含零概率的任意分布直接宣称三者等价。

\subsection{团、势函数与因子分解}
\label{subsec:lihang-crf-factorization}

\begin{definition}
\label{def:lihang-crf-clique}
图 $G$ 中任意两个不同结点均有边相连的非空结点子集称为\emph{团}（clique）。若一个团不真包含于其他团，则称为\emph{极大团}（maximal clique）。教材使用“最大团”指这一概念；它与在所有团中结点数最多的\emph{最大团}（maximum clique）应作区分。记全部极大团的集合为 $\mathcal C$。
\end{definition}

\begin{theorem}[有限离散情形的 Hammersley--Clifford 定理]
\label{thm:lihang-crf-hammersley-clifford}
设 $V\neq\varnothing$，每个 $\mathcal Y_v$ 有限，且对笛卡尔积 $\prod_{v\in V}\mathcal Y_v$ 中的每个配置 $y$ 都有 $p(y)>0$。则相对于 $G$ 的成对、局部、全局马尔可夫性质等价，并等价于极大团因子分解
\begin{equation}
\label{eq:lihang-crf-clique-factorization}
  p(y)=\frac1Z\prod_{C\in\mathcal C}\psi_C(y_C),\qquad
  Z=\sum_y\prod_{C\in\mathcal C}\psi_C(y_C),\qquad \psi_C(y_C)>0.
\end{equation}
这里 $Z$ 是\emph{配分函数}（partition function），亦称规范化因子；$\psi_C$ 是\emph{势函数}（potential function）。
\end{theorem}

\begin{proof}
先证成对性质推出因子分解。固定基准配置 $y^0$，令 $\ell(y)=\log p(y)$。对每个 $A\subseteq V$ 定义
\[
  h_A(y_A)=\sum_{B\subseteq A}(-1)^{|A|-|B|}
      \ell(y_B,y^0_{V\setminus B}).
\]
对 $A$ 求和并交换有限和，一个固定 $B$ 的系数为 $\sum_{A:B\subseteq A\subseteq V}(-1)^{|A|-|B|}$，该系数仅在 $B=V$ 时为 $1$，其余为 $0$，故 $\ell(y)=\sum_{A\subseteq V}h_A(y_A)$。若 $A$ 不是团，则存在不相邻的 $u,v\in A$。由成对条件独立及正性，对任意固定的其余配置 $z$ 都有
\[
  \ell(y_u,y_v,z)-\ell(y_u,y_v^0,z)
  -\ell(y_u^0,y_v,z)+\ell(y_u^0,y_v^0,z)=0.
\]
将 $h_A$ 中的项按是否包含 $u,v$ 分组，每组恰为上述混合差分乘以一个符号，因此 $h_A=0$。剩余的非空 $A$ 都是团，分别把 $h_A$ 归入包含 $A$ 的某个极大团，并把常数 $h_\varnothing$ 归入任一极大团。将各组求和后取指数，便得到严格正的团势分解；规范化给出式~\eqref{eq:lihang-crf-clique-factorization}。

再证因子分解推出全局性质。若 $C$ 分离 $A,B$，删去 $C$ 后，任何连通分支都不能同时与 $A,B$ 相交。将这些分支分成两组 $D_A,D_B$，使 $A\subseteq D_A$、$B\subseteq D_B$；其余分支可任意归入一组。没有团同时包含 $D_A,D_B$ 中的结点，故给定 $y_C$ 后，势函数乘积可写成 $g(y_{D_A},y_C)h(y_{D_B},y_C)$。其规范化因子也分成两个有限和的乘积，因此 $Y_{D_A}\perp Y_{D_B}\mid Y_C$，边缘化即得 $Y_A\perp Y_B\mid Y_C$。结合上节的全局性质蕴含局部性质、局部性质蕴含成对性质，得到全部等价关系。
\end{proof}

势函数给出局部配置的相对权重，不必逐项求和为 $1$，也一般不是某个边缘概率或条件概率。正势可以写成 $\psi_C(y_C)=\exp\{-E_C(y_C)\}$，其中 $E_C$ 称为局部能量（local energy）；总能量越小，配置的未归一化权重越大。有限取值及严格正势保证 $0<Z<\infty$。若允许零势来表达禁配，仍可以直接定义并规范化因子模型，但不能原样援用上述严格正版本的反向结论。

\begin{example}
\label{ex:lihang-crf-four-node-cliques}
教材图 11.3 及习题 11.1 的四结点图，其边为 $\{1,2\},\{1,3\},\{2,3\},\{2,4\},\{3,4\}$。写出严格正的图模型因子分解，并说明缺边的含义。
\end{example}
\begin{solve}
两个极大团是 $\{1,2,3\}$ 与 $\{2,3,4\}$，所以
\[
  p(y_1,y_2,y_3,y_4)
  =\frac{\psi_{123}(y_1,y_2,y_3)\psi_{234}(y_2,y_3,y_4)}
  {\displaystyle\sum_{z_1,z_2,z_3,z_4}
       \psi_{123}(z_1,z_2,z_3)\psi_{234}(z_2,z_3,z_4)}.
\]
没有 $\{1,4\}$ 这条边，表示 $Y_1\perp Y_4\mid(Y_2,Y_3)$；不能据此删去任一团内部的其他作用，也不能断言 $Y_1,Y_4$ 边缘独立。
\end{solve}

\section{条件随机场的定义与形式}
\label{sec:lihang-crf-representations}

\subsection{条件随机场与线性链结构}
\label{subsec:lihang-crf-definition}

\begin{definition}
\label{def:lihang-crf}
设 $X$ 是输入随机对象，$Y=(Y_v)_{v\in V}$ 是输出随机变量族。若对输入分布几乎处处的 $x$，条件分布 $p(y\mid x)$ 相对于同一无向图 $G$ 满足全局马尔可夫性质，则称该条件分布模型为\emph{条件随机场}。在严格正的有限输出空间中，等价地可要求每个结点满足
\begin{equation}
\label{eq:lihang-crf-conditional-markov}
  p(y_v\mid x,y_{V\setminus\{v\}})
  =p(y_v\mid x,y_{N(v)}).
\end{equation}
\end{definition}

这里约束的是给定整个 $X$ 后输出变量之间的依赖；不要求 $X$ 本身具有相同图结构，也不要求观测分量相互独立。模型直接指定 $p(y\mid x)$，无须同时估计 $p(x)$，因此属于判别模型。对连续输入，给定 $X=x$ 理解为使用条件分布的一个版本，模型公式可在所考虑的输入空间上直接定义。

本章主要讨论 $X=(X_1,\ldots,X_n)$、$Y=(Y_1,\ldots,Y_n)$，其中 $n\geq1$，每个 $Y_i$ 取值于同一有限标记集 $\mathcal Y=\{1,\ldots,m\}$。若输出图的边为 $\{i,i+1\}$，则得到线性链 CRF；在严格正情形，其局部性质为
\[
  p(y_i\mid x,y_1,\ldots,y_{i-1},y_{i+1},\ldots,y_n)
  =p(y_i\mid x,y_{i-1},y_{i+1}),
\]
端点只保留实际存在的邻居，$n=1$ 时没有其他输出变量。\emph[caution]{输出的局部马尔可夫性质不限制特征只能观察 $x_i$}：某个位置的势可以使用整个输入序列，只要它依赖的输出标记仍限于该结点或相邻结点。

\subsection{参数化形式与局部特征}
\label{subsec:lihang-crf-features}

严格正的链式条件分布可以由相邻输出结点上的正势表示。实际建模时选定 $K_1$ 个\emph{转移特征}（transition features）$t_k(a,b,x,i)$ 与 $K_2$ 个\emph{状态特征}（state features）$s_l(b,x,i)$，再学习对应权重 $\lambda_k,\mu_l\in\mathbb R$，得到对数线性模型（log-linear model）
\begin{align}
\label{eq:lihang-crf-feature-score}
  S_w(y,x)
  &=\sum_{i=2}^{n}\sum_{k=1}^{K_1}\lambda_k t_k(y_{i-1},y_i,x,i)
    +\sum_{i=1}^{n}\sum_{l=1}^{K_2}\mu_l s_l(y_i,x,i),\\
\label{eq:lihang-crf-feature-model}
  p_w(y\mid x)&=\frac{\exp S_w(y,x)}{Z_w(x)},\qquad
  Z_w(x)=\sum_{z\in\mathcal Y^n}\exp S_w(z,x).
\end{align}
特征与权重均为有限实数时，输出空间有限，故条件概率严格正且和为 $1$。常用特征为取值 $0,1$ 的指示函数，例如 $t_k=\mathbf1\{y_{i-1}=a_0,y_i=b_0\}$ 或 $s_l=\mathbf1\{y_i=b_0,\ x_i\text{满足某条件}\}$；一般对数线性形式也允许实值特征。相同权重可以在多个位置共享，位置是否启用某个特征由其函数定义决定。选定的有限特征族规定了一个参数模型族，不能据此声称任意 CRF 都可由这一组预先固定的特征精确表示。

与隐马尔可夫模型（hidden Markov model，HMM）相比，CRF 无须假定给定标记后各位置观测独立。例如，HMM 的联合概率形如
\[
  p(x,y)=\pi(y_1)b_{y_1}(x_1)
         \prod_{i=2}^n a(y_{i-1},y_i)b_{y_i}(x_i).
\]
在 $p(x)>0$ 时，将它除以 $p(x)$ 就得到关于 $y$ 的链式条件分布；严格正因子取对数后可写成 CRF 的局部得分。反之，CRF 的特征可以任意依赖输入上下文，其势也无须满足 HMM 转移概率与发射概率各自的规范化约束，因此一般不必对应上述生成过程。

\begin{example}
\label{ex:lihang-crf-feature-example}
教材例 11.1。固定输入 $x$，取 $n=3$、$\mathcal Y=\{1,2\}$。下表给出全部非零特征的触发条件与权重；不满足条件时特征为 $0$，满足时为 $1$。求 $y=(1,2,2)$ 的得分与未归一化权重。
\begin{center}
\begin{tabular}{ccc}
\toprule
特征 & 取值为 $1$ 的条件 & 权重\\
\midrule
$t_1(a,b,x,i)$ & $i\in\{2,3\},\ a=1,\ b=2$ & $\lambda_1=1$\\
$t_2(a,b,x,i)$ & $i=2,\ a=1,\ b=1$ & $\lambda_2=0.6$\\
$t_3(a,b,x,i)$ & $i=3,\ a=2,\ b=1$ & $\lambda_3=1$\\
$t_4(a,b,x,i)$ & $i=2,\ a=2,\ b=1$ & $\lambda_4=1$\\
$t_5(a,b,x,i)$ & $i=3,\ a=2,\ b=2$ & $\lambda_5=0.2$\\
$s_1(b,x,i)$ & $i=1,\ b=1$ & $\mu_1=1$\\
$s_2(b,x,i)$ & $i\in\{1,2\},\ b=2$ & $\mu_2=0.5$\\
$s_3(b,x,i)$ & $i\in\{2,3\},\ b=1$ & $\mu_3=0.8$\\
$s_4(b,x,i)$ & $i=3,\ b=2$ & $\mu_4=0.5$\\
\bottomrule
\end{tabular}
\end{center}
\end{example}
\begin{solve}
转移特征在 $i=2$ 触发 $t_1$、在 $i=3$ 触发 $t_5$；状态特征依次触发 $s_1,s_2,s_4$。因此
\[
  S_w((1,2,2),x)=1+0.2+1+0.5+0.5=3.2,
  \qquad \exp S_w((1,2,2),x)=e^{3.2}.
\]
其中 $3.2$ 是对数得分，$e^{3.2}$ 是正权重，二者都不是已经规范化的条件概率；条件概率还需除以 $Z_w(x)$。特别地，$s_2$ 只在前两个位置定义为可能非零，不能因为 $y_3=2$ 就再计入一次 $0.5$。
\end{solve}

\subsection{全局特征与向量形式}
\label{subsec:lihang-crf-global-features}

令 $K=K_1+K_2$、$w=(\lambda_1,\ldots,\lambda_{K_1},\mu_1,\ldots,\mu_{K_2})^{\mathsf T}$。为统一边界，固定虚拟标记 $y_0=\mathrm{start}$、$y_{n+1}=\mathrm{stop}$，它们不是额外待预测的随机变量。将转移与状态特征合并为局部特征 $f_k(y_{i-1},y_i,x,i)$：转移特征仅在 $2\leq i\leq n$ 启用，状态特征在 $1\leq i\leq n$ 启用并忽略第一个输出参数，其余位置取零。定义
\begin{equation}
\label{eq:lihang-crf-global-features}
  F_k(y,x)=\sum_{i=1}^{n+1}f_k(y_{i-1},y_i,x,i),\qquad
  F(y,x)=(F_1(y,x),\ldots,F_K(y,x))^{\mathsf T}.
\end{equation}
于是局部特征沿序列的累积构成全局特征，式~\eqref{eq:lihang-crf-feature-model} 化为
\begin{equation}
\label{eq:lihang-crf-vector-model}
  p_w(y\mid x)=\exp\{w^{\mathsf T}F(y,x)-\log Z_w(x)\},\qquad
  Z_w(x)=\sum_y\exp\{w^{\mathsf T}F(y,x)\}.
\end{equation}
后续也允许显式加入依赖 $y_1$ 的起始特征及依赖 $y_n$ 的终止特征，分别放在 $i=1,n+1$；若不加入，终止位置的局部得分为 $0$。给所有输出序列同时增加一个只依赖 $x$ 的得分，会在分子分母中抵消，因而参数表示不必唯一。

\subsection{矩阵形式与路径权重}
\label{subsec:lihang-crf-matrix}

记 $\mathcal Y_0=\{\mathrm{start}\}$、$\mathcal Y_i=\mathcal Y$（$1\leq i\leq n$）、$\mathcal Y_{n+1}=\{\mathrm{stop}\}$，定义局部对数得分与转移势矩阵
\begin{equation}
\label{eq:lihang-crf-matrices}
  W_i(a,b\mid x)=\sum_{k=1}^K w_kf_k(a,b,x,i),\qquad
  M_i(a,b\mid x)=\exp W_i(a,b\mid x).
\end{equation}
矩阵 $M_i$ 的行对应前一个标记、列对应当前标记；$M_1$ 为 $1\times m$，$M_2,\ldots,M_n$ 为 $m\times m$，$M_{n+1}$ 为 $m\times1$。给定 $x,w$ 后这些都是确定的正矩阵，一般不是行和为 $1$ 的随机转移矩阵。其路径乘积与规范化关系为
\begin{equation}
\label{eq:lihang-crf-matrix-product}
  p_w(y\mid x)=\frac{\displaystyle\prod_{i=1}^{n+1}M_i(y_{i-1},y_i\mid x)}{Z_w(x)},\qquad
  Z_w(x)=M_1M_2\cdots M_{n+1}.
\end{equation}
右端矩阵乘积是一个 $1\times1$ 标量。矩阵乘法对全部中间指标求和，恰好枚举所有标记路径，因此该标量就是配分函数。教材把起止状态嵌入统一大小的方阵后取乘积的 $(\mathrm{start},\mathrm{stop})$ 元素，与这里的矩形矩阵表示等价；方阵中为屏蔽非法起止边而填入的零不属于合法路径上的正势。

\begin{example}
\label{ex:lihang-crf-symbolic-paths}
教材例 11.2 的矩形表示为
\[
  M_1=(a_1,a_2),\quad
  M_2=\begin{pmatrix}b_{11}&b_{12}\\b_{21}&b_{22}\end{pmatrix},\quad
  M_3=\begin{pmatrix}c_{11}&c_{12}\\c_{21}&c_{22}\end{pmatrix},\quad
  M_4=\begin{pmatrix}1\\1\end{pmatrix}.
\]
求全部路径的权重与配分函数，其中各 $a,b,c$ 均为正数。
\end{example}
\begin{solve}
对任意 $y\in\{1,2\}^3$，路径权重统一为 $a_{y_1}b_{y_1y_2}c_{y_2y_3}$。例如 $(1,1,2)$ 对应 $a_1b_{11}c_{12}$，$(2,2,1)$ 对应 $a_2b_{22}c_{21}$。对最后一个标记先求和，再对前两个标记求和，得到
\[
  Z_w(x)=(a_1b_{11}+a_2b_{21})(c_{11}+c_{12})
        +(a_1b_{12}+a_2b_{22})(c_{21}+c_{22}).
\]
展开正好有八项，每条路径各出现一次；这也直接说明矩阵乘法为何能代替显式枚举。
\end{solve}

\section{条件随机场的概率计算}
\label{sec:lihang-crf-probabilities}

\subsection{前向与后向递推}
\label{subsec:lihang-crf-forward-backward}

以下固定 $x,w$，必要时省略其记号。\emph{前向量}（forward quantity）$\alpha_i(b)$ 是从起点到位置 $i$、以 $b$ 结尾的所有前缀路径的权重和；\emph{后向量}（backward quantity）$\beta_i(b)$ 是固定位置 $i$ 的标记为 $b$ 后，从 $i+1$ 到终点的所有后缀路径的权重和。它们都是未归一化的量，后向量不再包含到达位置 $i$ 的那条边。初始化与递推为
\begin{align}
\label{eq:lihang-crf-forward}
  \alpha_0(\mathrm{start})&=1,&
  \alpha_i(b)&=\sum_{a\in\mathcal Y_{i-1}}\alpha_{i-1}(a)M_i(a,b),
       &&1\leq i\leq n+1,\\
\label{eq:lihang-crf-backward}
  \beta_{n+1}(\mathrm{stop})&=1,&
  \beta_i(a)&=\sum_{b\in\mathcal Y_{i+1}}M_{i+1}(a,b)\beta_{i+1}(b),
       &&0\leq i\leq n.
\end{align}
若 $\alpha_i,\beta_i$ 均记为列向量，矩阵写法是 $\alpha_i^{\mathsf T}=\alpha_{i-1}^{\mathsf T}M_i$、$\beta_i=M_{i+1}\beta_{i+1}$。第 $i$ 步两者的维度均为 $|\mathcal Y_i|$；端点维度为 $1$，内部维度为 $m$。

\begin{proposition}
\label{prop:lihang-crf-partition-split}
对任意 $i=0,1,\ldots,n+1$，有
\begin{equation}
\label{eq:lihang-crf-partition-split}
  Z_w(x)=\sum_{b\in\mathcal Y_i}\alpha_i(b)\beta_i(b)
        =\alpha_i^{\mathsf T}\beta_i
        =\alpha_{n+1}(\mathrm{stop})=\beta_0(\mathrm{start}).
\end{equation}
\end{proposition}
\begin{proof}
每条完整路径在位置 $i$ 有唯一标记 $b$，并唯一分成到达 $b$ 的前缀和从 $b$ 出发的后缀。其边权乘积也相应分成两部分；固定 $b$ 后，分别对两部分求和，得到 $\alpha_i(b)\beta_i(b)$。再对 $b$ 求和即遍历全部完整路径，端点等式由初始化得到。
\end{proof}

\begin{notealert}[type=caution,title={首尾势与教材习题 11.2 的边界式}]
当没有终止特征时，$M_{n+1}=\mathbf1$，故 $Z_w(x)=\alpha_n^{\mathsf T}\mathbf1$；起点处始终有 $Z_w(x)=M_1\beta_1=\alpha_1^{\mathsf T}\beta_1$。在本节的后向量定义下，直接写 $Z_w(x)=\mathbf1^{\mathsf T}\beta_1$ 一般会漏掉 $M_1$，只有起始权重全为 $1$ 或另行把首势吸收到后向量时才成立。例如 $n=1$、$M_1=(1,2)$、$M_2=(1,1)^{\mathsf T}$ 时，$Z_w(x)=3$，而 $\mathbf1^{\mathsf T}\beta_1=2$。因此教材习题 11.2 应连同端点与索引约定理解，其通用形式是式~\eqref{eq:lihang-crf-partition-split}。
\end{notealert}

\subsection{结点与相邻结点的边缘概率}
\label{subsec:lihang-crf-marginals}

将命题~\ref{prop:lihang-crf-partition-split} 的路径划分限制在固定标记上，就得到结点边缘分布（node marginal）
\begin{equation}
\label{eq:lihang-crf-node-marginal}
  p_w(Y_i=b\mid x)=\frac{\alpha_i(b)\beta_i(b)}{Z_w(x)},\qquad 1\leq i\leq n.
\end{equation}
固定相邻两点及连接它们的边，得到边缘分布（edge marginal）
\begin{equation}
\label{eq:lihang-crf-edge-marginal}
  p_w(Y_{i-1}=a,Y_i=b\mid x)
  =\frac{\alpha_{i-1}(a)M_i(a,b)\beta_i(b)}{Z_w(x)},
  \qquad 1\leq i\leq n+1.
\end{equation}
边界中的起止标记固定，因而 $i=1,n+1$ 的公式也有确定含义。对边缘分布求和应恢复相邻结点的边缘分布，而全部边缘概率之和应为 $1$；这些关系可用来检查递推是否漏乘或重复计算了某个局部势。已知完整 $y$ 时，其概率则由路径乘积除以 $Z_w(x)$ 直接计算。

\subsection{特征期望、计算量与数值实现}
\label{subsec:lihang-crf-expectations}

全局特征是局部特征的和。利用期望的线性性及式~\eqref{eq:lihang-crf-edge-marginal}，得到
\begin{equation}
\label{eq:lihang-crf-feature-expectation}
  \mathbb E_{p_w(\cdot\mid x)}[F_k(Y,x)]
  =\sum_{i=1}^{n+1}\sum_{a\in\mathcal Y_{i-1}}\sum_{b\in\mathcal Y_i}
     f_k(a,b,x,i)
     \frac{\alpha_{i-1}(a)M_i(a,b)\beta_i(b)}{Z_w(x)}.
\end{equation}
若 $f_k$ 是仅依赖 $y_i$ 的状态特征，可以直接用结点边缘概率求和，避免对不相关的前一标记重复计算。设训练输入的经验分布为 $\widehat p_X(x)$，则模型与经验输入分布组合成联合分布 $q_w(x,y)=\widehat p_X(x)p_w(y\mid x)$，并有
\[
  \mathbb E_{q_w}[F_k(Y,X)]
  =\sum_x\widehat p_X(x)\mathbb E_{p_w(\cdot\mid x)}[F_k(Y,x)].
\]
这只是用经验输入分布对条件期望加权，不等于未知真实总体联合分布下的期望，也不等于训练输出的经验特征平均。

在局部势已经计算、每次标量运算视为常数时间的条件下，稠密链的一次前向与后向计算各需 $O(nm^2)$ 时间，而显式枚举需要考察 $m^n$ 条序列。只求 $Z_w(x)$ 时前向递推可滚动保存 $O(m)$ 个量；保存全部前后向量需要 $O(nm)$ 空间。构造势或计算所有特征期望还需计入特征求值成本，朴素地逐一处理 $K$ 个特征时可达 $O(Knm^2)$，不能将这部分成本也忽略。

直接相乘指数势可能溢出或下溢。可在对数域定义 $A_i(b)=\log\alpha_i(b)$，将前向递推写成
\[
  A_i(b)=\operatorname{LSE}_{a}\{A_{i-1}(a)+W_i(a,b\mid x)\},\qquad
  \operatorname{LSE}(z_1,\ldots,z_r)
  =c+\log\sum_{j=1}^r e^{z_j-c},\quad c=\max_j z_j.
\]
这就是稳定的 log-sum-exp 计算，后向量同理；最后得到 $\log Z_w(x)$，再以局部对数权重减去它计算边缘概率。禁配边可设对数势为 $-\infty$，但须存在至少一条合法完整路径；若没有合法路径，$Z_w(x)=0$，条件分布本身便没有定义。

\section{条件随机场的学习算法}
\label{sec:lihang-crf-learning}

设训练集 $D=\{(x^{(r)},y^{(r)})\}_{r=1}^N$，$N\geq1$，样本对独立同分布，各样本内部的标记依赖由 CRF 表示。样本长度可以不同，式~\eqref{eq:lihang-crf-global-features} 中的求和上限随相应输入长度变化。假设特征族固定，所有训练输入的输出空间有限、特征值有限，且训练标注是合法输出；本节优化变量仅为 $w\in\mathbb R^K$。经验分布及经验特征均值记为
\[
  \widehat p(x,y)=\frac1N\sum_{r=1}^N
          \mathbf1\{(x^{(r)},y^{(r)})=(x,y)\},\quad
  \widehat p_X(x)=\sum_y\widehat p(x,y),\quad
  \widehat F_k=\frac1N\sum_{r=1}^N F_k(y^{(r)},x^{(r)}).
\]
这里经验分布是有限支持的概率测度，即使输入本身连续，上述对经验支持的求和仍然有意义。平均条件对数似然（conditional log-likelihood）与带二次正则项的平均负对数似然分别为
\begin{align}
\label{eq:lihang-crf-log-likelihood}
  \ell(w)&=\frac1N\sum_{r=1}^N
    \left[w^{\mathsf T}F(y^{(r)},x^{(r)})-\log Z_w(x^{(r)})\right],\\
\label{eq:lihang-crf-objective}
  J_\rho(w)&=-\ell(w)+\frac\rho2\lVert w\rVert_2^2,\qquad \rho\geq0.
\end{align}
极大化条件似然等价于最小化 $J_0$；$\rho>0$ 则给出正则化条件极大似然估计。使用平均量时，梯度、经验期望和正则化系数都应采用这一尺度，不能把总和与平均的公式混用。

\subsection{改进的迭代尺度法}
\label{subsec:lihang-crf-iis}

\emph{改进的迭代尺度法}（improved iterative scaling，IIS）以增量 $\delta=(\delta_1,\ldots,\delta_K)^{\mathsf T}$ 更新 $w$。以下讨论无正则项的 $\ell(w)$，并额外假设全局特征 $F_k(y,x)\geq0$；指示型局部特征满足这一要求，但任意实值 CRF 特征未必满足。记
\[
  T(x,y)=\sum_{k=1}^K F_k(y,x),\qquad
  \mathbb E_w[H]=\sum_x\widehat p_X(x)\sum_y p_w(y\mid x)H(x,y).
\]
与样本上的经验均值 $\widehat F_k$ 不同，$\mathbb E_w$ 中的输出按当前模型分布取平均。

\begin{proposition}
\label{prop:lihang-crf-iis-update}
在上述非负特征假设下，若每个增量方程
\begin{equation}
\label{eq:lihang-crf-iis-equation}
  \mathbb E_w\!\left[F_k(Y,X)e^{\delta_kT(X,Y)}\right]
  =\widehat F_k,\qquad k=1,\ldots,K,
\end{equation}
都有有限解，则以同一个旧参数 $w$ 计算全部增量并同时更新为 $w+\delta$，不会降低平均对数似然。
\end{proposition}
\begin{proof}
由配分函数之比，
\[
  \ell(w+\delta)-\ell(w)
  =\delta^{\mathsf T}\widehat F
   -\sum_x\widehat p_X(x)\log
        \mathbb E_{p_w(\cdot\mid x)}e^{\delta^{\mathsf T}F(Y,x)}.
\]
用 $\log u\leq u-1$，并在 $T>0$ 时以 $F_k/T$ 为权重应用指数函数的凸性，得到
\[
  e^{\sum_k\delta_kF_k}
  \leq\sum_k\frac{F_k}{T}e^{\delta_kT}.
\]
当 $T=0$ 时全部 $F_k=0$，该配置对似然增量下界的贡献为零。由此得到可分离的下界
\[
  G(\delta)=\sum_k\delta_k\widehat F_k
    -\mathbb E_w\!\left[
       \mathbf1\{T>0\}\sum_k\frac{F_k}{T}
                     (e^{\delta_kT}-1)\right],
  \qquad \ell(w+\delta)-\ell(w)\geq G(\delta),
\]
其中 $T=0$ 的被积项定义为 $0$，且 $G(0)=0$。对 $\delta_k$ 求导得到
\[
  \frac{\partial G}{\partial\delta_k}
    =\widehat F_k-\mathbb E_w[F_ke^{\delta_kT}],\qquad
  \frac{\partial^2G}{\partial\delta_k^2}
    =-\mathbb E_w[F_kT e^{\delta_kT}]\leq0.
\]
式~\eqref{eq:lihang-crf-iis-equation} 因而逐坐标极大化这个凹下界，给出 $G(\delta)\geq G(0)=0$，从而似然不降。
\end{proof}

算法从 $w=0$ 等有限初值开始，在每一轮计算当前模型，求解全部标量方程，再同时更新；达到似然增量或参数增量容差时停止，并设置最大迭代次数。若改为逐坐标立刻更新，则下一坐标应使用已经更新的模型，不能混用两种约定。对活跃特征，方程左端关于 $\delta_k$ 严格递增，其导数为 $\mathbb E_w[F_kT e^{\delta_kT}]>0$，可采用带区间保护的牛顿法求根。若 $\widehat F_k=0$ 而模型期望始终为正，有限根一般不存在，增量会趋向 $-\infty$；需要删除无效特征、施加参数限制或改用适当的正则化目标，不能把无限增量当作普通更新。该推导也不保证未经修改的公式能优化带二次正则项的目标。

基本 IIS 方程中 $T(x,y)$ 随输出变化，$F_k e^{\delta_kT}$ 已不一定是原有局部特征的线性和，因而仅靠普通结点、边缘分布一般不能直接计算它。教材进一步引入松弛特征，使指数项与输出序列脱离。选择常数
\[
  S>\max_{x\in\operatorname{supp}(\widehat p_X),\ y}T(x,y),\qquad
  F_0(y,x)=S-T(x,y)\geq0,
\]
则扩充后全部特征之和恒为 $S$。$F_0$ 的权重带来的得分为 $w_0S-\sum_{k=1}^K w_0F_k$，其中 $w_0S$ 在条件规范化时抵消，因此加入它不扩大原有条件分布族。对扩充后的全部特征应用迭代尺度更新，得到教材的\emph{算法 S}：
\begin{equation}
\label{eq:lihang-crf-iis-S}
  \delta_k=\frac1S\log\frac{\widehat F_k}{\mathbb E_w[F_k]},
  \qquad k=0,1,\ldots,K,
\end{equation}
这里要求分子分母均为正，退化特征需另行处理。原有特征的期望由前向--后向计算，松弛特征的期望为 $S-\sum_{k=1}^K\mathbb E_w[F_k]$。在期望比相同的情况下，$S$ 越大，增量的绝对值越小，可能导致更新缓慢。

\emph{算法 T} 对每个训练输入使用较紧的上界 $S(x)=\max_yT(x,y)$，并加入 $F_0(y,x)=S(x)-T(x,y)$。求这个最大值需要对局部特征总和运行最大加法动态规划，而不是把求和型前向递推原样使用。扩充后总特征数为 $S(x)$，更新方程变为
\begin{equation}
\label{eq:lihang-crf-iis-T}
  \widehat F_k
  =\sum_x\widehat p_X(x)
     \mathbb E_{p_w(\cdot\mid x)}[F_k(Y,x)]e^{\delta_kS(x)}.
\end{equation}
若特征是非负整数计数，则 $S(x)$ 为整数。令 $u=e^{\delta_k}>0$，按 $S(x)$ 的取值分组，方程成为
\[
  \sum_{s=0}^{S_{\max}}a_{k,s}u^s=\widehat F_k,\qquad
  a_{k,s}=\sum_{x:S(x)=s}\widehat p_X(x)
                   \mathbb E_{p_w(\cdot\mid x)}[F_k(Y,x)].
\]
在目标期望为正、该特征活跃时，左端在 $u>0$ 上严格递增，因而有唯一正根；不是声称该多项式只有一个实根。非整数特征时仍可对 $\delta_k$ 解单调方程，但通常不再是多项式。以上两种松弛方式使模型输出期望可以通过原链的前向--后向递推计算，代价则是额外特征及可能较小的更新幅度。

\subsection{梯度、凸性与拟牛顿法}
\label{subsec:lihang-crf-gradient-bfgs}

对有限输出空间，配分函数的和可以逐项微分，得到
\begin{equation}
\label{eq:lihang-crf-logZ-gradient}
  \nabla_w\log Z_w(x)=\mathbb E_{p_w(\cdot\mid x)}[F(Y,x)].
\end{equation}
因此平均负对数似然的梯度为“模型期望减经验期望”：
\begin{equation}
\label{eq:lihang-crf-gradient}
  \nabla J_\rho(w)=\frac1N\sum_{r=1}^N
    \left[\mathbb E_{p_w(\cdot\mid x^{(r)})}F(Y,x^{(r)})
           -F(y^{(r)},x^{(r)})\right]+\rho w.
\end{equation}
每个条件期望都由式~\eqref{eq:lihang-crf-feature-expectation} 计算，无须枚举所有标记序列。无正则且有限最优点存在时，梯度为零意味着模型特征期望与经验特征平均匹配；加入正则项后，驻点满足两者之差为 $-\rho w$，不再是逐项完全相等。

\begin{proposition}
\label{prop:lihang-crf-convexity}
在固定特征、有限训练输出空间与有限特征值的条件下，$J_\rho$ 在 $\mathbb R^K$ 上可微且凸，其 Hessian 为
\begin{equation}
\label{eq:lihang-crf-hessian}
  \nabla^2J_\rho(w)
  =\frac1N\sum_{r=1}^N
       \operatorname{Cov}_{p_w(\cdot\mid x^{(r)})}
                    [F(Y,x^{(r)})]+\rho I_K.
\end{equation}
若 $\rho>0$，则 $J_\rho$ 强凸，存在唯一有限的最小点。
\end{proposition}
\begin{proof}
继续对式~\eqref{eq:lihang-crf-logZ-gradient} 求导，第 $(k,l)$ 个元素为
\[
  \mathbb E_w[F_kF_l\mid x]
     -\mathbb E_w[F_k\mid x]\mathbb E_w[F_l\mid x],
\]
从而得到式~\eqref{eq:lihang-crf-hessian}。任意方向 $a$ 上的二次型为各条件方差 $\operatorname{Var}_w(a^{\mathsf T}F\mid x^{(r)})$ 的平均加上 $\rho\lVert a\rVert_2^2$，故非负；$\rho>0$ 时至少为 $\rho\lVert a\rVert_2^2$。又因负对数概率非负，$J_\rho(w)\geq\rho\lVert w\rVert_2^2/2$，所以目标连续且随 $\lVert w\rVert_2\to\infty$ 趋向无穷，在某个有界闭集内达到最小值；强凸性保证该点唯一。
\end{proof}

无正则时，凸性只保证任意有限驻点都是全局最小点，不保证最小点存在或参数唯一。例如某个方向上的 $a^{\mathsf T}F(y,x)$ 对每个训练输入都与 $y$ 无关，则沿该方向改变参数不改变条件概率；若训练标注可被得分持续拉开，也可能只有参数趋向无穷时才能逼近似然上确界。上述凸性依赖特征固定；若特征本身由待训练的非线性网络产生，对全部参数联合训练时不能沿用这个凸性结论。

教材习题 11.3 可直接由梯度式得到\emph{梯度下降法}（gradient descent）。输入为训练集、特征、$\rho\geq0$、初值 $w^{(0)}$、容差 $\varepsilon>0$ 及最大迭代次数；每轮先对全部样本计算局部势、前后向量、$\log Z$ 和特征期望，再组成 $g^{(t)}=\nabla J_\rho(w^{(t)})$。若 $\lVert g^{(t)}\rVert_\infty\leq\varepsilon$ 则停止，否则选择步长 $\eta_t>0$ 并更新
\begin{equation}
\label{eq:lihang-crf-gradient-descent}
  w^{(t+1)}=w^{(t)}-\eta_tg^{(t)}.
\end{equation}
例如固定 $c,\gamma\in(0,1)$，从一个正试探步长开始反复乘以 $\gamma$，直到满足 Armijo 充分下降条件
\[
  J_\rho(w^{(t)}-\eta_tg^{(t)})
  \leq J_\rho(w^{(t)})-c\eta_t\lVert g^{(t)}\rVert_2^2.
\]
非零梯度处足够小的步长会满足该条件。有限特征下 Hessian 有统一上界；在 $\rho>0$、每轮从固定正试探步长开始回溯且计算精确时，这给出到唯一最小点的收敛。达到次数上限时应返回当前近似值及未满足容差的状态，不能声称已经求得精确最优参数。

\emph{拟牛顿法}（quasi-Newton method）利用梯度变化近似曲率，避免显式构造高维 Hessian。教材采用 BFGS 方法：以正定对称矩阵 $B_t$ 近似 Hessian，解
\[
  B_t p_t=-g^{(t)},\qquad
  w^{(t+1)}=w^{(t)}+\eta_t p_t,
\]
并记 $s_t=w^{(t+1)}-w^{(t)}$、$z_t=g^{(t+1)}-g^{(t)}$，更新
\begin{equation}
\label{eq:lihang-crf-bfgs}
  B_{t+1}=B_t-
      \frac{B_ts_ts_t^{\mathsf T}B_t}{s_t^{\mathsf T}B_ts_t}
      +\frac{z_tz_t^{\mathsf T}}{z_t^{\mathsf T}s_t}.
\end{equation}
非零梯度且 $B_t$ 正定时，$g^{(t)\mathsf T}p_t=-g^{(t)\mathsf T}B_t^{-1}g^{(t)}<0$，故 $p_t$ 是下降方向。线搜索应同时保证充分下降及曲率条件，例如采用 Wolfe 条件，从而得到 $z_t^{\mathsf T}s_t>0$ 并保持更新后的正定性；曲率退化时不能盲目使用含零分母的公式。算法重复计算目标、梯度、搜索方向、线搜索和矩阵更新，直到梯度足够小或达到迭代上限。完整 $B_t$ 需 $O(K^2)$ 存储，有限记忆的 L-BFGS 形式则只保存若干对 $s_t,z_t$。这些优化方法与 IIS 求解同一类条件似然目标，但不要求特征非负。

\section{条件随机场的预测算法}
\label{sec:lihang-crf-decoding}

给定输入 $x$ 后，序列级最大后验预测（maximum a posteriori decoding）求
\begin{equation}
\label{eq:lihang-crf-map}
  y^*\in\arg\max_y p_w(y\mid x)
      =\arg\max_y w^{\mathsf T}F(y,x)
      =\arg\max_y\sum_{i=1}^{n+1}W_i(y_{i-1},y_i\mid x).
\end{equation}
这是因为 $Z_w(x)$ 与待比较的 $y$ 无关，指数函数又严格递增。只求最优序列时可以省去配分函数；若还要报告该序列的条件概率，则仍须计算 $\log Z_w(x)$，不能把最优得分本身解释为置信概率。

\subsection{维特比递推与回溯}
\label{subsec:lihang-crf-viterbi}

\emph{维特比算法}（Viterbi algorithm）保存到达每个末端标记的最高前缀得分 $d_i(b)$，并记录取得该值的前驱标记 $\pi_i(b)$。初始化 $d_0(\mathrm{start})=0$，对 $i=1,\ldots,n$ 依次计算
\begin{align}
\label{eq:lihang-crf-viterbi-score}
  d_i(b)&=\max_{a\in\mathcal Y_{i-1}}
       \{d_{i-1}(a)+W_i(a,b\mid x)\},\\
\label{eq:lihang-crf-viterbi-backpointer}
  \pi_i(b)&\in\arg\max_{a\in\mathcal Y_{i-1}}
       \{d_{i-1}(a)+W_i(a,b\mid x)\}.
\end{align}
最后须计入终止势，令
\[
  y_n^*\in\arg\max_b\{d_n(b)+W_{n+1}(b,\mathrm{stop}\mid x)\},\qquad
  S^*=\max_b\{d_n(b)+W_{n+1}(b,\mathrm{stop}\mid x)\}.
\]
再按 $y_{i-1}^*=\pi_i(y_i^*)$，从 $i=n$ 回溯至 $2$，得到 $y^*=(y_1^*,\ldots,y_n^*)$。所有并列选择采用预先固定的规则，例如取标记编号最小者。没有终止特征时最后一项为零，才可直接对 $d_n$ 取最大值。算法时间为 $O(nm^2)$；只求最优得分可用 $O(m)$ 滚动空间，恢复路径还需保存 $O(nm)$ 个回溯指针。

\begin{proof}[递推正确性]
对给定末端 $b$，任意长度 $i$ 的前缀均由某个以 $a$ 结尾的长度 $i-1$ 前缀和边 $(a,b)$ 组成。如果前一段不是到达 $a$ 的最优前缀，替换成更优者不会改变最后一条边，却会提高总得分，因此最优前缀必须由式~\eqref{eq:lihang-crf-viterbi-score} 给出。由起点得分为零开始归纳，再加上终止势，就得到完整路径的最高得分；回溯指针记录了每一步实际选取的前驱，因而恢复出达到该值的路径。
\end{proof}

前向算法在每个末端对全部前驱路径的权重求和，维特比算法则在对数域取最大值；二者利用相同的链式分解，却回答不同问题。此外，逐位置选择 $\arg\max_b p_w(Y_i=b\mid x)$ 一般不等于式~\eqref{eq:lihang-crf-map}：前者在允许任意标记序列作为决策时最小化逐位置错误数之和的条件期望，后者最小化整条序列的零一损失。即使使用同一模型，也应按评价损失区分这两种预测。

\subsection{教材例题与数值练习}
\label{subsec:lihang-crf-decoding-examples}

\begin{example}
\label{ex:lihang-crf-viterbi-example}
教材例 11.3。对例~\ref{ex:lihang-crf-feature-example} 的特征模型，使用维特比算法求最优标记序列。
\end{example}
\begin{solve}
把到达结点的状态得分合并到相应边上，起始得分与两个内部得分矩阵为
\[
  W_1=(1,0.5),\qquad
  W_2=\begin{pmatrix}1.4&1.5\\1.8&0.5\end{pmatrix},\qquad
  W_3=\begin{pmatrix}0.8&1.5\\1.8&0.7\end{pmatrix},
\]
并且终止得分为零。于是 $d_1=(1,0.5)$，第二步为
\[
  d_2(1)=\max\{1+1.4,\ 0.5+1.8\}=2.4,\quad \pi_2(1)=1,
\]
\[
  d_2(2)=\max\{1+1.5,\ 0.5+0.5\}=2.5,\quad \pi_2(2)=1.
\]
第三步为
\[
  d_3(1)=\max\{2.4+0.8,\ 2.5+1.8\}=4.3,\quad \pi_3(1)=2,
\]
\[
  d_3(2)=\max\{2.4+1.5,\ 2.5+0.7\}=3.9,\quad \pi_3(2)=1.
\]
因此末端为 $y_3^*=1$，依次回溯得 $y_2^*=2$、$y_1^*=1$，即 $y^*=(1,2,1)$、$S^*=4.3$。作为独立核对，按字典序列出八条路径的得分：
\[
\begin{array}{c|rrrrrrrr}
 y&111&112&121&122&211&212&221&222\\\hline
 S_w(y,x)&3.2&3.9&4.3&3.2&3.1&3.8&2.8&1.7
\end{array}
\]
这里三位数只是三元标记序列的简写。由此 $Z_w(x)=\sum_y e^{S_w(y,x)}\approx260.9850$，最优路径的条件概率为 $e^{4.3}/Z_w(x)\approx0.282391$，而例~\ref{ex:lihang-crf-feature-example} 的路径概率约为 $0.0939998$。最高得分并不意味着该路径已经具有接近 $1$ 的概率。
\end{solve}

\begin{example}
\label{ex:lihang-crf-exercise-paths}
教材习题 11.4。给定方阵
\[
  \widetilde M_1=\begin{pmatrix}0&0\\0.5&0.5\end{pmatrix},\quad
  \widetilde M_2=\begin{pmatrix}0.3&0.7\\0.7&0.3\end{pmatrix},\quad
  \widetilde M_3=\begin{pmatrix}0.5&0.5\\0.6&0.4\end{pmatrix},\quad
  \widetilde M_4=\begin{pmatrix}0&1\\0&1\end{pmatrix},
\]
固定 $\mathrm{start}=2$、$\mathrm{stop}=2$。求八条内部状态序列的概率及其中概率最大的序列。
\end{example}
\begin{solve}
只取 $\widetilde M_1$ 的第 $2$ 行及 $\widetilde M_4$ 的第 $2$ 列，得到本章约定的 $M_1=(0.5,0.5)$、$M_4=(1,1)^{\mathsf T}$，内部矩阵保持不变。矩阵乘法给出
\[
  M_1M_2=(0.5,0.5),\qquad
  M_1M_2M_3=(0.55,0.45),\qquad Z_w(x)=1.
\]
因此此例中路径权重恰好已经规范化，按 $M_1(y_1)M_2(y_1,y_2)M_3(y_2,y_3)$ 求得
\[
\begin{array}{c|rrrrrrrr}
 y&111&112&121&122&211&212&221&222\\\hline
 p_w(y\mid x)&0.075&0.075&0.210&0.140&0.175&0.175&0.090&0.060
\end{array}
\]
八项相加为 $1$，唯一最大项对应 $y^*=(1,2,1)$，其概率为 $0.21$。本例的起始行和内部矩阵恰好都满足相应行和为 $1$，所以配分函数等于 $1$；一般 CRF 的势不具备这种性质，不能据此省略规范化。
\end{solve}

本章的计算围绕同一个局部得分展开：对所有路径求和得到配分函数，固定结点或边后求和得到边缘概率，对局部特征加权求和得到学习梯度所需的期望，对路径得分取最大值则得到维特比预测。教材习题 11.1 的因子分解已在例~\ref{ex:lihang-crf-four-node-cliques} 给出，习题 11.2 的通用边界恒等式由命题~\ref{prop:lihang-crf-partition-split} 证明，习题 11.3 的学习更新为式~\eqref{eq:lihang-crf-gradient-descent}，习题 11.4 的完整计算见例~\ref{ex:lihang-crf-exercise-paths}。应同时记住三类条件：无向图定理的严格正性、迭代尺度法的非负特征假设，以及前向--后向和维特比递推中一致的首尾边界约定。


\fi

\selectchapter{12}
\ifcompilechapter
\chapter{监督学习方法总结}
\label{ch:lh-supervised-summary}

监督学习方法的差别，可以依次从四个问题理解：预测什么对象，允许使用什么模型，用什么准则评价模型，以及怎样从数据中求出模型。本章沿用教材的“适用问题、模型、学习策略、学习算法”四条线索，串联感知机、$k$ 近邻法、朴素贝叶斯法、决策树、逻辑斯谛回归与最大熵模型、支持向量机、提升方法、EM 算法、隐马尔可夫模型和条件随机场。这里把逻辑斯谛回归与最大熵模型合为一组，沿用教材的十组方法口径；其中 EM 是适用于多种模型的参数估计算法，并不单独规定一种预测模型。

本章统一以 $X,Y$ 表示随机输入与输出，以 $x,y$ 表示其取值；训练集记为 $\mathcal D=\{(x_i,y_i)\}_{i=1}^N$，需要用样本平均估计总体期望时，假定样本独立同分布。$\mathcal X,\mathcal Y$ 分别表示输入、输出空间，$\mathcal H$ 表示候选函数集合，$\mathbf1\{A\}$ 表示条件 $A$ 成立时为 $1$、否则为 $0$ 的示性函数。除特别说明外，对数均为自然对数；这一区分会影响损失函数之间的数值比较。

\section{适用问题}
\label{sec:lh-summary-tasks}

\emph{分类}（classification）预测有限集合中的类别，\emph{回归}（regression）预测实数或实向量，\emph{标注}（sequence labeling）则为输入序列预测相应的标签序列。教材本篇主要讨论分类与标注，不能由此认为监督学习只包含这两种任务。设每个位置有 $K\geq2$ 种标签，序列长度为 $n$，且不额外限制合法标签序列，则标注的输出空间为 $\mathcal Y^n$，共有 $K^n$ 个候选结果；普通分类可视为长度为 $1$ 的情形。标注的困难不仅来自类别较多，还来自相邻标签之间的依赖，逐位置选择最大概率标签未必得到概率最大的完整序列。

感知机、$k$ 近邻法（$k$-nearest neighbors，$k$-NN）、朴素贝叶斯法、决策树、逻辑斯谛回归与最大熵模型、支持向量机、提升方法都可用于分类。教材中的基本感知机、基本支持向量机和 AdaBoost 先处理二分类，多分类需要另行规定模型和策略，例如一对其余、一对一或直接定义多类目标。$k$ 近邻法、决策树与提升方法也有回归版本，但其输出和损失函数须随任务改变，不能直接照搬二分类公式。隐马尔可夫模型（hidden Markov model，HMM）与线性链条件随机场（conditional random field，CRF）主要用于序列标注；在一阶链结构、每个位置 $K$ 个状态且局部得分已知时，维特比动态规划可用 $O(nK^2)$ 次量级的运算寻找最优标签序列，因而不必枚举全部 $K^n$ 个候选。

\emph{期望极大算法}（expectation-maximization algorithm，EM）的适用条件是概率模型含有未观测变量，而不是任务必须带有类别标签。例如，朴素贝叶斯模型在类别已知时可以通过计数学习，在类别未知时可将类别作为隐变量；HMM 的状态序列已知时属于完全数据估计，未知时可用 EM 学习。两者的概率结构可以保持相同，变化的是训练数据中哪些变量可见。因此，\emph{生成或判别描述建模方式，监督或无监督描述可用数据与学习任务}，两组概念不能相互替代。

教材关于方法准确率的比较应理解为经验倾向。条件随机场可以直接利用丰富的观测特征，支持向量机与提升方法也具有较强的表示或拟合能力，但没有一种方法仅凭名称就必然优于另一种。数据规模、特征表示、噪声、正则化和评价指标都会影响结果；实际选择应在相同数据划分与指标下验证，并把训练成本、预测成本和概率解释的需要一起考虑。

\section{模型}
\label{sec:lh-summary-models}

\subsection{概率描述、决策函数与生成判别之分}
\label{subsec:lh-summary-model-types}

\emph{概率模型}为输出指定概率分布，例如 $p_\theta(y\mid x)=P_\theta(Y=y\mid X=x)$；$\theta$ 表示参数。\emph{决策函数}则直接给出预测 $d(x)\in\mathcal Y$。两者并不互斥：在有限类别、误判代价相等的条件下，可以从条件分布导出
\begin{equation}
d_\theta(x)\in\operatorname*{arg\,max}_{y\in\mathcal Y}p_\theta(y\mid x),
\label{eq:lh-summary-bayes-decision}
\end{equation}
并为并列最大值固定一种选择规则。若误判代价不同，应最小化条件期望代价，最大概率规则便未必适用。反过来，只知道 $d(x)$ 或一个实值得分，通常不能唯一恢复条件概率；支持向量机的间隔、提升分类器的加权得分都不能未经建模或校准就当作概率。

\begin{definition}[生成模型与判别模型]
\label{def:lh-summary-generative-discriminative}
在带类别或标签的预测问题中，\emph{生成模型}（generative model）对联合分布 $p_\theta(x,y)$ 建模，再在边缘概率或密度为正的输入处，通过
\begin{equation}
p_\theta(y\mid x)=\frac{p_\theta(x,y)}{\sum_{u\in\mathcal Y}p_\theta(x,u)}
\label{eq:lh-summary-generative-conditional}
\end{equation}
获得条件分布；此式针对离散输出，连续输出时将求和换为相应积分。\emph{判别模型}（discriminative model）直接对条件分布 $p_\theta(y\mid x)$ 或决策函数建模，其学习不要求同时拟合输入的边缘分布。
\end{definition}

朴素贝叶斯法与 HMM 属于生成方法；感知机、$k$ 近邻法、决策树、逻辑斯谛回归与最大熵模型、支持向量机、提升方法和 CRF 属于判别方法。逻辑斯谛回归、最大熵模型与 CRF 本身有明确的条件概率解释，也可以通过式~\eqref{eq:lh-summary-bayes-decision} 输出类别或序列。分类树的叶节点既可保存多数类，也可保存经验类别分布；$k$ 近邻法的邻域类别比例也可用作局部概率估计。这些不同输出方式不改变它们是否学习联合分布的区别。

生成与判别、分类与标注是两个相互独立的比较方向，表~\ref{tab:lh-summary-model-relations} 用教材讨论的四种概率模型说明这一点。同一行比较联合建模与条件建模，同一列比较分类与序列标注；这些对应关系并不意味着生成模型与判别模型具有相同的参数约束或训练目标。
\begin{table}[htbp]
\centering
\caption{分类与标注中的生成模型和判别模型}
\label{tab:lh-summary-model-relations}
\begin{tabularx}{\textwidth}{lYY}
\toprule
预测任务 & 生成模型：联合分布 & 判别模型：条件分布 \\
\midrule
分类 & 朴素贝叶斯法 $p(x,y)$ & 逻辑斯谛回归、最大熵模型 $p(y\mid x)$ \\
序列标注 & HMM $p(x_{1:n},y_{1:n})$ & 线性链 CRF $p(y_{1:n}\mid x_{1:n})$ \\
\bottomrule
\end{tabularx}
\end{table}

以二分类朴素贝叶斯为例，设 $Y\in\{-1,+1\}$，$X=(X_1,\ldots,X_d)$，并假定各输入分量在给定 $Y$ 后条件独立。在所涉及的先验与条件概率或密度均为正时，
\begin{equation}
\log\frac{p(Y=+1\mid x)}{p(Y=-1\mid x)}
=\log\frac{p(Y=+1)}{p(Y=-1)}
+\sum_{j=1}^d\log\frac{p(x_j\mid Y=+1)}{p(x_j\mid Y=-1)}.
\label{eq:lh-summary-naive-bayes-odds}
\end{equation}
右端是各分量贡献之和。对于离散特征，可用各取值的示性函数把它写成对数线性形式，其系数由满足归一化条件的生成模型参数导出。逻辑斯谛回归则直接估计条件模型的系数，不要求上述条件独立假设。因此，相似的判别边界表达式不代表相同的建模假设或估计目标；比较模型族是否等价，还须检查具体的特征与参数化。

在序列情形，记 $x_{1:n}=(x_1,\ldots,x_n)$、$y_{1:n}=(y_1,\ldots,y_n)$。一阶 HMM 的联合分布写为
\begin{equation}
p_\theta(x_{1:n},y_{1:n})
=\pi_{y_1}\prod_{t=2}^n a_{y_{t-1},y_t}\prod_{t=1}^n b_{y_t}(x_t),
\label{eq:lh-summary-hmm-factorization}
\end{equation}
其中 $\pi_k$ 为初始状态概率，$a_{jk}$ 为转移概率，$b_k(x)$ 为发射概率或密度。它要求状态满足一阶马尔可夫性，且给定状态序列后各观测相互独立、每个观测只依赖相应状态。线性链 CRF 则直接设
\begin{equation}
p_\theta(y_{1:n}\mid x_{1:n})
=\frac{\exp\!\bigl(\theta^{\mathsf T}\Phi(x_{1:n},y_{1:n})\bigr)}
{\sum_{u_{1:n}}\exp\!\bigl(\theta^{\mathsf T}\Phi(x_{1:n},u_{1:n})\bigr)},
\label{eq:lh-summary-crf-model}
\end{equation}
其中 $\Phi$ 是由相邻标签与观测序列的局部特征求和得到的固定特征向量，分母对合法标签序列求和。它可以让局部特征依赖整段已知观测，不要求观测满足 HMM 的生成假设；但标签之间仍须保持指定的链结构，才能直接使用相应动态规划。长度为 $1$ 时，上式具有分类最大熵模型的形式。

\subsection{特征空间与线性、非线性的含义}
\label{subsec:lh-summary-representations}

感知机与线性支持向量机通常在欧氏空间中使用仿射得分 $s(x)=w^{\mathsf T}x+b$。逻辑斯谛回归将这一得分映射为概率；更一般的\emph{对数线性模型}（log-linear model）写为
\begin{equation}
p_\theta(y\mid x)=\frac{\exp(\theta^{\mathsf T}\Phi(x,y))}{Z_\theta(x)},
\qquad Z_\theta(x)=\sum_{u\in\mathcal Y}\exp(\theta^{\mathsf T}\Phi(x,u)).
\label{eq:lh-summary-log-linear}
\end{equation}
这里对数概率关于参数的表达式含有归一化项，准确地说，是未归一化的对数得分对参数线性。概率本身通常非线性；若采用非线性特征，分类边界在原输入空间也可以非线性。因此，“线性”必须说明是相对于原始输入、特征表示，还是模型参数而言。

核支持向量机以 $K(x,x')=\langle\phi(x),\phi(x')\rangle$ 隐式使用希尔伯特空间中的特征映射，其得分在该特征空间中线性，含偏置时为仿射，在原空间中可以非线性；内积、范数或距离的计算须与所用核和空间一致。$k$ 近邻法依靠选定距离进行局部比较，通常用于欧氏空间，也可用于具有适当距离的一般对象，并不在定义上要求一定存在内积。决策树按特征与切分规则划分输入空间，可以处理数值或离散属性，在叶区域内输出常数或类别分布，通常形成分段的非线性预测。

提升方法使用\emph{加法模型}（additive model）$F_M(x)=\sum_{m=1}^M\alpha_m h_m(x)$，其中 $h_m$ 为基学习器、$\alpha_m$ 为系数。它对给定基学习器的系数是线性的，但整体函数可以随基学习器而非线性。若所有 $h_m$ 都是同一特征表示下的仿射实值函数，其线性组合仍为仿射函数；若 $h_m$ 是树或阈值分类器，则不再有此结论。模型的输入类型由基学习器支持的表示决定，组合本身不会自动消除基学习器的限制。

\section{学习策略}
\label{sec:lh-summary-strategies}

\subsection{从分类错误到代理损失与正则化}
\label{subsec:lh-summary-surrogate-losses}

对二分类 $Y\in\{-1,+1\}$，设 $f(x)\in\mathbb R$ 为分类得分，以其符号预测类别。定义\emph{函数间隔}（functional margin）$m=yf(x)$；正间隔表示分类正确，负间隔表示分类错误。为统一比较，下文把零间隔也记为错误，使用 $\ell_{01}(m)=\mathbf1\{m\leq0\}$；若实际预测在零分处固定选择某一类别，其真正的 0--1 损失不会超过这个保守的记法。0--1 损失不连续且不凸，直接优化往往困难，因此使用较易计算的\emph{代理损失}（surrogate loss）：
\begin{align}
\ell_{\mathrm h}(m)&=[1-m]_+=\max\{0,1-m\},
&&\text{合页损失（hinge loss）},\label{eq:lh-summary-hinge}\\
\ell_{\log}(m)&=\log(1+e^{-m}),
&&\text{逻辑斯谛损失（logistic loss）},\label{eq:lh-summary-logistic-loss}\\
\ell_{\exp}(m)&=e^{-m},
&&\text{指数损失（exponential loss）}.\label{eq:lh-summary-exp-loss}
\end{align}
三者均为间隔的非增凸函数。合页损失在 $m\geq1$ 时为零，并在 $m=1$ 处不可微；其余两者光滑且在有限间隔处严格为正。指数损失对很大的负间隔增长最快，因此被严重误分的样本可能产生很大影响；这解释了其对某些噪声较敏感的原因，但并非对任意数据的性能结论。

\begin{proposition}[代理损失与分类错误的上界]
\label{prop:lh-summary-loss-bounds}
对任意实数 $m$，有
\begin{equation}
\ell_{01}(m)\leq\ell_{\mathrm h}(m),\qquad
\ell_{01}(m)\leq\ell_{\exp}(m),\qquad
\ell_{01}(m)\leq\frac{\ell_{\log}(m)}{\log2}.
\label{eq:lh-summary-loss-bounds}
\end{equation}
自然对数形式的 $\ell_{\log}$ 本身并不总是 $\ell_{01}$ 的上界。
\end{proposition}
\begin{proof}
当 $m>0$ 时，左端为零，三个右端均非负。当 $m\leq0$ 时，$1-m\geq1$、$e^{-m}\geq1$，且 $\log(1+e^{-m})\geq\log2$，于是得到三个不等式。若取 $m=-0.1$，则分类错误，但 $\ell_{\log}(-0.1)=\log(1+e^{0.1})\approx0.7444<1$，故最后一个结论也成立。
\end{proof}

因此，教材把三种损失都称为 0--1 损失上界时，需要为逻辑斯谛损失补上底数或尺度约定：以 $2$ 为底的对数可直接给出上界；采用自然对数时用式~\eqref{eq:lh-summary-loss-bounds}。损失乘以正数不会改变无正则化问题的最优解，但若存在正则化项，保持同一优化问题还须让正则化系数乘以相同常数。代理损失之间的上界关系也不意味着其训练最优解必然同时最小化有限样本上的 0--1 损失。

给定损失 $L$，\emph{总体风险}为 $R(f)=\mathbb E[L(Y,f(X))]$，假定相应期望存在；训练数据给出的\emph{经验风险}为 $\widehat R_N(f)=N^{-1}\sum_{i=1}^N L(y_i,f(x_i))$。前者依赖未知数据分布，后者可以计算。教材所用的正则化风险形式为
\begin{equation}
\min_{f\in\mathcal H}\left\{
\frac1N\sum_{i=1}^N L(y_i,f(x_i))+\lambda J(f)\right\},
\qquad\lambda\geq0,
\label{eq:lh-summary-regularized-risk}
\end{equation}
其中 $J(f)$ 衡量或约束模型复杂度。令 $\lambda=0$ 得到经验风险最小化；$\lambda>0$ 时在拟合与复杂度之间权衡。这里的正则化目标表达的是学习策略，优化算法如何求解以及预测能否泛化，仍是另外两个问题。

对软间隔线性支持向量机，一种常用参数形式为
\begin{equation}
\min_{w,b}\left\{\frac12\|w\|_2^2
+C\sum_{i=1}^N[1-y_i(w^{\mathsf T}x_i+b)]_+\right\},\qquad C>0.
\label{eq:lh-summary-svm-objective}
\end{equation}
这里惩罚的是权向量的平方范数，通常不惩罚偏置；换成经验平均形式时须相应调整系数。逻辑斯谛回归和最大熵模型常用负对数似然，也可以加入二次正则化。基本 AdaBoost 以指数损失进行前向分步拟合，没有必须显式出现的 $\lambda J$，但可用轮数、基学习器复杂度和缩减步长控制拟合程度。\emph{早停止}（early stopping）通常依据验证结果选择迭代轮数，属于隐式控制复杂度的做法，不能未经证明就等同于某个固定的范数惩罚。

\subsection{似然、后验与树的学习准则}
\label{subsec:lh-summary-likelihood}

对条件概率模型，\emph{极大似然估计}（maximum likelihood estimation，MLE）在给定输入后最大化标签的条件似然，等价于
\begin{equation}
\min_\theta\left\{-\frac1N\sum_{i=1}^N\log p_\theta(y_i\mid x_i)\right\}.
\label{eq:lh-summary-conditional-mle}
\end{equation}
该式要求给定输入后的标签按样本分解；独立同分布抽样是充分条件。生成模型的完全数据 MLE 则使用 $p_\theta(x_i,y_i)$，并同时拟合输入与输出的联合分布。两种似然的目标通常不同，即使最后都通过条件分布做预测，也不能把训练目标直接互换。

对于 $y\in\{-1,+1\}$，若 $p_\theta(y\mid x)=1/(1+e^{-yf_\theta(x)})$，则
\[
-\log p_\theta(y\mid x)=\log(1+e^{-yf_\theta(x)}),
\]
所以二分类逻辑斯谛损失恰为条件负对数似然。对一般的对数线性分类模型或 CRF，式~\eqref{eq:lh-summary-log-linear} 给出
\begin{equation}
-\log p_\theta(y\mid x)=\log Z_\theta(x)-\theta^{\mathsf T}\Phi(x,y).
\label{eq:lh-summary-loglinear-nll}
\end{equation}
它是多类或结构化输出的对数损失；只有特定二分类参数化才化为单个标量间隔的 $\log(1+e^{-m})$。CRF 的归一化因子对整个输出序列求和，不能用逐位置独立的二分类损失替代。

\emph{最大后验估计}（maximum a posteriori estimation，MAP）加入参数先验密度 $\pi(\theta)$，在后验可定义且最大值存在的条件下，等价于最小化
\begin{equation}
-\sum_{i=1}^N\log p_\theta(y_i\mid x_i)-\log\pi(\theta).
\label{eq:lh-summary-map-objective}
\end{equation}
因此先验负对数可以解释为正则化项。若 $\theta\in\mathbb R^q$ 且先验为均值零、协方差 $\tau^2 I_q$ 的正态分布，$\tau>0$、$I_q$ 为 $q$ 阶单位矩阵，则除去与参数无关的常数，并把目标除以 $N$，得到
\begin{equation}
-\frac1N\sum_{i=1}^N\log p_\theta(y_i\mid x_i)
+\frac{\|\theta\|_2^2}{2N\tau^2}.
\label{eq:lh-summary-gaussian-map}
\end{equation}
若把 $J(\theta)$ 定义为 $\|\theta\|_2^2/2$，此处 $\lambda=1/(N\tau^2)$。系数随使用似然总和还是平均而改变，且这里只对真正赋予该正态先验的参数施加相应惩罚。MAP 给出参数点估计，并不等同于对参数后验积分得到的贝叶斯预测。

分类树在使用叶节点类别分布与对数损失时，也有似然解释。设树 $\mathcal T$ 有叶节点 $\ell$，其中第 $k$ 类有 $n_{\ell k}$ 个样本，$n_\ell=\sum_k n_{\ell k}>0$。在每个叶内用频率 $\widehat p_{\ell k}=n_{\ell k}/n_\ell$ 估计类别概率，完全数据负对数似然为
\begin{equation}
-\sum_\ell\sum_k n_{\ell k}\log\widehat p_{\ell k}
=\sum_\ell n_\ell H(\widehat p_\ell),
\qquad H(p)=-\sum_k p_k\log p_k,
\label{eq:lh-summary-tree-entropy}
\end{equation}
其中约定 $0\log0=0$。再加叶节点数惩罚 $\alpha|\mathcal L(\mathcal T)|$，$\alpha\geq0$、$\mathcal L(\mathcal T)$ 为叶集合，就得到一种正则化似然准则。但树结构的搜索仍是离散组合问题，贪心切分不保证得到其全局最优树；采用基尼指数的分类 CART 或平方损失的回归 CART，也不能不加区分地说成使用了上述对数似然损失。

当训练变量不完全可见时，似然须先边缘化隐变量。例如对独立样本 $(x_i,z_i)$，仅观测 $x_i$，观测数据对数似然为 $\sum_i\log\sum_{z_i}p_\theta(x_i,z_i)$，而不是把未知 $z_i$ 当作已经观测到的固定标签。内层求和使优化结构发生变化，这正是下一节讨论 EM 的原因。

\section{学习算法}
\label{sec:lh-summary-algorithms}

\subsection{解析估计、数值优化与结构搜索}
\label{subsec:lh-summary-optimization-methods}

模型与准则确定后，学习通常归结为估计参数、搜索结构或保存供预测使用的数据。完全观测的离散朴素贝叶斯模型，其 MLE 可以由类别频数与类内特征频数归一化得到；完全观测的离散 HMM，同样可以由初始状态、状态转移和发射事件的计数估计相应概率。这里需要被归一化的计数总和为正；未出现的条件下参数未必由数据唯一确定，平滑还会改变估计公式。若 HMM 的状态不可见，就不能直接使用状态计数，须改为后验期望计数等方法。维特比算法用于给定模型后的最优路径预测，Baum--Welch 算法是用于参数学习的 EM，两者解决的问题不同。

$k$ 近邻法通常没有一个必须预先求出的有限维参数最优解，训练主要保存样本并可建立近邻搜索索引，预测时根据距离和投票规则计算输出；$k$、距离及权重规则可用验证数据选择。感知机按误分类样本更新权值，其经典有限步收敛结论要求有限训练集线性可分等条件；不可分时不能据此保证算法停止。逻辑斯谛回归、最大熵模型与固定特征的 CRF 则常用梯度法、拟牛顿法等数值优化方法，算法输出达到给定容差的近似参数，而不是自动得到某个解析公式。

对正半定核与标准软间隔准则，支持向量机可以通过凸二次规划学习；\emph{序列最小最优化}（sequential minimal optimization，SMO）每次选择少量对偶变量，在维持约束的条件下优化子问题。决策树采用局部特征选择、递归生成及剪枝，属于利用问题结构的搜索方法。AdaBoost 则在已有加法模型上逐轮选择基学习器和系数，使当前指数经验损失下降；它并非每一轮都重新联合优化全部基学习器。弱学习器集合的限制与前向选择方式意味着，损失关于函数值凸，也不能直接推出这一训练过程已找到整个候选模型集合的全局最优解。

\subsection{EM 的单调性与凸优化的边界}
\label{subsec:lh-summary-guarantees}

\begin{proposition}[EM 的似然单调性]
\label{prop:lh-summary-em-monotonicity}
设观测数据为 $D$，隐变量 $Z$ 取有限多个值，模型在所讨论参数处满足 $p_\theta(D,z)>0$。记观测对数似然为 $\ell(\theta)=\log\sum_z p_\theta(D,z)$，在当前参数 $\theta^{(r)}$ 下定义
\[
q_r(z)=p_{\theta^{(r)}}(z\mid D),\qquad
Q(\theta\mid\theta^{(r)})=\sum_z q_r(z)\log p_\theta(D,z).
\]
若更新满足 $Q(\theta^{(r+1)}\mid\theta^{(r)})\geq Q(\theta^{(r)}\mid\theta^{(r)})$，则 $\ell(\theta^{(r+1)})\geq\ell(\theta^{(r)})$。精确最大化 $Q$ 的 M 步是这一条件的充分条件。
\end{proposition}
\begin{proof}
定义熵 $H(q_r)=-\sum_zq_r(z)\log q_r(z)$，以及相对熵
$D_{\mathrm{KL}}(q_r\|p_\theta)=\sum_zq_r(z)\log(q_r(z)/p_\theta(z\mid D))$。由 $p_\theta(D,z)=p_\theta(z\mid D)p_\theta(D)$，有恒等式
\[
\ell(\theta)=Q(\theta\mid\theta^{(r)})+H(q_r)
+D_{\mathrm{KL}}(q_r\|p_\theta).
\]
相对熵非负：对正概率分布使用 $\log u\leq u-1$，可得
$-D_{\mathrm{KL}}(q\|p)=\sum_zq(z)\log(p(z)/q(z))\leq\sum_z(p(z)-q(z))=0$。
在 $\theta=\theta^{(r)}$ 时，相对熵为零。将两个参数处的恒等式相减，便得
\[
\ell(\theta^{(r+1)})-\ell(\theta^{(r)})
\geq Q(\theta^{(r+1)}\mid\theta^{(r)})-Q(\theta^{(r)}\mid\theta^{(r)})\geq0.\qedhere
\]
\end{proof}

EM 的输入是含隐变量的概率模型与观测数据，初始化选择使后验可定义的参数；E 步计算当前隐变量后验，M 步提高上述 $Q$，然后迭代。实现时可按相邻对数似然差小于容差或达到迭代上限终止，输出当前参数，并检查 M 步是否真正满足所需改进条件。命题~\ref{prop:lh-summary-em-monotonicity} 在有限正概率情形下给出证明；连续隐变量或允许零概率时，需要另外检查积分、支持及有限性条件。若似然值还存在有限上界，单调性保证似然值序列收敛，但\emph[caution]{不单凭这一点保证参数序列收敛，也不保证收敛到全局最大点}。驻点性质需要进一步的正则条件，实际结果还可能依赖初始化。

固定特征的条件对数线性模型具有更强的优化结构。设参数可取整个 $\mathbb R^q$，候选输出集合有限且特征有限，则 $\log Z_\theta(x)$ 的梯度为模型下的特征期望，海森矩阵为
\begin{equation}
\nabla_\theta^2\log Z_\theta(x)
=\operatorname{Cov}_\theta\bigl(\Phi(x,Y)\mid X=x\bigr)\succeq0.
\label{eq:lh-summary-logpartition-hessian}
\end{equation}
这是因为对有限和逐项求导，二阶导数等于特征外积的期望减去均值的外积；任意向量 $v$ 对应的二次型是 $v^{\mathsf T}\Phi(x,Y)$ 的条件方差，故非负。因此式~\eqref{eq:lh-summary-loglinear-nll} 关于参数凸，加入凸正则化后仍凸。逻辑斯谛回归、最大熵模型与无隐标签的固定特征 CRF 属于这一情形；标准支持向量机也有凸目标与凸约束。但改变为含隐变量的模型或同时学习非线性特征参数后，不能直接沿用这一结论。

\begin{example}
\label{ex:lh-summary-convex-nonattainment}
取两条二分类样本 $(x_1,y_1)=(-1,-1)$、$(x_2,y_2)=(1,+1)$，使用无偏置得分 $f_w(x)=wx$，$w\in\mathbb R$。判断无正则化逻辑斯谛经验损失是否存在有限最优参数。
\end{example}
\begin{solve}
两条样本的间隔都为 $w$，故平均损失 $J(w)=\log(1+e^{-w})$。它满足
\[
J'(w)=-\frac1{1+e^w}<0,\qquad
J''(w)=\frac{e^w}{(1+e^w)^2}>0.
\]
因此 $J$ 严格凸且处处递减，有限 $w$ 时总为正，而 $w\to+\infty$ 时趋于零。其下确界为零，却不存在有限最小点。若改为 $J(w)+\lambda w^2/2$、$\lambda>0$，目标连续、强凸且在 $|w|\to\infty$ 时趋于无穷，于是存在唯一有限最小点。这说明凸性、最小值能否达到与最优参数是否唯一必须分别判断。
\end{solve}

由例~\ref{ex:lh-summary-convex-nonattainment} 可知，教材“凸优化问题的全局最优解保证存在”的表述需要补上存在性条件。凸性保证凸可行集上已有局部最小点是全局最小点；对无约束可微凸函数，已有驻点也是全局最小点。连续目标在非空紧可行集上可以达到最小值；在非空有限维闭可行集上，若连续目标在可行变量的范数趋于无穷时也趋于无穷，则可由有界子水平集保证存在性。对固定特征的非负负对数似然，惩罚全部参数的正系数平方范数给出一种充分条件；若有未惩罚方向，则须另行检查。硬间隔支持向量机还要求训练数据可分以保证可行。其他方法也不宜统一说成“都不是凸问题”：例如完全观测的离散概率模型在概率单纯形上有凸负对数似然，而 $k$ 近邻法通常不以同一类参数优化问题来定义。

表~\ref{tab:lh-summary-ten-methods} 按教材涉及的基本版本重新归纳十组方法。它同时列出任务、模型、策略与算法，便于区分“选了什么函数”“以什么为好”以及“怎样求解”；不同变体的损失、约束和算法仍须单独说明。
\begin{longtable}{L{0.15\textwidth}L{0.25\textwidth}L{0.27\textwidth}L{0.21\textwidth}}
\caption{十组学习方法的模型、准则与主要学习方式}\label{tab:lh-summary-ten-methods}\\
\toprule
方法 & 任务与模型特点 & 典型学习准则 & 主要学习方式 \\
\midrule
\endfirsthead
\toprule
方法 & 任务与模型特点 & 典型学习准则 & 主要学习方式 \\
\midrule
\endhead
\bottomrule
\endfoot
感知机 & 基本版本为二分类；线性判别函数 & 按误分类点的负间隔构造感知机准则，寻找分离面 & 误分类驱动的迭代更新；可分时有经典收敛保证 \\
$k$ 近邻法 & 分类；基于距离与邻域投票的局部判别 & 基本形式不显式求解统一的参数风险最小化 & 保存样本、搜索近邻；验证选择 $k$ 等超参数 \\
朴素贝叶斯法 & 分类；条件独立假设下的生成概率模型 & 完全数据极大似然或带先验的估计 & 离散情形用频数归一化；可加平滑 \\
决策树 & 分类；递归区域划分与叶输出 & 熵、基尼指数等局部准则，并以剪枝控制复杂度 & 贪心特征选择、生成与剪枝 \\
逻辑斯谛回归与最大熵模型 & 分类；条件对数线性模型 & 条件负对数似然，可加凸正则化 & 梯度法、拟牛顿法；最大熵模型也可用迭代尺度法 \\
支持向量机 & 基本版本为二分类；线性或核判别函数 & 合页损失与平方范数惩罚，或硬间隔约束 & 凸二次规划、SMO 等 \\
提升方法 & 分类；基学习器的加法组合 & 基本 AdaBoost 使用指数损失 & 前向分步拟合；控制轮数与基学习器复杂度 \\
EM 算法 & 含隐变量概率模型的参数估计；无独立的预测模型 & 提高观测数据似然，或相应的带先验目标 & 后验期望的 E 步与提高辅助目标的 M 步 \\
隐马尔可夫模型 & 标注；马尔可夫状态与条件独立观测的联合模型 & 完全数据或观测数据似然 & 状态已知时计数；未知时可用 Baum--Welch \\
条件随机场 & 标注；固定特征下的条件对数线性模型 & 结构化条件负对数似然，可加正则化 & 梯度法、拟牛顿法；链上用动态规划算所需期望 \\
\end{longtable}

% 教材来源：机器学习/textbooks/机器学习（分章原文）/机器学习-13.pdf。
% 对应教材第 13 章，原书第 245--254 页；本文件沿用自动章节编号。
% 原文部分图表的文本层不完整，以下不转录其数值，计算例题均另行构造。

\fi

\selectchapter{13}
\ifcompilechapter
\chapter{无监督学习概论}
\label{ch:lh-unsupervised-overview}

监督学习借助输入与标签之间的对应关系学习预测规则；当数据中只有观测样本，没有预先给定的目标标签时，仍然可以研究样本如何分组、哪些特征可以压缩、数据可能由什么机制生成。\emph{无监督学习}（unsupervised learning）正是从无标注数据中学习统计规律或潜在结构的机器学习问题。本章对应教材第13章，依次讨论基本原理、基本问题、机器学习三要素和典型方法，为后续的聚类、矩阵分解、概率话题模型与图分析建立共同语言。

本章需要向量与矩阵、条件概率和基本微分知识；主成分分析中的正交投影和图分析中的概率转移在使用时说明。以下以 $X$ 表示观测随机向量，$x$ 表示其取值；$Z$ 表示潜在变量，$z$ 表示其取值；$\theta$ 表示待学习的参数。离散事件的概率记为 $P$，连续变量的密度或统一书写的概率模型记为 $p$。本章的计算例题另行构造，用于展开教材中的基本思想。

\section{无监督学习基本原理}
\label{sec:lh-unsup-principles}

设无标注训练集为 $U=\{x_1,\ldots,x_N\}$，其中 $x_j\in\mathbb R^M$。沿用教材按列存放样本的约定，将数据写成
\begin{equation}
\mathbf X=(x_1,\ldots,x_N)
=\begin{pmatrix}
x_{11}&\cdots&x_{1N}\\
\vdots&&\vdots\\
x_{M1}&\cdots&x_{MN}
\end{pmatrix}\in\mathbb R^{M\times N}.
\label{eq:lh-unsup-data-matrix}
\end{equation}
这里 $x_{ij}$ 是第 $j$ 个样本的第 $i$ 个特征，行对应特征，列对应样本。数据中没有附带的监督标签 $y_j$，并不意味着学习过程没有目标函数，也不意味着模型不含类别变量；类别可以作为需要从数据中推断的潜在变量。

\begin{definition}[无监督学习的模型形式]
\label{def:lh-unsup-models}
无监督学习根据数据选择参数及必要的模型结构，其分析结果可以用确定性映射 $z=g_\theta(x)$、条件分布 $p_\theta(z\mid x)$，或生成方向的条件分布 $p_\theta(x\mid z)$ 表示。映射给出类别或低维表示，$p_\theta(z\mid x)$ 表示给定观测后的潜在状态不确定性；若用 $p_\theta(x\mid z)$ 描述数据生成，还需要指定潜在变量的分布，才能确定完整的观测模型。
\end{definition}

教材用“压缩”概括发现结构的共同思路：聚类用较少的类别概括大量样本，降维用较少的坐标概括高维特征，概率模型用共享的结构与参数概括样本的分布。按式~\eqref{eq:lh-unsup-data-matrix} 的约定，这分别对应样本之间的结构、特征之间的结构，以及两者的联合结构；教材称其为纵向、横向及纵横向结构。这是一种组织问题的直观方式，而不是对所有算法的严格分类，例如协同聚类（co-clustering）会同时对样本和特征分组。

\emph[caution]{压缩所得的结构取决于模型、损失和数据表示。}若不限制表示的维数、类别数或模型复杂度，逐点记忆样本也可能取得很小的训练损失，却没有解释或推广能力；反过来，过度压缩又会抹去有用差异。因此，“损失较小的压缩揭示本质结构”应理解为建模动机，不能视为保证发现真实语义类别的定理。距离的选择、特征尺度、复杂度限制和领域知识共同决定哪些差异被保留。

无监督学习既可直接用于已有数据的探索、可视化和结构分析，也可作为监督学习的前处理，例如先学到低维表示，再训练分类器。若模型还定义了新样本的映射或概率分布，就可以用于新数据的表示、归类或概率评价；但某些方法只给出训练样本上的分组或坐标，不能未经额外定义就当作适用于任意输入的预测函数。没有外部标签使评价更依赖目标与假设，增加数据量也不能代替对这些假设的检查。

\section{基本问题}
\label{sec:lh-unsup-problems}

\subsection{聚类}
\label{subsec:lh-unsup-clustering}

\emph{聚类}（clustering）依据给定的相似性或距离，将样本组织成若干组，使同组样本在所选标准下较相似。类别不是训练数据附带的答案，而是学习结果；类别数 $K$ 可以预先指定，也可以按模型选择规则决定。对欧氏数据，距离常由特征向量计算，但“距离近”并不自动等同于某种任务中的“语义相同”。

\begin{definition}[硬聚类与概率软聚类]
\label{def:lh-unsup-hard-soft}
给定 $K$ 个类别，\emph{硬聚类}（hard clustering）将每个样本分配给唯一类别 $z_j\in\{1,\ldots,K\}$。若要求各类非空，则得到样本索引集的划分 $\mathcal C_1,\ldots,\mathcal C_K$，满足
\[
\mathcal C_k\ne\varnothing,\qquad
\mathcal C_k\cap\mathcal C_\ell=\varnothing\ (k\ne\ell),\qquad
\bigcup_{k=1}^K\mathcal C_k=\{1,\ldots,N\}.
\]
本章用\emph{软聚类}（soft clustering）主要指概率形式：对第 $j$ 个样本输出隶属概率
\begin{equation}
r_{kj}=P_\theta(Z=k\mid X=x_j),\qquad
r_{kj}\geq0,\qquad \sum_{k=1}^Kr_{kj}=1.
\label{eq:lh-unsup-responsibility}
\end{equation}
它描述类别的不确定性，而不是同时断言样本具有多个已知的真实标签。取 $\arg\max_k r_{kj}$ 并指定平局规则，可以进一步得到硬分配。
\end{definition}

硬聚类可在模型定义了新样本规则时写成 $z=g_\theta(x)$，软聚类则写成条件分布 $P_\theta(Z\mid X=x)$。聚类与监督分类虽然都可能输出类别编号，但前者没有训练标签用于规定编号的语义；同时置换全部类别编号不会改变分组本身。层次聚类和 $K$ 均值聚类（$K$-means clustering）是硬聚类的典型方法，利用高斯混合模型后验概率进行分配则属于软聚类。

\subsection{降维}
\label{subsec:lh-unsup-dimension-reduction}

\begin{definition}[降维]
\label{def:lh-unsup-dimension-reduction}
\emph{降维}（dimensionality reduction）将高维样本 $x\in\mathbb R^M$ 转换为较低维表示 $z\in\mathbb R^d$，其中 $d<M$，通常写成
\begin{equation}
z=g_\theta(x),\qquad
\mathbf Z=(z_1,\ldots,z_N)\in\mathbb R^{d\times N}.
\label{eq:lh-unsup-representation}
\end{equation}
学习的目标是在所选评价标准下，尽可能保留原数据中有用的信息或关系。$g_\theta$ 可以是线性映射，也可以是非线性映射；若先中心化再作线性投影，对原坐标而言得到的是仿射映射。
\end{definition}

降维常基于数据近似位于低维线性子空间或低维流形（manifold）的假设：样本的坐标很多，但独立变化的自由度较少。这个假设需要由数据与任务支持；任意高维数据都能被映射到低维空间，却不能因此保证信息没有损失。所谓保留信息可以具体指重构误差小、保留较大方差、保持局部邻域，或近似保持某类距离；这些目标一般并不等价。

降维得到的坐标通常是多个原始特征的组合，因而不同于只挑选部分原坐标的特征选择。它也不同于聚类：降维输出连续表示，聚类输出分组或分组概率。低维图上出现分离的点群可以为后续分析提供线索，但不构成原数据必然含有相应真实类别的证明。

\subsection{概率模型估计}
\label{subsec:lh-unsup-probability-estimation}

\emph{概率模型估计}（probability model estimation）假设观测数据来自某个概率模型，在预先给定的候选模型族内学习其参数，必要时还要选择具体结构。为区分结构与样本的潜在状态，记 $s$ 为全局模型结构，$\theta$ 为该结构下的参数，$Z$ 为单个样本的潜在变量。对离散潜在变量，生成模型满足
\begin{equation}
p_{\theta,s}(x,z)=P_{\theta,s}(Z=z)\,p_{\theta,s}(x\mid z),\qquad
p_{\theta,s}(x)=\sum_zP_{\theta,s}(Z=z)\,p_{\theta,s}(x\mid z).
\label{eq:lh-unsup-latent-model}
\end{equation}
若 $Z$ 连续，则相应地对潜在状态积分。观测数据的概率或密度来自边缘化，不能只指定一个条件分布 $p_{\theta,s}(x\mid z)$ 而省略潜在状态的分布。

教材在概率模型估计的概述中还用隐式结构表示混合成分数或图结构。这里将这类模型层面的选择统一放入 $s$，避免与“某个样本属于哪个成分”的潜在变量 $Z$ 混淆。例如，高斯混合模型的结构可以包含成分数 $K$，潜在变量取值 $Z=k$ 表示样本来自第 $k$ 个成分，参数则包括混合权重、各成分均值与协方差。在本章的频率学派估计中，$\theta$ 是未知但固定的量；只有另行引入参数先验时，才对其作贝叶斯随机建模。

对固定结构 $s$，假设观测 $x_1,\ldots,x_N$ 来自独立同分布抽样，极大似然策略为
\begin{equation}
\widehat\theta_s\in\underset{\theta\in\Theta_s}{\arg\max}\;
\ell_s(\theta),\qquad
\ell_s(\theta)=\sum_{j=1}^N\log p_{\theta,s}(x_j),
\label{eq:lh-unsup-likelihood}
\end{equation}
其中 $\Theta_s$ 是容许的参数集合，上式以最大值存在为前提。模型可以是混合模型，也可以是用有向图或无向图表达依赖关系的概率图模型；学到模型后，边缘分布与条件分布在数学上随之确定，但高维模型中的具体推断仍可能需要近似计算。若没有潜在变量，则直接估计 $p_\theta(x)$，例如估计一个高斯分布的均值和协方差。

\emph[caution]{模型结构选择与固定结构下的参数估计是不同问题。}增加成分或参数通常会提高模型拟合数据的能力，不能据此直接认定结构更好；还需要复杂度限制、正则化、先验或适当的验证准则。即使结构固定，有限的极大似然解也未必存在。例如，对至少含两个成分且不限制协方差的高斯混合模型，可以保持一个正权重成分对全部观测给出正密度，同时令另一个正权重成分的均值等于某个观测点、协方差趋于零；该点的密度趋于无穷，其他点的密度仍有正下界，故观测对数似然无上界。因此，写下极大化目标并不等于已经证明估计问题有良好解。

概率估计与软聚类通过贝叶斯公式联系起来。固定结构并省略下标 $s$，设 $\pi_k=P_\theta(Z=k)$、$p_k(x)=p_\theta(x\mid Z=k)$，则在分母为正时
\begin{equation}
P_\theta(Z=k\mid X=x)
=\frac{\pi_kp_k(x)}{\sum_{\ell=1}^K\pi_\ell p_\ell(x)}.
\label{eq:lh-unsup-bayes-clustering}
\end{equation}
\begin{proof}
由联合分布分解有 $p_\theta(x,Z=k)=\pi_kp_k(x)$，再对 $k$ 求和得到观测边缘分布，二者之比即为后验概率。若各类先验均为 $1/K$，公共因子约去，后验与类条件密度成正比；若先验不均匀，则不能丢掉权重 $\pi_k$。
\end{proof}
因此，先估计生成模型再计算类别后验，是实现软聚类的一条路线；并不是所有软分配方法都必须先构建生成模型。

\section{机器学习三要素}
\label{sec:lh-unsup-three-elements}

无监督学习同样可以从\emph{模型、策略、算法}三个层面理解。模型规定允许怎样表示数据结构，策略规定怎样比较候选模型，算法则规定怎样从初值和训练数据求出具体结果。三者相互联系，但不能混为一谈：选定一个目标函数不等于给出了求解步骤，算法使目标值下降也不等于证明已经找到全局最优解。

在聚类中，模型可以是 $K$ 个中心及最近中心分配规则，策略可以是最小化样本到所属中心的平方距离和，算法可以交替更新分配与中心。在降维中，模型可以是 $d$ 维正交投影，策略可以是最小化平方重构误差，算法可以使用特征值分解或奇异值分解（singular value decomposition，SVD）。在概率模型估计中，模型可以是给定成分数的混合分布，策略可以是极大化观测对数似然，算法则可以采用期望最大化（expectation-maximization，EM）等迭代方法。

用一般记号表示，学习策略常写成
\begin{equation}
\min_{\theta\in\Theta}\;
\frac1N\sum_{j=1}^NL(x_j;\theta)+\lambda\Omega(\theta),
\qquad \lambda\geq0,
\label{eq:lh-unsup-general-objective}
\end{equation}
其中 $L$ 是由建模目标决定的单样本损失，$\Omega$ 是可选的复杂度惩罚，$\Theta$ 可以包含维数、正交性或概率归一化等约束。这里没有监督标签并不妨碍定义损失，例如 $L=-\log p_\theta(x)$ 就给出负对数似然。某些方法还需要同时优化样本的潜在分配或表示，这些量也属于优化变量。

教材列举的算法并不都归结为梯度下降：层次聚类逐步合并或分裂样本组，$K$ 均值采用交替更新，主成分分析可用矩阵分解求解，概率潜在语义分析和潜在狄利克雷分配则涉及潜在变量推断。是否需要迭代、是否有闭式步骤、能保证怎样的最优性或收敛性，都要结合具体模型与目标判断。

\section{无监督学习方法}
\label{sec:lh-unsup-methods}

\subsection{聚类：用中心概括样本}
\label{subsec:lh-unsup-kmeans}

层次聚类用一系列嵌套的分组描述不同尺度的相似关系；$K$ 均值则直接用 $K$ 个中心代表样本。对 $1\leq K\leq N$，以 $\mu_k\in\mathbb R^M$ 表示第 $k$ 类中心，$z_j\in\{1,\ldots,K\}$ 表示第 $j$ 个样本的分配，$K$ 均值的典型目标为
\begin{equation}
\min_{\mu_1,\ldots,\mu_K,\,z_1,\ldots,z_N}
J(\mu,z)=\sum_{j=1}^N\|x_j-\mu_{z_j}\|_2^2.
\label{eq:lh-unsup-kmeans-objective}
\end{equation}
它用欧氏平方距离定义组内差异，并不是对任意形状、任意相似性标准都合适的聚类准则。特征单位的改变会改变距离和结果，因而是否需要标准化应由任务决定。

给定初始中心，每轮先固定中心，按最近中心分配样本；再固定分配，取各组样本的均值：
\begin{equation}
z_j^{(t+1)}\in\underset{1\leq k\leq K}{\arg\min}\,
\|x_j-\mu_k^{(t)}\|_2^2,\qquad
\mu_k^{(t+1)}=\frac1{n_k^{(t+1)}}
\sum_{j:z_j^{(t+1)}=k}x_j,
\label{eq:lh-unsup-kmeans-updates}
\end{equation}
其中 $n_k^{(t+1)}$ 是新分配中第 $k$ 类的样本数，均值公式要求其大于 $0$。实现时需指定平局规则和空类处理方式，例如重新初始化空类的中心；还应输入迭代上限，并以分配不再变化或目标改变量足够小作为停止条件。达到上限时返回当前近似结果，而不能声称已得到最优解。

\begin{proposition}[$K$ 均值正常更新的单调性]
\label{prop:lh-unsup-kmeans-descent}
在每一类均非空、两步均按式~\eqref{eq:lh-unsup-kmeans-updates} 执行时，一次分配更新和一次中心更新均不增加目标~\eqref{eq:lh-unsup-kmeans-objective}。因此连续正常更新所得目标值单调不增且有下界，从而收敛；这个结论本身不保证达到全局最小值。
\end{proposition}
\begin{proof}
固定中心时，各样本对目标的贡献相互独立，选择最近中心逐项最小化平方距离。固定某一非空类 $\mathcal C_k$，记其均值为 $\bar x_k$，则对任意 $\mu$，
\[
\sum_{j\in\mathcal C_k}\|x_j-\mu\|_2^2
=\sum_{j\in\mathcal C_k}\|x_j-\bar x_k\|_2^2
+|\mathcal C_k|\,\|\mu-\bar x_k\|_2^2.
\]
交叉项因 $\sum_{j\in\mathcal C_k}(x_j-\bar x_k)=0$ 而消失，故取均值最小化该类误差。对全部类求和即得第二步的单调性；目标非负，故单调目标值有极限。
\end{proof}

\begin{example}
\label{ex:lh-unsup-kmeans}
取五个二维样本 $A=(0,0)$、$B=(0,2)$、$C=(1,1)$、$D=(4,3)$、$E=(4,5)$，令 $K=2$，初始中心为 $\mu_1^{(0)}=A$、$\mu_2^{(0)}=D$。求一次完整更新后的分组、中心与组内平方误差。
\end{example}
\begin{solve}
比较到两中心的距离，得到第一组为 $A,B,C$，第二组为 $D,E$，在旧中心下的目标值为 $0+4+2+0+4=10$。更新后的中心为 $\mu_1^{(1)}=(1/3,1)$、$\mu_2^{(1)}=(4,4)$，故
\[
J=\frac{10}{9}+\frac{10}{9}+\frac49+1+1
=\frac{14}{3}<10.
\]
再次分配时分组不变。这给出了当前初始化得到的稳定分组，但这一计算没有穷尽全部分组，不能单凭停止就宣称全局最优。一次稠密的分配与均值更新需要 $O(NKM)$ 次量级的运算，若执行 $T$ 轮，总时间量级为 $O(TNKM)$。
\end{solve}

\subsection{降维：用低维坐标保留主要变化}
\label{subsec:lh-unsup-pca}

\emph{主成分分析}（principal component analysis，PCA）提供了一个目标明确的线性降维例子。记样本均值为 $\bar x=N^{-1}\sum_jx_j$，中心化数据矩阵为 $\mathbf X_c=(x_1-\bar x,\ldots,x_N-\bar x)$，经验协方差矩阵取
\[
S=\frac1N\mathbf X_c\mathbf X_c^{\mathsf T}\in\mathbb R^{M\times M}.
\]
这里使用分母 $N$，改用 $N-1$ 只改变整体尺度，不改变主方向。令 $W=(w_1,\ldots,w_d)\in\mathbb R^{M\times d}$，约束 $W^{\mathsf T}W=I_d$，以
\begin{equation}
z_j=W^{\mathsf T}(x_j-\bar x),\qquad
\widehat x_j=\bar x+Wz_j
\label{eq:lh-unsup-pca-encoding}
\end{equation}
分别表示低维坐标和正交投影重构。

\begin{proposition}[PCA 的两个等价目标]
\label{prop:lh-unsup-pca-objective}
在上述欧氏平方损失与正交约束下，最小化重构误差
\begin{equation}
\min_{W^{\mathsf T}W=I_d}
\sum_{j=1}^N\|(x_j-\bar x)-WW^{\mathsf T}(x_j-\bar x)\|_2^2
\label{eq:lh-unsup-pca-reconstruction}
\end{equation}
等价于最大化保留的总方差 $\operatorname{tr}(W^{\mathsf T}SW)$。若 $S$ 的特征值为 $\lambda_1\geq\cdots\geq\lambda_M\geq0$，取前 $d$ 个特征向量作为 $W$ 的列即可达到最优，最小重构误差为 $N\sum_{i=d+1}^M\lambda_i$。
\end{proposition}
\begin{proof}
$WW^{\mathsf T}$ 是正交投影矩阵，由勾股分解可得总误差为
\[
\sum_j\|x_j-\bar x\|_2^2
-\sum_j\|W^{\mathsf T}(x_j-\bar x)\|_2^2
=N\operatorname{tr}(S)-N\operatorname{tr}(W^{\mathsf T}SW).
\]
第一项与 $W$ 无关，故得到两个目标的等价性。设 $u_i$ 是 $S$ 的一组标准正交特征向量，令 $a_i=\sum_{q=1}^d(u_i^{\mathsf T}w_q)^2$，则 $0\leq a_i\leq1$、$\sum_i a_i=d$，并且
\[
\operatorname{tr}(W^{\mathsf T}SW)=\sum_{i=1}^M\lambda_i a_i
\leq\sum_{i=1}^d\lambda_i.
\]
最后的不等式可由把权重从较小特征值移向尚未取满的较大特征值得到。取 $w_q=u_q$ 时达到等号，将其代回误差表达式即得结论。
\end{proof}

若截断处有特征值重合，最优子空间可能不唯一；即使子空间唯一，其中的正交坐标基也不唯一。对 $\mathbf X_c$ 作 SVD 时，其左奇异向量给出同样的主方向，因此矩阵分解是实现这一目标的算法工具。非线性降维则需要另选映射族与评价准则，不能直接套用以上线性投影结论。

\begin{example}
\label{ex:lh-unsup-pca}
样本为 $x_j=(t_j,2t_j)^{\mathsf T}$，其中 $t_j\in\{-2,-1,1,2\}$。求一个无重构损失的一维表示，并说明两个原特征之间的关系。
\end{example}
\begin{solve}
样本均值为零，协方差矩阵为
\[
S=\frac52\begin{pmatrix}1&2\\2&4\end{pmatrix}.
\]
其特征值为 $25/2$ 和 $0$，第一主方向可取 $w_1=(1,2)^{\mathsf T}/\sqrt5$。于是 $z_j=w_1^{\mathsf T}x_j=\sqrt5\,t_j$，且 $w_1z_j=x_j$，重构误差为零。虽然观测有两个坐标，但它们满足第二坐标恒为第一坐标的两倍，样本实际位于一维子空间中。
\end{solve}

这个例子说明了“高维观测、低维结构”的含义。实际数据有噪声时，舍弃方向通常带来非零重构误差；保留最大方差也不保证保留某个监督任务中最有区分力的信息。教材用高维样本投影到二维后的点群说明探索性分析的用途，应将这种可视化观察与已经验证的类别结论区分开。

\subsection{话题分析：文本的低维与概率表示}
\label{subsec:lh-unsup-topic-analysis}

\emph{话题分析}（topic analysis）希望从文档集合中发现语义相关的词语组，并描述每个文档由哪些话题构成。设词表大小为 $V$、文档数为 $D$，用词频矩阵 $\mathbf C\in\mathbb R^{V\times D}$ 表示语料，其中 $c_{vd}$ 是第 $v$ 个词在第 $d$ 个文档中的出现次数。这是式~\eqref{eq:lh-unsup-data-matrix} 的一个实例：特征是词，样本是文档；采用词频后，模型主要分析词的出现与共现，并未保留完整词序。

从降维角度，\emph{潜在语义分析}（latent semantic analysis，LSA）利用矩阵分解构造较低维的文档和词语表示，例如用截断 SVD 得到
\[
\mathbf C\approx U_K\Sigma_KV_K^{\mathsf T},
\qquad K<\min\{V,D\}.
\]
这种线性表示的坐标可以带符号，不能未经额外约束就解释为话题概率。从概率估计角度，则可将话题表示为词上的分布，将文档表示为话题上的分布。

具体地，设话题数为 $K$，记 $\varphi_{kv}=P(\text{词}=v\mid\text{话题}=k)$、$\vartheta_{dk}=P(\text{话题}=k\mid\text{文档}=d)$，满足 $\varphi_{kv},\vartheta_{dk}\geq0$、$\sum_v\varphi_{kv}=1$、$\sum_k\vartheta_{dk}=1$。在这类混合表示中，文档内一个词的边缘分布为
\begin{equation}
P(\text{词}=v\mid\text{文档}=d)
=\sum_{k=1}^K\vartheta_{dk}\varphi_{kv}.
\label{eq:lh-unsup-topic-mixture}
\end{equation}
\emph{概率潜在语义分析}（probabilistic latent semantic analysis，PLSA）使用这种话题混合结构建模词与文档的关系；\emph{潜在狄利克雷分配}（latent Dirichlet allocation，LDA）进一步以狄利克雷分布（Dirichlet distribution）对文档的话题比例作生成建模，常见的贝叶斯形式也对话题的词分布设置先验。这里仅交代模型之间的联系，参数估计和后验推断留待相应专章；马尔可夫链蒙特卡罗法（Markov chain Monte Carlo，MCMC）是相关概率推断工具之一，并非话题模型本身。

\begin{example}
\label{ex:lh-unsup-topic-mixture}
取六个词 $w_1,\ldots,w_6$，用两个理想化话题分布
\[
\varphi_1=(1/3,1/3,1/3,0,0,0),\qquad
\varphi_2=(0,0,0,1/3,1/3,1/3)
\]
说明文档的混合表示。若某文档的话题比例为 $(3/4,1/4)$，求其词分布。
\end{example}
\begin{solve}
由式~\eqref{eq:lh-unsup-topic-mixture}，前三个词各以概率 $(3/4)(1/3)=1/4$ 出现，后三个词各以概率 $(1/4)(1/3)=1/12$ 出现，总概率为 $3/4+1/4=1$。文档同时含有两个话题，只是比例不同，因此话题分析一般不是把每篇文档硬分配到唯一类别。这里的稀疏分布是人为给定的理想例子，不是声称从某个语料训练 LDA 后必然得到这些参数。
\end{solve}

可靠的话题发现需要语料提供足够的词语共现信息；少量、过短或缺乏差异的文档往往难以支持稳定的结构估计，但不存在脱离模型与数据条件的统一最少文档数。即使拟合较好，话题编号仍可置换，词语组的语义解释也需要结合内容核查，不能把模型分量直接当成人工定义的标准类别。

\subsection{图分析：随机游走与结点重要性}
\label{subsec:lh-unsup-graph-analysis}

\emph{图分析}（graph analytics）研究以图表示的实体关系，数据可以是有向图、无向图或超图。\emph{链接分析}（link analysis）是其中一类问题，关注链接结构所反映的结点关系与重要性。PageRank 将有向图上的访问过程表示为随机游走（random walk），用长期访问概率定义一种结点重要性，因此不需要预先标注哪些结点重要。

设有向图有 $n$ 个结点，暂设每个结点的出度 $d_j^+>0$。采用列表示出发结点的约定，定义转移矩阵
\begin{equation}
P_{ij}=\begin{cases}
1/d_j^+,&\text{存在从结点 $j$ 到结点 $i$ 的边},\\
0,&\text{否则}.
\end{cases}
\qquad \sum_{i=1}^nP_{ij}=1.
\label{eq:lh-unsup-graph-transition}
\end{equation}
若 $q_j^{(t)}$ 是第 $t$ 步处于结点 $j$ 的概率，$q^{(t)}$ 是相应列向量，则全概率公式给出
\begin{equation}
q^{(t+1)}=Pq^{(t)},\qquad
q^{(0)}\geq0,\qquad \mathbf1^{\mathsf T}q^{(0)}=1,
\label{eq:lh-unsup-random-walk}
\end{equation}
其中 $\mathbf1$ 是全 $1$ 列向量。\emph{平稳分布}（stationary distribution）是满足 $q^*=Pq^*$、$q^*\geq0$、$\mathbf1^{\mathsf T}q^*=1$ 的概率向量；一旦访问分布等于它，下一步仍保持不变。

\emph[caution]{存在平稳分布与从任意初值迭代收敛是不同结论。}有限状态链有平稳分布；不可约时平稳分布唯一，进一步非周期时，式~\eqref{eq:lh-unsup-random-walk} 才保证从任意初始分布收敛到它。不可约指任意两个状态之间均可沿正概率转移到达；非周期指返回某状态的所有可能正步数的最大公约数为 $1$，不可约链的各状态周期相同。这些是有限状态马尔可夫链的标准结论，证明留待相关专章。作为反例，只有 $A\to B$、$B\to A$ 的二结点图有唯一平稳分布 $(1/2,1/2)^{\mathsf T}$，但从 $(1,0)^{\mathsf T}$ 出发的分布在两个状态间交替，并不收敛。

\begin{example}
\label{ex:lh-unsup-pagerank}
另取三结点图，边为 $A\to B$、$A\to C$、$B\to C$、$C\to A$。按结点顺序 $A,B,C$ 写出转移矩阵并求平稳分布。
\end{example}
\begin{solve}
从 $A$ 出发以各 $1/2$ 的概率访问 $B,C$，从 $B,C$ 出发分别必到 $C,A$，故
\[
P=\begin{pmatrix}0&0&1\\1/2&0&0\\1/2&1&0\end{pmatrix}.
\]
平稳方程给出 $q_A^*=q_C^*$、$q_B^*=q_A^*/2$；再由概率和为 $1$，得到 $q^*=(2/5,1/5,2/5)^{\mathsf T}$。该图强连通，且从 $A$ 返回 $A$ 有长度为 $2$ 和 $3$ 的回路，所以对应链不可约且非周期，迭代从任意概率初值收敛。均匀初值的下一步为 $(1/3,1/6,1/2)^{\mathsf T}$，已经不同于均匀分布，说明链接关系会改变长期访问权重。
\end{solve}

真实链接图可能含有出度为零的悬挂结点，也可能不强连通或有周期。为补足教材概述中的收敛条件，可先用一个严格正的概率向量 $v$ 替代悬挂结点对应的转移列，再引入阻尼系数 $0<\alpha<1$，每步以概率 $\alpha$ 沿图转移，以概率 $1-\alpha$ 按 $v$ 重新选择结点。此时
\begin{equation}
G=\alpha P+(1-\alpha)v\mathbf1^{\mathsf T},\qquad
q^{(t+1)}=\alpha Pq^{(t)}+(1-\alpha)v.
\label{eq:lh-unsup-damped-pagerank}
\end{equation}
这是对基本随机游走模型的补充，并非假定任意原图已经满足收敛条件。

\begin{proposition}[阻尼迭代的唯一性与收敛性]
\label{prop:lh-unsup-pagerank-contraction}
设 $P$ 非负且每列之和为 $1$，$v$ 是概率向量，$0<\alpha<1$。则式~\eqref{eq:lh-unsup-damped-pagerank} 在概率向量集合中有唯一不动点 $q^*$，从任意概率初值出发均收敛，且
\begin{equation}
\|q^{(t)}-q^*\|_1\leq\alpha^t\|q^{(0)}-q^*\|_1
\leq2\alpha^t,
\label{eq:lh-unsup-pagerank-rate}
\end{equation}
其中 $\|a\|_1=\sum_i|a_i|$。此结论只需 $v$ 为概率向量；取其各分量严格正，还使 $G$ 的全部元素为正。
\end{proposition}
\begin{proof}
对任意向量 $a$，由三角不等式和列和约束有
\[
\|Pa\|_1\leq\sum_i\sum_jP_{ij}|a_j|
=\sum_j|a_j|=\|a\|_1.
\]
令 $T(q)=\alpha Pq+(1-\alpha)v$，则 $T$ 将概率向量映为概率向量，且 $\|T(q)-T(r)\|_1\leq\alpha\|q-r\|_1$。连续迭代满足 $\|q^{(t+1)}-q^{(t)}\|_1\leq\alpha^t\|q^{(1)}-q^{(0)}\|_1$，几何级数的尾和趋于零，所以迭代是柯西列。概率向量集合闭，故极限仍在其中；由连续性极限是不动点。若存在两个不动点，其距离不超过自身的 $\alpha$ 倍，只能相等。最后反复应用压缩不等式，并用任意两个概率向量的 $1$ 范数距离至多为 $2$，得到误差界。
\end{proof}

实际算法输入图、$\alpha$、$v$、误差容限 $\varepsilon>0$ 与迭代上限，从均匀分布等概率初值开始反复更新。由同一压缩不等式可得 $\|q-q^*\|_1\leq\|q-T(q)\|_1/(1-\alpha)$，因此当一步残差不超过 $(1-\alpha)\varepsilon$ 时，当前结果与不动点的距离不超过 $\varepsilon$；否则继续迭代或在达到上限时报告近似结果。若按稀疏边表计算并单独累计悬挂结点概率，每轮运算量为 $O(|E|+n)$，其中 $|E|$ 为边数，无须显式存储稠密矩阵 $G$。

在没有阻尼且满足平稳方程时，一个结点接收的权重为各入链来源结点权重除以其出度后的总和；在阻尼模型中，还要乘以 $\alpha$ 并加上重启项。因此，重要性取决于来源结点的重要程度与其分配给该边的份额，而不仅是入度，也不能保证多一条入边就一定在整个图中排名更高。PageRank 最初的网页链接解释可由此理解：网页是结点，超链接是有向边，长期访问概率提供一种结构性排序信号，而不是网页内容质量的绝对度量。

本章的三类基本问题与四类应用并非互斥分类：概率模型可以导出软聚类，话题分析可以采用降维或概率建模，图分析可以利用概率转移刻画关系结构。阅读后续具体方法时，应始终对应其数据表示、待学习对象、优化或估计准则以及算法适用条件，判断所得结构在何种意义下概括了原始数据。

% 教材来源：机器学习/textbooks/机器学习（分章原文）/机器学习-14.pdf。
% 对应教材第 14 章，原书第 255--269 页；本文件沿用自动章节编号。

\fi

\selectchapter{14}
\ifcompilechapter
\chapter{聚类方法}
\label{ch:lh-clustering}

\emph{聚类}（clustering）根据样本之间的相似性，把没有类别标签的数据归入若干个\emph{簇}（cluster），以发现数据的结构。类别由所选表示、相似性度量和算法共同确定，因而聚类结果并不是脱离这些选择而唯一存在的“真实分类”。本章先说明样本间与簇间的距离，再介绍产生嵌套划分的层次聚类，以及直接优化簇内平方误差的 $k$ 均值聚类；重点是理解度量、目标函数和求解算法各自承担的作用。

本章采用教材的列向量约定：给定 $n$ 个观测样本，每个样本有 $m$ 个实值属性，记
\begin{equation}
X=(x_1,\ldots,x_n)=[x_{rj}]_{m\times n},\qquad
x_j=(x_{1j},\ldots,x_{mj})^{\mathsf T}\in\mathbb R^m.
\label{eq:lh-cluster-data}
\end{equation}
这里 $X$ 是观测数据矩阵；即使两个样本的坐标相同，也保留不同的样本索引。因此下文严格地对索引集合 $\mathcal I=\{1,\ldots,n\}$ 进行划分，$G\subseteq\mathcal I$ 表示一个簇，$i\in G$ 表示样本 $x_i$ 属于该簇。这样可以避免把重复观测误删成集合中的同一个元素。

\section{聚类的基本概念}
\label{sec:lh-cluster-basics}

\subsection{相似度或距离}
\label{subsec:lh-cluster-similarity}

\emph{相似度}（similarity）通常越大表示越相似，\emph{距离}（distance）通常越小表示越相似。严格的度量 $d$ 应满足非负性、$d(x,y)=0$ 当且仅当 $x=y$、对称性和三角不等式；聚类也可以使用不满足全部度量公理的相异度，但此时不能直接使用只对度量成立的推论。属性的单位、尺度和相关性会影响距离，故在选择算法以前就应确定哪些差异应被视为重要。

\begin{definition}[闵可夫斯基距离]
\label{def:lh-cluster-minkowski}
对 $x_i,x_j\in\mathbb R^m$ 及 $1\leq p<\infty$，\emph{闵可夫斯基距离}（Minkowski distance）定义为
\begin{equation}
d_p(x_i,x_j)=\left(\sum_{r=1}^m|x_{ri}-x_{rj}|^p\right)^{1/p}.
\label{eq:lh-cluster-minkowski}
\end{equation}
$p=1$ 时为\emph{曼哈顿距离}（Manhattan distance），$p=2$ 时为\emph{欧氏距离}（Euclidean distance）；其极限形式为\emph{切比雪夫距离}（Chebyshev distance）
\begin{equation}
d_\infty(x_i,x_j)=\max_{1\leq r\leq m}|x_{ri}-x_{rj}|.
\label{eq:lh-cluster-chebyshev}
\end{equation}
\end{definition}

令 $a=\max_r|x_{ri}-x_{rj}|$，则 $a\leq d_p(x_i,x_j)\leq m^{1/p}a$，由夹逼可得 $p\to\infty$ 时的极限。$p$ 增大时，较大的坐标差异受到更强的相对强调；当 $0<p<1$ 时，同样的表达式一般不满足三角不等式，不属于这里的距离族。若不同属性的单位不同，直接使用欧氏距离会使数值尺度大的属性占主导；按属性标准差缩放是一种处理方式，但标准差为零的属性必须另行处理，而且缩放本身也体现了对各属性重要性的选择。

\begin{definition}[马哈拉诺比斯距离]
\label{def:lh-cluster-mahalanobis}
设 $S\in\mathbb R^{m\times m}$ 为对称正定矩阵，\emph{马哈拉诺比斯距离}（Mahalanobis distance），简称马氏距离，定义为
\begin{equation}
d_S(x_i,x_j)=\sqrt{(x_i-x_j)^{\mathsf T}S^{-1}(x_i-x_j)}.
\label{eq:lh-cluster-mahalanobis}
\end{equation}
通常以数据的样本协方差矩阵作为 $S$。具体地，当 $n\geq2$ 时，令 $\bar x=n^{-1}\sum_{j=1}^n x_j$，取 $S=(n-1)^{-1}\sum_{j=1}^n(x_j-\bar x)(x_j-\bar x)^{\mathsf T}$，并要求它可逆。
\end{definition}

正定性保证 $S^{-1/2}$ 存在，故 $d_S(x_i,x_j)=\|S^{-1/2}x_i-S^{-1/2}x_j\|_2$，即先对数据进行\emph{白化}（whitening），再计算欧氏距离。若 $S=\operatorname{diag}(s_1^2,\ldots,s_m^2)$，它就是按标准差缩放后的欧氏距离；若 $S=I_m$，则退化为普通欧氏距离。\emph[caution]{协方差矩阵为单位矩阵只说明各分量不相关且方差为 $1$，一般不能推出独立。}对于可逆线性变换 $y=Ax+b$，若同时把 $S$ 变为 $ASA^{\mathsf T}$，则所得马氏距离不变；这里的尺度不变性要求数据与协方差作相应变换，而不是任意改变数据后仍固定原来的 $S$。

样本协方差矩阵的秩至多为 $\min\{m,n-1\}$，所以 $n\leq m$ 时必然不可逆，属性之间有线性依赖时也会出现奇异性。此时不能直接代入式~\eqref{eq:lh-cluster-mahalanobis}；可在适当子空间中计算，或用 $S+\lambda I_m$、$\lambda>0$ 正则化。使用伪逆则可能使不同点的距离为零，因而一般只得到\emph{伪度量}（pseudometric）；这些处理改变或限制了所采用的几何结构。

\begin{definition}[相关系数与夹角余弦]
\label{def:lh-cluster-correlation-cosine}
令 $a_i=m^{-1}\sum_{r=1}^m x_{ri}$、$u_i=x_i-a_i\mathbf1_m$，其中 $\mathbf1_m$ 为全 $1$ 向量。当 $u_i,u_j$ 均非零时，两个样本向量的\emph{相关系数}（correlation coefficient）为
\begin{equation}
r_{ij}=\frac{u_i^{\mathsf T}u_j}{\|u_i\|_2\|u_j\|_2}
=\frac{\sum_{r=1}^m(x_{ri}-a_i)(x_{rj}-a_j)}
{\sqrt{\sum_{r=1}^m(x_{ri}-a_i)^2}\sqrt{\sum_{r=1}^m(x_{rj}-a_j)^2}}.
\label{eq:lh-cluster-correlation}
\end{equation}
当 $x_i,x_j$ 均非零时，\emph{夹角余弦}（cosine similarity）为
\begin{equation}
s_{ij}=\frac{x_i^{\mathsf T}x_j}{\|x_i\|_2\|x_j\|_2}.
\label{eq:lh-cluster-cosine}
\end{equation}
两者都位于 $[-1,1]$，相关系数是将两个向量分别按其坐标均值中心化之后的夹角余弦。
\end{definition}

这里的相关系数是在两个已观测向量的对应坐标之间计算的描述量，不能仅凭该式就把属性坐标视为独立同分布的随机样本。两种相似度都忽略正比例缩放，而相关系数还忽略向量各坐标共同增加的常数。$r_{ij}$ 或 $s_{ij}$ 接近 $1$ 表示同向，接近 $-1$ 表示反向；只有在应用明确把反向模式也看成相似时，才宜使用 $|r_{ij}|$。不能把“相关系数绝对值大”与“相关系数大”混为同一种规则。常量向量的相关系数、零向量的余弦都未定义；对于含负坐标的数据，余弦接近 $0$ 表示近乎正交，并不涵盖全部不相似情形。

\begin{example}
\label{ex:lh-cluster-metric-choice}
取 $u=(1,2,3)^{\mathsf T}$、$v=(2,4,6)^{\mathsf T}$、$w=(1,2,4)^{\mathsf T}$，比较欧氏距离与余弦相似度给出的关系。
\end{example}
\begin{solve}
有 $\|u-w\|_2=1<\sqrt{14}=\|u-v\|_2$，故按位置距离看，$w$ 更接近 $u$。但 $v=2u$，所以 $s(u,v)=1$，而 $s(u,w)=17/\sqrt{294}<1$，按方向相似度看，$v$ 更接近 $u$。两种结论分别保留绝对数值与相对方向，选择哪一种应由问题决定。
\end{solve}

若希望把余弦用于最小相异度的合并规则，可以使用 $1-s_{ij}$，但它一般不满足三角不等式。对单位化向量 $\widehat x_i=x_i/\|x_i\|_2$，有 $\|\widehat x_i-\widehat x_j\|_2^2=2(1-s_{ij})$，所以 $\sqrt{2(1-s_{ij})}$ 是单位球面上的欧氏距离。相似度到相异度的转换应与后续规则相配合，不能一面把较大的值解释为更相似，一面仍选择数值最小的一对。

\subsection{类或簇}
\label{subsec:lh-cluster-clusters}

\begin{definition}[硬聚类与软聚类]
\label{def:lh-cluster-partition}
\emph{硬聚类}（hard clustering）把 $\mathcal I$ 划分为若干非空子集 $G_1,\ldots,G_k$，满足
\begin{equation}
G_\ell\cap G_q=\varnothing\quad(\ell\ne q),\qquad
\bigcup_{\ell=1}^kG_\ell=\mathcal I.
\label{eq:lh-cluster-partition}
\end{equation}
每个样本恰属于一个簇。\emph{软聚类}（soft clustering）则允许以多个隶属程度或归属概率描述一个样本，例如用 $r_{i\ell}\geq0$、$\sum_{\ell=1}^k r_{i\ell}=1$ 表示其对各簇的归属；允许类别重叠的描述也比硬划分更宽泛。本章两个主要算法都输出硬划分。
\end{definition}

簇是否“足够紧密”还需要额外标准。设 $G\subseteq\mathcal I$，$n_G=|G|\geq2$，$d_{ij}=d(x_i,x_j)$。给定 $T>0$，教材依次考虑四类约束：所有点对都近、每个点至少有一个近邻、每个点到其余点的平均距离小，以及全簇的平均距离小且极端距离受控。将不同点对明确写出，分别为
\begin{align}
&d_{ij}\leq T,\qquad &&\forall i,j\in G,
\label{eq:lh-cluster-diameter-condition}\\
&\forall i\in G\ \exists j\in G\setminus\{i\},\qquad &&d_{ij}\leq T,
\label{eq:lh-cluster-neighbor-condition}\\
&\frac1{n_G-1}\sum_{\substack{j\in G\\j\ne i}}d_{ij}\leq T,
&&\forall i\in G,
\label{eq:lh-cluster-point-average}\\
&\frac1{n_G(n_G-1)}\sum_{i\in G}\sum_{\substack{j\in G\\j\ne i}}d_{ij}\leq T,
\qquad &&d_{ij}\leq V\quad(\forall i,j\in G),
\label{eq:lh-cluster-global-average}
\end{align}
其中最后一种另给 $V>0$。这些是不同的紧密性标准，并非互相等价的簇定义，也不会单独决定整个数据集应如何划分。单元素簇可约定直径为零，但上述含 $n_G-1$ 的平均距离条件不直接适用。

\begin{proposition}[紧密性条件之间的关系]
\label{prop:lh-cluster-conditions}
当 $n_G\geq2$ 时，在同一阈值 $T$ 下，式~\eqref{eq:lh-cluster-diameter-condition} 蕴含式~\eqref{eq:lh-cluster-point-average}，后者又蕴含式~\eqref{eq:lh-cluster-neighbor-condition}。若还要求 $V\geq T$，则式~\eqref{eq:lh-cluster-diameter-condition} 也蕴含式~\eqref{eq:lh-cluster-global-average}。
\end{proposition}
\begin{proof}
若每个 $d_{ij}\leq T$，则固定 $i$ 后，其余 $n_G-1$ 个距离的平均值不超过 $T$；再对 $i$ 平均，全簇的平均距离也不超过 $T$。若某个点到其他所有点的距离都大于 $T$，其平均值也大于 $T$，故平均值约束必保证至少存在一个不超过 $T$ 的距离。最后，$d_{ij}\leq T\leq V$ 给出额外的最大距离约束。
\end{proof}

不能把上述结论理解为对任意预先指定的 $V$ 都成立：若 $V<T$，第一个条件不保证距离小于 $V$。逆向蕴含也一般不成立。例如直线上的四点 $0,0.1,10,10.1$ 在 $T=0.2$ 时各有近邻，却不满足直径约束；三点 $0,1,2$ 在 $T=1.5$ 时，每点到其余点的平均距离都不超过 $1.5$，直径却为 $2$。前者也说明“每点有近邻”甚至不能保证阈值邻接图连通。

\begin{definition}[簇的中心、直径与散布]
\label{def:lh-cluster-statistics}
对非空簇 $G$，其\emph{中心}（centroid）、\emph{直径}（diameter）和\emph{散布矩阵}（scatter matrix）分别为
\begin{align}
\mu_G&=\frac1{n_G}\sum_{i\in G}x_i,
\label{eq:lh-cluster-centroid}\\
D_G&=\max_{i,j\in G}d_{ij},
\label{eq:lh-cluster-diameter}\\
A_G&=\sum_{i\in G}(x_i-\mu_G)(x_i-\mu_G)^{\mathsf T}.
\label{eq:lh-cluster-scatter}
\end{align}
当 $n_G\geq2$ 时，本笔记采用通常的样本协方差归一化 $S_G=A_G/(n_G-1)$；若把簇视为等权经验分布，其协方差则为 $A_G/n_G$。两种归一化应加以区分，分母由样本数决定，$A_G,S_G$ 的维度均为 $m\times m$。
\end{definition}

中心是坐标均值，不一定是实际样本，也不一定处于高密度区域。散布矩阵保留各方向的变化及其联系，欧氏情形下其迹满足 $\operatorname{tr}(A_G)=\sum_{i\in G}\|x_i-\mu_G\|_2^2$，恰好是该簇对 $k$ 均值目标的贡献。这里把 $S_G$ 称为样本协方差是代数上的约定；若要进一步声称它是某个总体协方差的无偏估计，还需相应的随机抽样假设，数据依赖的聚类结果本身不自动保证该性质。

\subsection{类与类之间的距离}
\label{subsec:lh-cluster-linkage}

对于两个非空簇 $G_p,G_q$，簇间相异度也称\emph{连接}（linkage），它把样本间的比较扩展到簇的比较。教材列出的四种规则为
\begin{align}
D_{\mathrm{single}}(G_p,G_q)&=\min_{i\in G_p,\,j\in G_q}d_{ij},
\label{eq:lh-cluster-single-linkage}\\
D_{\mathrm{complete}}(G_p,G_q)&=\max_{i\in G_p,\,j\in G_q}d_{ij},
\label{eq:lh-cluster-complete-linkage}\\
D_{\mathrm{centroid}}(G_p,G_q)&=d(\mu_{G_p},\mu_{G_q}),
\label{eq:lh-cluster-centroid-linkage}\\
D_{\mathrm{average}}(G_p,G_q)&=\frac1{n_pn_q}\sum_{i\in G_p}\sum_{j\in G_q}d_{ij},
\label{eq:lh-cluster-average-linkage}
\end{align}
其中 $n_p=|G_p|$、$n_q=|G_q|$，依次称为\emph{单连接}（single linkage，最短距离）、\emph{完全连接}（complete linkage，最长距离）、中心距离和平均距离。中心距离要求样本处于可以计算均值的空间；只给出一般相异度矩阵时，单连接、完全连接以及平均连接仍可通过汇总跨簇点对计算，中心距离则未必有定义。

单连接只需存在一对近邻就可能合并，容易沿一串近邻形成长链；完全连接要求最远的跨簇点对也足够接近，倾向于限制簇的直径。平均连接综合所有跨簇点对，中心距离只比较两个均值，可能忽略形状与内部离散程度。即使点间 $d$ 是度量，这些连接规则也不一定在所有簇的集合上构成度量，例如不同簇的均值可以相同。连接方式的改变会改变后续合并次序，并非仅改变距离的计算速度。

\section{层次聚类}
\label{sec:lh-hierarchical-clustering}

\emph{层次聚类}（hierarchical clustering）构造一系列彼此嵌套的划分。\emph{聚合聚类}（agglomerative clustering）从单元素簇出发，逐次合并；\emph{分裂聚类}（divisive clustering）从包含全部样本的一个簇出发，逐次拆分。每一层仍是硬划分，但不同层之间保留包含关系；只需在某一层截断，就可得到指定簇数的聚类结果。分裂方法还必须明确选择哪个簇以及如何二分，仅说“取最远的两个点”并不能决定其余样本的归属。

聚合方法需要确定样本间相异度、簇间连接和停止条件。以下用最小连接值作为合并规则，输入 $n$ 个样本及其距离矩阵 $D=[d_{ij}]_{n\times n}$，输出合并记录及最终划分：先置 $\mathcal G^{(0)}=\{\{1\},\ldots,\{n\}\}$；第 $t$ 步在当前不同簇中选择
\begin{equation}
(A,B)\in\operatorname*{arg\,min}_{\substack{G_p,G_q\in\mathcal G^{(t)}\\p<q}}
D(G_p,G_q),\qquad
\mathcal G^{(t+1)}=(\mathcal G^{(t)}\setminus\{A,B\})\cup\{A\cup B\}.
\label{eq:lh-cluster-agglomerative-step}
\end{equation}
记录被合并的两个簇及其连接值，并更新新簇与其余簇的连接；并列时按固定规则选择。若目标是完整层次结构，则执行 $n-1$ 次合并；若只需 $k$ 个簇，则在 $n-k$ 次合并后停止。若限制合并后簇的直径，则应在接受合并前检查该约束；单连接值小不保证新簇的直径小。

\begin{proposition}[连接值的递推更新]
\label{prop:lh-cluster-linkage-update}
设 $U=A\cup B$，$A,B,C$ 为互不相交的非空簇，则
\begin{align}
D_{\mathrm{single}}(U,C)&=\min\{D_{\mathrm{single}}(A,C),D_{\mathrm{single}}(B,C)\},
\notag\\
D_{\mathrm{complete}}(U,C)&=\max\{D_{\mathrm{complete}}(A,C),D_{\mathrm{complete}}(B,C)\},
\notag\\
D_{\mathrm{average}}(U,C)&=\frac{|A|D_{\mathrm{average}}(A,C)+|B|D_{\mathrm{average}}(B,C)}{|A|+|B|}.
\label{eq:lh-cluster-linkage-updates}
\end{align}
中心距离可先用 $\mu_U=(|A|\mu_A+|B|\mu_B)/(|A|+|B|)$ 更新中心，再计算 $d(\mu_U,\mu_C)$。
\end{proposition}
\begin{proof}
$U$ 与 $C$ 的点对集合是 $A$ 与 $C$、$B$ 与 $C$ 的点对集合的并，故最小值与最大值分别按上述方式更新。平均连接的距离总和拆成两部分，分别为 $|A||C|D_{\mathrm{average}}(A,C)$ 与 $|B||C|D_{\mathrm{average}}(B,C)$，再除以 $(|A|+|B|)|C|$ 即得结论。中心更新则直接来自均值定义。
\end{proof}

初始欧氏距离矩阵的计算需要 $O(n^2m)$ 时间和 $O(n^2)$ 存储。若每轮扫描当前所有簇对寻找最小值，总扫描量为 $\sum_{q=2}^nO(q^2)=O(n^3)$；结合式~\eqref{eq:lh-cluster-linkage-updates}，单连接、完全连接或平均连接的一种直接实现的总时间为 $O(n^2m+n^3)$。教材给出的 $O(n^3m)$ 可理解为在合并过程中反复以 $m$ 维运算重新计算距离时的粗上界，不能视为所有实现必需的复杂度。以上区分了输入距离的构造、最小值搜索及连接值更新三个开销来源。

\begin{example}
\label{ex:lh-cluster-agglomerative-five}
教材的五样本例采用单连接。为突出合并所需的信息，记 $G_i=\{i\}$，已知最小的非对角距离为 $d_{35}=1$；合并 $G_3,G_5$ 后，$G_6=\{3,5\}$ 到 $G_1,G_2,G_4$ 的单连接值依次为 $2,5,5$，其余三对距离为 $d_{12}=7$、$d_{14}=9$、$d_{24}=4$。求后续合并次序。
\end{example}
\begin{solve}
第一步在连接值 $1$ 处得到 $G_6$。随后最小连接值为 $2$，故把 $G_1$ 与 $G_6$ 合并为 $G_7=\{1,3,5\}$。利用最小值更新，有
\[
D_{\mathrm{single}}(G_7,G_2)=\min\{7,5\}=5,\qquad
D_{\mathrm{single}}(G_7,G_4)=\min\{9,5\}=5.
\]
由于 $D_{\mathrm{single}}(G_2,G_4)=4$，下一步合并为 $G_8=\{2,4\}$，最后在连接值 $5$ 处把 $G_7,G_8$ 合并。完整合并值依次为 $1,2,4,5$；若需要两个簇，应停在 $\{\{1,3,5\},\{2,4\}\}$。这些记录即可表达教材树状图中的包含关系，无须另行假设各点的坐标。
\end{solve}

单连接的合并值不会下降，因为新连接是两个旧连接的最小值，而被选中的合并值已经是当前全局最小值；完全连接和平均连接也有相应的单调性。但中心连接可能出现新合并值低于先前合并值的情形，不能不加条件地把所有层次聚类都理解为合并高度单调的树。聚合过程一旦合并便不再拆开，早期选择会影响整个层次结构。

作为教材分裂聚类习题的一种具体构造，可先计算所有点间距离，从 $\{\mathcal I\}$ 出发，每轮选取直径最大的非单元素簇 $G$，找出其中最远的一对样本作为两个种子，把其他样本指派给距离较近的种子，按固定规则处理并列，并使两个种子分别保留在两个子簇中；重复至簇数达到指定的 $k\leq n$。若所选簇直径为零，则按索引直接分成两个非空部分。对于含 $s$ 个点的簇，扫描最远点对需要 $O(s^2)$ 次距离查表；在最不均衡的连续分裂中，$\sum_{s=2}^nO(s^2)=O(n^3)$，故缓存距离的一种朴素实现总时间为 $O(n^2m+n^3)$、存储为 $O(n^2)$。这是一种明确的分裂启发式，复杂度依赖所选规则，不能作为所有分裂聚类算法的统一结论。

\section{\texorpdfstring{$k$}{k} 均值聚类}
\label{sec:lh-kmeans}

\subsection{模型}
\label{subsec:lh-kmeans-model}

\emph{$k$ 均值聚类}（$k$-means clustering）给定簇数 $k$，通过中心描述各簇，得到平坦的硬划分。设 $1\leq k\leq n$，用满射
\begin{equation}
C:\mathcal I\longrightarrow\{1,\ldots,k\},\qquad G_\ell=\{i:C(i)=\ell\}
\label{eq:lh-kmeans-assignment}
\end{equation}
表示划分；满射条件保证每个簇非空。不同的簇编号可以描述同一个无标号划分，故模型的内容是哪些样本同属一簇，而不是编号本身。$k=1$ 时全部样本属于一个簇，$k=n$ 时可令每个观测独成一簇；通常研究 $1<k<n$。

\subsection{策略}
\label{subsec:lh-kmeans-objective}

对划分 $C$，令 $n_\ell=|G_\ell|$、$\mu_\ell=n_\ell^{-1}\sum_{C(i)=\ell}x_i$。$k$ 均值以\emph{簇内平方和}（within-cluster sum of squares）为目标，求解
\begin{equation}
\min_{\substack{C:\mathcal I\to\{1,\ldots,k\}\\C\text{ 为满射}}}
W(C),\qquad
W(C)=\sum_{\ell=1}^k\sum_{C(i)=\ell}\|x_i-\mu_\ell\|_2^2.
\label{eq:lh-kmeans-wcss}
\end{equation}
这一目标衡量同簇样本围绕均值的离散程度，不是所有簇间距离的直接最大化。虽然教材在目标中使用欧氏距离平方，\emph[caution]{平方欧氏距离本身一般不满足三角不等式}：在直线上，$|0-2|^2=4>|0-1|^2+|1-2|^2=2$。因此应把它看作此处的平方损失，而不能自动赋予普通欧氏距离的全部性质。

\begin{proposition}[均值的最优性与成对距离表示]
\label{prop:lh-kmeans-centroid-optimum}
对非空簇 $G$ 及任意中心 $c\in\mathbb R^m$，有
\begin{align}
\sum_{i\in G}\|x_i-c\|_2^2
&=\sum_{i\in G}\|x_i-\mu_G\|_2^2+n_G\|c-\mu_G\|_2^2,
\label{eq:lh-kmeans-variance-decomposition}\\
\sum_{i\in G}\|x_i-\mu_G\|_2^2
&=\frac1{2n_G}\sum_{i\in G}\sum_{j\in G}\|x_i-x_j\|_2^2.
\label{eq:lh-kmeans-pairwise}
\end{align}
因而固定簇成员时，平方损失的唯一最优中心为 $\mu_G$。
\end{proposition}
\begin{proof}
将 $x_i-c=(x_i-\mu_G)+(\mu_G-c)$ 展开平方并求和，交叉项因 $\sum_{i\in G}(x_i-\mu_G)=0$ 消失，即得式~\eqref{eq:lh-kmeans-variance-decomposition}。另一方面，展开所有点对的平方距离可得
\[
\sum_{i,j\in G}\|x_i-x_j\|_2^2
=2n_G\sum_{i\in G}\|x_i\|_2^2-2\left\|\sum_{i\in G}x_i\right\|_2^2
=2n_G\sum_{i\in G}\|x_i-\mu_G\|_2^2.
\]
除以 $2n_G$ 即得第二式。由于 $n_G>0$，第一式的附加项仅在 $c=\mu_G$ 时为零。
\end{proof}

为说明为什么不能直接枚举最优划分，设 $S(n,k)$ 为把 $n$ 个有区别的观测划分成 $k$ 个无标号非空簇的数量，即第二类斯特林数（Stirling number of the second kind），则
\begin{equation}
S(n,k)=\frac1{k!}\sum_{q=0}^k(-1)^{k-q}\binom{k}{q}q^n.
\label{eq:lh-kmeans-stirling}
\end{equation}
其推导是先数到 $k$ 个有标号簇的满射：全部映射有 $k^n$ 个，按哪些簇为空使用容斥原理，得到右端求和；每个无标号非空划分恰对应 $k!$ 种簇标号，再除以 $k!$。对固定 $k\geq2$、$n\to\infty$，主项为 $k^n/k!$，其余项相对主项趋于零，故 $S(n,k)\sim k^n/k!$。这里的指数增长有明确的参数范围；$k=1$ 或 $k=n$ 时都只有一个无标号划分。教材指出一般欧氏平方和聚类的全局优化是 NP 困难问题，而候选划分多这一计数事实本身并不是 NP 困难性的证明。

\subsection{算法}
\label{subsec:lh-kmeans-algorithm}

引入独立的中心变量 $M=(c_1,\ldots,c_k)$，把目标写为
\begin{equation}
J(C,M)=\sum_{i=1}^n\|x_i-c_{C(i)}\|_2^2.
\label{eq:lh-kmeans-joint-objective}
\end{equation}
若固定 $M$ 且暂不强制非空约束，每个样本的归属可以单独优化，最优选择是最近中心；若固定一个非空划分 $C$，由命题~\ref{prop:lh-kmeans-centroid-optimum}，最优中心就是各簇均值。这给出教材采用的交替更新，也常称为 Lloyd 迭代。

算法的输入为 $X$、$k$、初始中心 $M^{(0)}$、最大迭代次数及可选容差，输出为近似划分与中心。可随机选择 $k$ 个样本作为初始中心；随机性只来自初始化，给定初始化及并列处理规则后，后续精确更新是确定的。第 $t$ 轮执行
\begin{align}
C^{(t)}(i)&\in\operatorname*{arg\,min}_{1\leq\ell\leq k}
\|x_i-c_\ell^{(t)}\|_2^2,\qquad i=1,\ldots,n,
\label{eq:lh-kmeans-assignment-step}\\
c_\ell^{(t+1)}&=\frac1{n_\ell^{(t)}}\sum_{C^{(t)}(i)=\ell}x_i,
\qquad n_\ell^{(t)}=|\{i:C^{(t)}(i)=\ell\}|>0.
\label{eq:lh-kmeans-center-step}
\end{align}
首次分配遇到等距时可选编号最小的中心；以后若原归属仍是最近中心之一，则保持原归属，否则按固定编号选择。更新均值后，若划分与上轮相同则停止；若采用中心位移或目标改变量小于容差的规则，则得到的是容差意义下的近似结果。达到迭代上限时也返回当前划分及其均值，并注明尚未满足所选收敛条件。

\begin{notealert}[type=caution,title={空簇必须单独处理}]
最近中心分配可能产生空簇，随机抽取不同索引也不能排除重复坐标或后续空簇。若 $n_\ell^{(t)}=0$，式~\eqref{eq:lh-kmeans-center-step} 无定义。一种保持 $k$ 个非空簇的修复是：从至少含两个样本的簇中取一个样本，单独组成空缺簇，并把新中心设为该样本，再更新各簇均值。$n\geq k$ 保证有空缺时总能找到这样的供体簇。此移动把被移样本对旧中心的非负损失改为零，其他样本先保留原中心，随后重算均值又不增加损失，故该修复不增大 $J$；它不再是固定全部中心后的纯最近中心分配，下面的有限终止证明则明确以无空簇情形为条件。
\end{notealert}

\begin{example}
\label{ex:lh-kmeans-five}
教材给出五个二维样本
\[
X=\begin{pmatrix}0&0&1&5&5\\2&0&0&0&2\end{pmatrix}.
\]
取 $k=2$，先以 $x_1=(0,2)^{\mathsf T}$、$x_2=(0,0)^{\mathsf T}$ 为初始中心，求算法的最终划分。
\end{example}
\begin{solve}
$x_1,x_2$ 分别属于自己的中心；$x_3,x_4,x_5$ 到两个初始中心的平方距离依次为 $(5,1)$、$(29,25)$、$(25,29)$。因此第一轮得到
\[
G_1=\{1,5\},\qquad G_2=\{2,3,4\},\qquad
c_1=(5/2,2)^{\mathsf T},\quad c_2=(2,0)^{\mathsf T}.
\]
到这两个新中心的平方距离按样本索引排列为
\[
\begin{pmatrix}
25/4&8\\41/4&4\\25/4&1\\41/4&9\\25/4&13
\end{pmatrix}.
\]
逐行选择较小值后归属不变，故算法停止。最终簇内平方和为 $W=25/2+(4+1+9)=53/2$。停止的原因是该划分与均值相互一致，不能仅据此声称它在全部二簇划分中最优。
\end{solve}

稠密数据的一次分配需要 $O(nkm)$ 时间，重算均值需要 $O(nm)$ 时间，因此每轮为 $O(nkm)$；运行 $T_{\mathrm{iter}}$ 轮的总时间为 $O(T_{\mathrm{iter}}nkm)$，还需计入所采用的初始化与空簇处理开销。输入数据占 $O(nm)$ 存储，若不保存全部点到中心的距离，额外只需 $O(n+km)$ 存储。教材给出的 $O(nkm)$ 应理解为单轮量级，而不是与迭代轮数无关的完整运行时间。

\subsection{算法特性}
\label{subsec:lh-kmeans-properties}

\begin{proposition}[目标单调性与有条件的有限终止]
\label{prop:lh-kmeans-convergence}
若每轮分配后各簇均非空，并使用精确算术、精确最近中心分配及均值更新，则交替迭代的目标值不增，且因非负而收敛。若并列时保持原归属，则以划分不变为停止条件的算法有限步终止。这个结论不保证全局最优。
\end{proposition}
\begin{proof}
固定中心时，式~\eqref{eq:lh-kmeans-assignment-step} 对每项分别取最小值；固定划分时，式~\eqref{eq:lh-kmeans-center-step} 由均值的最优性最小化平方和。因此
\begin{equation}
J(C^{(t+1)},M^{(t+2)})
\leq J(C^{(t+1)},M^{(t+1)})
\leq J(C^{(t)},M^{(t+1)}).
\label{eq:lh-kmeans-descent}
\end{equation}
其中首尾两项分别是 $W(C^{(t+1)})$ 与 $W(C^{(t)})$，故目标单调且下有界。进一步，保持等距点原归属时，只要有样本改变归属，它到新中心的距离就严格小于到原中心的距离，第二个不等式便严格成立。无空簇划分的数量有限，而每个划分经均值更新后对应唯一的 $W(C)$；严格下降时不可能重复访问旧划分，故最终必出现归属不变并停止。
\end{proof}

终止时，每个样本属于一个最近中心，每个中心是其簇均值，因而分别固定两块变量时都不能通过单独优化另一块来降低目标。这是交替优化的固定点条件；它不意味着同时改变划分与中心也不能改进。仅凭“非负目标单调下降”本身也只能推出目标值收敛，不能自动推出划分有限步稳定，后者还使用了有限划分数及严格下降规则。

\begin{example}
\label{ex:lh-kmeans-initialization}
对例~\ref{ex:lh-kmeans-five} 的同一组数据，改以 $x_1$ 与 $x_5$ 为初始中心，比较所得结果。
\end{example}
\begin{solve}
第一次分配得到 $G_1=\{1,2,3\}$、$G_2=\{4,5\}$，其均值为 $c_1=(1/3,2/3)^{\mathsf T}$、$c_2=(5,1)^{\mathsf T}$。再次按最近中心分配，归属不变，故这一划分也为固定点。它的平方和为
\[
W=\left(\frac{17}{9}+\frac59+\frac89\right)+(1+1)
=\frac{16}{3}<\frac{53}{2}.
\]
因此例~\ref{ex:lh-kmeans-five} 的固定点确实不是全局最优；本例给出了一个更优划分，但尚未通过枚举或其他下界证明本例划分就是全局最优解。
\end{solve}

初始化决定算法从哪里开始下降。可多次随机初始化，比较所得 $W$ 并保留较小者，也可先用层次聚类得到 $k$ 个簇，再取各簇中心或选取距其中心最近的样本作为初始中心。多次运行增加寻找较好固定点的机会，额外计算量约随运行次数增加，但有限次重启仍不构成全局最优保证。中心从某一初值移动的幅度也没有普遍的小量保证，不能由最近中心分配规则推出中心始终靠近初始位置。

平方损失使离中心很远的点具有较大影响，均值也可能被少量离群样本拉动。固定中心后，第 $\ell$ 个簇的判定区域由对所有 $q\ne\ell$ 成立的 $\|x-c_\ell\|_2^2\leq\|x-c_q\|_2^2$ 确定；若各中心互异，这些条件展开后给出以超平面分界的\emph{沃罗诺伊划分}（Voronoi partition），边界上的并列再按指定规则处理。它适合用欧氏空间中围绕中心的紧密程度刻画数据；对于弯曲、细长或密度差异很大的结构，较小的平方和未必对应希望发现的簇。更换相异度后也不能机械保留均值更新，例如一维绝对误差之和由中位数最小化，原有更新公式必须由新目标重新推导。

簇数 $k$ 需要结合数据与任务选择。记 $W_k^*$ 为恰有 $k$ 个非空簇时的全局最小平方和，则
\begin{equation}
W_{k+1}^*\leq W_k^*\quad(1\leq k<n),\qquad W_n^*=0.
\label{eq:lh-kmeans-k-monotone}
\end{equation}
因为总能把一个至少含两个样本的簇再拆成两个，先让两个新簇共用旧中心，损失保持不变，再分别取均值，损失不会增加。于是不能仅靠最小训练平方和来确定有意义的 $k$，否则总能偏好更细的划分；不同 $k$ 下的有限次算法输出甚至未必服从上述全局最优值的单调性。

教材建议观察簇数增加时平均直径等紧密性指标的变化。下降由快转缓的“肘部”可作为探索线索，但其存在、唯一性及与任务所需簇数的一致性都没有一般保证；任意运行所得平均直径也不必单调。二分查找要求事先存在明确且单调的判定条件，不能仅凭一条可能有拐点的经验曲线就保证找到“最优 $k$”。实际比较还需结合簇的可解释性、重采样或重启后的稳定性，以及领域中需要保留的结构。

\begin{remark}[与高斯混合模型及 EM 的联系]
\label{rem:lh-kmeans-gmm}
为比较硬归属与概率归属，设潜在类别 $Z\in\{1,\ldots,k\}$，$P(Z=\ell)=1/k$，且观测随机向量 $U$ 满足 $U\mid Z=\ell\sim\mathcal N(c_\ell,\sigma^2I_m)$，其中各混合分量共享固定的 $\sigma>0$。给定观测 $x_i$，EM 的 E 步计算责任度
\begin{equation}
r_{i\ell}=P(Z=\ell\mid U=x_i)
=\frac{\exp(-\|x_i-c_\ell\|_2^2/(2\sigma^2))}
{\sum_{q=1}^k\exp(-\|x_i-c_q\|_2^2/(2\sigma^2))},
\label{eq:lh-kmeans-gmm-responsibility}
\end{equation}
M 步对均值使用加权平均 $c_\ell=\sum_i r_{i\ell}x_i/\sum_i r_{i\ell}$。有限 $\sigma>0$ 且中心有限时，每个样本对各簇都有严格为正的责任度，而 $k$ 均值使用取 $0$ 或 $1$ 的硬归属。固定观测与中心，当 $\sigma\to0$ 且每个样本的最近中心唯一时，责任度趋于对应最近中心的示性值；若最近中心并列，极限不必是单一硬归属。

这一联系也可以从目标看出。令 $a_i=\min_\ell\|x_i-c_\ell\|_2^2$，则
\[
a_i\leq-2\sigma^2\log\left(\frac1k\sum_{\ell=1}^k
e^{-\|x_i-c_\ell\|_2^2/(2\sigma^2)}\right)
\leq a_i+2\sigma^2\log k.
\]
左边利用每个指数项至多为 $e^{-a_i/(2\sigma^2)}$，右边利用至少有一项等于这个值。因 $n,k$ 有限，求和并令 $\sigma\to0$，便得到 $\sum_i a_i$：混合模型负对数似然去掉与中心无关的项并乘以 $2\sigma^2$ 后，其极限是最近中心平方误差。因此 $k$ 均值可与这一受限球形高斯模型的小方差极限联系起来，但一般的高斯混合模型还可估计混合权重和各分量协方差，普通 EM 并不等同于 $k$ 均值。两者都可能受初始化影响，输出形式与所优化的目标则不同。
\end{remark}


% 教材来源：机器学习/textbooks/机器学习（分章原文）/机器学习-15.pdf。
% 对应教材第 15 章，原书第 271--295 页；沿用本文的自动章节编号。

\fi

\selectchapter{15}
\ifcompilechapter
\chapter{奇异值分解}
\label{ch:lh-svd}

把一个大矩阵表示成若干简单矩阵的乘积，可以揭示其主要变化方向，也为近似计算和数据压缩提供依据。\emph{奇异值分解}（singular value decomposition，SVD）把任意实矩阵分解为两个正交矩阵与一个非负矩形对角矩阵的乘积：正交矩阵确定输入、输出空间中的方向，奇异值刻画各方向的伸缩程度。保留全部非零奇异值可精确重构原矩阵，只保留最大的若干个则得到平方损失意义下的最优低秩近似；这也是主成分分析和潜在语义分析的重要线性代数基础。

本章按教材的“定义与性质、计算、矩阵近似”展开，并补充关键证明与适用条件。除另有说明，设 $A\in\mathbb R^{m\times n}$，$m,n$ 为正整数，$p=\min\{m,n\}$，$r=\operatorname{rank}(A)$；向量均为列向量，$\|x\|_2=(x^{\mathsf T}x)^{1/2}$，$I_d$ 表示 $d$ 阶单位矩阵。记 $\mathcal R(A)=\{Ax:x\in\mathbb R^n\}$ 为值域，即列空间，$\mathcal N(A)=\{x\in\mathbb R^n:Ax=0\}$ 为零空间；$\operatorname{span}\{x_1,\ldots,x_s\}$ 表示这些向量张成的子空间，$\mathcal S^\perp$ 表示子空间 $\mathcal S$ 的正交补。

\section{奇异值分解的定义与性质}
\label{sec:lh-svd-definition-properties}

\subsection{定义与定理}
\label{subsec:lh-svd-definition}

\begin{definition}[完全奇异值分解]
\label{def:lh-full-svd}
若 $A$ 可以写成
\begin{equation}
A=U\Sigma V^{\mathsf T},\qquad
U\in\mathbb R^{m\times m},\quad
\Sigma\in\mathbb R^{m\times n},\quad
V\in\mathbb R^{n\times n},
\label{eq:lh-full-svd}
\end{equation}
其中 $U,V$ 为\emph{正交矩阵}（orthogonal matrix），即 $U^{\mathsf T}U=UU^{\mathsf T}=I_m$、$V^{\mathsf T}V=VV^{\mathsf T}=I_n$；$\Sigma$ 为\emph{矩形对角矩阵}（rectangular diagonal matrix），满足
\begin{equation}
\Sigma_{ij}=
\begin{cases}
\sigma_i,&i=j\leq p,\\
0,&\text{其他情形},
\end{cases}
\qquad \sigma_1\geq\sigma_2\geq\cdots\geq\sigma_p\geq0,
\label{eq:lh-svd-rectangular-diagonal}
\end{equation}
则称式~\eqref{eq:lh-full-svd} 为 $A$ 的\emph{完全奇异值分解}（full singular value decomposition）。$\sigma_i$ 称为\emph{奇异值}（singular value）；写 $U=(u_1,\ldots,u_m)$、$V=(v_1,\ldots,v_n)$，其列分别称为\emph{左奇异向量}（left singular vector）与\emph{右奇异向量}（right singular vector）。
\end{definition}

矩形对角矩阵共有 $p$ 个对角位置，不能在 $m<n$ 时仍把 $n$ 个数全部写到其对角线上。两个正交矩阵保持各自空间中的内积和长度；$\Sigma$ 则允许伸缩、消去某些方向，以及在不同维度的空间之间进行映射。因此 SVD 不要求 $A$ 为方阵，也不要求它对称或可逆。

\begin{theorem}[奇异值分解的存在性]
\label{thm:lh-svd-existence}
任意实矩阵 $A\in\mathbb R^{m\times n}$ 均有完全奇异值分解，恰有 $r$ 个正奇异值。其余奇异值为零；当 $r=0$ 时 $A=0$。
\end{theorem}
\begin{proof}
从输入空间中的对称矩阵 $A^{\mathsf T}A$ 出发。对任意 $x\in\mathbb R^n$，有 $x^{\mathsf T}A^{\mathsf T}Ax=\|Ax\|_2^2\geq0$，故 $A^{\mathsf T}A$ 为半正定矩阵。由实对称矩阵的谱定理，存在一组标准正交特征向量 $v_1,\ldots,v_n$，对应特征值可按 $\lambda_1\geq\cdots\geq\lambda_n\geq0$ 排列。又因为
\[
A^{\mathsf T}Ax=0\ \Longleftrightarrow\ Ax=0,
\]
其中正向由 $\|Ax\|_2^2=0$ 得到，反向由直接代入得到，所以两矩阵的零空间相同，由秩与零度之和等于列数可知 $\operatorname{rank}(A^{\mathsf T}A)=r$。因此前 $r$ 个特征值为正，其余为零。

令 $\sigma_i=\sqrt{\lambda_i}$（$1\leq i\leq p$），并对 $1\leq i\leq r$ 定义
\begin{equation}
u_i=\frac{Av_i}{\sigma_i}.
\label{eq:lh-svd-left-construction}
\end{equation}
这些向量满足
\[
u_i^{\mathsf T}u_j
=\frac{v_i^{\mathsf T}A^{\mathsf T}Av_j}{\sigma_i\sigma_j}
=\frac{\lambda_j}{\sigma_i\sigma_j}v_i^{\mathsf T}v_j
=\begin{cases}1,&i=j,\\0,&i\ne j.\end{cases}
\]
因而它们是 $\mathcal R(A)$ 中的 $r$ 个标准正交向量，构成该空间的一组基。利用
\[
y\in\mathcal R(A)^\perp
\ \Longleftrightarrow\ y^{\mathsf T}Ax=0\ \text{对一切 }x\in\mathbb R^n
\ \Longleftrightarrow\ A^{\mathsf T}y=0,
\]
可在 $\mathcal N(A^{\mathsf T})$ 中选取标准正交基 $u_{r+1},\ldots,u_m$，把前面的向量补成 $\mathbb R^m$ 的标准正交基。令 $U=(u_1,\ldots,u_m)$、$V=(v_1,\ldots,v_n)$，并按式~\eqref{eq:lh-svd-rectangular-diagonal} 构造 $\Sigma$。由于 $Av_i=\sigma_i u_i$ 对 $i\leq r$ 成立，而 $Av_i=0$ 对 $i>r$ 成立，逐列比较得到 $AV=U\Sigma$，右乘 $V^{\mathsf T}$ 即得所需分解。若 $r=0$，取任意正交矩阵 $U,V$ 与零矩阵 $\Sigma$ 即可。
\end{proof}

这个构造明确区分了两类方向：正奇异值对应的左右奇异向量由 $Av_i=\sigma_i u_i$ 联系，零奇异值对应的方向通过零空间补齐。\emph[caution]{式~\eqref{eq:lh-svd-left-construction} 只能用于正奇异值}，不能在 $\sigma_i=0$ 时用除法求左奇异向量。

\subsection{紧奇异值分解与截断奇异值分解}
\label{subsec:lh-compact-truncated-svd}

完全分解包含了零空间的基，而重构 $A$ 只需要非零奇异值及其对应向量。下面先考虑 $r\geq1$，令 $U_r=(u_1,\ldots,u_r)\in\mathbb R^{m\times r}$、$V_r=(v_1,\ldots,v_r)\in\mathbb R^{n\times r}$，并令 $\Sigma_r=\operatorname{diag}(\sigma_1,\ldots,\sigma_r)\in\mathbb R^{r\times r}$。

\begin{definition}[紧分解与截断分解]
\label{def:lh-compact-truncated-svd}
精确等式
\begin{equation}
A=U_r\Sigma_rV_r^{\mathsf T}
\label{eq:lh-compact-svd}
\end{equation}
称为 $A$ 的\emph{紧奇异值分解}（compact singular value decomposition）。对整数 $1\leq k<r$，仅保留前 $k$ 个奇异值及其对应向量，得到
\begin{equation}
A_k=U_k\Sigma_kV_k^{\mathsf T},\qquad
U_k\in\mathbb R^{m\times k},\quad
\Sigma_k\in\mathbb R^{k\times k},\quad
V_k\in\mathbb R^{n\times k},
\label{eq:lh-truncated-svd}
\end{equation}
称为 $A$ 的\emph{截断奇异值分解}（truncated singular value decomposition）所给出的秩 $k$ 近似。其中 $\Sigma_k=\operatorname{diag}(\sigma_1,\ldots,\sigma_k)$，$U_k,V_k$ 分别取相应矩阵的前 $k$ 列。
\end{definition}

紧分解与原矩阵的秩相同，且保持精确等式；当 $k<r$ 时 $A_k\ne A$，应写 $A\approx A_k$，不能把近似关系当作精确分解。由于 $U_k^{\mathsf T}U_k=V_k^{\mathsf T}V_k=I_k$、$\Sigma_k$ 可逆，式~\eqref{eq:lh-truncated-svd} 的秩恰为 $k$。为统一后面的结论，约定 $A_0=0$、$A_r=A$。当 $r=0$ 时只需使用 $A_0=0$，无须另写含正奇异值的分解。

这里 $U_r,V_r$ 的列标准正交，但当它们不是方阵时，并非定义中的正交矩阵：一般只有 $U_r^{\mathsf T}U_r=I_r$，而 $U_rU_r^{\mathsf T}$ 是到 $\mathcal R(A)$ 的\emph{正交投影}（orthogonal projection），把向量映到该子空间并使残差与子空间正交，未必等于 $I_m$。若以稠密数组存储矩阵，并只存对角线上的奇异值，秩 $k$ 因子需要约 $k(m+n+1)$ 个数，而原矩阵需要 $mn$ 个数；只有前者更小时才有这种存储意义上的压缩。因此紧分解在代数上无损，却不必对任意矩阵节省存储。

\begin{example}
\label{ex:lh-svd-five-by-four}
教材例 15.1--15.3 使用同一个矩阵说明完全、紧与截断分解。设
\[
A=\begin{pmatrix}
1&0&0&0\\
0&0&0&4\\
0&3&0&0\\
0&0&0&0\\
2&0&0&0
\end{pmatrix}.
\]
求其紧奇异值分解，并写出一个完全分解以及保留两个奇异值的近似。
\end{example}
\begin{solve}
计算 $A^{\mathsf T}A=\operatorname{diag}(5,9,0,16)$，按降序排列可得奇异值 $4,3,\sqrt5,0$，所以 $r=3$。取
\[
U_3=\begin{pmatrix}
0&0&1/\sqrt5\\
1&0&0\\
0&1&0\\
0&0&0\\
0&0&2/\sqrt5
\end{pmatrix},\qquad
\Sigma_3=\operatorname{diag}(4,3,\sqrt5),\qquad
V_3^{\mathsf T}=\begin{pmatrix}0&0&0&1\\0&1&0&0\\1&0&0&0\end{pmatrix}.
\]
三列 $u_i$ 和三列 $v_i$ 各自标准正交，逐项相乘得到 $A=U_3\Sigma_3V_3^{\mathsf T}$。在 $U_3$ 后补上
\[
u_4=(0,0,0,1,0)^{\mathsf T},\qquad
u_5=(2,0,0,0,-1)^{\mathsf T}/\sqrt5,
\]
在 $V_3$ 后补上 $v_4=(0,0,1,0)^{\mathsf T}$，并令 $\Sigma\in\mathbb R^{5\times4}$ 的对角线为 $4,3,\sqrt5,0$，即得一个完全分解。仅保留前两项则得到
\begin{equation}
A_2=\begin{pmatrix}
0&0&0&0\\
0&0&0&4\\
0&3&0&0\\
0&0&0&0\\
0&0&0&0
\end{pmatrix}.
\label{eq:lh-svd-five-by-four-rank-two}
\end{equation}
这个近似舍去了第三个奇异方向上的贡献，原矩阵中的 $1$ 和 $2$ 因此同时消失。
\end{solve}

\subsection{几何解释}
\label{subsec:lh-svd-geometry}

把 $A$ 看作从 $\mathbb R^n$ 到 $\mathbb R^m$ 的线性映射。式~\eqref{eq:lh-full-svd} 把它分成三个连续作用：$V^{\mathsf T}$ 将输入改写到右奇异向量构成的正交坐标系，$\Sigma$ 沿坐标方向按奇异值缩放并处理维度差异，$U$ 再把结果变换到输出空间。正交变换可以包含旋转和反射，不能一概称为旋转；当 $m\ne n$ 时，中间一步也不能只理解成同一空间内的缩放。

对任意输入 $x$，令 $c_i=v_i^{\mathsf T}x$，便有
\begin{equation}
Ax=\sum_{i=1}^r\sigma_i c_i u_i,\qquad
\|Ax\|_2^2=\sum_{i=1}^r\sigma_i^2c_i^2.
\label{eq:lh-svd-directional-action}
\end{equation}
因此沿 $v_i$ 的单位输入被变成沿 $u_i$、长度为 $\sigma_i$ 的输出；属于 $\mathcal N(A)$ 的输入分量完全消失。单位球 $\{x:\|x\|_2\leq1\}$ 的像是在 $\mathcal R(A)$ 中、以 $u_i$ 为主轴方向、以 $\sigma_i$ 为半轴长的椭球体。若 $r<p$，它在相应方向上退化；若 $r=0$，整个像只有原点。由式~\eqref{eq:lh-svd-directional-action} 还可得 $\max_{\|x\|_2=1}\|Ax\|_2=\sigma_1$，因为上界由 $\sum_i c_i^2=1$ 得到，取 $x=v_1$ 可达到。

\begin{example}
\label{ex:lh-svd-geometric-two-by-two}
教材例 15.4 的线性变换由矩阵 $B=\begin{pmatrix}3&1\\2&1\end{pmatrix}$ 给出。求单位圆在该变换下所得椭圆的半轴长，并解释其与矩阵两列的区别。
\end{example}
\begin{solve}
有 $B^{\mathsf T}B=\begin{pmatrix}13&5\\5&2\end{pmatrix}$，特征方程为 $\lambda^2-15\lambda+1=0$，所以两个半轴长为
\[
\sigma_1=\sqrt{\frac{15+\sqrt{221}}2}\approx3.8643,\qquad
\sigma_2=\sqrt{\frac{15-\sqrt{221}}2}\approx0.2588.
\]
先取这两个特征值对应的单位特征向量 $v_i$，再令 $u_i=Bv_i/\sigma_i$，即可得到输入、输出的主轴方向。矩阵两列则是标准基的像 $Be_1=(3,2)^{\mathsf T}$、$Be_2=(1,1)^{\mathsf T}$，其内积为 $5$，并不正交；它们通常不是椭圆主轴。SVD 正是通过改变输入与输出两侧的正交基，把一般方向上的变换改写成互不混合的主轴伸缩。
\end{solve}

\subsection{主要性质}
\label{subsec:lh-svd-properties}

\begin{proposition}[奇异值与两侧特征分解]
\label{prop:lh-svd-eigen-relations}
若 $A=U\Sigma V^{\mathsf T}$，则
\begin{equation}
A^{\mathsf T}A=V(\Sigma^{\mathsf T}\Sigma)V^{\mathsf T},\qquad
AA^{\mathsf T}=U(\Sigma\Sigma^{\mathsf T})U^{\mathsf T}.
\label{eq:lh-svd-two-gram-matrices}
\end{equation}
两矩阵具有相同的正特征值 $\sigma_1^2,\ldots,\sigma_r^2$，重数也相同，但零特征值的个数分别为 $n-r$、$m-r$。对 $1\leq i\leq p$ 有
\begin{equation}
Av_i=\sigma_i u_i,\qquad A^{\mathsf T}u_i=\sigma_i v_i.
\label{eq:lh-svd-paired-vectors}
\end{equation}
若 $n>p$，余下的 $v_{p+1},\ldots,v_n$ 满足 $Av_i=0$；若 $m>p$，余下的 $u_{p+1},\ldots,u_m$ 满足 $A^{\mathsf T}u_i=0$。
\end{proposition}
\begin{proof}
对式~\eqref{eq:lh-full-svd} 转置并相乘，利用 $U^{\mathsf T}U=I_m$、$V^{\mathsf T}V=I_n$ 即得式~\eqref{eq:lh-svd-two-gram-matrices}。两个对角矩阵中非零的对角元均为 $\sigma_1^2,\ldots,\sigma_r^2$。再从 $AV=U\Sigma$、$A^{\mathsf T}U=V\Sigma^{\mathsf T}$ 逐列读取，就得到配对关系和多余方向上的零映射。
\end{proof}

\begin{proposition}[秩与四个基本子空间]
\label{prop:lh-svd-four-subspaces}
正奇异值的个数等于 $\operatorname{rank}(A)$，并且
\begin{equation}
\begin{aligned}
\mathcal R(A)&=\operatorname{span}\{u_1,\ldots,u_r\},&
\mathcal N(A^{\mathsf T})&=\operatorname{span}\{u_{r+1},\ldots,u_m\},\\
\mathcal R(A^{\mathsf T})&=\operatorname{span}\{v_1,\ldots,v_r\},&
\mathcal N(A)&=\operatorname{span}\{v_{r+1},\ldots,v_n\}.
\end{aligned}
\label{eq:lh-svd-four-subspaces}
\end{equation}
空向量组张成零子空间。这些基均标准正交，特别地 $\mathcal R(A)^\perp=\mathcal N(A^{\mathsf T})$、$\mathcal R(A^{\mathsf T})^\perp=\mathcal N(A)$。
\end{proposition}
\begin{proof}
左右乘可逆矩阵不改变秩，所以 $\operatorname{rank}(A)=\operatorname{rank}(\Sigma)=r$。式~\eqref{eq:lh-svd-directional-action} 给出 $\mathcal R(A)$ 的基，转置后同理得到 $\mathcal R(A^{\mathsf T})$ 的基。又由 $\|Ax\|_2^2=\sum_{i=1}^r\sigma_i^2(v_i^{\mathsf T}x)^2$ 可知，$Ax=0$ 当且仅当 $v_i^{\mathsf T}x=0$ 对所有 $i\leq r$ 成立，从而确定 $\mathcal N(A)$；对 $A^{\mathsf T}$ 重复同一论证即可确定另一个零空间。
\end{proof}

\emph{奇异值唯一，奇异向量一般不唯一。}排序后的奇异值由 $A^{\mathsf T}A$ 的特征值唯一确定；每一对正奇异值对应的向量同时变号，不改变 $\sigma_i u_iv_i^{\mathsf T}$。若同一正奇异值 $\sigma$ 的重数大于 $1$，在对应的左右子空间中同时使用相同的正交矩阵 $Q$ 换基，也不改变乘积，因为 $(U_*Q)(\sigma I)(V_*Q)^{\mathsf T}=\sigma U_*V_*^{\mathsf T}$。零空间的左右基则可分别任意正交换基。因此不能把一次数值计算得到的奇异向量符号或重根子空间中的具体基看作唯一答案。

与对称矩阵的正交对角化相比，SVD 允许矩形矩阵，并在定义域和值域分别选基；对称矩阵的分解 $A=Q\Lambda Q^{\mathsf T}$ 则在同一个空间中使用同一组基。若 $A$ 为实对称矩阵，把其特征值按绝对值降序排列，令 $\Sigma$ 的对角元为这些绝对值，并把负特征值的符号吸收到相应左奇异向量中，就得到一个 SVD；零特征值处可任选符号。故对称矩阵的奇异值是特征值的绝对值，只有半正定时才可直接令 $U=V=Q$、$\Sigma=\Lambda$。

\section{奇异值分解的计算}
\label{sec:lh-svd-computation}

定理~\ref{thm:lh-svd-existence} 已给出一种通过特征分解构造 SVD 的方法，适于手算小矩阵和理解原理。其输入为实矩阵 $A$，输出为满足式~\eqref{eq:lh-full-svd} 的 $U,\Sigma,V$；若只需紧分解，则不必求零空间的全部基。在精确算术下可以按以下有限步骤完成。
\begin{enumerate}
\item 计算 $G=A^{\mathsf T}A$，求其全部特征值，并按降序排列为 $\lambda_1,\ldots,\lambda_n$；取相应的标准正交特征向量 $v_1,\ldots,v_n$，组成 $V$。重复特征值的特征子空间内也要选取标准正交基，单纯把任意一组特征向量分别单位化不足以保证正交。
\item 以正特征值的个数确定 $r$，取 $\sigma_i=\sqrt{\lambda_i}$（$1\leq i\leq p$），构造 $m\times n$ 的矩形对角矩阵 $\Sigma$。当 $m<n$ 时，余下的 $n-m$ 个特征值必为零。
\item 对 $1\leq i\leq r$，计算 $u_i=Av_i/\sigma_i$；求解 $A^{\mathsf T}y=0$，在其解空间中取标准正交基补足 $u_{r+1},\ldots,u_m$，组成 $U$。若 $r=0$，直接取任意正交矩阵 $U$。
\item 返回 $U,\Sigma,V$；可用 $U^{\mathsf T}U=I_m$、$V^{\mathsf T}V=I_n$ 与 $U\Sigma V^{\mathsf T}=A$ 核对结果。只需紧分解时，返回前 $r$ 列和 $r$ 阶正对角块；若只需秩 $k$ 近似，则返回相应前 $k$ 项。
\end{enumerate}

\begin{example}
\label{ex:lh-svd-rank-one-computation}
求教材例 15.5 中矩阵 $C=\begin{pmatrix}1&1\\2&2\\0&0\end{pmatrix}$ 的完全奇异值分解。
\end{example}
\begin{solve}
有 $C^{\mathsf T}C=\begin{pmatrix}5&5\\5&5\end{pmatrix}$，其特征值为 $10,0$，可取标准正交特征向量
\[
v_1=\frac1{\sqrt2}(1,1)^{\mathsf T},\qquad
v_2=\frac1{\sqrt2}(1,-1)^{\mathsf T}.
\]
于是 $\sigma_1=\sqrt{10}$、$\sigma_2=0$，且 $u_1=Cv_1/\sqrt{10}=(1,2,0)^{\mathsf T}/\sqrt5$。零空间 $\mathcal N(C^{\mathsf T})$ 由方程 $y_1+2y_2=0$ 确定，可取标准正交基 $u_2=(-2,1,0)^{\mathsf T}/\sqrt5$、$u_3=(0,0,1)^{\mathsf T}$。因此
\[
U=\begin{pmatrix}
1/\sqrt5&-2/\sqrt5&0\\
2/\sqrt5&1/\sqrt5&0\\
0&0&1
\end{pmatrix},\quad
\Sigma=\begin{pmatrix}\sqrt{10}&0\\0&0\\0&0\end{pmatrix},\quad
V=\frac1{\sqrt2}\begin{pmatrix}1&1\\1&-1\end{pmatrix}.
\]
直接相乘得到 $C=U\Sigma V^{\mathsf T}$。由于只有一个正奇异值，紧分解进一步写成 $C=\sqrt{10}\,u_1v_1^{\mathsf T}$，不需要保留补入的三个零空间方向。
\end{solve}

上面的特征分解方法是原理性的构造。对 $m\geq n$ 的稠密矩阵，显式形成 $A^{\mathsf T}A$ 约需 $O(mn^2)$ 运算，其稠密特征分解约需 $O(n^3)$ 运算；构造全部正奇异值对应的左向量还需 $O(mnr)$，这些估计未计入补齐完整 $U$ 的工作。当 $m<n$ 时，也可先对较小的 $AA^{\mathsf T}$ 求特征分解，再用 $v_i=A^{\mathsf T}u_i/\sigma_i$ 构造正奇异值对应的右向量。

数值计算中不宜据此直接形成 Gram 矩阵后再求全部特征值。例如 $A$ 列满秩时，若以最大、最小奇异值之比定义条件数 $\kappa_2(A)=\sigma_1/\sigma_n$，则 $\kappa_2(A^{\mathsf T}A)=\kappa_2(A)^2$，小奇异值的信息更容易受到舍入误差影响。常用数值 SVD 通过正交变换先把原矩阵化为双对角形式，再求相应分解，避免显式构造 $A^{\mathsf T}A$。浮点计算还必须用与数据尺度、精度相称的容差判断数值秩，而不能简单把计算结果是否严格等于零作为标准；数值秩随容差变化，应与精确代数秩区别。教材中的手算过程不提供某一种迭代算法的停止准则或误差保证。

\section{奇异值分解与矩阵近似}
\label{sec:lh-svd-approximation}

\subsection{弗罗贝尼乌斯范数}
\label{subsec:lh-svd-frobenius}

\begin{definition}[弗罗贝尼乌斯范数]
\label{def:lh-frobenius-norm}
对实矩阵 $A=(a_{ij})\in\mathbb R^{m\times n}$，其\emph{弗罗贝尼乌斯范数}（Frobenius norm）定义为
\begin{equation}
\|A\|_{\mathrm F}=\left(\sum_{i=1}^m\sum_{j=1}^n a_{ij}^2\right)^{1/2}.
\label{eq:lh-frobenius-definition}
\end{equation}
相应内积为 $\langle B,C\rangle_{\mathrm F}=\sum_{i,j}b_{ij}c_{ij}=\operatorname{tr}(B^{\mathsf T}C)$，其中 $B,C$ 同型，$\operatorname{tr}$ 表示方阵对角元之和。
\end{definition}

把矩阵所有元素依次排列成向量，式~\eqref{eq:lh-frobenius-definition} 就是该向量的欧氏范数。因此 $\|A-B\|_{\mathrm F}^2$ 是逐元素平方误差的总和，$\|A-B\|_{\mathrm F}$ 则是其平方根；二者具有相同的最小化解，但数值和单位并不相同。这里所有元素具有相同权重，并假定都已观测；若只在部分位置计算误差，或对不同位置加权，后面的最优近似结论不能直接套用。

\begin{lemma}[正交不变性与奇异值表示]
\label{lem:lh-frobenius-orthogonal-invariance}
若 $Q\in\mathbb R^{m\times m}$、$P\in\mathbb R^{n\times n}$ 为正交矩阵，则
\begin{equation}
\|QAP^{\mathsf T}\|_{\mathrm F}=\|A\|_{\mathrm F},\qquad
\|A\|_{\mathrm F}^2=\sum_{i=1}^p\sigma_i^2=\sum_{i=1}^r\sigma_i^2.
\label{eq:lh-frobenius-singular-values}
\end{equation}
\end{lemma}
\begin{proof}
设 $A$ 的列为 $a_1,\ldots,a_n$，则 $\|QA\|_{\mathrm F}^2=\sum_j\|Qa_j\|_2^2=\sum_j\|a_j\|_2^2=\|A\|_{\mathrm F}^2$。对各行或对转置矩阵使用同一论证，得到右乘 $P^{\mathsf T}$ 也不改变范数。将 $A=U\Sigma V^{\mathsf T}$ 代入，便有 $\|A\|_{\mathrm F}^2=\|\Sigma\|_{\mathrm F}^2=\sum_{i=1}^p\sigma_i^2$；零奇异值不贡献平方和。
\end{proof}

\subsection{矩阵的最优近似}
\label{subsec:lh-svd-best-approximation}

低秩近似的目标是，在限制模型秩的同时尽量减小原矩阵各元素的重构误差。对给定整数 $0\leq k\leq r$，考虑
\begin{equation}
\min_{B\in\mathbb R^{m\times n},\ \operatorname{rank}(B)\leq k}
\|A-B\|_{\mathrm F}^2.
\label{eq:lh-svd-low-rank-objective}
\end{equation}
这里优化变量是整个矩阵 $B$，秩是约束，SVD 是求解这个目标的工具；不能仅从“奇异值按大小排列”就认定最优性，还需证明任何其他秩不超过 $k$ 的矩阵都不可能取得更小误差。

\begin{theorem}[截断 SVD 的最优性]
\label{thm:lh-svd-best-approximation}
设 $A$ 的秩为 $r$，$A_k$ 按式~\eqref{eq:lh-truncated-svd} 及 $A_0=0$、$A_r=A$ 定义。则对每个整数 $0\leq k\leq r$，式~\eqref{eq:lh-svd-low-rank-objective} 的最小值均可达到，且
\begin{equation}
\min_{\operatorname{rank}(B)\leq k}\|A-B\|_{\mathrm F}^2
=\|A-A_k\|_{\mathrm F}^2
=\sum_{i=k+1}^r\sigma_i^2.
\label{eq:lh-svd-optimal-tail-error}
\end{equation}
相应的最小范数误差为 $\bigl(\sum_{i=k+1}^r\sigma_i^2\bigr)^{1/2}$。若 $1\leq k<r$，则最优解的秩恰为 $k$；若还满足 $\sigma_k>\sigma_{k+1}$，最优近似矩阵唯一。
\end{theorem}
\begin{proof}
当 $k=0$ 时可行矩阵只有零矩阵，当 $k=r$ 时 $A$ 本身可行，两种情形直接成立。以下设 $1\leq k<r$。任取可行矩阵 $B$，设 $P_B$ 是到其列空间的正交投影，$d=\operatorname{rank}(B)\leq k$。分解
\[
A-B=(I_m-P_B)A+(P_BA-B).
\]
等式右端两个矩阵的每一对对应列相互正交，因此
\begin{equation}
\|A-B\|_{\mathrm F}^2
=\|(I_m-P_B)A\|_{\mathrm F}^2+\|P_BA-B\|_{\mathrm F}^2
\geq\|A\|_{\mathrm F}^2-\|P_BA\|_{\mathrm F}^2.
\label{eq:lh-svd-projection-lower-bound}
\end{equation}
由紧分解及右侧正交变换的不变性，令 $t_i=\|P_Bu_i\|_2^2$，可得
\[
\|P_BA\|_{\mathrm F}^2=\sum_{i=1}^r\sigma_i^2t_i,\qquad
0\leq t_i\leq1,\qquad
\sum_{i=1}^r t_i\leq\sum_{i=1}^m\|P_Bu_i\|_2^2=d\leq k.
\]
最后一个等式可由 $\|P_BU\|_{\mathrm F}^2=\|P_B\|_{\mathrm F}^2=d$ 得到；正交投影的特征值只有 $0,1$，其中有 $d$ 个 $1$。利用奇异值的降序排列，进一步有
\begin{align*}
\sum_{i=1}^r\sigma_i^2t_i
&\leq\sum_{i=1}^k\sigma_i^2t_i+
\sigma_k^2\sum_{i=k+1}^r t_i\\
&\leq\sum_{i=1}^k\sigma_i^2t_i+
\sigma_k^2\left(k-\sum_{i=1}^k t_i\right)\\
&=\sum_{i=1}^k\bigl[\sigma_k^2+(\sigma_i^2-\sigma_k^2)t_i\bigr]
\leq\sum_{i=1}^k\sigma_i^2.
\end{align*}
代入式~\eqref{eq:lh-svd-projection-lower-bound} 及引理~\ref{lem:lh-frobenius-orthogonal-invariance}，得到任意可行矩阵都满足 $\|A-B\|_{\mathrm F}^2\geq\sum_{i=k+1}^r\sigma_i^2$。另一方面，$A_k$ 的秩为 $k$，且 $A-A_k$ 在奇异坐标中仅有尾部奇异值，故其误差恰为该平方和，证明了最优解的存在性和最优值公式。秩小于 $k$ 时，同一结论在其实际秩处给出的下界严格更大，因为前 $k$ 个奇异值均为正，所以最优解必须有秩 $k$。

若 $\sigma_k>\sigma_{k+1}$，上述加权和达到上界时必须有 $t_1=\cdots=t_k=1$，且 $t_{k+1}=\cdots=t_r=0$。理由是只要前 $k$ 个位置尚有未分配的权重，把尾部权重移到这些位置，或在总权重不足 $k$ 时补入权重，都会严格增大加权和。于是最优矩阵的列空间只能是 $\operatorname{span}\{u_1,\ldots,u_k\}$；式~\eqref{eq:lh-svd-projection-lower-bound} 取等又要求 $B=P_BA$，所以 $B=U_kU_k^{\mathsf T}A=A_k$，唯一性得证。
\end{proof}

严格的间隙条件约束的是截断位置，并不要求前 $k$ 个奇异值各不相同。若 $\sigma_k=\sigma_{k+1}>0$，可以在跨越截断位置的重根子空间内选择不同方向，得到不同的最优矩阵；例如 $A=I_2$、$k=1$ 时，每个单位向量 $q\in\mathbb R^2$ 给出的 $qq^{\mathsf T}$ 都是最优近似，误差平方均为 $1$。这也说明，奇异向量不唯一与最优近似矩阵不唯一需要分别判断。

\begin{example}
\label{ex:lh-svd-optimal-error}
续例~\ref{ex:lh-svd-five-by-four}，求秩至多为 $1$、$2$ 的最优近似误差，并判断式~\eqref{eq:lh-svd-five-by-four-rank-two} 是否为唯一的秩至多为 $2$ 的最优近似。
\end{example}
\begin{solve}
三个正奇异值为 $4,3,\sqrt5$，故 $\|A\|_{\mathrm F}^2=16+9+5=30$。秩至多为 $1$ 时，最小误差平方为 $9+5=14$，最小范数误差为 $\sqrt{14}$；秩至多为 $2$ 时，二者分别为 $5$ 与 $\sqrt5$。因为 $\sigma_2=3>\sqrt5=\sigma_3$，式~\eqref{eq:lh-svd-five-by-four-rank-two} 是唯一的最优秩 $2$ 近似。这正是教材例 15.6 所求的近似；若将 $\sqrt5$ 误写成误差平方，就混淆了范数与平方损失。
\end{solve}

\subsection{矩阵的外积展开式}
\label{subsec:lh-svd-outer-products}

\begin{proposition}[外积展开]
\label{prop:lh-svd-outer-product-expansion}
矩阵 $A$ 及其截断近似可写为
\begin{equation}
A=\sum_{i=1}^r\sigma_i u_iv_i^{\mathsf T},\qquad
A_k=\sum_{i=1}^k\sigma_i u_iv_i^{\mathsf T}.
\label{eq:lh-svd-outer-products}
\end{equation}
每个 $u_iv_i^{\mathsf T}$ 是秩为 $1$ 的 $m\times n$ 矩阵，它们在弗罗贝尼乌斯内积下标准正交。
\end{proposition}
\begin{proof}
当 $r=0$ 时以空和表示零矩阵；当 $r\geq1$ 时，把 $U_r\Sigma_r$ 按列、$V_r^{\mathsf T}$ 按行分块相乘，即得两个求和式。\emph{外积}（outer product）$u_iv_i^{\mathsf T}$ 的各列是 $u_i$ 的标量倍数，且两向量都非零，故其秩为 $1$。又有
\[
\langle u_iv_i^{\mathsf T},u_jv_j^{\mathsf T}\rangle_{\mathrm F}
=(u_i^{\mathsf T}u_j)(v_i^{\mathsf T}v_j)
=\begin{cases}1,&i=j,\\0,&i\ne j,\end{cases}
\]
因此这些秩一矩阵相互正交且范数均为 $1$。
\end{proof}

外积 $u_iv_i^{\mathsf T}$ 的第 $s$ 行第 $t$ 列元素为 $(u_i)_s(v_i)_t$，所以它同时描述一个输出方向与一个输入方向之间的配合。一般而言，任意非零秩一矩阵也能写成两个非零向量的外积：取其非零列 $a$，其余各列均为 $a$ 的倍数，记这些倍数组成向量 $b$，就得到 $ab^{\mathsf T}$。例如例~\ref{ex:lh-svd-rank-one-computation} 中的矩阵为 $(1,2,0)^{\mathsf T}(1,1)$；SVD 进一步把两向量单位化，并用非负奇异值记录整体尺度。这给出了教材习题 15.4 所要求的秩一表示及实例。

由外积的正交性，每加入一项 $\sigma_i u_iv_i^{\mathsf T}$，重构误差平方恰好减少 $\sigma_i^2$。若 $A\ne0$，可定义保留的\emph{能量比例}
\begin{equation}
\rho_k=\frac{\sum_{i=1}^k\sigma_i^2}{\sum_{i=1}^r\sigma_i^2},\qquad
\frac{\|A-A_k\|_{\mathrm F}}{\|A\|_{\mathrm F}}=\sqrt{1-\rho_k}.
\label{eq:lh-svd-retained-energy}
\end{equation}
给定目标比例 $\rho\in(0,1]$，可以选取满足 $\rho_k\geq\rho$ 的最小整数 $k$；给定相对范数误差容许值 $\varepsilon\in[0,1)$，则选择满足 $\sum_{i=k+1}^r\sigma_i^2\leq\varepsilon^2\sum_{i=1}^r\sigma_i^2$ 的 $k$。例~\ref{ex:lh-svd-five-by-four} 中，保留两项的能量比例为 $25/30=5/6$，相对范数误差为 $1/\sqrt6$。奇异值是否快速衰减取决于数据；只有尾部平方和较小时，小 $k$ 才能得到高精度近似。能量比例描述本矩阵的重构程度，本身不保证分类准确率或预测效果。

教材习题 15.5 用查询与网页之间的点击关系说明矩阵分解的含义。一般地，令 $A\in\mathbb R^{M\times N}$ 的行对应 $M$ 条查询、列对应 $N$ 个网页，$a_{ij}$ 为对应点击次数。截断分解可以把查询表示为 $Z=U_k\Sigma_k\in\mathbb R^{M\times k}$，把网页表示为 $H=V_k\in\mathbb R^{N\times k}$，从而
\[
a_{ij}\approx (A_k)_{ij}=Z_{i,:}H_{j,:}^{\mathsf T}.
\]
这里 $Z_{i,:}$、$H_{j,:}$ 分别表示矩阵的第 $i$ 行、第 $j$ 行。也可将尺度对称分配，使用 $U_k\Sigma_k^{1/2}$ 与 $V_k\Sigma_k^{1/2}$ 作为两侧表示。奇异值刻画各潜在方向对重构的贡献，但方向中的坐标可为负，近似矩阵也未必逐元素非负，不能直接把它们解释为概率或唯一的语义主题。SVD 本身不要求数据中心化；若用于主成分分析，则还需结合该模型对中心化与样本协方差的定义说明其关系。

% 教材来源：机器学习/textbooks/机器学习（分章原文）/机器学习-16.pdf。
% 对应教材第 16 章，原书第 297--320 页；沿用本文件的自动章节编号。

\fi

\selectchapter{16}
\ifcompilechapter
\chapter{主成分分析}
\label{ch:lh-pca}

当多个观测变量含有相近的变化趋势时，直接保留全部变量会掩盖数据的主要结构。\emph{主成分分析}（principal component analysis，PCA）用一组正交方向重新表达数据，并依次选取投影方差最大的方向，以少量线性组合概括原变量的主要变化。本章中的“信息保留”特指\emph{方差的保留}，并非信息论中的熵，也不保证保留所有与分类、预测有关的信息。其核心问题是：在给定维数的线性子空间中，怎样使投影后的总方差最大，等价地，使重构的均方误差最小？

本章按教材的总体主成分分析与样本主成分分析两部分展开。前者以随机向量的均值和协方差为已知量，给出主成分的定义、最优化性质和解释方法；后者用有限数据估计这些量，再通过特征值分解或奇异值分解计算主成分。为区分随机对象与数据，以下用 $X=(X_1,\ldots,X_m)^{\mathsf T}$ 表示随机向量，$x_j\in\mathbb R^m$ 表示第 $j$ 个观测值，$\mathbf X=(x_1,\ldots,x_n)\in\mathbb R^{m\times n}$ 表示以样本为列的数据矩阵。$I_d$ 表示 $d$ 阶单位矩阵，$\operatorname{tr}$ 表示矩阵的迹，$\operatorname{Cov}$ 与 $\operatorname{Var}$ 分别表示协方差与方差。

\section{总体主成分分析}
\label{sec:lh-pca-population}

\subsection{基本想法}
\label{subsec:lh-pca-idea}

设 $\mathbb E\|X\|_2^2<\infty$，记均值向量与协方差矩阵为
\begin{equation}
\mu=\mathbb E X,\qquad
\Sigma=\mathbb E\bigl[(X-\mu)(X-\mu)^{\mathsf T}\bigr]
=(\sigma_{ij})_{m\times m}.
\label{eq:lh-pca-population-cov}
\end{equation}
矩阵 $\Sigma$ 对称半正定，因为对任意 $a\in\mathbb R^m$，$a^{\mathsf T}\Sigma a=\operatorname{Var}(a^{\mathsf T}X)\geq0$。将数据中心移到原点后，沿单位方向 $a$ 的坐标为 $a^{\mathsf T}(X-\mu)$，其方差就是这个二次型。若不约束 $\|a\|_2$，把 $a$ 放大即可任意放大正方差，因而单位长度约束是比较方向所必需的。

以二维数据为例，若点云大致沿某条直线延展，则沿这条直线的投影变化大，沿垂直方向的变化小。把前者作为第一坐标轴、后者作为第二坐标轴，便把主要变化与较小变化分开；只保留第一坐标，就得到一维近似。对于任意单位向量 $a$ 和中心化的观测值 $z$，勾股分解给出
\begin{equation}
\|z\|_2^2=(a^{\mathsf T}z)^2
+\|z-aa^{\mathsf T}z\|_2^2.
\label{eq:lh-pca-pythagoras}
\end{equation}
因此，在样本已经中心化时，最大化投影坐标的平方和等价于最小化点到该直线的距离平方和；取期望后，同一关系把投影方差最大化与均方重构误差最小化联系起来。单个坐标的平方只是对方差的一项贡献，方差还须取期望或按相应分母作样本汇总。

\emph[caution]{保留全部坐标的正交变换不改变总方差；PCA 优化的是前若干个坐标保留的方差。}在 $m$ 维空间中使用完整的正交基，仅仅改变坐标表示，没有降维；只有舍弃部分坐标后，所选方向才决定损失多少方差。中心化是这一直观与方差相对应的前提，而把各变量进一步缩放到单位方差是另一项选择，将在第~\ref{subsec:lh-pca-standardized-population} 小节讨论。

\subsection{定义和导出}
\label{subsec:lh-pca-definition}

对于 $Y_k=a_k^{\mathsf T}(X-\mu)$，有
\begin{equation}
\mathbb E Y_k=0,\qquad
\operatorname{Var}(Y_k)=a_k^{\mathsf T}\Sigma a_k,\qquad
\operatorname{Cov}(Y_j,Y_k)=a_j^{\mathsf T}\Sigma a_k.
\label{eq:lh-pca-linear-cov}
\end{equation}
教材也使用不中心化的形式 $a_k^{\mathsf T}X$，其均值为 $a_k^{\mathsf T}\mu$，与这里的形式只相差常数，因而求出的方向、方差与协方差相同。以下统一使用中心化的主成分，使降维和重构公式中的平移明确可见。

\begin{definition}[总体主成分]
\label{def:lh-pca-population}
对满足式~\eqref{eq:lh-pca-population-cov} 的随机向量 $X$，依次选择单位向量 $a_1,\ldots,a_m$，使
\begin{equation}
a_k\in\mathop{\arg\max}_{a\in\mathbb R^m}
\left\{a^{\mathsf T}\Sigma a:
a^{\mathsf T}a=1,\quad a^{\mathsf T}a_j=0\ (1\leq j<k)\right\}.
\label{eq:lh-pca-sequential}
\end{equation}
称 $a_k$ 为第 $k$ 个\emph{主成分方向}，称随机变量 $Y_k=a_k^{\mathsf T}(X-\mu)$ 为第 $k$ 个\emph{总体主成分}（population principal component）。记 $A=(a_1,\ldots,a_m)$，则 $A^{\mathsf T}A=AA^{\mathsf T}=I_m$，完整主成分向量为 $Y=A^{\mathsf T}(X-\mu)$。
\end{definition}

这里显式要求方向两两正交，以涵盖奇异协方差矩阵的情况。教材以“单位系数且主成分互不相关”逐次定义：当已选方向满足 $\Sigma a_j=\lambda_j a_j$ 且 $\lambda_j>0$ 时，$\operatorname{Cov}(a^{\mathsf T}X,a_j^{\mathsf T}X)=\lambda_j a^{\mathsf T}a_j=0$ 正好等价于方向正交；若 $\lambda_j=0$，不相关约束自动成立，便不能据此推出正交。这时仍可在零特征空间内补成标准正交基。\emph[caution]{各系数向量长度为 $1$ 本身并不保证它们构成正交矩阵。}

\begin{theorem}[主成分与协方差矩阵的谱分解]
\label{thm:lh-pca-eigen}
设 $\Sigma$ 的特征值按 $\lambda_1\geq\cdots\geq\lambda_m\geq0$ 排列，$a_1,\ldots,a_m$ 是对应的一组标准正交特征向量，则它们满足定义~\ref{def:lh-pca-population}，且
\begin{equation}
\operatorname{Var}(Y_k)=\lambda_k,\qquad
\Sigma=A\Lambda A^{\mathsf T},\qquad
\Lambda=\operatorname{diag}(\lambda_1,\ldots,\lambda_m).
\label{eq:lh-pca-spectral}
\end{equation}
反之，逐次最优化得到的方向可组成这样一组按特征值排序的标准正交特征向量。
\end{theorem}
\begin{proof}
由实对称矩阵的谱定理，存在上述标准正交特征向量。将任意单位向量写成 $a=\sum_{j=1}^m c_j a_j$，则
\[
\sum_{j=1}^m c_j^2=1,\qquad
a^{\mathsf T}\Sigma a=\sum_{j=1}^m\lambda_j c_j^2
\leq\lambda_1.
\]
取 $a=a_1$ 达到上界，故它给出第一主成分。若还要求 $a\perp a_1,\ldots,a_{k-1}$，则 $c_1=\cdots=c_{k-1}=0$，上式的最大值变为 $\lambda_k$，取 $a=a_k$ 达到。等号成立时，$a$ 只能在尚未选取的、特征值为 $\lambda_k$ 的特征子空间内有非零分量，因此也是对应特征向量；逐次应用即得反向结论。最后，$a_k^{\mathsf T}\Sigma a_k=\lambda_k$，而 $\Sigma A=A\Lambda$ 与 $AA^{\mathsf T}=I_m$ 给出谱分解。
\end{proof}

从约束优化角度看，第一步的拉格朗日函数为 $L(a,\lambda)=a^{\mathsf T}\Sigma a-\lambda(a^{\mathsf T}a-1)$，驻点满足 $\Sigma a=\lambda a$。这说明候选方向是特征向量；要断定最大值对应最大特征值，还需要证明中的谱展开比较，而不能仅靠驻点方程。后续步骤是在已选方向的正交补中重复同一比较。

方向的符号没有唯一性：$a_k$ 与 $-a_k$ 都有效，只使主成分坐标整体变号。若特征值重合，相应特征子空间内任意标准正交基都可采用。对 $1\leq q<m$，若 $\lambda_q>\lambda_{q+1}$，前 $q$ 个方向张成的子空间唯一，但其中的基仍可能因重特征值而不唯一；若截断位置满足 $\lambda_q=\lambda_{q+1}$，最优 $q$ 维子空间也可能不唯一。零特征值对应的中心化主成分方差为零，因而几乎处处等于零，不携带随机变化。

\subsection{主要性质}
\label{subsec:lh-pca-properties}

由式~\eqref{eq:lh-pca-spectral} 可直接求出完整主成分的协方差和总方差：
\begin{equation}
\operatorname{Cov}(Y)=A^{\mathsf T}\Sigma A=\Lambda,\qquad
\sum_{k=1}^m\operatorname{Var}(Y_k)
=\sum_{k=1}^m\lambda_k
=\operatorname{tr}(\Sigma)=\sum_{i=1}^m\sigma_{ii}.
\label{eq:lh-pca-trace}
\end{equation}
后一等式也可由 $\operatorname{tr}(A\Lambda A^{\mathsf T})=\operatorname{tr}(\Lambda A^{\mathsf T}A)$ 得到。由于 $X-\mu=AY$，完整变换保留全部数据；前 $q$ 个主成分只保留这一展开的前 $q$ 项。

主成分之间的零协方差表示\emph{互不相关}（uncorrelated），一般不表示\emph{相互独立}（independent）。例如，令 $U$ 以相同概率取 $-1,0,1$，取 $X_1=U$、$X_2=U^2-2/3$，则两者均值为零，方差分别为 $2/3$、$2/9$，且 $\operatorname{Cov}(X_1,X_2)=\mathbb E U^3=0$。协方差矩阵已经对角化，它们本身即可作为主成分，但 $X_2$ 由 $X_1$ 决定，且两者都不是常量，所以并不独立。若 $X$ 联合服从多元正态分布，线性变换后仍联合正态，此时对角协方差才可推出各主成分相互独立。

\begin{definition}[因子负荷量]
\label{def:lh-pca-loading}
当 $\sigma_{ii}>0$ 且 $\lambda_k>0$ 时，称主成分 $Y_k$ 与原变量 $X_i$ 的相关系数为\emph{因子负荷量}（factor loading），记为
\begin{equation}
\ell_{ik}=\operatorname{Corr}(X_i,Y_k)
=\frac{\sqrt{\lambda_k}\,a_{ik}}{\sqrt{\sigma_{ii}}},
\qquad a_k=(a_{1k},\ldots,a_{mk})^{\mathsf T}.
\label{eq:lh-pca-loading}
\end{equation}
\end{definition}

为验证这一公式，令 $e_i$ 为第 $i$ 个标准基向量，则 $\operatorname{Cov}(X_i,Y_k)=e_i^{\mathsf T}\Sigma a_k=\lambda_k a_{ik}$；除以两个标准差即可。方向系数 $a_{ik}$ 与相关系数 $\ell_{ik}$ 是不同的量，不能混用；“负荷”在不同文献中也可能指其他经过缩放的系数，本章沿用教材的相关系数定义。若原变量或主成分方差为零，其相关系数无定义，不把零方差主成分代入式~\eqref{eq:lh-pca-loading}。

若各 $\sigma_{ii}>0$，并令 $r=\operatorname{rank}(\Sigma)$，则正方差主成分满足
\begin{equation}
\sum_{i=1}^m\sigma_{ii}\ell_{ik}^2=\lambda_k
\quad(1\leq k\leq r),\qquad
\sum_{k=1}^r\ell_{ik}^2=1
\quad(1\leq i\leq m).
\label{eq:lh-pca-loading-identities}
\end{equation}
第一式使用 $\sum_i a_{ik}^2=1$；第二式使用 $\sigma_{ii}=\sum_{k=1}^r\lambda_k a_{ik}^2$。这些恒等式说明，全部非零主成分可以解释每个变量的全部方差，但保留其中一部分时，不同变量的保留程度可能相差很大。

\subsection{主成分的个数}
\label{subsec:lh-pca-component-count}

逐次最大化方差还需回答一个整体问题：前 $q$ 个主成分是否在所有 $q$ 维正交投影中保留最多方差？下面的结论同时解释应保留哪些方向、应舍弃哪些方向。

\begin{theorem}[给定维数下的方差最优化]
\label{thm:lh-pca-variance-optimal}
设 $\Sigma$ 的特征值和方向如定理~\ref{thm:lh-pca-eigen}，$1\leq q\leq m$。对任意 $B\in\mathbb R^{m\times q}$ 且 $B^{\mathsf T}B=I_q$，有
\begin{equation}
\sum_{j=m-q+1}^m\lambda_j
\leq\operatorname{tr}(B^{\mathsf T}\Sigma B)
\leq\sum_{j=1}^q\lambda_j.
\label{eq:lh-pca-variance-bounds}
\end{equation}
取最后 $q$ 个特征向量达到左端，取 $A_q=(a_1,\ldots,a_q)$ 达到右端。
\end{theorem}
\begin{proof}
令 $C=A^{\mathsf T}B$，记第 $j$ 行的平方范数为 $w_j=\sum_{k=1}^q c_{jk}^2$。由 $C^{\mathsf T}C=I_q$，有 $\sum_jw_j=q$；又 $w_j=\|B^{\mathsf T}a_j\|_2^2\leq\|a_j\|_2^2=1$，故 $0\leq w_j\leq1$。于是
\[
\operatorname{tr}(B^{\mathsf T}\Sigma B)=\sum_{j=1}^m\lambda_jw_j.
\]
当 $q<m$ 时，由 $\sum_{j\leq q}(1-w_j)=\sum_{j>q}w_j$，得到
\[
\sum_{j=1}^q\lambda_j-\sum_{j=1}^m\lambda_jw_j
\geq\lambda_q\sum_{j\leq q}(1-w_j)
-\lambda_q\sum_{j>q}w_j=0.
\]
这一上界论证只用到了特征值的排序，故也可应用于 $-\Sigma$，其降序特征值为 $-\lambda_m,\ldots,-\lambda_1$，得到 $-\operatorname{tr}(B^{\mathsf T}\Sigma B)\leq-\sum_{j=m-q+1}^m\lambda_j$，即所需下界。$q=m$ 时两个界均等于 $\operatorname{tr}(\Sigma)$。取相应特征向量时，$w_j$ 在所选位置为 $1$、其余为 $0$，从而达到相应界。
\end{proof}

对中心化向量的 $q$ 维投影重构为 $\widehat X_q=\mu+A_qA_q^{\mathsf T}(X-\mu)$。更一般地，对任意上述 $B$，$P=BB^{\mathsf T}$ 是正交投影矩阵，满足 $P^{\mathsf T}=P$、$P^2=P$，因此
\begin{align}
\mathbb E\|X-\mu-BB^{\mathsf T}(X-\mu)\|_2^2
&=\operatorname{tr}\bigl((I_m-P)\Sigma\bigr)\notag\\
&=\operatorname{tr}(\Sigma)-\operatorname{tr}(B^{\mathsf T}\Sigma B).
\label{eq:lh-pca-population-reconstruction}
\end{align}
这里使用 $\mathbb E\|M(X-\mu)\|_2^2=\operatorname{tr}(M\Sigma M^{\mathsf T})$ 和迹的循环性质，所需期望因有限二阶矩而存在。由定理~\ref{thm:lh-pca-variance-optimal}，PCA 使这一误差最小，最小值为 $\sum_{j=q+1}^m\lambda_j$。所以最大投影方差与最小均方重构误差是同一优化问题的两种表达；该结论针对欧氏距离下的正交投影，并不声称任意非线性压缩都不能更好。

\begin{definition}[方差贡献率与累计方差贡献率]
\label{def:lh-pca-explained-variance}
若 $\operatorname{tr}(\Sigma)>0$，第 $k$ 个主成分的\emph{方差贡献率}（explained variance ratio）与前 $q$ 个主成分的\emph{累计方差贡献率}分别为
\begin{equation}
\eta_k=\frac{\lambda_k}{\sum_{j=1}^m\lambda_j},\qquad
C_q=\sum_{k=1}^q\eta_k
=\frac{\sum_{k=1}^q\lambda_k}{\operatorname{tr}(\Sigma)}.
\label{eq:lh-pca-explained-variance}
\end{equation}
\end{definition}

给定阈值 $\tau\in(0,1]$，可取满足 $C_q\geq\tau$ 的最小 $q$。教材举出的 $70\%$ 或 $80\%$ 是选择阈值的例子，并无普遍适用的固定标准。由式~\eqref{eq:lh-pca-population-reconstruction}，$1-C_q$ 正是最优重构误差占总方差的比例；若总方差为零，所有变量都几乎处处等于均值，贡献率的分母为零，应直接按常量数据处理。

累计贡献率概括全部变量，不能说明某一个变量保留得是否充分。对 $\sigma_{ii}>0$，定义前 $q$ 个主成分对变量 $X_i$ 的贡献率为
\begin{equation}
v_i(q)=\frac{\sum_{k=1}^q\lambda_k a_{ik}^2}{\sigma_{ii}}
=\sum_{\substack{1\leq k\leq q\\\lambda_k>0}}\ell_{ik}^2,
\qquad 0\leq v_i(q)\leq1.
\label{eq:lh-pca-variable-contribution}
\end{equation}
事实上，$\widehat X_{i,q}-\mu_i=\sum_{k=1}^qa_{ik}Y_k$，各 $Y_k$ 互不相关，所以重构变量的方差为分子，而残差方差为 $\sigma_{ii}(1-v_i(q))$。这些方向上的残差与已保留的主成分也互不相关，故该重构同时是用这些主成分作线性最小二乘预测的结果。因而 $v_i(q)$ 可解释为线性回归的决定系数，即复相关系数的平方；在 $q$ 包含零方差方向时，只使用正方差主成分定义复相关。这一解释涉及线性可解释方差，不声称等于条件均值所解释的全部非线性变化。

\subsection{规范化变量的总体主成分}
\label{subsec:lh-pca-standardized-population}

若变量的量纲或尺度不同，原始协方差会使方差较大的变量占据较大权重。假设各 $\sigma_{ii}>0$，令 $D=\operatorname{diag}(\sqrt{\sigma_{11}},\ldots,\sqrt{\sigma_{mm}})$，作\emph{标准化}（standardization，教材称“规范化”）
\begin{equation}
Z=D^{-1}(X-\mu),\qquad
Z_i=\frac{X_i-\mu_i}{\sqrt{\sigma_{ii}}},\qquad
R=\operatorname{Cov}(Z)=D^{-1}\Sigma D^{-1}.
\label{eq:lh-pca-standardization}
\end{equation}
$R$ 是原变量的相关矩阵，主对角元均为 $1$。对 $R$ 作主成分分析称为相关矩阵 PCA，对 $\Sigma$ 作分析称为协方差矩阵 PCA；两者优化的是不同尺度下的方差，结果通常不同。若某变量方差为零，不能除以其标准差，应先移除该常量变量，并在重构时补回其常量值。

设 $R$ 的特征值为 $\theta_1\geq\cdots\geq\theta_m\geq0$，对应标准正交特征向量为 $b_1,\ldots,b_m$，$b_k=(b_{1k},\ldots,b_{mk})^{\mathsf T}$，标准化后的主成分为 $W_k=b_k^{\mathsf T}Z$。将前述结论中的 $\Sigma$ 换成 $R$，便得到
\begin{equation}
\sum_{k=1}^m\theta_k=\operatorname{tr}(R)=m,\qquad
\eta_k^{(R)}=\frac{\theta_k}{m},\qquad
\operatorname{Corr}(Z_i,W_k)=\sqrt{\theta_k}\,b_{ik}
\quad(\theta_k>0).
\label{eq:lh-pca-standardized-properties}
\end{equation}
因此，对正特征值有 $\sum_i\operatorname{Corr}(Z_i,W_k)^2=\theta_k$，对每个变量有 $\sum_{k:\theta_k>0}\operatorname{Corr}(Z_i,W_k)^2=1$。由于正比例缩放和平移不改变相关系数，同一负荷量也等于 $\operatorname{Corr}(X_i,W_k)$。

标准化不只是消去量纲的记号变换。例如 $\Sigma=\operatorname{diag}(100,1)$ 时，协方差 PCA 唯一的第一主轴为第一个坐标方向，保留方差比例为 $100/101$；标准化后 $R=I_2$，所有单位方向的方差均为 $1$，任何第一主轴的贡献率都为 $1/2$。前者重视原尺度上的绝对变化，后者使每个变量的边际方差具有相同权重。若把标准化空间的重构变回原单位，它最小化的是 $\|D^{-1}(X-\widehat X)\|_2^2$ 的期望，而非一般意义下的原单位欧氏误差。因此是否标准化应由变量含义和分析目标决定，不能把两种结果直接当作同一个最优解。

\section{样本主成分分析}
\label{sec:lh-pca-sample}

\subsection{样本主成分的定义和性质}
\label{subsec:lh-pca-sample-definition}

总体协方差通常未知，需要用观测数据代替。给定 $n\geq2$ 个 $m$ 维样本 $x_1,\ldots,x_n$，仍以样本为列，定义
\begin{equation}
\mathbf X=(x_1,\ldots,x_n),\qquad
\bar x=\frac1n\sum_{j=1}^n x_j,\qquad
\mathbf X_c=\mathbf X-\bar x\mathbf 1_n^{\mathsf T},\qquad
S=\frac1{n-1}\mathbf X_c\mathbf X_c^{\mathsf T}.
\label{eq:lh-pca-sample-cov}
\end{equation}
这里 $\mathbf 1_n$ 为全部分量等于 $1$ 的 $n$ 维向量，$x_{ij}$ 为 $\mathbf X$ 第 $i$ 行第 $j$ 列元素，即第 $j$ 个样本的第 $i$ 个变量，$\bar x_i=n^{-1}\sum_jx_{ij}$。样本协方差矩阵 $S=(s_{i\ell})_{m\times m}$ 的元素为
\[
s_{i\ell}=\frac1{n-1}\sum_{j=1}^n(x_{ij}-\bar x_i)(x_{\ell j}-\bar x_\ell).
\]
对任意方向 $a,b\in\mathbb R^m$，样本得分 $a^{\mathsf T}x_j$ 的均值为 $a^{\mathsf T}\bar x$，样本方差为 $a^{\mathsf T}Sa$；两组得分的样本协方差为 $a^{\mathsf T}Sb$。这些都是给定数据后的代数等式，无须先假设样本服从正态分布。

\begin{definition}[样本主成分]
\label{def:lh-pca-sample}
将定义~\ref{def:lh-pca-population} 中的 $\Sigma$ 换成样本协方差矩阵 $S$，逐次最大化 $a^{\mathsf T}Sa$，所得标准正交方向称为\emph{样本主成分方向}。若其前 $q$ 个方向组成 $\widehat A_q$，则
\begin{equation}
\mathbf Y_q=\widehat A_q^{\mathsf T}\mathbf X_c\in\mathbb R^{q\times n}
\label{eq:lh-pca-sample-scores}
\end{equation}
称为主成分\emph{得分矩阵}（score matrix），第 $j$ 列为样本 $x_j$ 的 $q$ 维表示，第 $k$ 行为全部样本在第 $k$ 个主成分上的得分。
\end{definition}

系数向量描述投影规则，得分是把具体样本代入后得到的坐标；二者也不等于因子负荷量。由定理~\ref{thm:lh-pca-eigen}，$\widehat A_q$ 的列就是 $S$ 的前 $q$ 个标准正交特征向量。若相应特征值为 $\widehat\lambda_1\geq\cdots\geq\widehat\lambda_q$，则
\begin{equation}
\frac1{n-1}\mathbf Y_q\mathbf Y_q^{\mathsf T}
=\operatorname{diag}(\widehat\lambda_1,\ldots,\widehat\lambda_q).
\label{eq:lh-pca-score-cov}
\end{equation}
总体部分关于迹、方差贡献率、负荷量与重构误差的代数结论都可用 $S$ 和样本统计量重写，但样本得分互不相关并不保证用这些估计方向变换未来观测时，总体协方差也恰好对角化。

中心化矩阵满足 $\mathbf X_c\mathbf 1_n=0$，故
\begin{equation}
\operatorname{rank}(S)=\operatorname{rank}(\mathbf X_c)
\leq\min\{m,n-1\}.
\label{eq:lh-pca-sample-rank}
\end{equation}
因此当 $m\geq n$ 时，样本协方差必然奇异，至多只有 $n-1$ 个正方差主成分。若秩为 $r$，保留前 $r$ 个主成分就能精确重构当前中心化样本，但这不等于证明总体只具有 $r$ 个变化方向，也不意味着这些方向都能稳定推广到新数据。

\begin{proposition}[样本协方差的无偏性]
\label{prop:lh-pca-unbiased-cov}
设随机样本 $X^{(1)},\ldots,X^{(n)}$ 独立同分布，具有有限二阶矩、均值 $\mu$ 和协方差 $\Sigma$，以随机样本构造式~\eqref{eq:lh-pca-sample-cov} 中的 $S$，则 $\mathbb E S=\Sigma$。
\end{proposition}
\begin{proof}
令 $\bar X=n^{-1}\sum_jX^{(j)}$，对每组样本都有
\[
\sum_{j=1}^n(X^{(j)}-\bar X)(X^{(j)}-\bar X)^{\mathsf T}
=\sum_{j=1}^n(X^{(j)}-\mu)(X^{(j)}-\mu)^{\mathsf T}
-n(\bar X-\mu)(\bar X-\mu)^{\mathsf T}.
\]
由独立性，$\operatorname{Cov}(\bar X)=\Sigma/n$。取期望，右端为 $n\Sigma-n(\Sigma/n)=(n-1)\Sigma$，再除以 $n-1$ 即得结论。
\end{proof}

这个结论解释了样本协方差中的分母 $n-1$，也给出教材相应习题的证明。若统一改用 $n$ 作分母，协方差只乘以 $(n-1)/n$，因而方向与方差贡献率不变，但特征值及其与奇异值的换算必须同步改变。\emph[caution]{样本协方差的无偏性不能直接传递给样本相关系数、特征值或特征向量。}相关系数含随机标准差的比值，其期望一般不等于总体相关系数；讨论这些量的极大似然估计还需另给分布模型，不能仅凭 PCA 的代数分解得出。

对所有 $s_{ii}>0$ 的变量，令 $D_s=\operatorname{diag}(\sqrt{s_{11}},\ldots,\sqrt{s_{mm}})$，样本标准化矩阵与样本相关矩阵为
\begin{equation}
\mathbf Z=D_s^{-1}\mathbf X_c,\qquad
R_s=D_s^{-1}SD_s^{-1}
=\frac1{n-1}\mathbf Z\mathbf Z^{\mathsf T}.
\label{eq:lh-pca-sample-standardization}
\end{equation}
$\mathbf Z$ 的每一行均值为零、样本方差为 $1$，所以 $R_s$ 的对角元为 $1$、迹为 $m$。相关矩阵 PCA 对 $R_s$ 求特征向量，并用它们左乘 $\mathbf Z$ 得到得分；协方差矩阵 PCA 则对 $S$ 求特征向量，并左乘 $\mathbf X_c$。不能把其中一种矩阵的特征向量与另一种尺度的数据混合使用。

\subsection{相关矩阵的特征值分解算法}
\label{subsec:lh-pca-eigendecomposition}

教材的传统算法先标准化，再分解样本相关矩阵。为同时说明两种尺度，以下用 $\mathbf M$ 表示选定的中心化分析矩阵：相关矩阵 PCA 取 $\mathbf M=\mathbf Z$，协方差矩阵 PCA 取 $\mathbf M=\mathbf X_c$。记
\[
C_M=\frac1{n-1}\mathbf M\mathbf M^{\mathsf T}
=\widehat A\operatorname{diag}(\gamma_1,\ldots,\gamma_m)\widehat A^{\mathsf T},
\qquad\gamma_1\geq\cdots\geq\gamma_m\geq0.
\]
这里 $\widehat A=(\widehat a_1,\ldots,\widehat a_m)$ 为标准正交特征向量矩阵，$\gamma_i$ 为所选尺度下的样本主成分方差；取相关矩阵时 $\sum_i\gamma_i=m$。算法的输入为 $\mathbf X\in\mathbb R^{m\times n}$、尺度选择，以及目标维数 $q$ 或累计贡献率阈值 $\tau$；输出包括预处理参数、投影方向、主成分得分与方差贡献率，具体过程为：
\begin{enumerate}
\item 计算 $\bar x$ 和 $\mathbf X_c$；若选择相关矩阵 PCA，先移除零样本方差变量，再计算 $D_s$ 与 $\mathbf Z$，确定 $\mathbf M$。此后 $m$ 指保留的变量数。
\item 形成 $C_M$，使用实对称矩阵的特征值分解，按降序同步排列特征值与特征向量。理论上求解 $\det(C_M-\gamma I_m)=0$ 可刻画特征值，数值计算不必显式展开特征多项式。
\item 若总方差 $\operatorname{tr}(C_M)=0$，返回常量数据结果；否则计算 $\eta_i=\gamma_i/\operatorname{tr}(C_M)$。按给定 $q$ 截断，或取累计贡献率首次达到 $\tau$ 的维数，通常只保留正特征值对应的方向。
\item 令 $\widehat A_q=(\widehat a_1,\ldots,\widehat a_q)$，计算 $\mathbf Y_q=\widehat A_q^{\mathsf T}\mathbf M$。如需解释主成分，再按式~\eqref{eq:lh-pca-loading} 与式~\eqref{eq:lh-pca-variable-contribution} 的样本版本计算负荷量和各变量贡献率；相关矩阵情形的负荷量为 $\sqrt{\gamma_k}\,\widehat a_{ik}$。
\item 保存 $\bar x$、必要的 $D_s$、保留变量的索引与 $\widehat A_q$。新观测 $x_{\mathrm{new}}$ 必须沿用这些参数：协方差 PCA 的得分为 $\widehat A_q^{\mathsf T}(x_{\mathrm{new}}-\bar x)$，相关矩阵 PCA 的得分为 $\widehat A_q^{\mathsf T}D_s^{-1}(x_{\mathrm{new}}-\bar x)$，其中先取已保留变量的坐标。
\end{enumerate}
这是矩阵分解算法，取得分解并完成截断后即可输出，无须另设类似梯度下降的外层更新过程；若底层数值分解未达到容差，应报告数值求解未完成。稠密情况下，形成 $C_M$ 的计算量为 $O(m^2n)$，完整特征分解约为 $O(m^3)$，形成得分为 $O(qmn)$；除存储数据外，协方差矩阵需 $O(m^2)$ 存储。若用于后续监督学习，均值、标准差、方向与维数选择都应在训练数据上确定，再用于验证或测试数据，避免引入测试数据的信息。

\begin{example}
\label{ex:lh-pca-four-courses}
教材例题给出语文、外语、数学、物理四门课程标准化成绩的样本相关矩阵
\begin{equation}
R_s=\begin{pmatrix}
1&0.44&0.29&0.33\\
0.44&1&0.35&0.32\\
0.29&0.35&1&0.60\\
0.33&0.32&0.60&1
\end{pmatrix}.
\label{eq:lh-pca-course-correlation}
\end{equation}
要求累计方差贡献率至少为 $75\%$，求保留的主成分并解释其含义。
\end{example}
\begin{solve}
直接按式~\eqref{eq:lh-pca-course-correlation} 所列矩阵重新计算，四个特征值约为 $2.170165,\ 0.871005,\ 0.566179,\ 0.392650$，其和为 $4$。所以各方差贡献率约为 $0.542541,\ 0.217751,\ 0.141545,\ 0.098163$。第一主成分尚未达到阈值，前两个的累计贡献率约为 $0.760293$，故取 $q=2$。这里特征值表示主成分方差，贡献率还必须除以总方差 $4$。

记四门课程标准化后的取值为 $z_1,z_2,z_3,z_4$，选择第一方向系数均为正、第二方向语文系数为正的一组符号，得到得分公式
\begin{align*}
y_1&\approx0.4599z_1+0.4763z_2+0.5287z_3+0.5311z_4,\\
y_2&\approx0.5679z_1+0.4909z_2-0.4756z_3-0.4586z_4.
\end{align*}
第一主成分把四门课程同向组合，可解释为一种综合成绩方向；第二主成分把两门语言课程与数学、物理作反向组合，可解释为文理成绩的相对差异。整体改变某一方向的符号不会改变它表达的相对关系，也不会改变所保留的子空间。

按未舍入的特征值和特征向量计算，相关负荷量与各变量贡献率见表~\ref{tab:lh-pca-course-loadings}。前两个主成分保留约 $76\%$ 的总方差，但对各门课程保留的方差比例分别约为 $74\%$、$70\%$、$80\%$、$80\%$。负荷量说明变量与主成分的相关关系，不能仅凭其大小断言某门课程具有因果意义上的重要性。
\begin{table}
\centering
\caption{四门课程的前两个主成分负荷量与变量贡献率}
\label{tab:lh-pca-course-loadings}
\begin{tabular}{lrrrr}
\toprule
项目 & 语文 & 外语 & 数学 & 物理\\
\midrule
$\ell_{i1}$ & $0.6775$ & $0.7017$ & $0.7789$ & $0.7823$\\
$\ell_{i2}$ & $0.5300$ & $0.4582$ & $-0.4438$ & $-0.4280$\\
$v_i(2)$ & $0.7399$ & $0.7023$ & $0.8037$ & $0.7953$\\
\bottomrule
\end{tabular}
\end{table}
本例系数、负荷量与贡献率统一由题给相关矩阵计算，部分数值与教材表中的近似值不同。由于只给出了相关矩阵而没有逐人的原始成绩，能够确定投影规则及其方差性质，不能据此还原每个学生的主成分得分。
\end{solve}

\subsection{数据矩阵的奇异值分解算法}
\label{subsec:lh-pca-svd}

形成协方差矩阵并非计算主成分的必要步骤。设选定尺度下的中心化矩阵 $\mathbf M\in\mathbb R^{m\times n}$ 的秩为 $r>0$，其\emph{紧致奇异值分解}（compact singular value decomposition）为
\begin{equation}
\mathbf M=U_r\Delta_rV_r^{\mathsf T},\qquad
\Delta_r=\operatorname{diag}(d_1,\ldots,d_r),\qquad
d_1\geq\cdots\geq d_r>0,
\label{eq:lh-pca-data-svd}
\end{equation}
其中 $U_r\in\mathbb R^{m\times r}$、$V_r\in\mathbb R^{n\times r}$ 的列分别标准正交。于是
\begin{equation}
C_M=\frac1{n-1}\mathbf M\mathbf M^{\mathsf T}
=U_r\operatorname{diag}\left(\frac{d_1^2}{n-1},\ldots,
\frac{d_r^2}{n-1}\right)U_r^{\mathsf T}.
\label{eq:lh-pca-svd-covariance}
\end{equation}
因而\emph{以样本为列时，直接分解数据矩阵应取左奇异向量作为主成分方向}，主成分方差为 $\gamma_k=d_k^2/(n-1)$，其余特征值为零。对于 $1\leq q\leq r$，记 $U_q,V_q$ 为前 $q$ 列，$\Delta_q=\operatorname{diag}(d_1,\ldots,d_q)$，则
\begin{equation}
\widehat A_q=U_q,\qquad
\mathbf Y_q=U_q^{\mathsf T}\mathbf M=\Delta_qV_q^{\mathsf T},\qquad
\widehat{\mathbf M}_q=U_q\mathbf Y_q
=U_q\Delta_qV_q^{\mathsf T}.
\label{eq:lh-pca-svd-scores}
\end{equation}
方向矩阵为 $m\times q$，得分矩阵为 $q\times n$，重构矩阵仍为 $m\times n$，这三个维度也能帮助识别转置错误。

教材算法先将数据转置并缩放，定义
\begin{equation}
\mathbf T=\frac1{\sqrt{n-1}}\mathbf M^{\mathsf T}
\in\mathbb R^{n\times m},\qquad
\mathbf T^{\mathsf T}\mathbf T=C_M.
\label{eq:lh-pca-transposed-data}
\end{equation}
此时 $\mathbf T$ 的每一列对应一个变量，该列在 $n$ 个样本上的均值为零。若 $\mathbf T=\widetilde U_r\widetilde\Delta_r\widetilde V_r^{\mathsf T}$，则应取其\emph{右奇异向量} $\widetilde V_r\in\mathbb R^{m\times r}$ 作为方向。可选择 $\widetilde U_r=V_r$、$\widetilde V_r=U_r$、$\widetilde\Delta_r=\Delta_r/\sqrt{n-1}$；若单独求解两次分解，符号或重奇异值子空间内的基可能不同，但对应子空间相同。记 $\widetilde\Delta_r$ 的对角元为 $s_1\geq\cdots\geq s_r>0$，便有
\begin{equation}
\gamma_k=s_k^2,\qquad
\mathbf Y_q=\widetilde V_q^{\mathsf T}\mathbf M
=\sqrt{n-1}\,\widetilde\Delta_q\widetilde U_q^{\mathsf T}.
\label{eq:lh-pca-transposed-scores}
\end{equation}
所以教材取右奇异向量与直接分解 $\mathbf M$ 时取左奇异向量完全一致；关键是分解对象已经从 $m\times n$ 变成 $n\times m$。奇异向量给出投影方向，不能直接把方向矩阵本身称为样本得分矩阵。

采用教材的转置形式，输入预处理后的 $\mathbf M$ 及目标维数 $q$ 或累计贡献率阈值 $\tau$，可将算法写为：
\begin{enumerate}
\item 构造式~\eqref{eq:lh-pca-transposed-data} 中的 $\mathbf T$；若 $\|\mathbf T\|_F=0$，返回常量数据结果。这里 $\|A\|_F=(\sum_{i,j}a_{ij}^2)^{1/2}$ 为弗罗贝尼乌斯范数（Frobenius norm）。
\item 计算所需的奇异值和右奇异向量，按奇异值降序排列，确定数值秩；保留的正方差主成分个数取 $1\leq q\leq r$。若按阈值选维数，则取满足 $\sum_{k=1}^q s_k^2/\|\mathbf T\|_F^2\geq\tau$ 的最小 $q$，并求出足够多的奇异值以达到该阈值。
\item 返回方向矩阵 $\widetilde V_q$、得分 $\mathbf Y_q=\widetilde V_q^{\mathsf T}\mathbf M$、方差 $s_k^2$ 和贡献率 $s_k^2/\|\mathbf T\|_F^2$；需要重构时计算 $\widehat{\mathbf M}_q=\widetilde V_q\mathbf Y_q$。
\item 将重构恢复到原始尺度：协方差 PCA 取 $\widehat{\mathbf X}_q=\bar x\mathbf 1_n^{\mathsf T}+\widehat{\mathbf M}_q$；相关矩阵 PCA 取 $\widehat{\mathbf X}_q=\bar x\mathbf 1_n^{\mathsf T}+D_s\widehat{\mathbf M}_q$，并补回预先移除的常量变量。
\end{enumerate}
若只计算截断分解，贡献率的分母仍须为全部数据的平方范数，不能用已保留奇异值的平方和代替，否则累计贡献率会被错误地算成 $1$。数值秩需依据奇异值容差判断；将很小但非零的奇异值视为零属于数值近似，应与精确代数秩区分。

稠密完整 SVD 的计算量约为 $O(mn\min\{m,n\})$；只求少数分量时可以采用截断算法，但具体复杂度和误差依赖所用算法，不能统一宣称为 $O(qmn)$。直接 SVD 无须显式形成 $m\times m$ 协方差矩阵，特别适合变量多于样本的情形；在有限精度下还可避免形成 Gram 矩阵所带来的条件数平方效应。具体说，在非零奇异值对应的子空间上，$\mathbf M\mathbf M^{\mathsf T}$ 的最大、最小正特征值之比为 $(d_1/d_r)^2$，因此很小的分量在显式形成协方差时更容易受到舍入误差影响。

\begin{theorem}[样本 PCA 与最佳低秩近似]
\label{thm:lh-pca-low-rank}
设 $\mathbf M$ 的紧致 SVD 如式~\eqref{eq:lh-pca-data-svd}，$1\leq q\leq r$。在所有秩不超过 $q$ 的矩阵中，截断重构 $\widehat{\mathbf M}_q$ 是弗罗贝尼乌斯范数下的一个最佳近似，且
\begin{equation}
\min_{\substack{L\in\mathbb R^{m\times n}\\\operatorname{rank}(L)\leq q}}
\|\mathbf M-L\|_F^2
=\|\mathbf M-\widehat{\mathbf M}_q\|_F^2
=\sum_{j=q+1}^r d_j^2
=(n-1)\sum_{j=q+1}^m\gamma_j.
\label{eq:lh-pca-low-rank}
\end{equation}
其中令 $\gamma_j=0$ 对 $j>r$。去掉目标中的平方不改变最优解。
\end{theorem}
\begin{proof}
对任意可行的 $L$，令 $P_L$ 为其列空间上的正交投影，则 $P_LL=L$。按每列的正交分解，
\[
\|\mathbf M-L\|_F^2
=\|(I_m-P_L)\mathbf M\|_F^2
+\|P_L\mathbf M-L\|_F^2
\geq\|\mathbf M\|_F^2-\operatorname{tr}(P_L\mathbf M\mathbf M^{\mathsf T}).
\]
投影 $P_L$ 的秩不超过 $q$，故由定理~\ref{thm:lh-pca-variance-optimal} 及特征值非负性，$\operatorname{tr}(P_L\mathbf M\mathbf M^{\mathsf T})\leq\sum_{j=1}^q d_j^2$。又 $\|\mathbf M\|_F^2=\sum_{j=1}^r d_j^2$，得到误差下界 $\sum_{j>q}d_j^2$。取 $L=U_qU_q^{\mathsf T}\mathbf M=\widehat{\mathbf M}_q$，由于各奇异向量标准正交，残余项 $\sum_{j>q}d_ju_jv_j^{\mathsf T}$ 的平方范数恰好等于该下界，因而达到最优值。最后用 $d_j^2=(n-1)\gamma_j$ 即得末式。
\end{proof}

这给出教材关于 PCA 等价于低秩逼近的习题所需依据：对中心化或进一步标准化的分析矩阵，先投影到前 $q$ 个主方向，再由得分重构，正是最佳秩不超过 $q$ 的近似。约束应写为 $\operatorname{rank}(L)\leq q$；若误写为严格小于 $q$，最多只能保留 $q-1$ 个非零方向，问题便发生改变。未经中心化就直接截断原始数据矩阵，求出的仍是该原始矩阵的低秩近似，但方向反映的是关于原点的二阶矩，通常不等于围绕样本均值的协方差 PCA。

保留全部正奇异值时，当前样本可被精确重构；保留少数奇异值时，$\sum_{j>q}d_j^2$ 量化所损失的平方误差。但只有在额外的信号与噪声假设下，才能把小方差方向解释为噪声，普通 PCA 本身不提供这一判别。它也对异常值和尺度敏感，且只寻找线性子空间；是否适合作为后续任务的预处理，应结合变量含义、验证表现与重构需求判断，而不能只看累计方差贡献率。

% 教材来源：机器学习/textbooks/机器学习（分章原文）/机器学习-17.pdf。
% 对应教材第 17 章，原书第 321--338 页；章节、公式和定理均自动编号。

\fi

\selectchapter{17}
\ifcompilechapter
\chapter{潜在语义分析}
\label{ch:lh-lsa}

文本检索首先需要回答两个问题：怎样把一篇文本表示成可以计算的对象，以及怎样衡量两篇文本讨论的内容是否接近。直接按单词计数得到的表示简单而稀疏，但同义词可能被视为彼此无关的坐标，多义词又可能把不同内容混在同一坐标中。\emph{潜在语义分析}（latent semantic analysis, LSA）从单词在文本集合中的共同出现模式出发，通过低秩矩阵分解构造较低维的表示，尝试在单词与文本之间引入潜在话题这一层关系；用于信息检索时也称\emph{潜在语义索引}（latent semantic indexing, LSI）。

本章沿用教材“单词向量空间与话题向量空间、潜在语义分析算法、非负矩阵分解算法”的路线，先说明矩阵两侧因子的含义，再讨论截断奇异值分解与非负矩阵分解。前者以正交方向描述主要变化，后者以非负成分的叠加描述文本，二者都是本章所讨论的非概率话题分析方法。阅读需要矩阵乘法、正交投影、特征值和基本微分；所用奇异值分解及优化步骤在本章给出必要说明。全文的 $\log$ 为自然对数，向量范数 $\|\cdot\|_2$ 为欧氏范数，矩阵范数 $\|\cdot\|_F$ 为弗罗贝尼乌斯范数，矩阵不等式 $A\geq0$ 表示所有元素非负。

\section{单词向量空间与话题向量空间}
\label{sec:lh-lsa-spaces}

\subsection{单词向量空间}
\label{subsec:lh-lsa-word-space}

\begin{definition}[单词向量空间模型]
\label{def:lh-lsa-vsm}
给定文本集合 $\mathcal D=\{d_1,\ldots,d_n\}$ 和词表 $\mathcal V=\{w_1,\ldots,w_m\}$，以
\begin{equation}
X=(x_{ij})_{m\times n}=[x_1,\ldots,x_n]\in\mathbb R_{\geq0}^{m\times n}
\label{eq:lh-lsa-word-document}
\end{equation}
表示\emph{单词-文本矩阵}（word-document matrix）。其中 $x_{ij}$ 是单词 $w_i$ 在文本 $d_j$ 中的频数或非负权值，列向量 $x_j\in\mathbb R^m$ 是该文本的表示。以这些向量及其内积、夹角等度量刻画文本关系，称为\emph{向量空间模型}（vector space model, VSM）。
\end{definition}

这里每一维对应一个单词，每一列对应一篇文本，不能把列向量误认为某个单词的词向量。该表示忽略词序，属于词袋（bag-of-words）表示；它把单词出现的模式作为文本内容的可计算近似，并没有给出自然语言语义的严格定义。词表通常很大，而每篇文本只使用其中一小部分单词，因此 $X$ 往往是稀疏矩阵。

除直接计数外，教材使用\emph{单词频率-逆文本频率}（term frequency-inverse document frequency, TF-IDF）加权。记 $c_{ij}$ 为单词 $w_i$ 在文本 $d_j$ 中出现的次数，$L_j=\sum_{i=1}^m c_{ij}$ 为文本的词数，$n_i=\#\{j:c_{ij}>0\}$ 为含该单词的文本数。对 $L_j>0$、$1\leq n_i\leq n$，教材采用的权值为
\begin{equation}
x_{ij}=\frac{c_{ij}}{L_j}\log\frac{n}{n_i}.
\label{eq:lh-lsa-tfidf}
\end{equation}
第一项衡量词在本篇文本中的相对频率，第二项降低普遍出现的词的权重；当 $n_i=n$ 时，该词的逆文本频率为 $0$。这是 TF-IDF 的一种具体约定，带平滑项、对数词频等其他版本会产生不同的矩阵，使用时应先固定定义。空文本不能代入式~\eqref{eq:lh-lsa-tfidf}，加权后全零的文本向量也需要单独处理。

对非零向量 $x_p,x_q$，可以用内积 $x_p^{\mathsf T}x_q$ 或\emph{余弦相似度}（cosine similarity）
\begin{equation}
\operatorname{sim}_{\mathrm{word}}(d_p,d_q)
=\frac{x_p^{\mathsf T}x_q}{\|x_p\|_2\|x_q\|_2}
\label{eq:lh-lsa-word-cosine}
\end{equation}
比较文本。内积同时受方向和向量长度影响，余弦消除了整体正比例缩放的影响；对本节的非负向量，余弦在 $[0,1]$ 内。查询也可以使用同一词表和语料统计量表示成向量，再按与各文本的相似度排序。稀疏内积只需在共同非零坐标上求和，但相似度的语义解释仍依赖表示是否合适。

\begin{example}
\label{ex:lh-lsa-synonym-polysemy}
教材的单词-文本矩阵例子包含四篇文本。依次按 airplane、aircraft、computer、apple、fruit、produce 排列行，频数矩阵为
\begin{equation}
X=\begin{pmatrix}
2&0&0&0\\
0&2&0&0\\
0&0&1&0\\
0&0&2&3\\
0&0&0&1\\
1&2&2&1
\end{pmatrix}.
\label{eq:lh-lsa-four-documents}
\end{equation}
比较前两篇文本以及后两篇文本的词空间余弦，说明表示的局限。
\end{example}
\begin{solve}
前两列内积为 $2$，范数分别为 $\sqrt5$、$\sqrt8$；后两列内积为 $8$，范数分别为 $3$、$\sqrt{11}$，故
\[
\operatorname{sim}_{\mathrm{word}}(d_1,d_2)=\frac1{\sqrt{10}}\approx0.3162,
\qquad
\operatorname{sim}_{\mathrm{word}}(d_3,d_4)=\frac8{3\sqrt{11}}\approx0.8040.
\]
airplane 与 aircraft 表达相近内容，却分属两个坐标，因而前两篇文本的联系没有被充分表达。apple 在后两篇文本中可分别指计算机品牌和水果，共享坐标却不能区分词义。这分别体现了\emph{多词一义性}（synonymy）与\emph{一词多义性}（polysemy）带来的问题；produce 等共有词的权重也会影响比较结果。
\end{solve}

\subsection{话题向量空间}
\label{subsec:lh-lsa-topic-space}

话题（topic）在这里指由一组相关单词共同反映的内容方向。希望不同的同义词能参与相同话题，而同一个多义词能参与不同话题，再用一篇文本在各话题上的权值描述其内容。这是建模动机，并不意味着矩阵分解一定能恢复人的语义分类，也不意味着每个分解方向都能获得清晰的话题名称。

设选取 $k$ 个向量 $t_\ell=(t_{1\ell},\ldots,t_{m\ell})^{\mathsf T}\in\mathbb R^m$，构成\emph{单词-话题矩阵}（word-topic matrix）$T=[t_1,\ldots,t_k]\in\mathbb R^{m\times k}$。它们张成的子空间为 $\mathcal S=\operatorname{span}\{t_1,\ldots,t_k\}$，维数是 $\operatorname{rank}(T)\leq k$；只有列线性无关时维数才等于 $k$。若用 $y_j\in\mathbb R^k$ 表示文本 $d_j$ 的话题系数，便希望
\begin{equation}
x_j\approx Ty_j=\sum_{\ell=1}^k y_{\ell j}t_\ell,
\qquad
X\approx TY,\qquad
Y=[y_1,\ldots,y_n]\in\mathbb R^{k\times n}.
\label{eq:lh-lsa-topic-factorization}
\end{equation}
$Y$ 称为\emph{话题-文本矩阵}（topic-document matrix）。$T$ 的列描述话题在单词坐标中的形式，$Y$ 的列描述文本在这些话题上的系数，$Ty_j$ 则是返回原词空间后的重构向量。

\begin{proposition}[坐标、投影与内积]
\label{prop:lh-lsa-projection}
若 $T$ 列满秩，则对任意 $x\in\mathbb R^m$，最小二乘问题 $\min_y\|x-Ty\|_2^2$ 有唯一解
\begin{equation}
y^*=(T^{\mathsf T}T)^{-1}T^{\mathsf T}x,
\qquad
\widehat x=Ty^*=T(T^{\mathsf T}T)^{-1}T^{\mathsf T}x.
\label{eq:lh-lsa-general-projection}
\end{equation}
其中 $\widehat x\in\mathcal S$ 是正交投影，$y^*\in\mathbb R^k$ 是其坐标。若 $T^{\mathsf T}T=I_k$，则 $y^*=T^{\mathsf T}x$，且对任意 $a,b\in\mathbb R^k$ 有 $(Ta)^{\mathsf T}(Tb)=a^{\mathsf T}b$。
\end{proposition}
\begin{proof}
目标对 $y$ 的梯度为 $2T^{\mathsf T}(Ty-x)$。列满秩使 $T^{\mathsf T}T$ 正定，解正规方程便得到唯一最小点，且 $T^{\mathsf T}(x-\widehat x)=0$，即残差与 $\mathcal S$ 正交。最后，$(Ta)^{\mathsf T}(Tb)=a^{\mathsf T}(T^{\mathsf T}T)b$，代入正交条件即得结论。
\end{proof}

因此，把词空间向量变成坐标的映射一般不是 $T$ 本身，$T$ 的作用是把系数重构为词空间向量。若话题列不正交，重构向量的内积为 $y_p^{\mathsf T}(T^{\mathsf T}T)y_q$，未必等于 $y_p^{\mathsf T}y_q$；直接采用后一个内积，是在系数空间选用了另一种度量。潜在语义分析还需要从数据中学习 $T$，而不是事先知道话题后只求 $Y$，这正是下两节引入矩阵分解的原因。

\section{潜在语义分析算法}
\label{sec:lh-lsa-svd}

\subsection{矩阵奇异值分解算法}
\label{subsec:lh-lsa-truncated-svd}

设 $X\in\mathbb R^{m\times n}$ 的秩为 $r>0$，其紧奇异值分解（singular value decomposition, SVD）为
\begin{equation}
X=U_r\Sigma_rV_r^{\mathsf T}
=\sum_{\ell=1}^r\sigma_\ell u_\ell v_\ell^{\mathsf T},
\qquad \sigma_1\geq\cdots\geq\sigma_r>0.
\label{eq:lh-lsa-svd}
\end{equation}
这里 $U_r\in\mathbb R^{m\times r}$、$V_r\in\mathbb R^{n\times r}$ 的列分别是正交归一的左、右奇异向量，$\Sigma_r=\operatorname{diag}(\sigma_1,\ldots,\sigma_r)$。左奇异向量是 $XX^{\mathsf T}$ 的特征向量，右奇异向量是 $X^{\mathsf T}X$ 的特征向量，对应特征值均为 $\sigma_\ell^2$。不要求 $m\geq n$；只需 $r\leq\min(m,n)$。

给定 $1\leq k\leq r$，保留最大的 $k$ 个奇异值及其对应向量，得到\emph{截断奇异值分解}（truncated SVD）
\begin{equation}
X_k=U_k\Sigma_kV_k^{\mathsf T},\qquad
T=U_k,\qquad Y=\Sigma_kV_k^{\mathsf T}=U_k^{\mathsf T}X.
\label{eq:lh-lsa-svd-factors}
\end{equation}
各因子的维度依次为 $m\times k$、$k\times k$、$k\times n$。$U_k$ 的列张成所学话题子空间，文本 $d_j$ 的话题表示及其重构分别为
\begin{equation}
y_j=U_k^{\mathsf T}x_j
=(\sigma_1v_{j1},\ldots,\sigma_kv_{jk})^{\mathsf T},
\qquad
\widehat x_j=U_ky_j=\sum_{\ell=1}^k\sigma_\ell v_{j\ell}u_\ell.
\label{eq:lh-lsa-document-coordinates}
\end{equation}
其中 $v_{j\ell}$ 是右奇异向量 $v_\ell$ 的第 $j$ 个分量。坐标等式对原始 $X$ 精确成立，而 $x_j\approx\widehat x_j$ 一般只是近似；只有保留了全部非零奇异方向时，才能把整个重构等式写成 $X=X_k$。

由\crosschapterref{15}{定理~\ref{thm:lh-svd-best-approximation}}{截断 SVD 的最优性定理}，截断 SVD 在给定秩约束下使平方重构误差最小，具体为
\begin{equation}
\min_{\operatorname{rank}(B)\leq k}\|X-B\|_F^2
=\|X-X_k\|_F^2=\sum_{\ell=k+1}^r\sigma_\ell^2.
\label{eq:lh-lsa-best-rank-k}
\end{equation}
当 $k=r$ 时，右端的空和为 $0$。该结论的证明已在前文给出，本章将其用于单词-文本矩阵；残差的各奇异方向相互正交，使误差可由丢弃的奇异值平方和直接计算。

这个最优性针对给定矩阵的平方重构误差，不能直接推出检索或话题解释的效果最优。可以考察保留比例 $\sum_{\ell=1}^k\sigma_\ell^2/\sum_{\ell=1}^r\sigma_\ell^2$ 来描述重构信息量，但 $k$ 仍应结合下游任务选择。$k$ 太小会合并不同内容，太大则保留更多原始词差异及噪声；这种取舍不能由“维数越低越好”概括。

\emph[tip]{SVD 话题分析的计算路线}是：先固定分词、词表和加权方法，构造 $X$；再给定 $k$，计算前 $k$ 个奇异三元组，输出 $U_k$ 与 $Y=\Sigma_kV_k^{\mathsf T}$；最后对非零的 $y_j$ 使用内积或余弦比较文本。新查询 $q\in\mathbb R^m$ 使用相同预处理后，以 $y_q=U_k^{\mathsf T}q$ 投影到同一空间，不能为查询单独重新学习一套坐标。即使 $q\ne0$，也可能有 $y_q=0$，此时话题空间余弦仍未定义。

由于 $U_k$ 的列正交归一，$y_p^{\mathsf T}y_q=\widehat x_p^{\mathsf T}\widehat x_q$、$\|y_j\|_2=\|\widehat x_j\|_2$。又有 $Y^{\mathsf T}Y=V_k\Sigma_k^2V_k^{\mathsf T}$，说明本文的内积保留了奇异值的尺度。若改用 $V_k^{\mathsf T}$ 的列作为表示，就相当于把坐标再乘 $\Sigma_k^{-1}$，会改变相似度，不能与式~\eqref{eq:lh-lsa-document-coordinates} 混用。奇异向量可能含负分量，其符号本身也不唯一；同时把 $u_\ell,v_\ell$ 变号不会改变重构与文本相似度。

\subsection{例子}
\label{subsec:lh-lsa-examples}

\begin{example}
\label{ex:lh-lsa-four-documents-svd}
对例~\ref{ex:lh-lsa-synonym-polysemy} 的矩阵取 $k=2$ 做潜在语义分析，观察词空间相似度的变化。这也给出教材习题 17.1 的一种具体计算。
\end{example}
\begin{solve}
对式~\eqref{eq:lh-lsa-four-documents} 的矩阵数值求解奇异值分解，得到
\[
(\sigma_1,\sigma_2,\sigma_3,\sigma_4)
\approx(4.4770,\ 2.7520,\ 2.0000,\ 1.1762).
\]
选取一组符号约定后，前两个话题上的文本坐标为
\begin{equation}
Y=\Sigma_2V_2^{\mathsf T}\approx
\begin{pmatrix}
0.7861&1.5721&2.8586&2.9635\\
1.0770&2.1540&-0.1045&-1.3276
\end{pmatrix}.
\label{eq:lh-lsa-four-documents-coordinates}
\end{equation}
用未舍入的坐标计算，$\operatorname{sim}_{\mathrm{topic}}(d_1,d_2)=1$，$\operatorname{sim}_{\mathrm{topic}}(d_3,d_4)\approx0.9269$；平方重构误差为 $\sigma_3^2+\sigma_4^2\approx5.3835$。前两篇文本在此二维表示中同向，原来的同义词差异被压缩；但后两篇文本的相似度也提高了，多义词的混淆没有自动消失。这说明低秩表示能够汇总共现信息，却不能保证把每一种语义差异都正确区分，尤其在只有少量文本时更应核查结果。
\end{solve}

\begin{example}
\label{ex:lh-lsa-textbook-nine-documents}
教材还给出含 $11$ 个单词、$9$ 篇文本的频数矩阵，并取 $k=3$。说明分解后各矩阵及其列的含义，读取两篇文本的坐标。
\end{example}
\begin{solve}
分解为 $X\approx U_3(\Sigma_3V_3^{\mathsf T})$，其中 $U_3$ 为 $11\times3$，$\Sigma_3V_3^{\mathsf T}$ 为 $3\times9$。前者的三列是三个潜在方向，后者的九列是九篇文本的坐标。教材将后一矩阵显示为两位小数：
\begin{equation}
Y\approx\begin{pmatrix}
1.37&0.86&1.33&1.02&0.86&1.92&1.09&1.13&1.72\\
-0.84&-0.39&-1.20&-0.63&-0.37&1.44&0.18&-0.81&1.15\\
-0.82&0.28&-0.32&0.50&0.44&-1.02&1.10&0.00&0.68
\end{pmatrix}.
\label{eq:lh-lsa-nine-documents-coordinates}
\end{equation}
例如 $y_1\approx(1.37,-0.84,-0.82)^{\mathsf T}$、$y_2\approx(0.86,-0.39,0.28)^{\mathsf T}$，由这些已舍入数据计算的余弦约为 $0.7182$。这里每个坐标是带符号的线性系数，负数不表示“负概率”；也不能把第一行误读成第一篇文本。显示的两位小数仅用于阅读，精确重构与相似度计算应使用分解所得的完整数值。
\end{solve}

\section{非负矩阵分解算法}
\label{sec:lh-lsa-nmf}

\subsection{非负矩阵分解}
\label{subsec:lh-lsa-nmf-definition}

SVD 允许正负系数相互抵消，而词频本身是非负的。若希望每篇文本由若干非负话题成分相加构成，可以同时约束两个因子非负。

\begin{definition}[非负矩阵分解]
\label{def:lh-lsa-nmf}
给定非负矩阵 $X\in\mathbb R_{\geq0}^{m\times n}$ 和整数 $k\geq1$，寻找
\begin{equation}
W=[w_1,\ldots,w_k]\in\mathbb R_{\geq0}^{m\times k},\qquad
H=[h_1,\ldots,h_n]\in\mathbb R_{\geq0}^{k\times n},\qquad
X\approx WH,
\label{eq:lh-lsa-nmf-factorization}
\end{equation}
称为\emph{非负矩阵分解}（non-negative matrix factorization, NMF）。$W$ 称为基矩阵，$H$ 称为系数矩阵，近似程度由选定的损失函数衡量。
\end{definition}

对第 $j$ 篇文本有 $x_j\approx Wh_j=\sum_{\ell=1}^k h_{\ell j}w_\ell$，各项只有非负叠加，没有正负抵消。这里“基矩阵”是惯用名称，并不保证列向量线性无关。常取 $k<\min(m,n)$ 以获得低秩表示，但若按存储的标量个数衡量压缩，还必须有 $k(m+n)<mn$；仅有两个因子各自的某一维变小，不足以保证总参数更少。

\subsection{潜在语义分析模型}
\label{subsec:lh-lsa-nmf-topic-model}

在文本分析中，$W$ 的第 $\ell$ 列表示第 $\ell$ 个话题的单词权重，$H$ 的第 $j$ 列表示第 $j$ 篇文本的话题权重。由于两者都非负，可以用每列中权重较大的词解释话题，再看各文本由哪些话题共同构成。可重构向量属于非负锥
\begin{equation}
\mathcal C_W=\{Wh:h\in\mathbb R_{\geq0}^k\},
\label{eq:lh-lsa-nmf-cone}
\end{equation}
它一般不是线性子空间。固定 $W$ 后，求非负系数是带约束的拟合问题，通常既不是 $h=W^{\mathsf T}x$，也不是词空间上的固定线性变换。

非负性便于“部分相加构成整体”的解释，但不自动赋予概率意义。若 $c_\ell=\sum_iW_{i\ell}>0$，可令 $C=\operatorname{diag}(c_1,\ldots,c_k)$，采用
\begin{equation}
\widetilde W=WC^{-1},\qquad
\widetilde H=CH,\qquad
\widetilde W\widetilde H=WH.
\label{eq:lh-lsa-nmf-rescaling}
\end{equation}
此时 $\widetilde W$ 的每列和为 $1$，可以当作归一化的词权重阅读；$\widetilde H$ 的列和仍未必为 $1$，一般保留着文本的总量尺度。没有进一步的随机生成假设，整个分解仍是非概率模型。

式~\eqref{eq:lh-lsa-nmf-rescaling} 也说明分解至少存在正对角缩放的不唯一性，话题顺序还可同时置换。比较不同文本的系数余弦前，应固定话题列的尺度；否则某一行 $H$ 的缩放会改变夹角，即使 $WH$ 完全相同。$W$ 一般不正交，因此 $h_p^{\mathsf T}h_q$ 与重构文本的内积 $h_p^{\mathsf T}W^{\mathsf T}Wh_q$ 也须区分。

\subsection{非负矩阵分解的形式化}
\label{subsec:lh-lsa-nmf-objectives}

教材考虑两种损失。第一种是平方损失，采用不改变最小点的 $1/2$ 因子，将问题写成
\begin{equation}
\min_{\substack{W\in\mathbb R_{\geq0}^{m\times k}\\H\in\mathbb R_{\geq0}^{k\times n}}}
J_E(W,H)=\frac12\|X-WH\|_F^2
=\frac12\sum_{i=1}^m\sum_{j=1}^n\bigl(x_{ij}-(WH)_{ij}\bigr)^2.
\label{eq:lh-lsa-nmf-euclidean}
\end{equation}
第二种是\emph{广义 Kullback--Leibler 散度}（generalized Kullback--Leibler divergence）
\begin{equation}
D(A\|B)=\sum_{i,j}\left(a_{ij}\log\frac{a_{ij}}{b_{ij}}-a_{ij}+b_{ij}\right),
\label{eq:lh-lsa-generalized-kl}
\end{equation}
其中 $A,B$ 同型且非负。按连续延拓约定，$a=0$ 时单项为 $b$，$a>0,b=0$ 时为 $+\infty$，$a=b=0$ 时为 $0$。相应优化问题为
\begin{equation}
\min_{\substack{W\in\mathbb R_{\geq0}^{m\times k}\\H\in\mathbb R_{\geq0}^{k\times n}}}
J_D(W,H)=D(X\|WH).
\label{eq:lh-lsa-nmf-divergence}
\end{equation}
这一定义允许原始频数矩阵不归一化，不能省去 $-a_{ij}+b_{ij}$ 两项。

\begin{proposition}[散度的基本性质]
\label{prop:lh-lsa-kl-properties}
对非负同型矩阵，$D(A\|B)\geq0$，且等号当且仅当 $A=B$。一般有 $D(A\|B)\ne D(B\|A)$；若 $\sum_{i,j}a_{ij}=\sum_{i,j}b_{ij}=1$，则式~\eqref{eq:lh-lsa-generalized-kl} 化为两个离散概率分布之间的 KL 散度。
\end{proposition}
\begin{proof}
对 $a,b>0$，令 $t=a/b$，单项为 $b(t\log t-t+1)$。函数 $t\log t-t+1$ 的导数为 $\log t$，在 $t=1$ 取得唯一最小值 $0$；边界情形按上述约定仍非负，且仅在 $a=b$ 时为 $0$。逐项求和得到非负性和等号条件。交换 $a,b$ 一般改变数值，故不对称；两矩阵总和均为 $1$ 时，线性项求和相消，剩下通常的 KL 表达式。
\end{proof}

两种损失对 $W,H$ 之一固定后的另一变量都是凸的，但联合优化一般非凸。例如标量 $X=1$ 时，$(w,h)=(1,1)$ 与 $(2,1/2)$ 的两种损失都为 $0$，它们的中点 $(3/2,3/4)$ 却给出 $wh=9/8$ 和正损失，违反联合凸性。因此不能把每个子问题的凸性当作整个分解的全局最优保证，数值迭代的结果可能随初值而不同。

\subsection{算法}
\label{subsec:lh-lsa-nmf-algorithm}

先对平方损失求梯度。因 $(WH)_{ij}=\sum_\ell W_{i\ell}H_{\ell j}$，逐项求导得
\begin{align}
\nabla_WJ_E&=(WH-X)H^{\mathsf T}=WHH^{\mathsf T}-XH^{\mathsf T},
\label{eq:lh-lsa-nmf-gradient-w}\\
\nabla_HJ_E&=W^{\mathsf T}(WH-X)=W^{\mathsf T}WH-W^{\mathsf T}X.
\label{eq:lh-lsa-nmf-gradient-h}
\end{align}
例如 $\partial J_E/\partial W_{i\ell}=\sum_j((WH)_{ij}-x_{ij})H_{\ell j}$。普通梯度下降可能越出非负约束；若分别取元素步长 $W_{i\ell}/(WHH^{\mathsf T})_{i\ell}$ 和 $H_{\ell j}/(W^{\mathsf T}WH)_{\ell j}$，则加法更新化成教材的\emph{乘法更新规则}（multiplicative update rules）。

\begin{theorem}[平方损失的乘法更新]
\label{thm:lh-lsa-nmf-euclidean-update}
给定 $X\geq0$，对严格正的当前因子，在各分母为正时，固定另一因子并作下列任一整块更新，均不增大 $J_E$：
\begin{align}
W_{i\ell}&\leftarrow W_{i\ell}
\frac{(XH^{\mathsf T})_{i\ell}}{(WHH^{\mathsf T})_{i\ell}},
\label{eq:lh-lsa-nmf-update-w}\\
H_{\ell j}&\leftarrow H_{\ell j}
\frac{(W^{\mathsf T}X)_{\ell j}}{(W^{\mathsf T}WH)_{\ell j}}.
\label{eq:lh-lsa-nmf-update-h}
\end{align}
每一块右端都使用该块更新前的值；完成一块后，再以更新后的因子计算下一块。
\end{theorem}
\begin{proof}
固定 $W$，各列的目标可分别写成 $f(h)=\frac12\|x-Wh\|_2^2$。令 $A=W^{\mathsf T}W$、$b=W^{\mathsf T}x$，则 $\nabla f(h)=Ah-b$。对当前 $h>0$，定义对角矩阵 $K_{\ell\ell}=(Ah)_\ell/h_\ell$。因 $A$ 对称且逐项非负，对任意 $z$ 有
\[
z^{\mathsf T}(K-A)z
=\frac12\sum_{\ell,s}A_{\ell s}h_\ell h_s
\left(\frac{z_\ell}{h_\ell}-\frac{z_s}{h_s}\right)^2\geq0.
\]
故辅助函数
\[
G(h'\mid h)=f(h)+(h'-h)^{\mathsf T}\nabla f(h)
+\frac12(h'-h)^{\mathsf T}K(h'-h)
\]
满足 $G(h'\mid h)\geq f(h')$，且 $G(h\mid h)=f(h)$。其最小点为
\[
h'_\ell=h_\ell-\frac{h_\ell}{(Ah)_\ell}\bigl((Ah)_\ell-b_\ell\bigr)
=h_\ell\frac{b_\ell}{(Ah)_\ell}\geq0.
\]
于是 $f(h')\leq G(h'\mid h)\leq G(h\mid h)=f(h)$，逐列相加得到式~\eqref{eq:lh-lsa-nmf-update-h} 的结论。对转置问题使用同一论证，即得 $W$ 更新的结论。这里的辅助函数论证也说明，单凭“使用某种梯度步长”本身还不足以证明下降。
\end{proof}

对散度损失，令 $Z=WH$，在 $Z_{ij}>0$ 的位置定义 $R_{ij}=x_{ij}/Z_{ij}$，并令 $E=\mathbf1_m\mathbf1_n^{\mathsf T}$ 为全 $1$ 矩阵。由 $\partial(x\log(x/z)-x+z)/\partial z=1-x/z$ 得
\begin{equation}
\nabla_WJ_D=(E-R)H^{\mathsf T},\qquad
\nabla_HJ_D=W^{\mathsf T}(E-R).
\label{eq:lh-lsa-nmf-kl-gradient}
\end{equation}
对应的更新不再使用平方损失中的分母，而是下述规则。

\begin{theorem}[散度损失的乘法更新]
\label{thm:lh-lsa-nmf-kl-update}
给定 $X\geq0$，在当前 $W,H$ 严格正、各分母为正时，交替作以下整块更新，均不增大 $J_D$：
\begin{align}
W_{i\ell}&\leftarrow W_{i\ell}
\frac{\sum_{j=1}^n H_{\ell j}x_{ij}/(WH)_{ij}}
{\sum_{j=1}^n H_{\ell j}},
\label{eq:lh-lsa-nmf-kl-update-w}\\
H_{\ell j}&\leftarrow H_{\ell j}
\frac{\sum_{i=1}^m W_{i\ell}x_{ij}/(WH)_{ij}}
{\sum_{i=1}^m W_{i\ell}}.
\label{eq:lh-lsa-nmf-kl-update-h}
\end{align}
同样，每块内的全部比值使用该块开始时的因子计算，第二块使用第一块已经更新的结果。
\end{theorem}
\begin{proof}
仍固定 $W$ 并只看一列 $x,h$。记 $z_i=(Wh)_i$、$\alpha_{i\ell}=W_{i\ell}h_\ell/z_i$，则 $\alpha_{i\ell}>0$、$\sum_\ell\alpha_{i\ell}=1$。由 $-\log$ 的凸性，对 $h'>0$ 有
\[
-\log\!\left(\sum_\ell W_{i\ell}h'_\ell\right)
\leq-\sum_\ell\alpha_{i\ell}
\log\frac{W_{i\ell}h'_\ell}{\alpha_{i\ell}},
\]
且在 $h'=h$ 时等号成立。将其乘以 $x_i\geq0$ 并加上线性项，便得到散度目标的一个相切上界；去掉与 $h'$ 无关的常数后，上界为
\[
\sum_\ell\left[
\left(\sum_iW_{i\ell}\right)h'_\ell
-\left(\sum_ix_i\alpha_{i\ell}\right)\log h'_\ell\right].
\]
逐坐标最小化可得 $h'_\ell=(\sum_ix_i\alpha_{i\ell})/(\sum_iW_{i\ell})$，代入 $\alpha_{i\ell}$ 即为式~\eqref{eq:lh-lsa-nmf-kl-update-h}。当分子为零时，最小点按连续延拓取 $h'_\ell=0$。由“目标不超过上界，上界经最小化不增”得到下降结论，转置后得到 $W$ 的更新。
\end{proof}

两种损失可使用同一迭代流程。输入非负矩阵 $X$、话题数 $k$、最大迭代次数 $t_{\max}$、容差 $\varepsilon>0$ 和选定的损失，输出数值近似因子 $W,H$ 及最终损失。若 $X=0$，直接取零乘积即可；否则先移除全零行、全零列并记录位置，最后补回对应的零权重，以免无信息的行列产生 $0/0$。对剩余的矩阵执行以下步骤。
\begin{enumerate}
\item 初始化严格正的 $W,H$，固定一种话题列的尺度约定，并记录初始损失 $J^{(0)}$。只要求非负的初始化可能把某些元素永久锁在零值，因而通常从正数开始。
\item 在第 $t$ 轮，平方损失先按式~\eqref{eq:lh-lsa-nmf-update-w} 更新整个 $W$，再按式~\eqref{eq:lh-lsa-nmf-update-h} 更新整个 $H$；散度损失分别改用式~\eqref{eq:lh-lsa-nmf-kl-update-w} 和式~\eqref{eq:lh-lsa-nmf-kl-update-h}。每一块更新前统一计算其右端，避免在同一块中混用新旧元素。
\item 如需对 $W$ 的列归一化，按式~\eqref{eq:lh-lsa-nmf-rescaling} 同时反向缩放 $H$ 的对应行，以保持乘积及损失不变；不能只缩放 $W$。此处也可选正的列欧氏范数作为 $c_\ell$。
\item 计算 $J^{(t)}$。若 $|J^{(t)}-J^{(t-1)}|/\max\{1,J^{(t-1)}\}\leq\varepsilon$，可按损失变化停止；否则继续，直到 $t_{\max}$。返回因子时说明所用停止条件，达到上限并不表示已满足最优性条件。
\end{enumerate}
移除全零行列后，严格正初值使上述精确更新保持正性，因而每轮分母有定义；补回全零行列不会改变有效部分的拟合。浮点计算中若出现下溢、零分母或死话题，应检查并重置相应成分。直接在分母加小常数会改变更新映射，此时不能不加说明地沿用原规则的精确单调性结论。

\begin{notealert}[type=caution,title={损失收敛不等于得到全局最优分解}]
在更新有定义的条件下，两种损失均非增且下界为 $0$，因此损失数列收敛；这不足以推出因子数列一定收敛，更不足以推出全局最优。在严格正的内部，乘法更新不改变因子等价于相应梯度为零；在边界上，零元素即使梯度方向要求它增加，也可能因乘法形式保持为零，不能把所有不动点都称为约束最小点。
\end{notealert}

更具体地，在目标可微的位置，局部最小点必须满足非负约束的一阶 KKT 条件：$W,H\geq0$，$\nabla_WJ,\nabla_HJ\geq0$，且
\begin{equation}
W\odot\nabla_WJ=0,\qquad H\odot\nabla_HJ=0,
\label{eq:lh-lsa-nmf-kkt}
\end{equation}
其中 $\odot$ 表示逐元素乘法。在正元素上梯度必须为零，在零元素上梯度必须非负；这些只是一般非凸问题的必要条件。实际可结合损失变化与一阶条件残差检查停止结果，并以多个初值比较拟合，而不把某次迭代输出视为唯一的话题结构。

\section{本章小结与习题联系}
\label{sec:lh-lsa-summary}

本章的共同模型是 $X\approx TY$，关键在于同时学习潜在方向与文本系数。SVD 取 $T=U_k$、$Y=\Sigma_kV_k^{\mathsf T}$，在平方重构误差下给出全局最优的秩至多 $k$ 的近似，允许正负系数；NMF 取 $T=W$、$Y=H$，要求两因子非负，得到成分叠加的表示，但需要求解联合非凸优化。前者的系数由正交投影得到，后者的非负系数通常要通过约束优化得到，二者的“话题空间”几何含义并不完全相同。

教材习题 17.1 的计算见例~\ref{ex:lh-lsa-four-documents-svd}，结果同时展示同义词联系的改善与多义词混淆仍可能存在；习题 17.2 所需的散度损失算法，由式~\eqref{eq:lh-lsa-nmf-kl-update-w}、式~\eqref{eq:lh-lsa-nmf-kl-update-h} 及共同迭代流程给出。对习题 17.3，必须说明所用算法与稀疏性假设：稠密矩阵完整 SVD 的常规复杂度为 $O(mn\min\{m,n\})$，截断方法的开销则依赖所用数值算法和迭代次数，不能仅凭保留 $k$ 个方向就把完整分解的代价改写成 $O(mnk)$。例如采用块子空间迭代、每轮进行与 $X,X^{\mathsf T}$ 的乘法和正交化时，$q$ 轮的主要开销为 $O(q\,\operatorname{nnz}(X)k+q(m+n)k^2)$，其中 $\operatorname{nnz}(X)$ 是非零元素数，$q$ 还依赖谱间隙和精度要求。

对平方损失 NMF，按 $XH^{\mathsf T}$、$W^{\mathsf T}X$、$HH^{\mathsf T}$、$W^{\mathsf T}W$ 及相应乘积计算，一轮稠密交替更新的开销为 $O(mnk+(m+n)k^2)$，$t$ 轮再乘以 $t$；若 $X$ 稀疏，前一项可改为 $\operatorname{nnz}(X)k$。散度更新若显式形成 $WH$ 和逐元素比值，稠密实现每轮为 $O(mnk)$，并需要存储或分块处理 $m\times n$ 的重构量。因子本身占用 $O((m+n)k)$ 存储；复杂度描述须与实际矩阵乘法次序及中间量的存储方式一致。

对习题 17.4，\emph{主成分分析}（principal component analysis, PCA）与 SVD 型 LSA 都可通过奇异值分解寻找低维正交方向，但数据预处理和解释目标不同。以列为样本时，标准 PCA 先求均值 $\bar x=n^{-1}\sum_jx_j$ 并中心化，$X_c=X-\bar x\mathbf1_n^{\mathsf T}$，再对 $X_c$ 做 SVD；其方向最大化相对均值的方差，重构为 $\bar x+U_kU_k^{\mathsf T}(x-\bar x)$，得到仿射子空间上的近似。本文的 LSA 通常直接分解词频或 TF-IDF 矩阵 $X$，重构为 $U_kU_k^{\mathsf T}x$，得到经过原点的线性子空间近似。中心化可能破坏词矩阵的稀疏性并引入负值，故不能把标准 PCA 的预处理无条件移到本章模型中；若对同一个中心化矩阵使用相同的截断 SVD，两者在线性代数计算上才一致。

% 教材来源：机器学习/textbooks/机器学习（分章原文）/机器学习-19.pdf。
% 对应教材第 19 章，原书第 351--383 页；保留五节及其主要小节。

\fi

\selectchapter{18}
\ifcompilechapter
\chapter{马尔可夫链蒙特卡罗法}
\label{ch:lh-mcmc}

复杂概率模型常能写出未归一化的密度，却难以计算归一化常数、边缘分布或后验期望。\emph{蒙特卡罗法}（Monte Carlo method）以随机样本的平均近似积分；\emph{马尔可夫链蒙特卡罗法}（Markov chain Monte Carlo，MCMC）进一步构造以目标分布为平稳分布的马尔可夫链，用一条相关的样本轨迹完成计算。关键是分别回答两个问题：所构造的链是否保持目标分布，以及从指定初值出发是否能够充分探索该分布。前者由不变性或细致平衡证明，后者还需要不可约性等遍历条件。

本章依次讨论一般蒙特卡罗方法、马尔可夫链、MCMC 的计算框架、Metropolis--Hastings 算法和吉布斯抽样。记 $S$ 为状态空间，$X_t$ 为第 $t$ 步的随机状态，$x^{(t)}$ 为其一次实现；$p$ 表示目标密度或概率质量函数，$\pi$ 表示相应概率分布。积分在离散空间中相应改为求和；$\mathbb E_\pi$ 与 $\operatorname{Var}_\pi$ 分别表示在 $\pi$ 下的期望与方差。除特别说明外，所有抽样步骤使用新的、相互独立的辅助随机数。

\section{蒙特卡罗法}
\label{sec:lh-mcmc-monte-carlo}

\subsection{随机抽样}
\label{subsec:lh-mcmc-random-sampling}

随机抽样（random sampling）的任务是在给定分布的条件下生成随机变量的实现，与根据已有数据估计未知分布的任务不同。直接抽样可利用已知分布的专门生成方法；例如连续严格递增的分布函数 $F$ 满足 $F^{-1}(U)\sim F$，其中 $U\sim\operatorname{Unif}(0,1)$。复杂目标未必能直接抽样，此时可以借助容易生成的建议分布。

\begin{definition}[接受--拒绝抽样]
\label{def:lh-mcmc-rejection}
设 $p$ 与 $q$ 是同一参考测度下的概率密度，$q$ 可直接抽样，且存在已知有限常数 $c$，使 $p(x)\leq c q(x)$ 几乎处处成立。独立生成 $Y\sim q$ 和 $U\sim\operatorname{Unif}(0,1)$，若
\begin{equation}
U\leq\frac{p(Y)}{c q(Y)},
\label{eq:lh-mcmc-rejection-rule}
\end{equation}
则接受 $Y$，否则重新生成一对 $(Y,U)$；重复至获得所需的 $N$ 个接受值。这称为\emph{接受--拒绝抽样}（accept--reject sampling）。
\end{definition}

\begin{proposition}[接受值的分布与抽样代价]
\label{prop:lh-mcmc-rejection-correctness}
在定义~\ref{def:lh-mcmc-rejection} 的条件下，单次尝试的接受概率为 $1/c$，接受值服从 $p$，重复独立尝试得到的接受值相互独立；生成 $N$ 个接受值所需尝试次数的期望为 $Nc$。
\end{proposition}
\begin{proof}
对任意可测集合 $A$，由条件期望及 $U$ 的独立性，
\[
P(Y\in A,\text{接受})
=\int_A q(y)\frac{p(y)}{c q(y)}\,\mathrm dy
=\frac1c\int_A p(y)\,\mathrm dy.
\]
取 $A=S$ 得到接受概率 $1/c$，再除以此概率即得接受条件下的目标分布。每次接受后继续使用独立的新尝试，故各接受值独立；等待一次成功的尝试数服从成功概率为 $1/c$ 的几何分布，其期望为 $c$，累加便得 $Nc$。
\end{proof}

包络条件同时要求 $q$ 覆盖目标的支撑；$c$ 越大，拒绝越多。若仅知未归一化函数 $\widetilde p$，但已知 $0<Z=\int_S\widetilde p(x)\,\mathrm dx<\infty$ 且 $\widetilde p\leq Mq$，则接受概率改为 $\widetilde p(Y)/(Mq(Y))$，仍生成 $p=\widetilde p/Z$ 的样本，整体接受率为 $Z/M$。算法不必知道 $Z$，但仍需一个有效的全局包络 $M$；高维问题中这常是主要困难。

\subsection{数学期望估计}
\label{subsec:lh-mcmc-expectation}

设 $X_1,\ldots,X_N$ 独立同分布于 $\pi$，$f:S\to\mathbb R$ 可测且 $\mathbb E_\pi|f(X)|<\infty$。强大数定律给出
\begin{equation}
\widehat\mu_N=\frac1N\sum_{i=1}^N f(X_i)
\xrightarrow[N\to\infty]{\text{几乎必然}}
\mu=\mathbb E_\pi f(X).
\label{eq:lh-mcmc-iid-law}
\end{equation}
这里 $\mu$ 是总体量，$\widehat\mu_N$ 是由抽样引入随机性的估计量，$N^{-1}\sum_i f(x_i)$ 则是一次运行得到的数值。若还满足 $\sigma_f^2=\operatorname{Var}_\pi(f(X))<\infty$，则
\begin{equation}
\mathbb E\widehat\mu_N=\mu,\qquad
\operatorname{Var}(\widehat\mu_N)=\frac{\sigma_f^2}{N}.
\label{eq:lh-mcmc-iid-variance}
\end{equation}
当 $\sigma_f^2>0$ 时，独立同分布中心极限定理进一步给出 $\sqrt N(\widehat\mu_N-\mu)/\sigma_f$ 依分布趋于标准正态。因此有限方差下典型随机误差按 $N^{-1/2}$ 缩小；误差缩小到原来的十分之一，通常需要约一百倍样本。该阶数不显含维度，但函数方差与生成一个样本的成本都可能随维度急剧增长，不能据此断言高维计算没有困难。

\subsection{积分计算}
\label{subsec:lh-mcmc-integration}

设 $h$ 可积，选取满足 $h(x)\ne0$ 的区域上几乎处处为正、且能够抽样的密度 $p$，令 $f=h/p$，在 $h=p=0$ 处可令 $f=0$，则
\begin{equation}
I=\int_S h(x)\,\mathrm dx
=\int_S\frac{h(x)}{p(x)}p(x)\,\mathrm dx
=\mathbb E_p\!\left[\frac{h(X)}{p(X)}\right]
\approx\frac1N\sum_{i=1}^N\frac{h(X_i)}{p(X_i)}.
\label{eq:lh-mcmc-integral}
\end{equation}
这称为\emph{蒙特卡罗积分}（Monte Carlo integration）。支撑覆盖保证比值有意义，$\int |h|<\infty$ 保证强大数定律适用；若要使用式~\eqref{eq:lh-mcmc-iid-variance} 的误差分析，还需 $\int h(x)^2/p(x)\,\mathrm dx<\infty$。改变抽样分布会改变估计方差，因此选择 $p$ 也是计算方法的一部分。

\begin{example}
\label{ex:lh-mcmc-integrals}
分别把 $I_1=\int_0^1e^{-x^2}\,\mathrm dx$ 与 $I_2=\int_{\mathbb R}x\varphi(x)\,\mathrm dx$ 写成蒙特卡罗估计，其中 $\varphi(x)=(2\pi)^{-1/2}e^{-x^2/2}$。
\end{example}
\begin{solve}
对第一个积分，取 $U_i$ 独立同分布于 $\operatorname{Unif}(0,1)$，则 $\widehat I_{1,N}=N^{-1}\sum_i e^{-U_i^2}$，被平均的变量有界，故估计几乎必然收敛且具有有限方差。对第二个积分，取 $Z_i$ 独立同分布于 $N(0,1)$，则 $\widehat I_{2,N}=N^{-1}\sum_i Z_i$，其期望为 $0$、方差为 $1/N$；积分本身由对称性等于 $0$。一次有限抽样的估计一般不恰为 $0$，也不保证每增加一个样本就更接近真值。
\end{solve}

教材还提到\emph{重要性抽样}（importance sampling）。它使用容易抽样的密度 $q$，把目标期望写成 $\mathbb E_p f(X)=\mathbb E_q[f(X)p(X)/q(X)]$；要求 $q$ 覆盖相关支撑及相应可积性。若目标只知到比例常数，可用自归一化比值 $\sum_i w_i f(Y_i)/\sum_i w_i$，其中 $Y_i\sim q$、$w_i=\widetilde p(Y_i)/q(Y_i)$。该比值在适当大数定律条件下一致，有限样本通常有偏；它与接受--拒绝法以及后文的 MCMC 是不同的抽样计算方案。

\section{马尔可夫链}
\label{sec:lh-mcmc-markov-chain}

\subsection{基本定义}
\label{subsec:lh-mcmc-markov-definition}

\begin{definition}[马尔可夫链与时间齐次性]
\label{def:lh-mcmc-markov}
设随机过程 $(X_t)_{t\geq0}$ 取值于同一状态空间 $S$。若对任意可测集合 $A$，有
\begin{equation}
P(X_{t+1}\in A\mid X_0,\ldots,X_t)
=P(X_{t+1}\in A\mid X_t)\quad\text{几乎必然},
\label{eq:lh-mcmc-markov-property}
\end{equation}
则称其为\emph{马尔可夫链}（Markov chain）。若右侧可写为与 $t$ 无关的 $K(X_t,A)$，则称链为\emph{时间齐次的}（time homogeneous），$K$ 为其转移核，$X_0$ 的分布 $\mu_0$ 为初始分布。本章均讨论时间齐次的一阶链。
\end{definition}

马尔可夫性表达的是给定当前状态后的条件独立关系，不表示相邻状态独立。给定初始分布与转移规则便确定了整条链的有限维联合分布。若一个过程的下一步只依赖最近 $r$ 步，则可把 $(X_{t-r+1},\ldots,X_t)$ 视为一个增广状态，转成一阶链；代价是状态空间增大。把天气或词序列建模为马尔可夫链时，这种依赖范围是需要检验的模型假设，并非数据自动具有的性质。

\subsection{离散状态马尔可夫链}
\label{subsec:lh-mcmc-discrete-chain}

设 $S$ 有限或可数。沿用教材的\emph{列随机矩阵}约定，令
\begin{equation}
P_{ij}=P(X_{t+1}=i\mid X_t=j),\qquad
P_{ij}\geq0,\qquad\sum_iP_{ij}=1.
\label{eq:lh-mcmc-column-matrix}
\end{equation}
因此列索引是出发状态，行索引是到达状态。设列向量 $\mu_t$ 的第 $i$ 个分量为 $P(X_t=i)$，由全概率公式，
\begin{equation}
(\mu_{t+1})_i=\sum_jP_{ij}(\mu_t)_j,
\qquad\mu_{t+1}=P\mu_t,\qquad\mu_t=P^t\mu_0.
\label{eq:lh-mcmc-distribution-evolution}
\end{equation}
$P^t$ 表示矩阵的 $t$ 次幂，转置另写为 $P^{\mathsf T}$。条件于中间状态还得到 Chapman--Kolmogorov 关系 $P^{r+s}_{ij}=\sum_kP^r_{ik}P^s_{kj}$。若改用行随机矩阵，状态分布应相应写为行向量并从右侧乘矩阵，不能混用两套约定。

\begin{definition}[平稳分布]
\label{def:lh-mcmc-stationarity}
若概率列向量 $\pi$ 满足
\begin{equation}
P\pi=\pi,\qquad\pi_i\geq0,\qquad\sum_i\pi_i=1,
\label{eq:lh-mcmc-stationary-system}
\end{equation}
则称 $\pi$ 为链的\emph{平稳分布}（stationary distribution）。以 $\pi$ 初始化时，每个 $X_t$ 都具有分布 $\pi$，但这些随机变量仍可能高度相关。
\end{definition}

求平稳分布就是解式~\eqref{eq:lh-mcmc-stationary-system}，其中非负性与归一化缺一不可。平稳分布的存在、唯一性，以及从其他初值出发是否收敛到它，是三个不同问题。

\begin{example}
\label{ex:lh-mcmc-stationary-matrices}
求下列两个列随机矩阵的平稳分布，并比较唯一性：
\[
P=\begin{pmatrix}1/2&1/2&1/4\\1/4&0&1/4\\1/4&1/2&1/2\end{pmatrix},
\qquad
Q=\begin{pmatrix}1&1/3&0\\0&1/3&0\\0&1/3&1\end{pmatrix}.
\]
\end{example}
\begin{solve}
对 $P$，第二行方程给出 $\pi_2=(\pi_1+\pi_3)/4$，与归一化联立得 $\pi_2=1/5$；第一行方程再给出 $2\pi_1-\pi_3=2/5$，故唯一解为 $\pi=(2/5,1/5,2/5)^{\mathsf T}$。对 $Q$，第二行方程为 $\pi_2=\pi_2/3$，故 $\pi_2=0$，其余条件只要求 $\pi_1+\pi_3=1$。全部平稳分布为 $(a,0,1-a)^{\mathsf T}$，$0\leq a\leq1$；状态 $1$ 与 $3$ 分别吸收，长期行为保留了初始质量如何进入两个吸收类的信息。
\end{solve}

\subsection{连续状态马尔可夫链}
\label{subsec:lh-mcmc-continuous-chain}

对 $S\subseteq\mathbb R^d$，转移核 $K(x,A)$ 要求：固定 $x$ 时它是关于可测集合 $A$ 的概率测度，固定 $A$ 时它关于 $x$ 可测。它表示 $P(X_{t+1}\in A\mid X_t=x)$。若存在转移密度 $k(x,y)$，则
\begin{equation}
K(x,A)=\int_A k(x,y)\,\mathrm dy,\qquad
k(x,y)\geq0,\qquad\int_S k(x,y)\,\mathrm dy=1.
\label{eq:lh-mcmc-transition-density}
\end{equation}
分布 $\pi$ 的不变性要求
\begin{equation}
\pi(A)=\int_S K(x,A)\,\pi(\mathrm dx)
\quad\text{对所有可测集合 }A.
\label{eq:lh-mcmc-invariant-kernel}
\end{equation}
当 $\pi$ 具有密度 $p$ 且 $K$ 具有密度 $k$ 时，上式化为 $p(y)=\int_Sp(x)k(x,y)\,\mathrm dx$，几乎处处成立。转移核未必有普通的密度：例如拒绝候选后停在原地，会产生位于 $x$ 的点质量。因此式~\eqref{eq:lh-mcmc-invariant-kernel} 比仅用密度的表达更一般，也是后文证明的基本形式。

\subsection{马尔可夫链的性质}
\label{subsec:lh-mcmc-markov-properties}

\begin{definition}[不可约、周期与正常返]
\label{def:lh-mcmc-ergodic-conditions}
对有限或可数状态链，若任意 $i,j\in S$ 都存在 $n\geq0$ 使 $(P^n)_{ij}>0$，称链\emph{不可约}（irreducible）；这允许多步可达，不要求每个一步转移概率都为正。对可返回的状态 $i$，定义周期
\[
d(i)=\gcd\{n\geq1:(P^n)_{ii}>0\}.
\]
若每个状态周期为 $1$，称链\emph{非周期}（aperiodic）。从 $i$ 出发定义首次返回时间 $\tau_i^+=\inf\{t\geq1:X_t=i\}$，约定空集的下确界为 $+\infty$；若 $P_i(\tau_i^+<\infty)=1$，称 $i$ 常返，若进一步 $\mathbb E_i\tau_i^+<\infty$，称 $i$ 为\emph{正常返}状态（positive recurrent）。不可约链的状态共享常返类型，故可直接称该链正常返。
\end{definition}

不可约链中的状态具有相同周期：若存在从 $i$ 到 $j$ 的 $a$ 步路径和从 $j$ 到 $i$ 的 $b$ 步路径，则 $a+b$ 是 $i$ 的可返回时长，且对 $j$ 的任一可返回时长 $n$，$a+n+b$ 也是 $i$ 的可返回时长。因此 $d(i)$ 整除所有这样的 $n$，故 $d(i)\mid d(j)$；交换 $i,j$ 得周期相等。特别地，不可约链只要某状态有正自环便非周期。正常返则控制返回的平均等待时间；有限不可约链自动正常返，无限链需要另行检查。

\begin{notealert}[type=caution,title={正常返不能用首次返回概率的正极限定义}]
设 $f_i(n)=P_i(\tau_i^+=n)$，则 $\sum_{n\geq1}f_i(n)\leq1$，所以必有 $f_i(n)\to0$。因此教材关于首次到达概率具有正极限的表述不能作为正常返的定义；应使用定义~\ref{def:lh-mcmc-ergodic-conditions} 中的有限平均返回时间条件。常返但平均返回时间无穷的状态称为零常返（null recurrent）。
\end{notealert}

\begin{theorem}[离散链的平稳分布与遍历结论]
\label{thm:lh-mcmc-ergodic}
设有限或可数状态链不可约且正常返，则存在唯一平稳概率分布 $\pi$，且每个 $\pi_i>0$。对任意初始状态和满足 $\sum_i\pi_i|f(i)|<\infty$ 的函数，
\begin{equation}
\frac1N\sum_{t=1}^N f(X_t)
\xrightarrow[N\to\infty]{\text{几乎必然}}
\sum_i\pi_i f(i).
\label{eq:lh-mcmc-ergodic-average}
\end{equation}
若链还非周期，则对任意 $i,j$，
\begin{equation}
\lim_{t\to\infty}(P^t)_{ij}=\pi_i.
\label{eq:lh-mcmc-marginal-limit}
\end{equation}
对不可约可数状态链，存在平稳概率分布也反过来蕴含正常返；有限不可约链总满足该条件。
\end{theorem}

这里引用马尔可夫链的标准遍历结论，对应教材的平稳分布与遍历定理；其完整证明依赖返回周期分解和更新理论，本章不展开。\emph[caution]{时间平均的收敛与单个时刻分布的收敛应分开。}例如确定性循环 $1\to2\to3\to1$ 的平稳分布唯一且均匀，但从固定状态出发的分布始终轮换，不收敛；沿轨迹的三个状态所占比例却趋于 $1/3$，因此时间平均仍收敛。连续状态下单点通常具有零概率，不能直接照搬逐点命中的定义；使用一般状态空间的遍历结论时，需检查相应条件，例如正 Harris 常返（positive Harris recurrence，即对适当非零参考测度下每个正测度集合，从任意初始状态都以概率 $1$ 最终进入，并且具有不变概率分布）及用于分布收敛的非周期性。

\begin{example}
\label{ex:lh-mcmc-birth-death}
考虑 $S=\{0,1,2,\ldots\}$ 上的链：从 $i\geq1$ 以概率 $a$ 移到 $i-1$，以概率 $b$ 移到 $i+1$；从 $0$ 以概率 $a$ 留在原地，以概率 $b$ 移到 $1$，其中 $a,b>0$、$a+b=1$。判断何时存在平稳概率分布。
\end{example}
\begin{solve}
平稳时对有限集合 $\{0,\ldots,i\}$ 的流入与流出相等，故 $b\pi_i=a\pi_{i+1}$，于是 $\pi_i=\pi_0(b/a)^i$。几何级数仅在 $a>b$ 时可归一化，此时 $\pi_i=(1-b/a)(b/a)^i$；该不可约链存在平稳概率分布，因而正常返。$a\leq b$ 时任何非零候选序列都不可归一化，因此没有平稳概率分布。边界 $a=b$ 的链实际为零常返：若 $h_i$ 为从 $i$ 出发先到 $0$ 而非 $m>i$ 的概率，则 $h_0=1$、$h_m=0$ 及 $h_i=(h_{i-1}+h_{i+1})/2$ 给出 $h_i=1-i/m$，令 $m\to\infty$ 得到最终到达 $0$ 的概率为 $1$；结合各状态互通及不存在平稳概率分布可知其为零常返。存在正自环虽然保证非周期，却不能替代正常返条件。
\end{solve}

\begin{definition}[细致平衡与可逆性]
\label{def:lh-mcmc-reversibility}
在列随机矩阵约定下，若概率分布 $\pi$ 满足
\begin{equation}
\pi_jP_{ij}=\pi_iP_{ji}\qquad\text{对所有 }i,j,
\label{eq:lh-mcmc-detailed-balance}
\end{equation}
则称链关于 $\pi$ \emph{可逆}（reversible），上式称\emph{细致平衡方程}（detailed balance equation）。一般转移核的对应条件是测度 $\pi(\mathrm dx)K(x,\mathrm dy)$ 关于交换 $x,y$ 对称。
\end{definition}

\begin{proposition}[细致平衡蕴含不变性]
\label{prop:lh-mcmc-balance-invariance}
满足细致平衡的概率分布是平稳分布，但该条件本身不保证不可约、非周期或平稳分布唯一。
\end{proposition}
\begin{proof}
在离散情形，对式~\eqref{eq:lh-mcmc-detailed-balance} 关于出发状态 $j$ 求和，得
\[
(P\pi)_i=\sum_jP_{ij}\pi_j
=\pi_i\sum_jP_{ji}=\pi_i,
\]
最后一项为 $P$ 的第 $i$ 列之和。一般情形对对称测度的第一坐标积分，同样得到式~\eqref{eq:lh-mcmc-invariant-kernel}。反面例子中，$P=I_2$ 对每个概率向量都满足细致平衡，却可约且平稳分布不唯一；矩阵 $\left(\begin{smallmatrix}0&1\\1&0\end{smallmatrix}\right)$ 关于均匀分布可逆，却有周期 $2$。
\end{proof}

可逆性表示平稳状态下正向路径与反向路径具有相同概率，强于单纯的平稳性。前述三状态确定性循环保持均匀分布，却不满足细致平衡，说明细致平衡只是证明平稳性的充分条件。教材由可逆性直接推出唯一平稳分布及遍历性、以及相关习题要求证明可逆必不可约的说法，均需加上额外条件；上述反例已说明不能无条件使用。

\section{马尔可夫链蒙特卡罗法}
\label{sec:lh-mcmc-framework}

\subsection{基本想法}
\label{subsec:lh-mcmc-idea}

设目标分布为 $\pi$，密度为 $p=\widetilde p/Z$，其中 $\widetilde p\geq0$、$0<Z<\infty$，而 $Z$ 难以直接计算。MCMC 不要求每一步直接生成一个独立的目标样本，而是选择转移核 $K$，使 $\pi$ 不变，并让链在目标支撑内充分探索。若适用的遍历定理成立，则沿轨迹平均可以近似 $\mathbb E_\pi f$；若再满足分布收敛条件，晚期单个状态的边缘分布也接近 $\pi$。这两种结论不要求轨迹中的状态相互独立。

选定初始状态 $x^{(0)}$ 后，前 $m$ 步常被称为\emph{燃烧期}或预热期（burn-in），之后保留 $N$ 个状态，定义
\begin{equation}
\widehat\mu_{m,N}=\frac1N\sum_{t=m+1}^{m+N}f(X_t),
\qquad
\widehat\pi_{m,N}(A)=\frac1N\sum_{t=m+1}^{m+N}\mathbf1\{X_t\in A\}.
\label{eq:lh-mcmc-estimators}
\end{equation}
在遍历条件及可积性成立时，固定 $m$ 后令 $N\to\infty$，这两个量分别几乎必然收敛到 $\mathbb E_\pi f$ 与 $\pi(A)$；删去固定有限项不改变时间平均极限。教材用总迭代数 $n$ 表示终点，其均值分母为 $n-m$；这里设 $n=m+N$，以区别预热长度与实际保留数。

\emph[caution]{预热是控制初始分布影响的计算安排，不能把有限样本自动变成精确、独立、无偏的目标样本。}若 $X_0\sim\pi$，则由不变性每个时刻均服从 $\pi$，可积函数的样本均值无偏；若从一般固定初值出发，有限长度均值通常有初始化偏差。渐近遍历性没有直接给出充分的 $m$，不同目标和转移核的混合速度可能差别很大。

\subsection{基本步骤}
\label{subsec:lh-mcmc-steps}

MCMC 的输入是目标未归一化密度 $\widetilde p$、待求期望的函数 $f$、初值、预热长度 $m\geq0$ 与保留数 $N\geq1$。先构造转移核并核查不变性及遍历条件，再模拟 $m+N$ 步，最后输出后 $N$ 个状态、式~\eqref{eq:lh-mcmc-estimators} 的均值和相应误差诊断。终止于预先规定的计算预算，或在可信的误差评估达到要求后终止；仅观察到均值暂时稳定，不能据此证明已访问所有重要区域。相距较远的初值所产生的多条链、轨迹和自相关的比较，可以帮助发现探索不足，但也不是有限运行已经收敛的数学证明。

相关性改变了误差公式。设链以 $\pi$ 初始化，$\sigma_f^2=\operatorname{Var}_\pi f(X)>0$ 且有限，记
$\rho_k=\operatorname{Cov}(f(X_0),f(X_k))/\sigma_f^2$。展开和的方差得到精确等式
\begin{equation}
\operatorname{Var}(\widehat\mu_{0,N})
=\frac{\sigma_f^2}{N}
\left[1+2\sum_{k=1}^{N-1}\left(1-\frac{k}{N}\right)\rho_k\right].
\label{eq:lh-mcmc-correlated-variance}
\end{equation}
其中滞后 $k$ 的协方差出现 $N-k$ 次。若 $\sum_{k\geq1}|\rho_k|<\infty$，则支配收敛给出
\begin{equation}
N\operatorname{Var}(\widehat\mu_{0,N})\longrightarrow
\sigma_f^2\tau_f,\qquad
\tau_f=1+2\sum_{k=1}^{\infty}\rho_k.
\label{eq:lh-mcmc-autocorrelation-time}
\end{equation}
当 $0<\tau_f<\infty$ 时，称 $\tau_f$ 为该函数的积分自相关时间（integrated autocorrelation time），并常用\emph{有效样本量}（effective sample size）$N_{\mathrm{eff}}\approx N/\tau_f$ 描述该函数均值的方差相当于多少独立样本。它依赖所考察的函数，不能给一条链赋予对所有统计量通用的精确有效样本数。式~\eqref{eq:lh-mcmc-autocorrelation-time} 只陈述方差渐近关系；正态置信区间还需另外满足马尔可夫链中心极限定理的条件。

每隔 $\ell$ 步保留一个状态称为间隔取样或稀疏化（thinning），其新转移核是 $K^\ell$。这可能减弱相关性并节省存储，却一般不能使样本精确独立，也会丢弃计算得到的信息。MCMC 的效率应按单位计算成本的误差或有效样本量评价；不能只看接受率，也不能笼统断言它总优于接受--拒绝法。

\subsection{马尔可夫链蒙特卡罗法与统计学习}
\label{subsec:lh-mcmc-statistical-learning}

设 $D$ 为观测数据，$\Theta$ 为未知参数，先验密度为 $p(\theta)$、似然为 $p(D\mid\theta)$。贝叶斯公式给出
\begin{equation}
p(\theta\mid D)
=\frac{p(D\mid\theta)p(\theta)}{p(D)},
\qquad
p(D)=\int p(D\mid\vartheta)p(\vartheta)\,\mathrm d\vartheta.
\label{eq:lh-mcmc-bayes}
\end{equation}
分母是归一化常数，也称模型证据；它须为有限正数，后验才是概率分布。对后验抽样，可以把 $\widetilde p(\theta)=p(D\mid\theta)p(\theta)$ 作为目标核。若还含隐变量 $Z$，则可对联合后验 $p(\theta,z\mid D)$ 抽样，忽略样本的 $z$ 分量即可获得参数边缘分布的相关样本，不必显式计算积分 $p(\theta\mid D)=\int p(\theta,z\mid D)\,\mathrm dz$。

后验均值、方差以及预测分布都可由函数平均近似。例如若未来观测 $Y_*$ 给定参数后与 $D$ 条件独立，则后验预测密度满足
\begin{equation}
p(y_*\mid D)=\int p(y_*\mid\theta)p(\theta\mid D)\,\mathrm d\theta
\approx\frac1N\sum_{t=m+1}^{m+N}p(y_*\mid\theta^{(t)}).
\label{eq:lh-mcmc-posterior-prediction}
\end{equation}
这些是关于后验分布的积分计算，与寻找后验众数的最大后验估计不同。MH 中归一化常数相消使抽样成为可能，却不意味着该常数本身已经被估计；若目标是模型证据，还需要额外的估计方法。

\section{Metropolis--Hastings 算法}
\label{sec:lh-mcmc-mh}

\subsection{基本原理}
\label{subsec:lh-mcmc-mh-principle}

设当前状态为 $x$，由建议密度 $q(y\mid x)$ 生成候选 $y$。若总接受建议，目标分布一般不会保持不变；Metropolis--Hastings（MH）算法用\emph{接受概率}（acceptance probability）修正正反两个方向的概率流。对 $\widetilde p(x)>0$ 且实际可能提出的候选，令
\begin{equation}
\alpha(x,y)=\min\left\{1,
\frac{\widetilde p(y)q(x\mid y)}{\widetilde p(x)q(y\mid x)}\right\}.
\label{eq:lh-mcmc-mh-acceptance}
\end{equation}
这里 $q(y\mid x)$ 是从 $x$ 提出 $y$ 的密度，$q(x\mid y)$ 是反向建议密度。若分子为零而分母为正，则接受概率为零；分母为零的候选在当前建议下不会以正概率出现，可在这些无关点任意约定。初值选在目标正密度区域内，可避免从零目标密度状态计算比值。正反式中的归一化常数 $Z$ 相消。

对连续建议，完整转移核为
\begin{equation}
K(x,A)=\int_Aq(y\mid x)\alpha(x,y)\,\mathrm dy
+r(x)\mathbf1\{x\in A\},\qquad
r(x)=1-\int_Sq(y\mid x)\alpha(x,y)\,\mathrm dy.
\label{eq:lh-mcmc-mh-full-kernel}
\end{equation}
第二项是拒绝后留在原地的概率。离散情形则对 $y\ne x$ 定义 $K_{yx}=q(y\mid x)\alpha(x,y)$，并令 $K_{xx}=1-\sum_{y\ne x}K_{yx}$，得到列和为 $1$ 的随机矩阵。单独写 $q\alpha$ 只表示接受移动的部分，通常不能作为完整转移核。

\begin{proposition}[MH 的细致平衡]
\label{prop:lh-mcmc-mh-balance}
设 $p=\widetilde p/Z$ 是正规概率密度，建议和接受规则如上，则式~\eqref{eq:lh-mcmc-mh-full-kernel} 关于 $\pi$ 可逆，因而 $\pi$ 是其平稳分布。
\end{proposition}
\begin{proof}
在有定义的接受移动上，
\begin{equation}
p(x)q(y\mid x)\alpha(x,y)
=\min\{p(x)q(y\mid x),p(y)q(x\mid y)\}
=p(y)q(x\mid y)\alpha(y,x).
\label{eq:lh-mcmc-mh-flow}
\end{equation}
该式关于 $x,y$ 对称。拒绝部分集中在对角线 $x=y$，其测度 $\pi(\mathrm dx)r(x)\delta_x(\mathrm dy)$ 同样对称，其中 $\delta_x$ 为位于 $x$ 的单位点质量。因此完整转移核满足细致平衡，由命题~\ref{prop:lh-mcmc-balance-invariance} 得不变性。
\end{proof}

上述证明只解决平稳性。建议的可达性、目标支撑是否连通，以及接受修正是否阻断路径，仍需检查；建议矩阵不可约并不总能使 MH 链不可约，因为只有单向建议而没有反向建议时，接受率可能为零。在有限目标支撑上，若接受后的链不可约，则自动正常返；若还有正自环或其他非周期依据，才可使用分布收敛结论。连续情形需检查相应的一般状态空间条件。必要时可将转移核改为 $\tfrac12 I+\tfrac12 K$，其中 $I(x,A)=\mathbf1\{x\in A\}$，加入留步保持不变性；它能处理周期性，却不能修复可约性。

建议分布有两种常用形式。若 $q(y\mid x)=q(x\mid y)$，接受率化为
\begin{equation}
\alpha(x,y)=\min\{1,\widetilde p(y)/\widetilde p(x)\}.
\label{eq:lh-mcmc-metropolis}
\end{equation}
这称为 Metropolis 选择，例如 $Y=x+\varepsilon$，其中 $\varepsilon\sim N(0,\Sigma)$、$\Sigma$ 为固定正定矩阵。步幅过小会使接受率高却移动缓慢，步幅过大则可能频繁拒绝；高接受率本身不代表误差小。若 $q(y\mid x)=q(y)$，则称\emph{独立建议}（independence proposal），有
\begin{equation}
\alpha(x,y)=\min\{1,w(y)/w(x)\},\qquad
w(x)=\widetilde p(x)/q(x).
\label{eq:lh-mcmc-independence-proposal}
\end{equation}
候选值的生成不依赖当前状态，但接受与否仍依赖它，所得链通常不独立。建议必须覆盖目标相关支撑，并应尽量避免在目标重要区域给出过小概率。

对于多元状态 $x=(x_1,\ldots,x_d)$，记 $x_{-j}$ 为除第 $j$ 个坐标外的其余坐标。\emph{满条件分布}（full conditional distribution）指给定其余所有分量后的分布，例如
\begin{equation}
p(x_j\mid x_{-j})
=\frac{\widetilde p(x_j,x_{-j})}
{\int\widetilde p(u,x_{-j})\,\mathrm du},
\label{eq:lh-mcmc-full-conditional}
\end{equation}
要求分母有限且为正。把一组坐标 $B$ 留作变量同样得到块满条件分布。若 $x,y$ 仅在第 $j$ 个坐标不同，即 $y_{-j}=x_{-j}$，则
\begin{equation}
\frac{p(y)}{p(x)}
=\frac{p(y_j\mid x_{-j})}{p(x_j\mid x_{-j})}.
\label{eq:lh-mcmc-conditional-ratio}
\end{equation}
这里固定的其余坐标必须相同；一般两个任意点之间不能省略边缘密度比。

\begin{notealert}[type=caution,title={形式上的条件密度不保证联合分布正规}]
以 $\widetilde p(x_1,x_2)=\exp\{-\tfrac12(x_1-1)^2(x_2-1)^2\}$ 为例，固定 $x_2\ne1$ 时可形式上识别出 $N(1,(x_2-1)^{-2})$ 的条件形状；但在整个 $\mathbb R^2$ 上，Tonelli 定理给出
\[
\int_{\mathbb R^2}\widetilde p\,\mathrm dx_1\mathrm dx_2
=\sqrt{2\pi}\int_{\mathbb R}\frac{\mathrm dx_2}{|x_2-1|}
=\infty.
\]
因此教材例 19.9 的这一乘积平方型表达式，在未补充支撑限制或其他因子的情况下不能归一化为联合概率密度；$x_2=1$ 时相应的一维核也不可归一化。识别满条件形式前，应先确认联合目标及所用条件分布均正规。
\end{notealert}

\subsection{Metropolis--Hastings 算法}
\label{subsec:lh-mcmc-mh-algorithm}

输入 $\widetilde p$、建议密度 $q(\cdot\mid\cdot)$、函数 $f$、满足 $\widetilde p(x^{(0)})>0$ 的初值、$m\geq0$ 与 $N\geq1$。对 $t=1,\ldots,m+N$，依次执行：在当前状态 $x=x^{(t-1)}$ 下生成候选 $y\sim q(\cdot\mid x)$，计算式~\eqref{eq:lh-mcmc-mh-acceptance}，独立生成 $u\sim\operatorname{Unif}(0,1)$，再令
\begin{equation}
x^{(t)}=\begin{cases}
y,&u\leq\alpha(x,y),\\
x^{(t-1)},&u>\alpha(x,y).
\end{cases}
\label{eq:lh-mcmc-mh-update}
\end{equation}
迭代完成后输出 $x^{(m+1)},\ldots,x^{(m+N)}$ 及函数均值。与接受--拒绝法不同，\emph[caution]{MH 每次拒绝后仍记下一次原状态，重复状态必须参与平均}；只保留接受时刻通常会改变状态权重，产生另一条跳链，不能直接用同一均值公式估计目标期望。

计算时宜使用对数接受比
\[
\Delta=\log\widetilde p(y)-\log\widetilde p(x)
+\log q(x\mid y)-\log q(y\mid x),
\]
并比较 $\log u\leq\min\{0,\Delta\}$，以减少指数上溢或下溢；零密度对应直接拒绝或相应的负无穷对数。若一次建议生成和建议密度计算的成本为 $C_q$、一次目标核计算成本为 $C_p$，缓存当前目标值后每步通常需要 $O(C_p+C_q)$ 的计算，实际代价由目标结构决定。

\begin{example}
\label{ex:lh-mcmc-beta-posterior}
十次条件独立的伯努利试验中观测到四次成功。令成功率 $\Theta$ 的先验为 $\operatorname{Beta}(1,1)$，构造 MH 算法估计后验均值与方差，并给出可核查的解析值。
\end{example}
\begin{solve}
似然正比于 $\theta^4(1-\theta)^6$，先验在 $(0,1)$ 上为常数，故后验为 $\operatorname{Beta}(5,7)$。其中 $\operatorname{Beta}(a,b)$ 的密度为 $\theta^{a-1}(1-\theta)^{b-1}/B(a,b)$，$a,b>0$，$B(a,b)=\int_0^1u^{a-1}(1-u)^{b-1}\,\mathrm du$。初值取 $1/2$，提出 $y=x+\varepsilon$、$\varepsilon\sim N(0,s^2)$，其中 $s>0$；$y\notin(0,1)$ 时直接拒绝，否则使用
\[
\alpha(x,y)=\min\left\{1,
\frac{y^4(1-y)^6}{x^4(1-x)^6}\right\}.
\]
保留状态的均值估计 $\mathbb E(\Theta\mid D)$，二阶矩均值减去均值平方估计后验方差；后者在一般有限运行中不无偏。作为计算核对，由 Beta 积分的递推关系得
\[
\mathbb E(\Theta\mid D)=\frac5{12},\qquad
\operatorname{Var}(\Theta\mid D)=\frac{5\cdot7}{12^2\cdot13}=\frac{35}{1872}.
\]
这些是解析后验量，并非一次模拟所得数值。若改为不断重抽直到 $y\in(0,1)$，建议会变成依赖 $x$ 的截断正态，需重新计算正反建议密度，不能继续直接使用上述对称建议接受率。
\end{solve}

\subsection{单分量 Metropolis--Hastings 算法}
\label{subsec:lh-mcmc-component-mh}

若联合建议难以构造，可在每轮内依次更新各坐标。设第 $t$ 轮更新第 $j$ 个坐标之前，当前向量为
\begin{equation}
z^{(t,j-1)}=(x_1^{(t)},\ldots,x_{j-1}^{(t)},
x_j^{(t-1)},\ldots,x_d^{(t-1)}).
\label{eq:lh-mcmc-partial-sweep}
\end{equation}
记其中的第 $j$ 坐标为 $a=x_j^{(t-1)}$，其余坐标为 $z_{-j}$。从 $q_j(b\mid a,z_{-j})$ 提出新坐标 $b$，接受率为
\begin{equation}
\alpha_j=\min\left\{1,
\frac{p(b\mid z_{-j})q_j(a\mid b,z_{-j})}
{p(a\mid z_{-j})q_j(b\mid a,z_{-j})}\right\}.
\label{eq:lh-mcmc-component-acceptance}
\end{equation}
按此接受率令 $x_j^{(t)}=b$ 或 $a$，其他坐标暂不改变；完成 $j=1,\ldots,d$ 称为一次完整扫描（sweep），得到 $x^{(t)}$。式~\eqref{eq:lh-mcmc-conditional-ratio} 说明这就是将一般 MH 应用于条件目标，每一步可只计算随该坐标变化的因子。

每个坐标更新核都保持联合分布不变，故依次复合后仍保持该分布：若 $\pi K_j=\pi$，则 $\pi K_1\cdots K_d=\pi$，这里 $\pi K$ 表示式~\eqref{eq:lh-mcmc-invariant-kernel} 右侧的分布作用。\emph[caution]{保持同一分布的可逆核按固定次序复合，整体一般不再可逆}，但不变性依然成立。若在每个微步以与状态无关的固定概率随机选择坐标，则混合核保持各核共有的细致平衡。两种扫描方式仍分别需要检查遍历条件，逐坐标更新可能受支撑结构限制而无法跨越不同区域。

\section{吉布斯抽样}
\label{sec:lh-mcmc-gibbs}

\subsection{基本原理}
\label{subsec:lh-mcmc-gibbs-principle}

\emph{吉布斯抽样}（Gibbs sampling）适用于联合分布难以直接抽样，而各满条件分布容易抽样的情形。固定其余坐标 $z_{-j}$，直接从 $p(\cdot\mid z_{-j})$ 生成新的第 $j$ 个坐标。将其作为单分量 MH 的建议，即取
\[
q_j(b\mid a,z_{-j})=p(b\mid z_{-j}),\qquad
q_j(a\mid b,z_{-j})=p(a\mid z_{-j}),
\]
代入式~\eqref{eq:lh-mcmc-component-acceptance}，分子分母相等，故在相关正概率区域内 $\alpha_j=1$。这说明 Gibbs 是单分量 MH 的特例；其特点是用精确满条件抽样消除了接受判定，而不是生成了独立的联合样本。

\begin{proposition}[Gibbs 更新保持联合分布]
\label{prop:lh-mcmc-gibbs-invariance}
设 $p(x_j,x_{-j})$ 是正规联合密度，相关满条件分布存在。固定其余坐标而从该满条件分布重新抽取第 $j$ 坐标的核，关于联合分布可逆；任意固定顺序的完整扫描均保持该联合分布不变。
\end{proposition}
\begin{proof}
令其余坐标的共同值为 $v$，更新前后的第 $j$ 坐标为 $a,b$。正向联合概率流的密度为
\[
p(v)p(a\mid v)p(b\mid v),
\]
交换 $a,b$ 后不变，故单坐标核满足细致平衡。由不变核复合仍不变，完整扫描也保持 $p$；这里没有断言复合后的核必可逆。
\end{proof}

接受率为 $1$ 并不保证链能够探索整个支撑。例如目标只在两个离散点 $(0,0)$、$(1,1)$ 上各有质量 $1/2$，给定一个坐标就完全确定另一个坐标，逐坐标 Gibbs 会永远停留在初始点。这是可约链；同时更新两个坐标的块抽样才可能跨越两点。在离散或退化条件分布中，Gibbs 也可能抽回原值，因此“无拒绝”不能解释为“每步必移动”。

\subsection{吉布斯抽样算法}
\label{subsec:lh-mcmc-gibbs-algorithm}

给定联合目标的满条件抽样器、函数 $f$、有效初值 $x^{(0)}$、预热轮数 $m$ 与保留轮数 $N$。对 $t=1,\ldots,m+N$，按顺序生成
\begin{equation}
\begin{aligned}
x_1^{(t)}&\sim p(x_1\mid x_2^{(t-1)},\ldots,x_d^{(t-1)}),\\
x_j^{(t)}&\sim p(x_j\mid x_1^{(t)},\ldots,x_{j-1}^{(t)},
x_{j+1}^{(t-1)},\ldots,x_d^{(t-1)}),\quad 1<j<d,\\
x_d^{(t)}&\sim p(x_d\mid x_1^{(t)},\ldots,x_{d-1}^{(t)}).
\end{aligned}
\label{eq:lh-mcmc-gibbs-sweep}
\end{equation}
每轮保留完整向量 $x^{(t)}$，最终输出后 $N$ 轮向量与函数均值。更新第 $j$ 个坐标时必须使用本轮已更新的前面坐标；把所有坐标同时按上一轮的条件值独立抽样，通常不保持目标联合分布。完整扫描的成本为各满条件生成成本之和。若某个满条件无法直接抽样，可以在该坐标内使用一次保持条件分布不变的 MH 更新，得到 Metropolis-within-Gibbs 方法；其不变性论证相同，仍需检查整体遍历条件。

\begin{example}
\label{ex:lh-mcmc-gaussian-gibbs}
设二元正态目标为
\[
\begin{pmatrix}X_1\\X_2\end{pmatrix}
\sim N\!\left(\begin{pmatrix}0\\0\end{pmatrix},
\begin{pmatrix}1&\rho\\\rho&1\end{pmatrix}\right),\qquad |\rho|<1.
\]
推导满条件分布与 Gibbs 更新，并分析相关参数如何影响计算。
\end{example}
\begin{solve}
联合密度的指数部分为
\[
-\frac{x_1^2-2\rho x_1x_2+x_2^2}{2(1-\rho^2)}
=-\frac{(x_1-\rho x_2)^2}{2(1-\rho^2)}-\frac{x_2^2}{2}.
\]
固定 $x_2$ 后第二项与 $x_1$ 无关，故 $X_1\mid X_2=x_2\sim N(\rho x_2,1-\rho^2)$；交换两坐标即得另一个满条件。因此使用每轮相互独立的 $\varepsilon_1^{(t)},\varepsilon_2^{(t)}\sim N(0,1)$，可写为
\begin{equation}
\begin{aligned}
x_1^{(t)}&=\rho x_2^{(t-1)}+\sqrt{1-\rho^2}\,\varepsilon_1^{(t)},\\
x_2^{(t)}&=\rho x_1^{(t)}+\sqrt{1-\rho^2}\,\varepsilon_2^{(t)}.
\end{aligned}
\label{eq:lh-mcmc-gaussian-gibbs}
\end{equation}
这里第二个条件均值使用新生成的 $x_1^{(t)}$。代入消去该量，得到
\[
x_2^{(t)}=\rho^2x_2^{(t-1)}
+\sqrt{1-\rho^2}\bigl(\rho\varepsilon_1^{(t)}+\varepsilon_2^{(t)}\bigr).
\]
创新项独立于过去，方差为 $1-\rho^4$；初值的影响按 $(\rho^2)^t$ 衰减。以目标分布初始化后，第二坐标的滞后 $k$ 自相关为 $\rho^{2k}$，故估计该坐标均值时的 $\tau_f=(1+\rho^2)/(1-\rho^2)$。当 $|\rho|$ 接近 $1$ 时，尽管每次都接受，链的均值估计仍可能需要很多轮。$|\rho|=1$ 时协方差矩阵奇异，所用普通二元密度与正条件方差均不再成立，不能把这一边界情况直接代入算法。
\end{solve}

\subsection{抽样计算}
\label{subsec:lh-mcmc-gibbs-computation}

满条件计算的重要技巧是只保留与待更新变量有关的联合因子。设 $D$ 为固定观测，$A$ 为随机超参数，$\Theta$ 为模型参数，$Z$ 为隐变量，模型分解为
\begin{equation}
p(a,\theta,z,D)=p(a)p(\theta\mid a)p(z,D\mid\theta).
\label{eq:lh-mcmc-hierarchical-model}
\end{equation}
若 $a_i$、$\theta_j$、$z_k$ 分别是各向量的一个分量，则
\begin{equation}
\begin{aligned}
p(a_i\mid a_{-i},\theta,z,D)
&\propto p(\theta\mid a)p(a),\\
p(\theta_j\mid\theta_{-j},a,z,D)
&\propto p(z,D\mid\theta)p(\theta\mid a),\\
p(z_k\mid z_{-k},a,\theta,D)
&\propto p(z,D\mid\theta).
\end{aligned}
\label{eq:lh-mcmc-hierarchical-conditionals}
\end{equation}
每行的比例都仅相对于左侧待抽变量而言，其他变量固定；右侧还可以继续删除与该分量无关的因子。例如更新 $a_i$ 时，$p(z,D\mid\theta)$ 为常数，因而不必反复计算。若模型进一步按图结构分解，只需保留含该变量的局部因子；这解释了 Gibbs 在稀疏概率模型中的计算优势。

上述化简免去的是与当前变量无关的常数，并不自动提供一个可用的抽样器。识别出正态、Beta 等标准条件分布后，可直接使用相应生成方法；仅知道“正比于某函数”，还需确认其正规性并构造合适的抽样步骤。各满条件也必须来自同一个一致的联合目标，不能随意拼凑。若坐标间强依赖导致探索缓慢，可考虑块更新或重新参数化；若支撑使逐坐标路径断开，则必须改变可达路径，而不能单靠增加拒绝后的留步解决。

本章的计算逻辑由三层组成：用函数平均表示待求积分，用不变转移核保证目标分布被保持，再用遍历条件与相关性分析解释平均的极限和有限运行的误差。MH 通过正反建议概率之比校正移动，Gibbs 则在满条件可直接抽样时使接受率等于 $1$。实际使用时，目标正规性、转移可达性、初始化影响与样本相关性都应在各自的位置核查，不能由一个细致平衡等式统一替代。

% 教材来源：机器学习/textbooks/机器学习（分章原文）/机器学习-18.pdf。
% 对应教材第 18 章，原书第 339--349 页；沿用本文件的自动章节编号。

\fi

\selectchapter{19}
\ifcompilechapter
\chapter{概率潜在语义分析}
\label{ch:lh-plsa}

一个文本可以同时谈论多个话题，同一个单词也可能出现在不同话题中。因此，只把文本或单词划分到互不相交的类别，往往不足以表达它们的语义联系。\emph{概率潜在语义分析}（probabilistic latent semantic analysis，PLSA），又称\emph{概率潜在语义索引}（probabilistic latent semantic indexing，PLSI），用不可直接观测的话题连接文本与单词：每个文本对应一个话题分布，每个话题对应一个单词分布，从可观测的单词频数反推这两类分布。这是一种无监督学习方法，话题数由建模者预先指定，话题的语义名称则需要根据学到的分布解释，不能把名称本身当作训练标签。

本章先建立生成模型及其等价的共现模型，再从参数、几何与矩阵分解三个角度解释模型，最后推导极大似然估计的期望极大算法（expectation-maximization algorithm，EM）。所需的条件概率、有限混合分布、拉格朗日乘子与对数函数凹性都将在具体推导中使用；全文采用自然对数，明确区分观测频数、模型概率与隐变量的后验概率。

\section{概率潜在语义分析模型}
\label{sec:lh-plsa-model}

\subsection{基本想法}
\label{subsec:lh-plsa-idea}

设词表为 $\mathcal W=\{w_1,\ldots,w_M\}$，文本指标集合为 $\mathcal D=\{d_1,\ldots,d_N\}$，预先给定的话题集合为 $\mathcal Z=\{z_1,\ldots,z_K\}$。这里 $M,N,K$ 都是正整数，$M$ 计算不同词的种数，而非语料中单词出现的总次数。用 $n_{ij}=n(w_i,d_j)\in\{0,1,2,\ldots\}$ 表示单词 $w_i$ 在文本 $d_j$ 中出现的次数，得到\emph{单词--文本矩阵}
\begin{equation}
X=(n_{ij})\in\mathbb R^{M\times N},\qquad
L_j=\sum_{i=1}^M n_{ij},\qquad
n=\sum_{j=1}^N L_j>0.
\label{eq:lh-plsa-counts}
\end{equation}
$L_j$ 为文本长度，$n$ 为语料中词项出现的总次数。以下删去空文本和全零词行，因而 $L_j>0$，且词表中的每个词至少出现一次；文本可以有不同长度。

模型只使用一个词出现在哪个文本以及出现多少次，不使用词序，这称为\emph{词袋假设}（bag-of-words assumption）。话题并非预先列出的关键词集合，而是整个词表上的概率分布；一个词可以在多个话题下具有正概率。相应地，文本也不必只属于一个话题，而可以具有不同的话题混合比例。这种表示实现了对文本和单词的软聚类（soft clustering），但“同属一个话题”表示模型给出的统计关联，不直接等同于词义相同。

为了区分随机变量与其取值，以下用 $D$、$W$、$Z$ 分别表示文本指标、单词和话题随机变量，其取值分别为 $d_j$、$w_i$、$z_k$。一次观测对应一个词项出现所给出的二元组 $(w_i,d_j)$，与之相关的话题 $z_k$ 不可直接观测。\emph[caution]{隐变量是每次词项出现所对应的话题，而不是整个文本唯一的话题。}同一文本内不同词项出现可以由不同话题产生。

\subsection{生成模型}
\label{subsec:lh-plsa-generative}

\begin{definition}[概率潜在语义分析的生成模型]
\label{def:lh-plsa-generative}
设
\[
\pi_j=P(D=d_j),\qquad
\theta_{kj}=P(Z=z_k\mid D=d_j),\qquad
\phi_{ik}=P(W=w_i\mid Z=z_k).
\]
记 $\pi=(\pi_1,\ldots,\pi_N)^{\mathsf T}$。各参数非负，并满足
\begin{equation}
\sum_{j=1}^N\pi_j=1,\qquad
\sum_{k=1}^K\theta_{kj}=1\quad(1\leq j\leq N),\qquad
\sum_{i=1}^M\phi_{ik}=1\quad(1\leq k\leq K).
\label{eq:lh-plsa-normalization}
\end{equation}
假设给定话题后单词与文本条件独立，即 $W\perp D\mid Z$，并规定
\begin{equation}
P(D=d_j,Z=z_k,W=w_i)=\pi_j\theta_{kj}\phi_{ik}.
\label{eq:lh-plsa-complete-joint}
\end{equation}
将话题边缘化，得到生成模型的观测概率
\begin{equation}
\begin{aligned}
p_{ij}&=P(W=w_i,D=d_j)
=\pi_j\sum_{k=1}^K\phi_{ik}\theta_{kj},\\
p_{i\mid j}&=P(W=w_i\mid D=d_j)
=\sum_{k=1}^K\phi_{ik}\theta_{kj},
\end{aligned}
\label{eq:lh-plsa-generative}
\end{equation}
其中条件概率式用于 $\pi_j>0$ 的文本。
\end{definition}

对于一次词项出现，生成过程依次为：按 $\pi$ 选取文本指标 $d_j$，按该文本的话题分布 $\theta_{\cdot j}$ 选取话题 $z_k$，再按该话题的词分布 $\phi_{\cdot k}$ 选取单词 $w_i$；观测时只保留文本与单词。图模型的依赖关系可写成 $D\to Z\to W$。条件独立假设具体意味着
\[
P(W=w_i\mid Z=z_k,D=d_j)=\phi_{ik},
\]
即给定话题后，同一个话题的词分布在所有文本中共享，而不是每个文本另有一套话题词分布。

也可以固定每个文本的长度 $L_j$，对该文本的每个词项出现独立抽取话题和单词。这时该文本的词频向量服从多项分布，单次抽取所用的是词表上的类别分布（categorical distribution）。将单次类别抽取与多次抽取得到的多项计数区分开，有助于理解概率参数和观测频数的关系。两种观察方式分别给出联合似然和给定文本指标后的条件似然；在忽略与参数无关的多项系数后，它们为
\begin{equation}
\mathcal L_{\mathrm{joint}}=\prod_{i=1}^M\prod_{j=1}^N
\left(\pi_j p_{i\mid j}\right)^{n_{ij}},\qquad
\mathcal L_{\mathrm{cond}}=\prod_{i=1}^M\prod_{j=1}^N
p_{i\mid j}^{n_{ij}}.
\label{eq:lh-plsa-likelihoods}
\end{equation}
若把有顺序的独立词项样本汇总为计数，还要乘以相应的多项系数；该系数不影响极大似然估计。这里的乘积分解依赖词项独立抽样的建模假设，不能仅由条件独立 $W\perp D\mid Z$ 单独推出。

\subsection{共现模型}
\label{subsec:lh-plsa-cooccurrence}

生成模型从文本出发选择话题，共现模型则先选择话题，再分别生成文本和单词。在后一种表示下，文本与单词的地位是对称的，话题解释了两者的共现关联。

\begin{definition}[概率潜在语义分析的共现模型]
\label{def:lh-plsa-cooccurrence}
设 $\rho_k=P(Z=z_k)$、$\psi_{jk}=P(D=d_j\mid Z=z_k)$，并保留 $\phi_{ik}=P(W=w_i\mid Z=z_k)$，记 $\rho=(\rho_1,\ldots,\rho_K)^{\mathsf T}$。各参数非负，满足
\[
\sum_{k=1}^K\rho_k=1,\qquad
\sum_{j=1}^N\psi_{jk}=1,\qquad
\sum_{i=1}^M\phi_{ik}=1.
\]
在 $W\perp D\mid Z$ 的假设下，共现概率为
\begin{equation}
p_{ij}=\sum_{k=1}^K\rho_k\phi_{ik}\psi_{jk}.
\label{eq:lh-plsa-cooccurrence}
\end{equation}
其生成过程为先按 $\rho$ 选取话题，再在给定话题后独立地按 $\psi_{\cdot k}$、$\phi_{\cdot k}$ 选取文本与单词，对应依赖关系 $D\leftarrow Z\to W$。
\end{definition}

\begin{proposition}[两种模型的等价性]
\label{prop:lh-plsa-equivalence}
生成模型与共现模型描述相同的一族单词--文本联合分布。对于正概率的条件事件，两者可由贝叶斯公式相互转换；零概率条件事件上的条件分布可以任意指定，不改变联合分布。
\end{proposition}
\begin{proof}
给定生成模型的参数，令
\begin{equation}
\rho_k=\sum_{j=1}^N\pi_j\theta_{kj},\qquad
\psi_{jk}=\frac{\pi_j\theta_{kj}}{\rho_k}\quad(\rho_k>0).
\label{eq:lh-plsa-generative-to-cooccurrence}
\end{equation}
由非负性与归一化条件，$\sum_k\rho_k=1$，且每个正权重话题满足 $\sum_j\psi_{jk}=1$。若 $\rho_k=0$，则所有 $\pi_j\theta_{kj}=0$，可任取一个归一化的 $\psi_{\cdot k}$。于是总有 $\rho_k\psi_{jk}=\pi_j\theta_{kj}$，代入式~\eqref{eq:lh-plsa-cooccurrence} 即得式~\eqref{eq:lh-plsa-generative}。

反过来，给定共现模型，令
\begin{equation}
\pi_j=\sum_{k=1}^K\rho_k\psi_{jk},\qquad
\theta_{kj}=\frac{\rho_k\psi_{jk}}{\pi_j}\quad(\pi_j>0).
\label{eq:lh-plsa-cooccurrence-to-generative}
\end{equation}
同理可验证归一化条件；若 $\pi_j=0$，任取归一化的 $\theta_{\cdot j}$，仍有 $\pi_j\theta_{kj}=\rho_k\psi_{jk}$。因此两种表示给出相同的联合分布。\qedhere
\end{proof}

这也完成了教材中证明两种模型等价的习题。等价的是可表示的概率分布，而不是各个参数数组逐项相同：$\theta_{kj}$ 描述给定文本后的话题概率，$\psi_{jk}$ 描述给定话题后的文本概率，不能交换两者的条件方向。生成模型便于解释文本的混合组成，共现模型便于解释文本与单词如何通过话题形成关联。

\subsection{模型性质}
\label{subsec:lh-plsa-properties}

直接为 $MN$ 个单词--文本对指定联合概率，需要 $MN-1$ 个自由参数。生成模型的参数化在概率单纯形内部需要
\begin{equation}
(N-1)+N(K-1)+K(M-1)=K(M+N-1)-1
\label{eq:lh-plsa-parameter-count}
\end{equation}
个自由坐标；共现模型相应为 $(K-1)+K(M-1)+K(N-1)$，结果相同。若文本边缘概率 $\pi$ 已固定，则生成模型只需后两组参数。当 $K$ 远小于 $M$ 和 $N$ 时，参数规模由 $O(MN)$ 降为 $O(K(M+N))$。这是参数化的坐标数，不表示不同参数一定对应不同的联合分布；例如，统一置换话题指标不会改变模型。较少的参数限制了模型容量，但不能据此断言不会过拟合。

几何上，用\emph{概率单纯形}（probability simplex）
\begin{equation}
\Delta^{M-1}=\left\{v\in\mathbb R^M:
v_i\geq0,\ \sum_{i=1}^M v_i=1\right\}
\label{eq:lh-plsa-simplex}
\end{equation}
表示词表上的全部概率分布。每个话题对应点 $\phi_{\cdot k}\in\Delta^{M-1}$，每个文本的模型词分布满足
\begin{equation}
p_{\cdot\mid j}=\sum_{k=1}^K\theta_{kj}\phi_{\cdot k}
\in\operatorname{conv}\{\phi_{\cdot1},\ldots,\phi_{\cdot K}\},
\label{eq:lh-plsa-convex-hull}
\end{equation}
其中 $\operatorname{conv}$ 表示凸包（convex hull）。非负且和为 $1$ 的系数使这里的线性组合成为凸组合，所以所有文本的模型分布都被限制在话题点的凸包中；观测相对频率 $n_{\cdot j}/L_j$ 则未必恰好落在其中，学习是在寻找合适的近似。$K=1$ 时凸包为一点；$K=2$ 且两个话题不同，凸包为线段。一般地，凸包维数至多为 $\min\{K-1,M-1\}$，只有 $K$ 个话题点仿射无关，即差向量 $\phi_{\cdot2}-\phi_{\cdot1},\ldots,\phi_{\cdot K}-\phi_{\cdot1}$ 线性无关时，才是 $(K-1)$ 维单纯形。此时混合系数是唯一的重心坐标；若存在仿射相关性，同一个词分布可能有不同的话题混合表示。教材所称的“话题单纯形”应结合这个条件理解。

矩阵形式使模型与潜在语义分析的联系更清楚。令
\[
\Phi=(\phi_{ik})\in\mathbb R^{M\times K},\quad
\Theta=(\theta_{kj})\in\mathbb R^{K\times N},\quad
\Psi=(\psi_{jk})\in\mathbb R^{N\times K},
\]
并令 $B=(p_{i\mid j})$、$A=(p_{ij})$，则
\begin{equation}
B=\Phi\Theta,\qquad
A=\Phi\Theta\operatorname{diag}(\pi)
=\Phi\operatorname{diag}(\rho)\Psi^{\mathsf T}.
\label{eq:lh-plsa-matrix-factorization}
\end{equation}
$\operatorname{diag}(v)$ 表示以向量 $v$ 为对角元的对角矩阵。$\Phi$、$\Theta$、$\Psi$ 都是非负且列归一化的矩阵，$\operatorname{diag}(\rho)$ 的非负对角元之和为 $1$，因此 $B$ 的各列是条件分布，$A$ 的全部元素之和为 $1$。二者的秩都不超过 $K$，但它们分别是条件概率矩阵与联合概率矩阵，不能直接与未归一化的频数矩阵 $X$ 混同。

\emph{潜在语义分析}（latent semantic analysis，LSA）通常对单词--文本矩阵作截断奇异值分解：取 $K\leq\min\{M,N\}$，保留最大的 $K$ 个奇异值及其左右奇异向量，得到 $X\approx U_K\Sigma_KV_K^{\mathsf T}$。其中左右奇异向量分别正交，但其分量可以为负。PLSA 则在非负、归一化约束下估计概率因子，不要求正交，学习目标来自似然。因此，两者虽然都通过低维因子描述文本与单词的关联，却不能把奇异向量直接当作话题概率，也不能把截断奇异值分解直接当作 PLSA 的参数估计。

\section{概率潜在语义分析的算法}
\label{sec:lh-plsa-algorithm}

\subsection{极大似然目标与隐变量}
\label{subsec:lh-plsa-objective}

根据式~\eqref{eq:lh-plsa-likelihoods}，联合对数似然在略去数据常数后分解为
\begin{equation}
\begin{aligned}
\ell_{\mathrm{joint}}(\pi,\Phi,\Theta)
&=\sum_{j=1}^N L_j\log\pi_j+\ell(\Phi,\Theta),\\
\ell(\Phi,\Theta)&=\sum_{i=1}^M\sum_{j=1}^N
n_{ij}\log\left(\sum_{k=1}^K\phi_{ik}\theta_{kj}\right).
\end{aligned}
\label{eq:lh-plsa-loglikelihood}
\end{equation}
以下对零频数项约定 $0\log p=0$；若 $n_{ij}>0$ 而模型概率为 $0$，则对数似然为 $-\infty$。学习是在式~\eqref{eq:lh-plsa-normalization} 的非负、归一化约束下极大化该目标。对 $\pi$ 的优化与其他两组参数分离，拉格朗日方程为 $L_j/\pi_j=\lambda$，求和得 $\lambda=n$，从而
\begin{equation}
\widehat\pi_j=\frac{L_j}{n}.
\label{eq:lh-plsa-document-marginal}
\end{equation}
由于各 $L_j>0$，目标 $\sum_jL_j\log\pi_j$ 在单纯形内部严格凹，并在边界趋于 $-\infty$，该解是唯一最大点。因此 $\widehat\pi_j$ 是随机词项出现来自文本 $d_j$ 的经验概率，只有文本等长时才为 $1/N$；后续只需优化 $\Phi$ 与 $\Theta$。

似然准则也可以从分布拟合理解。定义经验联合分布 $\widehat p_{ij}=n_{ij}/n$，则经验分布到模型分布的\emph{库尔贝克--莱布勒散度}（Kullback--Leibler divergence，KL divergence）为
\begin{equation}
D_{\mathrm{KL}}(\widehat p\|p)
=\sum_{i,j:\,\widehat p_{ij}>0}\widehat p_{ij}
\log\frac{\widehat p_{ij}}{p_{ij}}.
\label{eq:lh-plsa-kl}
\end{equation}
当经验正概率处的模型概率均为正时，由 $-\log x\geq1-x$ 可得该散度非负；若某处模型概率为零，则散度为 $+\infty$。由于
\[
\frac1n\ell_{\mathrm{joint}}
=\sum_{i,j}\widehat p_{ij}\log\widehat p_{ij}
-D_{\mathrm{KL}}(\widehat p\|p),
\]
极大似然估计等价于在 PLSA 可表示的分布族中最小化该散度。这个等价关系说明了概率拟合与平方误差矩阵近似的区别，而不提供整个问题的凸性或全局解保证。

直接对 $\ell$ 求极值会遇到“对数内含求和”的耦合。若知道每次词项出现来自哪个话题，令 $n_{ijk}$ 表示单词--文本对 $(w_i,d_j)$ 中来自话题 $z_k$ 的次数，则 $\sum_kn_{ijk}=n_{ij}$，完全数据的参数相关对数似然为
\begin{equation}
\ell_c=\sum_jL_j\log\pi_j+
\sum_{i,j,k}n_{ijk}(\log\phi_{ik}+\log\theta_{kj}).
\label{eq:lh-plsa-complete-loglikelihood}
\end{equation}
这时不同概率列的估计可以分开完成。EM 的思路是在当前参数下求出未知计数的条件期望，再以这些期望计数重新估计参数。隐藏的数据是话题取值或对应的计数 $n_{ijk}$，后验概率只是计算条件期望的工具，不能把一个概率分布本身说成观测数据或隐数据。

\subsection{E 步与 M 步的推导}
\label{subsec:lh-plsa-em-derivation}

记 $\eta=(\Phi,\Theta)$，第 $t$ 轮参数为 $\eta^{(t)}$。\emph{E 步}使用当前参数计算\emph{责任度}（responsibility），即每个已观测词项出现来自话题 $z_k$ 的后验概率：
\begin{equation}
r_{ijk}^{(t)}=P_{\eta^{(t)}}(Z=z_k\mid W=w_i,D=d_j)
=\frac{\phi_{ik}^{(t)}\theta_{kj}^{(t)}}
{\sum_{h=1}^K\phi_{ih}^{(t)}\theta_{hj}^{(t)}}.
\label{eq:lh-plsa-e-step}
\end{equation}
其中 $\pi_j$ 在分子分母中约去，责任度同时依赖单词和文本，通常并不等于文本先验话题比例 $\theta_{kj}$。只需对 $n_{ij}>0$ 的位置计算它；这些位置的分母必须为正。由给定观测后的条件分布，
\begin{equation}
\mathbb E_{\eta^{(t)}}[n_{ijk}\mid X]=n_{ij}r_{ijk}^{(t)}.
\label{eq:lh-plsa-expected-count}
\end{equation}
同一个单词--文本对可能出现多次，故这里必须乘以频数 $n_{ij}$，而不能把每个非零矩阵元都只计为一次样本。

将式~\eqref{eq:lh-plsa-complete-loglikelihood} 对隐数据取上述条件期望，并去掉与 $\eta$ 无关的项，得到
\begin{equation}
Q(\eta\mid\eta^{(t)})=
\sum_{i,j,k}n_{ij}r_{ijk}^{(t)}
(\log\phi_{ik}+\log\theta_{kj}).
\label{eq:lh-plsa-q-function}
\end{equation}
\emph{M 步}在非负和列归一化约束下极大化 $Q$。注意 $r^{(t)}$ 在此步骤中已经固定，求导时不能再将它视为候选参数 $\eta$ 的函数。

为简化推导，记当前轮的期望计数为
\begin{equation}
\begin{aligned}
a_{ik}^{(t)}&=\sum_{j=1}^N n_{ij}r_{ijk}^{(t)},\qquad
b_{kj}^{(t)}=\sum_{i=1}^M n_{ij}r_{ijk}^{(t)},\\
c_k^{(t)}&=\sum_{i=1}^M a_{ik}^{(t)}
=\sum_{j=1}^N b_{kj}^{(t)}.
\end{aligned}
\label{eq:lh-plsa-count-statistics}
\end{equation}
$a_{ik}^{(t)}$ 是话题 $k$ 产生单词 $i$ 的期望次数，$b_{kj}^{(t)}$ 是文本 $j$ 使用话题 $k$ 的期望次数，$c_k^{(t)}$ 是整个语料中分配给话题 $k$ 的期望总次数。于是
\[
Q=\sum_{i,k}a_{ik}^{(t)}\log\phi_{ik}
+\sum_{k,j}b_{kj}^{(t)}\log\theta_{kj}.
\]
引入列归一化约束的乘子 $\lambda_k$ 与 $\mu_j$，拉格朗日函数为
\[
\mathcal A=Q+
\sum_k\lambda_k\left(1-\sum_i\phi_{ik}\right)
+\sum_j\mu_j\left(1-\sum_k\theta_{kj}\right).
\]
若期望计数均为正，内点一阶条件给出
\[
\frac{a_{ik}^{(t)}}{\phi_{ik}}=\lambda_k,\qquad
\frac{b_{kj}^{(t)}}{\theta_{kj}}=\mu_j.
\]
将两式分别乘以 $\phi_{ik}$、$\theta_{kj}$ 并沿相应列求和，得到 $\lambda_k=c_k^{(t)}$ 与 $\mu_j=L_j$。因此更新式为
\begin{align}
\phi_{ik}^{(t+1)}
&=\frac{\sum_jn_{ij}r_{ijk}^{(t)}}
{\sum_{m=1}^M\sum_jn_{mj}r_{mjk}^{(t)}}
=\frac{a_{ik}^{(t)}}{c_k^{(t)}},
\label{eq:lh-plsa-m-step-phi}\\
\theta_{kj}^{(t+1)}
&=\frac{\sum_in_{ij}r_{ijk}^{(t)}}{L_j}
=\frac{b_{kj}^{(t)}}{L_j}.
\label{eq:lh-plsa-m-step-theta}
\end{align}
其中 $\sum_kb_{kj}^{(t)}=\sum_in_{ij}\sum_kr_{ijk}^{(t)}=L_j$，解释了第二个分母为何就是文本长度。$Q$ 是各概率列上凹函数的和，故上述解实现全局的 M 步最大化。若某些期望计数为零而对应总计数仍为正，归一化计数公式依然成立，零计数位置取零即可；这可由单列的 KL 散度非负性验证，无须在零点使用对数的内点求导。若 $c_k^{(t)}=0$，该轮 $Q$ 与 $\phi_{\cdot k}$ 无关，可保留原列，不能计算 $0/0$。

\begin{proposition}[EM 的似然单调性]
\label{prop:lh-plsa-em-monotonicity}
若当前观测对的模型概率为正，E 步精确计算式~\eqref{eq:lh-plsa-e-step}，M 步在约束下极大化式~\eqref{eq:lh-plsa-q-function}，则 $\ell(\eta^{(t+1)})\geq\ell(\eta^{(t)})$。由有限初值出发，对数似然值序列有有限极限，但这一结论不保证参数序列收敛，也不保证达到全局最大值。
\end{proposition}
\begin{proof}
对每个 $n_{ij}>0$ 的位置，任取话题分布 $q_{ij}=(q_{ijk})_{k=1}^K$，由对数函数的凹性，
\[
\log\sum_k\phi_{ik}\theta_{kj}
\geq\sum_kq_{ijk}\log\frac{\phi_{ik}\theta_{kj}}{q_{ijk}}.
\]
先在正概率情况下使用上式，零项再按极限约定处理。将右端乘以 $n_{ij}$ 并求和，记所得下界为 $F(\eta,q)$。当 $q_{ijk}=r_{ijk}^{(t)}$ 时，各正权重项中的比值在 $\eta=\eta^{(t)}$ 处相同，故下界取等号；固定 $q=r^{(t)}$ 后，$F$ 与 $Q$ 只差一个与 $\eta$ 无关的熵项。因此
\[
\ell(\eta^{(t+1)})\geq F(\eta^{(t+1)},r^{(t)})
\geq F(\eta^{(t)},r^{(t)})=\ell(\eta^{(t)}).
\]
又因每个 $p_{i\mid j}\leq1$，式~\eqref{eq:lh-plsa-loglikelihood} 中的条件对数似然不超过 $0$，故由单调有界性得到似然值的收敛。该论证只约束目标值，既没有证明参数的唯一性，也没有证明整个优化问题的凹性。
\end{proof}

分别固定 $\Phi$ 或 $\Theta$ 时，$\ell$ 是另一组参数的凹函数，但这不推出其关于两组参数联合凹。若 EM 参数序列收敛到所有参数均为正的内点，由更新映射的连续性，极限是 EM 不动点；将更新式代入似然的一阶导数可验证其满足归一化约束下的驻点条件。边界解需要单独检查相应的约束最优性条件，不能由似然值收敛直接断言是局部最大值。

\subsection{算法实施与手算例题}
\label{subsec:lh-plsa-em-implementation}

算法输入为非负整数矩阵 $X$、话题数 $K$、容差 $\varepsilon>0$ 和最大迭代次数 $T_{\max}$；输出为话题词分布 $\Phi$、文本话题分布 $\Theta$，以及需要联合模型时的 $\widehat\pi_j=L_j/n$。在移除空文本与全零词行后，可按以下步骤计算。
\begin{enumerate}
\item 为 $\Phi^{(0)}$ 和 $\Theta^{(0)}$ 选取严格正的初值，并分别按列归一化，计算初始条件对数似然。可以先随机产生正数再归一化；随机性来自初始化，随后的标准 EM 更新本身是确定的。
\item 对每个 $n_{ij}>0$ 的位置执行 E 步，累计式~\eqref{eq:lh-plsa-count-statistics} 的期望计数，再用式~\eqref{eq:lh-plsa-m-step-phi} 与式~\eqref{eq:lh-plsa-m-step-theta} 同时更新全部参数。当前轮的全部责任度均应使用旧参数，不能在累计过程中混用部分更新后的列。
\item 重新计算 $\ell^{(t+1)}=\ell(\Phi^{(t+1)},\Theta^{(t+1)})$。若增量非负且满足
\[
\frac{\ell^{(t+1)}-\ell^{(t)}}{\max\{1,|\ell^{(t)}|\}}
\leq\varepsilon,
\]
则返回当前参数；否则继续迭代。达到 $T_{\max}$ 时也返回当前近似解，并注明是否满足停止条件。若出现超出舍入误差的似然下降，应检查实现，不能将其当作正常收敛。
\end{enumerate}
停止条件检测的是似然增量，不能作为全局最优的证书。对不同初始化重复运行并比较最终似然，可以减少选到较差解的机会，但仍不能提供全局最优保证。若全部话题以完全相同的词分布和对称权重初始化，更新可能始终保持这种对称性，因此应避免用这样的初值期待算法自动区分话题。

令 $S=\#\{(i,j):n_{ij}>0\}$ 为非零计数位置的个数。一次 E 步及计数累计需要 $O(SK)$ 运算，归一化需要 $O(K(M+N))$；稠密情况下 $S=MN$。责任度可以逐个非零位置计算并立即累计，无须保存 $M\times N\times K$ 的数组，除存储输入外工作空间可为 $O(K(M+N))$。极小概率的计算可使用 $\log\phi_{ik}+\log\theta_{kj}$ 和 log-sum-exp 技巧，并检查每列的非负性与归一化。直接在更新后加任意常数并重新归一化会改变标准 EM 的 M 步，不能未经分析就沿用原来的单调性结论。

% 自拟例题：用精确分数演示一次 EM 更新，不是教材习题 18.3 的数值结果。
\begin{example}
\label{ex:lh-plsa-one-em-step}
设有两个单词、两个文本和两个话题，频数及初值为
\[
X=\begin{pmatrix}2&1\\1&2\end{pmatrix},\qquad
\Phi^{(0)}=\begin{pmatrix}3/4&1/4\\1/4&3/4\end{pmatrix},\qquad
\Theta^{(0)}=\begin{pmatrix}3/5&2/5\\2/5&3/5\end{pmatrix}.
\]
计算一轮 EM 更新，并检查词分布与文本话题分布的归一化。
\end{example}
\begin{solve}
两个文本长度都为 $3$，故 $\widehat\pi=(1/2,1/2)$。初始条件概率矩阵为
\[
B^{(0)}=\Phi^{(0)}\Theta^{(0)}
=\begin{pmatrix}11/20&9/20\\9/20&11/20\end{pmatrix}.
\]
只列出话题 $1$ 的责任度，话题 $2$ 的责任度逐项取补数：
\[
R_1^{(0)}=(r_{ij1}^{(0)})_{i,j}
=\begin{pmatrix}9/11&2/3\\1/3&2/11\end{pmatrix}.
\]
例如，$r_{111}^{(0)}=((3/4)(3/5))/(11/20)=9/11$，大于文本 $1$ 的先验话题比例 $3/5$，因为单词 $1$ 在话题 $1$ 中也更常见。将责任度分别乘以频数，话题 $1$ 的期望计数矩阵为
\[
\begin{pmatrix}18/11&2/3\\1/3&4/11\end{pmatrix}.
\]
它的行和为 $76/33$、$23/33$，列和为 $65/33$、$34/33$，总和为 $3$。归一化后得到
\[
\Phi^{(1)}=\begin{pmatrix}76/99&23/99\\23/99&76/99\end{pmatrix},\qquad
\Theta^{(1)}=\begin{pmatrix}65/99&34/99\\34/99&65/99\end{pmatrix}.
\]
两矩阵各列之和都为 $1$，而计数中的分数正是条件期望，并不要求取整数。新的 $B^{(1)}$ 对角元为 $5722/9801$，非对角元为 $4079/9801$，故条件对数似然由
\[
4\log(11/20)+2\log(9/20)\approx-3.98836
\]
增加至
\[
4\log(5722/9801)+2\log(4079/9801)\approx-3.90593.
\]
这验证了本轮更新的似然上升；它只是一次迭代的结果，尚不能视为收敛解。
\end{solve}

\begin{notealert}[type=caution,title={文本指标与新文本}]
基本 PLSA 为训练集合中的每个文本单独保存一列 $\theta_{\cdot j}$，并未给新文本的话题比例规定先验分布，也未直接给出生成未见文本指标的机制。若要表示一个新文本，可以固定学到的 $\Phi$，仅对新文本的混合比例 $\theta_{\cdot\mathrm{new}}$ 最大化其词项条件似然；责任度仍用式~\eqref{eq:lh-plsa-e-step}，再以新文本词频和长度更新该列。这是附加的参数拟合步骤，词表外单词也需预先处理，不能认为原模型已自动解决新文本推断。
\end{notealert}

\subsection{共现模型的 EM 算法}
\label{subsec:lh-plsa-cooccurrence-em}

教材还要求推导共现模型的 EM 算法。其联合对数似然为
\[
\ell_{\mathrm{co}}(\rho,\Phi,\Psi)
=\sum_{i,j}n_{ij}\log\left(\sum_k\rho_k\phi_{ik}\psi_{jk}\right).
\]
E 步由贝叶斯公式得到
\begin{equation}
\widetilde r_{ijk}^{(t)}=
\frac{\rho_k^{(t)}\phi_{ik}^{(t)}\psi_{jk}^{(t)}}
{\sum_{h=1}^K\rho_h^{(t)}\phi_{ih}^{(t)}\psi_{jh}^{(t)}}.
\label{eq:lh-plsa-cooccurrence-e-step}
\end{equation}
记 $\widetilde a_{ik}^{(t)}=\sum_jn_{ij}\widetilde r_{ijk}^{(t)}$、$\widetilde b_{kj}^{(t)}=\sum_in_{ij}\widetilde r_{ijk}^{(t)}$、$\widetilde c_k^{(t)}=\sum_{i,j}n_{ij}\widetilde r_{ijk}^{(t)}$。完全数据对数似然的条件期望为
\[
Q_{\mathrm{co}}=\sum_k\widetilde c_k^{(t)}\log\rho_k
+\sum_{i,k}\widetilde a_{ik}^{(t)}\log\phi_{ik}
+\sum_{k,j}\widetilde b_{kj}^{(t)}\log\psi_{jk}.
\]
三组参数分别在各自的概率单纯形上优化。由于 $\sum_k\widetilde c_k^{(t)}=n$，对三组归一化约束分别使用前述乘子法，得到 M 步
\begin{equation}
\rho_k^{(t+1)}=\frac{\widetilde c_k^{(t)}}{n},\qquad
\phi_{ik}^{(t+1)}=\frac{\widetilde a_{ik}^{(t)}}{\widetilde c_k^{(t)}},\qquad
\psi_{jk}^{(t+1)}=\frac{\widetilde b_{kj}^{(t)}}{\widetilde c_k^{(t)}}.
\label{eq:lh-plsa-cooccurrence-m-step}
\end{equation}
当 $\widetilde c_k^{(t)}=0$ 时，更新 $\rho_k^{(t+1)}=0$，而 $\phi_{\cdot k}$ 与 $\psi_{\cdot k}$ 可保留原来的归一化分布。其余停止条件、零概率处理与似然单调性的论证与生成模型相同。

若两种模型以命题~\ref{prop:lh-plsa-equivalence} 中相互对应的正参数初始化，式~\eqref{eq:lh-plsa-cooccurrence-e-step} 与式~\eqref{eq:lh-plsa-e-step} 给出相同责任度。共现模型完成一轮 M 步后，
\[
\pi_j^{(t+1)}=\sum_k\rho_k^{(t+1)}\psi_{jk}^{(t+1)}
=\frac{\sum_k\widetilde b_{kj}^{(t)}}{n}=\frac{L_j}{n},
\]
再由式~\eqref{eq:lh-plsa-cooccurrence-to-generative} 转换得到
\[
\theta_{kj}^{(t+1)}=
\frac{\rho_k^{(t+1)}\psi_{jk}^{(t+1)}}{\pi_j^{(t+1)}}
=\frac{\widetilde b_{kj}^{(t)}}{L_j}.
\]
这正是生成模型的更新式，说明两套算法的不同形式来自参数化的不同，而不是在学习两种不相容的概率规律。实际使用时，$\Phi$ 的高概率词可用于解释话题，$\Theta$ 可用于表示训练文本的混合组成；具体话题名称及话题数的选择仍需结合任务与拟合结果判断。


% 教材来源：机器学习/textbooks/机器学习（分章原文）/机器学习-20.pdf。
% 对应教材第 20 章，原书第 385--414 页；保留本文的自动章节编号。
% 变分推导补充核对：Blei, Ng, Jordan, Latent Dirichlet Allocation,
% JMLR 3 (2003), 993--1022, §§5.1、5.3--5.4 及附录 A。
% https://www.jmlr.org/papers/volume3/blei03a/blei03a.pdf

\fi

\selectchapter{20}
\ifcompilechapter
\chapter{潜在狄利克雷分配}
\label{ch:lh-lda}

一篇文档可以同时讨论多个话题，同一话题也可以出现在不同文档中。\emph{潜在狄利克雷分配}（latent Dirichlet allocation，LDA）用两层概率向量描述这种混合结构：每篇文档具有一个话题分布，每个话题具有一个词分布；每个词位先抽取话题，再由该话题抽取词。与只估计混合权重的模型相比，LDA 进一步为这些概率向量指定狄利克雷先验，使文档之间能够共享统计结构，并通过先验约束缓和数据稀疏带来的估计问题。先验提供的是建模与正则化机制，其效果仍取决于数据、超参数与推断方法，不能理解为自动消除过拟合。

本章依次讨论狄利克雷分布及其共轭性、LDA 的生成模型、收缩的吉布斯抽样和变分 EM。前一种推断方法对后验构造马尔可夫链，用样本近似后验期望；后一种在受限分布族中优化证据下界，用一个易处理的分布近似后验。二者都要区分模型本身、推断目标和有限计算得到的近似结果。以下均使用自然对数，文档长度视为已给定；词汇表大小、话题数与文档数分别记为 $V$、$K$、$M$，总词位数记为 $I=\sum_{m=1}^M N_m$。

\section{狄利克雷分布}
\label{sec:lh-lda-dirichlet}

\subsection{分布定义}
\label{subsec:lh-lda-dirichlet-definition}

\begin{definition}[类别分布与多项分布]
\label{def:lh-lda-categorical-multinomial}
设 $d\geq2$，概率向量 $p=(p_1,\ldots,p_d)$ 满足 $p_j\geq0$、$\sum_jp_j=1$。若随机变量 $Y\in\{1,\ldots,d\}$ 满足 $P(Y=j)=p_j$，则称其服从\emph{类别分布}（categorical distribution），记为 $Y\sim\operatorname{Cat}(p)$。独立重复 $n$ 次这种试验，令 $X_j$ 为结果 $j$ 出现的次数，则计数向量 $X=(X_1,\ldots,X_d)$ 服从\emph{多项分布}（multinomial distribution），记为 $X\sim\operatorname{Mult}(n,p)$，其概率质量函数为
\begin{equation}
P(X_1=x_1,\ldots,X_d=x_d)
=\frac{n!}{\prod_{j=1}^d x_j!}\prod_{j=1}^d p_j^{x_j},
\quad x_j\in\{0,1,\ldots,n\},\quad\sum_{j=1}^d x_j=n.
\label{eq:lh-lda-multinomial}
\end{equation}
其中零次幂按 $p_j^0=1$ 解释，包括 $p_j=0$ 的情形。
\end{definition}

类别变量记录一次试验的类别，多项变量记录多次试验的计数。$n=1$ 时，多项计数向量只有一个分量为 $1$，与类别变量的独热编码相对应。教材在 LDA 中有时将单次话题或单词抽样也称为多项抽样，以下用 $\operatorname{Cat}$ 明确这一点。若 $d=2$，只保留第一类的计数就得到二项分布；再取 $n=1$，便得到伯努利分布。

\begin{definition}[狄利克雷分布]
\label{def:lh-lda-dirichlet}
令 $a=(a_1,\ldots,a_d)$，各 $a_j>0$，并记 $a_0=\sum_{j=1}^d a_j$。定义单纯形
\[
\Delta^{d-1}=\left\{\vartheta\in\mathbb R^d:
\vartheta_j\geq0,\ \sum_{j=1}^d\vartheta_j=1\right\}.
\]
若随机概率向量 $\Theta=(\Theta_1,\ldots,\Theta_d)$ 在单纯形内部的密度为
\begin{equation}
p(\vartheta\mid a)=\frac1{B(a)}\prod_{j=1}^d\vartheta_j^{a_j-1},
\qquad
B(a)=\frac{\prod_{j=1}^d\Gamma(a_j)}{\Gamma(a_0)},
\label{eq:lh-lda-dirichlet-density}
\end{equation}
则称 $\Theta$ 服从\emph{狄利克雷分布}（Dirichlet distribution），记为 $\Theta\sim\operatorname{Dir}(a)$。这里 $B(a)$ 称为多元贝塔函数，$\Gamma(s)=\int_0^\infty t^{s-1}e^{-t}\,\mathrm dt$，$s>0$，满足 $\Gamma(s+1)=s\Gamma(s)$。
\end{definition}

由于最后一个坐标由其余坐标决定，式~\eqref{eq:lh-lda-dirichlet-density} 的密度是相对于前 $d-1$ 个坐标的勒贝格测度定义的：积分区域为 $\vartheta_j>0$、$\sum_{j<d}\vartheta_j<1$，并令 $\vartheta_d=1-\sum_{j<d}\vartheta_j$。以下简写 $\mathrm d\vartheta$ 都采用这一约定，而不是在 $\mathbb R^d$ 中直接积分。边界的概率为 $0$，但当某些 $a_j<1$ 时，密度可以在靠近边界处趋于无穷。

\begin{proposition}[归一化积分]
\label{prop:lh-lda-dirichlet-normalization}
对任意 $a_j>0$，有
\begin{equation}
\int_{\Delta^{d-1}}\prod_{j=1}^d\vartheta_j^{a_j-1}\,\mathrm d\vartheta
=B(a).
\label{eq:lh-lda-beta-integral}
\end{equation}
\end{proposition}
\begin{proof}
从 $\prod_j\Gamma(a_j)$ 的积分表示出发，令 $t_j=s\vartheta_j$，其中 $s=\sum_jt_j>0$、$\vartheta\in\Delta^{d-1}$。该变量替换的雅可比绝对值为 $s^{d-1}$，因而
\[
\prod_{j=1}^d\Gamma(a_j)
=\int_0^\infty e^{-s}s^{a_0-1}\,\mathrm ds
\int_{\Delta^{d-1}}\prod_{j=1}^d\vartheta_j^{a_j-1}\,\mathrm d\vartheta.
\]
被积函数非负，可由 Tonelli 定理交换积分次序；第一个积分等于 $\Gamma(a_0)$，相除即得结论。
\end{proof}

当 $d=2$，取 $\vartheta_1=t$、$\vartheta_2=1-t$，就得到\emph{贝塔分布}（beta distribution）$\operatorname{Beta}(a_1,a_2)$，密度为 $t^{a_1-1}(1-t)^{a_2-1}/B(a_1,a_2)$，$0<t<1$。因此“狄利克雷先验与多项观测”的二分类特例是“贝塔先验与二项观测”。先验描述的是未知概率向量的不确定性，类别分布或多项分布描述的是给定该向量后的观测，两种随机性位于不同层次。

\subsection{共轭先验与基本性质}
\label{subsec:lh-lda-dirichlet-conjugacy}

\begin{proposition}[矩与对数矩]
\label{prop:lh-lda-dirichlet-moments}
若 $\Theta\sim\operatorname{Dir}(a)$，则
\begin{align}
\mathbb E[\Theta_j]&=\frac{a_j}{a_0},&
\operatorname{Var}(\Theta_j)&=\frac{a_j(a_0-a_j)}{a_0^2(a_0+1)},
\label{eq:lh-lda-dirichlet-moments}\\
\operatorname{Cov}(\Theta_j,\Theta_\ell)&=-\frac{a_ja_\ell}{a_0^2(a_0+1)},\quad j\ne\ell,&
\mathbb E[\log\Theta_j]&=\psi(a_j)-\psi(a_0),
\label{eq:lh-lda-dirichlet-logmoment}
\end{align}
其中 $\psi(s)=\frac{\mathrm d}{\mathrm ds}\log\Gamma(s)$ 称为双伽马函数（digamma function）。
\end{proposition}
\begin{proof}
令 $e_j$ 为第 $j$ 个单位向量。将 $\vartheta_j$ 乘入密度后，利用式~\eqref{eq:lh-lda-beta-integral} 得
\[
\mathbb E[\Theta_j]=\frac{B(a+e_j)}{B(a)}=\frac{a_j}{a_0},\qquad
\mathbb E[\Theta_j^2]=\frac{B(a+2e_j)}{B(a)}
=\frac{a_j(a_j+1)}{a_0(a_0+1)}.
\]
同理，$j\ne\ell$ 时 $\mathbb E[\Theta_j\Theta_\ell]=B(a+e_j+e_\ell)/B(a)=a_ja_\ell/(a_0(a_0+1))$，减去均值乘积即可得到方差和协方差。对式~\eqref{eq:lh-lda-beta-integral} 关于 $a_j$ 求导再除以 $B(a)$，得到 $\mathbb E[\log\Theta_j]=\partial\log B(a)/\partial a_j$，即所述对数矩。交换微分与积分是合法的，因为在任意给定正参数附近，可用仍具有正指数参数的可积幂函数乘 $|\log\vartheta_j|$ 控制导数。
\end{proof}

把 $a$ 写成 $a_0\mu$，其中 $\mu_j>0$、$\sum_j\mu_j=1$，则 $\mu$ 决定均值，$a_0$ 称为浓度参数（concentration parameter）；固定均值时，浓度越大，方差越小。对称情形 $a_j=a_\ast$ 的均值均为 $1/d$，但随机抽出的向量通常不等于均匀向量。$a_\ast=1$ 时密度在单纯形上均匀；$a_\ast<1$ 时偏好靠近边界的稀疏组成；$a_\ast>1$ 时密度在中心取得最大值。这里“稀疏”表示若干分量接近 $0$，并不表示它们以正概率恰好等于 $0$。

取对数可将密度写成
\[
p(\vartheta\mid a)
=\exp\left\{\sum_j(a_j-1)\log\vartheta_j-\log B(a)\right\}.
\]
因此它属于指数分布族，自然参数为 $(a_j-1)_j$，充分统计量为 $(\log\vartheta_j)_j$。这也解释了变分推断中反复出现的 $\mathbb E[\log\Theta_j]$，它一般不等于 $\log\mathbb E[\Theta_j]$。

\begin{theorem}[狄利克雷与多项分布的共轭性]
\label{thm:lh-lda-conjugacy}
设 $\Theta\sim\operatorname{Dir}(a)$，给定 $\Theta=\vartheta$ 后，$Y_1,\ldots,Y_n$ 条件独立且均服从 $\operatorname{Cat}(\vartheta)$。观测序列 $y$ 的计数为 $c_j=\sum_{r=1}^n\mathbf1\{y_r=j\}$，则
\begin{equation}
\Theta\mid Y_1=y_1,\ldots,Y_n=y_n\sim\operatorname{Dir}(a+c).
\label{eq:lh-lda-conjugate-posterior}
\end{equation}
下一次同机制观测的后验预测概率为
\begin{equation}
P(Y_{n+1}=j\mid y,a)=\frac{c_j+a_j}{n+a_0}.
\label{eq:lh-lda-posterior-predictive}
\end{equation}
\end{theorem}
\begin{proof}
有序序列的似然为 $\prod_j\vartheta_j^{c_j}$，与先验相乘得到 $\prod_j\vartheta_j^{a_j+c_j-1}$，归一化即为 $\operatorname{Dir}(a+c)$。若只观测多项计数，似然还含 $n!/\prod_jc_j!$，但该因子不依赖 $\vartheta$，所以后验相同。预测概率为 $\int\vartheta_jp(\vartheta\mid y,a)\,\mathrm d\vartheta$，用命题~\ref{prop:lh-lda-dirichlet-moments} 的均值公式即可。
\end{proof}

先验与后验属于同一分布族，称为共轭；$a$ 是先验的参数，也是整个层次模型的\emph{超参数}（hyperparameter）。从后验均值看，$a_j$ 像预先加入的计数，故称先验伪计数（prior pseudo-count）。当 $n>0$ 时，式~\eqref{eq:lh-lda-posterior-predictive} 等于 $(n/(n+a_0))(c_j/n)+(a_0/(n+a_0))(a_j/a_0)$，明确显示了经验频率与先验均值的加权折中。伪计数不必为整数，也不能把后验均值与最大后验估计混用：只有各 $a_j+c_j>1$ 时，内部众数才为 $(a_j+c_j-1)/(a_0+n-d)$。

\begin{example}
\label{ex:lh-lda-beta-update}
一个二类别概率 $\Theta_1$ 的先验为 $\operatorname{Beta}(2,1)$，四次观测中第一类出现一次、第二类出现三次，求后验及下一次第一类的预测概率。
\end{example}
\begin{solve}
后验为 $\operatorname{Beta}(3,4)$，预测概率等于后验均值 $3/7$，也可写成 $(4/7)(1/4)+(3/7)(2/3)$。此时内部众数为 $(3-1)/(3+4-2)=2/5$，与均值不同。未出现某一类时，正超参数仍给出正预测概率，这正是先验平滑的作用。
\end{solve}

\section{潜在狄利克雷分配模型}
\label{sec:lh-lda-model}

\subsection{基本想法}
\label{subsec:lh-lda-model-idea}

LDA 是混合成员模型（mixed-membership model）：一个文档的不同词位可以由不同话题生成，而不是整篇文档只能归入一个类别。话题本身是词汇表上的概率分布，“体育”“经济”等名称是学习后依据高概率词给出的解释，并不是模型预先观察到的标签。同一词也可以在不同话题下具有不同概率，因此词与话题之间不存在固定的一一对应。

概率潜在语义分析（probabilistic latent semantic analysis，PLSA）同样用“文档的话题混合、话题的词混合”表示词概率，但通常为训练文档分别估计混合权重。LDA 为文档权重引入共享的先验，因而对新文档也有明确的推断机制；本章完整模型还为话题的词分布设置先验。不能仅将 PLSA 与 LDA 的差别概括为“一个没有先验、一个有先验”，因为对文档混合权重的生成假设及其推断方式也随之改变。

\subsection{模型定义与生成过程}
\label{subsec:lh-lda-model-definition}

令词汇表为 $\mathcal V=\{v_1,\ldots,v_V\}$，话题用整数 $1,\ldots,K$ 表示。第 $m$ 篇文档有 $N_m$ 个词位；随机变量 $W_{mn}\in\{1,\ldots,V\}$ 表示词的索引，$Z_{mn}\in\{1,\ldots,K\}$ 表示该词位的话题，观测值分别记为 $w_{mn}$、$z_{mn}$。记文档的话题概率向量为 $\Theta_m$，话题的词概率向量为 $\Phi_k$，其取值为 $\theta_m=(\theta_{m1},\ldots,\theta_{mK})$ 和 $\phi_k=(\phi_{k1},\ldots,\phi_{kV})$。全体取值组成矩阵 $\theta\in\mathbb R^{M\times K}$、$\phi\in\mathbb R^{K\times V}$，每行均为概率向量。

\begin{definition}[完整 LDA 模型]
\label{def:lh-lda-model}
给定 $K,V\geq2$、$M\geq1$、各 $N_m\geq1$，以及正超参数向量 $\alpha\in(0,\infty)^K$、$\beta\in(0,\infty)^V$，LDA 按下述过程生成语料：
\begin{enumerate}
\item 对每个话题 $k$，独立抽取 $\Phi_k\sim\operatorname{Dir}(\beta)$；这些词分布在所有文档间共享。
\item 对每篇文档 $m$，独立抽取 $\Theta_m\sim\operatorname{Dir}(\alpha)$，且这些变量与全体 $\Phi_k$ 先验独立。
\item 给定上述概率向量，对每个词位 $(m,n)$，先抽取 $Z_{mn}\mid\Theta_m\sim\operatorname{Cat}(\Theta_m)$，再抽取 $W_{mn}\mid Z_{mn}=k,\Phi\sim\operatorname{Cat}(\Phi_k)$；给定这些父变量后，各次抽样独立。
\end{enumerate}
\end{definition}

记 $\alpha_0=\sum_k\alpha_k$、$\beta_0=\sum_v\beta_v$。$\alpha_k/\alpha_0$ 是先验平均话题比例，$\alpha_0$ 控制文档话题组成围绕该均值的集中程度；$\beta_v/\beta_0$ 与 $\beta_0$ 对话题词分布起对应作用。取所有分量为 $1$ 得到单纯形上的均匀先验，不是强制所有文档、所有话题采用均匀分布。$K$ 是预先给定的模型规模，可以根据验证集表现或具体分析目标选择；该模型不会在一次推断中自动确定话题数。

模型把真实词序中的语法、相邻依赖等信息舍去，采用词袋（bag of words）假设。文档长度在本章作为条件给定；若另建长度模型，其参数和独立性假设需要单独说明，不能从上述生成过程推出长度分布。生成算法输入词汇表、$K$、长度及超参数，输出词序列和隐变量；面对实际语料时，词序列已知，任务转为反向推断隐变量，两者的输入输出不同。

\subsection{概率图模型}
\label{subsec:lh-lda-graphical-model}

用有向概率图描述时，主要依赖为 $\alpha\to\Theta_m\to Z_{mn}\to W_{mn}$ 和 $\beta\to\Phi_k\to W_{mn}$，最后一条表示 $Z_{mn}$ 选择相应话题的词分布。板块表示法（plate notation）用重复区域分别表示 $K$ 个话题、$M$ 篇文档和第 $m$ 篇的 $N_m$ 个词位；被重复的是不同随机变量，它们共享分布形式或父变量，并不是同一个随机变量被多次复制。

可观测量只有 $W$，而 $Z$、$\Theta$、$\Phi$ 都是隐变量。尽管习惯上也称 $\theta$、$\phi$ 为“参数”，在完整贝叶斯模型中它们是带有先验和后验的随机概率向量；$\alpha$、$\beta$ 是本章暂时固定的超参数。先验中的独立性也不代表后验独立性：观察一个词会同时影响它的话题归属和相关概率向量，由此产生耦合，这正是精确推断困难的来源。

\subsection{随机变量序列的可交换性}
\label{subsec:lh-lda-exchangeability}

\begin{definition}[可交换序列]
\label{def:lh-lda-exchangeability}
有限随机变量序列 $X_1,\ldots,X_N$ 若对任意排列 $\pi$ 都满足 $(X_1,\ldots,X_N)\overset{\mathrm d}=(X_{\pi(1)},\ldots,X_{\pi(N)})$，则称其可交换（exchangeable）。无限序列的任意有限子序列都可交换时，称其无限可交换（infinitely exchangeable）。
\end{definition}

独立同分布蕴含可交换，反向一般不成立。对有限类别空间，de Finetti 表示定理说明：一个无限可交换序列的有限维联合分布可表示为类别独立同分布的混合，即存在概率向量上的概率测度 $\mu$，使
\begin{equation}
P(X_1=x_1,\ldots,X_N=x_N)
=\int_{\Delta^{d-1}}\prod_{n=1}^N\vartheta_{x_n}\,\mu(\mathrm d\vartheta).
\label{eq:lh-lda-definetti}
\end{equation}
这一定理需要无限可交换性；只知道某个固定长度的序列可交换，不能直接得到同样结论。定理也不规定混合测度必须是狄利克雷分布，LDA 对先验族的选择是额外的建模假设。

LDA 中，同一文档的 $Z_{mn}$ 给定 $\Theta_m$ 后独立同分布，积分消去 $\Theta_m$ 后仍可交换，却通常不独立。例如对 $n\ne r$，由条件独立性与全协方差公式可得
\begin{equation}
\operatorname{Cov}\bigl(\mathbf1\{Z_{mn}=k\},\mathbf1\{Z_{mr}=k\}\bigr)
=\operatorname{Var}(\Theta_{mk})>0.
\label{eq:lh-lda-topic-dependence}
\end{equation}
同样，给定 $\theta_m,\phi$ 后词位独立同分布，边缘化后文档内词序列可交换。词序无关来自模型假设，并不表示自然语言中的词序本来无关。

\subsection{概率公式}
\label{subsec:lh-lda-probability-formulas}

对全体隐变量取值 $z,\theta,\phi$ 与观测 $w$，联合密度或质量函数为
\begin{equation}
p(w,z,\theta,\phi\mid\alpha,\beta)
=\prod_{k=1}^K\operatorname{Dir}(\phi_k\mid\beta)
\prod_{m=1}^M\left[
\operatorname{Dir}(\theta_m\mid\alpha)
\prod_{n=1}^{N_m}\theta_{m,z_{mn}}\phi_{z_{mn},w_{mn}}
\right].
\label{eq:lh-lda-full-joint}
\end{equation}
这里 $\operatorname{Dir}(\cdot\mid\cdot)$ 表示式~\eqref{eq:lh-lda-dirichlet-density} 的密度。求和消去每个词位的话题，得到
\begin{equation}
p(w_m\mid\theta_m,\phi)
=\prod_{n=1}^{N_m}\left(\sum_{k=1}^K\theta_{mk}\phi_{k,w_{mn}}\right).
\label{eq:lh-lda-document-conditional}
\end{equation}
再积分消去文档权重，可定义给定共享词分布时的文档概率
\begin{equation}
p(w_m\mid\alpha,\phi)
=\int_{\Delta^{K-1}}\operatorname{Dir}(\theta_m\mid\alpha)
\prod_{n=1}^{N_m}\left(\sum_{k=1}^K\theta_{mk}\phi_{k,w_{mn}}\right)\,\mathrm d\theta_m.
\label{eq:lh-lda-document-marginal}
\end{equation}
完整模型的语料概率为
\begin{equation}
p(w\mid\alpha,\beta)
=\int\left[\prod_{k=1}^K\operatorname{Dir}(\phi_k\mid\beta)\right]
\prod_{m=1}^Mp(w_m\mid\alpha,\phi)\,\mathrm d\phi.
\label{eq:lh-lda-corpus-marginal}
\end{equation}
所有文档共享同一组 $\Phi_k$，所以给定 $\Phi$ 时文档独立，积分消去 $\Phi$ 后一般不独立；不能将式~\eqref{eq:lh-lda-corpus-marginal} 改成各文档完全边缘概率的乘积。

推断目标为
\[
p(z,\theta,\phi\mid w,\alpha,\beta)
=\frac{p(w,z,\theta,\phi\mid\alpha,\beta)}{p(w\mid\alpha,\beta)}.
\]
分母涉及积分和话题指派的组合求和，直接枚举 $I$ 个词位的话题共有 $K^I$ 种可能。共轭性可以帮助消去连续隐变量或计算局部更新，但不会自动消除所有词位话题之间的耦合，因而仍需要下面的近似推断方法。

\section{LDA 的吉布斯抽样算法}
\label{sec:lh-lda-gibbs}

\subsection{基本想法}
\label{subsec:lh-lda-gibbs-idea}

\emph{收缩的吉布斯抽样}（collapsed Gibbs sampling，也称折叠吉布斯抽样）先对式~\eqref{eq:lh-lda-full-joint} 中的 $\theta,\phi$ 积分，得到 $p(w,z\mid\alpha,\beta)$，再固定观测 $w$，以 $p(z\mid w,\alpha,\beta)$ 为目标分布轮流更新各个话题指派。此处“收缩”指解析边缘化连续变量，保留的抽样对象是 $z$；不是生成新的观测词，也不是直接抽样条件分布 $p(w\mid z,\alpha,\beta)$。

为整理积分，定义文档与话题计数、话题与词计数及其总数：
\begin{align}
c_{mk}&=\sum_{n=1}^{N_m}\mathbf1\{z_{mn}=k\},&
d_{kv}&=\sum_{m=1}^M\sum_{n=1}^{N_m}\mathbf1\{z_{mn}=k,\ w_{mn}=v\},\notag\\
c_{m\cdot}&=\sum_kc_{mk}=N_m,&
d_{k\cdot}&=\sum_vd_{kv}.
\label{eq:lh-lda-counts}
\end{align}
记 $c_m=(c_{m1},\ldots,c_{mK})$、$d_k=(d_{k1},\ldots,d_{kV})$。计数由当前话题指派和固定词序列确定；$\sum_kd_{k\cdot}=I$，且 $d_{k\cdot}=\sum_mc_{mk}$，这些恒等式也可用于检查实现。

\subsection{算法的主要部分}
\label{subsec:lh-lda-gibbs-main}

\begin{proposition}[收缩联合分布]
\label{prop:lh-lda-collapsed-joint}
对有序词序列和话题序列，有
\begin{equation}
p(w,z\mid\alpha,\beta)
=\prod_{m=1}^M\frac{B(\alpha+c_m)}{B(\alpha)}
\prod_{k=1}^K\frac{B(\beta+d_k)}{B(\beta)}.
\label{eq:lh-lda-collapsed-joint}
\end{equation}
因此 $p(z\mid w,\alpha,\beta)$ 正比于式~\eqref{eq:lh-lda-collapsed-joint}。
\end{proposition}
\begin{proof}
式~\eqref{eq:lh-lda-full-joint} 中涉及第 $m$ 个文档权重的部分为 $B(\alpha)^{-1}\prod_k\theta_{mk}^{\alpha_k+c_{mk}-1}$，积分得到 $B(\alpha+c_m)/B(\alpha)$。涉及第 $k$ 个词分布的部分同理得到 $B(\beta+d_k)/B(\beta)$，将这些因子相乘即可。这里描述的是具体有序序列的概率，不能再乘入多项计数的排列系数。
\end{proof}

\begin{proposition}[单词位的满条件分布]
\label{prop:lh-lda-gibbs-conditional}
令 $i=(m,n)$ 为当前词位，$w_i=v$，$z_{-i}$ 表示其余词位的话题指派。上标 $-i$ 表示所有计数都先排除当前词位，则
\begin{equation}
P(Z_i=k\mid z_{-i},w,\alpha,\beta)
=\frac{r_k}{\sum_{\ell=1}^Kr_\ell},\qquad
r_k=\frac{c_{mk}^{-i}+\alpha_k}{N_m-1+\alpha_0}
\frac{d_{kv}^{-i}+\beta_v}{d_{k\cdot}^{-i}+\beta_0}.
\label{eq:lh-lda-gibbs-conditional}
\end{equation}
第一个分母与候选话题 $k$ 无关，可在计算正比权重时省略；第二个分母通常随 $k$ 改变，必须保留。
\end{proposition}
\begin{proof}
将当前词位设为候选话题 $k$，相对于删除该词位后的计数，只有 $c_{mk}$ 和 $d_{kv}$ 增加 $1$。利用 $\Gamma(t+1)=t\Gamma(t)$ 得
\[
\frac{B(\alpha+c_m^{-i}+e_k)}{B(\alpha+c_m^{-i})}
=\frac{c_{mk}^{-i}+\alpha_k}{N_m-1+\alpha_0},\qquad
\frac{B(\beta+d_k^{-i}+e_v)}{B(\beta+d_k^{-i})}
=\frac{d_{kv}^{-i}+\beta_v}{d_{k\cdot}^{-i}+\beta_0}.
\]
其他因子对所有候选值相同，由式~\eqref{eq:lh-lda-collapsed-joint} 相消，最后对 $k$ 归一化即得。
\end{proof}

两个因子分别是排除当前词位后的“该文档倾向选哪个话题”和“该话题倾向生成哪个词”的后验预测概率。\emph[caution]{更新前必须先减去当前词位的旧计数。}否则旧指派会参与自身的条件概率计算，所得更新不再是上述目标分布的吉布斯更新。抽样也不能换成直接取最大权重：后者是另一种确定性更新，不再产生同一个后验分布的样本。

\begin{example}
\label{ex:lh-lda-gibbs-step}
取 $K=2$、$V=3$、$\alpha=(1,1)$、$\beta=(1,1,1)$。某词位观测为第一个词，删除该词位后，其文档话题计数为 $(2,1)$，两个话题的词计数分别为 $(3,1,0)$ 和 $(0,2,2)$，求该词位的新话题分布。
\end{example}
\begin{solve}
删除后文档长度为 $3$，两个话题总计数均为 $4$。由式~\eqref{eq:lh-lda-gibbs-conditional}，$r_1=(3/5)(4/7)=12/35$，$r_2=(2/5)(1/7)=2/35$，归一化后为 $(6/7,1/7)$。第一个话题既更符合当前文档的组成，也更常生成当前词；算法仍以 $1/7$ 的概率选择第二个话题，而不是必然选择第一个。
\end{solve}

\subsection{算法的后处理}
\label{subsec:lh-lda-gibbs-postprocessing}

给定一次完整指派 $z$，共轭性给出
\begin{equation}
\Theta_m\mid z,w,\alpha,\beta\sim\operatorname{Dir}(\alpha+c_m),\qquad
\Phi_k\mid z,w,\alpha,\beta\sim\operatorname{Dir}(\beta+d_k).
\label{eq:lh-lda-conditional-posteriors}
\end{equation}
在这个条件下，各文档权重与各话题词分布相互独立，条件后验均值为
\begin{equation}
\widehat\theta_{mk}(z)=\frac{c_{mk}+\alpha_k}{N_m+\alpha_0},\qquad
\widehat\phi_{kv}(z)=\frac{d_{kv}+\beta_v}{d_{k\cdot}+\beta_0}.
\label{eq:lh-lda-gibbs-posterior-means}
\end{equation}
这些是给定该次指派后的条件均值，不是完整后验已经精确求出的证明。若保留了 $S$ 次指派样本 $z^{(1)},\ldots,z^{(S)}$，则由全期望公式可用
\begin{equation}
\overline\theta_{mk}=\frac1S\sum_{s=1}^S\widehat\theta_{mk}(z^{(s)}),\qquad
\overline\phi_{kv}=\frac1S\sum_{s=1}^S\widehat\phi_{kv}(z^{(s)})
\label{eq:lh-lda-rao-blackwell}
\end{equation}
近似完整后验均值。特别地，$\phi$ 的分母随样本变化，一般不能用“平均计数之比”代替“各次比值的平均”。需要完整联合后验样本时，还可以在抽到 $z$ 后按式~\eqref{eq:lh-lda-conditional-posteriors} 抽取 $\Theta,\Phi$。

对称话题先验下，交换话题编号不会改变模型，因而存在标签交换（label switching）。若样本或多条链发生整体标签置换，直接平均同一编号的词分布可能掩盖话题结构；解释具体话题时需要说明标签对齐方式，或报告不依赖话题编号的量。后验概率、后验均值、某次硬指派以及高概率词展示分别回答不同问题，不宜统称为唯一的“真实话题”。

\subsection{算法与计算代价}
\label{subsec:lh-lda-gibbs-algorithm}

输入语料 $w$、话题数 $K$、正超参数 $\alpha,\beta$、预热轮数 $B_0\geq0$、保留间隔 $L\geq1$ 和样本数 $S\geq1$；一轮指完整扫描所有词位。输出话题指派样本及式~\eqref{eq:lh-lda-rao-blackwell} 的后验均值近似。实现时维护 $M\times K$ 的计数矩阵 $C=(c_{mk})$、$K\times V$ 的矩阵 $D=(d_{kv})$、各话题总数和当前指派，过程如下。
\begin{enumerate}
\item 为每个词位独立均匀抽取初始话题，累加 $C,D$ 与话题总数；检查式~\eqref{eq:lh-lda-counts} 的计数恒等式。
\item 每轮按固定顺序扫描词位 $i=(m,n)$。先从旧话题的 $c_{mk}$、$d_{k,w_i}$、$d_{k\cdot}$ 中各减 $1$；再用排除后的计数计算式~\eqref{eq:lh-lda-gibbs-conditional}，抽取新话题并写回；最后给新话题对应的三个计数各加 $1$。
\item 完成 $B_0$ 轮预热后继续迭代，每隔 $L$ 轮保留一次完整指派或所需的条件均值，累计 $S$ 次后停止，并报告所用轮数及诊断结果。
\end{enumerate}

正超参数使每个词位选取任一话题的条件概率均为正。对于固定有限语料，完整扫描所构成的马尔可夫链可以通过逐位重选到达任意指派，也能保持原指派不变，故在有限状态空间上不可约、非周期，并以 $p(z\mid w,\alpha,\beta)$ 为唯一平稳分布；当迭代轮数趋于无穷时，遍历平均满足相应的大数定律。这是渐近结论：有限预热并不保证已经到达平稳分布，相邻样本也通常相关。多链比较、后验统计量轨迹和有效样本量可帮助诊断；增加保留间隔本身不会保证独立性。

朴素实现每个词位计算 $K$ 个权重，若维护话题总数使分母可在常数时间取得，一轮耗时 $O(IK)$，$T$ 轮抽样耗时 $O(TIK)$。稠密计数与指派占用 $O(MK+KV+I)$ 空间，存储全部 $S$ 次指派还需 $O(SI)$。若每次保留时显式计算所有条件均值，另需 $O(S(MK+KV))$ 时间；这些后处理和存储费用应与抽样费用分别计算。稀疏或其他加速实现可能改变实际代价，本章不将它们的复杂度归于上述朴素算法。

\section{LDA 的变分 EM 算法}
\label{sec:lh-lda-variational-em}

\subsection{变分推理}
\label{subsec:lh-lda-variational-inference}

\emph{变分推理}（variational inference）把后验近似转化为优化问题。设 $x$ 是固定观测，$h$ 为全部隐变量，$\xi$ 为待固定或估计的模型参数，联合分布为 $p(x,h\mid\xi)$。选择易计算的分布族 $\mathcal Q$，在其中寻找使 $D_{\mathrm{KL}}(q(h)\Vert p(h\mid x,\xi))$ 最小的 $q$。这里的 KL 散度有方向，既不是距离，也不能在不改变目标的情况下交换两个分布的位置。

\begin{proposition}[证据下界恒等式]
\label{prop:lh-lda-elbo-identity}
假设 $0<p(x\mid\xi)<\infty$，$q$ 在后验为零的集合上也为零，且以下期望有限。定义\emph{证据下界}（evidence lower bound，ELBO）
\begin{equation}
\mathcal L(q,\xi)=\mathbb E_q[\log p(x,h\mid\xi)]-\mathbb E_q[\log q(h)].
\label{eq:lh-lda-elbo-definition}
\end{equation}
则
\begin{equation}
\log p(x\mid\xi)=\mathcal L(q,\xi)
+D_{\mathrm{KL}}\bigl(q(h)\Vert p(h\mid x,\xi)\bigr)
\geq\mathcal L(q,\xi).
\label{eq:lh-lda-elbo-identity}
\end{equation}
等号成立当且仅当 $q$ 与后验几乎处处相同。离散变量的积分解释为求和，混合变量同时进行积分与求和。
\end{proposition}
\begin{proof}
将 $\log p(h\mid x,\xi)=\log p(x,h\mid\xi)-\log p(x\mid\xi)$ 代入 KL 散度定义，利用 $\int q(h)\,\mathrm dh=1$ 即得恒等式。KL 非负可由 $-\log t\geq1-t$ 积分得到；等号要求两个分布几乎处处一致。
\end{proof}

固定 $\xi$ 时，$\log p(x\mid\xi)$ 与 $q$ 无关，所以最小化上述 KL 散度等价于最大化 ELBO。观测数据的边缘概率称为证据（evidence），其对数为对数证据；下界约束的是对数证据。优化并未使真实后验变得简单，而是在所选分布族内寻找近似，近似偏差不会仅因增加迭代次数而消失。

常用的\emph{平均场}（mean-field）分解为 $q(h)=\prod_jq_j(h_j)$。固定其他因子后，将 ELBO 中依赖 $q_j$ 的部分写成
\[
\int q_j(h_j)\,\mathbb E_{q_{-j}}[\log p(x,h\mid\xi)]\,\mathrm dh_j
-\int q_j(h_j)\log q_j(h_j)\,\mathrm dh_j.
\]
若指数化后的函数可归一化，它等于某个常数减去 $q_j$ 到该归一化分布的 KL 散度。因此坐标最优解满足
\begin{equation}
\log q_j^*(h_j)=\mathbb E_{q_{-j}}[\log p(x,h\mid\xi)]+\text{与 $h_j$ 无关的常数}.
\label{eq:lh-lda-mean-field-coordinate}
\end{equation}
轮流更新各因子便得到坐标上升变分推理。独立性是对近似分布 $q$ 的限制；它并不声称给定观测后的真实隐变量相互独立，且平均场中的一个因子仍可包含多个相依坐标。

\subsection{变分 EM 算法与收敛含义}
\label{subsec:lh-lda-variational-em-principle}

模型参数也未知时，变分 EM 交替改进同一个 $\mathcal L(q,\xi)$：E 步固定 $\xi$，优化受限族内的 $q$；M 步固定 $q$，优化模型参数 $\xi$。不必要求每一步都求出全局最大值，但若要保证下界单调，所用更新必须使当前下界不减。精确 EM 是 E 步能够取到真实后验的特殊情形。

\begin{proposition}[下界的单调性]
\label{prop:lh-lda-vem-monotonicity}
若每次 E、M 更新都不降低下界，则
\begin{equation}
\mathcal L(q^{(t)},\xi^{(t)})
\leq\mathcal L(q^{(t+1)},\xi^{(t)})
\leq\mathcal L(q^{(t+1)},\xi^{(t+1)}).
\label{eq:lh-lda-vem-monotonicity}
\end{equation}
若这些下界值还有统一有限上界，则它们收敛。仅由该结论不能推出参数序列收敛、收敛到全局最优点，或真实边缘似然每轮不减。
\end{proposition}
\begin{proof}
第一个不等式来自 E 步，第二个来自 M 步；单调有上界的实数序列收敛。由式~\eqref{eq:lh-lda-elbo-identity}，真实对数似然还包含随迭代变化的 KL 间隙，因此下界的增加不能控制两次真实似然的差。
\end{proof}

\begin{notealert}[type=caution,title={不能把精确 EM 的等号直接用于变分 EM}]
教材式 (20.41) 用 $\log p(x\mid\xi^{(t)})=\mathcal L(q^{(t+1)},\xi^{(t)})$ 推出似然单调性。该等号要求 E 步得到真实后验；一般受限平均场分布只给出不等式，因此应使用式~\eqref{eq:lh-lda-vem-monotonicity}。目标值收敛也不等于每个迭代变量都收敛，局部驻点性质还需要相应的正则性与更新条件。
\end{notealert}

下面的简化 LDA 直接建模离散语料的概率，故 $p(w\mid\xi)\leq1$，其 ELBO 不超过 $0$；从有限初始下界出发，满足单调更新条件即可保证下界值收敛。达到最大迭代次数、下界相对增量很小和已经找到全局最优解是不同的停止描述，输出时应如实区分。

\subsection{算法推导}
\label{subsec:lh-lda-variational-derivation}

本节先沿用教材的简化模型：将话题词矩阵 $\phi$ 作为待估计的确定性参数，不为它设置 $\operatorname{Dir}(\beta)$ 先验；只保留文档权重的狄利克雷先验。因此模型参数为 $\xi=(\alpha,\phi)$，$\beta$ 不参与本节这部分更新，不能直接把完整模型的条件均值当作简化模型的 M 步。为使初始下界有限，可将 $\phi$ 各行初始化为严格正的概率向量。

暂略文档下标，令一篇文档为 $w=(w_1,\ldots,w_N)$，则
\begin{equation}
p(\theta,z,w\mid\alpha,\phi)
=\operatorname{Dir}(\theta\mid\alpha)\prod_{n=1}^N\theta_{z_n}\phi_{z_n,w_n}.
\label{eq:lh-lda-simplified-joint}
\end{equation}
选择近似分布
\begin{equation}
q(\theta,z\mid\gamma,\eta)
=\operatorname{Dir}(\theta\mid\gamma)
\prod_{n=1}^N\operatorname{Cat}(z_n\mid\eta_n),
\quad\gamma_k>0,\quad\eta_{nk}\geq0,\quad\sum_k\eta_{nk}=1.
\label{eq:lh-lda-variational-family}
\end{equation}
其中 $\gamma\in(0,\infty)^K$、$\eta\in\mathbb R^{N\times K}$ 是变分参数；$\eta_{nk}$ 是近似分布赋予词位 $n$ 的话题 $k$ 的概率。平均场把整向量 $\Theta$ 视为一个因子，并没有将其受和为 $1$ 约束的各坐标分成独立因子。记 $\gamma_0=\sum_k\gamma_k$、$h_k=\psi(\gamma_k)-\psi(\gamma_0)=\mathbb E_q[\log\Theta_k]$。

\emph[tip]{先展开联合分布与近似分布的对数，再分别求期望。}这样 ELBO 自然分成先验、话题指派、词生成、狄利克雷因子熵与类别因子熵五项：
\begin{equation}
\begin{aligned}
\mathcal L(\gamma,\eta;\alpha,\phi)
={}&\log\Gamma(\alpha_0)-\sum_k\log\Gamma(\alpha_k)
+\sum_k(\alpha_k-1)h_k\\
&+\sum_{n,k}\eta_{nk}h_k
+\sum_{n,k}\eta_{nk}\log\phi_{k,w_n}\\
&-\log\Gamma(\gamma_0)+\sum_k\log\Gamma(\gamma_k)
-\sum_k(\gamma_k-1)h_k\\
&-\sum_{n,k}\eta_{nk}\log\eta_{nk}.
\end{aligned}
\label{eq:lh-lda-expanded-elbo}
\end{equation}
类别熵使用 $0\log0=0$ 的连续延拓；若 $\phi_{k,w_n}=0$，则该位置的 $\eta_{nk}$ 必须为 $0$，否则下界为负无穷。多文档总下界为各文档式~\eqref{eq:lh-lda-expanded-elbo} 的和，每篇具有自己的 $\gamma_m,\eta_m$，共享 $\alpha,\phi$。

\begin{proposition}[E 步的坐标更新]
\label{prop:lh-lda-variational-e-step}
固定模型参数，在式~\eqref{eq:lh-lda-variational-family} 内交替优化各因子，可取
\begin{align}
\eta_{nk}&=\frac{\phi_{k,w_n}\exp\{\psi(\gamma_k)-\psi(\gamma_0)\}}
{\sum_{\ell=1}^K\phi_{\ell,w_n}\exp\{\psi(\gamma_\ell)-\psi(\gamma_0)\}},
\label{eq:lh-lda-eta-update}\\
\gamma_k&=\alpha_k+\sum_{n=1}^N\eta_{nk}.
\label{eq:lh-lda-gamma-update}
\end{align}
每次更新使用其他因子的当前值，并将 $\eta_n$ 归一化。$-\psi(\gamma_0)$ 对所有候选话题相同，可从式~\eqref{eq:lh-lda-eta-update} 的分子分母中约去。
\end{proposition}
\begin{proof}
固定 $\gamma$ 时，与 $\eta_n$ 有关的目标为 $\sum_k\eta_{nk}(h_k+\log\phi_{k,w_n}-\log\eta_{nk})$。加上约束 $\sum_k\eta_{nk}=1$ 的拉格朗日项，对正分量求导，得 $h_k+\log\phi_{k,w_n}-\log\eta_{nk}-1+\lambda_n=0$，指数化、归一化即得第一式。固定所有 $\eta_n$ 时，由式~\eqref{eq:lh-lda-mean-field-coordinate}，
\[
\log q^*(\theta)
=\sum_k\left(\alpha_k-1+\sum_n\eta_{nk}\right)\log\theta_k+\text{常数},
\]
故最优因子为参数 $\alpha_k+\sum_n\eta_{nk}$ 的狄利克雷分布，得到第二式。该推导直接确定整个因子的坐标最优解，无需把耦合的 $\gamma$ 坐标分别求导后错误地视作互不相关。
\end{proof}

式~\eqref{eq:lh-lda-gamma-update} 是“先验参数加期望话题计数”，对应共轭更新；式~\eqref{eq:lh-lda-eta-update} 则结合当前词的似然和话题对数概率的期望。\emph[caution]{其中使用的是 $\exp\{\mathbb E_q\log\Theta_k\}$，不是 $\mathbb E_q\Theta_k$。}例如 $\gamma=(2,1)$ 时，利用 $\psi(s+1)-\psi(s)=1/s$，两个指数化对数期望的比为 $e$，而两个均值的比为 $2$，二者不能互换。

初始化可取 $\eta_{nk}=1/K$、$\gamma_k=\alpha_k+N/K$。每轮固定 $\gamma$ 更新全部 $\eta_n$，再更新整个 $\gamma$，重复至下界相对增量不超过给定容差或达到内层迭代上限。每个块更新都是当前块的最大化，因此下界不减；总体优化仍可能依赖初值。数值上可计算 $s_{nk}=\log\phi_{k,w_n}+\psi(\gamma_k)$，先减去 $\max_ks_{nk}$ 再指数化归一化，以避免溢出。

恢复文档下标后，固定全部变分参数，定义期望词计数
\begin{equation}
\overline d_{kv}=\sum_{m=1}^M\sum_{n=1}^{N_m}
\eta_{mnk}\mathbf1\{w_{mn}=v\},\qquad
\overline d_{k\cdot}=\sum_v\overline d_{kv}.
\label{eq:lh-lda-expected-counts}
\end{equation}
关于 $\phi$ 的 M 步为最大化 $\sum_{k,v}\overline d_{kv}\log\phi_{kv}$，约束每行非负且和为 $1$。逐行使用拉格朗日乘子得到
\begin{equation}
\phi_{kv}^{\mathrm{new}}=\frac{\overline d_{kv}}{\overline d_{k\cdot}},
\qquad\overline d_{k\cdot}>0.
\label{eq:lh-lda-phi-m-step}
\end{equation}
若某行总期望计数为 $0$，目标与该行无关，可保留原行；若某个词在语料中没有出现，无先验的该更新可能给它概率 $0$。这与完整贝叶斯模型中正 $\beta_v$ 带来的平滑不同。

记 $h_{mk}=\psi(\gamma_{mk})-\psi(\gamma_{m0})$，$\gamma_{m0}=\sum_k\gamma_{mk}$。关于 $\alpha$ 的目标为
\begin{equation}
F(\alpha)=M\left[\log\Gamma(\alpha_0)-\sum_k\log\Gamma(\alpha_k)\right]
+\sum_{m,k}(\alpha_k-1)h_{mk},\qquad\alpha_k>0.
\label{eq:lh-lda-alpha-objective}
\end{equation}
令 $g$ 为梯度、$H$ 为海森矩阵，则
\begin{align}
g_k&=M\bigl[\psi(\alpha_0)-\psi(\alpha_k)\bigr]+\sum_mh_{mk},
\label{eq:lh-lda-alpha-gradient}\\
H_{k\ell}&=M\psi'(\alpha_0)-M\mathbf1\{k=\ell\}\psi'(\alpha_k).
\label{eq:lh-lda-alpha-hessian}
\end{align}
这里 $\psi'$ 为三伽马函数（trigamma function）。注意两项均含文档数 $M$；该式直接由梯度微分得到。由于 $\log B(\alpha)$ 是指数族的对数归一化函数，其海森矩阵是 $(\log\Theta_k)_k$ 的协方差矩阵，故 $F$ 为凹函数，但有限最优点的存在性及数值求解仍需具体检查。

牛顿方向可通过解线性方程 $H\delta=-g$ 得到，再令 $\alpha^{\mathrm{new}}=\alpha+\rho\delta$，$0<\rho\leq1$。实现中不必显式求逆，并应通过回溯保持所有分量为正、使 $F$ 不减；未经阻尼的完整牛顿步不自动满足这些条件。$H$ 具有“对角矩阵加秩一矩阵”的结构：令 $a_k^{\mathrm H}=-M\psi'(\alpha_k)$、$b^{\mathrm H}=M\psi'(\alpha_0)$，则 $H=\operatorname{diag}(a^{\mathrm H})+b^{\mathrm H}\mathbf1\mathbf1^{\mathsf T}$。在矩阵可逆且分母非零时，
\begin{equation}
(H^{-1}g)_k=\frac{g_k-c^{\mathrm H}}{a_k^{\mathrm H}},\qquad
c^{\mathrm H}=\frac{\sum_kg_k/a_k^{\mathrm H}}
{1/b^{\mathrm H}+\sum_k1/a_k^{\mathrm H}}.
\label{eq:lh-lda-alpha-newton-structure}
\end{equation}
因此给定梯度及特殊函数值后，可在线性于 $K$ 的时间内求出方向；接近奇异时仍要关注数值稳定性。若 $\alpha$ 事先固定，直接省略这一 M 步即可。

\subsection{算法、完整模型推广与复杂度}
\label{subsec:lh-lda-variational-algorithm}

简化 LDA 的变分 EM 输入语料、$K$、初始正 $\alpha$、各行归一化且严格正的初始 $\phi$、内外层容差和迭代上限，输出模型参数近似 $\widehat\alpha,\widehat\phi$ 与各文档的变分参数 $\gamma_m,\eta_m$。交替执行以下两步，监测总 ELBO 的变化。
\begin{enumerate}
\item E 步：固定 $\alpha,\phi$，对各文档用式~\eqref{eq:lh-lda-eta-update} 和式~\eqref{eq:lh-lda-gamma-update} 交替更新，达到内层停止条件后保留局部变分参数。
\item M 步：固定 $\gamma,\eta$，累加式~\eqref{eq:lh-lda-expected-counts} 的期望计数，用式~\eqref{eq:lh-lda-phi-m-step} 更新 $\phi$；需要学习 $\alpha$ 时，以式~\eqref{eq:lh-lda-alpha-gradient} 和式~\eqref{eq:lh-lda-alpha-hessian} 作阻尼牛顿更新。各步均检查下界不减，达到外层容差或迭代上限后停止。
\end{enumerate}

训练后，$\mathbb E_q[\Theta_{mk}]=\gamma_{mk}/\gamma_{m0}$ 是文档话题比例的变分近似，$q(Z_{mn}=k)=\eta_{mnk}$ 是词位话题概率的近似；$\arg\max_k\eta_{mnk}$ 只是一种额外的硬指派规则。对新文档固定学得的 $\alpha,\phi$，只执行局部 E 步即可得到话题表示，不必重新为整套模型拟合全局词分布。

为连接到定义~\ref{def:lh-lda-model} 的完整模型，可在近似分布中再加入 $\prod_k\operatorname{Dir}(\phi_k\mid\lambda_k)$，其中 $\lambda_k\in(0,\infty)^V$。固定超参数 $\alpha,\beta$ 后，同样由式~\eqref{eq:lh-lda-mean-field-coordinate} 得到
\begin{align}
\eta_{mnk}&\propto\exp\left\{
\psi(\gamma_{mk})-\psi(\gamma_{m0})
+\psi(\lambda_{k,w_{mn}})-\psi(\lambda_{k0})\right\},
\label{eq:lh-lda-full-vi-eta}\\
\gamma_{mk}&=\alpha_k+\sum_n\eta_{mnk},&
\lambda_{kv}&=\beta_v+\overline d_{kv},\qquad\lambda_{k0}=\sum_v\lambda_{kv}.
\label{eq:lh-lda-full-vi-dirichlet}
\end{align}
第一式仍需对 $k$ 归一化，词分布的变分均值为 $\lambda_{kv}/\lambda_{k0}$。这里 $\lambda$ 是随机词分布的变分参数，而简化模型中的 $\phi$ 是点估计；给定 $\alpha,\beta$ 的这些更新属于后验变分推理，若进一步学习超参数，才增加相应的外层 M 步。这个推广展示了共轭先验在两条推断路线中的共同作用：吉布斯抽样用它解析积分并恢复条件后验，平均场推断用它保持坐标更新的狄利克雷形式。

最后明确朴素算法的计算量。对长度 $N_m$ 的文档，一轮局部坐标更新需 $O(N_mK)$；若第 $m$ 篇在一次外层 E 步中使用 $J_m$ 轮，则 E 步耗时 $O(\sum_mJ_mN_mK)$。M 步累加期望计数和归一化词分布需 $O(IK+KV)$；关于 $\alpha$ 的统计量 $\sum_mh_{mk}$ 可用 $O(MK)$ 预先累加，之后每次利用式~\eqref{eq:lh-lda-alpha-newton-structure} 的牛顿方向计算需 $O(K)$，另计线搜索费用。若各文档统一使用 $J$ 轮、外层共 $T$ 轮，则主项为 $O(T(JIK+IK+KV+MK))$，并加上超参数优化费用；不能在未说明迭代轮数时给出固定的总复杂度。

稠密保存所有 $\eta,\gamma,\phi$ 需 $O(IK+MK+KV)$ 空间；固定全局参数进行 E 步时，可以逐文档处理并只累加充分统计量，将工作空间降为 $O(K\max_mN_m+MK+KV)$，不计输入语料的 $O(I)$ 存储。完整模型还需保存 $K\times V$ 的 $\lambda$，量级不变。抽样误差与变分近似误差不同：增加平稳链的样本量可以降低蒙特卡罗误差，而增加平均场迭代主要改善受限目标的优化，不能消除分布族限制带来的偏差；比较两种方法时应同时说明模型版本、超参数、计算预算和评价指标。


\fi

\selectchapter{21}
\ifcompilechapter
\chapter{无监督学习方法总结}
\label{chap:lh-unsupervised-summary}

% 研读来源：机器学习/textbooks/机器学习（分章原文）/机器学习-22.pdf，教材第 435--438 页。
% 沿用教材两节及第一节的三个小节；条件说明、公式推导与例题为学习补充。

无监督学习从没有任务标签的数据中寻找结构，但“结构”可以指样本的分组、低维表示、文本中的潜在话题，也可以指图上结点的重要程度。这些任务并不共享一个唯一的模型或损失函数。本章沿用教材的两条总结路线：先区分具体任务、模型、学习策略和基础求解方法，再从矩阵分解与概率建模两个角度比较话题模型。比较时尤其要区分\emph{待求的数学对象}与\emph{求解它的算法}，以及算法的极限性质与有限计算得到的近似结果。

\section{无监督学习方法的关系和特点}
\label{sec:lh-unsup-relations}

\subsection{各种方法之间的关系}
\label{subsec:lh-unsup-method-relations}

教材介绍的方法可以沿“任务---模型与目标---求解工具”理解。聚类要得到样本的组别或属于各组的程度，典型方法有层次聚类、$k$ 均值聚类与高斯混合模型；降维要得到保留某些信息的低维表示，主成分分析是线性降维方法；话题分析用较少的潜在成分解释词与文档的关联，包括潜在语义分析、非负矩阵分解、概率潜在语义分析与潜在狄利克雷分配；图分析中的 PageRank 则把结点的重要程度表示为一个马尔可夫链的平稳分布。这些任务可以交叉，例如话题表示同时具有软分组和降维的作用，但输出相似不意味着模型和目标相同。

\emph{奇异值分解}（singular value decomposition，SVD）、\emph{非负矩阵分解}（nonnegative matrix factorization，NMF）、矩阵特征值求解，以及含隐变量概率模型的学习与推断，是可复用的基础工具。SVD 可以求主成分，也可以构造潜在语义表示；特征向量求解可以用于 PageRank；期望最大化算法、变分推断和马尔可夫链蒙特卡罗方法可以用于不同的隐变量模型。因此，SVD 本身并不规定数据来自什么分布，EM 也不对应一个固定的概率模型；这些工具还可以出现在监督学习或半监督学习中。

需要保留三个层次的区别。\emph{模型}规定允许怎样表示数据，例如高斯混合分布或低秩矩阵；\emph{学习策略}规定评价表示的准则，例如似然最大、重构误差最小或后验推断；\emph{算法}规定怎样计算，例如交替更新、SVD、Gibbs 抽样或幂法。同一模型可有不同的估计策略，同一目标也可能由不同算法求解。NMF 这个名称既常指非负因子模型，也常统称求解该模型的一类方法，使用时应把具体损失函数和更新规则另行说明。

\subsection{无监督学习方法}
\label{subsec:lh-unsup-task-methods}

\emph{硬聚类}（hard clustering）把每个样本分配到一个组，\emph{软聚类}（soft clustering）则保留各组的隶属程度；在概率模型中，这通常是给定观测后的隐类别后验概率。层次聚类在选定层次切割后给出硬划分，$k$ 均值直接给出硬划分，高斯混合模型通常给出软分配，也可以再取最大后验类别得到硬标签。层次聚类依赖距离、簇间连接准则和合并或分裂规则，并不普遍对应同一个“类内距离最小”的全局目标。

设有 $n\geq1$ 个样本 $x_i\in\mathbb R^d$，$i=1,\ldots,n$，预先给定簇数 $1\leq K\leq n$。$k$ 均值的一个明确目标是
\begin{equation}
\label{eq:lh-unsup-summary-kmeans-objective}
    \min_{\substack{c_i\in\{1,\ldots,K\}\\ \mu_1,\ldots,\mu_K\in\mathbb R^d}}
    \sum_{i=1}^{n}\|x_i-\mu_{c_i}\|_2^2.
\end{equation}
这里 $c_i$ 是分组变量，$\mu_k$ 是簇中心。通常的交替分配与中心更新不会增加目标，但联合问题并不凸，初始化、并列分配及空簇处理会影响结果。高斯混合模型则最大化
\[
    \sum_{i=1}^{n}\log\left[\sum_{k=1}^{K}
    \pi_k\mathcal N(x_i\mid\mu_k,\Sigma_k)\right],
    \qquad \pi_k\geq0,\quad\sum_k\pi_k=1,\quad\Sigma_k\succ0,
\]
其中 $\mathcal N(\cdot\mid\mu,\Sigma)$ 表示高斯密度。它的隐类别通过后验概率参与 EM 更新；这与式~\eqref{eq:lh-unsup-summary-kmeans-objective} 的硬分配目标不同。若允许协方差任意趋于奇异，高斯混合似然还可能无上界，讨论最大值存在或算法收敛时需要额外约束或正则化。

\emph{主成分分析}（principal component analysis，PCA）在线性正交子空间中保留尽量大的样本方差。令中心化数据矩阵为 $X\in\mathbb R^{d\times n}$，样本均值已从每列减去，$1\leq r\leq\min(d,n)$。其主方向矩阵可由
\begin{equation}
\label{eq:lh-unsup-pca-objective}
    \max_{A\in\mathbb R^{d\times r},\ A^{\mathsf T}A=I_r}
    \operatorname{tr}(A^{\mathsf T}XX^{\mathsf T}A)
\end{equation}
求得，$I_r$ 为 $r$ 阶单位矩阵，$\operatorname{tr}$ 为矩阵的迹。由
\[
    \|X-AA^{\mathsf T}X\|_F^2
    =\|X\|_F^2-\operatorname{tr}(A^{\mathsf T}XX^{\mathsf T}A)
\]
可知，最大化保留方差等价于最小化正交投影的平方重构误差；$\|B\|_F^2$ 表示矩阵各元素平方之和。$A$ 可取 $X$ 的前 $r$ 个左奇异向量。中心化、是否按特征尺度标准化及保留维数都是问题设定的一部分，不能单由“使用 SVD”确定。

\emph{潜在语义分析}（latent semantic analysis，LSA）对词---文档矩阵作截断 SVD，用低秩近似表示词与文档之间的关联。NMF 改用非负因子，便于把表示解释为非负成分的加和。\emph{概率潜在语义分析}（probabilistic latent semantic analysis，PLSA）通过离散隐话题解释词与文档的共现，通常以 EM 最大化观测似然；\emph{潜在狄利克雷分配}（latent Dirichlet allocation，LDA）进一步给文档的话题比例等量设置先验，通过后验推断学习。LSA 与平方损失 NMF 本身不要求概率解释；PLSA 的极大似然形式也没有给参数设置 LDA 那样的狄利克雷先验。具体联系见第~\ref{sec:lh-unsup-topic-relations} 节。

PageRank 采用另一种结构。对含 $s$ 个结点的有向图，取列随机转移矩阵 $P\in\mathbb R^{s\times s}$，即 $P_{ij}\geq0$、$\sum_iP_{ij}=1$，其中第 $j$ 列表示从结点 $j$ 出发的转移概率；无出链结点需先补成合法概率列。给定阻尼参数 $0<\alpha<1$ 和概率向量 $v$，定义
\begin{equation}
\label{eq:lh-unsup-pagerank}
    G=\alpha P+(1-\alpha)v\mathbf1^{\mathsf T},\qquad
    p=Gp,\quad p\geq0,\quad\mathbf1^{\mathsf T}p=1,
\end{equation}
其中 $\mathbf1$ 为全 $1$ 列向量。$p$ 是修正后随机游走的平稳分布，可用 $p^{(t+1)}=Gp^{(t)}$ 求解。对两个概率向量 $a,b$，有 $\mathbf1^{\mathsf T}(a-b)=0$，因而
\[
    \|Ga-Gb\|_1=\alpha\|P(a-b)\|_1
    \leq\alpha\|a-b\|_1.
\]
概率单纯形在 $\ell_1$ 距离下完备，映射 $G$ 将其映到自身，故压缩映射定理保证平稳概率向量唯一，并且从任意概率初值出发几何收敛。若 $v$ 的每个分量都正，$G$ 还严格为正，平稳分布也严格为正。此处算法目标是求一个概率向量，不是拟合带标签的结点分数。

表~\ref{tab:lh-unsup-method-summary} 把以上任务按模型、策略与算法并列比较。表中的算法是常见选择，不能据此把某个模型限定为只能使用这一算法。
\begin{table}[htbp]
\centering
\caption{无监督学习方法的模型、策略与常见算法（据教材总结并补明条件）}
\label{tab:lh-unsup-method-summary}
\begin{tabularx}{\textwidth}{lYYY}
\toprule
方法 & 模型或输出 & 学习或求解策略 & 常见算法 \\
\midrule
层次聚类 & 层次化的样本划分 & 按指定距离与连接准则组织层次 & 合并或分裂启发式 \\
$k$ 均值 & $K$ 个中心与硬划分 & 类内平方距离和最小 & 分配与均值交替更新 \\
高斯混合 & 高斯成分的概率混合 & 观测数据似然最大 & EM \\
PCA & 低维正交子空间 & 投影方差最大或投影误差最小 & SVD 或对称特征值分解 \\
LSA & 词---文档矩阵的低秩近似 & 平方重构误差最小 & 截断 SVD \\
NMF & 两个非负因子的乘积 & 指定非负分解损失最小 & 交替优化等 \\
PLSA & 词与文档的隐话题模型 & 观测词频的似然最大 & EM \\
LDA & 带狄利克雷先验的话题模型 & 给定语料的后验推断 & Gibbs 抽样或变分推断 \\
PageRank & 图上随机游走的平稳分布 & 求满足规范化条件的平稳向量 & 幂迭代等 \\
\bottomrule
\end{tabularx}
\end{table}

\subsection{基础机器学习方法}
\label{subsec:lh-unsup-basic-methods}

矩阵分解的主要区别来自因子的约束和逼近误差的定义。对任意实矩阵 $D$，SVD 都存在；左右奇异向量的正交性是分解的数学性质，\emph[caution]{不要求原始特征或样本彼此正交}。用于学习时，真正需要选择的是截断维数、数据表示与误差准则。NMF 要求输入矩阵和两个因子逐元素非负，通常只能求近似分解；低维非负因子是否足以准确重构数据，属于表示假设。NMF 联合优化一般非凸，因子还存在正比例缩放与置换等不唯一性。

含隐变量概率模型的学习与推断则要处理对隐变量求和或积分。设 $x$ 为已观测数据，$Z$ 为隐变量，$\theta$ 为待估参数，联合概率或密度为 $p_\theta(x,z)$。对试探分布 $q(z)$，在下式各项有限且 $q$ 的支撑与模型相容的条件下，定义\emph{证据下界}（evidence lower bound，ELBO）
\begin{equation}
\label{eq:lh-unsup-elbo}
    \mathcal F(q,\theta)
    =\mathbb E_q[\log p_\theta(x,Z)]-\mathbb E_q[\log q(Z)].
\end{equation}
离散变量的期望为求和，连续变量的期望为积分。由 $p_\theta(x,z)=p_\theta(z\mid x)p_\theta(x)$，直接得到
\begin{equation}
\label{eq:lh-unsup-elbo-identity}
    \log p_\theta(x)=\mathcal F(q,\theta)
    +\operatorname{KL}\bigl(q\|p_\theta(\cdot\mid x)\bigr),
\end{equation}
其中 $\operatorname{KL}(q\|p)=\mathbb E_q[\log(q(Z)/p(Z))]\geq0$。该非负性可由 $\log t\leq t-1$ 推出；等号当且仅当两个分布几乎处处相等。因此下界在 $q$ 等于精确后验时取到观测对数似然。

\emph{期望最大化算法}（expectation-maximization algorithm，EM）的 E 步取 $q_t(z)=p_{\theta_t}(z\mid x)$，M 步选择使 $\mathcal F(q_t,\theta)$ 不下降的新参数。在 E 步可精确计算、更新合法且各项有限的条件下，
\[
    \log p_{\theta_{t+1}}(x)
    \geq\mathcal F(q_t,\theta_{t+1})
    \geq\mathcal F(q_t,\theta_t)
    =\log p_{\theta_t}(x).
\]
这是似然的单调性保证。若似然值另有有限上界，其数值序列收敛；参数序列收敛、极限点为驻点等结论还需要连续性、更新映射和有界性等条件，更不能由单调性推出全局最优。后验虽可能有解析式，整个 EM 优化仍可能需要多次迭代并依赖初值。

\emph{变分推断}（variational inference）把 $q$ 限制在便于计算的分布族 $\mathcal Q$ 内。在固定参数下，由式~\eqref{eq:lh-unsup-elbo-identity}，最大化 $\mathcal F(q,\theta)$ 等价于最小化 $\operatorname{KL}(q\|p_\theta(\cdot\mid x))$；若 $\mathcal Q$ 不包含真实后验，最优下界通常仍有间隙。变分 EM 交替更新 $q$ 与 $\theta$，精确的坐标更新保证下界不下降，\emph[caution]{一般不保证真实观测似然逐步不下降}。在完全贝叶斯模型中，也可以把未知参数并入隐变量，一同近似其后验分布。

\emph{马尔可夫链蒙特卡罗方法}（Markov chain Monte Carlo，MCMC）通过构造以目标分布为不变分布的马尔可夫链作随机近似。\emph{吉布斯抽样}（Gibbs sampling）是一种 MCMC 方法，逐个或分块从完整条件分布抽取隐变量或未知参数；Metropolis--Hastings 方法则利用提议与接受拒绝步骤构造转移。有限状态链不可约且非周期时，其分布趋近唯一平稳分布；一般状态空间可采用正 Harris 常返（positive Harris recurrence）及非周期性等条件保证相应的分布收敛。在遍历定理适用且 $f$ 可积时，样本平均
\begin{equation}
\label{eq:lh-unsup-mcmc-average}
    \frac1T\sum_{t=1}^{T}f(Z^{(t)})
    \longrightarrow\mathbb E_{\pi}[f(Z)]
\end{equation}
几乎必然成立，其中 $\pi$ 为目标分布。链中样本通常相关，有限运行的误差受混合速度和自相关影响。\emph[caution]{目标后验的渐近抽样保证不是参数全局最优的保证}；后验可能多峰，样本也不会因运行时间增加就全部集中到其众数。

因此，EM 适合精确后验期望可计算、参数更新可处理的模型；变分推断适合能设计可计算近似分布族的模型；Gibbs 抽样适合完整条件分布可抽样、且链能有效探索后验的模型。不能仅按“简单模型用 EM，复杂模型用抽样”机械选择，更不能把后验分布近似和点参数优化混为同一问题。

最后，\emph{幂法}（power method）是矩阵特征向量的基础求解方法。若方阵 $B$ 可对角化，有唯一模最大的特征值 $\lambda_1$，满足 $|\lambda_1|>|\lambda_2|\geq\cdots$，且初向量在第一特征向量方向上的展开系数不为零，则反复乘 $B$ 并规范化可以逼近主特征方向，误差渐近受 $|\lambda_2/\lambda_1|$ 控制。若主特征值为负或非实数，规范化向量还可能交替变号或旋转，应区分向量本身与其方向的收敛。PageRank 的正概率向量迭代具有式~\eqref{eq:lh-unsup-pagerank} 的特殊结构，可由压缩性直接论证，无须把任意矩阵上的幂法都视为必然收敛。

\section{话题模型之间的关系和特点}
\label{sec:lh-unsup-topic-relations}

话题模型的共同出发点是用少量潜在成分表示较大的词---文档关联矩阵。统一记数据矩阵为 $D\in\mathbb R^{M\times N}$，$M$ 为词的数量，$N$ 为文档的数量；选取 $1\leq r\leq\min(M,N)$ 个成分，以
\begin{equation}
\label{eq:lh-unsup-factorization}
    D\approx UV,\qquad U\in\mathbb R^{M\times r},\quad
    V\in\mathbb R^{r\times N}
\end{equation}
表示。$U$ 的第 $k$ 列描述第 $k$ 个成分在词空间中的方向或分布，$V$ 的第 $n$ 列描述文档 $n$ 的成分系数。通常 $r$ 远小于 $M,N$，因此乘积的秩至多为 $r$，但不要求原矩阵恰好具有秩 $r$。不同模型对 $D$ 的预处理、误差和因子约束作不同规定，因而式~\eqref{eq:lh-unsup-factorization} 只给出共同形式，不给出模型等价性。

\begin{definition}[矩阵上的 Bregman 散度]
\label{def:lh-unsup-bregman}
设 $\Phi$ 是某个凸矩阵域上可微的严格凸函数，$D$、$Q$ 属于其可微定义域。对应的\emph{Bregman 散度}（Bregman divergence）定义为
\begin{equation}
\label{eq:lh-unsup-bregman}
    B_\Phi(D\|Q)=\Phi(D)-\Phi(Q)-\langle\nabla\Phi(Q),D-Q\rangle,
\end{equation}
其中 $\langle A,B\rangle=\sum_{m,n}a_{mn}b_{mn}$ 是矩阵内积。严格凸性保证其非负，且仅在 $D=Q$ 时为零；它一般不对称，也不满足三角不等式，故不应称为距离。边界元素为零时，需针对所选函数定义适当的极限延拓。
\end{definition}

取 $\Phi(D)=\tfrac12\|D\|_F^2$，由 $\nabla\Phi(Q)=Q$ 可得 $B_\Phi(D\|Q)=\tfrac12\|D-Q\|_F^2$。若矩阵逐元素为正，取 $\Phi(D)=\sum_{m,n}(d_{mn}\log d_{mn}-d_{mn})$，则
\begin{equation}
\label{eq:lh-unsup-generalized-kl}
    B_\Phi(D\|Q)=\sum_{m,n}
    \left[d_{mn}\log\frac{d_{mn}}{q_{mn}}-d_{mn}+q_{mn}\right].
\end{equation}
这一\emph{广义 KL 散度}（generalized Kullback--Leibler divergence）允许非负矩阵：当 $d_{mn}=0$ 时对数项按极限取 $0$，当 $d_{mn}>0,q_{mn}=0$ 时相应项为 $+\infty$。若 $D$、$Q$ 的元素和同为 $1$，线性项抵消，得到概率分布之间的 KL 散度。若两者未等质量规范化，则不能随意删去式~\eqref{eq:lh-unsup-generalized-kl} 的线性项。

在这一框架下，LSA 对选定词权重矩阵最小化平方误差，NMF 对非负矩阵最小化指定损失并约束因子非负，PLSA 则对规范化词频分布最小化 KL 散度并约束因子具有概率解释。表~\ref{tab:lh-unsup-topic-factorizations} 采用下面推导中的参数化；其目的是把损失与约束写清，并非给三个模型强加相同的数据预处理。
\begin{table}[htbp]
\centering
\caption{矩阵分解视角下的话题模型}
\label{tab:lh-unsup-topic-factorizations}
\begin{tabularx}{\textwidth}{lYY}
\toprule
方法 & 目标函数 & 因子约束与解释 \\
\midrule
LSA & $\|D-UV\|_F^2$ & 可取 $U^{\mathsf T}U=I_r$，$V$ 为实系数，吸收奇异值 \\
平方损失 NMF & $\|D-UV\|_F^2$ & $D,U,V$ 逐元素非负 \\
PLSA & $\sum_{m,n}d_{mn}\log\bigl(d_{mn}/(UV)_{mn}\bigr)$ & $U,V\geq0$，$U$ 各列和为 $1$，$V$ 第 $n$ 列和为文档概率 $p_n$ \\
\bottomrule
\end{tabularx}
\end{table}

对 LSA，设 $D=L\Sigma R^{\mathsf T}$ 为 SVD，取前 $r$ 个奇异向量，得到
\begin{equation}
\label{eq:lh-unsup-lsa-factors}
    D_r=L_r\Sigma_rR_r^{\mathsf T}=UV,\qquad
    U=L_r,\quad V=\Sigma_rR_r^{\mathsf T}.
\end{equation}
由\crosschapterref{15}{定理~\ref{thm:lh-svd-best-approximation}}{截断 SVD 的最优性定理} 的最佳低秩近似性质，当 $r$ 不超过矩阵秩时直接应用该定理；当 $r$ 超过矩阵秩时补零奇异值并以 $D$ 本身重构，统一得到
\[
    \min_{\operatorname{rank}(Q)\leq r}\|D-Q\|_F^2
    =\sum_{j>r}\sigma_j^2,
\]
其中 $\sigma_j$ 是降序排列的奇异值，求和遍及被舍弃的奇异值。因而式~\eqref{eq:lh-unsup-lsa-factors} 是平方误差下的全局最优低秩近似。这一参数化满足 $U^{\mathsf T}U=I_r$、$VV^{\mathsf T}=\Sigma_r^2$，所以 $V$ 的行两两正交，却一般不是单位长度。\emph[caution]{写成两因子 $UV$ 时，不能一面把奇异值吸收到 $V$，一面又要求 $V$ 的行全部单位正交}。若要把两个正交因子单独列出，就应保留中间的 $\Sigma_r$。重奇异值还可能导致最优子空间的基不唯一。

对平方损失 NMF，约束是 $U\geq0,V\geq0$，各不等式均按元素理解。这里允许零元素，不应把“非负”误写成“严格正”。每个文档向量由非负成分加和得到，因子无需正交，联合优化一般无全局最优保证。对任意正对角矩阵 $S$，$UV=(US)(S^{-1}V)$，说明系数的尺度不能脱离规范化约定解释。也可以对 NMF 选择式~\eqref{eq:lh-unsup-generalized-kl} 的损失；“采用非负因子”与“采用平方损失”是两个分别作出的选择。

PLSA 的概率约束需要先说明矩阵怎样规范化。设词频为非负整数 $c_{mn}$，总词数为 $C=\sum_{m,n}c_{mn}>0$，只保留非空文档，定义
\begin{equation}
\label{eq:lh-unsup-plsa-data}
    d_{mn}=\frac{c_{mn}}C,\qquad p_n=\sum_m d_{mn}>0,
    \qquad\sum_n p_n=1.
\end{equation}
记 $W$、$J$、$Z$ 分别为一个词出现的词类型、所属文档编号与隐话题，取值分别为 $m,n,k$。PLSA 设给定话题后词类型与文档条件独立，其一种等价参数化为
\begin{equation}
\label{eq:lh-unsup-plsa-factorization}
    q_{mn}=P(W=m,J=n)=p_n\sum_{k=1}^{r}\phi_{mk}\theta_{kn},
    \quad \sum_m\phi_{mk}=1,\quad\sum_k\theta_{kn}=1,
    \quad\phi_{mk},\theta_{kn}\geq0.
\end{equation}
取 $u_{mk}=\phi_{mk}$、$v_{kn}=p_n\theta_{kn}$，就有 $Q=UV$、$\sum_m u_{mk}=1$、$\sum_k v_{kn}=p_n$，从而 $\sum_{m,n}q_{mn}=1$。这里 $U$ 的列是词分布，而 $V$ 的列带有文档概率质量；若把 $V$ 直接定义成条件话题比例 $\theta_{kn}$，便须另乘 $p_n$，不能混用两种参数化。

\begin{proposition}[PLSA 似然与 KL 目标]
\label{prop:lh-unsup-plsa-kl}
在词频数据和文档概率按式~\eqref{eq:lh-unsup-plsa-data} 固定、词出现按式~\eqref{eq:lh-unsup-plsa-factorization} 的模型计算似然时，PLSA 的极大似然估计等价于在上述概率约束下最小化 $\operatorname{KL}(D\|UV)$。
\end{proposition}

\begin{proof}
在模型所采用的条件独立词生成假设下，给定各文档词数的对数似然除去仅依赖计数的常数为
\[
    \ell(\phi,\theta)
    =\sum_{m,n}c_{mn}\log\left(\sum_k\phi_{mk}\theta_{kn}\right)
    =C\sum_{m,n}d_{mn}\log q_{mn}
     -C\sum_n p_n\log p_n.
\]
最后一项与待优化因子无关，而
\[
    \operatorname{KL}(D\|Q)
    =\sum_{m,n}d_{mn}\log d_{mn}
     -\sum_{m,n}d_{mn}\log q_{mn}.
\]
第一项由数据固定，故两个优化等价。若某个正词频对应 $q_{mn}=0$，对数似然为 $-\infty$，KL 散度为 $+\infty$，两种表述在边界上仍一致。这里用的是原始词频的规范化；任意实值加权矩阵或未经说明的 TF--IDF 权重不能直接解释成同一个词频似然。
\end{proof}

这一联系也说明，KL 损失 NMF 在附加相应概率规范化条件后可获得 PLSA 解释，但一般的 NMF 不能仅凭“存在两个非负因子”就与 PLSA 等同。PLSA 的训练通常由 EM 计算隐话题后验责任度，再更新词分布与文档话题比例；有限样本下学到的是每个训练文档的参数，模型本身没有为新文档的话题比例规定 LDA 那样的先验。对新文档可固定词分布重新估计其话题比例，这与给新文档定义完整的先验生成机制是不同的做法。

教材进一步从概率图模型的角度解释 LSA 和 NMF。把第 $n$ 个文档向量记为 $\boldsymbol d_n\in\mathbb R^M$，系数记为 $v_n\in\mathbb R^r$，在给定参数和系数时，假设各文档独立且
\begin{equation}
\label{eq:lh-unsup-gaussian-factor-model}
    \boldsymbol D_n\mid U,v_n\sim\mathcal N(Uv_n,\sigma^2I_M),
    \qquad \sigma^2>0.
\end{equation}
这里大写 $\boldsymbol D_n$ 为随机向量，小写 $\boldsymbol d_n$ 为观测值，$\sigma^2$ 先视为给定。对 $N$ 个文档，负对数似然为
\[
    -\log p(D\mid U,V)
    =\frac{MN}{2}\log(2\pi\sigma^2)
     +\frac{1}{2\sigma^2}\sum_{n=1}^{N}\|\boldsymbol d_n-Uv_n\|_2^2.
\]
因此，在相同的因子约束下最大化这一高斯似然，恰好等价于最小化平方重构误差。选正交基的约束可解释 LSA，选非负因子约束可解释平方损失 NMF；这是一种对文档数值向量的连续噪声模型解释，不会自动把向量的分量变成词概率，也不是对整数词频的多项分布建模。高斯噪声允许负观测值，非负约束限制的是分解因子和重构均值，不能混淆这两个层面。

LDA 与 PLSA 的主要差别则来自层次化先验。令 $\phi_k$ 为第 $k$ 个话题的词概率向量，$\theta_n$ 为第 $n$ 个文档的话题概率向量；完全贝叶斯的常见形式为
\begin{equation}
\label{eq:lh-unsup-lda-generative}
\begin{aligned}
    \phi_k&\sim\operatorname{Dir}(\beta), & k&=1,\ldots,r,\\
    \theta_n&\sim\operatorname{Dir}(\alpha), & n&=1,\ldots,N,\\
    Z_{ni}\mid\theta_n&\sim\operatorname{Cat}(\theta_n),&
    W_{ni}\mid Z_{ni}=k,\phi_k&\sim\operatorname{Cat}(\phi_k).
\end{aligned}
\end{equation}
其中 $i=1,\ldots,L_n$ 为文档 $n$ 内的词位置，$L_n$ 为给定文档长度，$\operatorname{Dir}$、$\operatorname{Cat}$ 分别表示狄利克雷分布与类别分布，超参数 $\alpha\in\mathbb R_{>0}^{r}$、$\beta\in\mathbb R_{>0}^{M}$ 给定；各文档话题比例及各话题词分布按先验独立抽取，各词在给定相应变量后条件独立生成。实际推断目标是给定词观测的 $p(\Theta,\Phi,Z\mid W,\alpha,\beta)$，其中 $\Theta=(\theta_n)_{n=1}^{N}$、$\Phi=(\phi_k)_{k=1}^{r}$，可用 Gibbs 抽样或变分方法近似。PLSA 主要估计固定参数，LDA 则把相关量纳入随机生成与后验推断；这种先验结构会改变估计及新文档的建模方式，但不提供无条件的准确率或全局最优保证。

\begin{example}
\label{ex:lh-unsup-factor-constraints}
设规范化词---文档矩阵为
\[
    D=\begin{pmatrix}1/2&0\\0&1/2\end{pmatrix},\qquad M=N=2,
\]
用一个成分 $r=1$ 比较平方误差低秩近似与 PLSA 的概率约束。
\end{example}

\begin{solve}
两个奇异值都是 $1/2$，故任意秩至多为 $1$ 的矩阵，其最小平方重构误差为 $1/4$。例如
\[
    U=\begin{pmatrix}1\\0\end{pmatrix},\qquad
    V=\begin{pmatrix}1/2&0\end{pmatrix}
\]
给出一个 LSA 最优解，同时也是平方损失 NMF 的可行最优解。但乘积矩阵的元素和仅为 $1/2$，第二列全为零，不能作为本例 PLSA 的联合概率分布；其对应的正词频还会产生无穷 KL 散度。

PLSA 中 $p_1=p_2=1/2$，只有一个话题时两篇文档的词分布必须相同。写成
\[
    U=\begin{pmatrix}a\\1-a\end{pmatrix},\qquad
    V=\begin{pmatrix}1/2&1/2\end{pmatrix},\qquad0\leq a\leq1,
\]
则有限 KL 目标要求 $0<a<1$，且
\[
    \operatorname{KL}(D\|UV)
    =-\frac12\log a-\frac12\log(1-a).
\]
由 $a(1-a)\leq1/4$，最小值在 $a=1/2$ 取得，最小 KL 散度为 $\log2$，重构矩阵的四个元素均为 $1/4$。它的平方误差也恰为 $1/4$，但两个目标的数值不可直接比较；这个例子说明，即使平方误差最优解不唯一，概率约束仍会排除其中一些解，不能把任意低秩分解直接当作话题概率模型。
\end{solve}

本章各方法的联系可以据此把握：矩阵分解提供低维表示，概率模型说明数据和隐变量怎样生成，学习策略规定需要优化或推断什么，算法决定如何近似完成计算。判断两种方法是否等价，应逐一核对数据表示、损失、约束和概率规范化；判断算法保证，应明确它针对目标值、参数、分布还是样本平均。共同使用低维成分或同一种求解工具，都不足以跳过这些条件。



\fi

\selectchapter{22}
\ifcompilechapter
\chapter{PageRank 算法}
\label{chap:lihang-pagerank}

% 研读来源：机器学习/textbooks/机器学习（分章原文）/机器学习-21.pdf，教材第 415--433 页。
% 沿用教材的两节、七个小节；补充悬挂结点修正、收敛证明、误差界与习题解答。
% 教材一般定义中的 M 须为列随机矩阵；零列修正不能仅由阻尼平滑项替代。

网页、引文和社交关系都可以表示为有向图。若只用入边数衡量结点的重要程度，就把所有来源看得同样重要，也没有考虑一个来源把注意力分给了多少结点。\emph{PageRank} 是一种无监督的链接分析（link analysis）方法：让游走者沿有向边随机移动，以长期访问各结点的概率刻画其相对重要程度。一个结点得到的权重取决于指向它的结点已有的权重，因此问题归结为概率分布的递推及其不动点（fixed point）。本章先定义图上的随机游走，说明基本 PageRank 的条件，再通过悬挂结点修正与随机跳转得到适用于任意有限有向图的一般定义，最后讨论迭代算法、幂法和线性方程组解法。

\section{PageRank 的定义}
\label{sec:lihang-pr-definition}

\subsection{基本想法}
\label{subsec:lihang-pr-idea}

把网页看作结点，把网页 $v_j$ 指向网页 $v_i$ 的超链接看作有向边 $v_j\to v_i$。设一个游走者到达 $v_j$ 后，在该结点的所有出边中等概率选择一条，并沿边访问下一个结点。如果这种访问过程的状态分布趋于稳定，就用稳定概率作为各结点的 PageRank 值。这里的概率由指定的图结构和游走机制决定，不是对真实用户行为频率的直接测量；在网页检索中，它给出与查询无关的链接重要度，实际排序还需要考虑网页与查询的相关程度。

记结点 $v_j$ 的出度为 $k_j$。若其 PageRank 值为 $R_j$，它沿每条出边贡献 $R_j/k_j$ 的权重，结点 $v_i$ 的权重则由所有入边的贡献相加得到。因此，\emph{来源的重要程度与来源的出度共同决定一条链接的贡献}：来自高权重结点的链接贡献可能较大，同一来源连出的链接越多，单条链接分得的权重越少。这一递归关系需要连同非负性和总和为 $1$ 的条件一起求解，不能靠一次入度统计完成。

\subsection{有向图和随机游走模型}
\label{subsec:lihang-pr-graph-walk}

\begin{definition}[有向图与连通性]
\label{def:lihang-pr-directed-graph}
有限\emph{有向图}（directed graph）记为 $G=(V,E)$，其中 $V=\{v_1,\ldots,v_n\}$，$n\geq1$，$E\subseteq V\times V$ 是有向边集合。本章允许自环，不考虑重边。按边的方向依次连接的边序列称为有向路径，其长度是所经过的边数；若任意两个不同结点之间均有双向可达的有向路径，则称图为\emph{强连通图}（strongly connected graph）。记
\[
    k_j=\bigl|\{i:(v_j,v_i)\in E\}\bigr|,\qquad
    \mathcal I_i=\{j:(v_j,v_i)\in E\},\qquad
    \mathcal D=\{j:k_j=0\},
\]
分别为出度、指向 $v_i$ 的结点索引集合和\emph{悬挂结点}（dangling node）的索引集合。
\end{definition}

先假定所有 $k_j>0$。令随机变量 $X_t\in\{1,\ldots,n\}$ 表示时刻 $t$ 所在结点的索引，$t=0,1,\ldots$。沿出边等概率移动给出时间齐次的一阶马尔可夫链（time-homogeneous first-order Markov chain）：在具有正概率的历史上，给定当前状态后，下一个状态的条件分布与更早的状态无关，且
\begin{equation}
\label{eq:lihang-pr-transition}
    P(X_{t+1}=i\mid X_t=j)=m_{ij},\qquad
    m_{ij}=\begin{cases}
        1/k_j,&(v_j,v_i)\in E,\\
        0,&(v_j,v_i)\notin E.
    \end{cases}
\end{equation}
矩阵 $M=(m_{ij})\in\mathbb R^{n\times n}$ 称为转移矩阵（transition matrix）。本章沿用教材的\emph{列随机矩阵}（column-stochastic matrix）约定：第 $j$ 列表示从 $v_j$ 出发的转移分布，因此
\begin{equation}
\label{eq:lihang-pr-stochastic}
    m_{ij}\geq0,\qquad \sum_{i=1}^{n}m_{ij}=1,\qquad
    \boldsymbol 1^{\mathsf T}M=\boldsymbol 1^{\mathsf T},
\end{equation}
其中 $\boldsymbol 1$ 是 $n$ 维全 $1$ 列向量。采用行分布向量的文献会使用转置后的约定，不能把两种方向的矩阵与递推公式混用。

定义概率单纯形（probability simplex）与时刻 $t$ 的状态分布为
\[
    \Delta_n=\{r\in\mathbb R^n:r_i\geq0,\ \boldsymbol 1^{\mathsf T}r=1\},\qquad
    R_t=\bigl(P(X_t=1),\ldots,P(X_t=n)\bigr)^{\mathsf T}\in\Delta_n.
\]
由全概率公式，$P(X_{t+1}=i)=\sum_jm_{ij}P(X_t=j)$，从而
\begin{equation}
\label{eq:lihang-pr-distribution-update}
    R_{t+1}=MR_t,\qquad R_t=M^tR_0.
\end{equation}
$X_t$ 是一次随机游走中的随机状态，$R_t$ 则是该状态的完整分布；矩阵迭代直接计算分布，并不需要实际抽样一条游走轨迹。

\begin{definition}[平稳分布与周期]
\label{def:lihang-pr-stationarity-period}
若 $R\in\Delta_n$ 满足 $MR=R$，则称 $R$ 为该链的\emph{平稳分布}（stationary distribution）。对至少能返回自身的状态 $j$，其周期定义为
\[
    h_j=\gcd\{t\geq1:(M^t)_{jj}>0\}.
\]
若 $h_j=1$，称该状态非周期；若 $h_j>1$，称其具有周期性。有限强连通图且各结点有出边时，所有状态可以相互到达，马尔可夫链称为\emph{不可约}（irreducible），各状态具有相同周期；该共同周期为 $1$ 时，称图及链\emph{非周期}（aperiodic）。
\end{definition}

由于各条已有边的转移概率为正，$(M^t)_{ij}>0$ 当且仅当存在从 $v_j$ 到 $v_i$ 的长度为 $t$ 的有向路径。周期是所有正长度返回路径长度的最大公约数，并不是任取一条闭合路径的长度。例如，单向三环的返回长度都是 $3$ 的倍数，周期为 $3$；强连通图中若有一个自环，则共同周期为 $1$；没有自环也可能非周期，例如同一结点存在长度为 $2$ 和 $3$ 的返回路径时，最大公约数也是 $1$。

\subsection{PageRank 的基本定义}
\label{subsec:lihang-pr-basic-definition}

\begin{theorem}[有限不可约非周期链的收敛]
\label{thm:lihang-pr-finite-chain}
设 $M$ 为有限维列随机矩阵，对应链不可约且非周期。则存在唯一平稳分布 $R\in\Delta_n$，其各分量严格为正，并且对任意初始分布 $R_0\in\Delta_n$，都有
\begin{equation}
\label{eq:lihang-pr-basic-limit}
    \lim_{t\to\infty}M^tR_0=R.
\end{equation}
\end{theorem}

\begin{proof}
先说明平稳分布的存在，这一步只需要有限性与列随机性。任取 $q\in\Delta_n$，令 $\bar q_N=N^{-1}\sum_{t=0}^{N-1}M^tq$，则 $\bar q_N\in\Delta_n$，且
\[
    M\bar q_N-\bar q_N=\frac{M^Nq-q}{N},\qquad
    \|M\bar q_N-\bar q_N\|_1\leq\frac2N,
\]
其中 $\|x\|_1=\sum_i|x_i|$。单纯形为紧集，可取收敛子列 $\bar q_{N_l}\to R\in\Delta_n$；利用矩阵乘法的连续性取极限，得到 $MR=R$。

不可约与非周期性进一步保证存在整数 $L\geq1$，使 $B=M^L$ 的全部元素严格为正。其图论依据是：每个结点的返回长度集合的最大公约数为 $1$，可以选出有限多个返回长度使其最大公约数仍为 $1$；这些长度的非负整数线性组合包含所有充分大的整数。将相应返回路径拼接，再接上一条到达目标结点的固定路径，就能得到任意两个结点之间所有充分大的长度的路径；结点对有限，故可选一个共同的 $L$。

令 $\eta=\min_{i,j}b_{ij}>0$、$\beta=n\eta\in(0,1]$，并记 $u=\boldsymbol1/n$。若 $\beta<1$，则
\[
    B=\beta u\boldsymbol1^{\mathsf T}+(1-\beta)C,
\]
其中 $C=(B-\eta\boldsymbol1\boldsymbol1^{\mathsf T})/(1-\beta)$ 仍为列随机矩阵。任意列随机矩阵 $C$ 都满足
\[
    \|Cz\|_1
    \leq\sum_i\sum_j c_{ij}|z_j|=\|z\|_1.
\]
对 $p,q\in\Delta_n$，因 $\boldsymbol1^{\mathsf T}(p-q)=0$，有
\[
    \|Bp-Bq\|_1\leq(1-\beta)\|p-q\|_1.
\]
若 $\beta=1$，则 $B=u\boldsymbol1^{\mathsf T}$，上述差直接为零。因而两个 $M$ 的平稳分布也必为 $B$ 的不动点，收缩不等式迫使它们相同。对已得到的平稳分布，$R=BR$ 给出 $R_i\geq\eta\sum_jR_j=\eta>0$。最后写 $t=kL+r$，其中 $0\leq r<L$，利用 $M^r$ 不增大 $1$ 范数，得到
\[
    \|M^tR_0-R\|_1
    \leq\|B^kR_0-R\|_1
    \leq(1-\beta)^k\|R_0-R\|_1\longrightarrow0.
\]
在 $\beta=1$ 时，$k\geq1$ 后该差已为零。
\end{proof}

\begin{definition}[基本 PageRank]
\label{def:lihang-pr-basic}
给定有限强连通、非周期且各结点均有出边的有向图，按式~\eqref{eq:lihang-pr-transition} 定义随机游走。其唯一平稳分布 $R$ 称为该图的基本 PageRank，结点 $v_i$ 的 PageRank 值记为 $\operatorname{PR}(v_i)=R_i$。它由
\begin{equation}
\label{eq:lihang-pr-basic-equations}
    MR=R,\qquad R_i>0,\qquad \sum_iR_i=1
\end{equation}
确定，分量形式为
\begin{equation}
\label{eq:lihang-pr-basic-components}
    R_i=\sum_{j\in\mathcal I_i}\frac{R_j}{k_j},\qquad i=1,\ldots,n.
\end{equation}
\end{definition}

\begin{example}
\label{ex:lihang-pr-basic-four}
考虑教材的四结点图，按 $(A,B,C,D)$ 排列结点，边为 $A\to B,C,D$，$B\to A,D$，$C\to A$，$D\to B,C$。从均匀初始分布出发，求基本 PageRank。
\end{example}

\begin{solve}
按出发结点逐列填写矩阵，得到
\begin{equation}
\label{eq:lihang-pr-four-matrix}
    M=\begin{pmatrix}
        0&1/2&1&0\\
        1/3&0&0&1/2\\
        1/3&0&0&1/2\\
        1/3&1/2&0&0
    \end{pmatrix}.
\end{equation}
图是强连通的，且从 $A$ 出发有长度为 $2$ 的返回路径 $A\to B\to A$ 和长度为 $3$ 的返回路径 $A\to D\to C\to A$，故非周期。取 $R_0=(1,1,1,1)^{\mathsf T}/4$，前两次迭代为
\[
    R_1=\frac1{24}(9,5,5,5)^{\mathsf T},\qquad
    R_2=\frac1{48}(15,11,11,11)^{\mathsf T}.
\]
为求精确极限，令 $R=(a,b,c,e)^{\mathsf T}$。方程 $MR=R$ 给出
\[
    a=b/2+c,\qquad b=a/3+e/2,\qquad
    c=a/3+e/2,\qquad e=a/3+b/2.
\]
由中间两式得 $b=c$，再比较关于 $b,e$ 的两式得 $b=e$，继而 $a=3b/2$。结合 $a+b+c+e=1$，得到
\begin{equation}
\label{eq:lihang-pr-four-basic-result}
    R=\frac19(3,2,2,2)^{\mathsf T}.
\end{equation}
迭代不必逐步单调接近各分量的极限；定理~\ref{thm:lihang-pr-finite-chain} 保证的是整个分布向量收敛。结点 $A$ 的权重最高，虽然其入度与其他结点相同，这也说明 PageRank 不等同于入度排序。
\end{solve}

基本定义的条件各有作用。不可约性使所有结点相互可达，并保证平稳分布唯一且严格为正；非周期性进一步保证从任意初态的分布收敛。若取单向三环，均匀分布仍然平稳，但从单个结点出发的分布不断循环，不能收敛；若 $M=I$，则每个分布都平稳，极限依赖初态。\emph[caution]{平稳分布存在、平稳分布唯一和从任意初态收敛是不同结论}。有限列随机矩阵总有平稳分布，一般有向图的直接链接矩阵却未必已经是列随机矩阵，悬挂结点正是需要先处理的情形。

\begin{example}
\label{ex:lihang-pr-dangling-loss}
在例~\ref{ex:lihang-pr-basic-four} 中去掉 $C\to A$，把 $C$ 对应的整列设为零，得到矩阵 $H$。说明为什么不能继续把 $H^tR_0$ 当作状态分布。
\end{example}

\begin{solve}
$H$ 的第三列和为零，其余列和为 $1$。对任意非负向量 $r$，有 $\boldsymbol1^{\mathsf T}Hr=\boldsymbol1^{\mathsf T}r-r_C$。因此，从均匀分布出发，
\[
    HR_0=\frac1{24}(3,5,5,5)^{\mathsf T},\qquad
    \boldsymbol1^{\mathsf T}HR_0=\frac34<1.
\]
丢失的质量来自没有规定下一步去向的结点 $C$。在这个具体图中，到达 $C$ 的质量在下一次乘以 $H$ 时被移出系统；从任一结点出发，在三次更新内被移出的概率至少为 $1/3$。例如，从 $B$ 出发前两步到达 $C$ 的概率为 $(1/2)(1/3)+(1/2)(1/2)=5/12$，第三步移出这些质量。于是 $\|H^{3k}R_0\|_1\leq(2/3)^k$，中间各步质量也不增加，故 $H^tR_0\to0$。这不是链的平稳概率分布，而是把游走提前终止造成的质量损失；并非所有带零列矩阵的幂都趋于零，若另有无法到达悬挂结点的封闭部分，就可能保留质量。
\end{solve}

\subsection{PageRank 的一般定义}
\label{subsec:lihang-pr-general-definition}

一般图既可能含悬挂结点，也可能不强连通或具有周期。为得到对所有初始分布都收敛的模型，先把链接规则补成完整的概率转移，再加入随机跳转。对任意有限有向图，定义
\begin{equation}
\label{eq:lihang-pr-dangling-correction}
    h_{ij}=\begin{cases}
        1/k_j,&k_j>0\text{ 且 }(v_j,v_i)\in E,\\
        0,&\text{其他情形},
    \end{cases}
    \qquad a_j=\begin{cases}1,&j\in\mathcal D,\\0,&j\notin\mathcal D,\end{cases}
    \qquad u=\frac1n\boldsymbol1,\qquad M=H+ua^{\mathsf T}.
\end{equation}
其中 $a=(a_1,\ldots,a_n)^{\mathsf T}$。非悬挂列保持原来的链接分布，悬挂列改为 $u$，表示到达无出边结点后均匀选择下一结点。这样定义的 $M$ 才是列随机矩阵；无悬挂结点时，它就是原来的链接转移矩阵。除非特别说明，下文的 $M$ 都已完成这一修正。

\begin{definition}[一般 PageRank]
\label{def:lihang-pr-general}
给定经式~\eqref{eq:lihang-pr-dangling-correction} 修正的列随机矩阵 $M$ 和\emph{阻尼因子}（damping factor）$d\in(0,1)$，定义一般转移矩阵
\begin{equation}
\label{eq:lihang-pr-google-matrix}
    A=dM+(1-d)u\boldsymbol1^{\mathsf T}.
\end{equation}
其唯一平稳分布称为一般 PageRank，记为 $R$。由于 $\boldsymbol1^{\mathsf T}R=1$，它等价地满足
\begin{equation}
\label{eq:lihang-pr-general-fixed-point}
    R=dMR+(1-d)u,\qquad R\in\Delta_n.
\end{equation}
\end{definition}

这一机制表示每一步以概率 $d$ 按 $M$ 移动，以概率 $1-d$ 重新均匀选择结点；在悬挂结点处，两种分支都采用均匀分布。$A$ 的每列和为 $1$，且 $A_{ij}\geq(1-d)/n>0$。由此所有结点一步即可相互到达，且每个结点都有正的自转移概率，定理~\ref{thm:lihang-pr-finite-chain} 适用。阻尼因子控制保留链接结构的程度；$d$ 越接近 $1$，连续沿链接游走的影响越强，$d$ 越接近 $0$，分布越接近均匀分布。教材采用 $0.85$ 作为常见取值示例，这不是对任意应用都最优的数学结论。

\begin{proposition}[一般 PageRank 的存在性、唯一性与误差界]
\label{prop:lihang-pr-contraction}
对任意列随机矩阵 $M$ 及 $0<d<1$，映射 $T(r)=dMr+(1-d)u$ 在 $\Delta_n$ 上有唯一不动点 $R$，并且
\begin{equation}
\label{eq:lihang-pr-contraction}
    \|T(p)-T(q)\|_1\leq d\|p-q\|_1,\qquad p,q\in\Delta_n.
\end{equation}
从任意 $R_0\in\Delta_n$ 迭代 $R_{t+1}=T(R_t)$，都有
\begin{equation}
\label{eq:lihang-pr-priori-bound}
    \|R_t-R\|_1\leq d^t\|R_0-R\|_1\leq2d^t,
\end{equation}
以及可计算的后验误差界
\begin{equation}
\label{eq:lihang-pr-posteriori-bound}
    \|R_t-R\|_1\leq\frac{\|R_{t+1}-R_t\|_1}{1-d},\qquad
    \|R_{t+1}-R\|_1\leq\frac{d}{1-d}\|R_{t+1}-R_t\|_1.
\end{equation}
此外，$R_i\geq(1-d)/n>0$。
\end{proposition}

\begin{proof}
列随机性保证 $T(\Delta_n)\subseteq\Delta_n$。由 $\|Mz\|_1\leq\|z\|_1$ 得到式~\eqref{eq:lihang-pr-contraction}。概率单纯形在 $1$ 范数下完备，因此压缩映射定理给出唯一不动点和迭代收敛；连续应用收缩不等式，并用两个概率向量的 $1$ 范数距离不超过 $2$，即得式~\eqref{eq:lihang-pr-priori-bound}。为得到不依赖未知 $R$ 的误差估计，注意
\[
    \|R_t-R\|_1
    \leq\|R_t-T(R_t)\|_1+\|T(R_t)-T(R)\|_1
    \leq\|R_{t+1}-R_t\|_1+d\|R_t-R\|_1.
\]
移项得到第一个后验界，再应用一次收缩得到第二个。最后由 $R_i=d(MR)_i+(1-d)/n$ 及 $(MR)_i\geq0$ 得到严格正性。
\end{proof}

将固定点方程展开，可以看到悬挂结点修正如何进入各结点的得分：
\begin{equation}
\label{eq:lihang-pr-general-components}
    R_i=d\sum_{j\in\mathcal I_i}\frac{R_j}{k_j}
        +\frac d n\sum_{j\in\mathcal D}R_j+\frac{1-d}{n},
    \qquad i=1,\ldots,n.
\end{equation}
有入边的来源必有正出度，所以第一项没有除以零的问题；第二项把全部悬挂质量均匀分给各结点。无悬挂结点时，第二项为零，即得到教材写出的通常分量公式。

\begin{notealert}[type=caution,title={阻尼平滑不能代替零列修正}]
教材先展示悬挂结点造成的零列，随后给出一般 PageRank 方程；要使后面的马尔可夫链解释与归一化同时成立，必须明确其中的 $M$ 已是列随机矩阵。若仍用含零列的 $H$，对概率向量 $r$ 有
\[
    \boldsymbol1^{\mathsf T}\bigl(dHr+(1-d)u\bigr)
    =1-d\sum_{j\in\mathcal D}r_j,
\]
只要悬挂结点有正质量，结果就不再是概率分布，不能把这一未经修正的迭代直接解释为分布的演化。把悬挂结点设为自环同样可以形成另一种完整模型，但其结果可能与均匀补列不同，因此需要明确所采用的转移规则。
\end{notealert}

\section{PageRank 的计算}
\label{sec:lihang-pr-computation}

\subsection{迭代算法}
\label{subsec:lihang-pr-iteration}

固定点方程本身给出一个确定性的分布迭代算法。输入为图的稀疏链接矩阵 $H$、悬挂指示向量 $a$、阻尼因子 $d\in(0,1)$、初始分布 $R_0\in\Delta_n$ 和误差要求 $\varepsilon>0$；输出为一般 PageRank 的近似值。具体计算不必显式形成通常稠密的 $M$ 或 $A$。
\begin{enumerate}
    \item 由各结点出度建立 $H$ 与 $a$，令 $u=\boldsymbol1/n$、$t=0$；可取 $R_0=u$，并设置最大迭代次数 $T_{\max}\geq1$。
    \item 计算当前悬挂质量 $s_t=a^{\mathsf T}R_t$，再更新
    \begin{equation}
    \label{eq:lihang-pr-sparse-update}
        R_{t+1}=dHR_t+\bigl(ds_t+1-d\bigr)u.
    \end{equation}
    令 $\delta_t=\|R_{t+1}-R_t\|_1$。
    \item 若 $d\delta_t/(1-d)\leq\varepsilon$，返回 $R_{t+1}$；由式~\eqref{eq:lihang-pr-posteriori-bound}，其 $1$ 范数误差不超过 $\varepsilon$。若未满足且尚未达到 $T_{\max}$，令 $t\gets t+1$ 并继续；达到次数上限时，返回当前近似值、后验误差界以及未达到精度要求的说明。
\end{enumerate}
上述误差界针对精确算术；浮点实现还应监测概率和与舍入误差。仅把“前后两次差很小”作为停止条件，需要同时说明采用的范数和阈值；当 $d$ 接近 $1$ 时，差值还要乘上 $d/(1-d)$ 才能控制输出误差，不能把差值本身等同于真实误差。

设边数为 $m=|E|$，采用稀疏存储时，一步矩阵向量乘法、悬挂质量求和与均匀向量更新共需 $O(m+n)$ 时间，存储量为 $O(m+n)$。由 $2d^t\leq\varepsilon$，
\[
    t\geq\left\lceil\frac{\log(2/\varepsilon)}{-\log d}\right\rceil
\]
是满足误差要求的充分迭代次数；它是适用于所有初态的上界，实际停止可能更早。$d$ 接近 $1$ 时，$-\log d$ 与 $1-d$ 同阶，这一最坏情形界说明收敛可能变慢。随机跳转项虽然对应一个稠密矩阵，却是秩一结构，不应为了实现迭代而存储 $n^2$ 个元素。

\begin{example}
\label{ex:lihang-pr-self-loop}
把例~\ref{ex:lihang-pr-basic-four} 中的 $C\to A$ 改为 $C\to C$，其余边不变，取 $d=0.8$，求一般 PageRank。这是教材用于演示迭代算法的图；$C$ 在这里具有自环，并非一个尚待补列的悬挂结点。
\end{example}

\begin{solve}
转移矩阵及平滑向量分别为
\[
    M=\begin{pmatrix}
        0&1/2&0&0\\
        1/3&0&0&1/2\\
        1/3&0&1&1/2\\
        1/3&1/2&0&0
    \end{pmatrix},\qquad (1-d)u=\frac1{20}\boldsymbol1.
\]
从 $R_0=\boldsymbol1/4$ 出发，前两步为
\[
    R_1=\frac1{60}(9,13,25,13)^{\mathsf T},\qquad
    R_2=\frac1{300}(41,53,153,53)^{\mathsf T}.
\]
为核对极限，写 $R=(a,b,c,e)^{\mathsf T}$。关于 $b,e$ 的固定点方程为
\[
    b=\frac4{15}a+\frac25e+\frac1{20},\qquad
    e=\frac4{15}a+\frac25b+\frac1{20},
\]
相减得 $b=e$。又 $a=2b/5+1/20$，代入上式可得 $b=e=19/148$、$a=15/148$；最后由归一化得到
\begin{equation}
\label{eq:lihang-pr-self-loop-result}
    R=\frac1{148}(15,19,95,19)^{\mathsf T}.
\end{equation}
$C$ 的权重超过一半，因为沿链接到达 $C$ 后会持续停留，只有随机跳转机制提供离开机会。该现象体现了封闭链接结构对权重的聚集作用；阻尼项仍保证其他结点的权重为正。
\end{solve}

\subsection{幂法}
\label{subsec:lihang-pr-power-method}

\emph{幂法}（power method）通过反复进行矩阵向量乘法，提取绝对值最大的特征值及对应特征向量的方向。对一般方阵 $B\in\mathbb R^{n\times n}$，其复特征值中绝对值最大的称为主特征值（dominant eigenvalue），相应非零向量称为主特征向量（dominant eigenvector）。具有最大绝对值的特征值可能不止一个，因此幂法收敛需要附加条件。

为说明原理，假定 $B$ 在复数域可对角化，存在一组特征向量基 $w_1,\ldots,w_n$，主特征值 $\lambda_1\neq0$，且
\[
    |\lambda_1|>|\lambda_2|\geq\cdots\geq|\lambda_n|,\qquad
    x_0=\sum_{j=1}^{n}c_jw_j,\qquad c_1\neq0.
\]
直接展开有
\begin{equation}
\label{eq:lihang-pr-power-expansion}
    B^kx_0=c_1\lambda_1^k\left[
        w_1+\sum_{j=2}^{n}\frac{c_j}{c_1}
        \left(\frac{\lambda_j}{\lambda_1}\right)^kw_j
    \right].
\end{equation}
由于各 $|\lambda_j/\lambda_1|<1$，括号中的向量趋于 $w_1$，从而得到主特征向量的方向。若 $c_1=0$，初态没有主方向分量，这一论证失效；若模最大的特征值不唯一，可能出现振荡。对一般实矩阵，负特征值还可能使规范化向量交替变号，因此需要区分方向趋近与带符号向量本身收敛。若用未规范化的 $z_k=B^kx_0$ 估计特征值，只有在 $(w_1)_j\neq0$ 的分量上，才可由式~\eqref{eq:lihang-pr-power-expansion} 得到 $(z_{k+1})_j/(z_k)_j\to\lambda_1$。

实际计算为避免向量过大或过小，可每步规范化：
\begin{equation}
\label{eq:lihang-pr-power-normalization}
    y_{t+1}=Bx_t,\qquad
    x_{t+1}=\frac{y_{t+1}}{\|y_{t+1}\|_\infty},\qquad
    \|x\|_\infty=\max_i|x_i|,
\end{equation}
前提是 $y_{t+1}\neq0$。这只是改变向量的尺度，没有把无穷范数等于 $1$ 的向量自动变成概率分布。

对 PageRank 的矩阵 $A=dM+(1-d)u\boldsymbol1^{\mathsf T}$，所有元素严格为正，且 $A$ 列随机，故 Perron--Frobenius 定理保证主特征值为 $1$、对应的正特征向量方向唯一，其余特征值的模小于 $1$。这里还可得到更具体的界：若 $Az=\lambda z$ 且 $\lambda\neq1$，则由 $\boldsymbol1^{\mathsf T}A=\boldsymbol1^{\mathsf T}$ 得 $\boldsymbol1^{\mathsf T}z=0$，于是
\[
    Az=dMz,\qquad |\lambda|\|z\|_1=d\|Mz\|_1\leq d\|z\|_1,
\]
即 $|\lambda|\leq d<1$；对复向量的 $1$ 范数，三角不等式同样适用。前面的可对角化假设只用于一般幂法的特征展开说明，PageRank 的收敛已由命题~\ref{prop:lihang-pr-contraction} 保证，不要求另外假定 $A$ 可对角化。

因此，可取任意非零非负 $x_0$，将式~\eqref{eq:lihang-pr-power-normalization} 中的 $B$ 换为 $A$ 进行迭代；每步 $y_{t+1}$ 严格为正，不会出现零分母。若采用教材的无穷范数规范化，最后必须返回
\begin{equation}
\label{eq:lihang-pr-power-probability}
    \widehat R=\frac{x_t}{\boldsymbol1^{\mathsf T}x_t}.
\end{equation}
为使停止条件直接控制概率误差，可每步先形成 $q_t=x_t/(\boldsymbol1^{\mathsf T}x_t)$，再检验 $\|Aq_t-q_t\|_1/(1-d)\leq\varepsilon$，达到次数上限仍未满足时报告未达到精度要求。计算 $Ax_t$ 同样不必形成稠密矩阵：
\begin{equation}
\label{eq:lihang-pr-power-sparse}
    Ax_t=dHx_t+
    \left[da^{\mathsf T}x_t+(1-d)\boldsymbol1^{\mathsf T}x_t\right]u.
\end{equation}
与式~\eqref{eq:lihang-pr-sparse-update} 相比，未归一的 $x_t$ 不一定满足 $\boldsymbol1^{\mathsf T}x_t=1$，所以不能把这里的最后一项直接写成 $(1-d)u$。

两种算法实际上在计算同一条概率分布序列：由于 $\boldsymbol1^{\mathsf T}A=\boldsymbol1^{\mathsf T}$，把 $x_{t+1}$ 再按分量和归一化，得到 $q_{t+1}=Aq_t=dMq_t+(1-d)u$。也就是说，概率向量迭代已经是对一般转移矩阵的幂法；无穷范数规范化只是中间尺度的另一种选择，不是另一套 PageRank 定义。

\begin{example}
\label{ex:lihang-pr-power-three}
考虑三结点图 $A\to B,C$，$B\to C$，$C\to A$。取 $d=0.85$，用幂法计算一般 PageRank，并区分无穷范数规范化结果与概率向量。
\end{example}

\begin{solve}
按 $(A,B,C)$ 排列，转移矩阵及一般转移矩阵为
\[
    M=\begin{pmatrix}0&0&1\\1/2&0&0\\1/2&1&0\end{pmatrix},\qquad
    A=\begin{pmatrix}
        0.05&0.05&0.9\\
        0.475&0.05&0.05\\
        0.475&0.9&0.05
    \end{pmatrix}.
\]
取 $x_0=(1,1,1)^{\mathsf T}$，第一步为
\[
    y_1=(1,0.575,1.425)^{\mathsf T},\qquad
    x_1=\frac1{57}(40,23,57)^{\mathsf T}
        \approx(0.701754,0.403509,1)^{\mathsf T}.
\]
继续计算，得到
\[
    x_2\approx(1,0.422406,0.781451)^{\mathsf T},\qquad
    x_3\approx(0.866018,0.598490,1)^{\mathsf T}.
\]
这些显示值只用于展示过程，继续迭代时应保留计算精度。由固定点方程和归一化可求得精确概率向量
\begin{equation}
\label{eq:lihang-pr-three-result}
    R=\frac1{1769}(686,380,703)^{\mathsf T}
      \approx(0.387790,0.214811,0.397400)^{\mathsf T}.
\end{equation}
无穷范数规范化的极限则为 $(686/703,380/703,1)^{\mathsf T}$，分量之和并不等于 $1$。两者给出相同的排序 $C>A>B$，但只有按分量和归一化的结果具有平稳概率的含义。
\end{solve}

\subsection{代数算法}
\label{subsec:lihang-pr-algebra}

把一般 PageRank 方程中的未知项移至左侧，得到
\begin{equation}
\label{eq:lihang-pr-linear-system}
    (I-dM)R=(1-d)u.
\end{equation}
其中 $I$ 为 $n$ 阶单位矩阵。列随机矩阵的诱导 $1$ 范数为 $\|M\|_1=\max_j\sum_i|m_{ij}|=1$，所以 $\|dM\|_1=d<1$，矩阵几何级数在该范数下绝对收敛。其部分和满足
\[
    (I-dM)\sum_{k=0}^{N}(dM)^k=I-(dM)^{N+1},\qquad
    \|(dM)^{N+1}\|_1\leq d^{N+1}\longrightarrow0.
\]
右乘时也得到相同恒等式，取极限证明 $I-dM$ 可逆，并给出
\begin{equation}
\label{eq:lihang-pr-neumann-series}
    R=(1-d)(I-dM)^{-1}u
      =(1-d)\sum_{k=0}^{\infty}d^kM^ku.
\end{equation}
这既证明线性方程组有唯一解，也说明它确为概率向量：各项非负，而取分量和后得到 $(1-d)\sum_{k=0}^{\infty}d^k=1$。有限维求和是连续线性运算，范数收敛保证了这里可以将它与级数求和交换。

式~\eqref{eq:lihang-pr-neumann-series} 还给出游走长度的解释。设非负整数随机变量 $K$ 满足 $P(K=k)=(1-d)d^k$，独立于初始结点和各次按 $M$ 的转移；从均匀分布出发走 $K$ 步，终点的分布就是 $R$。阻尼因子因此决定了不同长度游走的权重，而不是只给所有结点加上同一常数后就结束计算。保留 $k=0,\ldots,N$ 项时，被省略项的总质量恰为 $d^{N+1}$，这也是该截断近似的 $1$ 范数误差。

代数算法的输入是 $M,d$，输出是式~\eqref{eq:lihang-pr-linear-system} 的解。逆矩阵表达式用于理论推导，数值实现通常直接解线性方程组；对稠密 $n$ 阶矩阵，常规消元需 $O(n^3)$ 时间与 $O(n^2)$ 存储，不适于直接处理超大规模图。稀疏直接法的存储与计算还受消元填充影响，不能只按原图边数估算。无论采用哪种解法，对所得近似 $\widehat R\in\Delta_n$ 都可用残差
\begin{equation}
\label{eq:lihang-pr-linear-residual}
    b=(1-d)u-(I-dM)\widehat R,\qquad
    \|\widehat R-R\|_1\leq\frac{\|b\|_1}{1-d}
\end{equation}
检验误差，因为 $b=T(\widehat R)-\widehat R$，也可由 $\|(I-dM)^{-1}\|_1\leq\sum_{k\geq0}d^k=1/(1-d)$ 直接推出。

\section{本章概要与习题}
\label{sec:lihang-pr-summary-exercises}

本章的三个等价对象是随机游走的平稳分布、列随机矩阵 $A$ 的归一化正特征向量和线性方程组~\eqref{eq:lihang-pr-linear-system} 的解。基本模型依赖不可约与非周期条件，一般模型则通过完整的悬挂转移和随机跳转保证存在唯一正分布。迭代算法直接更新概率，幂法提取同一不动点的方向，代数算法直接求解固定点方程；三者的输出只有在总和为 $1$ 时才具有概率解释。当 $d=0$ 时，解直接为 $u$；当 $d=1$ 时，收缩保证消失，需要重新检查基本链的结构条件，不能沿用 $d<1$ 的唯一性与误差界。

\begin{example}
\label{ex:lihang-pr-exercise-powers}
教材习题 21.1：设 $S$ 为列随机矩阵，证明对任意非负整数 $k$，$S^k$ 仍为列随机矩阵。
\end{example}

\begin{proof}
$S^0=I$ 满足要求。若 $S^k$ 元素非负，则 $S^{k+1}=SS^k$ 的每个元素都是非负乘积之和，仍然非负；另一方面，
\[
    \boldsymbol1^{\mathsf T}S^{k+1}
    =(\boldsymbol1^{\mathsf T}S)S^k
    =\boldsymbol1^{\mathsf T}S^k
    =\boldsymbol1^{\mathsf T}.
\]
由归纳法结论成立。概率上，$S^k$ 的第 $j$ 列就是从结点 $j$ 出发经过 $k$ 步后的状态分布。
\end{proof}

\begin{example}
\label{ex:lihang-pr-exercise-initialization}
教材习题 21.2：在例~\ref{ex:lihang-pr-basic-four} 的图中，改用不同初始分布进行迭代，验证仍得到同一个基本 PageRank。
\end{example}

\begin{solve}
令 $e_A=(1,0,0,0)^{\mathsf T}$、$e_C=(0,0,1,0)^{\mathsf T}$，分别表示初始时确定地位于 $A$ 和 $C$。直接计算得
\[
    Me_A=(0,1/3,1/3,1/3)^{\mathsf T},\qquad
    M^2e_A=(1/2,1/6,1/6,1/6)^{\mathsf T},\qquad
    Me_C=e_A.
\]
由于 $M$ 不可约且非周期，定理~\ref{thm:lihang-pr-finite-chain} 保证两组分布均收敛到式~\eqref{eq:lihang-pr-four-basic-result}；第二组从第一步起就是第一组序列延迟一拍的结果。一般初始分布 $R_0=\sum_j(R_0)_je_j$ 是有限凸组合，故
\[
    M^tR_0=\sum_j(R_0)_jM^te_j
    \longrightarrow\sum_j(R_0)_jR=R.
\]
这证明了初态无关性，有限次迭代中看到相近数值本身并不能替代这一极限论证。
\end{solve}

\begin{example}
\label{ex:lihang-pr-exercise-general-existence}
教材习题 21.3：证明一般 PageRank 对应的马尔可夫链有平稳分布，并写出结点得分公式。
\end{example}

\begin{proof}
先按式~\eqref{eq:lihang-pr-dangling-correction} 补全悬挂列，保证 $M$ 列随机；再令 $A=dM+(1-d)u\boldsymbol1^{\mathsf T}$。其每个元素至少为 $(1-d)/n>0$，所以对应有限链不可约且非周期，由定理~\ref{thm:lihang-pr-finite-chain} 存在唯一平稳分布。也可直接应用命题~\ref{prop:lihang-pr-contraction}，得到唯一的 $R=T(R)\in\Delta_n$。逐分量展开就是式~\eqref{eq:lihang-pr-general-components}；在无悬挂结点时，它化为教材的公式
\[
    \operatorname{PR}(v_i)
    =d\sum_{j\in\mathcal I_i}\frac{\operatorname{PR}(v_j)}{k_j}
        +\frac{1-d}{n}.
\]
若没有列随机性的前提，这一证明不能成立。
\end{proof}

\begin{example}
\label{ex:lihang-pr-exercise-spectrum}
教材习题 21.4：证明随机矩阵的最大特征值为 $1$。这里将“最大”精确表述为谱半径（spectral radius）等于 $1$，且 $1$ 本身是特征值。
\end{example}

\begin{proof}
设 $S$ 为列随机矩阵。由 $S^{\mathsf T}\boldsymbol1=\boldsymbol1$，知 $1$ 是 $S^{\mathsf T}$ 的特征值；矩阵与其转置具有相同的特征多项式，所以 $1$ 也是 $S$ 的特征值。对任意复特征对 $Sz=\lambda z$，$z\neq0$，有
\[
    |\lambda|\|z\|_1=\|Sz\|_1
    \leq\sum_i\sum_js_{ij}|z_j|
    =\sum_j|z_j|=\|z\|_1.
\]
故所有特征值均满足 $|\lambda|\leq1$，结合 $1$ 是特征值，得到谱半径 $\rho(S)=\max_{\lambda}|\lambda|=1$。
\end{proof}

这一结论不保证 $1$ 是单特征值，也不排除其他模为 $1$ 的特征值。例如，单向三环的转移矩阵还具有两个非实的单位模特征值，其幂可能持续振荡。一般 PageRank 的严格正矩阵及其收缩性质排除了这些障碍，因此“随机矩阵的谱半径为 $1$”与“幂法从任意概率初态收敛到唯一方向”不能混为一谈。

教材还提及个性化 PageRank（personalized PageRank）、按主题调整的 PageRank 和其他链接分析方法。与本章直接相关的扩展是将均匀重启分布 $u$ 换成给定的概率向量 $v$，并明确悬挂列的补全分布。对任意已补成列随机矩阵的 $M$，方程 $R=dMR+(1-d)v$ 在 $0<d<1$ 时仍以同样的 $1$ 范数收缩率保证唯一概率解；若 $v$ 各分量严格为正，则该解也严格为正。允许 $v$ 含零分量时，唯一性依然成立，但不能无条件断言每个结点都获得正权重。这一扩展进一步说明，PageRank 取决于图结构和指定的游走规则，其“重要度”应在模型含义下理解。



\fi
\end{document}
